% SPDX-FileCopyrightText: 2026 Stefan Palan, Jürgen Fleiß
% SPDX-License-Identifier: LPPL-1.3c+
% SlideMage layout profile "fira": Fira Sans Light for text, Fira Mono for code
% (SlideMage v9.7.0).
%
% A layout profile holds everything that depends on the font: which font files
% load and how their weights map onto \textbf and \textit, the code fonts, the
% math handling that goes with them, and the metrics the framework measures
% against. SlideMage_setup.tex inputs SlideMage_layout_\SMLayoutProfile before it
% loads the theme, so that every length the theme and the catalog give in ex or
% em is evaluated after the fonts are set up. A template picks its profile with
%
%   \newcommand{\SMLayoutProfile}{fira}     % before % SPDX-FileCopyrightText: 2026 Stefan Palan, Jürgen Fleiß
% SPDX-License-Identifier: LPPL-1.3c+
% SlideMage setup layer and hook catalog (SlideMage v9.0.0).
%
% A template \input{}s this file first, sets whatever it wants to change from
% the catalog below, then \input{}s SlideMage_base:
%
%   % SlideMage template name: My University
%   % SPDX-FileCopyrightText: 2026 Stefan Palan, Jürgen Fleiß
% SPDX-License-Identifier: LPPL-1.3c+
% SlideMage setup layer and hook catalog (SlideMage v9.0.0).
%
% A template \input{}s this file first, sets whatever it wants to change from
% the catalog below, then \input{}s SlideMage_base:
%
%   % SlideMage template name: My University
%   % SPDX-FileCopyrightText: 2026 Stefan Palan, Jürgen Fleiß
% SPDX-License-Identifier: LPPL-1.3c+
% SlideMage setup layer and hook catalog (SlideMage v9.0.0).
%
% A template \input{}s this file first, sets whatever it wants to change from
% the catalog below, then \input{}s SlideMage_base:
%
%   % SlideMage template name: My University
%   \input{SlideMage_setup}
%   \definecolor{c_primary_highlight}{HTML}{...}
%   \renewcommand{\SMInstituteLogoPath}{schemes/my_logo.pdf}
%   \setlength{\SMLogoHeight}{3ex}
%   \input{SlideMage_base}
%   ... additions only: dark-mode beamer colors, citation format, project macros
%
% Rules for a template:
%   1. Before this file: only \SMDarkMode (written into the deck by SlideMage).
%   2. Between this file and SlideMage_base: \setlength, \renewcommand,
%      \definecolor and \colorlet on the names in this catalog. Every name
%      exists already, so a misspelt \renewcommand is an error instead of a
%      setting that silently does nothing.
%   3. After SlideMage_base: additions only.
%
% Decks never load the beamer theme themselves; this file does. A length given
% in ex or em is evaluated where the template sets it, in the normal body font,
% after the theme has loaded its fonts.

\newcommand{\SMSetupLoaded}{}

% The beamer theme, with the body fonts.
\usetheme{gotham}

% ---------------------------------------------------------------------------
% Palette
% ---------------------------------------------------------------------------
% The SlideMage look; a template redefines what it changes (usually all of
% them). The framework combines these pairs, so a template should keep them
% readable (WCAG 2.1): body text c_primary_dark on each c_secondary_* block
% title color >= 4.5:1, title/footer text on its box color >= 4.5:1, white text
% on c_primary_gray (highlightframe), link underline c_secondary_blue on the
% canvas >= 3:1.
%
% SlideMage's own values: wherever possible a color the app itself uses
% (src/slidemage/palette.py, qt_deck_table.py, assets/splash.png). White
% title/footer text on c_primary_dark is 16.7:1, white on c_primary_gray 5.0:1.
\definecolor{c_primary_dark}{HTML}{1e1e1e}      % body text, dark boxes; splash screen background (palette dark "bg")
\definecolor{c_primary_highlight}{HTML}{0066cc} % accent: bullets, highlights; SlideMage button blue (Create PDF, progress bar, links)
\definecolor{c_primary_gray}{HTML}{6f6f6f}      % highlightframe background; palette dark "disabled_fg", dark enough for white text
\definecolor{c_primary_white}{HTML}{FFFFFF}
\definecolor{c_secondary_violet}{HTML}{ba68c8}  % block titles (block); palette dark "tab_deck_editor"
\definecolor{c_tertiary_violet}{HTML}{dcb4e4}   % block bodies; derived (50% white-blend)
% Derived, not an app color: the button blue is too dark for the dark text
% an infoblock title puts on it, the app's dark-mode link blue (4da6ff) too
% light for a link underline on white; this one sits between them and passes
% both (4.9:1 and 3.4:1).
\definecolor{c_secondary_blue}{HTML}{3d8fe0}    % infoblock titles, link underline
\definecolor{c_tertiary_blue}{HTML}{9ec7f0}     % derived (50% white-blend)
\definecolor{c_secondary_green}{HTML}{4caf50}   % exampleblock titles; palette dark "status_pass"
\definecolor{c_tertiary_green}{HTML}{81c784}    % palette dark "tab_slide_editor"
\definecolor{c_secondary_red}{HTML}{f04848}     % alertblock titles; palette dark "status_fail"
\definecolor{c_tertiary_red}{HTML}{f8a4a4}      % derived (50% white-blend)
\definecolor{c_secondary_orange}{HTML}{ffb74d}  % reference block titles; palette dark "tab_archived"
\definecolor{c_tertiary_orange}{HTML}{ffdba6}   % derived (50% white-blend)
\definecolor{c_secondary_yellow}{HTML}{ffd500}  % logo/splash yellow
\definecolor{c_tertiary_yellow}{HTML}{ffe066}   % deck editor's highlight yellow
\definecolor{c_tertiary_gray}{HTML}{e4e4e4}     % block bodies, process boxes; palette light "panel_bg"
% Supporting colors ("custom": not part of the palette proper).
% c_custom_gray is hyperlink text, button fill and the quote attribution.
\definecolor{c_custom_gray}{HTML}{2d2d2d}       % palette dark "panel_bg"
\definecolor{c_custom_red}{HTML}{cc0000}        % palette light "status_fail"
\definecolor{c_custom_green}{HTML}{1a8a1a}      % palette light "status_pass"
\definecolor{c_custom_brown}{HTML}{996600}      % palette light "syntax_brace"

% ---------------------------------------------------------------------------
% Roles: which palette color a part of the slide uses (color names, read
% when used, so they follow a redefined palette color)
% ---------------------------------------------------------------------------
\newcommand{\SMBulletColor}{c_primary_highlight}       % itemize bullets
\newcommand{\SMTitleBoxColor}{c_primary_dark}          % frame title bar
\newcommand{\SMTitleTextColor}{c_primary_white}        % frame title text
\newcommand{\SMFooterBoxColorDefault}{c_primary_dark}  % footer box
\newcommand{\SMFooterTextColor}{c_primary_white}       % footer text, slide id

% Optional colors a template may define; SlideMage_base derives each one from
% the palette when the template leaves it undefined:
%   SMCanvasColor          slide background                  (white)
%   SMCardBackgroundColor  SMCard fill                       (white)
%   SMCardBorderColor      SMCard border                     (c_primary_dark)
%   SMProcessBorderColor   SMProcess border, dark mode       (c_tertiary_gray)
%   SMCodeKeywordColor     \SMCodeFile commands              (c_primary_highlight)
%   SMCodeCommentColor     \SMCodeFile comments              (c_primary_gray)
%   SMCodeBackgroundColor  \SMCodeFile panel                 (c_tertiary_gray)
%   SMCodeTextColor        \SMCodeFile text                  (c_primary_dark)

% ---------------------------------------------------------------------------
% Logo
% ---------------------------------------------------------------------------
% A vector PDF cropped to its ink, found on the TeX search path (a project
% path such as schemes/<name>/logo.pdf, or a file next to the template).
\newcommand{\SMInstituteLogoPath}{slidemage-logo.pdf}
\newlength{\SMLogoHeight}
\setlength{\SMLogoHeight}{3ex}                         % logo height in the title bar

% ---------------------------------------------------------------------------
% Frame title bar
% ---------------------------------------------------------------------------
% Space around the frame title (the theme's registers under their SlideMage
% names). The right padding is measured against the compiled PDF for the logo:
% see SlideMage_base.tex's comment on \gothamFrametitleRightpading.
\let\SMFrametitleTopPadding\gothamFrametitleToppading
\let\SMFrametitleBottomPadding\gothamFrametitleBottompading
\let\SMFrametitleLeftPadding\gothamFrametitleLeftpading
\let\SMFrametitleRightPadding\gothamFrametitleRightpading
\setlength{\SMFrametitleRightPadding}{0.28cm}

%   \definecolor{c_primary_highlight}{HTML}{...}
%   \renewcommand{\SMInstituteLogoPath}{schemes/my_logo.pdf}
%   \setlength{\SMLogoHeight}{3ex}
%   % SPDX-FileCopyrightText: 2026 Stefan Palan, Jürgen Fleiß
% SPDX-License-Identifier: LPPL-1.3c+
% A template \input{}s SlideMage_setup (the theme and every template setting
% with its default) before this file; a template from before SlideMage v9.0.0
% does not, and SlideMage's Project compatibility update fixes it.
\ifdefined\SMSetupLoaded\else
	\PackageError{SlideMage}{The template does not \string\input{SlideMage_setup}}{%
		Add \string\input{SlideMage_setup} at the top of the template, before its colors,
		or run Info -> Project compatibility in SlideMage.}
\fi
% \SMLuaOnce{<lua>}: run Lua now AND remember it, so a precompiled format
% (SlideMage's fast preview, see preview_format.py) can replay it. A LuaTeX
% format stores TeX state but not Lua state: every global function, table and
% callback a preamble sets up is gone when the format is loaded. The replay list
% is run by the format's \everyjob; in an ordinary build it is never used.
\newtoks\SMLuaReplayList
\newcommand{\SMLuaOnce}[1]{%
	\directlua{#1}%
	\global\SMLuaReplayList\expandafter{\the\SMLuaReplayList\directlua{#1}}%
}
% The packages this file uses, loaded here explicitly (the theme loads most of them already, which makes
% each \usepackage below a no-op): iftex for \ifluatex, fontspec for \newfontfamily, xcolor for the colors.
\usepackage{iftex}
\usepackage{fontspec}
\usepackage{xcolor}
\usepackage{etoolbox} % \AtEndEnvironment, used for the itemize level-drop spacing below
\usepackage{tikz} % For painting a white box over title bar
\usetikzlibrary{calc}
\usetikzlibrary{arrows.meta} % for e.g. arr/.style={-{Stealth[...]}...} on diagram slides
\usetikzlibrary{positioning} % for `right=<dist> of <node>`-style relative node placement
\usetikzlibrary{decorations.text} % text-along-a-path, for a circular photo's curved caption
\usetikzlibrary{decorations.pathreplacing} % the `brace` decoration shape, for \SMUnderbrace below
% TikZ reveal steps that keep the picture's size (2026-09-30): `visible on=<3->`
% draws an element on steps 3+ and draws it invisibly (opacity 0) before, so it
% still counts toward the bounding box -- unlike \only<3->{\draw ...;}, whose
% element is simply missing on steps 1-2 and can make the whole picture jump
% when it appears (the "diagram shifts downward in step 7" class of fix).
% Same key names as the aobs-tikz library (overlay-beamer-styles), defined here
% so no extra package is needed.
\tikzset{
	invisible/.style={opacity=0, text opacity=0},
	alt/.code args={<#1>#2#3}{\alt<#1>{\pgfkeysalso{#2}}{\pgfkeysalso{#3}}},
	visible on/.style={alt={#1{}{invisible}}},
}
\usepackage{booktabs} % \toprule/\midrule/\bottomrule for SMTable
\usepackage{colortbl} % \arrayrulecolor{}, used by SMTable to color those rules
\usepackage{pifont} % \ding{51}/\ding{55} checkmark/X-mark glyphs, e.g. slides/fr3-Forms-of-Capital-Market-Efficiency-Strong-Semistrong-Weak-current.tex's own validity ratings
\usepackage{pgfplots} % real-data line/bar charts (e.g. slides/3vp-Metallgesellschaft-Oil-Price-Chart-current.tex) -- confirmed safe alongside this deck's own \tableofcontents, see CONVERSION_NOTES.md's "beamer/TeX gotchas" entry
\pgfplotsset{compat=1.18}

% The math fonts (Computer Modern next to the sans text font) and the \text{} wrapper that
% goes with them are font handling: SlideMage_layout_\SMLayoutProfile.tex, input by SlideMage_setup.

% Colors, institute logo path, and the two logo-calibration lengths
% (\SMLogoHeight, \SMFrametitleRightPadding = \SMFrametitleRightPadding)
% are template settings: SlideMage_setup.tex defines each with its default
% (the hook catalog), the template changes what it wants, and this file only
% reads them -- shared machinery, no institution-specific branding. See e.g.
% SlideMage_template_uni-graz.tex for the Graz values. \SMLogoWidth is the one
% logo-derived length that *does* stay computed here (via \settowidth,
% further below) rather than hardcoded per institution, since it already
% adapts automatically to whatever image \SMInstituteLogoPath points at.

% Style hyperlinks (\href{url}{text}) consistently across all slides:
% underlined, so a link is visually recognizable as such against plain
% body text -- yellow (this project's usual accent color) was tried first
% for the *text* itself, but yellow-on-white body text is hard to read,
% unlike yellow used as an accent against dark/black text elsewhere in the
% theme. Went through a few more candidates on 2026-07-06 (all direct
% instruction) before settling on c_secondary_blue for both text
% and underline (a real color reads more clearly as "this is a link" than
% any gray does, without yellow's legibility problem).
%
% 2026-07-07, direct instruction: split text color and underline color
% instead of using one color for both -- link text back to
% c_custom_gray (the near-black dark gray also used for
% quoteframe's attribution text), underline in a separate accent color.
% Plain \underline can't do this: its rule is drawn in whatever color is
% current at that point, so a single \underline{...} call always
% underlines in its own text's color, never a second, different one.
% \SMHrefUnderline below reimplements it instead: measure the linked
% text's own width into a savebox, print that box in customgray, then
% back up by the box's own width and draw a plain \rule of the same
% width, raised to sit just below the text's baseline, in the accent
% color. Same single-line limitation as plain \underline (and the reason
% \url{} is used instead of \href on slides/acw-Some-Topics-in-Testing-the-EMH-current.tex --
% see that slide's own comment): a \rule, like \underline, can't span a
% line break either, so this only ever wraps normal, comparatively short
% \href display text, never something that needs mid-word breaking.
%
% First tried c_primary_highlight for the underline (same day) -- confirmed
% via the compiled PDF that the rule *was* being drawn (right color,
% position, width), so not a compile-time bug, but yellow (very high
% luminance, close to white) at this rule's 0.4pt thickness reads as
% essentially invisible against the white slide background, worst at
% footnote size (see e.g. slides/za2-Startup-Financing-Categorization-current.tex's
% own long inline URL) where it's easy to miss entirely. Switched
% 2026-07-07 (direct instruction) to c_secondary_blue instead --
% far more contrast against white at the same thickness -- and the rule
% was also thickened at the same time (0.4pt -> 0.6pt), on the reasoning
% that it might need it for the same contrast reason. It didn't: blue
% already reads clearly at the original thickness, and 0.6pt instead
% just looked heavier than intended (direct feedback 2026-07-08, e.g. on
% slides/6qi-Yield-Spreads-Corporate-Bonds-current.tex). Reverted to
% 0.4pt -- only the color was ever actually asked to change.
%
% `citecolor` (2026-07-23, direct instruction: "make sure the year in the
% short citation is printed in black", found on
% slides/64m-APT-Choosing-the-Factors-Chen-Roll-Ross-current.tex's own
% `\textcite{}` call, "and generally" -- a deck-wide fix, not a one-off):
% with `colorlinks=true` and no explicit `citecolor`, hyperref falls back
% to its own factory default (green) for every citation hyperlink --
% confirmed directly, rendering 64m with the real biblatex/biber setup
% (this deck's own `\hypersetup` alone, with no citecolor, has no other
% source that could be coloring it): `\textcite{chen1986}` printed as
% "Chen, Roll, Ross (\textcolor{green}{1986})", the year in green, nothing
% else on the line affected (`\textcite` only wraps the year/label itself
% in the actual hyperlink, not the surrounding author names, which are
% plain text). `citecolor=black` set explicitly here, alongside
% urlcolor/linkcolor, rather than left to hyperref's own default, so this
% applies deck-wide to every citation command (`\cite`, `\textcite`,
% `\parencite`, ...) that ends up as a colored link, not just 64m's own
% call site.
\hypersetup{colorlinks=true, urlcolor=c_custom_gray, linkcolor=c_custom_gray, citecolor=black}


\let\SMOrigHref\href
\newsavebox{\SMHrefTextBox}
%
% 2026-07-08 (direct instruction): the underline's own rule thickness now
% scales with the ambient font size instead of a single fixed 0.2pt for
% every link everywhere -- 0.2pt at this deck's footnote size, 0.4pt at its
% normal body-text size, interpolated (and extrapolated, for any other
% size a link happens to appear at) linearly in between.
%
% Deriving "current font size" from \f@size (the textbook NFSS way) was
% tried first and confirmed broken specifically *inside* a beamer frame --
% \typeout{\meaning\f@size} right after this preamble loads shows a normal,
% defined macro (`macro:->10.95`), but the exact same check inside
% `\begin{standardframe}...\end{standardframe}` reports `\f@size` as
% *undefined*, presumably from beamer's own overlay-aware font-switching
% machinery not keeping classic NFSS's own current-size bookkeeping in
% sync the way plain (non-beamer) LaTeX does. Since every real use of this
% macro happens inside a frame, that path was a dead end regardless of how
% well it worked in isolation.
%
% Measuring a reference glyph's own rendered height instead sidesteps this
% entirely: building a box and reading back its own height just measures
% whatever font is *actually* active at the call site, independent of
% whether beamer's own internal "current size" bookkeeping happens to
% agree -- and within one font
% family, every glyph's metrics scale by exactly the same ratio the point
% size does, so a measured height is just as valid a stand-in for "current
% size" as \f@size would have been, had it worked. "M" (present in every
% context this is used in) is the reference glyph; \SMHrefUnderlineFnHeight/
% \SMHrefUnderlineBodyHeight are its height at this deck's footnote font
% and normal body font respectively, each measured once, here, in the
% relevant font context.
% Deliberately built from raw TeX primitives (\newdimen/\setbox/\ht),
% *not* \newlength/\setlength/\settoheight -- found, while building this,
% that loading this deck's theme silently turns \setlength and
% \settoheight into something else entirely: `\meaning\setlength` reads
% `macro:->\protect \setlength` right after `\usetheme{gotham}` (still the
% ordinary kernel definition beforehand), and the same for `\settoheight`.
% Not chased to which file does this or why (bisected only as far as
% "plain beamer alone doesn't, beamer+gotham does" -- most likely
% `translator`, which beamer loads for localized strings and is known to
% robustify commands this way, but not confirmed) -- what matters
% practically is that a *global* assignment through either one silently
% no-ops instead of erroring, confirmed directly and in isolation (a
% minimal beamer+gotham file, no other content): `\global\setlength{\len}
% {\otherlen}` leaves `\len` at 0pt after the group closes, while the
% exact same pattern with plain `\newdimen`/raw `=` assignment propagates
% correctly. Every plain (non-global) `\setlength` elsewhere in this
% project's own files is unaffected -- it's specifically the *global*
% form that breaks, which is why this is the first place it's shown up
% despite `\setlength` being used constantly throughout this preamble.
\newdimen\SMHrefUnderlineFnHeight
\newdimen\SMHrefUnderlineBodyHeight
\newdimen\SMHrefUnderlineNowHeight
\newdimen\SMHrefUnderlineThickness
% A dedicated box register, not the commonly-borrowed box0 -- several
% slides in this deck (the "Long-term bonds" tables) already use box0
% themselves (`\sbox0{...}`/`\usebox0`) for their own, unrelated
% mid-table-layout purposes, and a link can appear anywhere, including
% inside one of those tables -- reusing box0 here would silently clobber
% whatever they were holding there.
\newbox\SMHrefUnderlineBox
% The two reference heights are (re)measured on every single \href call,
% not once at preamble-load time -- tried once-at-load first and found a
% second, independent problem with it: `\usebeamerfont{footnote}` right
% here (before `\begin{document}`) silently has no effect at all (footnote-
% size "M" and normal-size "M" measured identically, confirmed directly),
% presumably because beamer's font-theme machinery isn't fully installed
% until the document starts. Since every real `\href` call happens well
% inside the document (inside some frame), measuring fresh each time
% sidesteps that timing question entirely rather than needing to find the
% right hook to measure once, early enough to be free but late enough to
% be correct -- the extra cost is two throwaway one-glyph boxes per link,
% not something worth optimizing away in a slide deck.
\newcommand{\SMHrefUnderlineMeasureRefHeights}{%
	{\usebeamerfont{footnote}\setbox\SMHrefUnderlineBox=\hbox{M}\global\SMHrefUnderlineFnHeight=\ht\SMHrefUnderlineBox}%
	{\normalsize\setbox\SMHrefUnderlineBox=\hbox{M}\global\SMHrefUnderlineBodyHeight=\ht\SMHrefUnderlineBox}%
}
% A plain e-TeX \dimexpr can multiply a dimension by another dimension and
% then divide by a third, as one combined operation, as long as the
% *middle* operand of the chain is itself wrapped in its own nested
% \dimexpr...\relax -- confirmed directly (an expl3 \dim_set:Nn attempt
% first, using \seq_.../\fp_... the way \setfootnotesymbols above does,
% kept hitting "Missing ) inserted for expression"/type-mismatch errors:
% dividing one dimension by another is a *ratio*, not a dimension, and
% \dim_set:Nn's own expression grammar doesn't implicitly convert between
% the two the way this nested-\dimexpr form does). No expl3 needed here at
% all -- plain e-TeX arithmetic, already available in any current LaTeX.
\newcommand{\SMHrefUnderlineSetThickness}{%
	\SMHrefUnderlineThickness=\dimexpr 0.2pt +
		(\SMHrefUnderlineNowHeight - \SMHrefUnderlineFnHeight)
		* \dimexpr0.2pt\relax
		/ (\SMHrefUnderlineBodyHeight - \SMHrefUnderlineFnHeight)\relax
}
\newcommand{\SMHrefUnderline}[1]{%
	\sbox{\SMHrefTextBox}{#1}%
	\textcolor{c_custom_gray}{\usebox{\SMHrefTextBox}}%
	\hspace{-\wd\SMHrefTextBox}%
	\SMHrefUnderlineMeasureRefHeights
	\setbox\SMHrefUnderlineBox=\hbox{M}\SMHrefUnderlineNowHeight=\ht\SMHrefUnderlineBox
	\SMHrefUnderlineSetThickness
	\raisebox{-0.35ex}[0pt][0pt]{\textcolor{c_secondary_blue}{\rule{\wd\SMHrefTextBox}{\SMHrefUnderlineThickness}}}%
}
\renewcommand{\href}[2]{\SMOrigHref{#1}{\SMHrefUnderline{#2}}}
% \url{} (the `url` package hyperref depends on/patches) defaults to a
% monospace-ish "tt" style regardless of the above -- used deliberately in
% only one place in this deck,
% slides/acw-Some-Topics-in-Testing-the-EMH-current.tex, specifically
% *for* its line-breaking-at-slashes behavior (a plain \href with the
% URL as its own display text has no break points at all -- see that
% slide's own comment). Every other same-text-as-target URL in the deck
% (including slides/bb6-A-Perpetual-Bond-current.tex's own video link,
% which used to be a \url{} too) uses plain \href{url}{url} instead, since
% \url{} doesn't pick up this deck's custom colored underline below (its
% breakable content can't be measured into the savebox the underline
% needs -- see that mechanism's own comment) and most of these URLs are
% short enough to not need the breaking \url{} exists for anyway. 2026-07-07,
% direct instruction: that one genuinely-needed \url{} (on s-acw) should
% still look like this deck's other hyperlinks, i.e. the surrounding body
% font, not \url{}'s own default monospace.
% `\urlstyle{same}` (a standard url-package switch) does exactly that,
% with no effect on \url{}'s own breaking behavior. Set deck-wide, not
% only on that one slide, so any future \url{} use matches by default
% too, rather than needing to remember this each time. Coloring already
% matches automatically (hyperref's own \url respects
% \hypersetup{urlcolor=...} the same way \href does); the underline
% \href gets above doesn't carry over, since \underline's own argument
% can't contain line breaks -- wrapping \url{}'s necessarily-breakable
% text in \underline would silently undo the breaking this exists for,
% so that one difference is accepted rather than "fixed."
\urlstyle{same}



% The colour roles of the theme follow the template's palette: SMColorDark and SMColorA, the canvas
% default and the roles (beamerthemeSlideMage.sty). Needs to go after the template's color redefinitions.
\SMApplyColorRoles
% The hooks that go into registers of beamer's and LaTeX's own (text margins, footnote font
% and spacing): needs to go after the template as well.
\SMApplyLayoutHooks

% Derived from the template's hooks and the font (v9.7.0; they were hand-kept numbers of the catalog
% and the layout profile before). The setup layer holds a macro under each of these four names that
% stops with an error when a template or a profile sets one; a name still holding that macro is
% replaced by the length it stands for, and a name that was redefined is reported.
\makeatletter
\def\SM@derivedlength#1{%
	\expandafter\ifx\csname #1\expandafter\endcsname\csname SM@guard@#1\endcsname\else
		\PackageError{SlideMage}{\expandafter\string\csname #1\endcsname\space is derived, not set}{%
			A template or a layout profile redefined it. SlideMage_base derives it from the hooks and
			the font; remove that line from the template or the profile.}%
	\fi
	\expandafter\let\csname #1\endcsname\relax
	\expandafter\newlength\csname #1\endcsname}
\SM@derivedlength{SMContentTopAlign}
\SM@derivedlength{SMTitleBarBottomAlign}
\SM@derivedlength{SMFootnoteRightShift}
\SM@derivedlength{SMFooterRightShift}
% v9.8.0: the lengths below were numbers of the layout profile (measured once at Fira) or of the
% catalog; each is evaluated in the font where it is used. The footer's, the lists' and the section
% frame's follow with the code that uses them.
\SM@derivedlength{SMFooterTextDrop}
\SM@derivedlength{SMFootlineVOffset}
\SM@derivedlength{SMLabelItemWidth}
\SM@derivedlength{SMLabelItemSep}
\SM@derivedlength{SMEnumerateNestedMinNumber}
\SM@derivedlength{SMEnumerateNestedLabelSep}
\SM@derivedlength{SMEnumerateNestedLeftMargin}
\SM@derivedlength{SMEnumerateLabelBox}
\SM@derivedlength{SMSectionframeRefDepth}
\SM@derivedlength{SMSectionframeBottomGap}
\SM@derivedlength{SMSectionframeBottomPadCorrection}
\SM@derivedlength{SMSectionframeHeightCorrection}
% v9.9.0: the pull-ups of the frames the automatic Top Margin does not correct by their content were
% the last three numbers of the layout profiles. \SMContentTopShift is measured (the frame's origin,
% see AUTOMATIC TOP MARGIN; in a run without \jobname.top, from the first frame that counts on), the two columns
% corrections are measured where a columns environment starts such a frame.
\SM@derivedlength{SMContentTopShift}
\SM@derivedlength{SMColumnsTopCorrection}
\SM@derivedlength{SMColumnsTTopCompensation}
\makeatother
% - \SMContentTopAlign, the Top Margin as an overlay coordinate: the same line the automatic Top Margin
%   puts the first text line's box on (the Top Margin is in bp, the coordinate in pt: 41.58bp is 41.7359pt).
\setlength{\SMContentTopAlign}{\SMTopMargin}
% - \SMTitleBarBottomAlign, the visible bottom of the title bar: where the logo ends. The frame title
%   template centres the logo on the capital height of the title font (\raisebox{.5\fontcharht\font`E-0.5\height}
%   at the baseline, which is the top padding plus the height of an X below the bar's top), so the
%   bottom is the top padding, plus that X, minus half an E, plus half the logo.
\begingroup\usebeamerfont{frametitle}\settototalheight{\dimen0}{X}\dimen1=\fontcharht\font`E
\global\SMTitleBarBottomAlign=\dimexpr\SMFrametitleTopPadding+\dimen0-0.5\dimen1+0.5\SMLogoHeight\relax\endgroup
% - \SMFootnoteRightShift: a footnote ends where the text margin does; the logo ends its right padding
%   from the page's edge, so the footnotes are pushed on by the difference.
\setlength{\SMFootnoteRightShift}{\dimexpr\SMTextMarginRight-\SMFrametitleRightPadding\relax}
% (\SMFooterRightShift is derived with the footer, below.)
% - \SMFootlineVOffset (v9.8.0; the catalog's -0.5ex-0.77pt until then), the space the footline
%   template puts before the footline box. Beamer sets the frame text area to end where the footline
%   box begins (its height plus depth, and the offset, are its \footheight) plus a fixed 4pt
%   (beamerbaseframecomponents.sty, \beamer@calculateheadfoot), and the last footnote line's baseline
%   is the bottom of that area. So the baseline lies
%       \SMFootlineHeight + \SMFootlineDepth + offset + 4pt
%   above the page's bottom edge, and the offset that puts the descenders of that line (the depth of
%   "gjpqy" in the footnote font) \SMFootnoteFooterGap above the footer box follows.
\makeatletter
\begingroup\usebeamerfont{footnote}\settodepth{\dimen0}{gjpqy}
\global\SMFootlineVOffset=\dimexpr\SMFooterBoxHeight+\SMFootnoteFooterGap+\dimen0-4pt-\SMFootlineHeight-\SMFootlineDepth\relax
\xdef\SMPDFInfoFootnoteDescender{\strip@pt\dimexpr\dimen0*7200/7227\relax}\endgroup
\makeatother

% The effective margins, for tools/check_layout.py (SlideMage v9.5.0): custom keys of the PDF's
% info dictionary, in bp (PDF units, what the tool measures in). SMTopMargin is the Top Margin
% (\SMTopMargin), SMLeftMargin the frame title text's left edge (\SMFrametitleLeftPadding). A template
% has set its hooks by now; a slide cannot change them. SMTitleBarBottom (v9.6.0) is the bottom edge of
% the title bar (\SMTitleBarBottomAlign): the tool tells a frame with the ordinary title bar by it.
% SMFontAscender, the text font's ascender, follows once it is known (below, with the Lua file).
% v9.8.0: SMFootnoteFooterGap, SMFootnoteDescender (the depth of "gjpqy" in the footnote font) and
% SMFooterBoxHeight, which tools/calibrate_font_profile.py checks the footnotes' position against.
\ExplSyntaxOn
\edef\SMPDFInfoTopMargin{\dim_to_decimal_in_bp:n{\SMTopMargin}}
\edef\SMPDFInfoLeftMargin{\dim_to_decimal_in_bp:n{\SMFrametitleLeftPadding}}
\edef\SMPDFInfoTitleBarBottom{\dim_to_decimal_in_bp:n{\SMTitleBarBottomAlign}}
\edef\SMPDFInfoFootnoteGap{\dim_to_decimal_in_bp:n{\SMFootnoteFooterGap}}
\edef\SMPDFInfoFooterBoxHeight{\dim_to_decimal_in_bp:n{\SMFooterBoxHeight}}
\ExplSyntaxOff
\hypersetup{pdfinfo={SMTopMargin={\SMPDFInfoTopMargin}, SMLeftMargin={\SMPDFInfoLeftMargin},
	SMTitleBarBottom={\SMPDFInfoTitleBarBottom},
	SMFootnoteFooterGap={\SMPDFInfoFootnoteGap}, SMFootnoteDescender={\SMPDFInfoFootnoteDescender},
	SMFooterBoxHeight={\SMPDFInfoFooterBoxHeight}}}

% Shared diagram roles: templates set the colors once; slides use these styles.
\makeatletter
\@ifundefinedcolor{SMCardBackgroundColor}{\colorlet{SMCardBackgroundColor}{white}}{}
\@ifundefinedcolor{SMCardBorderColor}{\colorlet{SMCardBorderColor}{c_primary_dark}}{}
\@ifundefinedcolor{SMProcessBorderColor}{\colorlet{SMProcessBorderColor}{c_tertiary_gray}}{}
\makeatother
\tikzset{
	SMCard/.style={draw=SMCardBorderColor, fill=SMCardBackgroundColor, text=c_primary_dark},
	SMProcess/.style={draw=none, fill=c_tertiary_gray, text=c_primary_dark},
}

% The theme's divider frames for \part, \section, \subsection and \subsubsection
% (\SMDividerFrames in beamerthemeSlideMage.sty) are off, so \section{} (used
% for a deck's chapter grouping/table of contents) does not insert its own
% plain, unstyled divider frame plus a small "current section highlighted in
% the outline" frame -- redundant with (and visually inconsistent next to)
% this project's own hand-styled sectionframe (see README.md's "Special frame
% types" section), which every chapter in this deck already opens with as an
% explicit \slide{}. \section{} only feeds \tableofcontents (see
% slides/nkj-Agenda-current.tex).
%
% Same for \subsection{} (2026-07-15, direct instruction): \subsection{} calls
% still register in the PDF's own bookmark/outline structure (Acrobat's
% "chapters" pane -- a plain PDF sectioning feature, independent of whether
% beamer auto-inserts a frame for it) and still feed 	ableofcontents (see
% slides/nkj-Agenda-current.tex's own `	ableofcontents[hideallsubsections]`,
% which then hides them from that one slide specifically), without also
% spawning their own unstyled divider slide in the middle of a chapter every
% time one is added.


% A color leak from \titleframe/\quoteframe into whichever frame follows
% them is fixed in standardframe's own opening code below (search
% "Forces this frame's own body") rather than here -- an
% \AtBeginEnvironment{frame}{\color{c_primary_dark}} hook was tried first
% and confirmed, by isolated reproduction, to have *no effect at all*:
% beamer's own "frame" isn't a plain \newenvironment etoolbox's
% \AtBeginEnvironment can patch that way, so the hook silently never fired.


%\colorlet{colorPale}{gPaleYell} % BG in light/normal mode
%\colorlet{colorDark}{gDarkBlack} % FG in light/normal mode
%\colorlet{colorA}{gAnthracite} % frametitle, standin.out,
%\colorlet{colorAreversed}{gLightTeal} % frametitle, standin.in,
%\colorlet{colorB}{gMidGrey} % gray BG : progress bar, blocks
%\colorlet{colorC}{gDeepYellOr} % progress bar
%\colorlet{colorD}{gLightOrange} % alert
%\colorlet{colorE}{gLightGreen} % example

% Opt-in yellow border for an image placed via a tikz node (e.g.
% \node[SMImageBorder] {\includegraphics{...}};) -- matches the
% 3pt/accent1-yellow outline PowerPoint draws directly on a picture shape's
% own border (<a:ln w="38100"><a:solidFill><a:schemeClr val="accent1"/>,
% where accent1 turned out to already be this project's own
% c_primary_highlight, not a new color to map). A tikz *style* (a named
% bundle of node keys), not a wrapper command: adding or removing it is
% just adding or deleting one word from that one \node's own option list,
% no separate macro call/argument needed -- add `SMImageBorder` to draw
% the border, leave it off for a plain image, same as any other tikz key.
%
% Thickness is a single named length, \SMImageBorderWidth, not a bare
% "3pt" baked directly into the style (2026-07-15, direct instruction: "make
% sure the yellow image borders have consistent thickness throughout the
% project by defining a single variable") -- confirmed directly, by
% measuring the actual PDF content-stream "w" operator, that every
% SMImageBorder use already rendered at the exact same 3pt regardless
% (a genuinely shared style was already in effect, no numeric drift found),
% so this is a refactor for future adjustability/discoverability, not a fix
% for an actual inconsistency -- any perceived thickness difference between
% slides traces to something else instead (e.g. a border straddling flush
% image content on one slide and not another, which reads as "heavier" even
% at an identical stroke width -- see the SMImageBorder padding gotcha
% below).
% (\SMImageBorderWidth, 2.5pt, is a hook of SlideMage_setup.tex.)
% outer sep = half the stroke (2026-10-01, direct instruction: "if a figure
% has a border, that border should determine its position"): a bordered
% node's anchors (north east, west, ...) sit on the border's *outer* edge, so
% anchoring it at a margin puts the border on the margin. That is TikZ's own
% default (0.5\pgflinewidth, evaluated with the node's final line width, so a
% slide's line width=1.2pt override still gets 0.6pt); the style sets it
% explicitly so that an "outer sep=0pt" earlier in the option list
% (\cornerfigure's node, for an unbordered image) doesn't take it away. A
% slide that writes outer sep=0pt *after* the style opts out again -- ts8 did,
% and its border stuck out 1.25pt past the Right Margin.
\tikzset{SMImageBorder/.style={draw=c_primary_highlight, line width=\SMImageBorderWidth, outer sep=0.5\pgflinewidth}}

% Named styles for the drawing conventions slides repeat by hand
% (2026-09-30): a restyle -- "make the arrows a bit bolder", "all
% highlights dashed" -- is then one edit here instead of one per slide.
% Each is exactly the option list the slides used, so switching a slide to
% the name changes nothing on the page.
\tikzset{
	SMHighlight/.style={c_custom_red, very thick, rounded corners},   % red box around part of a figure/table (18x)
	SMDottedHighlight/.style={densely dotted, thick, c_custom_red},  % red dotted guide line (19x)
	SMArrow/.style={-{Stealth[length=2mm]}},                         % the usual arrow tip (54x)
}

% \SMTikzmark{name} -- a zero-size remember-picture node at the current
% text position (on the baseline, 2pt up), for arrows/braces that point at
% a word. Three slides defined their own identical copy (\tikzmark on 33r
% and ha9 -- two \newcommand's of one name, an error the moment both land in
% one deck -- and ebq's \ebqTikzmark).
\newcommand{\SMTikzmark}[1]{\tikz[remember picture, overlay, baseline] \node[inner sep=0pt, anchor=base, yshift=2pt] (#1) {};}

% \SMTikzword{name}{text} -- typesets text as usual, but as a named
% remember-picture node, so an overlay arrow can point at the word itself
% ((name.south), (name.north east), ...) and follows it when the layout
% moves (2026-10-01, direct instruction: 3ac's arrow at measured page
% coordinates had drifted off its words). Transcripts keep the text.
\newcommand{\SMTikzword}[2]{\tikz[remember picture, baseline=(#1.base)] \node[inner sep=0pt, outer sep=0pt] (#1) {#2};}

% \SMLineTop -- a zero-width strut as tall as a text line's *line box*
% (\SMFontAscender em above the baseline, the text font's ascender, 0.935 for Fira Sans: the box the Top Margin is
% measured on). Put at the start of the first line of a text column that sits
% next to a figure in top-aligned minipages,
%   \begin{minipage}[t]{0.4\linewidth}\vspace{0pt}
%   	\begin{itemize}\item \SMLineTop First item ...
%   \end{minipage}\hfill
%   \begin{minipage}[t]{0.56\linewidth}\vspace{0pt} <figure> \end{minipage}
% and the text's line box starts exactly where the figure starts -- so the
% automatic Top Margin puts both on it. Without it the first line's box is
% only as tall as its glyphs and the figure lands ~1.8pt below the margin
% (2026-10-01, m2k: its two minipages were vertically *centred*, LaTeX's
% default, and the text column started 12pt below the Top Margin;
% check_layout.py's `coltop` check now reports that pattern).
\newcommand{\SMLineTop}{\rule{0pt}{\SMFontAscender em}}

% \EUR -- the euro sign in Fira Sans, in text and in math alike. In math,
% \text{\texteuro} came out in the math font's style ("Euro symbols ...
% proper", 2026-07-25); slides wrote \text{\sffamily€} by hand (34x).
\newcommand{\EUR}{\ifmmode\text{\sffamily€}\else\texteuro\fi}

% \shrinklists{n} -- all list levels n font steps smaller, for a slide with
% more text than fits (put it inside the frame, like the \setbeamerfont
% lines it replaces; the frame's group resets it). The beamer defaults are
% normal/\small/\footnotesize for levels 1-3; n=1 gives
% \small/\footnotesize/\scriptsize, n=2 \footnotesize/\scriptsize/\scriptsize,
% n=3 \scriptsize on all three levels (the floor; v5.30.0 -- before, every
% n>=2 meant n=2). 13 slides wrote 2-3 \setbeamerfont lines for this, in five
% different combinations.
% The item gaps shrink with the font since 2026-10-01 -- see \SMNoListParsep
% next to the itemize code.
\newcommand{\shrinklists}[1]{%
	\ifcase#1\relax
	\or
		\setbeamerfont{itemize/enumerate body}{size=\small}%
		\setbeamerfont{itemize/enumerate subbody}{size=\footnotesize}%
		\setbeamerfont{itemize/enumerate subsubbody}{size=\scriptsize}%
	\or
		\setbeamerfont{itemize/enumerate body}{size=\footnotesize}%
		\setbeamerfont{itemize/enumerate subbody}{size=\scriptsize}%
		\setbeamerfont{itemize/enumerate subsubbody}{size=\scriptsize}%
	\else
		\setbeamerfont{itemize/enumerate body}{size=\scriptsize}%
		\setbeamerfont{itemize/enumerate subbody}{size=\scriptsize}%
		\setbeamerfont{itemize/enumerate subsubbody}{size=\scriptsize}%
	\fi
}

% The "padding gotcha" promised above, finally root-caused (2026-07-23,
% background-task follow-up to the IA2 fine-tuning batch #2 anomaly on
% slides/m2k-Exercise-CAPM-2-current.tex/
% slides/n8e-Tail-Risk-Variance-Clustering-current.tex/
% slides/fm6-Tail-Risk-Market-vs-Simulated-current.tex/
% slides/xqq-Fat-Tails-current.tex, all fixed there with a pragmatic
% `line width=1.2pt` override and an explicitly *unresolved* root cause at
% the time -- "declares an identical stroke width and CTM in the compiled
% PDF's own raw content stream, yet renders at a genuinely different pixel
% thickness", with CTM scale/clips/dash/cap/join all directly checked and
% ruled out). Confirmed via a controlled reproducer (four `SMImageBorder`
% nodes -- plain opaque image, nested pgfplots axis, plain image at
% `inner sep=3pt`, nested plain tikzpicture at `inner sep=3pt` -- rendered
% at 600dpi and measured by a half-max cross-section profile, not just
% eyeballed): the one variable that actually predicts the rendered width is
% `inner sep`, not content type at all. At `inner sep=0pt` a *fully opaque*
% image measured 1.2pt; the exact same image at `inner sep=3pt` measured
% 2.52pt (matching the declared 2.5pt); a nested pgfplots axis and a plain
% nested tikzpicture, both at `inner sep=3pt`, also measured 2.52pt. Two
% more points confirmed it's a smooth transition, not a binary switch:
% `inner sep=0.6pt` measured 1.8pt (= the 1.2pt floor plus that 0.6pt gap,
% exactly); `inner sep=1.25pt` (= \SMImageBorderWidth/2, the stroke's
% own half-width) already measured the full 2.52pt.
%
% The mechanism: tikz draws a node's border stroke centered on the node's
% own boundary path, then paints the node's *content* on top of it,
% afterwards. At `inner sep=0`, the content's own bounding box coincides
% exactly with that boundary path, so if the content is opaque all the way
% to its own edge, it paints directly over the stroke's inner half,
% visually leaving only the outer half (`\SMImageBorderWidth`/2,
% ~1.2pt) -- not a rendering bug, just tikz's normal draw order interacting
% with zero padding. A gap of at least `\SMImageBorderWidth`/2 (or
% content that's transparent right at its own edge -- see below) leaves the
% border with nothing opaque to cover it, so it renders at its true
% declared width instead.
%
% This also explains why the fix landed on exactly the 4 slides it did,
% each for a different reason: slides/m2k-Exercise-CAPM-2-current.tex
% used `inner sep=3pt` for unrelated layout reasons (room for the axis's
% own tick labels), well past the 1.25pt threshold, so nothing ever covered
% its border. slides/n8e-Tail-Risk-Variance-Clustering-current.tex/
% slides/fm6-Tail-Risk-Market-vs-Simulated-current.tex/
% slides/xqq-Fat-Tails-current.tex all use `inner sep=0pt` -- matching the
% deck's own overwhelming convention -- but their own chart PNGs
% (`chart-index-daily-returns-*.png`/`chart-index-fat-tails-histogram*.png`,
% matplotlib exports with a transparent background) measured alpha=0 at
% every edge pixel (confirmed directly, Pillow), so there was no opaque
% content to cover their own border either. Every other `SMImageBorder`
% use in this project that was never flagged uses an opaque image (a JPG,
% or a PNG with no transparent margin) at `inner sep=0pt` -- which
% *accidentally* renders at the ~1.2pt covered-half width instead of the
% full 2.5pt this constant nominally declares, and that accidental width is
% what reads as "the deck's normal border thickness" throughout, since it's
% by far the most common case.
%
% Practical upshot for a *new* slide: a `SMImageBorder` node with
% `inner sep=0pt` wrapping opaque content (the common case) needs no
% special handling -- it will match the rest of the deck automatically. A
% node that instead uses a nonzero `inner sep`, or wraps content with a
% transparent edge (a matplotlib export, an SVG-derived PDF with padding,
% ...), will render at the *full* \SMImageBorderWidth and look
% noticeably heavier next to the deck's other borders -- either add an
% explicit `line width=\SMImageBorderWidth/2` (matching the 4 fixed
% slides' own literal 1.2pt) to match the accidental deck-wide look, or
% reduce `inner sep` below \SMImageBorderWidth/2 so the same coverage
% effect kicks in naturally. The 4 existing per-slide overrides above are
% left as-is -- they're already the correct fix for what's now a fully
% understood cause, not a mystery workaround needing revisiting.

% Reusable "mugshot" -- a circular-cropped photo, bordered like any other
% SMImageBorder image, with a caption curving along its own bottom arc
% (2026-07-06, direct instruction, generalized from
% slides/4uq-Definition-of-Capital-Market-Efficiency-current.tex's own
% Eugene Fama photo -- see that slide's earlier comment for the source
% PPTX effect this reproduces, `<a:prstGeom prst="ellipse">` +
% `<a:prstTxWarp prst="textArchDown">`, and CONVERSION_NOTES.md's own
% "beamer/TeX gotchas" entry, search "circular photo crop", for the
% clip/border/decorations.text mechanics themselves -- unchanged here,
% just factored out so a second photo doesn't need to repeat them).
%
%   \coordinate (myMugshotCenter) at (...);
%   \SMMugshot{myMugshotCenter}{media/photo-example.jpeg}{2.6cm}{Some Person, NP 2013}
%
% #1 the *name* of a coordinate the caller has already placed (not a raw
% point expression -- needed as a plain name, not just a point, since the
% caption's own arc path re-references it by name via `++(...)`, which
% only works with a named coordinate/node, not an arbitrary `($...$)`
% calc expression); #2 the image path; #3 the desired final diameter;
% #4 the caption text.
%
% Scales the source photo so it fully covers the circle with no gap
% anywhere -- both dimensions end up at least as big as the diameter --
% while using as much of the source photo as possible (2026-07-07,
% direct instruction: "should not include whitespace... scale the photo
% such that its edges just touch the circle's inner border"). This
% briefly (same day) scaled the *other* way -- by the photo's own
% diagonal, the largest scale at which the *entire* rectangle fits
% inside the circle with nothing clipped at all -- which does avoid
% ever clipping real content, but leaves the circle's own four corners
% uncovered (visible as plain background there, since a non-circular
% photo can only touch an inscribing circle at its corners, not along
% the whole rim) -- exactly the "whitespace" this round's feedback
% ruled out in favor of full coverage. Covering scale is
% `diameter / min(w,h)`: whichever of the photo's own two dimensions is
% *smaller* becomes exactly the diameter (that edge pair touches the
% circle's own inner border, per the direct instruction), and the other,
% larger dimension overhangs the diameter and gets clipped by the
% circular `\clip` -- for a portrait photo (every one used here so far),
% that means width becomes the diameter and a sliver clips off the top
% and bottom, which is the origin of the clipping problem the diagonal
% approach above existed to solve in the first place. Reintroduced
% deliberately: showing as much of the photo as the "no whitespace"
% constraint allows is now the stated priority, over never clipping
% anything.
%
% Measuring the photo's own native size needs `\savebox` (into
% \SMMugshotBox, declared once below, not a fresh box register per
% call) -- `\includegraphics` with no size option at all renders at the
% image's native resolution, and `\wd`/`\ht` then read that box's actual
% dimensions. pgfmath (already loaded for tikz) converts any TeX
% dimension fed into it to a plain pt-count number automatically, so
% arithmetic on the two measured lengths just works. `min(1,ratio)`
% (ratio = height/width) picks out whichever of the two is smaller
% *relative to the other*, without needing a `sqrt`/`^` at all this time
% -- so this version doesn't hit the "! Dimension too large" pgfmath
% overflow the diagonal version's own arithmetic did (see this file's
% own comment history / CONVERSION_NOTES.md's "beamer/TeX gotchas" for
% that point (easy to misread as a harmless warning, since compilation
% continues and produces *a* number regardless). Reformulating via the
% aspect ratio sidesteps the whole problem: `scale = diameter /
% sqrt(1+ratio^2)` is the same formula algebraically (factor `w` out of
% `sqrt(w^2+h^2)` and it becomes `w*sqrt(1+(h/w)^2)`), but every
% the full story of that overflow, from when this macro's own arithmetic
% did use a `sqrt`).
%
% Caption font reverted to this deck's normal body font (2026-07-07,
% direct instruction: "change the mugshot template back to using regular
% font instead of monospace") -- `\ttfamily` (Fira Mono) was a direct
% instruction on 2026-07-06, one day earlier; \normalfont resets
% explicitly rather than just omitting \ttfamily, so this doesn't
% silently inherit whatever font happens to be active at the call site.
\newsavebox{\SMMugshotBox}
\newcommand{\SMMugshot}[4]{%
	\savebox{\SMMugshotBox}{\includegraphics{#2}}%
	\pgfmathsetmacro{\SMMugshotImgW}{\wd\SMMugshotBox}%
	\pgfmathsetmacro{\SMMugshotImgH}{\ht\SMMugshotBox}%
	\pgfmathsetmacro{\SMMugshotRatio}{\SMMugshotImgH/\SMMugshotImgW}%
	\pgfmathsetmacro{\SMMugshotScaledW}{(#3)/min(1,\SMMugshotRatio)}%
	\begin{scope}
		\clip (#1) circle (0.5*#3);
		\node[anchor=center, inner sep=0pt] at (#1)
			{\includegraphics[width=\SMMugshotScaledW pt]{#2}};
	\end{scope}
	\draw[SMImageBorder] (#1) circle (0.5*#3);
	{\normalfont\tiny
	\path[decoration={text along path, text={#4}, text align=center}, decorate]
		(#1) ++(200:0.5*#3+0.25cm) arc (200:340:0.5*#3+0.25cm);
	}
}

% Reusable "warning sign" -- a rotated, no-fill, thick-bordered box with
% centered bold text, generalized (2026-07-24, direct instruction: "build a
% GENERAL reusable command for placing such signs") from the red
% "Extremely skewed outcomes!" callout on slides/awe-Example-Investment-
% Return-Simulation-current.tex's own source PPTX slide (233) -- read
% directly from that slide's `ppt/slides/slide233.xml`: a rotated
% no-fill rectangle (`rot="20410887"`, i.e. 340.18 degrees clockwise =
% -19.82 degrees), a thick solid border in the same color as the text
% (`<a:ln w="57150">` = 4.5pt, color C00000 = this project's own
% `c_custom_red`), text centered inside.
%
% Usage: \SMWarningSign[color]{x}{y}{text}{aspect ratio}{rotation}
% #1 (optional) the border/text color, defaults to \c_custom_red; #2/#3
% the sign's own center, as a page-absolute (x,-y) offset from
% "current page.north west" in pt -- the same coordinate convention as
% this deck's other absolute-overlay figures (e.g.
% slides/ivw-Divestment-of-Polluting-Plants-Overview-current.tex's own
% photo nodes); #4 the sign's text; #5 the desired width:height aspect
% ratio of the box; #6 the rotation angle in degrees, tikz convention
% (counterclockwise positive) -- note OOXML's own `rot` attribute on a
% source PPTX shape is clockwise positive instead, so a value read
% directly off a source slide needs negating first.
%
% The box is sized to fit the text *wrapped* at the requested aspect
% ratio, not a caller-supplied width -- matching how the source PPTX's
% own sign isn't a single unwrapped line either (rendering the actual
% source deck's own slide 233 to a flat image for comparison confirmed
% its text wraps across 3 lines inside a squarish box, not one long
% strip). #5 drives that wrap width via a closed-form estimate (start
% from the unwrapped text's own natural width #4 and single-line height
% #4, i.e. `sqrt(ratio*width*height)`, the width at which a box of
% roughly that many lines would hit the target ratio if text-ink area
% stayed constant under rewrapping) plus one correction pass (rewrap at
% that estimate, remeasure the real height, then solve
% `sqrt(width*height*ratio)` again treating the *measured* width*height
% as the area instead of the rough estimate) -- not exact, since
% rewrapping to a handful of lines isn't a smooth function of width, but
% close enough for the short warning-sign text this is meant for; a
% caller can always nudge #5 if a specific case lands visibly off.
% (\SMWarningSignPadding, 6pt, and \SMWarningSignBorderWidth, 4.5pt, are hooks of SlideMage_setup.tex.)
\newsavebox{\SMWarningSignBox}
\newlength{\SMWarningSignW}
\newlength{\SMWarningSignH}
\newcommand{\SMWarningSign}[6][c_custom_red]{%
	\savebox{\SMWarningSignBox}{\bfseries\color{#1}#4}%
	\pgfmathsetlength{\SMWarningSignW}{sqrt((#5)*\wd\SMWarningSignBox*(\ht\SMWarningSignBox+\dp\SMWarningSignBox))}%
	\savebox{\SMWarningSignBox}{\parbox{\SMWarningSignW}{\centering\bfseries\color{#1}#4}}%
	\pgfmathsetlength{\SMWarningSignW}{sqrt((#5)*\SMWarningSignW*(\ht\SMWarningSignBox+\dp\SMWarningSignBox))}%
	\savebox{\SMWarningSignBox}{\parbox{\SMWarningSignW}{\centering\bfseries\color{#1}#4}}%
	\setlength{\SMWarningSignH}{\ht\SMWarningSignBox+\dp\SMWarningSignBox}%
	% \vbox to 0pt{\vss ...} keeps this at zero space in the normal vertical
	% list, the same reason every other absolute-page-coordinate tikz overlay
	% in this file uses it (footer/logo/title-bar boxes, y5z's/22s's own
	% corner images, ...): a bare tikzpicture here is otherwise sized from
	% its own local origin (wherever this macro happens to be called in the
	% flow) out to the far-away absolute page coordinate the node is drawn
	% at, so without this wrapper the "sign" silently pushes real, sizeable
	% blank space into the frame's own body -- confirmed directly on
	% slides/awe-Example-Investment-Return-Simulation-current.tex: with this
	% call gated behind `\deckreveal[3]{...}`, the reveal step where it first
	% turns on pushed the *unrelated* left-column text down by about 100pt
	% relative to the steps before it, a "content moves between reveal
	% steps" bug rather than anything wrong with that text itself.
	\vbox to 0pt{\vss
	\begin{tikzpicture}[overlay, remember picture]
		\node[rotate=#6, draw=#1, line width=\SMWarningSignBorderWidth,
		      minimum width=\dimexpr\SMWarningSignW+2\SMWarningSignPadding\relax,
		      minimum height=\dimexpr\SMWarningSignH+2\SMWarningSignPadding\relax,
		      inner sep=0pt] at ($(current page.north west)+(#2,-#3)$)
			{\usebox{\SMWarningSignBox}};
	\end{tikzpicture}}%
}

% A \frac/\dfrac-stacked formula reads as cramped against a plain-text
% neighbor in an itemize -- it's visibly taller than a normal line (both
% above *and* below the line, since the fraction bar sits on the text
% baseline with numerator/denominator stacked around it), so both the
% fraction's own item *and* the item that follows it need a bit of extra
% headroom that a plain-text \item doesn't. Started as a manually-repeated
% \vspace{3pt} before each such \item (first on
% slides/u74-Example-Dynamic-Mattress-Company-Periods-current.tex);
% \tallitem below is that same convention factored into a single drop-in
% \item replacement, so a future fraction-heavy slide (or a change to the
% spacing amount) doesn't mean re-editing several slide files by hand.
% Takes an optional leading multiplier (default 1, i.e. 3pt): \tallitem[1.5]
% for 4.5pt, for a formula tall enough that the default doesn't read as
% enough room. Usage: replace \item with \tallitem both on any itemize line
% containing a stacked fraction *and* on the item immediately following it
% (same nesting level; skip this second replacement if that next item is
% itself already a \tallitem for its own fraction, or if the fraction item
% is the last one in its list) -- it's a plain substitute, not an
% environment, so it doesn't change how the surrounding list itself is
% structured. (Named \SMFracItem originally; renamed 2026-07-21 to the
% more discoverable \tallitem, and given the optional multiplier and the
% following-item rule above at the same time -- no behavior change for any
% existing default-strength call site.)
% Ratio, not a frozen length (2026-07-22, direct instruction: "modify
% \tallitem and \tallitemnobullet to also use scaling relative to font
% size") -- the original \newlength+\setlength{...}{3pt} pair freezes to
% one fixed pt value the moment \setlength runs, the same issue (and same
% fix) as \SMItemsepRatioI/II/III/IV further below: converted to a
% bare ratio macro, with the "em" unit written back on inside \tallitem's
% and \tallitemnobullet's own bodies (not here), so it's recomputed fresh
% against whatever font is actually live at each call site -- correct
% regardless of nesting level or a \footnotesize/\scriptsize override,
% unlike the old fixed 3pt (calibrated against one specific level's font,
% silently wrong at any other). 0.3 reproduces the original 3pt exactly
% at level 2's own font (10pt) -- the level \tallitem was first used at,
% on slides/u74-Example-Dynamic-Mattress-Company-Periods-current.tex (see
% this macro's own comment above) -- so no visual change there, while now
% correctly scaling up/down at every other level and font override.
% (\SMTallItemBaseRatio, 0.3, is a hook of SlideMage_setup.tex.)
\newlength{\SMTallItemBaseUnit}
\newlength{\SMTallItemSpace}
% The optional argument is either a number (the multiplier above -- any
% argument made only of digits and dots) or anything else, which is taken as
% a bullet label and passed on as \item[<label>] with the default spacing:
% \tallitem[\SMBulletEllipsis] gives an ellipsis-labelled item with the
% usual headroom (2026-09-26, direct instruction), \tallitem[] a label-less
% one. (For a label-less item with a custom multiplier, use
% \tallitemnobullet[<multiplier>].)
\ExplSyntaxOn
\cs_new_protected:Npn \SMTallItemIfNumberTF #1
	{ \regex_match:nnTF { \A [0-9.]+ \Z } {#1} }
\ExplSyntaxOff
\newcommand{\tallitem}[1][1]{%
	\setlength{\SMTallItemBaseUnit}{\SMTallItemBaseRatio em}%
	\SMTallItemIfNumberTF{#1}%
		{\setlength{\SMTallItemSpace}{#1\SMTallItemBaseUnit}%
		 \vspace{\SMTallItemSpace}\item}%
		{\setlength{\SMTallItemSpace}{\SMTallItemBaseUnit}%
		 \vspace{\SMTallItemSpace}\item[#1]}}

% For a line that needs \tallitem's own extra headroom but no bullet at
% all -- a display equation continuing the previous \item's own sentence,
% rather than a new bulleted point of its own (see README.md's own
% \tallitem section, which already documents this exact combination as a
% rare, previously hand-written case: "\vspace{\SMTallItemBaseUnit}\item[]").
% \tallitemnobullet (2026-07-21, direct instruction) factors that hand-
% written idiom into the same kind of drop-in replacement \tallitem
% already is for its own case, same optional multiplier: \tallitemnobullet[1.5].
\newcommand{\tallitemnobullet}[1][1]{%
	% Recomputes \SMTallItemBaseUnit itself, fresh, rather than reusing
	% whatever \tallitem's own body last left it at -- since it's now a
	% per-call, live "em"-based value (see \tallitem's own comment above),
	% not the one-time preamble-load-time constant it used to be, relying
	% on call order/a preceding \tallitem to have set it correctly would
	% silently break the moment \tallitemnobullet is used on its own, or
	% at a different font size than whichever \tallitem last ran.
	\setlength{\SMTallItemBaseUnit}{\SMTallItemBaseRatio em}%
	\setlength{\SMTallItemSpace}{#1\SMTallItemBaseUnit}%
	\vspace{\SMTallItemSpace}\item[]}

% flalign* (used to left-align a display equation inside an itemize --
% see CONVERSION_NOTES.md's "flalign* for left-aligning..." entry) inherits
% plain LaTeX's \abovedisplayskip/\belowdisplayskip, sized for a full
% document page, not a compact slide itemize -- reads as far too much blank
% space above/below the equation next to the surrounding item text (found
% 2026-07-21, direct instruction; most visible on
% slides/bwi-Pearson-Correlation-Definition-current.tex's own single-line
% equation, but the same oversized gap is present everywhere flalign* is
% used this way, e.g. slides/7se-Portfolio-Risk-Three-Assets-current.tex
% and slides/rb3-Asset-Allocation-Example-CAL-current.tex). \SMTightDisplaySkips
% tightens both skips -- and their "short" variants, since amsmath picks
% between the two depending on how much text precedes the display, so
% setting only one pair left the gap unfixed whenever the preceding line
% happened to be short -- down to something that reads as part of the
% same list rather than a separate floating display. Usage: wrap the
% *real*, unmodified flalign* in a plain TeX group together with this
% command, `{\SMTightDisplaySkips\begin{flalign*} ... \end{flalign*}}`,
% rather than a new same-named environment -- tried that first
% (`\newenvironment{SMTightFlalign*}{...\begin{flalign*}}{\end{flalign*}}`)
% and it broke LaTeX's own environment-nesting tracking (`\@currenvir`),
% surfacing as "\begin{flalign*} ... ended by \end{SMTightFlalign*}"
% and a fatal, PDF-less compile on the very first slide that used it. A
% plain group sidesteps that class of bug entirely since it never
% introduces a same-named wrapper environment.
% (\SMTightDisplaySkip, 3pt, is a hook of SlideMage_setup.tex.)
\newcommand{\SMTightDisplaySkips}{%
	\abovedisplayskip=\SMTightDisplaySkip
	\abovedisplayshortskip=\SMTightDisplaySkip
	\belowdisplayskip=\SMTightDisplaySkip
	\belowdisplayshortskip=\SMTightDisplaySkip}


% \SMFrametitleRightPadding (left/right padding around the logo, chosen
% to match its top/bottom margins -- measured against the title colorbox's
% own top and bottom edges -- rather than the theme's 2.2ex default) is
% also set by the per-university preamble, not here: the top/bottom margins
% are a side effect of vertically centering a \SMLogoHeight-tall logo in
% the title bar, not something set directly, so this value has to be
% re-measured and set per institution whenever its own logo/height differs
% from another's.

% The logo macro \SMInstituteLogo, the frame title hook that paints the white patch behind the
% logo and redraws it (\SMFrametitleHookCode), \chapterlogo and the template without the inline
% logo (SMNoInlineLogo) are the theme's (beamerthemeSlideMage.sty). The logo width, which the
% hook uses to size the white patch, is measured here from the template's logo path and height.
\settowidth{\SMLogoWidth}{\includegraphics[height=\SMLogoHeight]{\SMInstituteLogoPath}}



% Define slideID as a simple value holder (initially empty)
\newcommand{\slideID}{}%

% Set (locally, by \slidestep) when this inclusion of a slide id must not define
% the id's \hypertarget again, because an earlier \slidestep of the same id
% already owns it (see \slidestep's own comment).
\newif\ifSMSuppressSlideTarget

% Define slideID as a simple value holder (not a command with arguments)
\makeatletter
\newcommand{\setSlideID}[1]{%
	% Grab everything up to the first '-' and store it in \slideID
	\edef\tempa{#1}%
	\expandafter\setSlideID@aux\tempa-\relax
}
\def\setSlideID@aux#1-#2\relax{%
	\renewcommand{\slideID}{#1}%
	% Non-linear navigation (direct instruction, 2026-09-15): every real
	% slide inclusion -- \slide{}, \slidefinal{} (delegates to \slide{}),
	% \slideversion{}, \omitslide{}, \omitslideversion{}, \conditionalslide{}
	% -- calls \setSlideID{} immediately before \input{}ing the slide file,
	% regardless of which frame type (standardframe, sideimageframe,
	% figureframe, sectionframe, highlightframe, ...) that file happens to
	% use internally. This is the one place all of them funnel through, so
	% it is the correct place for a single, frame-type-independent
	% \hypertarget{} -- giving every slide a free, permanent jump target
	% keyed to its own existing id, with no per-slide authoring and no new
	% bookkeeping. Deliberately \hypertarget (hyperref's own mechanism, not
	% plain LaTeX \label{}: slides/ZZM-EN-Slidemaster-colors-current.tex's
	% own comment documents a real, already-hit collision between \label{}
	% and pgffor's \foreach inside a tikz loop, which this sidesteps
	% entirely. Used by \hyperlink{<bare id>}{...} anywhere in the deck
	% (README.md's own "Non-linear navigation" section) to jump straight to
	% this slide -- e.g. an excursus link's own "Back" button, or a
	% generated appendix's thumbnail grid (SlideMage's "Generate
	% appendix TOC" checkbox, appendix_toc.py).
	%
	% Known, pre-existing limitation, not introduced by this: if the same
	% id is genuinely \slide{}'d twice in one deck (it happens -- e.g. a
	% \slidefinal{} recap of an earlier slide), hyperref just warns
	% ("multiply defined") and the *later* occurrence's target wins, so a
	% link to that id lands on whichever instance is last in the document
	% -- inherent to any permanent-id-keyed target scheme, not specific to
	% this one.
	\ifSMSuppressSlideTarget\else\hypertarget{\slideID}{}\SM@agenda@slide{#1}{#2}\fi%
}
\makeatother

% \setSlideVersion{X-Y-Z} -- footer's own version display (added
% 2026-08-20, direct instruction: show the slide's version next to its id
% in every compiled deck). Takes the same filename-safe hyphenated form a
% version file's own "-vX-Y-Z.tex" suffix and \slideversion{}'s own second
% argument already use, and reformats it with dots ("1.0.0") for display,
% the conventional way semver is written -- \setSlideVersion@aux splits it
% the same way \setSlideID@aux above splits an id-slug, just at both
% hyphens instead of the first one only. \ifSMSlideVersionSet tracks
% whether *this* frame actually has a tracked version to show at all --
% reset false at the top of every \slide{}/\omitslide{}/\conditionalslide{}
% call (current.tex may or may not embed a \setSlideVersion{} of its own,
% see SlideMage's own automatic-versioning save flow) so an untracked
% slide's footer never inherits whatever the immediately preceding frame
% last set. \slideversion{}/\omitslideversion{} instead set it directly
% from their own second argument -- already an explicit version number at
% the call site, so there is nothing to read from the pinned file itself.
\newcommand{\slideVersion}{}%
\newif\ifSMSlideVersionSet
\makeatletter
\newcommand{\setSlideVersion}[1]{%
	\edef\SMTempSlideVersion{#1}%
	% No trailing "-" appended before \relax here, unlike \setSlideID@aux's
	% own \tempa-\relax above -- that extra hyphen is harmless there since
	% \setSlideID@aux's own #2 (everything after the first "-") is thrown
	% away entirely, but this pattern's #3 is a *used* greedy match
	% ("everything up to \relax"), so an appended "-" would land inside it
	% verbatim (confirmed directly: produced "v2.3.1-" instead of
	% "v2.3.1" the first time this was compiled).
	\expandafter\setSlideVersion@aux\SMTempSlideVersion\relax
	\global\SMSlideVersionSettrue
}
\def\setSlideVersion@aux#1-#2-#3\relax{%
	\renewcommand{\slideVersion}{#1.#2.#3}%
}
\makeatother

% \setSlideLanguageFromSlug{ID-slug} -- multi-language support (added
% 2026-08-24). Language lives in the id-slug's own *second* hyphen segment
% rather than a separate in-body macro call, so it's visible directly in
% Explorer/grep and in every \slide{}/\slideversion{}/etc. call site, not
% just inside the file it names: slides/<id>-EN-<slug>-current.tex,
% slides/<id>-DE-<slug>-current.tex, slides/<id>-mult-<slug>-current.tex
% (mult = multilingual, i.e. a single file holding both languages via
% \multilingualtext{}{} below, picked at compile time by \ifSMGerman).
% A slide that predates this feature simply has no such segment (its own
% second hyphen segment is the first word of an ordinary slug, e.g.
% "neb-Market-Environment-...") -- \ifSMSlideLanguageSet stays false
% for it, exactly like an untracked \setSlideVersion{} today shows no
% "v..." footer segment, so the other ~665 files need no migration.
% Same three-way split idiom as \setSlideVersion@aux above, but the
% trailing padding ("-zzz-zzz") is load-bearing here in a way it isn't
% there: \setSlideVersion@aux's own caller (X-Y-Z) always has exactly two
% hyphens, but an id-slug can legally have as few as *one* (e.g. a
% single-word slug like "abc-Word"), which would leave #2-#3\relax with no
% closing delimiter to match at all ("Use of ... doesn't match its
% definition") without this padding. #3 (and the padding itself) is never
% read, so its exact content doesn't matter -- only that the two extra
% hyphens are guaranteed to exist.
\newcommand{\slideLanguage}{}%
\newif\ifSMSlideLanguageSet
% \ifSMSlideLanguageMult -- set alongside \slideLanguage whenever it's
% specifically "mult", so the footer tag below never has to run
% \ifstrequal against \slideLanguage itself: etoolbox's \ifstrequal
% compares via \detokenize, which does not expand its argument -- fed a
% bare control sequence like \slideLanguage it would stringify the
% *macro's name*, not its meaning, and always take the false branch (see
% \decklanguage's own comment above, where this was caught the same way).
% A dedicated boolean, set right here where the value is still a plain
% matched string (#2, not a macro), sidesteps the trap entirely.
\newif\ifSMSlideLanguageMult
\makeatletter
\newcommand{\setSlideLanguageFromSlug}[1]{%
	\global\SMSlideLanguageSetfalse
	\global\SMSlideLanguageMultfalse
	\edef\SMTempSlideLangSlug{#1}%
	\expandafter\setSlideLanguageFromSlug@aux\SMTempSlideLangSlug-zzz-zzz\relax
}
\def\setSlideLanguageFromSlug@aux#1-#2-#3\relax{%
	\ifstrequal{#2}{EN}{\renewcommand{\slideLanguage}{EN}\global\SMSlideLanguageSettrue}{%
	\ifstrequal{#2}{DE}{\renewcommand{\slideLanguage}{DE}\global\SMSlideLanguageSettrue}{%
	\ifstrequal{#2}{mult}{\renewcommand{\slideLanguage}{mult}\global\SMSlideLanguageSettrue\global\SMSlideLanguageMulttrue}{}}}%
}
\makeatother

% Define a new \slide command to automatically pick the current version of the slide
% 2026-07-08 (direct instruction): slide filenames dropped their leading
% "s-" (e.g. slides/ab3-Some-Slide-current.tex, not the old
% slides/ab3-Some-Slide-current.tex) so slide IDs are directly
% searchable in Explorer -- this macro (and the other three below) is the
% *only* place that ever needed to change for that: every deck's own
% \slide{}/\slideversion{}/\omitslide{}/\omitslideversion{} calls pass just
% "ID-slug", never the filename itself, so renaming every file on disk to
% match needed no changes anywhere else (not the deck files, not
% slide_ids.csv, which only ever stored bare IDs to begin with).
\newcommand{\slide}[1]{%
	\global\SMSlideVersionSetfalse
	\setSlideID{#1}%
	\setSlideLanguageFromSlug{#1}%
	\input{slides/#1-current.tex}%
}

% \slidefinal{ID-slug} -- like \slide{}, but renders the slide at its final,
% fully-revealed overlay state in every deck version, instructor included
% (2026-07-12, direct instruction: show jww's finished build in place of a
% separate recap slide, "irrespective of which deck version is being
% compiled" -- that recap slide has since been deleted entirely, its own ID
% no longer in use). A slide
% built with \conditionalpause[instructor] already renders fully built in
% the during-course/after-course compiles -- that pause is a no-op for
% them, so nothing there ever waits on a reveal step -- the only case
% actually needing help is the instructor deck, which would otherwise still
% step through the build. Locally shadowing \conditionalpause to a no-op
% (same optional-argument signature as the real one, so it still accepts
% e.g. \conditionalpause[instructor] without erroring) for the duration of
% one \slide{} call makes every \conditionalpause in that one inclusion
% inert, regardless of the actual compiled \SMDeckVersion -- scoped
% inside \begingroup/\endgroup, so it reverts immediately after, and has no
% effect on any *other* \slide{}/\slidefinal{} call to the same ID
% elsewhere in the same deck (e.g. an earlier, normally-building instance of
% the same slide). \deckonly/\deckuncover/\deckvisible are untouched, so
% deck-version-tagged content (as opposed to reveal-step-tagged content)
% still follows its usual per-version rules.
\newcommand{\slidefinal}[1]{%
	\begingroup
		\renewcommand{\conditionalpause}[1][instructor]{}%
		\slide{#1}%
	\endgroup
}

% Load a specific, pinned version of a slide instead of whichever is "current"
% Usage: \slideversion{ID-slug}{M-m-p} loads slides/ID-slug-vM-m-p.tex
\newcommand{\slideversion}[2]{%
	\setSlideID{#1}%
	\setSlideLanguageFromSlug{#1}%
	\setSlideVersion{#2}%
	\input{slides/#1-v#2.tex}%
}

% Student-version decks: \studentversion (called anywhere in a deck, e.g.
% right after \input{SlideMage_template_uni-graz}, similar to \sidefootnotes or
% \setfootnotesymbols) switches that compile into the "handout" variant --
% including, as of 2026-07-03, actually forcing beamer's own "handout" mode
% (\gdef\beamer@currentmode{handout}, normally only reachable via the
% \documentclass[handout]{beamer} *class option*, set once at class-loading
% time -- too early for a command like \studentversion, called later from a
% deck's own preamble, to use that route). This was checked, not assumed:
% until this point \studentversion only ever touched \omitslide{}, so
% nothing suppressed \pause/overlay specs -- true only by coincidence, since
% no slide converted so far happens to use them. Late-redefining
% \beamer@currentmode was verified (isolated minimal-beamer tests, plain
% \pause *and* inline specs like \item<2->) to reproduce the same collapse
% -to-one-page, show-the-final-revealed-state behavior as the class option,
% even set well after \documentclass{beamer} -- beamer's overlay decoder
% (beamerbasedecode.sty) reads \beamer@currentmode at frame-processing time,
% not just once at class-option time, so this is safe, not a hack that
% happens to work by luck of file ordering.
% \omitslide{}/\omitslideversion{} are drop-in replacements for
% \slide{}/\slideversion{} that tag a slide as excluded from student deck versions:
%
%   \omitslide{xyz-Answer-key}                 % instead of \slide{...}
%   \omitslideversion{xyz-Answer-key}{1-0-0}    % instead of \slideversion{...}{...}
%
% In a \studentversion compile or a deck version other than instructor,
% a tagged slide's \input{} is skipped entirely and its frame number still
% advances, keeping numbering aligned across deck versions. In an
% instructor compile (the default), a tagged slide loads exactly as it
% would via \slide{}/\slideversion{}, with
% one difference: \SMOmittedSlide is set for the duration of that one
% \input{}, so the footer template below can print a marker next to that
% frame's slide ID -- a deliberately visible reminder, while reviewing the
% instructor deck, of which slides won't make it to students. The same tag
% works for all deck versions because it lives at the call site in the *deck*
% file (not inside the slide's own content file): the identical slide list
% can serve the instructor deck and both student deck versions, e.g.
%
%   \input{SlideMage_template_uni-graz}
%   \FooterText{...}
%   \makeatletter\ifdefined\HANDOUT\studentversion\fi\makeatother  % or similar
%   \input{course-content.tex}   % the shared \slide{}/\omitslide{} list
%
% with two thin wrapper decks (one plain, one calling \studentversion first)
% both \input-ing the same shared content file -- see README.md.
\newif\ifSMStudentVersion
% \makeatletter must wrap the \newcommand itself, not sit *inside* the macro
% body: catcodes apply at tokenization time, which for a macro body happens
% once, immediately, when \newcommand reads it -- not deferred until the
% macro is later called. \makeatletter placed inside the body (tried first)
% was too late: \beamer@currentmode had already been mis-tokenized as
% "\beamer" + the six-way-otherwise-ordinary character "@" + eleven plain
% letters "currentmode", by the time that inner \makeatletter could act on
% it, so the \gdef silently created/redefined the wrong thing and never
% touched beamer's actual mode variable. Confirmed via \meaning\beamer@
% currentmode before/after \studentversion -- unchanged with the inside-the-
% body version, fixed once \makeatletter moved out here.
\makeatletter
\newcommand{\studentversion}{%
	\global\SMStudentVersiontrue
	\gdef\beamer@currentmode{handout}%
}
\makeatother

% Deck version system - new boolean flags and string for current version
\newif\ifSMInstructorVersion
\newif\ifSMDuringCourseVersion
\newif\ifSMAfterCourseVersion
\gdef\SMDeckVersion{instructor}
\gdef\SMDeckVersionText{}

% \deckversion[] - set the deck version(s)
% Usage: \deckversion[instructor], \deckversion[during-course], etc.
% Multiple versions: \deckversion[instructor,during-course,after-course]
% Parses #1 as a comma-separated list via etoolbox's \forcsvlist (rather
% than the single \ifstrequal{whole-list}{"instructor"} check this used to
% do, which only ever matched a single-version list -- a multi-version
% \deckversion[instructor,during-course] silently set none of the three
% booleans, since the compared string was the literal, unsplit
% "instructor,during-course"). Note \forcsvlist, not etoolbox's similarly-
% named \forlistloop: \forlistloop iterates etoolbox's own internal
% pipe-delimited list format (built via \listadd), not a plain comma string
% like this command's own argument -- confirmed directly (isolated test),
% passing a raw comma string to it doesn't error or split on the commas, it
% just treats the whole string as a single one-item "list".
\makeatletter
\newcommand{\deckversion}[1][instructor]{%
	% \xdef (not \gdef): #1 needs to be fully expanded here, not just
	% stored as whatever token(s) were passed -- \deckversion[\myvar]
	% (compile-all-versions.bat's own override mechanism, via a wrapper
	% file's \def\deckversionoverride{...}) passes a macro, not a literal
	% string, and everywhere else in this deck-version system that reads
	% \SMDeckVersion only ever applies one \expandafter to it (see
	% \ifdeckversioninlist's own comment) -- enough to unwrap \SMDeckVersion
	% itself, but not a further macro still sitting inside it. Confirmed
	% directly: with the old \gdef, \deckversion[\deckversionoverride]
	% left \SMDeckVersionText empty and every \ifdeckversioninlist
	% check false, since \SMDeckVersion's own stored content was the
	% single unexpanded token \deckversionoverride, not "during-course".
	\xdef\SMDeckVersion{#1}%
	\global\SMInstructorVersionfalse
	\global\SMDuringCourseVersionfalse
	\global\SMAfterCourseVersionfalse
	\gdef\SMDeckVersionText{}%
	\def\do@deckversions##1{%
		\ifstrequal{##1}{instructor}{\global\SMInstructorVersiontrue\gdef\SMDeckVersionText{instructor deck}}{}%
		\ifstrequal{##1}{during-course}{\global\SMDuringCourseVersiontrue\gdef\SMDeckVersionText{student course deck}}{}%
		\ifstrequal{##1}{after-course}{\global\SMAfterCourseVersiontrue\gdef\SMDeckVersionText{student final deck}}{}%
	}%
	% \forcsvlist itself splits on literal commas in its own raw,
	% unexpanded argument (confirmed directly against etoolbox's own
	% source, \etb@defparser@do@aux) -- passing #1 straight through here
	% has the exact same problem as the \gdef above did, one call site
	% later. Building the whole \forcsvlist{...}{...} call inside an
	% \edef (protecting \forcsvlist/\do@deckversions themselves with
	% \noexpand so they're invoked *later*, once this constructed call is
	% run, not expanded away here) expands only \SMDeckVersion --
	% already fully resolved by the \xdef above -- into the list argument.
	\edef\SM@deckversion@call{\noexpand\forcsvlist{\noexpand\do@deckversions}{\SMDeckVersion}}%
	\SM@deckversion@call
}
\makeatother

% \deckoutputfolder[version]{path} - declares where SlideMage should
% place this deck's own compiled PDF, read directly out of this file's own
% source text by SlideMage/deck_scan.py (OUTPUT_FOLDER_RE) at deck-scan
% time -- never expanded by LaTeX itself. A silent no-op here (2026-09-16,
% direct instruction) purely so a deck carrying this directive still
% compiles cleanly under a plain lualatex/TeXstudio run, the same reason
% \deckversionoverride's own \ifdefined door (used elsewhere in this file)
% never needs a real definition either. #1 (the optional [version]) is
% intentionally unused: unlike \deckversion above, this macro's own
% argument is never inspected by LaTeX at all, only by SlideMage's own
% Python-side regex -- see README.md's "Custom output folder" section for
% the actual directive syntax.
\newcommand{\deckoutputfolder}[2][]{}

% \decktags{tag1,tag2} - declares this deck's tags (2026-09-19), read
% directly out of this file's own source text by SlideMage/deck_scan.py
% (DECKTAGS_RE) and never expanded by LaTeX itself -- a silent no-op here
% for exactly the same reason \deckoutputfolder above is. SlideMage's deck
% tables group decks by their tag: work (the pinned workdeck), course,
% research, test; a deck with none of those is a course deck. Written and
% edited from SlideMage itself (the "Tags" button, and the Duplicate
% dialog), placed just before \begin{document}.
\newcommand{\decktags}[1]{}

% Deck language system (added 2026-08-24) -- fixed per deck, like
% \FooterText or the institution \input{} choice, not a per-compile-run
% toggle like \deckversion above: a German course is its own deck file
% with mostly \setSlideLanguageFromSlug-tagged German/multilingual slide
% IDs, not the same deck file recompiled twice. \decklanguage[en]/
% \SMDeckLanguage/\ifSMGerman mirror \deckversion/\SMDeckVersion
% mechanically (same \xdef for the same \deckversion[\deckversionoverride]-
% style macro-argument reason documented above), so every deck gets the
% same \ifdefined\decklanguageoverride ... \fi override door
% \deckversionoverride already has, for SlideMage's own one-off render
% preview (slide_preview.py) and for ad-hoc hand-testing via
% decks/#workdeck.tex -- never wired into the Active-decks table itself.
\newif\ifSMGerman
\gdef\SMDeckLanguage{en}
\makeatletter
\newcommand{\decklanguage}[1][en]{%
	\xdef\SMDeckLanguage{#1}%
	\global\SMGermanfalse
	% \ifstrequal{<macro>}{...} is a trap, not just here but anywhere in this
	% file: its etoolbox implementation compares via
	% \edef\tempX{\detokenize{<arg>}} -- \detokenize does NOT expand its
	% argument first, so handed a bare control sequence it stringifies the
	% *macro's name* ("\SMDeckLanguage"), never its meaning ("de"),
	% and the comparison silently always takes the false branch (confirmed
	% directly: \ifSMGerman never went true this way, even under
	% \decklanguage[de]). \deckversion's own \do@deckversions above sidesteps
	% this by building the comparison inside an \edef first, so \ifstrequal
	% only ever sees an already-expanded literal string -- same fix here.
	\def\do@decklanguage##1{%
		\gdef\SMPDFLanguage{##1}% any other language: its own code (fr, es, ...)
		\ifstrequal{##1}{en}{\gdef\SMPDFLanguage{en-US}}{}%
		\ifstrequal{##1}{de}{\global\SMGermantrue\gdef\SMPDFLanguage{de-DE}}{}}%
	\edef\SM@decklanguage@call{\noexpand\do@decklanguage{\SMDeckLanguage}}%
	\SM@decklanguage@call
	% Document-level PDF /Lang, independent of (and not blocked by) the
	% structural PDF/UA tagging this project can't use under beamer -- see
	% \SMArtifactBegin's own comment above for why full tagging is out;
	% a screen reader still uses this Catalog-level entry to pick
	% pronunciation for the whole document.
	\hypersetup{pdflang=\SMPDFLanguage}%
}
\makeatother

% \multilingualtext{EN text}{DE text} -- for the handful of light,
% image-only/title-only slides (id-slug's own language segment "mult",
% see \setSlideLanguageFromSlug above) that hold both languages in one
% file rather than being split into two separate slide IDs, switching on
% whichever language the *deck* is compiling as. The two-argument form is
% EN then DE, as it always was.
%
% \multilingualtext*{EN={Hello}, DE={Hallo}, FR={Bonjour}} -- the keyed form
% (2026-10-05, direct instruction), for any number of languages: one
% key=value pair per language, the key being the deck's \decklanguage code
% (upper or lower case: EN = en). The value is picked whose key is the
% deck's language; a deck in a language the list does not have gets the EN
% value, or the first one if there is no EN. Commas and equals signs inside a
% value need the braces shown. A deck in another language than en/de just says
% \decklanguage[fr] (fr, es, ...): \ifSMGerman stays false for it, and the
% PDF's /Lang is the code itself.
\ExplSyntaxOn
\tl_new:N \l__sm_mlt_pick_tl
\tl_new:N \l__sm_mlt_en_tl
\tl_new:N \l__sm_mlt_first_tl
\bool_new:N \l__sm_mlt_pick_bool
\bool_new:N \l__sm_mlt_en_bool
\bool_new:N \l__sm_mlt_first_bool
\cs_new:Npn \__sm_mlt_pair:nn #1#2 { \ifSMGerman #2 \else #1 \fi }
\cs_new_protected:Npn \__sm_mlt_item:nn #1#2
  {
    \bool_if:NF \l__sm_mlt_first_bool
      { \tl_set:Nn \l__sm_mlt_first_tl {#2} \bool_set_true:N \l__sm_mlt_first_bool }
    \str_if_eq:eeT { \str_casefold:n {#1} } { en }
      { \tl_set:Nn \l__sm_mlt_en_tl {#2} \bool_set_true:N \l__sm_mlt_en_bool }
    \str_if_eq:eeT { \str_casefold:n {#1} } { \str_casefold:V \SMDeckLanguage }
      { \tl_set:Nn \l__sm_mlt_pick_tl {#2} \bool_set_true:N \l__sm_mlt_pick_bool }
  }
\cs_new_protected:Npn \__sm_mlt_item_bad:n #1
  { \msg_warning:nnn { slidemage } { multilingual-key-without-value } {#1} }
\msg_new:nnn { slidemage } { multilingual-key-without-value }
  { \string\multilingualtext*:~"#1"~has~no~=~value;~ignored. }
\cs_new_protected:Npn \__sm_mlt_keyed:n #1
  {
    \group_begin:
      \bool_set_false:N \l__sm_mlt_pick_bool
      \bool_set_false:N \l__sm_mlt_en_bool
      \bool_set_false:N \l__sm_mlt_first_bool
      \keyval_parse:NNn \__sm_mlt_item_bad:n \__sm_mlt_item:nn {#1}
      \bool_if:NTF \l__sm_mlt_pick_bool
        { \exp_last_unbraced:NV \group_end: \l__sm_mlt_pick_tl }
        {
          \bool_if:NTF \l__sm_mlt_en_bool
            { \exp_last_unbraced:NV \group_end: \l__sm_mlt_en_tl }
            { \exp_last_unbraced:NV \group_end: \l__sm_mlt_first_tl }
        }
  }
\NewDocumentCommand{\multilingualtext}{s m}
  { \IfBooleanTF {#1} { \__sm_mlt_keyed:n {#2} } { \__sm_mlt_pair:nn {#2} } }
\ExplSyntaxOff

\newif\ifSMOmittedSlide
\newcommand{\SMOmittedSlideMarker}{ | omitted from student deck}
% \stepcounter{framenumber} in the skipped branches below keeps the student
% deck's slide numbers matching the instructor deck's, gap and all, so "see
% slide 12" means the same physical slide to both audiences: beamer steps
% "framenumber" exactly once per \begin{frame} (checked in
% beamerbaseframe.sty -- once per *frame*, not per overlay page, so this is
% right regardless of how many overlay steps the omitted slide would have
% had), and this is the one place that increment is skipped along with the
% frame itself, so it has to be replayed by hand.
% Both forms are instructor-only; \studentversion also skips them.
\newcommand{\omitslide}[1]{%
	\ifSMStudentVersion
		\stepcounter{framenumber}%
	\else
		\ifdeckversioninlist{\SMDeckVersion}{instructor}%
			{\global\SMOmittedSlidetrue
			 \global\SMSlideVersionSetfalse
			 \setSlideID{#1}%
			 \setSlideLanguageFromSlug{#1}%
			 \input{slides/#1-current.tex}%
			 \global\SMOmittedSlidefalse}%
			{\stepcounter{framenumber}}%
	\fi
}
\newcommand{\omitslideversion}[2]{%
	\ifSMStudentVersion
		\stepcounter{framenumber}%
	\else
		\ifdeckversioninlist{\SMDeckVersion}{instructor}%
			{\global\SMOmittedSlidetrue
			 \setSlideID{#1}%
			 \setSlideLanguageFromSlug{#1}%
			 \setSlideVersion{#2}%
			 \input{slides/#1-v#2.tex}%
			 \global\SMOmittedSlidefalse}%
			{\stepcounter{framenumber}}%
	\fi
}

% Deck version-aware overlay commands
% Syntax: \deckonly<overlay>[versions]{content}
% If [versions] omitted, defaults to all decks (instructor,during-course,after-course)
%
% Built on etoolbox's own \forcsvlist for the comma-splitting (rather than a
% hand-rolled parser: an earlier version tried to split a list at each comma
% using \newcommand's plain, non-delimited argument grabbing
% (\newcommand{\ifdeckversioninlist@split}[3]{...} called as
% \ifdeckversioninlist@split#1\@nil{#1}) -- but a plain \newcommand doesn't do
% delimited-parameter matching against \@nil the way a \def with a matching
% pattern would; it just grabs the next 3 tokens/groups one at a time. So
% e.g. "instructor" silently split into single characters ("i", "n", "s",
% ...) instead of being read as one list item, which eventually ran the
% parser off the end of its own input and surfaced as "Undefined control
% sequence" wherever the first \conditionalpause[...]/\deckonly[...] in a
% deck happened to be used -- confirmed directly against the compile log,
% traced to slide y5z-CoCo-Bonds's own \conditionalpause[instructor].
%
% NOT etoolbox's own \ifinlist{item}{list}, despite the name similarity: that
% one tests membership in etoolbox's own internal list format, built only via
% \listadd (items joined with a literal "|", not ","); fed a plain comma
% string instead, confirmed directly (isolated test) to silently return
% false for every item, including an exact single-item match.
%
% The membership test itself needs one more subtlety: #1 here is the deck
% version to search for -- always \SMDeckVersion at every call site
% below, i.e. a macro, not a literal string -- so comparing it via a plain
% \ifstrequal{##1}{#1} does not work either: etoolbox's \ifstrequal compares
% via \detokenize, which (confirmed directly, isolated test) stringifies an
% unexpanded macro token to its own control-sequence *name* ("SMDeckVersion"),
% not the current value that macro holds. \expandafter forces that expansion
% before \ifstrequal ever sees it.
\makeatletter

\let\SM@original@only\only
\let\SM@original@uncover\uncover
\let\SM@original@visible\visible

\newif\ifSM@deckversion@found

% Usage: \ifdeckversioninlist{version}{list}{true-branch}{false-branch}
\newcommand{\ifdeckversioninlist}[2]{%
	\global\SM@deckversion@foundfalse
	\def\SM@deckversion@check##1{%
		\expandafter\ifstrequal\expandafter{#1}{##1}{\global\SM@deckversion@foundtrue}{}%
	}%
	\forcsvlist{\SM@deckversion@check}{#2}%
	\ifSM@deckversion@found
		\expandafter\@firstoftwo
	\else
		\expandafter\@secondoftwo
	\fi
}

% \@ifnextchar-based parsing (the same idiom LaTeX2e's own kernel commands
% use for an optional argument that isn't a plain trailing [...]) so the
% overlay spec reads exactly like beamer's own \only<...>, with the deck-
% version list as a further, genuinely optional [...] afterwards:
%   \deckonly<2>{content}                              visible in all versions
%   \deckonly<2>[instructor]{content}                  instructor only
%   \deckonly<2>[during-course,after-course]{content}  during + after
\newcommand{\deckonly}{\@ifnextchar<{\SM@deckonly@overlay}{\SM@deckonly@overlay<1->}}
\def\SM@deckonly@overlay<#1>{\@ifnextchar[{\SM@deckonly@versions<#1>}{\SM@deckonly@versions<#1>[instructor,during-course,after-course]}}
\def\SM@deckonly@versions<#1>[#2]#3{%
	\ifSMInstructorVersion
		\SM@original@only<#1>{#3}%
	\else
		\ifdeckversioninlist{\SMDeckVersion}{#2}{\SM@original@only<#1>{#3}}{}%
	\fi
}

% Same interface as \deckonly, wrapping \uncover and \visible respectively
\newcommand{\deckuncover}{\@ifnextchar<{\SM@deckuncover@overlay}{\SM@deckuncover@overlay<1->}}
\def\SM@deckuncover@overlay<#1>{\@ifnextchar[{\SM@deckuncover@versions<#1>}{\SM@deckuncover@versions<#1>[instructor,during-course,after-course]}}
\def\SM@deckuncover@versions<#1>[#2]#3{%
	\ifSMInstructorVersion
		\SM@original@uncover<#1>{#3}%
	\else
		\ifdeckversioninlist{\SMDeckVersion}{#2}{\SM@original@uncover<#1>{#3}}{}%
	\fi
}

\newcommand{\deckvisible}{\@ifnextchar<{\SM@deckvisible@overlay}{\SM@deckvisible@overlay<1->}}
\def\SM@deckvisible@overlay<#1>{\@ifnextchar[{\SM@deckvisible@versions<#1>}{\SM@deckvisible@versions<#1>[instructor,during-course,after-course]}}
\def\SM@deckvisible@versions<#1>[#2]#3{%
	\ifSMInstructorVersion
		\SM@original@visible<#1>{#3}%
	\else
		\ifdeckversioninlist{\SMDeckVersion}{#2}{\SM@original@visible<#1>{#3}}{}%
	\fi
}

% \conditionalpause[versions] - \pause, but only takes effect in the listed
% deck versions (default: instructor only); a no-op, not even a skipped
% overlay step, in any other version
% Usage: \conditionalpause[instructor,during-course] or \conditionalpause[after-course]
\newcommand{\conditionalpause}[1][instructor]{%
	\ifdeckversioninlist{\SMDeckVersion}{#1}{\pause}{}%
}

% \deckreveal[step]{content} - the "hidden answer, click-revealed for the
% instructor, absent for during-course, pre-shown for after-course" pattern
% already documented by hand on several exercise slides (e.g. mbz/x7q/37s/
% 3cv's own header comments: "a slide/reveal-step that is hidden in the
% source PowerPoint should, by default, only be made visible on the
% instructor version"). A bare, unconditional \only<N->{content} looks like
% it does this, but doesn't: \only's overlay range registers with beamer's
% own overlay-counting pass regardless of \ifSMInstructorVersion/
% \conditionalpause -- neither of those touch \only itself, they only guard
% *other* content around it -- so a frame with nothing else establishing a
% second overlay step still gets a real second slide out of a bare
% \only<2->, in *every* deck version, not just the instructor one. Confirmed
% directly (2026-07-21, direct instruction) on 37s/3cv: both a) leaked the
% "hidden" content to during-course as a genuine second slide, the exact
% thing this pattern exists to prevent, and b) showed a blank second bullet/
% cell first on that leaked slide's own step 1 -- \only<2-> hides the
% *content*, not the space it would occupy, so during-course's own step 1
% still had nothing else deciding where the bullet/cell landed.
% \deckreveal fixes both by picking a genuinely different rendering per
% version rather than reusing the same overlay-restricted \only call
% everywhere: instructor gets the real overlay build (\only<step->,
% default step 2, matching \tallitem's own default-multiplier convention);
% during-course gets nothing at all (no token emitted, so beamer's overlay
% pass never sees an overlay spec to register); after-course gets the
% content unconditionally, with no overlay spec of its own either, so it
% doesn't silently force a second slide there either. Content-level only
% (wraps a value/expression, like \only itself) -- doesn't cover an
% itemize \item's own bullet marker, which needs its own overlay spec
% (`\item<step->`) directly on the \item, not wrapped from outside; see
% 37s's own comment for that variant, hand-written rather than factored
% into a second macro since only one call site needs it so far.
% (2026-09-30: that variant is now \deckitem<step->, below.)
\NewDocumentCommand{\deckreveal}{s O{2} m}{%
	\ifSMDuringCourseVersion
	\else
		\ifSMInstructorVersion
			\IfBooleanTF{#1}{\visible<#2->{#3}}{\only<#2->{#3}}%
		\else
			#3%
		\fi
	\fi
}

% \deckrevealfinal[step]{content} - reveal step for the instructor deck only,
% everything shown at once in every other deck version (added 2026-09-19,
% direct instruction: "make the student versions show the final slide
% version immediately"). Unlike \deckreveal, the during-course deck gets the
% content too (not omitted), and like it no overlay spec is emitted outside
% the instructor version, so a non-instructor compile has no extra pages.
% For plain \pause steps use \conditionalpause[instructor] instead.
\NewDocumentCommand{\deckrevealfinal}{s O{2} m}{%
	\ifSMInstructorVersion
		\IfBooleanTF{#1}{\visible<#2->{#3}}{\only<#2->{#3}}%
	\else
		#3%
	\fi
}

% Starred forms, \deckreveal*[step]{content} and \deckrevealfinal*[step]{content}
% (2026-09-30, direct instruction: systematic replacements for per-slide
% fixes): the instructor steps before `step` keep the content's space
% (\visible instead of \only), so nothing around it moves when it appears --
% the "the diagram/equation/text jumps between reveal steps" class of
% complaint. Replaces the per-slide helpers that did the same by hand
% (\only<N->{x}\only<1..N-1>{\phantom{x}}; vxd's \vxdRevealB/C/D, c4c's
% \cRevealVisible). Other versions behave exactly as the unstarred forms.

% \deckitem<spec> - a list item whose overlay spec only applies in the
% instructor version; every other version gets a plain \item (the
% "student versions show the final state" rule). Replaces the pattern
%   \ifSMInstructorVersion \item<3-> text \else \item text \fi
% which repeated each item's text twice (42 items in 12 slides by
% 2026-09-30). An item label may follow: \deckitem<2->[\SMBulletEllipsis].
% Without a <spec> it is a plain \item.
\newcommand{\deckitem}{\@ifnextchar<{\SM@deckitem}{\item}}
\def\SM@deckitem<#1>{%
	\ifSMInstructorVersion
		\def\SM@next{\item<#1>}%
	\else
		\let\SM@next\item
	\fi
	\SM@next
}

% \conditionalslide[versions]{ID-slug} - \slide{}, but only included in the
% listed deck versions; keeps the frame counter in step with the omitted
% version otherwise (same \stepcounter{framenumber} technique \omitslide{}
% already uses for the student deck), so "see slide N" still points at the
% same physical slide across every deck version.
% Usage: \conditionalslide[instructor,after-course]{xyz-Answer-key}
\newcommand{\conditionalslide}[2][instructor,during-course,after-course]{%
	\ifdeckversioninlist{\SMDeckVersion}{#1}%
		{\global\SMSlideVersionSetfalse
		 \setSlideID{#2}%
		 \setSlideLanguageFromSlug{#2}%
		 \input{slides/#2-current.tex}}%
		{\stepcounter{framenumber}}%
}

% \slidestep{ID-slug}[versions]{steps} - \conditionalslide{} that also restricts
% which reveal (overlay) steps of the slide are typeset (2026-09-20, direct
% instruction). One slide file with several reveal steps can then appear at
% several places of a deck, each time showing only the steps that belong
% there, and a different set of steps per deck version:
%   \slidestep{egk-EN-Structure-of-a-Scientific-Paper-Overview}[instructor]{1,2}
%   \slidestep{egk-EN-Structure-of-a-Scientific-Paper-Overview}[during-course,after-course]{1}
% [versions] defaults to instructor,during-course,after-course (same as
% \conditionalslide; a presenter compile runs as version "instructor", so it
% always shows what the instructor deck shows); {steps} is a beamer overlay
% specification without the angle brackets ("1,2", "3", "2-", "1-2,4").
% The slide file itself is written once, with plain overlay specifications
% (\textcolor<2>{...}{...}, \only<3>{...}) and no version conditionals: which
% steps end up in the PDF is decided here, by handing the specification to
% beamer's own frame-level overlay restriction, \begin{frame}<steps>. Every
% frame environment opens its frame through \SMOpenFrame below, which
% inserts it.
%
% - Versions not in the list: the call does nothing at all -- not even a
%   \stepcounter{framenumber}, unlike \conditionalslide{}. Two \slidestep calls
%   at one position (say [instructor]{1,2} and [during-course,after-course]{1})
%   are alternatives, so every deck version has exactly one frame there;
%   stepping the counter for the omitted one would shift every later slide
%   number by one.
% - Jump target: the first \slidestep{} of a given id-slug in the document
%   defines the id's \hypertarget (see \setSlideID); later ones do not, so
%   there is no "multiply defined" warning and \hyperlink{id}{...} lands on the
%   first place the slide is shown.
% - \SMSlideSteps is cleared again right after the slide file has been
%   read, so it can never leak into a following plain \slide{}.
\gdef\SMSlideSteps{}
% \begin{frame}<steps>[options]{title}: the <...> has to be the very next
% token after "\begin{frame}" (beamer looks for it with \@ifnextchar), so the
% whole opening is built by \edef and only then executed. #1 is everything that
% follows "\begin{frame}" in the caller ("[#1]{#2}", "[plain,t]", ...).
% \SMFinalStateForStudents -- put it in a slide file, before the frame environment: in the
% student deck versions (everything but "instructor") that frame then typesets only its last
% overlay page, i.e. the final state of all its reveal steps (2026-09-21, direct instruction:
% "student decks should only show the final state"). Instructor and presenter compiles are
% unaffected. Unlike \conditionalpause/\deckrevealfinal, which each turn one reveal
% construct off, this needs no change to the slide body, so it also covers overlay specs
% like \only<2>, \item<3->, \uncover, \textcolor<2>, \alt and reveal steps that hide things
% again (the final state is what is left at the last step).
% How: beamer's frame-level overlay restriction \begin{frame}<N> (see \slidestep) with N =
% the frame's overlay count, which a frame only knows after it was typeset once. Every
% student compile therefore records N per frame (keyed by beamer's frame number, which is
% stable) in \jobname.ovl and reads it back at the start of the next run; the first run of
% a fresh job still typesets all pages, latexmk's rerun then restricts them. Beamer reports
% the same N for a restricted frame, so the value is stable.
\newif\ifSMFinalOnly
\newcommand{\SMFinalStateForStudents}{\global\SMFinalOnlytrue}
\makeatletter
\newwrite\SMOverlayFile
\newcommand{\SMSetOverlays}[2]{\global\@namedef{SMOverlays@#1}{#2}}
% Build files written before v5.3.0 (2026-10-01: this framework's \gotham...
% names became \SM...) still use the old names: .ovl (\gotham...SetOverlays),
% .top (...SetTopCorrection). The old names are assembled from \SMOldPrefix so a
% rerun of slidemage.sm_rename leaves them alone. The next compile rewrites
% the files with the new names.
\def\SMOldPrefix{gotham}
\expandafter\def\csname\SMOldPrefix SetOverlays\endcsname{\SMSetOverlays}
\expandafter\def\csname\SMOldPrefix SetTopCorrection\endcsname{\SMSetTopCorrection}
% A .aux written before v9.1.0 holds each subsection title (\SMAgendaSubsection)
% as `\gotham@subsectiontitleformat {Title}`: the Gotham theme wrapped titles
% in that no-op switch, which the SlideMage theme no longer has.
\expandafter\let\csname\SMOldPrefix @subsectiontitleformat\endcsname\@empty
\newcommand{\SMRecordOverlays}{%
	\ifSMInstructorVersion\else
		\immediate\write\SMOverlayFile{\string\SMSetOverlays{\the\c@framenumber}{\the\beamer@slideinframe}}%
	\fi
}
\apptocmd{\beamer@writeslidentry}{\SMRecordOverlays}{}{}
% "No file" is the line LaTeX's own \@input prints for a missing .aux or .nav: latexmk reads it and then knows the
% file as an input from the first run, so a second run that reads an unchanged one asks for no third.
\AtBeginDocument{%
	\InputIfFileExists{\jobname.ovl}{}{\typeout{No file \jobname.ovl.}}%
	\immediate\openout\SMOverlayFile=\jobname.ovl\relax
}
% Expands to <N> when the frame that is about to open is flagged and N is known.
\newcommand{\SMFinalSpec}{%
	\ifSMFinalOnly
		\ifSMInstructorVersion\else
			\ifcsname SMOverlays@\the\numexpr\c@framenumber+1\relax\endcsname
				<\csname SMOverlays@\the\numexpr\c@framenumber+1\relax\endcsname>%
			\fi
		\fi
	\fi
}
\newcommand{\SMOpenFrame}[1]{%
	\edef\SMOpenFrame@cmd{%
		\noexpand\begin{frame}%
		\ifx\SMSlideSteps\empty\SMFinalSpec\else<\SMSlideSteps>\fi
		\unexpanded{#1}}%
	\global\SMFinalOnlyfalse
	\SMOpenFrame@cmd
}
\NewDocumentCommand{\slidestep}{m O{instructor,during-course,after-course} m}{%
	\ifdeckversioninlist{\SMDeckVersion}{#2}%
		{\ifcsname SMSlideStepSeen@#1\endcsname
			\global\SMSuppressSlideTargettrue
		 \else
			\global\SMSuppressSlideTargetfalse
			\expandafter\gdef\csname SMSlideStepSeen@#1\endcsname{}%
		 \fi
		 \gdef\SMSlideSteps{#3}%
		 \global\SMSlideVersionSetfalse
		 \setSlideID{#1}%
		 \global\SMSuppressSlideTargetfalse
		 \setSlideLanguageFromSlug{#1}%
		 \input{slides/#1-current.tex}%
		 \gdef\SMSlideSteps{}}%
		{}%
}

\makeatother

% \SMUnderbrace{math}{label} - seamless TikZ-drawn brace, replacing plain
% amsmath \underbrace (2026-07-23, direct instruction, after presenting
% options on slide ebq -- see SESSION_HANDOFF.md's own entry for that task):
% amsmath assembles \underbrace's glyph from separate piecewise brace-half
% characters, which visibly shows a seam/notch at narrow widths (ebq's own
% "risk-free rate" term is narrow enough to show it plainly); this instead
% draws one continuous vector brace path via TikZ's `decorate` +
% `decoration={brace}` (decorations.pathreplacing, loaded above), which has
% no such seam at any width. #1 is typeset in math mode automatically (do
% not wrap the call itself in $...$); #2 is the label, set in upright text
% (font=\small, align=center so a multi-line label via \\ stays centered
% under the brace, e.g. two-line labels converted from \substack{}). Each
% call gets its own auto-numbered internal node name (SMUnderbraceCounter)
% so multiple calls can coexist on one slide without a caller-supplied name
% argument colliding (unlike the original \tikzbrace{}{}{} demo macro this
% replaces, which took the node name as a third argument).
\newcounter{SMUnderbraceCounter}
\newcommand{\SMUnderbrace}[2]{%
	\stepcounter{SMUnderbraceCounter}%
	\tikz[baseline=(SMUnderbraceNode\theSMUnderbraceCounter.base)]{
		\node[inner sep=0pt] (SMUnderbraceNode\theSMUnderbraceCounter) {$#1$};
		\draw[decorate,decoration={brace,amplitude=4pt,mirror,raise=1pt}]
			(SMUnderbraceNode\theSMUnderbraceCounter.south west) -- (SMUnderbraceNode\theSMUnderbraceCounter.south east)
			node[midway,below=6pt,font=\small,align=center] {#2};
	}
}

% PDF accessibility: mark purely decorative/running content (footer, below;
% also reused on titleframe's/quoteframe's own decorative background boxes,
% and -- since 2026-08-14, via the \includegraphics wrapper further down --
% on every image the alt-text manifest classifies as decorative, the
% institute logo among them) as a PDF "Artifact", so screen readers skip it
% -- via
% LuaTeX's raw \pdfextension literal
% primitive. \ifluatex-guarded (iftex is already loaded, via beamer.cls
% itself) because \pdfextension is LuaTeX-only -- this project's README
% documents xelatex as a supported alternative engine, and compiling with it
% hit exactly this: "Undefined control sequence: \pdfextension", hard-failing
% every frame, since the first version of this didn't guard the primitive at
% all. Under any non-LuaTeX engine these two commands now silently do
% nothing instead: decorative content just isn't marked as an Artifact under
% xelatex/pdflatex, rather than breaking the compile -- a real loss of
% accessibility under those engines specifically, not a full equivalent fix,
% but not something to solve today (see also: compiling with plain pdflatex
% already silently mis-renders the whole deck's fonts, so xelatex is really
% the only *other* engine this needs to not break under; pdflatex was never
% viable here to begin with, per "Compiling" in README.md).
% Deliberately NOT tagpdf/kernel tagging (\DocumentMetadata):
% beamer 3.77 refuses to compile at all when that's active ("The beamer
% class is not compatible with \DocumentMetadata ... switch to the new
% ltx-talk class if you require tagged output"), and ltx-talk is a from-
% scratch class (not beamer-based, no theme/template system) that this
% framework can't realistically migrate to. Bypassing beamer's check by
% force-loading the underlying package before \documentclass also failed --
% it gets past hyperref but then hits a genuine internal error in tagpdf's
% LuaTeX backend from the non-standard load order, not something to build
% real coursework on.
%
% /Artifact is a standalone PDF marked-content tag (PDF spec 14.8.2.2):
% valid in *any* PDF, with or without a structure tree, no /Properties
% resource entry or MCID required -- unlike structure-tree tagging, this
% doesn't require kernel tagging support at all. It's a real, working
% partial fix, not full PDF/UA conformance: there's no structure tree, so
% no true semantic headings, and reading order elsewhere on the slide still
% depends on plain content-stream order (verified, not enforced, for
% standardframe -- see its own comment).
% #1 is the artifact's inline PDF dictionary content, e.g.
% "/Type /Pagination /Subtype /Footer" for a running footer, or
% "/Type /Background" for a decorative fill -- not hardcoded to one /Type,
% since not everything decorative here is "Pagination" in the PDF spec's own
% sense of that word (running headers/footers/page numbers/watermarks);
% titleframe's/quoteframe's plain background color fills are semantically
% "Background" artifacts instead, a distinct standard type.
\newcommand{\SMArtifactBegin}[1]{%
	\ifluatex\pdfextension literal direct{/Artifact <<#1>> BDC}\fi
}
\newcommand{\SMArtifactEnd}{\ifluatex\pdfextension literal direct{EMC}\fi}

% Per-frame "Image source(s):" auto-footnote (ALT_TEXT_NOTES.md phase 6 step
% 3). The \includegraphics wrapper further down calls
% slidemage_alt.recordSource() for every image (decorative or informative --
% source is independent of role, same reasoning as media/alt-text.json's own
% ImageEntry.source), accumulating each frame's distinct, order-of-first-
% -appearance source strings in Lua. These two commands are standardframe's
% own hooks into that: reset at frame open (same per-overlay-replay reset
% \SMClearSideFootnotesBox already does), emit at frame close.
%
% Defined here (outside/before the "Alt-text manifest" \ifluatex block
% further down that actually loads SlideMage_alt.lua and redefines
% \includegraphics) rather than inside it, matching \SMArtifactBegin/End
% just above: standardframe unconditionally calls both, so they need to
% exist regardless of engine, each independently \ifluatex-guarded inside
% -- under any non-LuaTeX engine \SMResetImageSources/
% \SMEmitImageSourceFootnote are both no-ops, same degradation as the
% rest of this accessibility system under a non-LuaTeX engine. Safe to
% define before SlideMage_alt.lua is actually loaded (\AtBeginDocument,
% further down): these are only ever *called* from inside a real frame,
% well after \begin{document} has fired.
\newif\ifSMImageSourcesHasText
\newcommand{\SMImageSourcesText}{}
\newcommand{\SMResetImageSources}{%
	\ifluatex\directlua{slidemage_alt.resetFrameSources()}\fi
}
% \source{} is itself implemented as an empty-marker \footnote{} (see its
% own definition further below) -- calling it here gives the auto-footnote
% the exact same rendering (font, spacing, right-alignment, stretch-to-
% -sit-above-the-footer, side-footnote redirection under \sidefootnotes) as
% a hand-authored \source{} line, for free, rather than reimplementing any
% of that. Coexists with a real hand-authored \source{} call elsewhere in
% the same frame -- both are just separate unmarked footnote-styled lines,
% stacking the same way multiple real \footnote{}s already do.
% \SMImageSourceLabel -- the "Image source(s):"/"Bildquelle(n):" label,
% driven by \ifSMGerman (set from the compiling deck's own
% \decklanguage) rather than hardcoded, so a German deck's auto-footnote
% reads naturally instead of showing an English label on German content.
\newcommand{\SMImageSourceLabel}{\ifSMGerman Bildquelle(n)\else Image source(s)\fi}
\newcommand{\SMEmitImageSourceFootnote}{%
	\ifluatex
	\SMImageSourcesHasTextfalse
	\directlua{slidemage_alt.emitSourceFootnote()}%
	\ifSMImageSourcesHasText
		\source{\SMImageSourceLabel: \SMImageSourcesText}%
	\fi
	\fi
}

% Invisible slide-id stamp for footer-less frame types. \titleframe and
% \fullbleedframe (both defined further down) render via beamer's "plain"
% frame option, which suppresses the footline entirely -- so the visible
% footer's own "(slide id: ...)" text (see the footer template below) never
% reaches those pages. SlideMage's own compile-time preview snapshot
% (slide_snapshot.py) identifies which slide a compiled page belongs to
% purely by scanning its extracted text for that "slide id: XXX" string, so
% a plain frame with none of that text never gets a snapshot captured at
% all (caught 2026-08-15: titleframe's own title slide and 25 straight
% fullbleedframe "Excursus" slides in the IA2 deck all had no preview).
% This stamps the same marker, invisibly (PDF text-render mode 3, the same
% technique \SMAltInvisible below uses for image alt text) and wrapped
% as a /Pagination/Footer artifact, same as the real footer's own slide-id
% text -- so it reaches neither a sighted reader nor a screen reader,
% purely a technical marker for the compile pipeline. Only \slideID itself
% is stamped, not the omitted-slide/version suffix the visible
% footer also carries: nothing here is ever read by a human, and
% slide_snapshot.py's own regex only looks for the id.
%
% \vbox to 0pt{\vss ...}: same reasoning as every other absolute/zero-space
% overlay in this file (see \titleframe's own comment on this) -- \sbox
% builds an hbox, and a bare hbox placed directly in a frame's vertical
% list silently opens an implicit paragraph, switching following material
% into horizontal mode. Wrapping in a vbox contains that entirely within
% its own internal vertical mode instead of leaking out to the frame body.
\newsavebox{\SMSlideIDStampBox}
\newcommand{\SMInvisibleSlideIDStamp}{%
	\ifluatex
	\vbox to 0pt{\vss
	\begingroup
	\hbadness=10000 \hfuzz=\maxdimen
	\sbox{\SMSlideIDStampBox}{%
		\parbox[t]{\textwidth}{%
			\tiny\raggedright
			\SMArtifactBegin{/Type /Pagination /Subtype /Footer}%
			\pdfextension literal direct{3 Tr}%
			% Trailing \space before the closing paren, not just "\slideID)":
			% slide_snapshot.py's own regex ("slide id:\s*(\S+)") reads
			% everything up to the next whitespace as the id, and without a
			% space there it would capture "gcs)" instead of "gcs" (caught
			% by testing the actual compiled+extracted text, not assumed --
			% the real footer's own "... \slideID ... | ..." format always
			% has a space after \slideID for the same reason).
			(slide id: \slideID\space)%
			\pdfextension literal direct{0 Tr}%
			\SMArtifactEnd
		}%
	}%
	\wd\SMSlideIDStampBox=0pt
	\ht\SMSlideIDStampBox=0pt
	\dp\SMSlideIDStampBox=0pt
	\usebox{\SMSlideIDStampBox}%
	\endgroup
	}%
	\fi
}

% ---------------------------------------------------------------------------
% Alt-text manifest: automatic /Artifact marking of decorative images
% ---------------------------------------------------------------------------
% See ALT_TEXT_NOTES.md for the whole design; this is phase 1 of it. Every
% image in the project has a row in media/alt-text.csv saying whether it is
% "decorative" (carries no information beyond the slide text -- Midjourney
% illustrations, icons, chapter logos, institute logos) or "informative" (a
% chart/screenshot/photo a reader actually needs described). Decorative ones
% get wrapped in /Artifact here, so a screen reader skips them instead of
% announcing a bare filename. Informative ones are left alone for now; their
% descriptions are authored per slide and emitted later (phases 4 and 5).
%
% Why wrap \includegraphics rather than change the 265 call sites: the call
% sites don't need to change at all this way, including the ones inside tikz
% \node{}s and the ones that pass the filename through a macro argument
% (\SMInstituteLogoPath, or a #1 inside another macro's definition) --
% those only know their filename at expansion time, which a source-level
% rewrite couldn't have handled anyway. An image with no manifest row simply
% renders exactly as before, so a missing or unreadable CSV degrades to
% today's behavior rather than breaking every image in the deck.
%
% LuaTeX-only, same \ifluatex reasoning as \SMArtifactBegin above -- and
% here the guard covers the whole redefinition, not just the marking, since
% the lookup itself is Lua. Under any other engine \includegraphics is left
% completely untouched.
\ifluatex
	% The Lua itself lives in SlideMage_alt.lua rather than inline here --
	% see that file's own header for why (short version: TeX tokenizes a
	% \directlua argument first, so every "%" in a Lua string pattern
	% silently comments out the rest of its line, and the first version of
	% this failed exactly that way). This loader snippet is deliberately
	% free of "%" and "#" so it is safe to keep inline. Found via kpathsea
	% (TEXINPUTS), since it lives in latex/, outside the project.
	\SMLuaOnce{
		local path = kpse.find_file("SlideMage_alt.lua", "lua")
			or kpse.find_file("SlideMage_alt.lua", "tex")
		if path then
			dofile(path)
		end
	}

	\newif\ifSMAltDecorative
	\newif\ifSMAltHasText
	\newsavebox{\SMAltTextBox}
	\newcommand{\SMAltShort}{}

	% Invisible alt text (ALT_TEXT_NOTES.md phase 5). An untagged PDF is read
	% by *text extraction*, so text placed in the content stream at the image's
	% position is what a screen reader encounters there -- text rendering mode
	% 3 ("neither fill nor stroke", PDF spec 9.3.6) makes it invisible on the
	% page while leaving it fully extractable.
	%
	% Two deliberate properties:
	%
	% - The box's three dimensions are zeroed after typesetting, so the text
	%   contributes nothing whatsoever to layout. Setting \wd/\ht/\dp of an
	%   already-built box is the only way to be sure of that: \rlap/\smash
	%   each zero only some of them, and a stray descender would shift a
	%   hand-positioned slide.
	% - \SMAltImagePrefix ("Image: "/"Bild: ") prefixes the description so a listener can tell it apart
	%   from the slide's own prose, which they otherwise cannot.
	%
	% Known and accepted side effect: this text is real text, so it appears in
	% copy-paste and in Ctrl-F. That was the explicit tradeoff for reaching a
	% screen reader at all without a structure tree, which beamer cannot give
	% us (see README.md's accessibility section).
	% \SMAltImagePrefix -- language-conditional counterpart to
	% \SMImageSourceLabel above, for this invisible layer's own "Image: "/
	% "Bild: " prefix.
	\newcommand{\SMAltImagePrefix}{\ifSMGerman Bild: \else Image: \fi}
	\newcommand{\SMAltInvisible}[1]{%
		\begingroup
		% A long alt string in an unbreakable hbox would otherwise flood the
		% log with overfull-box warnings on every image.
		\hbadness=10000 \hfuzz=\maxdimen
		% Wrapped in a \parbox, and set \tiny, so the text stays *inside the
		% page*. Neither is cosmetic: an unbreakable line starting at the
		% image's own x-position runs past the right page edge after about 80
		% characters, and every PDF text extractor tested (pdftotext with and
		% without -layout, PyMuPDF) silently drops whatever falls outside the
		% page rectangle -- so a longer alt was being truncated mid-sentence
		% with no error anywhere. \tiny is free here (the text is invisible)
		% and fits several times more characters per line, which keeps even a
		% long description to one or two lines.
		\sbox{\SMAltTextBox}{%
			\parbox[t]{\textwidth}{%
				\tiny\raggedright
				\pdfextension literal direct{3 Tr}%
				\SMAltImagePrefix#1%
				\pdfextension literal direct{0 Tr}%
			}%
		}%
		\wd\SMAltTextBox=0pt
		\ht\SMAltTextBox=0pt
		\dp\SMAltTextBox=0pt
		\usebox{\SMAltTextBox}%
		\endgroup
	}

	% Deferred to \AtBeginDocument so it cannot possibly run before beamer has
	% loaded graphicx, and so a deck that loads further graphics packages of
	% its own still ends up with this wrapper outermost.
	\AtBeginDocument{%
		\directlua{
			local loaded = slidemage_alt.load()
			if loaded == nil then
				texio.write_nl("slidemage_alt: media/alt-text.csv not found -- "
					.. "decorative images will not be marked as artifacts.")
			end
		}%
		% \NewCommandCopy (LaTeX 2020-10-01 and later) rather than \let: it
		% copies a robust command's full definition correctly, which \let
		% does not for commands whose top level is a \protected wrapper.
		\NewCommandCopy{\SMOrigIncludegraphics}{\includegraphics}%
		% Signature "s o o m": the starred form and the legacy two-optional
		% bounding-box form aren't used anywhere in this project today
		% (checked, 2026-08-14: zero occurrences of either), but graphicx
		% accepts them, so the wrapper accepts and forwards them rather than
		% silently breaking the day one gets used. "o" (not "O{}") because
		% only -NoValue- distinguishes an absent optional argument from an
		% empty one, and graphicx reads [a][b] and [a] as different things.
		\RenewDocumentCommand{\includegraphics}{s o o m}{%
			\begingroup
			% \edef so a filename reaching this through a macro
			% (\SMInstituteLogoPath, or a #1 in another macro's body) is
			% resolved to the actual path before the lookup. Every path-valued
			% macro in this project is a plain expandable string; a fragile one
			% would need \protected@edef here instead.
			\edef\SMAltReference{#4}%
			\SMAltDecorativefalse
			\directlua{slidemage_alt.mark("\luaescapestring{\SMAltReference}")}%
			% Unconditional -- unlike SMAltInvisible below (informative
			% images only), a source is recorded regardless of role.
			\directlua{slidemage_alt.recordSource("\luaescapestring{\SMAltReference}", "\luaescapestring{\SMDeckLanguage}")}%
			% /Type /Layout is the PDF spec's artifact type for cosmetic design
			% elements, which is what a decorative illustration is. Not
			% /Background: that one means a fill or image behind the whole
			% page, and is already used for exactly that by titleframe and
			% quoteframe above.
			\ifSMAltDecorative\SMArtifactBegin{/Type /Layout}\fi
			% Emitted *before* the image, so the description precedes it in the
			% content stream and therefore in extraction order -- a listener
			% hears what the image is before reaching it, rather than after.
			% slidemage_alt.altText prints nothing for a decorative image or one
			% with no short alt yet, so this line is inert until the manifest
			% is filled in.
			\SMAltHasTextfalse
			\ifSMAltDecorative\else
				\directlua{slidemage_alt.altText("\luaescapestring{\SMAltReference}", "\luaescapestring{\SMDeckLanguage}")}%
				\ifSMAltHasText\SMAltInvisible{\SMAltShort}\fi
			\fi
			\IfBooleanTF{#1}%
				{\IfNoValueTF{#2}%
					{\SMOrigIncludegraphics*{#4}}%
					{\IfNoValueTF{#3}%
						{\SMOrigIncludegraphics*[#2]{#4}}%
						{\SMOrigIncludegraphics*[#2][#3]{#4}}}}%
				{\IfNoValueTF{#2}%
					{\SMOrigIncludegraphics{#4}}%
					{\IfNoValueTF{#3}%
						{\SMOrigIncludegraphics[#2]{#4}}%
						{\SMOrigIncludegraphics[#2][#3]{#4}}}}%
			\ifSMAltDecorative\SMArtifactEnd\fi
			\endgroup
		}%
	}
\fi

% \tikzillustration{ID}: extends the alt-text system above -- which only
% ever sees a real \includegraphics file -- to a hand-drawn tikz diagram
% with no file behind it at all (ALT_TEXT_NOTES.md). #2 is the
% illustration's own opaque, SlideMage-assigned ID (never hand-invented
% -- same "permanent ID" convention slides themselves use, so a description
% survives edits to the drawing itself); #1 (optional) is passed straight
% through as the inner tikzpicture's own options. The body between
% \begin/\end is the tikzpicture's content, completely unmodified.
%
% Role (decorative/informative) is looked up exactly the way an image's is
% -- slidemage_alt.mark(), the same Lua call \includegraphics's own wrapper
% above makes, against the same media/alt-text.csv -- rather than a second
% bracket argument SlideMage would otherwise have to text-edit into this
% call in place every time the role changes. SlideMage writes this
% illustration's own row into media/alt-text.json (and its compile-time
% media/alt-text.csv projection) keyed "tikz:ID", never a real media/ path;
% slidemage_alt.mark() itself needed no changes, since it already keys purely
% on whatever string it is given (SlideMage/alt_images.py).
%
% A decorative illustration is wrapped in the same \SMArtifactBegin{/Type
% /Background}/\End pair sideimageframe's/fullbleedframe's own decorative
% filler content already uses, so a screen reader skips it entirely; an
% informative one is left completely unmarked, same as real content in
% figureframe/fullbleedframe. Unlike an informative *image*, it gets no
% invisible-text alt layer (phase 5) of its own -- that layer's positioning
% is built around \includegraphics's own bounding box, which a freeform
% tikzpicture doesn't have the same way. The per-deck accessible transcript
% (README.md) is this illustration's own route to a screen reader instead,
% the same status phase 5 already carries for every image.
%
% \ifluatex-guarded exactly like \SMArtifactBegin/\SMArtifactEnd
% above (both \SMAltDecorativefalse/\ifSMAltDecorative only exist
% inside that same guard) -- under any other engine this degrades to a
% plain, unmarked tikzpicture, the same "no accessibility marking, but
% still compiles and renders" degradation the rest of this system already
% has. Safe even though those two control sequences are otherwise undefined
% under a non-LuaTeX engine: TeX's own \iffalse-branch skipping never looks
% up a token's meaning, so an unrecognized name inside the untaken
% \ifluatex branch is simply skipped over, never executed.
\newenvironment{tikzillustration}[2][]{%
	\ifluatex
		\SMAltDecorativefalse
		\directlua{slidemage_alt.mark("\luaescapestring{tikz:#2}")}%
		\ifSMAltDecorative\SMArtifactBegin{/Type /Background}\fi
	\fi
	\begin{tikzpicture}[#1]
}{%
	\end{tikzpicture}%
	\ifluatex\ifSMAltDecorative\SMArtifactEnd\fi\fi
}

% Footer
% The footer box on the footline, the footnote template with its symbols, \sidefootnotes and the
% strip they are drawn in are the theme's (beamerthemeSlideMage.sty); the footer text, which
% depends on the deck version, is the "frame footer" template below.
% Smaller footer + a bit of padding under it
\setbeamerfont{frame footer}{size=\SMFooterTextSize} % the catalog: {font size}{baseline skip} in pt

% Shrinks the gap between footnotes and the footer box below them. The
% theme positions the footer's content via \vspace*{\SMFootlineVOffset}
% (a fixed, negative length) before it starts; making that more negative
% turns out to give footnotes (which are placed just before the footer, and
% stretch to fill whatever room remains via beamer's own
% \beamer@framebottomskip) more room to extend into, pulling them down
% towards the footer -- checked empirically against the PDF's own
% coordinates, this does *not* move the footer's own text at all (that
% still sits exactly where \SMFooterBoxHeight below was calibrated for),
% so the footer box itself needs no compensating adjustment.
%
% Partly given back (+1.4pt) 2026-07-05: a footnote/source containing an
% underlined \href (see the hyperlink-styling comment above) sits close
% enough to the footer at the original -1.3pt that the underline visually
% touches the footer box (checked directly against the compiled PDF's own
% pixel coordinates: a plain-text footnote's own gap to the footer measured
% ~1.9pt, but an underlined-link footnote's gap measured only ~0.6pt at
% the same setting, since the strut rebuilt for footnotes just below --
% correctly sized for plain text's own natural depth -- doesn't account
% for the underline rule extending further down). +1.4pt splits the
% difference: a linked footnote's gap now measures close to the plain-
% footnote reference (~1.9pt) instead of colliding with the footer, and a
% plain footnote's own gap just becomes a little more relaxed (~3.3pt)
% rather than stretched uncomfortably tight -- applied globally rather
% than only to footnotes with links, since there's no clean per-footnote
% hook for it and the plain-footnote side effect is minor.
%
% Nudged back down by 0.4pt the next day (direct instruction: move
% footnotes down "by the width of the underline line") -- the +1.4pt above
% left a bit more headroom than necessary; 0.4pt is the underline rule's
% own measured stroke thickness (a solid ~1px core at 600dpi, i.e.
% 1/600in \approx 0.12pt, plus antialiasing bleed either side consistent
% with the common 0.4pt LaTeX default rule width -- measured directly
% against the compiled PDF's own pixel values, not assumed).
%
% Reverted the very next round (direct instruction: move back *up* by that
% same underline thickness) -- net effect is back to the plain +1.4pt this
% section started with; kept as two separate historical adjustments below
% rather than collapsing them into one, so this back-and-forth stays
% visible rather than silently disappearing.
% (v9.4.0: the adjustments are one value now, the default of the catalog's
% \SMFootlineVOffset in SlideMage_setup.tex: -0.5ex-1.3pt+1.4pt-0.4pt+0.4pt-0.87pt
% = -0.5ex-0.77pt.)

% 2026-07-06 (direct instruction): move footnotes down by half the
% measured distance between a linked footnote's own underline and the
% footer box, checked on slides/3ic-Corporate-Venture-Capital-
% current.tex specifically (the underline sits at y=245.95pt, the footer
% box's own top edge at y=247.69pt, a 1.74pt gap) -- half of that is
% 0.87pt, applied as a further negative adjustment (more negative =
% footnotes get pulled further down, per this section's own comment
% above).

% \SMFooterTextDrop (v9.8.0: derived; it was a number of the layout profile since v9.4.0, -0.49pt for
% Fira Sans, -0.37pt for Latin Modern Sans) centres the ink of the frame number in the footer box. The
% template below starts the footer line with \vspace*{drop}, which in the footline box's horizontal
% mode is a \vadjust below the line, as is the \vspace*{0.4ex} that ends it; the box has the declared
% depth \SMFootlineDepth, so the number's baseline lies
%     \SMFootlineDepth + (the depth of the line) + 0.4ex + drop
% above the page's bottom edge. The depth of the line is that of the slide id's strut
% (\SMFooterIDStrut makes it the same on every slide). The ink of the digits is what luaotfload
% keeps as the height and depth of a digit's box (the glyph's bounding box: 676 and 19 units in Fira
% Sans Light), so centring it in the footer box of height \SMFooterBoxHeight gives the drop.
\makeatletter
\newlength{\SM@footerlinedepth}
\newcommand{\SM@footertextdrop}{%
	\setbox\z@\hbox{\SMFooterIDSize\SMFooterIDStrut}\setlength{\SM@footerlinedepth}{\dp\z@}%
	\setbox\z@\hbox{0123456789}%
	\setlength{\SMFooterTextDrop}{\dimexpr(\SMFooterBoxHeight-\ht\z@+\dp\z@)/2-\SMFootlineDepth-\SM@footerlinedepth-0.4ex\relax}}
\makeatother
%
% \SMFooterRightShift (v9.7.0: derived; it was a number of the layout profile, 4.6pt for Fira Sans,
% -0.33pt for Latin Modern Sans) puts the frame number's right edge on the logo's right edge above
% it, as the footnotes' lines end there. The footer line is typeset in the footline box, which is
% \hsize wide and ends its lines \rightskip before that; after the number come, in this order,
% \hspace*{-4em} and the shift (set here), then the box's \SMRightFooterPadding and its
% \SMFootlineHRightOffset. So the number's right edge is
%     \hsize - \rightskip - \SMRightFooterPadding - \SMFootlineHRightOffset + 4em + shift
% and the shift that puts it at \paperwidth - \SMFrametitleRightPadding follows. Every term is the
% one in force where the footer is typeset (the \rightskip is 3ex of the footline font, the 4em is
% the footer's), which is why the shift is computed in the frame footer template and not once: the
% two fonts differ in size, and a profile with another x-height has other paddings in ex.

% The footer no longer shows a compile date/time (removed in v8.25.0: a
% changing stamp made every wrapper and PDF differ between compiles). The
% two names below remain as harmless no-ops so decks that still call
% \SMShowCompileTimestampfalse/true or print \SMCompileTimestamp compile.
\newif\ifSMShowCompileTimestamp
\newcommand{\SMCompileTimestamp}{}

% \SMPresenterModeOverride: SlideMage's own hook for a "Presenter"
% compile -- a separate output alongside the plain instructor PDF, never a
% toggle on it (see SlideMage/compiler.py's own CompileJob.presenter
% field and SlideMage/qt_deck_table.py's own "Presenter" checkbox,
% independent of the three version checkboxes the same way "After-course,
% per chapter" already is). Same \ifdefined-hook pattern as
% \deckversionoverride: compiler.py's wrapper file
% \def's it, before \input{}ing the real deck, whenever that checkbox is
% ticked for a job.
%
% Turns on beamer's own built-in dual-screen presenter mechanism:
% \note{...} content (see "Speaker notes" below -- already authored
% throughout this project's own slides, previously never actually emitted
% anywhere at all) becomes a facing "notes page" after every real slide,
% which presenter software (pympress is this project's own recommended
% choice -- see README.md's "Presenting" section) auto-detects and splits
% across the projector (plain slides only) and the presenter's own screen
% (current slide, next slide, and the note text). \usepackage{pgfpages} is
% loaded explicitly rather than trusted to \setbeameroption's own
% auto-load, since which beamer versions actually do that on their own
% isn't worth depending on.
%
% A presenter compile's own CompileJob always carries version="instructor"
% (or the deck's own hardcoded content, for a deck with no
% \deckversionoverride at all) regardless of whether the plain "Instructor"
% checkbox is *also* ticked in the same batch -- notes and full \pause
% reveals only make sense against the instructor's own content, so
% Presenter never needs (or offers) During-course/After-course versions of
% itself. slide_snapshot.py's own compile-time preview refresh is
% deliberately skipped for a presenter job (compiler.py's own
% _relocate_pdf_and_finish()) -- its page count and per-page layout are
% nothing like a normal compile's, so reading it the normal way would
% silently corrupt the real preview cache with wrong or missing pages.
\ifdefined\SMPresenterModeOverride
	\usepackage{pgfpages}
	\setbeameroption{show notes on second screen=right}
	% Presenter-mode note page (direct instruction, 2026-09-10): beamer's
	% own default "note page" template (beamerouterthemedefault.sty,
	% \defbeamertemplate*{note page}{default}) reserves its top quarter for
	% a thumbnail of the current slide (\insertslideintonotes) plus a
	% rotated compile-date printed beside it, then gives the title band
	% (short title + section/subsection/frametitle breadcrumb) the same
	% 0.25\paperheight, leaving only the bottom 0.75 for \insertnote. Both
	% the thumbnail and date are redundant here -- pympress (this project's
	% own recommended presenter software, see README.md) already shows the
	% current and next slide on its own -- so this redefinition drops that
	% whole top block (thumbnail + date) outright and shrinks the title
	% band from 0.25\paperheight to 0.08\paperheight (still one line for
	% the short title plus room for the breadcrumb, per beamer's own
	% layout below, copied unchanged apart from widening each line from
	% beamer's own 0.71-0.75\textwidth, sized to leave room for the now-
	% removed thumbnail, to 0.85-0.95\textwidth), giving \insertnote the
	% remaining ~0.92\paperheight instead of 0.75. No color-theme "note
	% title"/"note page" background is set anywhere in this project (a
	% plain grep confirms it), so beamer's own \ifbeamercolorempty guards
	% around both of those bands are dropped too -- they were already
	% pure no-ops.
	%
	% Registered as a new named variant ("SMPresenter"), not a
	% redefinition of beamer's own "default" -- confirmed directly that
	% \defbeamertemplate*{note page}{default}{...} here errors ("Command
	% \beamer@@tmpop@note page@default already defined"): both the
	% starred and unstarred \defbeamertemplate forms register a variant
	% name via a plain \newcommand internally (beamerbasetemplates.sty),
	% which always errors on a name beamer's own outer theme already
	% registered, star or not. Same pattern this file's own footnote
	% template already uses just above (\defbeamertemplate{footnote}
	% {SMFootnote}{...} then \setbeamertemplate{footnote}[SMFootnote]) --
	% define-then-select instead of overwrite-in-place.
	%
	% \makeatletter/\makeatother wraps this specific block (not the whole
	% \ifdefined) because \beamer@tempbox below is an @-command: this
	% point in the preamble sits between two of this file's own
	% \makeatletter/\makeatother pairs (the previous one closes well
	% above, the next opens well below), so @ is still catcode-other here
	% otherwise -- confirmed directly rather than assumed.
	\makeatletter
	\defbeamertemplate{note page}{SMPresenter}{%
		{%
			\scriptsize
			% The breadcrumb typesets the frame title (twice: once to measure, once
			% to show it), and a \footnote{}/\source{} in a frame title would run each
			% time: the slide's last reveal step then carried its footnote three
			% times (direct report 2026-10-13, mbz). Footnotes are switched off in
			% this group; the slide itself keeps its own, single one.
			\renewcommand*{\footnote}[2][]{}%
			\renewcommand*{\footnotemark}[1][]{}%
			\usebeamerfont{note title}\usebeamercolor[fg]{note title}%
			% Natural-height \vbox (no "to <len>"/\vfil): this project's own
			% real decks call \section{}/subsection ones sometimes not (a
			% plain grep of decks/*.tex confirms both), so a fixed-height
			% band sized for the *shortest* case overlaps \insertnote below
			% on any slide needing all three breadcrumb lines -- confirmed
			% directly: a first attempt at a fixed 0.08\paperheight did
			% exactly that against a real slide's own \note{}. A natural-
			% height box is exactly as tall as whichever of these three
			% lines are non-empty for *this* slide, never more, so it can
			% never collide with the notes below it, and hands \insertnote
			% every bit of space a section/subsection-free slide doesn't
			% need for a breadcrumb. \insertshorttitle (beamer's own first
			% line here, normally the deck's own \title{}) is dropped
			% outright rather than shrunk -- confirmed via grep that
			% \title{} is never called anywhere in this project, so that
			% line has always been silently empty.
			\vbox{\vskip0.3em
				\setbox\beamer@tempbox=\hbox{\insertsection}%
				\hbox{\ifdim\wd\beamer@tempbox>1pt{\hskip4pt\raise3pt\hbox{\vrule
							width0.4pt height7pt\vrule width 9pt
							height0.4pt}}\hskip1pt\hbox{\begin{minipage}[t]{0.95\textwidth}\def\breakhere{}\insertsection\end{minipage}}\fi%
				}%
				\setbox\beamer@tempbox=\hbox{\insertsubsection}%
				\hbox{\ifdim\wd\beamer@tempbox>1pt{\hskip17.4pt\raise3pt\hbox{\vrule
							width0.4pt height7pt\vrule width 9pt
							height0.4pt}}\hskip1pt\hbox{\begin{minipage}[t]{0.95\textwidth}\def\breakhere{}\insertsubsection\end{minipage}}\fi%
				}%
				\setbox\beamer@tempbox=\hbox{\insertshortframetitle}%
				\hbox{\ifdim\wd\beamer@tempbox>1pt{\hskip30.8pt\raise3pt\hbox{\vrule
							width0.4pt height7pt\vrule width 9pt
							height0.4pt}}\hskip1pt\hbox{\insertshortframetitle[width=0.9\textwidth]}\fi%
				}%
				\vskip0.3em}%
		}%
		\nointerlineskip
		\insertnote
	}
	\setbeamertemplate{note page}[SMPresenter]
	\makeatother
\fi
% \SMHideFooterText: opt-in flag suppressing just \FooterText (the left
% part of the footer) -- default false. Set true by the standout template's
% own hook below, so `\begin{frame}[standout]` hides only that left part,
% keeping the "(slide id: ...)" block and frame number visible (2026-07-16,
% direct instruction). Reset immediately after being read, inside the
% footer template itself, rather than relying on a per-environment closing
% hook (standout frames are bare `\begin{frame}[standout]...\end{frame}`,
% not one of this project's own environments like standardframe, so there's
% no custom closing code to reset it in) -- safe because \SMHideFooterText
% is only ever consumed here, once, immediately before being cleared.
\newif\ifSMHideFooterText
% The theme's own \defbeamertemplate*{standout}{SMDefault}
% (beamerthemeSlideMage.sty) is what `\begin{frame}[standout]` actually
% activates; \addtobeamertemplate's "after" hook appends to it without
% needing to touch the theme.
\addtobeamertemplate{standout}{}{\global\SMHideFooterTexttrue}

% \SMSlideLanguageFooterTag -- resolved once here rather than inlined
% into the footer template below, same reason \SMDeckVersionText is
% precomputed rather than built inline in \deckversion. Only ever appears
% appended directly onto the version digits ("v1.1.0EN"), so it's a no-op
% whenever \ifSMSlideVersionSet is false -- ties the language tag's
% rollout to the same version-bump touchpoint slides already go through,
% instead of needing a one-time pass over every slide. For an "mult"
% slide (\multilingualtext above) this shows whichever language the deck
% itself is compiling as, not a fixed value -- \MakeUppercase{en} = "EN".
\newcommand{\SMSlideLanguageFooterTag}{%
	\ifSMSlideLanguageSet
		\ifSMSlideLanguageMult\MakeUppercase{\SMDeckLanguage}\else\slideLanguage\fi
	\fi
}

% \SMFooterIDStrut -- fixed-size strut for the footer's "(slide id: ...)"
% block (2026-09-19, direct instruction: footer text moved by ~0.27pt between
% slides). Cause: the footline's beamercolorbox only has a *minimum* height and
% depth, and the id block is the only footer text that differs from slide to
% slide -- an id containing a descender letter (g, j, p, q, y: "egk", "y89")
% is deeper than the box's minimum depth and lifts the whole footer, an id
% without one ("kaz", "c2t") leaves it at the minimum. The strut spans the
% largest ascender and descender the block can ever contain (measured with the
% block's own font size, at use time, so it follows a font change), which makes
% the footline's height and depth identical for every slide id. Must be
% called after the block's \fontsize.
\newsavebox{\SMFooterIDMeasureBox}
\newcommand{\SMFooterIDStrut}{%
	\sbox{\SMFooterIDMeasureBox}{gjpqy()[]|bdfhklt0123456789}%
	\rule[-\dp\SMFooterIDMeasureBox]{0pt}{\ht\SMFooterIDMeasureBox}%
}

% \SMFooterHide: which footer elements a deck leaves out, a comma list of
% text (the left \FooterText), id ("slide id: ... | v..."), deck ("instructor
% deck" etc.). Empty by default = everything shown. SlideMage's deck
% settings write it as \def\SMFooterHide{id,deck} in the deck's preamble;
% it is read once at \begin{document}, so the deck may set it before or
% after \input{SlideMage_template_...}. The frame number is always shown.
\providecommand{\SMFooterHide}{}
\newif\ifSMFooterShowText
\newif\ifSMFooterShowID
\newif\ifSMFooterShowDeck
\newcommand{\SMFooterHideItem}[1]{%
	\ifstrequal{#1}{text}{\SMFooterShowTextfalse}{}%
	\ifstrequal{#1}{id}{\SMFooterShowIDfalse}{}%
	\ifstrequal{#1}{deck}{\SMFooterShowDeckfalse}{}%
}
\AtBeginDocument{%
	\SMFooterShowTexttrue\SMFooterShowIDtrue\SMFooterShowDecktrue
	\expandafter\forcsvlist\expandafter{\expandafter\SMFooterHideItem\expandafter}\expandafter{\SMFooterHide}%
}

\newif\ifSMNavigationBar % \SMNavigationBar, see "Navigation for readers" at the end
\setbeamertemplate{frame footer}{%
	\csname SM@footertextdrop\endcsname%
	\vspace*{\SMFooterTextDrop}%
	\ifSMHideFooterText
		\global\SMHideFooterTextfalse
	\else
		\ifSMFooterShowText
			\FooterText\ifSMNavigationBar\SMNavigationLinks\fi
		\fi
	\fi
	\hfill%
	{\SMFooterIDSize \SMFooterIDStrut
		\ifSMFooterShowID(slide id: \slideID\ifSMOmittedSlide\SMOmittedSlideMarker\fi\ifSMSlideVersionSet\ | v\slideVersion\SMSlideLanguageFooterTag\fi\ifSMFooterShowDeck\ | \SMDeckVersionText\fi)%
		\else\ifSMFooterShowDeck(\SMDeckVersionText)\fi\fi}\hspace*{2em}\insertframenumber%
	\hspace*{-4em}%
	\setlength{\SMFooterRightShift}{\dimexpr\paperwidth-\SMFrametitleRightPadding-\hsize+\rightskip+\SMRightFooterPadding+\SMFootlineHRightOffset-4em\relax}%
	\hspace*{-\SMFooterRightShift}%
	\vspace*{0.4ex}%
}




% \source{text} -- for a "Source: ..."/"Image source: ..."/etc. citation line
% (common at the bottom/edge of a converted PowerPoint slide) that should look
% and sit *exactly* like the footnotes above -- same font, spacing,
% right-alignment, and the same stretch-to-sit-just-above-the-footer
% behaviour beamer already gives real footnotes -- but with no marker, neither
% inline in the body text nor next to the note itself, since a source line
% isn't anchored to one specific word. Implemented as a real \footnote{} (so
% it gets that positioning for free, rather than reimplementing it) with
% \thefootnote emptied locally; since \footnote{} still advances the footnote
% counter even with an empty marker, \addtocounter reverses that so a real
% \footnote{} elsewhere on the same slide doesn't skip a symbol in the
% \fnsymbol sequence. Unlike a real footnote, no label is added automatically
% here -- the label varies ("Source:", "Image source:", ...), so write it as
% part of #1.
\newcommand{\source}[1]{%
	\begingroup
	\renewcommand{\thefootnote}{}%
	\footnote{#1}%
	\endgroup
	\addtocounter{footnote}{-1}%
}

% A plain tabular with yellow \toprule/\midrule/\bottomrule, matching the
% frametitle/footer accent color. Takes the same column-spec argument as
% tabular itself, e.g. \begin{SMTable}{lccc} ... \end{SMTable}. Only
% booktabs rules are colored (via \arrayrulecolor, scoped to this table by
% \begingroup/\endgroup) -- if a table ever needs a plain black \hline too,
% that's unaffected since \arrayrulecolor only recolors rules drawn *after*
% it takes effect, and booktabs' rules read it fresh each time.
%
% RULE-BLEED WARNING (2026-07-12, direct instruction: "make sure the table
% lines do not extend beyond the text"), read this before writing a new
% SMTable: a plain column spec like `{lccc}` above leaves the default
% \tabcolsep (6pt, or whatever a slide-local \setlength last set it to) as
% *implicit* padding on the table's own outermost edges -- invisible to the
% table's own text (which starts/ends flush with each column's actual
% content), but very visible to booktabs' \toprule/\midrule/\bottomrule and
% to colortbl's \rowcolor/\cellcolor, both of which span the *declared*
% column width including that implicit padding. The result, confirmed
% directly and independently on at least three slides so far
% (slides/u5h-Long-term-Bonds-Provisions-current.tex's own v1-1-0/v1-4-0,
% slides/c4c-Macbeth-Spot-Removers-Proposed-Structure-current.tex's own
% v1-2-2, slides/tw5-Exercise-Leverage-and-ROE-Solution-current.tex's own
% v1-4-0): every rule/highlight bleeds \tabcolsep past the visible text on
% both the left and right edges, most noticeable on a left-aligned
% (non-\centering) table where the mismatch is easy to see against the
% title's own left margin.
%
% (2026-09-30: write new tables as \begin{SMTable}[gap=12pt]{...}
% instead -- see "FLUSH EDGES BUILT IN" below; the rest of this comment
% describes the hand-written form that option now produces.)
% Not fixed here in the environment definition itself: doing so would
% change every existing SMTable's own column spacing out from under it
% (this environment has no way to tell "plain default spacing was intended"
% apart from "the bleed was never noticed"). Fix per table instead, the
% same three-part pattern every slide above uses -- worth reaching for by
% default on any *new* SMTable, not just when the bleed is spotted:
%   {\setlength{\tabcolsep}{0pt}
%   \begin{SMTable}{@{}l<{\hspace{<gap>pt}}c<{\hspace{<gap>pt}}c@{}}
%   ...
%   \end{SMTable}}
% -- \tabcolsep zeroed (scoped locally, e.g. by wrapping the whole call in
% a `{...}` group, so it doesn't leak into later tables), explicit `@{}` at
% both outer edges of the column spec (so booktabs/colortbl have no
% implicit padding left to bleed into), and array's `<{\hspace{<gap>pt}}>`
% column modifier appended after every column but the last, to put back an
% explicit, *content-column* gap (matching whatever `2*\tabcolsep` the
% table used to have, if preserving its exact look, e.g. `<{\hspace{12pt}}>`
% for the ordinary 6pt-\tabcolsep default) -- `<{...}>` puts the gap
% *inside* the preceding column's own content, which `\rowcolor` (if used)
% colors along with everything else, unlike a plain `@{...}` separator,
% which `\rowcolor` deliberately leaves untinted (see u5h's own v1-6-0
% comment for that exact failure mode if a highlighted table ever needs
% this too).
%
% RE-CONFIRMED 2026-07-16, direct instruction, after this exact bleed
% reappeared on essentially every SMTable in a freshly-converted
% chapter (Capital Budgeting Techniques) despite the warning above already
% existing at the time -- the rule above is not optional/only-if-noticed,
% it's the default way to write *any* new SMTable, full stop. Check
% every new SMTable against it before considering a slide done, not
% just when a rendered page happens to make the bleed visually obvious.
%
% BLOCK-RULE-COLOR CONVENTION (2026-07-16, direct instruction): a
% SMTable placed *inside* any block environment (`infoblock`,
% `exampleblock`, `referenceblock`, `definitionblock`, plain `block`, or
% their `SMFloating*Block` wrappers) should use black rules, not this
% environment's own default yellow -- yellow-on-the-block's-own-tinted-
% background reads poorly, unlike yellow-on-plain-white outside a block.
% Override locally, *inside* `\begin{SMTable}{...}`, before `\toprule`
% (not before `\begin{SMTable}` itself -- the environment's own
% opening code sets yellow *after* its `\begingroup`, so a color set
% before that point is simply overwritten; confirmed directly, first
% attempt this same day):
%   \begin{SMTable}{...}
%   	\arrayrulecolor{c_primary_dark}
%   	\toprule
%   	...
%   \end{SMTable}
% Booktabs rules read `\arrayrulecolor` fresh at each rule, and
% `\endgroup` (SMTable's own closing code) resets it back to yellow
% afterward, so this never leaks into a later table. First applied on
% slides/6ux-Example-51-BMA-Payback-Table-current.tex's own
% `SMFloatingInfoBlock`.
%
% FLUSH EDGES BUILT IN (2026-09-30, direct instruction: "propose a more
% systematic solution" for per-slide fixes -- 53 tables carried the
% three-part pattern above by hand): give the column gap as an option and
% write the plain column spec,
%   \begin{SMTable}[gap=12pt]{lccc} ... \end{SMTable}
% which produces exactly the hand-written form -- \tabcolsep 0pt, @{} at
% both edges, <{\hspace{12pt}} after every column that is followed by
% another one (not before an explicit @{...}/!{...} separator), *{n}{...}
% expanded -- so a table rewritten this way renders identically. After a
% fixed-width p/m/b/w column the gap is an @{\hspace{...}} separator
% instead: array places a <{...} suffix inside such a column's box, where
% it only pads the end of the cell text and leaves the columns touching
% (the hand-written pattern had exactly that flaw for p columns). The
% expansion is done by SlideMage_measure.lua (slidemage_measure.gap_spec).
% gap=12pt is the ordinary look (2 x beamer's 6pt \tabcolsep). Without the
% option the environment behaves as it always has (plain tabular with
% yellow rules), so every existing table is untouched.
\makeatletter
\pgfqkeys{/SMTable}{
	gap/.store in=\SMTable@gap,
}
\ifluatex
	\SMLuaOnce{
		local path = kpse.find_file("SlideMage_measure.lua", "lua")
			or kpse.find_file("SlideMage_measure.lua", "tex")
		if path then
			dofile(path)
		end
	}
	% The two numbers the Top Margin measurement needs. The ascender is read from the text font in
	% force (luaotfload's OpenType metrics) and replaces the layout profile's \SMFontAscender where
	% the font gives one; the Top Margin is the template's (SlideMage_setup.tex). Not replayed by
	% a precompiled format: the macro it sets lives in the format.
	\begingroup\normalfont\selectfont\directlua{slidemage_measure.derive_ascender()}\endgroup
	\SMLuaOnce{
		slidemage_measure.configure{top_margin = "\SMTopMargin", ascender = \SMFontAscender}
	}
\fi
\hypersetup{pdfinfo={SMFontAscender={\SMFontAscender}}}
\def\SMTable@begin#1{\begin{tabular}{#1}}
% BLOCK-RULE-COLOR, automatic (2026-09-30): inside any block (the
% convention above) the rules default to c_primary_dark -- all 11 tables
% placed in a block so far set that by hand. An \arrayrulecolor inside the
% table still overrides it, as before.
\newif\ifSM@inblock
\AtBeginEnvironment{block}{\SM@inblocktrue}
\AtBeginEnvironment{exampleblock}{\SM@inblocktrue}
\AtBeginEnvironment{alertblock}{\SM@inblocktrue}
\AtBeginEnvironment{infoblock}{\SM@inblocktrue}
\AtBeginEnvironment{referenceblock}{\SM@inblocktrue}
\AtBeginEnvironment{definitionblock}{\SM@inblocktrue}
\NewDocumentEnvironment{SMTable}{O{}m}{%
	\begingroup
	\ifSM@inblock
		\arrayrulecolor{c_primary_dark}%
	\else
		\arrayrulecolor{c_primary_highlight}%
	\fi
	\let\SMTable@gap\@empty
	\pgfqkeys{/SMTable}{#1}%
	\ifx\SMTable@gap\@empty
		\def\SMTable@spec{#2}%
	\else
		\setlength{\tabcolsep}{0pt}%
		\directlua{slidemage_measure.set_gap_spec(
			"\luaescapestring{\detokenize{#2}}",
			"\luaescapestring{\detokenize\expandafter{\SMTable@gap}}")}%
	\fi
	\expandafter\SMTable@begin\expandafter{\SMTable@spec}%
}{%
	\end{tabular}%
	\endgroup
}
\makeatother

% For a side-corner image overlay (the usual \vbox to 0pt{...} tikz overlay
% pattern used throughout this deck's converted slides) whose *top edge*
% should line up with the content text's own top edge -- distance from the
% page's own top edge to the cap-height top of a standardframe's first body
% line, measured directly against the compiled PDF's own text coordinates,
% not derived from \SMContentTopShift or guessed. Previous per-slide
% image placements each guessed their own downward shift (e.g. a bare
% "-1.3cm") independently -- close, but not exact, landing the image's top
% a couple mm above the text's own top; use this constant instead of a
% fresh guess for any new "text top edge"-aligned image.
%
% 41.58pt (2026-07-22, re-measured): the previous value, 44.0pt, went stale
% the moment \SMContentTopShift itself was increased (7.33pt -> 11.5964pt,
% "reduce the Top Margin", see that length's own comment above) without a
% matching update here -- that change pulls plain body text up, but this
% length is a separate, hand-measured constant with no formula tying it to
% \SMContentTopShift, so it kept anchoring images at the *old*, lower
% position. Confirmed directly (pymupdf against this same freshly-compiled
% preview deck, not assumed from the \SMContentTopShift delta alone):
% three independent plain-first-line standardframes with no per-slide
% top-shift correction (slides/5v4-CML-vs-SML-current.tex,
% slides/u5j-Exercise-CAPM-4-current.tex,
% slides/nt4-Exercise-CAPM-5-current.tex) all land their first line's own
% top edge at y=41.58pt post-change, not 44.0pt. Re-measure and adjust
% again if \SMContentTopShift ever changes in the future.
% (v9.7.0: \SMContentTopAlign is derived from \SMTopMargin, further up; it was a hook of SlideMage_setup.tex, 41.58pt, in v9.4.0-v9.6.0.)

% For a "side filler image" -- a decorative image stuck flush to the page's
% own right edge (zero margin) with its top edge flush against the title
% bar's own bottom edge (zero gap), rather than sitting in the usual
% content-aligned corner inset -- distance from the page's own top edge
% down to the standardframe title bar's bottom edge, measured directly
% against the compiled PDF (a pixel scan for the yellow-to-white
% transition, at this deck's 16cm-wide/9cm-tall page size). First used on
% slides/dfc-Angel-Investors-current.tex.
%
% 25.43pt -- the raw measured value (originally 25.44pt; re-measured
% 2026-07-07 against the current compiled PDF via the title bar's own
% white-patch rectangle, which came back 25.427pt, close enough to the
% original measurement to be the same underlying edge, not a drift).
% Briefly set to 24.94pt instead (2026-07-06-ish), a deliberate 0.5pt
% overlap *into* the title bar, after a report of a visible hairline
% white gap right at this seam in one render -- not reproduced in this
% project's own pdftoppm-based rendering checks at the time (every pixel
% scan at the raw value showed a clean cut, no gap), so assumed to be a
% renderer-specific sub-pixel rounding artifact rather than a real error
% in the constant, and "fixed" by deliberately overlapping instead of
% chasing the specific viewer. That overlap turned out to be the wrong
% call: direct feedback 2026-07-07 reported the *opposite* symptom on
% `sideimageframe` (the image's own top edge visibly overlapping the
% title box) -- exactly what a real 0.5pt overlap looks like once it's
% actually noticed, rather than a rendering artifact. Reverted back to
% the precise raw measurement: a real, small, deliberate overlap chosen
% to paper over an unreproduced gap in an unspecified other renderer
% turns out to be a worse trade than just using the measured value and
% letting a genuine future gap report (if one ever recurs, in a
% specific, identified renderer) be diagnosed on its own terms.
%
% Re-measured 2026-07-16, 25.43pt -> 25.32pt: the frametitle hook's
% outer sep fix (see its own comment) moved the visible yellow bottom
% edge up by the 0.2pt the trim patch's fill had previously been inset
% below its SMLogoAnchor.south anchor -- the yellow now genuinely
% ends flush with the logo's bottom edge, as always intended. New raw
% measurement against the rebuilt PDF: the trim patch's fill top edge
% (= the visible yellow bottom) sits 25.228 PDF units below the page
% top, i.e. 25.3226pt (earlier measurements in this comment quoted PDF
% coordinate values directly as "pt"; PDF units are big points, 1bp =
% 1.00375pt -- a systematic ~0.1pt understatement at this magnitude,
% harmless then, but corrected now that the constant is being reset
% anyway). Note this value is also what the overlay-logo nodes
% (titleframe/quoteframe/the hook's fixed-position redraw) derive their
% vertical inset from: visible-bar-bottom minus \SMLogoHeight is
% exactly the inline logo's own vertically-centered top margin.
% (v9.7.0: \SMTitleBarBottomAlign is derived from the paddings, the frame title font and the logo height, further up; it was a hook of SlideMage_setup.tex, 25.32pt, in v9.4.0-v9.6.0.)

% Note on the [t] default below: it top-aligns content within the frame
% (beamer's own default is vertical centering, "c"), so slides don't have
% their content floating in the middle. It also sidesteps a real beamer bug:
% \begin{frame}[]{Title} (a literal empty options bracket) fails outright,
% even in plain beamer with no theme at all -- so the optional argument must
% never default to an empty string. The title is always passed as an
% explicit argument (rather than left for beamer to pick up from a trailing
% brace group), since anything placed between \begin{frame} and its title --
% like our \setbeamertemplate calls -- breaks that trailing-brace mechanism.

% Pulls the frame's content up closer to the title bar: measured (via the
% PDF's own coordinates) at about 22pt of gap by default between the title
% bar's bottom edge and the first line of body content: this is a third of
% that, shrinking the gap to about two thirds of its original size.
%
% Called "Top Margin" in conversation from 2026-07-12 onward (direct
% instruction) -- the deck-wide name for the resulting gap between the
% title bar and a standardframe's first line of body content, i.e. the
% vertical position a plain-text frame's first baseline naturally lands at
% (measured directly against a reference slide's own compiled PDF, e.g.
% slides/aii-Leverage-Effects-current.tex's first line). Any first content
% that isn't plain text (a table, a minipage, an itemize with an unusual
% first item, ...) needs its own per-slide \SMIncreaseContentTopShift/
% \SMReduceContentTopShift correction to land on this same baseline --
% not portable between slides even for superficially similar content (see
% slides/zgz-EBIT-Definition-current.tex's own v1-2-0 comment for a case
% where an earlier, wrong correction was carried over uncritically).
% (\SMContentTopShift was a number of the layout profile from v9.4.0 to v9.8.0 and is measured
% since v9.9.0, see AUTOMATIC TOP MARGIN; its history is below.)
% 11.5964pt (2026-07-22, direct instruction: "reduce the Top Margin"
% generally, tested at 20%/30%/40% smaller against
% slides/8s8-Unit-Investment-Trusts-current.tex before settling on 40% --
% up from the original 7.33pt so the *resulting* gap (this length pulls
% content up by, so a bigger value means a smaller visible margin) comes
% out to 60% of its old size: measured directly on 8s8 at 10.67pt before,
% 6.42pt after (target was 6.40pt, within 0.02pt). Since every
% standardframe applies this same pull-up before any per-slide
% \SMIncreaseContentTopShift/\SMReduceContentTopShift correction,
% changing it here shifts the baseline project-wide, not just on 8s8.

% A `columns` environment that is the first content of a frame starts its first line
% ~10.17pt lower than a plain list does (measured on 5s6/vbb/yxv/c2t/n83/tj4/eey/yvw: first
% line box at 51.7pt against the 41.58pt Top Margin). It used to be corrected slide by
% slide with \SMIncreaseContentTopShift{10.13..10.17pt}, each value re-measured (and
% re-tuned when \SMContentTopShift changed); now the environment corrects itself
% (2026-09-21, direct instruction: "if this is a general problem for two-column slides, fix
% it generally"). "First content" = the frame's vertical list holds no box taller than
% 0.5pt yet -- the frame's own zero-size stamp boxes, glue and a zero-height overlay picture
% (\vbox to 0pt{...}) in front of the columns do not count, so a slide that puts such a
% picture first needs no special treatment either. (\prevdepth cannot tell: those zero-size
% boxes already set it.) Checked in LuaTeX by walking the current list.
% (\SMColumnsTopCorrection was a number of the layout profile from v9.4.0 to v9.8.0.)
\SMLuaOnce{
	slidemage_list_has_no_visible_box = function()
		local head = tex.nest[tex.nest.ptr].head
		for n in node.traverse(head) do
			if (n.id == node.id("hlist") or n.id == node.id("vlist")) and n.height > 32768 then
				return false
			end
		end
		return true
	end
}
% Since 2026-09-30 only for frames the automatic Top Margin doesn't manage
% (see AUTOMATIC TOP MARGIN below: [c] frames, sectionframe,
% \SMManualTopMargin, a deck with \SMAutoTopMarginfalse) -- in every
% other frame the measurement puts a columns-first slide on the margin by
% itself, so this correction (and the [T] compensation below) stays out.
% v9.9.0: the two corrections are measured where they are used, not numbers of the layout
% profile (10.17pt and 13.76pt, fitted on Fira; with Latin Modern Sans under another title bar a
% [T] columns landed 5.2pt high). Three reference lists are typeset in boxes at this point -- a
% top-level itemize, the same itemize in the first column of a columns[t], and of a columns[T]
% -- and \SMColumnsTopCorrection is how much lower the [t] one starts, \SMColumnsTTopCompensation
% how much of that the [T] one has taken back; a list in a column of either kind then starts
% where the list would start without the columns. Both names stay readable (the last values
% measured, 0pt before the first columns of the deck).
\makeatletter
\newcommand{\SM@columnsprobe}{%
	\SM@toprefprobe{\begin{itemize}\item\SM@topreftext\end{itemize}}\let\SM@columnsref\SM@topref
	\SM@toprefprobe{\begin{columns}[t]\begin{column}{0.46\linewidth}\begin{itemize}\item\SM@topreftext\end{itemize}\end{column}\end{columns}}%
	\global\SMColumnsTopCorrection=\dimexpr\SM@topref sp-\SM@columnsref sp\relax
	\SM@toprefprobe{\begin{columns}[T]\begin{column}{0.46\linewidth}\begin{itemize}\item\SM@topreftext\end{itemize}\end{column}\end{columns}}%
	\global\SMColumnsTTopCompensation=\dimexpr\SMColumnsTopCorrection+\SM@columnsref sp-\SM@topref sp\relax
}
\BeforeBeginEnvironment{columns}{%
	\ifSM@topprobing\else
	\ifnum0\ifdefined\SMTopAutoActive\SMTopAutoActive\fi=1
		\def\SMColumnsAtTop{0}%
	\else
		\edef\SMColumnsAtTop{\directlua{tex.sprint(slidemage_list_has_no_visible_box() and "1" or "0")}}%
	\fi
	\ifnum\SMColumnsAtTop=1
		\ifluatex\ifvmode\SM@columnsprobe\fi\fi
		\vspace*{-\SMColumnsTopCorrection}%
	\fi
	\fi
}
\makeatother
% The correction above is measured for `columns[t]`. `columns[T]` (top-aligned, used on
% most text-next-to-image slides) already starts its columns higher by itself -- beamer's
% own T key adds \vskip-1ex\nointerlineskip at every column head -- so the same pull-up left
% its first line 13.70pt (itemize) / 13.82pt (plain text) above the Top Margin, flush
% against the title bar (2026-09-30, direct instruction, seen on c58-DE-Unikat-Tipps).
% Measured against a plain standardframe: itemize 41.58pt, [t] 41.58pt, [T] 27.89pt;
% plain text 38.60pt, [T] 24.78pt. The T key gives the difference back while the columns
% environment is still in vertical mode, i.e. only for the columns[T] call itself (the
% per-column \setkeys calls run inside the columns' hbox, where \ifvmode is false) and only
% when the hook above applied its pull-up. A block as first content of a [T] column lands
% ~7pt below its plain-frame position with this -- extra clearance, not a collision.
% The column head also sets LaTeX's "top of a minipage" flag again (v5.14.1, found by
% check_layout.py's `coltop` in the SlideMage manual): beamer opens each column with
% \leavevmode, which clears it, so a list at the top of a [T] column kept its \topsep --
% zero at normal size, but a smaller list (\shrinklists, \small) loads \@listi's 4pt and
% started that much below a figure beside it that sat on the Top Margin.
% (\SMColumnsTTopCompensation was a number of the layout profile from v9.4.0 to v9.8.0; it is
% measured since v9.9.0, see above, and the figures in this paragraph are history.)
\makeatletter
\define@key{beamer@col}{T}[true]{%
	\def\beamer@colmode{T}\def\beamer@colalign{t}\def\beamer@colheadskip{\vskip-1ex\nointerlineskip\@minipagetrue}%
	\ifvmode
		\ifnum0\ifdefined\SMColumnsAtTop\SMColumnsAtTop\fi=1
			\vspace*{\SMColumnsTTopCompensation}%
		\fi
	\fi
}
\makeatother


% "Top Margin" (above) is one of four deck-wide margin names introduced in
% conversation from 2026-07-12 onward (direct instruction) for the plain
% page-edge/title-bar/footer gaps every slide is expected to obey by
% default. The other three, for reference (each maps to a length already
% declared elsewhere, not a new one introduced here):
% - "Left Margin": the default gap from the page's own left edge to the
%   title text and body content -- \SMFrametitleLeftPadding
%   (beamerthemeSlideMage.sty, 2.2ex, measured
%   directly against the compiled PDF at 12.552pt in this deck's usual
%   fonts -- see e.g. slides/tw5-Exercise-Leverage-and-ROE-Solution-current.tex's
%   own v1-1-0 comment for that same measured value).
% - "Right Margin": the default gap from the page's own right edge to the
%   institute logo (and, by the same convention, to anything else
%   deliberately flush-right, e.g. a corner-anchored image) --
%   \SMFrametitleRightPadding (0.28cm/7.93pt, this file's own
%   \SMInstituteLogo placement above uses it directly as the
%   logo's own inner sep).
% - "Bottom Margin": the default gap between a frame's own body content
%   and the \source{}/footnote line below it. Unlike the other three, this
%   one has no single settable length behind it -- \source{}'s own
%   vertical position isn't fixed the way the footer box is; it tracks the
%   frame's actual content height once that content gets tall enough (see
%   slides/u5h-Long-term-Bonds-Provisions-current.tex's own v1-3-0 comment,
%   "past a certain table height it starts tracking the table's own growth
%   roughly 1:1"). "Respecting Bottom Margin" in practice means leaving
%   comparable clearance to whatever a reference slide with unremarkable
%   content already shows, the same empirical, measure-against-a-reference-
%   slide approach "Top Margin" itself uses, rather than a fixed constant.

% Compensates a gap that opens up specifically when a standardframe's body
% starts with a \begin{minipage}{...} (used to width-limit text next to an
% image/diagram overlay, e.g. slides/igr-Budgeting-current.tex,
% slides/bdh-Working-Capital-Example-current.tex -- see README.md's
% "Available layouts" section for why minipage instead of columns there) --
% a minipage is a *box*, and TeX's vertical glue insertion (\topskip
% specifically) treats the first box in a vertical list differently from
% the first line of plain paragraph text, so the same \SMContentTopShift
% pull-up above lands closer to the title for minipage-wrapped content than
% for plain content. Cannot be fixed by \vspace(*) right after
% \begin{minipage}{...}: both plain and starred forms were tried and
% measurably had zero effect (the star only protects glue from page-break
% discard, not from \topskip's own separate glue-insertion calculation,
% which absorbs leading glue at the start of a vertical list regardless).
% Fixed instead by shrinking standardframe's own pull-up for the affected
% frame -- see \SMContentTopShiftReduction below. The value below is
% calibrated empirically, not derived from first principles: an isolated
% two-frame test (single-level itemize, one plain, one minipage-wrapped)
% measured only a 4.4pt gap, but real slides (nested itemize inside the
% minipage) showed a larger ~7.5-7.6pt gap against a real plain-itemize
% slide (slides/jyn-Categorization-current.tex) -- the extra few points
% likely come from itemize's own topsep/partopsep being computed differently
% depending on whether the list opens a box or a paragraph, on top of the
% topskip/box effect alone. Tuned so slides/igr-Budgeting-current.tex and
% slides/bdh-Working-Capital-Example-current.tex's first content line
% lands within 0.1pt of slides/jyn-Categorization-current.tex's own
% (200.917pt vs. 200.871pt/200.980pt in the compiled PDF).
% v9.4.0: the four corrections below (SMMinipageTopCorrection, SMMinipageWithCornerImageExtraCorrection,
% SMSideImageFrameMinipageCorrection, SMFlatItemizeTopCorrection) are 0pt. The automatic Top Margin (below) moves
% every frame it manages onto the Top Margin, so the slides that pass them to \SMReduceContentTopShift or
% \SMIncreaseContentTopShift land where they did: a copy of a real project's 16 decks (5274 pages, 122 slides
% that read SMMinipageTopCorrection, 123 SMFlatItemizeTopCorrection) rendered pixel-identically with the
% values at 0. The names stay defined because slides read them. The numbers in the comments are history.
\newlength{\SMMinipageTopCorrection}
% 9.38pt (2026-07-22, re-measured): the previous 7.6pt value went stale
% the moment \SMContentTopShift itself was increased (7.33pt ->
% 11.5964pt, "reduce the Top Margin", see that length's own comment above)
% without a matching update here -- confirmed directly (pymupdf against a
% freshly-compiled preview deck, not assumed from the \SMContentTopShift
% delta alone, which doesn't translate 1:1 into a minipage's own \topskip
% behavior): four independent slides all using this constant completely
% unmodified (slides/fr3-Forms-of-Capital-Market-Efficiency-Strong-Semistrong-Weak-current.tex,
% slides/za2-Startup-Financing-Categorization-current.tex,
% slides/yke-Implications-of-Homemade-Dividends-current.tex,
% slides/6j4-Modigliani-Miller-Assumptions-current.tex) all land their
% own first line at an identical y=39.80pt post-change, 1.78pt short of
% this deck's current 41.58pt Top Margin reference -- not slide-specific
% noise, the same deterministic amount on every one of them. Re-measure
% and adjust again if \SMContentTopShift ever changes in the future.
\setlength{\SMMinipageTopCorrection}{0pt}% obsolete since v9.4.0 (9.38pt before): the automatic Top Margin corrects it; slides still read the name

% A *third* minipage-top correction, on top of \SMMinipageTopCorrection
% above rather than replacing it, for a minipage whose own first internal
% content is a zero-height `\vbox to 0pt{...}` tikz overlay (a
% corner-anchored decorative image, placed there specifically so it stays
% unconditionally visible even when the minipage's own later content
% includes a `\pause` -- see e.g.
% slides/y5z-CoCo-Bonds-current.tex's own "v1-2-0" note for the fuller
% pause-inheritance story behind why the image has to sit there at all).
% 2026-07-08, found while chasing "the top margin of the text on y5z/22s
% is too low" -- measured directly against the compiled PDF, y5h/22s both
% landed 16.5pt too far down despite already using
% \SMMinipageTopCorrection, with the zero-height image overlay sitting
% *before* the minipage entirely as the actual root cause (that ordering
% makes the image, not the minipage, the frame body's true first box,
% which \topskip's own special first-box glue handling reacts to
% differently than a plain first minipage). Moving the image inside the
% minipage as its own first content closed most of the gap on its own
% (restoring the minipage as the frame's true first box) but left a
% further 10.35pt uncorrected -- confirmed the same on two independent
% slides (y5z, 22s) with identical structure, not assumed transferable
% from one to the other. Call via
% `\SMReduceContentTopShift{\dimexpr\SMMinipageTopCorrection-\SMMinipageWithCornerImageExtraCorrection\relax}`
% for any slide combining this exact pattern (corner image as a minipage's
% own first content, ahead of a `\pause`) -- confirmed transferable as-is
% to two more slides with the same structure
% (slides/qbj-Bank-Loans-Credit-Lines-current.tex,
% slides/jkf-Bank-Loans-Term-Loan-current.tex), neither of which
% actually uses `\pause` itself, so this isn't specific to that -- the
% image-inside-minipage-as-first-content shape alone is what triggers it.
\newlength{\SMMinipageWithCornerImageExtraCorrection}
\setlength{\SMMinipageWithCornerImageExtraCorrection}{0pt}% obsolete since v9.4.0 (10.35pt before), see above

% A *different* minipage-top correction, specifically for a minipage+
% itemize sitting as the first content of a `sideimageframe` (2026-07-06,
% found while chasing "the top margin of the text on ezz/acw is again
% off" -- see CONVERSION_NOTES.md's own "beamer/TeX gotchas" entry for the
% full story). `\SMMinipageTopCorrection` above is calibrated for a
% *plain* standardframe, and simply does not transfer to `sideimageframe`:
% measured directly on three real sideimageframe slides
% (slides/xmp-Implications-of-Capital-Market-Efficiency-Considerations-current.tex,
% slides/dfc-Angel-Investors-current.tex,
% slides/ezz-Examples-of-Tests-for-Market-Efficiency-current.tex --
% one flat itemize, one flat itemize, one nested-first-item itemize
% respectively, so this isn't a flat-vs-nested effect either), all three
% landed within 0.6pt of each other at their own *uncorrected* natural
% position (52.03-52.57pt against this deck's usual 44.0pt first-line
% reference) regardless of which of the two *wrong* corrections
% (`\SMReduceContentTopShift{\SMMinipageTopCorrection}` or
% `\SMIncreaseContentTopShift{\SMFlatItemizeTopCorrection}`) had
% been applied to each beforehand -- i.e. `sideimageframe` itself pushes
% its own first-content minipage down by a consistent ~8.2pt relative to
% a plain standardframe, on top of whatever generic itemize-shape
% correction would otherwise apply, most likely from the extra
% \sidefootnotes/\settowidth setup work `sideimageframe`'s own opening
% code does before `\begin{standardframe}` ever runs. Use
% `\SMIncreaseContentTopShift{\SMSideImageFrameMinipageCorrection}`
% (note: *increase*, not reduce) as the *only* topshift call on any
% future sideimageframe slide whose first content is a minipage+itemize --
% not the generic `\SMMinipageTopCorrection`/
% `\SMFlatItemizeTopCorrection` this deck uses everywhere else.
\newlength{\SMSideImageFrameMinipageCorrection}

% 6.26pt (2026-07-22, re-measured): stale for the same reason
% \SMMinipageTopCorrection above went stale -- \SMContentTopShift
% itself changed (7.33pt -> 11.5964pt) without a matching update here.
% Re-measured directly against a freshly-compiled preview deck across six
% independent sideimageframe slides using this constant unmodified: four
% (ezz, xmp, hzx, u6n) land their own first line at an identical
% y=39.64pt, 1.94pt short of this deck's current 41.58pt Top Margin
% reference (i.e. already pulled up *too far*); the other two (c7w, dfc)
% are close but not identical (39.53pt/-2.05, 40.18pt/-1.4) -- no single
% value fixes all six exactly, so this recalibrates to the four-way
% cluster (the majority case) and accepts a small (<0.6pt) residual on the
% two outliers. Since this constant feeds
% `\SMIncreaseContentTopShift{}` (more pull-up, i.e. *smaller* y), and
% the observed y was already short of target, the fix is a *decrease*
% (8.2pt - 1.94pt), not an increase -- caught by re-measuring after a
% first, wrong-direction attempt at 10.14pt drove the four-slide cluster
% to y=37.70pt, further from target instead of closer.
\setlength{\SMSideImageFrameMinipageCorrection}{0pt}% obsolete since v9.4.0 (6.26pt before, 8.2pt originally), see above

% Opt-in override: reduces standardframe's \SMContentTopShift pull-up by
% this amount for the frame currently being opened. Empty-default-style flag
% (mirrors \SMChapterLogoPath etc.): 0pt by default (no effect), set via
% \SMReduceContentTopShift{<length>} called *before* \begin{standardframe},
% reset back to 0pt in standardframe's own closing code so it never leaks
% into the next frame. Used by minipage-based slides to counteract
% \SMMinipageTopCorrection above.
\newcommand{\SMContentTopShiftReduction}{0pt}
\newcommand{\SMReduceContentTopShift}[1]{\global\def\SMContentTopShiftReduction{#1}}

% A *second*, separate first-line-position quirk, found and fixed
% 2026-07-05 after a run of slides (batch 4 plus several earlier ones) were
% found to land at two visibly different heights depending purely on
% content shape, with no minipage involved at all: a plain (non-minipage)
% standardframe whose top-level itemize's very *first* \item is immediately
% followed by its own nested \begin{itemize}/\begin{enumerate} lands 1.834pt
% higher than one whose first \item has no such immediate child (checked
% against every content slide converted so far via each compiled page's own
% first-line PDF baseline coordinate, not eyeballed -- 13 "nested-first-item"
% slides all landed within 0.05pt of each other, and 9 "flat-first-item"
% slides all landed within 0.05pt of each other, at a clean +1.834pt offset
% from the first group; isolated two-frame test cases, differing only in
% presence/absence of that immediate nested list, reproduced the exact same
% 1.834pt gap independent of which text/list type/item count was used).
% Not chased to one named TeX primitive (most likely some interaction of
% \addvspace's glue-coalescing behaviour with \partopsep/\topsep at the
% \begin{itemize}/\begin{frame} boundary, given the existing
% \SMMinipageTopCorrection comment above already flags topsep/partopsep
% as varying with box vs. paragraph context) -- calibrated empirically like
% \SMMinipageTopCorrection, not derived from first principles.
% Interacts with minipage differently, not additively: inside a minipage,
% this nested-vs-flat distinction disappears entirely (both land within
% 0.05pt of each other) -- so this correction is for *plain* itemize-first
% standardframes only, never combined with \SMMinipageTopCorrection.
\newlength{\SMFlatItemizeTopCorrection}

% 0.01pt (2026-07-22, re-measured): stale for the same reason
% \SMMinipageTopCorrection above went stale -- \SMContentTopShift
% itself changed (7.33pt -> 11.5964pt) without a matching update here.
% Re-measured directly against three independent flat-itemize slides using
% this constant unmodified (3ce, 5cr, c8e): all three landed their own
% first line at an identical y=39.76pt, 1.82pt short of this deck's
% current 41.58pt Top Margin reference -- coincidentally, the recalibrated
% \SMContentTopAlign reference now already accounts for almost all of
% what this correction used to supply on its own, leaving next to nothing
% left to add. Effectively zeroed out rather than removing the call
% itself, so every flat-itemize slide's own
% `\SMIncreaseContentTopShift{\SMFlatItemizeTopCorrection}` line
% keeps working unchanged.
\setlength{\SMFlatItemizeTopCorrection}{0pt}% obsolete since v9.4.0 (0.01pt before, 1.834pt originally), see above

% Opposite-direction counterpart to \SMReduceContentTopShift: *increases*
% the pull-up instead of reducing it (needed for \SMFlatItemizeTopCorrection
% above, since a flat first item needs *more* pull-up, not less). Same
% storage variable, just negated -- \SMContentTopShiftReduction is
% allowed to go negative, since the frame-opening code only ever computes
% \SMContentTopShift-\SMContentTopShiftReduction, never checks its
% sign.
\newcommand{\SMIncreaseContentTopShift}[1]{\global\def\SMContentTopShiftReduction{-#1}}

% AUTOMATIC TOP MARGIN (2026-09-30, direct instruction: "propose a more
% systematic solution" -- 195 slides carried a hand-measured
% \SMIncrease/ReduceContentTopShift, 89 of them a raw pt value, every one
% of them re-measured by hand whenever \SMContentTopShift or a frame type
% changed, and a stale one survived a central change unnoticed more than
% once, e.g. six slides double-corrected after the 2026-09-21 columns fix).
% Where a frame's first line lands depends on what comes first -- a plain
% paragraph sits 3pt higher than a list, an image 3.5pt higher, a minipage,
% columns, a table each differently -- so no constant pull-up can be right
% for all. Instead every top-aligned standardframe (and figureframe,
% highlightframe, sideimageframe) now measures where its first visible
% content actually landed (at shipout, SlideMage_measure.lua; overlays at
% absolute positions don't count, they are placed on the margin already) and
% moves its body onto the Top Margin line on the next run, via \jobname.top
% -- the same record-then-apply-on-rerun pattern as \jobname.ovl above;
% latexmk reruns while the file changes, and one rerun converges.
% Any older \SMIncrease/ReduceContentTopShift call still works but no
% longer matters: the measurement corrects whatever is left.
% Exceptions: [c] frames (centred on purpose: no pull-up at all), sectionframe
% (own geometry), and any frame opened after
%   \SMManualTopMargin
% which keeps the frame's default position: no correction that depends on what
% the frame starts with. (A \progressstrip frame is measured since v5.30.5.)
% The default position (v9.9.0) is \SMContentTopShift, the pull-up that puts the
% first line of a top-level list on the Top Margin in a standard frame. It used to
% be a number of the layout profile (11.5964pt, fitted on Fira under the default
% title bar: with another title bar or font the frames above landed up to 2.9pt
% off); now every top-aligned frame records where its body begins (the origin
% marker) and how far below that a top-level itemize's first line would start
% (typeset in a box there), and the measurement writes the pull-up most frames
% share into \jobname.top. The corrections of the measured frames are relative
% to it. In a run without \jobname.top the first titled frame's own value is used
% from the next frame on, so the first run already places the deck as the next one
% will (no deck needs a run more than with the old stored number); a frame that
% started that run elsewhere (the first one) is corrected from its default position. \SMAutoTopMarginfalse in a deck's preamble places every frame like
% \SMManualTopMargin does (the default position, measured; no content-dependent
% correction anywhere in the deck); it does not switch the measurement off.
\newif\ifSMAutoTopMargin
\SMAutoTopMargintrue
\makeatletter
\newif\ifSM@topmanual
\newif\ifSM@topsection
\newcommand{\SMManualTopMargin}{\global\SM@topmanualtrue}
\newcount\SM@topindex
\newwrite\SM@topfile
\newcommand{\SMSetTopCorrection}[2]{\global\@namedef{SM@topcorr@#1}{#2}}
\newcommand{\SMSetContentTopShift}[1]{\global\SMContentTopShift=#1\relax
	\ifluatex\directlua{slidemage_measure.top_shift_read = \strip@pt\SMContentTopShift}\fi}
\ifluatex
	\SMLuaOnce{
		luatexbase.add_to_callback("pre_shipout_filter", function(head)
			if slidemage_measure and slidemage_measure.top_measure_page then
				slidemage_measure.top_measure_page(head)
			end
			return true
		end, "slidemage top margin")
	}
	\AtBeginDocument{%
		\InputIfFileExists{\jobname.top}{}{\typeout{No file \jobname.top.}}%     % "No file": see \jobname.ovl
		\immediate\openout\SM@topfile=\jobname.top\relax
	}
	\AtEndDocument{\directlua{slidemage_measure.top_emit()}\immediate\closeout\SM@topfile}
\fi
% Before \begin{frame} (runs once per frame, not once per overlay page):
% number the frame and fetch its correction from the previous run.
\newcommand{\SM@topprepare}{%
	\ifdefined\SM@topshiftrun% a first run: measured on the first frame that counts
		\global\SMContentTopShift=\SM@topshiftrun\relax
		\global\let\SM@topshiftrun\undefined
	\fi
	\global\advance\SM@topindex by 1
	\xdef\SM@topthis{\the\SM@topindex}%
	\ifcsname SM@topcorr@\SM@topthis\endcsname
		\xdef\SM@topcorr{\csname SM@topcorr@\SM@topthis\endcsname}%
	\else
		\gdef\SM@topcorr{0pt}%
	\fi
	\ifSMAutoTopMargin\ifSM@topmanual\gdef\SM@topcorr{0pt}\fi\else\gdef\SM@topcorr{0pt}\fi
}
% In the frame body, right after the pull-up: the start marker, then the origin marker.
\newcommand{\SM@topstart}{%
	\ifluatex
		\directlua{slidemage_measure.top_begin(\SM@topthis, \strip@pt\dimexpr\SM@topcorr\relax,
			\number\dimexpr\SMContentTopShift-\SMContentTopShiftReduction+\SM@topcorr\relax,
			\strip@pt\SMContentTopShift, \ifSMAutoTopMargin\ifSM@topmanual false\else true\fi\else false\fi)}%
		\ifSMAutoTopMargin\ifSM@topmanual\else
			\def\SMTopAutoActive{1}% this frame is measured: the columns hook stays out
		\fi\fi
		\SM@toprefprobe{\begin{itemize}\item\SM@topreftext\end{itemize}}%
		\directlua{slidemage_measure.top_origin(\SM@topthis, \SM@topref,
			\ifSM@topsection false\else\ifSM@toptitled true\else false\fi\fi)}%
	\fi
}
% The reference list's first line box, in sp below this point (\SM@topref): #1 typeset in
% a box that starts as the body does here (the same \prevdepth, font, \parskip and list
% settings), measured by SlideMage_measure.lua like a frame's first content. Nothing of it
% reaches the page; a deck's default overlay specification stays out (an \item<+-> would
% count a step).
% The reference line's text: a capital and an ascender, as tall as a line of running text (a
% [T] column aligns its first line by the line's height, not by \baselineskip).
\newcommand{\SM@topreftext}{Xh}
\newbox\SM@topprobebox
\newif\ifSM@topprobing
\newif\ifSM@toptitled
\newcommand{\SM@toprefprobe}[1]{%
	\begingroup
	\SM@topprobingtrue
	\def\beamer@defaultospec{<*>}% beamer's own default: every step
	\def\SMColumnsAtTop{0}%
	\ifvmode\dimen@=\prevdepth\else\dimen@=\z@\fi
	\setbox\SM@topprobebox=\vbox{\prevdepth=\dimen@#1\par}%
	\xdef\SM@topref{\directlua{tex.sprint(slidemage_measure.box_content_top(\number\SM@topprobebox))}}%
	\endgroup
}
% End of the body, or anywhere a frame wants to stop the measurement
% (\progressstrip): nothing after this marker counts as first content.
\newcommand{\SM@topstop}{%
	\ifluatex
		\directlua{slidemage_measure.top_end(\SM@topthis)}%
	\fi
}
\makeatother

% Define custom frame type with en-dash bullets
\newenvironment{standardframe}[2][t]{%
	% Restart footnote numbering (the *, dagger, ... symbols) fresh on every
	% slide, rather than letting it run on cumulatively across the whole deck
	% the way beamer does by default. Two reasons: \fnsymbol only defines 9
	% symbols before erroring out, which a real multi-slide deck would run
	% through in no time if every \footnote{}/\source{} call across the whole
	% talk drew from one shared counter; and each slide's footnotes are read
	% and referenced only on that one slide anyway, so there's no reason for
	% their numbering to depend on how many footnotes earlier slides happened
	% to have. This must run *before* \begin{frame}: beamer saves the
	% counter's current value as that frame's baseline as part of
	% \begin{frame}'s own internal setup (restoring it before every overlay
	% replay, so overlays within one frame don't double-count) -- resetting
	% after \begin{frame} would be reset too late, after that baseline was
	% already captured.
	\setcounter{footnote}{0}%
	\csname SM@topprepare\endcsname
	\ifstrempty{#2}{\global\csname SM@toptitledfalse\endcsname}{\global\csname SM@toptitledtrue\endcsname}%
	\SMOpenFrame{[#1]{#2}}%
	% Forces this frame's own body to start with the correct default text
	% color, regardless of what any *previous* frame left active
	% (2026-07-16, direct instruction, root-caused via an isolated minimal
	% reproduction). \titleframe's/\quoteframe's own
	% \color{\SMTitleTextColor} calls -- set inside content built via a
	% stored box register or a zero-space overlay, used/shipped out *later*
	% rather than immediately -- don't reliably get "undone": xcolor's
	% \color doesn't emit a PDF color operator immediately when called, only
	% when the *next* real glyph/path actually needs it, so whatever color
	% one of those frame types last set can end up as the active PDF color
	% well into the *next* frame's own first draw call. Invisible for Graz
	% (the leaked color, c_primary_dark, happens to match ordinary body text
	% anyway) but catastrophic for MCI/Innsbruck (c_primary_white, rendering
	% plain non-itemize body text invisible against its own white
	% background -- itemize items escape it, since beamer's own itemize
	% template sets color fresh per bullet). An \AtBeginEnvironment{frame}
	% hook (etoolbox) was tried first and confirmed *not* to fire at all --
	% beamer's own "frame" isn't a plain \newenvironment etoolbox can patch
	% that way -- so this resets directly in standardframe's own opening
	% code instead, exactly where this project already has full control.
	\color{c_primary_dark}%
	% Defensive: the tokens after \begin{frame} are collected into the frame
	% body and replayed once per overlay slide, so this clears the side-
	% footnote collector at the start of every slide's pass -- normally a
	% no-op (the previous slide's footline already consumed the box), but it
	% stops stale footnotes leaking in if a previous frame ever collected
	% without shipping a footline.
	\SMClearSideFootnotesBox%
	\SMResetImageSources%
	% The itemize bullets are the theme's (\SMBulletDot, beamerthemeSlideMage.sty).
	% Only for top-aligned frames (the default, #1=t): a [c]-centered frame
	% (2026-07-22, direct instruction, found while centering
	% slides/bw9-Pearson-Correlation-Illustration-current.tex/agc/crj) was
	% landing its content almost flush against the title bar instead of
	% truly centered -- root cause was this same pull-up firing
	% unconditionally regardless of #1. beamer's own [c] option already
	% centers the frame body within the available space; stacking this
	% deck's *own* top-alignment pull-up on top of that just drags the
	% already-centered box further upward by the same fixed amount, which
	% is exactly the asymmetry that was measured (a [c] figureframe's image
	% landing ~0pt from the title bar and ~10pt from the footer instead of
	% symmetric). Guarded with \ifstrequal (etoolbox, already loaded above)
	% rather than \ifx, matching this file's own \ifstrequal{##1}{...}
	% pattern elsewhere (\ifx would need #1 fully expanded first).
	\ifstrequal{#1}{c}{}{%
		\vspace*{-\dimexpr\SMContentTopShift-\SMContentTopShiftReduction+\csname SM@topcorr\endcsname\relax}%
		\csname SM@topstart\endcsname}%
}{%
	\csname SM@topstop\endcsname
	% Must run *before* \end{frame}: like any other frame content, a
	% \footnote{}/\source{} call has to happen while the frame body is
	% still open -- everything in this closing block from \end{frame}
	% onward runs *after* the frame has already closed (see README.md's
	% own \titleframe comment on beamer's \beamer@collect@@body for the
	% general shape of this constraint).
	\SMEmitImageSourceFootnote%
	\end{frame}%
	\global\def\SMChapterLogoPath{}% opt-in only for the frame just closed
	\global\def\SMContentTopShiftReduction{0pt}% opt-in only for the frame just closed
	\global\csname SM@topmanualfalse\endcsname% opt-in only for the frame just closed
	\global\csname SM@topsectionfalse\endcsname% set by sectionframe for the frame just closed
	\global\SMSideFootnotesfalse% opt-in only for the frame just closed
	\global\def\SMFrametitleFullHeightYellow{}% opt-in only for the frame just closed
	\global\def\SMFrametitleLogoFixed{}% opt-in only for the frame just closed
	\global\def\SMFooterBoxColor{\SMFooterBoxColorDefault}% opt-in only for the frame just closed
}


% "a)"/"b)"/"c)"-labeled list, for exercise sub-tasks that read better
% lettered than bulleted (2026-07-20, direct instruction, first used on
% slides/y3a-Exercise-Diversification-2-current.tex's own bottom three
% items).
%
% First attempt (superseded, kept as a documented wrong turn -- see
% README.md's own "lualatex vs latex" section for why this project keeps
% wrong write-ups rather than deleting them): reuse itemize's own top-level
% `\llap`-based label positioning by locally overriding
% `\setbeamertemplate{itemize item}` inside a plain `\begin{itemize}` --
% the reasoning being that `\llap` draws its content flush against the
% text's own left edge "regardless of the label's width." That reasoning
% turned out to be wrong in a way that only showed up once actually
% measured against the compiled PDF (not just visually eyeballed): a plain
% bullet's own `\llap`-drawn footprint (title/Left-Margin-relative-left-
% edge to text-start) measured only ~9.05pt end-to-end on this deck's real
% slides, matching the Left-Margin-to-text-start gap exactly -- but "a)"
% rendered through the *exact same* mechanism measured a ~13.97pt
% footprint for the same gap, not the ~8.5-9pt `\settowidth` reported for
% the bare glyphs in isolation. The extra ~5pt was never fully root-caused
% (this custom itemize's own `\list` call, further below, only overrides
% `\makelabel` -- it never sets `\leftmargin`/`\labelwidth` itself, so
% those are left at whatever this deck's global itemize defaults happen to
% be, and `\llap`'s actual box-width behavior inside that inherited,
% unexamined list-parameter state is where the extra width most likely
% comes from) -- rather than keep guessing at *why*, measured the actual
% failure directly (the compiled "a)" sat ~4.9pt left of the Left Margin,
% confirmed via exact PDF text-position extraction, not a visual guess)
% and abandoned the approach once even `\tiny` plus manual per-letter
% kerning only just barely closed the gap while reading as illegibly small
% -- a real dead end, not a tuning problem.
%
% Working fix: don't reuse itemize's `\llap` label mechanism at all -- use
% a real, width-reserving `\makebox[W][r]{...}` label instead (the same
% *kind* of mechanism `\SMEnumerateLabelWidth` above already uses for
% a nested enumerate, just with an explicit `\leftmargin`/`\labelwidth`/
% `\labelsep` this environment sets itself rather than inheriting), inside
% its own plain `\list`, not a reused `\itemize`. `\leftmargin` (label box
% + gap) is set to 11.1pt -- 9.1pt label box (fits "a)"/"b)"/"c)" at this
% deck's normal `\small` itemize size with a little clearance, confirmed
% directly: measured label-left-edge stays right of the 12.552pt Left
% Margin for all three letters, closest case (widest letter, "c)") 1.37pt
% clear) plus 2pt `\labelsep` for a visible gap before the text. Trade-off,
% accepted deliberately rather than found by accident: item text lands
% ~2.2pt right of where a plain bulleted item's text would (23.8pt vs.
% 21.6pt) -- fitting a legible "a)"/"b)"/"c)" at normal reading size plus
% a real gap simply doesn't fit in the ~9pt a plain bullet's tiny dot
% leaves, so *exact* text-position parity with plain bullets and a legible
% label are mutually exclusive here; legibility won given the two-point
% shift is visually negligible and there's no plain-bulleted item left on
% the same slide to compare against directly anyway (y3a's own first item
% moved into its infoblock). Re-derive `\leftmargin`/`\labelwidth` the same
% way (render, measure with exact text-position extraction, adjust) if
% this is ever used at a different font size than this deck's normal
% itemize `\small`.
% The two lengths of every list with "a)" / "1." / "1." / "a." labels (v9.8.0: derived; they were
% 9.1pt and 4.75pt in the Fira profile and 9.16pt and 5.20pt in the lmsans one). \SMLabelItemWidth
% is the box that holds the widest label of the first items, a)-c) and 1)-9) with a bracket or a
% period, in the font in force where the list begins, plus 0.1pt of clearance; \SMLabelItemSep is the
% hook \SMLabelItemGap (in em of that font). Called first in each list's setup, so that
% the lists follow \small, \footnotesize and \shrinklists, and once at \begin{document}
% (in \small) for a slide that reads the two lengths itself.
\makeatletter
\newcommand{\SM@labelitemmetrics}{%
	\setlength{\SMLabelItemWidth}{0pt}%
	\@for\SM@labelprobe:=a),b),c),a.,b.,c.,1),2),3),4),5),6),7),8),9),1.,2.,3.,4.,5.,6.,7.,8.,9.\do{%
		\settowidth{\dimen@}{\SM@labelprobe}%
		\ifdim\dimen@>\SMLabelItemWidth\SMLabelItemWidth=\dimen@\fi}%
	\addtolength{\SMLabelItemWidth}{0.1pt}%
	\setlength{\SMLabelItemSep}{\SMLabelItemGap}}
\makeatother
\newcounter{SMAlphaItemCounter}
\newcommand{\SMAlphaItem}{\stepcounter{SMAlphaItemCounter}\alph{SMAlphaItemCounter})}
\newenvironment{SMAlphaItemize}{%
	\setcounter{SMAlphaItemCounter}{0}%
	\list{\SMAlphaItem}{%
		% 4.75pt (2026-07-20, direct instruction: "half the width of the
		% 'C' in 'Calculate'" -- \settowidth{C} at this deck's normal
		% itemize \small measured 5.5pt, so 2pt original + 2.75pt half-C).
		% Only \labelsep changes the text's own start position -- the
		% label's own box stays exactly where it was (its left edge is
		% pinned to this list's origin regardless of \labelsep, since
		% \leftmargin=\labelwidth+\labelsep always cancels \labelsep back
		% out of the label box's own position), so this widens the gap
		% without reopening the Left-Margin question already solved above.
		\csname SM@labelitemmetrics\endcsname\setlength{\labelsep}{\SMLabelItemSep}%
		\setlength{\labelwidth}{\SMLabelItemWidth}%
		\setlength{\leftmargin}{\labelwidth}%
		\addtolength{\leftmargin}{\labelsep}%
		\setlength{\itemindent}{0pt}%
		% \itemsep/\parsep left at plain \list's own defaults otherwise --
		% since this environment is built directly on \list rather than
		% reusing beamer's own itemize template, it never got beamer's own
		% itemize spacing hooks at all, only whatever \list falls back to.
		% Measured directly (2026-07-21, direct instruction, found on
		% slides/nt4-Exercise-CAPM-5-current.tex): that default reads as
		% *less* space between two items (13.7pt, item-to-item) than
		% between two wrapped lines within one *single* item (15.6pt,
		% pure \baselineskip) -- backwards from how a list should read
		% (item boundaries should be at least as open as a same-item line
		% wrap, not tighter). \itemsep set explicitly so the total gap at
		% an item boundary (baselineskip + itemsep) reads as clearly more
		% open than a plain wrapped line, not accidentally tighter.
		% Reduced from the original 6pt to 2.5pt (2026-07-22, found during
		% a full-project check after the itemize-rules work, not part of
		% that work itself): 6pt overshot by more than needed to clear the
		% 1.9pt gap (15.6 vs. 13.7), and slides/kmf-Exercise-CAPM-6-current.tex's
		% own 5-item list -- longer than nt4's own -- compounded that
		% excess across 4 item boundaries into a ~25pt overfull frame, its
		% own last item cut off by the footer (confirmed directly,
		% rendered and compared against a pre-regression PDF). 2.5pt
		% keeps a real, visible margin over the 1.9pt needed (not just
		% barely crossing it) while re-measuring clean (no overfull) on
		% kmf; re-verify against both nt4 and kmf directly if this is
		% ever changed again -- they're this environment's own tightest
		% (most items) and originally-tuned-against cases respectively.
		\setlength{\itemsep}{\SMLabelItemItemsep}%
	}
	\raggedright% plain \list defaults to justified+hyphenated text, unlike
	% this deck's own itemize (which calls this too, further below) --
	% missed on the first pass, only caught by comparing a rendered a)/b)/c)
	% item side by side against a plain bulleted one ("stan-\ndard deviation"
	% hyphenating where a plain-bulleted version of the same sentence
	% elsewhere in this deck doesn't).
}{%
	\endlist
}

% Three siblings of \SMAlphaItemize above, same measured \list metrics
% (labelsep/labelwidth/leftmargin/itemindent/itemsep), only the label text
% itself differs -- added 2026-08-09 for slides whose numbering has to
% match a real external scheme (statute paragraph numbers, PowerPoint's
% own "1./2." or "a./b." auto-numbering) rather than plain letters-with-
% parenthesis. Not merged into one parameterized environment: each name is
% called directly from a slide file, and duplicating six short \setlength
% lines three times is cheaper to read at the call site than a macro
% layer over \newenvironment for four near-identical bodies.
\newcounter{SMNumItemCounter}
\newcommand{\SMNumItem}{\stepcounter{SMNumItemCounter}\arabic{SMNumItemCounter})}
\newenvironment{SMNumItemize}{%
	\setcounter{SMNumItemCounter}{0}%
	\list{\SMNumItem}{%
		\csname SM@labelitemmetrics\endcsname\setlength{\labelsep}{\SMLabelItemSep}%
		\setlength{\labelwidth}{\SMLabelItemWidth}%
		\setlength{\leftmargin}{\labelwidth}%
		\addtolength{\leftmargin}{\labelsep}%
		\setlength{\itemindent}{0pt}%
		\setlength{\itemsep}{\SMLabelItemItemsep}%
	}
	\raggedright%
}{%
	\endlist
}

\newcounter{SMAlphaPeriodItemCounter}
\newcommand{\SMAlphaPeriodItem}{\stepcounter{SMAlphaPeriodItemCounter}\alph{SMAlphaPeriodItemCounter}.}
\newenvironment{SMAlphaPeriodItemize}{%
	\setcounter{SMAlphaPeriodItemCounter}{0}%
	\list{\SMAlphaPeriodItem}{%
		\csname SM@labelitemmetrics\endcsname\setlength{\labelsep}{\SMLabelItemSep}%
		\setlength{\labelwidth}{\SMLabelItemWidth}%
		\setlength{\leftmargin}{\labelwidth}%
		\addtolength{\leftmargin}{\labelsep}%
		\setlength{\itemindent}{0pt}%
		\setlength{\itemsep}{\SMLabelItemItemsep}%
	}
	\raggedright%
}{%
	\endlist
}

\newcounter{SMNumPeriodItemCounter}
\newcommand{\SMNumPeriodItem}{\stepcounter{SMNumPeriodItemCounter}\arabic{SMNumPeriodItemCounter}.}
\newenvironment{SMNumPeriodItemize}{%
	\setcounter{SMNumPeriodItemCounter}{0}%
	\list{\SMNumPeriodItem}{%
		\csname SM@labelitemmetrics\endcsname\setlength{\labelsep}{\SMLabelItemSep}%
		\setlength{\labelwidth}{\SMLabelItemWidth}%
		\setlength{\leftmargin}{\labelwidth}%
		\addtolength{\leftmargin}{\labelsep}%
		\setlength{\itemindent}{0pt}%
		\setlength{\itemsep}{\SMLabelItemItemsep}%
	}
	\raggedright%
}{%
	\endlist
}

% Itemize indentation: level 1 bullets should sit flush with the frame title
% (no extra left margin), and each deeper level should be indented just
% enough for its bullet to line up under the *text* (not the bullet) of the
% level above. A nested list's \leftmarginN is added on top of its parent's
% (indentation is cumulative), so giving every level the same "step" --
% label width + labelsep, the width a bullet plus its gap to text takes up --
% makes level N's text sit at N*step and level N's bullet at (N-1)*step, i.e.
% exactly at level (N-1)'s text (0 for level 1, flush with the title).
\makeatletter
\newlength{\SMBulletWidth}
\settowidth{\SMBulletWidth}{\SMBulletDot}
\newlength{\SMItemStep}
\setlength{\SMItemStep}{\labelsep}
\addtolength{\SMItemStep}{\SMBulletWidth}
\setlength{\leftmargini}{\SMItemStep}
\setlength{\leftmarginii}{\SMItemStep}
\setlength{\leftmarginiii}{\SMItemStep}
\setlength{\leftmarginiv}{\SMItemStep}
\makeatother

% \item[\SMBulletEllipsis] -- an ellipsis used as the item label (the source's
% "...how you contribute" continuation items). A plain \item[\ldots] right-aligns the
% label to the bullet dot's right edge, so the wider ellipsis pokes out to the left
% of where a bullet starts; this one starts at the dot's left edge instead
% (2026-09-20, direct instruction: the ellipses were too far left).
\newcommand{\SMBulletEllipsis}{\hbox to \SMBulletWidth{\rlap{\ldots}\hss}}

% \setbeamertemplate{itemize subitem}{\SMBulletEllipsisScript} -- the same for a list whose *bullet itself* is a small ellipsis (the source's
% "..." sub-points; slides quy, 3a8, unt). The bare \scriptsize… bullet used to sit 1.6pt left of where a dot starts (2026-09-20).
\newcommand{\SMBulletEllipsisScript}{\hbox to \SMBulletWidth{\rlap{\scriptsize…}\hss}}

% A 4th itemize level -- beamer's own \itemize (beamerbaselocalstructure.sty)
% hard-caps nesting at 3 levels (\ifnum \@itemdepth >2\relax\@toodeep),
% independently of anything this deck sets: a 4th \begin{itemize} always
% raised "! LaTeX Error: Too deeply nested." even though plain LaTeX's own
% kernel itemize supports 4. Found 2026-07-15, direct instruction, on
% slides/hrg-Exchange-Rate-Quotation-Forward-current.tex -- genuinely
% 4 levels deep in the source pptx itself (confirmed via a PowerPoint COM
% render, not assumed), not a conversion artifact to design around.
% Since this deck already styles every existing level identically (same
% \SMBulletDot, same \SMItemStep margin -- see the comments right
% above/below), extending to a 4th level needs no new visual design, just
% patching beamer's own depth limit up by one and telling it what to call
% the new level ("subsubsub", the next name in beamer's own
% sub/subsub/... naming scheme) so its template/color lookups
% (\usebeamertemplate{itemize subsubsubitem}, \usebeamercolor{itemize/
% enumerate subsubsub body}, etc.) resolve to *something* instead of
% erroring -- undefined beamer templates/colors are expected to fall back
% to sensible defaults (no extra spacing hook, plain text color) rather
% than raising an error, same as any other beamer template/color that a
% given theme never bothers to set; the actual bullet symbol still needs
% its own \SMBulletDot template (in the theme), added explicitly the same way
% as the existing 3 levels, not left to whatever fallback beamer has for
% an unset "itemize subsubsubitem" template.
% \beamer@computepref is patched wholesale (not via \patchcmd) since its
% body is one single \ifcase chain, not a stable string to match against.
% \itemize itself was *first* tried via \patchcmd{\itemize}{...}
% (matching this deck's own established practice for beamer/theme
% internals, e.g. beamerthemeSlideMage.sty's own
% \patchcmd{\beamer@title}{...}) -- this failed outright, confirmed
% directly (\patchcmd's own failure branch fired, raising a clean
% "Patching failed" error rather than silently doing nothing): recent
% LaTeX2e kernels implement an optional-argument environment like
% beamer's own `\renewcommand{\itemize}[1][]{...}` via
% `\@protected@testopt`, which stores the *real* body not under `\itemize`
% itself but under a differently-named internal control sequence (`\show
% \itemize` confirms this directly: it prints as
% `\@protected@testopt \itemize \\itemize {}`, i.e. plain `\itemize` is
% only a thin dispatcher, and the actual `\ifnum \@itemdepth >2...` body
% \patchcmd was trying to match against was never there to find). Rather
% than chase that internal storage name (liable to differ across LaTeX
% kernel versions, which \patchcmd's own documented usage doesn't need to
% care about), \itemize is fully re-declared here instead, copied
% verbatim from beamer's own beamerbaselocalstructure.sty (this project's
% installed version) with only the depth bound itself changed from 2 to
% 3 -- the same declaration form beamer itself used
% (\renewcommand{\itemize}[1][]{...}), which regenerates both the
% dispatcher and the real body correctly under whatever internal name
% the current kernel expects, sidestepping the naming question entirely.
% Trade-off, stated plainly: if a future beamer release changes this
% macro's own body for an unrelated reason, this override will silently
% keep using today's (2026-07) copy instead of picking up that change --
% acceptable here since the override is one line different from upstream
% and clearly marked, but re-diff against beamerbaselocalstructure.sty's
% own \itemize if beamer is ever upgraded and itemize behavior looks off.
% \SMNoListParsep (2026-10-01, direct instruction: "Implement a general
% fix"): a level-1 list typeset at \small or \footnotesize -- an ambient size
% switch before it, \setbeamerfont, \shrinklists -- loads LaTeX's
% size-specific \@listi, which brings \parsep 3pt plus 2pt (\footnotesize:
% similar) and \parskip with it; a normal-size list has neither. So a smaller
% list's items sat *further* apart than at normal size (6j4: 19.9pt per item
% at normal size, 20.8pt at \small), against this deck's rule that list
% spacing is em-based and shrinks with the font (\SMItemsepRatioI etc.).
% Removed wherever \parsep has a natural (non-stretch) part, i.e. exactly in
% those lists; deeper levels (0pt plus 1pt) are left as they were.
\makeatletter
\newcommand{\SMNoListParsep}{%
	\ifdim\parsep>\z@
		\setlength{\parsep}{\z@}\setlength{\parskip}{\z@}%
	\fi
}
\long\def\beamer@computepref#1{%
	\let\beameritemnestingprefix\@empty%
	\ifcase#1\or\or\def\beameritemnestingprefix{sub}\or\def\beameritemnestingprefix{subsub}\or\def\beameritemnestingprefix{subsubsub}\or\@toodeep\fi%
}
\renewcommand{\itemize}[1][]{%
	\ifblank{#1}{}{\def\beamer@defaultospec{#1}}%
	\ifnum \@itemdepth >3\relax\@toodeep\else
		\advance\@itemdepth\@ne
		\beamer@computepref\@itemdepth% sets \beameritemnestingprefix
		\usebeamerfont{itemize/enumerate \beameritemnestingprefix body}%
		\usebeamercolor[fg]{itemize/enumerate \beameritemnestingprefix body}%
		\usebeamertemplate{itemize/enumerate \beameritemnestingprefix body begin}%
		\list
			{\usebeamertemplate{itemize \beameritemnestingprefix item}}
			{\def\makelabel##1{%
					{%
						\hss\llap{{%
								\usebeamerfont*{itemize \beameritemnestingprefix item}%
								\usebeamercolor[fg]{itemize \beameritemnestingprefix item}##1}}%
					}%
				}%
			}
	\fi%
	\beamer@cramped%
	% \beamer@cramped (beamerbaselocalstructure.sty, unmodified upstream
	% code) unconditionally does `\itemsep0pt\topsep0.5pt`, every single
	% time -- called right here, in this project's own copy of beamer's
	% own \itemize body, same position as upstream. Setting \itemsep any
	% earlier (first tried as a separate \AtBeginEnvironment{itemize} hook,
	% run *before* this whole macro body even starts) gets silently wiped
	% out by this same line a moment later -- confirmed directly via
	% \typeout debugging (not assumed): the hook's own \setlength really
	% did run with the right values, \itemsep still read back as 0pt
	% immediately after. Setting it *here*, right after \beamer@cramped
	% already ran (not before), is what actually survives.
	\ifcase\@itemdepth
	\or\setlength{\itemsep}{\SMItemsepRatioI em}% level 1
	\or\setlength{\itemsep}{\SMItemsepRatioII em}% level 2
	\or\setlength{\itemsep}{\SMItemsepRatioIII em}% level 3
	\or\setlength{\itemsep}{\SMItemsepRatioIV em}% level 4
	\fi%
	\SMNoListParsep
	\raggedright%
	\beamer@firstlineitemizeunskip%
}
\makeatother

% A nested \begin{enumerate} inherits the same \leftmarginN step as itemize
% above for *where its text starts* (matching itemize's own "one step past
% the parent level's text" rule -- confirmed directly against the compiled
% PDF: an enumerate's own body text, e.g. "Senior secured debt", already
% lands exactly one \SMItemStep past its parent item's text), but its
% *label* ("1.", "2.", ...) doesn't follow \labelwidth the way plain LaTeX's
% list mechanism would -- beamer draws enumerate labels through its own
% template machinery, which turned out not to respect \labelwidth at all:
% \setlength{\labelwidth}{...} right after \begin{enumerate} measurably took
% effect (checked with \typeout{\the\labelwidth} immediately after) but the
% compiled numeral's own position never moved. What *does* move it: directly
% overriding \makelabel with a fixed-width \makebox -- empirically
% calibrated (render, measure the numeral's own x-position in the compiled
% PDF, adjust, repeat) since the box's own left edge turned out to sit at a
% fixed position unrelated to \labelwidth or \leftmarginN, with the
% right-aligned numeral inside it moving in lockstep with the box's width
% (confirmed linear: +1pt width -> +~1pt numeral position, across several
% test widths). 19.23pt (at this deck's default itemize font size) lands a
% single-digit numeral flush with the parent item's own text -- only valid
% at that font size and for single-digit lists (1-9); re-calibrate the same
% way (not by recomputing from first principles) for a differently-sized or
% 10+-item nested enumerate. Call right after \begin{enumerate}, before the
% first \item.
%
% 20.33pt (2026-07-08, direct instruction: "the numbers still extend a
% little farther to the left than the text above them does" --
% slides/f6m-Recovery-Rates-current.tex's own "Pecking order:" list).
% 19.23pt landed *close* but not quite flush: measured directly against
% the compiled PDF, every numeral's own left edge (20.55pt-20.95pt,
% varying slightly by digit width even though all are right-aligned to
% the same box) still sat up to ~1.05pt left of "Pecking order:"'s own
% text start (21.6pt) -- the widest digit tried, "6.", was the binding
% case. +1.1pt of box width (confirmed +1pt width -> +1pt numeral
% position, per this comment's own note above) covers that with a small
% margin, re-confirmed against the recompiled PDF.
\newcommand{\SMEnumerateLabelWidth}{\renewcommand{\makelabel}[1]{\makebox[\SMEnumerateLabelBox][r]{##1}}}

% Automatic version of the above (2026-07-23, background-task follow-up
% flagged in slides/kfh-Asset-Allocation-Process-current.tex's own v1-1-0
% comment: "fix the enumerate behavior generally, if possible", task
% `task_bb3fbb66`). That comment's first attempt, an
% `\AtBeginEnvironment{enumerate}{...}` hook, failed because the hook fires
% *before* `\enumerate`'s own body runs, and that body's own internal
% `\list{...}{\def\makelabel...}` setup (the beamer copy directly below the
% `\itemize` override above) clobbers whatever a hook that early leaves
% behind -- the exact same "template runs its own setting, external pre-set
% gets overwritten" problem `\itemize`'s own `\itemsep` comment above
% already documents. The real fix, same idea `\itemize` itself used above:
% override beamer's own low-level piece directly rather than hooking
% \AtBeginEnvironment. Unlike `\itemize` (declared via
% `\@protected@testopt`, needing a full re-declaration to reach -- see the
% `\itemize` override's own comment for the full story), `\enumerate`'s
% relevant piece, `\beamer@enum@` (beamerbaselocalstructure.sty), is a
% plain `\def` with no optional-argument indirection, so it can be
% `\renewcommand`'d directly, verbatim from upstream except for the
% `\ifnum\@itemdepth>1` branch added below.
%
% Only fires for a *nested* enumerate (`\@itemdepth>1`, i.e. already inside
% at least one `\itemize`/`\enumerate` when this one opens) -- a top-level
% enumerate (`\@itemdepth=1`, e.g.
% slides/t9z-Outland-Steel-Hedging-Borrowing-Investing-current.tex/
% slides/g4c-Outland-Steel-Hedging-Forward-Sale-current.tex/
% slides/j84-Outland-Steel-No-Hedging-current.tex's own outermost numbered
% list, or slides/nkj-Agenda-current.tex's own \tableofcontents-driven one)
% is left at beamer's stock label position entirely untouched, matching
% this project's own established behavior for that shape (never reported
% as misaligned, unlike the nested case) -- and, for nkj specifically,
% doesn't matter either way, since that slide's own
% `\SMTocLabelFlushLeft` call (right after its own `\begin{enumerate}`)
% overrides `\makelabel` again regardless of whatever this default sets.
%
% Reuses the exact same 20.33pt calibration as the manual
% `\SMEnumerateLabelWidth` above, confirmed (pymupdf, compiled PDF) to
% land correctly not just at `\@itemdepth=2` (one `\itemize` above the
% enumerate -- slides/kfh-Asset-Allocation-Process-current.tex/
% slides/f6m-Recovery-Rates-current.tex/
% slides/6pz-Outland-Steel-Hedging-Strategies-Overview-current.tex, all
% previously needing the manual call) but also at `\@itemdepth=3` (two
% `\itemize`s above it -- slides/bdh-Working-Capital-Example-current.tex's
% own "Development of balance sheet:" list, also previously needing the
% manual call, with the identical 20.33pt value) -- this deck's itemize/
% enumerate templates don't shrink font size by nesting depth (only the
% bullet glyph/margin step does), so the numeral's own glyph width -- the
% only thing \SMEnumerateLabelWidth's own box width actually
% compensates for -- stays the same at every depth, and one constant covers
% both cases without needing a depth-dependent formula. Not verified beyond
% depth 3 (no slide currently nests an enumerate that deep); if one ever
% does, re-measure before trusting this constant that far down, same as any
% other empirical calibration in this file.
% (slides/scs-Political-Risk-Stability-Index-current.tex also matches
% "`\SMEnumerateLabelWidth`" in a grep -- but only inside its own prose
% comments describing the general technique; its actual call is the
% bespoke `\SMEnumerateLabelFlushLeftShifted{16pt}{10.2pt}` for its own
% alphabetic, left-aligned, narrower-box labels, unaffected by this
% override either way since that call runs after `\begin{enumerate}` opens,
% same ordering as every other post-open `\makelabel` override here.)
%
% Former per-slide manual `\SMEnumerateLabelWidth` calls removed from
% the 4 slides above now that this covers them automatically; the macro
% itself is kept defined, both for this override to reuse and as a manual
% escape hatch (e.g. a differently-calibrated width some future slide might
% need).
% (v9.8.0: \SMEnumerateTextExtra, 6.3pt, is a hook of the catalog; the four lengths below were numbers of the
% layout profile: 7.85pt, 13.1pt, 5.25pt and 20.33pt for Fira Sans, 8.18pt, 13.43pt, 5.25pt and 20.33pt for
% Latin Modern Sans.)
% The numeral column of a nested enumerate, from the widest numeral of the list's own font (the digits
% 1-9 with the dot, which is what the item template sets). The slot a numeral sits in is that width plus
% \SMEnumerateNestedSlotExtra, the text starts \SMEnumerateNestedLabelGap after it
% (\SMEnumerateNestedLeftMargin is both), and the box the numeral is right-aligned in only has to hold a
% numeral of three digits ("999."), so it takes the width of that. Called first in the list's setup of
% every enumerate, and once at \begin{document} (in \small) for a slide that reads the lengths itself.
\makeatletter
\newcommand{\SM@enumeratemetrics}{%
	\setlength{\SMEnumerateNestedMinNumber}{0pt}%
	\setlength{\SMEnumerateLabelBox}{0pt}%
	\@for\SM@numprobe:=1,2,3,4,5,6,7,8,9\do{%
		\settowidth{\dimen@}{\SM@numprobe.}%
		\ifdim\dimen@>\SMEnumerateNestedMinNumber\SMEnumerateNestedMinNumber=\dimen@\fi
		\settowidth{\dimen@}{\SM@numprobe\SM@numprobe\SM@numprobe.}%
		\ifdim\dimen@>\SMEnumerateLabelBox\SMEnumerateLabelBox=\dimen@\fi}%
	\addtolength{\SMEnumerateNestedMinNumber}{\SMEnumerateNestedSlotExtra}%
	\setlength{\SMEnumerateNestedLabelSep}{\SMEnumerateNestedLabelGap}%
	\setlength{\SMEnumerateNestedLeftMargin}{\dimexpr\SMEnumerateNestedMinNumber+\SMEnumerateNestedLabelSep\relax}}
\makeatother
\makeatletter
\renewcommand{\beamer@enum@}{%
	\beamer@computepref\@itemdepth% sets \beameritemnestingprefix
	\usebeamerfont{itemize/enumerate \beameritemnestingprefix body}%
	\usebeamercolor[fg]{itemize/enumerate \beameritemnestingprefix body}%
	\usebeamertemplate{itemize/enumerate \beameritemnestingprefix body begin}%
	% Level-1 enumerate text sits further right than bullet text, as in the source
	% (numeral at the bullet dot's left edge, item text 27.9pt from the page's left edge
	% against 21.6pt for a bullet's text; 2026-09-20). Nested lists follow it, as they do for
	% plain bullets: a level-2 bullet starts under the level-1 text (2026-09-21, direct
	% instruction -- an earlier version gave the extra back to the nested list, which left
	% level-2 bullets 6.3pt left of the level-1 text above them). Local to this list.
	\ifnum\@itemdepth=1
		\addtolength{\leftmargini}{\SMEnumerateTextExtra}%
	\fi
	\expandafter
		\list
			{\usebeamertemplate{\beamer@enumtempl}}
			{\usecounter\@enumctr%
				\SM@enumeratemetrics% the numeral column from this list's font (v9.8.0)
				\ifnum\@itemdepth>1
					% Nested numerals: same rule as level 1 -- left edge at the left
					% edge of this level's bullet dot (2026-09-20). The box below is
					% the old 20.33pt right-aligned one, whose right end the numerals
					% used to fill flush; giving the numeral a 7.85pt left-aligned
					% slot at that end puts every numeral's left edge exactly at the
					% dot's, instead of up to 1pt right of it for narrow digits
					% ("1." vs "4."). A wider numeral keeps the old right alignment.
					% Hanging indent (2026-09-21, direct instruction: the second line of
					% an item did not start under its first line's text). The 20.33pt label
					% box is wider than beamer's \labelwidth, so it pushed the first
					% line's text 4.3pt right of \leftmargin, where every later line
					% starts. Sizing \labelwidth to the box and \leftmargin to
					% "numeral slot + \labelsep" puts the numeral's left edge at the dot's
					% and the text of every line at the same x (13.1pt from the parent's
					% text, i.e. where the first line always was).
					\setlength{\labelwidth}{\SMEnumerateLabelBox}%
					\setlength{\labelsep}{\SMEnumerateNestedLabelSep}%
					\setlength{\leftmargin}{\SMEnumerateNestedLeftMargin}%
					\def\makelabel##1{\makebox[\SMEnumerateLabelBox][r]{{%
							\setbox\@tempboxa\hbox{##1}%
							\ifdim\wd\@tempboxa>\SMEnumerateNestedMinNumber
								\box\@tempboxa
							\else
								\hbox to \SMEnumerateNestedMinNumber{\unhbox\@tempboxa\hss}%
							\fi}}}%
				\else
					% Level-1 numerals start at the left edge of the level-1 bullet
					% dot, like the source's numerals (2026-09-20, direct
					% instruction: the numbers were too far left). They used to
					% be right-aligned to the dot's right edge, which for a
					% numeral wider than the dot (every "1.") puts its left edge
					% left of where a bullet would start. A numeral too wide to
					% keep 1pt clear of the item text ("10.") keeps the old
					% right alignment instead of running into the text.
					\def\makelabel##1{{\hss\llap{{%
							\usebeamerfont*{enumerate \beameritemnestingprefix item}%
							\usebeamercolor[fg]{enumerate \beameritemnestingprefix item}%
							\setbox\@tempboxa\hbox{##1}%
							\ifdim\dimexpr\wd\@tempboxa+1pt\relax>\dimexpr\SMItemStep+\SMEnumerateTextExtra\relax
								\box\@tempboxa
							\else
								\hbox to \dimexpr\SMBulletWidth+\SMEnumerateTextExtra\relax{\rlap{\unhbox\@tempboxa}\hss}%
							\fi}}}}%
				\fi}%
	\beamer@cramped%
	\SMNoListParsep
	\raggedright%
	\beamer@firstlineitemizeunskip%
}
% The lists' lengths as a slide that reads them outside a list sees them: those of the \small font.
\AtBeginDocument{%
	\begingroup\small\SM@labelitemmetrics\SM@enumeratemetrics
	\xdef\SM@listmetrics{%
		\global\SMLabelItemWidth=\the\SMLabelItemWidth\relax
		\global\SMLabelItemSep=\the\SMLabelItemSep\relax
		\global\SMEnumerateNestedMinNumber=\the\SMEnumerateNestedMinNumber\relax
		\global\SMEnumerateNestedLabelSep=\the\SMEnumerateNestedLabelSep\relax
		\global\SMEnumerateNestedLeftMargin=\the\SMEnumerateNestedLeftMargin\relax
		\global\SMEnumerateLabelBox=\the\SMEnumerateLabelBox\relax}%
	\expandafter\endgroup\SM@listmetrics}
\makeatother

% Extra breathing room when a nested itemize closes and control returns to
% its *parent* list, right before that parent's next \item -- e.g. sub-items
% under "Mid-Term Exam" followed by the next top-level "Final Exam" -- so a
% new top-level group reads as visually distinct from the one before it,
% rather than the same uniform gap used between every other item regardless
% of level. \@itemdepth (LaTeX's own nesting-level counter, already
% decremented back to the parent's level by the time \end{itemize}'s normal
% closing has run) says which level control returned to -- 0 means the
% outermost list itself just closed (nothing left to add space before,
% deliberately excluded), 1 means back to a level-1 item, 2 means back to a
% level-2 item.
%
% Level-dependent (2026-07-06, direct instruction) -- a flat, single
% constant regardless of *which* level control returned to (the original
% version of this hook) reads as "extra"/inconsistent precisely because a
% return-to-level-1 and a return-to-level-2 transition both got the exact
% same gap, rather than level 1 getting a visibly bigger gap than level 2 --
% found on slides/57r-Reasons-for-Leasing-current.tex, where a level-2
% sub-list closing back to a level-1 item looked like a random extra gap
% rather than "level 1 items are consistently more separated than level 2
% ones." \addvspace (not plain \vspace) matters too: it coalesces with
% whatever glue itemize's own \itemsep/\topsep already contributes at the
% closing boundary (taking whichever is larger) rather than *adding* on top
% of it unconditionally -- \vspace's own stacking behavior was the other
% half of why the old flat version could look like too much space.
\makeatletter
% Ratios, not lengths (2026-07-21, direct instruction: "all of these
% spacings should scale with font size") -- a \newlength+\setlength{...}{0.5em}
% pair, the way this used to be written, freezes to a fixed *pt* value the
% one time \setlength runs (at preamble-load time, using whatever font is
% active then), despite the "em" in its own definition -- it does NOT
% re-track whichever font size is actually active on a given slide/level
% later. Confirmed directly: a \typeout{\f@size} placed at the point each
% of these actually gets used (inside \AtEndEnvironment{itemize}'s own
% hook body, further below) showed the *closing* list's own font still
% active there (9pt/10pt/10.95pt for a level-4/3/2 list closing,
% respectively -- not yet restored to the parent's own font), confirming
% this hook fires early enough in `\end{itemize}`'s own processing that a
% live "0.5em" computed *there* would correctly track whichever level is
% actually involved. Fixed by storing the bare ratio as a plain macro
% instead, and writing the literal "em" back on at the point of use
% (`\addvspace{\SMItemizeLevelDropSpaceRatioI em}`), so it's
% recomputed fresh, against whatever font is live, every time.
% Level 3's own drop-space, for when a closed *4th*-level list returns
% control to a level-3 item -- added 2026-07-15 alongside the 4th itemize
% level itself (search "subsubsub" further above). Interpolated from the
% existing I/II progression (0.5, 0.3 -- each step roughly 60% of the one
% before) rather than empirically re-measured against a compiled PDF the
% way I/II themselves were: no slide used a 4th level before
% slides/hrg-Exchange-Rate-Quotation-Forward-current.tex to measure
% against. Re-derive properly (render, measure, adjust) if this ever
% looks visibly off on a real slide.
% Level 4's own drop-space, for when a closed level-4 list returns control
% to a level-3 item -- added 2026-07-21 alongside \SMItemsepRatioIV.
% \@itemdepth at this hook's fire time is the depth of the list that *just
% closed*, confirmed directly via a temporary \typeout{\the\@itemdepth}
% (not assumed): closing the level-4 list on
% slides/ZZO-Itemize-Spacing-Reference-current.tex measured itemdepth=4
% there, closing level-3 measured 3, level-2 measured 2, and the outermost
% level-1 list itself closing measured 1 (never 0 -- \@itemdepth's own
% restore to the parent's depth happens *after* this hook runs, so the "0:
% outermost list itself closed" case below is dead code, kept only because
% removing a genuinely unreachable \ifcase branch isn't worth the risk of
% getting the numbering wrong a second time). So this is very much a real,
% reachable case -- not decorative symmetry -- and needed its own
% \addvspace branch (previously missing entirely, silently leaving the
% level-4-to-level-3 closing transition with no drop-space at all).
% Interpolated from the existing I/II/III progression as a starting guess,
% then re-measured directly against the compiled ZZO reference and bumped:
% 0.13 measured at 12.6pt, still *less* than level 4's own same-level gap
% of 13.4pt (violates rule III -- a level's own same-level gap must be
% smaller than that level's own closing-transition gap); raised until
% re-measurement cleared it with a comfortable margin again, matching the
% margin every other level's own pair already has.
\AtEndEnvironment{itemize}{%
	% Force vertical mode before \addvspace below (2026-09-22, found while
	% isolating a "Something's wrong--perhaps a missing \item" error that
	% fired at the *frame's* own \end{...}, not the itemize's, on every
	% slide using a plain, single-level \begin{itemize}...\end{itemize}
	% with nothing else in the frame -- confirmed via a from-scratch minimal
	% reproduction outside this project's own slides entirely, then
	% bisected to this exact hook by diffing against the same reproduction
	% with this \AtEndEnvironment call disabled). \AtEndEnvironment inserts
	% this code at the very *start* of \enditemize, before beamer's own
	% `\ifhmode\unskip\fi\endlist` -- so without this guard, \addvspace
	% below runs while TeX is often still in horizontal mode (right after
	% the last \item's own text, mid-paragraph). \addvspace's own hmode
	% branch defers via `\vadjust{\addvspace{...}}` instead of inserting
	% the glue directly, and that deferred insertion colliding with
	% beamer's own frame-body/overlay-replay collection is what actually
	% produced the corruption -- not reproducible as a plain "wrong
	% \item" typo anywhere in the affected slides' own source, only via
	% this exact hmode-vs-vmode timing. \ifhmode\unskip\par\fi here forces
	% vertical mode *before* \addvspace ever runs, matching the standard,
	% documented way to call \addvspace (kernel doc: "should only be used
	% between paragraphs") -- same idiom beamer's own \enditemize already
	% uses \unskip for, just moved earlier so it also covers this hook.
	% Doesn't disturb the \@itemdepth-based \ifcase dispatch just below:
	% \par doesn't touch \@itemdepth, so the closing list's own depth is
	% still read correctly afterward, same as before this guard existed.
	\ifhmode\unskip\par\fi
	\ifcase\@itemdepth
		% 0: dead code -- @itemdepth is never actually 0 when this hook
		% fires (see \SMItemizeLevelDropSpaceRatioIV's own comment
		% just above); kept only for alignment with the \ifcase branches
		% below, each keyed on the depth of the list that just closed.
	\or\addvspace{\SMItemizeLevelDropSpaceRatioI em}% level-1 (outermost) list itself closed -- a pre-existing quirk (not introduced by this comment's own edit): there's no parent item above level 1 for this to visually separate from, so this only actually matters if something follows the outermost list directly on the same frame; harmless on every slide seen so far since nothing does. Out of scope to redesign here.
	\or\addvspace{\SMItemizeLevelDropSpaceRatioII em}% level-2 list closed, back to a level-1 item
	\or\addvspace{\SMItemizeLevelDropSpaceRatioIII em}% level-3 list closed, back to a level-2 item
	\or\addvspace{\SMItemizeLevelDropSpaceRatioIV em}% level-4 list closed, back to a level-3 item
	\fi
}

% Per-level \itemsep -- the item-to-item gap *within* one list, at one
% fixed nesting depth (as opposed to the drop-space mechanism just above,
% which only fires when a nested list *closes* back up to its parent).
% Ratios, not lengths -- same "scale with font size" reasoning and fix as
% \SMItemizeLevelDropSpaceRatioI/II/III above (that comment has the
% full story on why a plain \newlength+\setlength pair here would freeze
% to one fixed pt value instead of actually tracking each level's own
% font size). Values only, defined here for discoverability alongside the
% drop-space ones above -- actually *applied* inside the `\itemize`
% redefinition itself (search "\beamer@cramped (beamerbaselocalstructure.sty"
% further above), not via a separate hook here; see that comment for why.
% Found missing entirely (2026-07-21, direct instruction, see
% slides/ZZO-Itemize-Spacing-Reference-current.tex's own header comment
% for the rules this and the drop-space mechanism together need to
% satisfy): this deck's own itemize never set \itemsep itself at any
% level -- confirmed directly, grepping this whole file -- so every depth
% shared whatever bare default beamer's own theme happens to leave
% \itemsep at, uniformly. That produced two visible problems, both
% measured directly (pymupdf against slides/2c2-Topics-Around-the-CAPM-Empirical-Validity-current.tex
% and an isolated bisection test, not guessed): (a) level-2 and level-3
% same-level item-to-item gaps came out *identical* (12.6pt each,
% confirmed against an isolated reproduction), since nothing distinguished
% them by depth at all -- only a closing transition ever got a
% level-dependent gap, never a same-level one; (b) that shared default
% \itemsep read as *smaller* than a single item's own wrapped-line
% \baselineskip (13.7-13.8pt within a wrap vs. 12.6pt between items),
% backwards from how a list should read -- an item boundary should never
% be tighter than a same-item line wrap.
% Ratios themselves are not arbitrary -- each chosen, then re-measured
% against the compiled ZZO reference slide and adjusted, specifically to
% satisfy all three of that slide's own header-comment rules together
% (extended to level 4 the same way, 2026-07-21), not just "level 1
% biggest" in isolation:
% - Rule I (level-1 same-level gap > level-2's > level-3's > level-4's)
%   falls straight out of \SMItemsepRatioI > ...II > ...III > ...IV,
%   added on top of the same \baselineskip every level shares (itself
%   already smaller per level here, since this theme's own itemize/
%   enumerate body font shrinks with depth too -- confirmed directly,
%   10.95pt/10pt/9pt/9pt for levels 1-4 respectively, level 4 sharing
%   level 3's own font since this theme has no distinct "subsubsub" size
%   of its own).
% - Rule II (every level's own same-level gap > that level's own wrapped-
%   line gap) needs \SMItemsepRatioIV alone (the smallest, hence the
%   binding case) to still be positive and large enough that it clearly
%   reads as more open than a plain wrapped line once rendered, not just
%   bigger on paper.
% - Rule III (a level's own same-level gap < that same level's own
%   *closing*-transition gap) is *not* automatic just because
%   \SMItemizeLevelDropSpaceRatio* was already bigger than the old
%   shared default -- \addvspace's own coalesce-to-the-larger-value
%   behavior (documented above) means the closing gap is
%   max(natural glue, that *drop-space* ratio), not a sum of the two. Each
%   drop-space ratio's own \ifcase branch is keyed on the depth of the
%   list that *just closed* (see \SMItemizeLevelDropSpaceRatioIV's own
%   comment for the full story on why, confirmed via \typeout, not
%   assumed), so it lines up with the same-numbered itemsep ratio despite
%   describing a different transition -- \SMItemizeLevelDropSpaceRatioIV
%   governs the level-4-list-closing-back-to-level-3 transition and is the
%   one that has to individually stay bigger than \SMItemsepRatioIV.
%   This branch was missing entirely at first (only I/II/III existed, a
%   leftover from before level 4 existed) -- confirmed directly on ZZO: the
%   level-4-to-level-3 closing transition measured 12.6pt with no drop-space
%   branch at all, less than level 4's own same-level gap of 13.4pt, a real
%   rule III violation. Adding the branch (\SMItemizeLevelDropSpaceRatioIV,
%   starting at 0.5) fixed it. Re-verify rule III directly against ZZO's own
%   compiled PDF (don't assume it still holds on paper) if any of these ratios is
%   ever changed again.

% Per-level "opening" space -- the gap between a level-n item and the
% first item of a level-(n+1) list that immediately follows it, as
% opposed to \SMItemizeLevelDropSpaceRatio* (a nested list *closing*
% back up) or \SMItemsepRatio* (two items at the *same* depth).
% Added 2026-07-21, direct instruction, replacing the original three
% separate rules on slides/ZZO-Itemize-Spacing-Reference-current.tex's
% own header comment with one multi-part inequality stated per level n:
%   level-n -> level-(n-1) close  >  level-n same-level  >
%   level-n -> level-(n+1) open  >  level-n's own wrap gap
% The "open" term was entirely unhandled before this -- confirmed
% directly on ZZO: every opening transition (level 1 into level 2, level
% 2 into level 3, level 3 into level 4) measured *smaller* than that same
% level's own wrap gap, backwards from what the rule needs, since nothing
% but beamer's own bare default \topsep governed it.
% \AtBeginEnvironment{itemize}, not \topsep -- \topsep was tried first
% (set via the same \ifcase\@itemdepth block, right after \itemsep,
% inside the `\itemize` redefinition above) and measured *unreliable*:
% two of the three levels barely moved despite a sizeable ratio, while
% the third moved by roughly the right amount -- \topsep's own value is
% evidently consulted by `\list` at some point this override comes too
% late for, at least for some levels, not evenly. \addvspace via
% \AtBeginEnvironment instead measured consistently and predictably
% (each level's own change matched its own ratio's calculated pt value
% directly), the same mechanism \SMItemizeLevelDropSpaceRatio* already
% relies on for the closing side. This hook fires with \@itemdepth
% reflecting the *parent* item's own depth (confirmed via \typeout, not
% assumed: opening the outermost level-1 list itself measured 0, opening
% a level-2 list from within a level-1 item measured 1, level-3-from-
% level-2 measured 2, level-4-from-level-3 measured 3) -- the opposite of
% \SMItemizeLevelDropSpaceRatio*'s own hook, which keys on the depth
% of the list that *just closed*. That makes the \ifcase numbering here
% match the "n" in the rule directly: case 1 is level-1's own open-to-
% level-2 ratio, case 2 is level-2's own open-to-level-3, case 3 is
% level-3's own open-to-level-4; case 0 (opening the outermost list
% itself) is a genuine no-op, correctly -- there's no parent item above
% level 1 for that transition to separate from.
% Ratios chosen, then re-measured against the compiled ZZO reference and
% adjusted, the same way as every other ratio here: each has to
% individually land strictly between that level's own \SMItemsepRatio*
% (the ceiling) and its own wrap gap (the floor). Re-verify directly
% against ZZO's own compiled PDF (don't assume it still holds on paper)
% if any of these, or the itemsep/drop-space ratios they're measured
% against, are ever changed again.
% \SMItemizeLevelOpenSpaceRatioI trimmed from 0.28 to 0.02 (2026-07-21,
% direct instruction: "the level-1 to level-2 distance...seems a bit
% much"). Confirmed directly: even at 0, this transition still measured
% ~16.3pt -- \addvspace's own coalesce-to-the-larger-value behavior means
% the ratio can only ever push this *above* whatever glue the list
% mechanism inserts here regardless (that natural floor, not the ratio),
% so 0.28 down to 0 only ever moved the actual gap by ~2pt (18.5->16.3);
% 0.02 splits the difference, keeping the mechanism nominally "in effect"
% rather than reading as accidentally disabled. \SMItemizeLevelOpenSpaceRatioIII
% (level-3-to-level-4) had to come down too, from 0.08 to 0.02, purely to
% restore its own margin -- changing RatioI shifted that unrelated
% transition's own measured gap closer to level 3's own same-level gap
% (12.1 vs. 12.0, too thin), confirmed directly on ZZO, not assumed;
% ratios elsewhere in this file have shown the same small-ripple behavior
% before (see \SMItemizeLevelDropSpaceRatioIV's own comment) --
% re-measure the *whole* table after touching any single one of these,
% not just the transition being changed.
\AtBeginEnvironment{itemize}{%
	% Same \ifhmode\unskip\par\fi guard as \AtEndEnvironment{itemize} above,
	% and for the identical reason -- see that hook's own comment for the
	% full story. This hook fires when a *nested* itemize opens (\@itemdepth
	% here is the *parent* item's depth, per the comment below), i.e. right
	% after "\item ...text..." mid-paragraph, immediately before the
	% \addvspace calls further down -- the same hmode-vs-vmode timing that
	% broke the closing hook, just on the opening side. Case 0 (opening the
	% outermost level-1 list itself) never reaches an \addvspace call below
	% regardless, so a plain single-level itemize was never affected by
	% this half of the bug -- only a *nested* \begin{itemize} (depth >=1 at
	% open time) actually needed this guard, which is why the "missing
	% \item" error survived this fix's own first pass specifically on the
	% two demo slides with a nested itemize (sectionframe's own ZZR,
	% standardframe's own ZZZ) while every plain single-level-itemize slide
	% was already fixed by the closing-hook guard alone.
	\ifhmode\unskip\par\fi
	\ifcase\@itemdepth
		% 0: opening the outermost (level-1) list itself -- no parent item
		% precedes it, nothing to add.
	\or\addvspace{\SMItemizeLevelOpenSpaceRatioI em}% opening a level-2 list from a level-1 item
	\or\addvspace{\SMItemizeLevelOpenSpaceRatioII em}% opening a level-3 list from a level-2 item
	\or\addvspace{\SMItemizeLevelOpenSpaceRatioIII em}% opening a level-4 list from a level-3 item
	\fi
}
\makeatother


% Even with the plain-enumerate template (SMPlainEnumerate, defined in the theme), \tableofcontents still
% doesn't match this deck's own normal itemize spacing -- beamer's own
% \beamer@sectionintoc (beamerbasetoc.sty), which wraps *every* section
% entry regardless of which "section in toc" template is active,
% unconditionally inserts a fixed \vskip1.5em before each entry (an
% unbounded \vfill, even worse, if subsections aren't hidden -- neither is
% reachable via a public \tableofcontents option, just this show/hide
% binary), *and* wraps the entry itself in \hbox{\vbox{...}} -- box
% wrapping that (like \SMMinipageTopCorrection's own minipage case
% elsewhere in this file) changes how surrounding baselineskip/topskip
% glue gets inserted compared to plain paragraph text. Both together added
% roughly 13pt of extra gap
% between entries beyond this deck's own normal ~18.6pt level-1 item
% spacing (measured directly against
% slides/57r-Reasons-for-Leasing-current.tex's own compiled PDF).
% Neither is controllable via a public \tableofcontents option, so this
% redefines \beamer@sectionintoc itself -- scoped safely since
% \tableofcontents is only ever called on the one Agenda slide in this
% deck -- keeping its numbering/hyperlink/part-number/hide-checking logic
% intact but dropping the vskip/vfill and the hbox/vbox wrapping, so an
% entry behaves exactly like a plain \item once printed.
\makeatletter
\def\beamer@sectionintoc#1#2#3#4#5{%
	\ifnum\c@tocdepth>0%
	\ifnum#4=\beamer@showpartnumber%
	{%
	\beamer@saveanother%
	\gdef\beamer@todo{}%
	\beamer@slideinframe=#1\relax%
	\expandafter\only\beamer@tocsections{\gdef\beamer@todo{%
			\beamer@tempcount=#5\relax%
			\advance\beamer@tempcount by\beamer@sectionadjust%
			\ifnum\beamer@tempcount>0
				\ifnum\beamer@tempcount>\beamer@toclastsection
					\def\inserttocsectionnumber{}%
				\else
					\edef\inserttocsectionnumber{\the\beamer@tempcount}%
				\fi
			\else
				\def\inserttocsectionnumber{}%
			\fi%
			\def\inserttocsection{\hyperlink{Navigation#3}{#2}}%
			\beamer@tocifnothide{\ifnum\c@section=#1\beamer@toc@cs\else\beamer@toc@os\fi}%
			{%
				\beamer@tocact{\ifnum\c@section=#1\beamer@toc@cs\else\beamer@toc@os\fi}{section in toc}%
			}%
		}%
	}%
	\beamer@restoreanother%
	}%
	\beamer@todo%
	\fi\fi%
}
\makeatother

% Left-aligns an enumerate's labels within a caller-chosen \labelwidth --
% used by the \tableofcontents Agenda slide (see
% slides/nkj-Agenda-current.tex's own comment for why its labels need
% this at 3.6pt). Defined here as a named helper rather than written as a
% literal \renewcommand{\makelabel}[1]{...#1...} inside the slide's own
% frame body, because that literal form -- specifically its raw `#1`
% parameter token -- breaks inside any frame that \tableofcontents forces
% through beamer's overlay reprocessing: beamer replays such a frame's
% collected body via \def-based machinery (\loop's own \iterate, and
% \beamer@doifinframe), and a `#` in the material of any \def must be
% doubled -- so the compile hits "! Illegal parameter number in definition
% of \iterate." (twice) plus the same for \beamer@doifinframe, once per
% compile, at that frame's own \end{frame}. TeX's error recovery ("I'll
% assume you meant ##") happens to produce the intended definition anyway,
% which is why the rendered PDF was always correct despite the errors --
% but lualatex still exits nonzero, which is exactly what kept breaking
% latexmk's automated lualatex/biber/rerun sequencing on the main deck
% (the "Known open issue" long tracked in CONVERSION_NOTES.md, originally
% misattributed to a biblatex interaction -- see its "Resolved" entry).
% Hoisting the definition out here means the frame body carries only this
% macro's *call* token, with no `#` in it, which sidesteps the whole
% mechanism. The inner \makelabel parameter is written `##1` below since
% it sits inside another \newcommand's body -- standard nesting, and also,
% incidentally, exactly the doubling rule the inline version fell foul of.
\newcommand{\SMTocLabelFlushLeft}[1]{%
	\setlength{\labelwidth}{#1}%
	\renewcommand{\makelabel}[1]{\makebox[\labelwidth][l]{##1}}%
}

% Same idea as \SMTocLabelFlushLeft above, plus an extra fixed
% \hspace shift before the label (for a list that needs to start further
% right than the plain nested-list indent, e.g. to line up with a
% sibling item's own text rather than its bullet -- see
% slides/scs-Political-Risk-Stability-Index-current.tex's own v1-5-0/
% v1-6-0 comments for the worked example this was hoisted out of).
% Found 2026-07-20 that the raw-`#1`-in-frame-body bug documented above
% is *not* actually \tableofcontents-specific, contrary to that comment's
% own framing -- slides/scs-Political-Risk-Stability-Index-current.tex has
% no \tableofcontents anywhere in the deck-wide overlay chain leading to
% it, yet hit the identical "Illegal parameter number in definition of
% \iterate" at its own \end{standardframe} from the exact same literal-
% \renewcommand{\makelabel}[1]{...#1...}-in-frame-body pattern. Whatever
% the precise trigger condition actually is, the same hoist-to-preamble
% fix (a bare macro *call*, no `#` token, left in the frame body) sidesteps
% it either way.
\newcommand{\SMEnumerateLabelFlushLeftShifted}[2]{%
	\setlength{\labelwidth}{#2}%
	\renewcommand{\makelabel}[1]{\hspace{#1}\makebox[\labelwidth][l]{##1}}%
}

% Width of a single digit in this deck's normal content font, measured once
% here (safe, plain preamble-load-time \settowidth) rather than inline in
% the Agenda slide's own frame body -- \tableofcontents there puts the whole
% frame through beamer's overlay-replay machinery, which the comment above
% \SMTocLabelFlushLeft already found chokes on raw tokens (like a bare
% \settowidth call) sitting directly in a replayed frame body, the same
% "Missing number"/"Illegal unit of measure" class of error as that `#`
% problem, confirmed directly the same way. Used by
% slides/nkj-Agenda-current.tex's own "widen the number-to-title gap"
% fix, which only ever needs to *reference* this precomputed value inside
% the frame, not measure anything itself there.
\newlength{\SMTocLabelExtraGap}
\settowidth{\SMTocLabelExtraGap}{0}

% For a two-column layout (e.g. bullets next to a chart/formula), nest
% beamer's own "columns" directly inside standardframe -- do NOT bake
% \begin{columns}/\end{columns} into a custom environment's begin/end code:
% beamer collects each frame's body as a token list to replay it once per
% overlay, and a columns environment that is opened in the "before" half and
% closed in the "after" half of a *different* environment confuses that
% collection (it never finds a properly matched \end{frame}). Nesting it in
% the document body, as below, avoids the issue entirely:
%
%   \begin{standardframe}{Title}
%   	\begin{columns}[t,onlytextwidth]
%   		\begin{column}{0.48\linewidth} ... \end{column}
%   		\begin{column}{0.48\linewidth} ... \end{column}
%   	\end{columns}
%   \end{standardframe}

% For automatic section dividers and single-statement "standout" slides, use
% what the gotham theme already provides instead of a custom environment:
%   \section{...}                                                (section divider, automatic)
%   \begin{frame}[standout]{Key takeaway}\end{frame}             (big single statement)
%   \begin{block}{...}/\begin{alertblock}{...}/\begin{exampleblock}{...}/\begin{referenceblock}{...}  (definitions, formulas, examples, references)
% For a course's own title slide, section-topic overview, quote, figure-only,
% and presenter-attention slides, see the five environments below.

% The block templates (block, alerted, example, reference, info and definition begin and end), their
% colors and the helpers they use (\SMBlockSideMargin, \SMBlockTitleWithMarker, \SMSetBlockBullets,
% \SMAlertMarker) are the theme's (beamerthemeSlideMage.sty); the three block types the beamer
% environments do not have are defined here.
\newenvironment<>{referenceblock}[1]{%
	\begin{actionenv}#2%
		\def\insertblocktitle{#1}%
		\par%
		\usebeamertemplate{block reference begin}}
	{\par%
		\usebeamertemplate{block reference end}%
	\end{actionenv}}

\newenvironment<>{infoblock}[1]{%
	\begin{actionenv}#2%
		\def\insertblocktitle{#1}%
		\par%
		\usebeamertemplate{block info begin}}
	{\par%
		\usebeamertemplate{block info end}%
	\end{actionenv}}

\newenvironment<>{definitionblock}[1]{%
	\begin{actionenv}#2%
		\def\insertblocktitle{#1}%
		\par%
		\usebeamertemplate{block definition begin}}
	{\par%
		\usebeamertemplate{block definition end}%
	\end{actionenv}}


% A narrower, freely-positionable version of \begin{block}{...} (2026-07-06,
% direct instruction) -- a plain block always spans the *current*
% \linewidth (that's what block begin/end's own \hspace+minipage wrapper
% above is measured against), so it can't be dropped into a fixed-width
% spot (e.g. beside a table, inside a tikz overlay \node) without first
% narrowing \linewidth some other way. \SMFloatingBlock just wraps a
% plain block in one more, *outer* minipage of the caller's own chosen
% width -- block begin/end's own inner minipage then measures itself
% against that narrower width instead of the frame's full one, nesting
% correctly with no changes needed to the block templates themselves.
% Title/body colors, margins, and font are exactly \block's own (same
% "block title"/"block body" beamer color names) -- this only changes the
% overall width and (by taking it out of the normal document flow, left to
% the caller) the position, not the look. Font size is also left to the
% caller (wrap the whole environment in a size declaration, e.g.
% {\footnotesize \begin{SMFloatingBlock}{5cm}{Excel} ... \end{SMFloatingBlock}})
% rather than fixed here, since the right size depends entirely on how
% much room the block ends up placed in.
%
% INNER PADDING (2026-09-30, direct instruction: systematic replacements for
% per-slide fixes): 20 of 27 floating blocks opened with a hand-written
% \vspace{2mm} and 17 closed with \par\vspace{1.2mm} (or 1mm/3.4pt), to give
% the body text room from the title bar and the block's bottom edge. All
% three SMFloating*Block environments now do that themselves, via an
% optional key list:
%   \begin{SMFloatingInfoBlock}[top=2mm, bottom=1.2mm]{6cm}{Title}
% Defaults top=2mm, bottom=1.2mm (the common values); 0pt inserts nothing
% at all. Blocks that had no or other padding carry their own values
% explicitly (e.g. [top=0pt]), so every block kept its exact look.
\makeatletter
\pgfqkeys{/SMFloatingBlock}{
	top/.store in=\SM@fb@top,
	bottom/.store in=\SM@fb@bottom,
}
\newcommand{\SM@fb@begin}[1]{%
	\pgfqkeys{/SMFloatingBlock}{top=2mm, bottom=1.2mm, #1}%
	\ifdim\SM@fb@top=0pt \else\vspace{\SM@fb@top}\fi
}
\newcommand{\SM@fb@end}{%
	\ifdim\SM@fb@bottom=0pt \else\par\vspace{\SM@fb@bottom}\fi
}
\NewDocumentEnvironment{SMFloatingBlock}{O{}mm}{%
	\begin{minipage}{#2}
	\begin{block}{#3}
	\SM@fb@begin{#1}%
}{%
	\SM@fb@end
	\end{block}
	\end{minipage}
}
\makeatother

% Same idea, but wrapping \exampleblock instead of \block (2026-07-06,
% direct instruction, first used on
% slides/fr3-Forms-of-Capital-Market-Efficiency-Strong-Semistrong-Weak-current.tex
% and slides/ha9-Monacorp-Truck-Lease-Breakeven-Rent-current.tex, both
% switching an existing \SMFloatingBlock to the "example" green header
% instead of the plain violet one) -- not worth generalizing into a single
% macro parameterized over the block name: \SMFloatingBlock's own body
% is only two lines, so a third variant is cheaper to just write out than
% to build a block-name-dispatch layer for.
\makeatletter
\NewDocumentEnvironment{SMFloatingExampleBlock}{O{}mm}{%
	\begin{minipage}{#2}
	\begin{exampleblock}{#3}
	\SM@fb@begin{#1}%
}{%
	\SM@fb@end
	\end{exampleblock}
	\end{minipage}
}
\makeatother

% Same idea again, wrapping \infoblock (2026-07-06, direct instruction,
% first used on
% slides/fr3-Forms-of-Capital-Market-Efficiency-Strong-Semistrong-Weak-current.tex's
% own "Random walk" formula box, switched from the green
% \SMFloatingExampleBlock to this blue one the same day).
\makeatletter
\NewDocumentEnvironment{SMFloatingInfoBlock}{O{}mm}{%
	\begin{minipage}{#2}
	\begin{infoblock}{#3}
	\SM@fb@begin{#1}%
}{%
	\SM@fb@end
	\end{infoblock}
	\end{minipage}
}
\makeatother

% Every SMFloatingInfoBlock anchored top-right against the institute
% logo (the deck's standard
% "(current page.north east)+(-\SMFrametitleRightPadding,-\SMContentTopAlign)"
% corner pattern) lands its own visible right edge short of the logo's,
% by a fixed amount independent of the block's own chosen width (#1) --
% \infoblock's underlying template insets its colored fill by
% \SMBlockSideMargin from *both* sides of whatever width it's given
% (that margin exists to match a first-level bullet's own left-hand
% position, see its own comment above; the fact that it also narrows the
% *right* side is just a side effect of reusing one symmetric
% hspace+minipage trick for both, not a deliberate right-margin choice),
% partly offset by beamer's own beamercolorbox "colsep*" background bleed
% (colsep* pads the colored fill outward past its own nominal box edge
% without also insetting the text, unlike unstarred "colsep" -- see
% beamerbasecolor.sty's own \beamercolorbox definition). Root-caused and
% measured directly (pymupdf against a compiled probe deck, not derived
% purely on paper) 2026-07-21, after
% slides/j8y-Asset-Allocation-Example-Weights-current.tex's own one-off
% "+2.2pt" nudge (measured that same day against just the institute
% logo's yellow chip) turned out to disagree with
% slides/fr3-Forms-of-Capital-Market-Efficiency-Strong-Semistrong-Weak-current.tex's
% own older "+3.165pt" nudge (2026-07-06, measured against the logo's
% *full* gray-bordered footprint, shadow/card included) -- both slides'
% own nudges were each individually correct for their own chosen
% reference edge, they just didn't agree on what "the logo's own edge"
% means. Direct instruction, 2026-07-21: the *full* gray-bordered
% footprint is the deck's intended reference, so 3.165pt is the value
% kept here; j8y's own newer 2.2pt nudge was the one revised to match,
% not this one.
%
% Use "-\SMFloatingInfoblockRightAlign" (not the bare
% "-\SMFrametitleRightPadding") for a SMFloatingInfoBlock's own
% anchor node whenever its right edge is meant to align with the logo --
% every current use of this pattern across the deck has been updated to
% match; new slides using this same corner pattern should use this
% length from the start rather than re-measuring their own one-off nudge.
\newlength{\SMFloatingInfoblockRightAlign}
\setlength{\SMFloatingInfoblockRightAlign}{\dimexpr\SMFrametitleRightPadding-3.165pt\relax}

% An infoblock's colored title bar starts 6.02pt below the top of its own
% box (the block template's top padding); a floating infoblock anchored at
% the Top Margin therefore sits 6.02pt low unless its anchor is lifted by
% this much. Measured 2026-07-23 (y3a/gf2: title-bar rect at 47.60pt vs.
% the 41.58pt Top Margin), written as a literal "+6.02pt" on 15 anchors
% until 2026-09-30.
\newlength{\SMFloatingInfoblockTopAlign}
\setlength{\SMFloatingInfoblockTopAlign}{6.02pt}

% "Bottom Margin" as a length (2026-09-30): the gap between the footer box
% and content anchored at the bottom of a slide, the mirror image of the Top
% Margin gap between the title bar and the first line (direct instruction on
% bwi, 2026-07-2x: "the same clearance from the footer that the text has from
% the title bar"). Measured from the page's bottom edge, like
% \SMContentTopAlign is from its top edge.
\newlength{\SMContentBottomAlign}
\setlength{\SMContentBottomAlign}{\dimexpr\SMFooterBoxHeight+\SMContentTopAlign-\SMTitleBarBottomAlign\relax}

% \cornerfigure[keys]{image file or content} (2026-09-30, direct instruction:
% systematic replacements for per-slide fixes) -- a figure in a corner of the
% content area, as a zero-space overlay: the construct 105 slides wrote out
% by hand,
%   \vbox to 0pt{\vss\begin{tikzpicture}[overlay, remember picture]
%     \node[anchor=north east, inner sep=0pt] at ($(current page.north east)
%       +(-\SMFrametitleRightPadding,-\SMContentTopAlign)$) {...};
%   \end{tikzpicture}}
% The usual call (2026-10-01, direct instruction: "the width or height of
% the figure has to be specified in the code, that's all that should be
% needed") names the image and one size; the other follows from the aspect
% ratio:
%   \cornerfigure[height=5.6cm]{media/chart-x.png}
%   \cornerfigure[width=4cm, border]{media/photo-x.jpg}
% Keys:
%   width=, height=      the image's size; with either one the argument is an
%                        image file (both: fitted inside, aspect ratio kept).
%                        Without them the argument is arbitrary content (a
%                        block, a \parbox, an \includegraphics with options).
%   border               the yellow image border (SMImageBorder); the
%                        border's outer edge sits on the margins
%   corner=ne|nw|se|sw   which corner (default ne): Top/Bottom Margin
%                        (\SMContentTopAlign/\SMContentBottomAlign)
%                        and Right/Left Margin (\SMFrametitleRightPadding,
%                        \SMFrametitleLeftPadding)
%   corner=e|w           on the Right/Left Margin, centred vertically between
%                        Top and Bottom Margin
%   block                the content is a SMFloatingInfoBlock: lifts it by
%                        \SMFloatingInfoblockTopAlign and aligns its
%                        visible right edge with the logo
%                        (\SMFloatingInfoblockRightAlign)
%   xshift, yshift       a deliberate offset from that corner (right/up +)
%   node={...}           extra TikZ node options
% The node has no inner or outer sep, so the figure's own edge -- or, with a
% border, the border's outer edge (see SMImageBorder) -- is what lands on
% the margins. Like any remember-picture overlay it needs the usual second
% compile. Text that must not run under the figure: besidefigure below.
\makeatletter
\newif\ifSM@cf@block
\newif\ifSM@cf@border
\pgfqkeys{/SMCorner}{
	corner/.store in=\SM@cf@corner,
	xshift/.store in=\SM@cf@x,
	yshift/.store in=\SM@cf@y,
	node/.store in=\SM@cf@node,
	width/.store in=\SM@cf@w,
	height/.store in=\SM@cf@h,
	gap/.store in=\SM@bf@gap,
	block/.code={\SM@cf@blocktrue},
	border/.code={\SM@cf@bordertrue},
}
% all keys back to their defaults, then the caller's
\def\SM@cf@setkeys#1{%
	\SM@cf@blockfalse\SM@cf@borderfalse
	\pgfqkeys{/SMCorner}{corner=ne, xshift=0pt, yshift=0pt, node={},
		width={}, height={}, gap=12pt, #1}%
}
% \SM@cf@content: the figure itself -- \includegraphics[width/height]{#1}
% when a size key was given, #1 unchanged otherwise (\long: content such as
% a floating block may hold paragraphs)
\long\def\SM@cf@makecontent#1{%
	\def\SM@cf@gkeys{}%
	\ifx\SM@cf@w\@empty\else
		\edef\SM@cf@gkeys{width=\unexpanded\expandafter{\SM@cf@w}}%
	\fi
	\ifx\SM@cf@h\@empty\else
		\edef\SM@cf@gkeys{\ifx\SM@cf@gkeys\@empty\else\unexpanded\expandafter{\SM@cf@gkeys},keepaspectratio,\fi
			height=\unexpanded\expandafter{\SM@cf@h}}%
	\fi
	\ifx\SM@cf@gkeys\@empty
		\def\SM@cf@content{#1}%
	\else
		\edef\SM@cf@content{\noexpand\includegraphics[\unexpanded\expandafter{\SM@cf@gkeys}]{\unexpanded{#1}}}%
	\fi
}
\NewDocumentCommand{\cornerfigure}{O{} +m}{%
	\begingroup
	\SM@cf@setkeys{#1}%
	\SM@cf@makecontent{#2}%
	\ifSM@cf@block
		\edef\SM@cf@right{\SMFloatingInfoblockRightAlign}%
		\edef\SM@cf@top{\dimexpr\SMContentTopAlign-\SMFloatingInfoblockTopAlign\relax}%
	\else
		\edef\SM@cf@right{\SMFrametitleRightPadding}%
		\edef\SM@cf@top{\SMContentTopAlign}%
	\fi
	\def\SM@cf@anchor{north east}%
	\edef\SM@cf@at{(current page.north east)+(-\SM@cf@right+\SM@cf@x,-\SM@cf@top+\SM@cf@y)}%
	\ifdefstring{\SM@cf@corner}{nw}{%
		\def\SM@cf@anchor{north west}%
		\edef\SM@cf@at{(current page.north west)+(\SMFrametitleLeftPadding+\SM@cf@x,-\SM@cf@top+\SM@cf@y)}}{}%
	\ifdefstring{\SM@cf@corner}{se}{%
		\def\SM@cf@anchor{south east}%
		\edef\SM@cf@at{(current page.south east)+(-\SM@cf@right+\SM@cf@x,\SMContentBottomAlign+\SM@cf@y)}}{}%
	\ifdefstring{\SM@cf@corner}{sw}{%
		\def\SM@cf@anchor{south west}%
		\edef\SM@cf@at{(current page.south west)+(\SMFrametitleLeftPadding+\SM@cf@x,\SMContentBottomAlign+\SM@cf@y)}}{}%
	% e/w: on the Right/Left Margin, centred vertically between Top and Bottom Margin
	\ifdefstring{\SM@cf@corner}{e}{%
		\def\SM@cf@anchor{east}%
		\edef\SM@cf@at{(current page.north east)+(-\SM@cf@right+\SM@cf@x,-0.5\paperheight-0.5\SMContentTopAlign+0.5\SMContentBottomAlign+\SM@cf@y)}}{}%
	\ifdefstring{\SM@cf@corner}{w}{%
		\def\SM@cf@anchor{west}%
		\edef\SM@cf@at{(current page.north west)+(\SMFrametitleLeftPadding+\SM@cf@x,-0.5\paperheight-0.5\SMContentTopAlign+0.5\SMContentBottomAlign+\SM@cf@y)}}{}%
	% TikZ's calc parser must see the coordinate itself, not a macro: build
	% the \node call with everything but the content already expanded
	\edef\SM@cf@cmd{\noexpand\node[anchor=\SM@cf@anchor, inner sep=0pt, outer sep=0pt%
		\ifSM@cf@border,SMImageBorder\fi
		\ifx\SM@cf@node\@empty\else,\unexpanded\expandafter{\SM@cf@node}\fi]
		at ($\SM@cf@at$)}%
	\vbox to 0pt{\vss
		\begin{tikzpicture}[overlay, remember picture]
			\SM@cf@cmd {\SM@cf@content};
		\end{tikzpicture}%
	}%
	\endgroup
}

% \begin{besidefigure}[keys]{image file or content} text \end{besidefigure}
% (2026-09-30; 2026-10-01 the figure's width is measured, not passed) -- text
% next to a figure in the top-right corner: the figure goes into \cornerfigure
% (same keys: height=/width=, border, ...), the text into a top-aligned
% minipage exactly as wide as the room left of the figure minus `gap`
% (default 12pt). 84 slides built this by hand, each with its own guessed
% minipage width (9.3cm, 0.58\linewidth, ...).
%   \begin{besidefigure}[height=5.6cm, border]{media/chart-x.png}
%   	\begin{itemize} ... \end{itemize}
%   \end{besidefigure}
% Optical top alignment (v5.27.0): the text column's first line sits on the Top
% Margin by its line box (the ascender height, 0.935em), so its visible letters
% start about 0.2em lower; a figure placed on the margin itself looked that
% much too high next to it. besidefigure therefore lowers its figure by this
% much (\SMBesideFigureTopGap, default 0.2em of the surrounding text), so the
% figure's top edge meets the top of the first line's letters. Give
% yshift=0pt among the keys to put the figure on the margin again.
\newlength{\SM@bf@width}
\NewDocumentEnvironment{besidefigure}{O{} +m}{%
	\SM@cf@setkeys{#1}%
	\SM@cf@makecontent{#2}%
	\settowidth{\SM@bf@width}{\SM@cf@content}%
	\ifSM@cf@border\addtolength{\SM@bf@width}{\SMImageBorderWidth}\fi
	\cornerfigure[yshift=-\SMBesideFigureTopGap,#1]{#2}%
	% text column: from the text's left edge to the figure's left edge minus gap
	\setlength{\SM@bf@width}{\dimexpr\linewidth+\beamer@rightmargin-\SMFrametitleRightPadding-\SM@bf@width-\SM@bf@gap\relax}%
	\begin{minipage}[t]{\SM@bf@width}%
}{%
	\end{minipage}%
}

% \fitfigure[keys]{image file} (2026-09-30) -- an image as large as it fits,
% keeping its aspect ratio, in the box from a top line down to the Bottom
% Margin (\SMContentBottomAlign) and from the Left to the Right Margin,
% placed as a zero-space overlay. Replaces the per-slide
% "scale up ... as much as possible" arithmetic (hand-computed widths and
% \vspace offsets, often converged in two passes).
%   top=here|margin   where the box starts: 'here' (default) = the current
%                     position in the slide (below the text so far), plus
%                     `gap`; 'margin' = the Top Margin
%   gap=<length>      space below the text for top=here (default 6pt)
%   valign=top|center|bottom, halign=left|center|right  (default center, center)
%   bottom=<length>   the box's bottom, from the page's bottom edge (default
%                     \SMContentBottomAlign)
%   node={...}        extra TikZ node options (e.g. SMImageBorder)
% The position is a remember-picture coordinate, so like every overlay here
% it needs the second compile.
\newcount\SM@fit@count
\def\SM@fit@settop#1{\coordinate (SMFit@top) at #1;}
\pgfqkeys{/SMFit}{
	top/.store in=\SM@fit@top,
	gap/.store in=\SM@fit@gap,
	valign/.store in=\SM@fit@valign,
	halign/.store in=\SM@fit@halign,
	bottom/.store in=\SM@fit@bottom,
	node/.store in=\SM@fit@node,
}
\NewDocumentCommand{\fitfigure}{O{} m}{%
	\begingroup
	\pgfqkeys{/SMFit}{top=here, gap=6pt, valign=center, halign=center,
		bottom=\SMContentBottomAlign, node={}, #1}%
	\global\advance\SM@fit@count by 1
	\edef\SM@fit@name{SMFit\the\SM@fit@count}%
	\ifdefstring{\SM@fit@top}{here}{%
		\vbox to 0pt{\hbox{\tikz[remember picture, overlay]\coordinate (\SM@fit@name) at (0,0);}\vss}%
	}{}%
	% alignment choices, decided here (pgfmath can't evaluate \ifstrequal)
	\def\SM@fit@xexpr{(\x2+\x3)/2}\def\SM@fit@h@anchor{}%
	\ifdefstring{\SM@fit@halign}{left}{\def\SM@fit@xexpr{\x2}\def\SM@fit@h@anchor{west}}{}%
	\ifdefstring{\SM@fit@halign}{right}{\def\SM@fit@xexpr{\x3}\def\SM@fit@h@anchor{east}}{}%
	\def\SM@fit@yexpr{(\y1+\y2)/2}\def\SM@fit@v@anchor{}%
	\ifdefstring{\SM@fit@valign}{top}{\def\SM@fit@yexpr{\y1}\def\SM@fit@v@anchor{north}}{}%
	\ifdefstring{\SM@fit@valign}{bottom}{\def\SM@fit@yexpr{\y2}\def\SM@fit@v@anchor{south}}{}%
	\edef\SM@fit@anchor{\SM@fit@v@anchor
		\ifx\SM@fit@v@anchor\@empty\else\ifx\SM@fit@h@anchor\@empty\else\space\fi\fi
		\SM@fit@h@anchor}%
	\ifx\SM@fit@anchor\@empty\def\SM@fit@anchor{center}\fi
	\ifdefstring{\SM@fit@top}{here}{%
		\edef\SM@fit@topcoord{(\SM@fit@name)}%
		\edef\SM@fit@gapexpr{\SM@fit@gap}%
	}{%
		\def\SM@fit@topcoord{($(current page.north)+(0,-\SMContentTopAlign)$)}%
		\def\SM@fit@gapexpr{0pt}%
	}%
	\edef\SM@fit@nodeopts{anchor=\SM@fit@anchor, inner sep=0pt%
		\ifx\SM@fit@node\@empty\else,\unexpanded\expandafter{\SM@fit@node}\fi}%
	\vbox to 0pt{\vss
		\begin{tikzpicture}[overlay, remember picture]
			\expandafter\SM@fit@settop\expandafter{\SM@fit@topcoord}%
			\coordinate (SMFit@bl) at ($(current page.south west)+(\SMFrametitleLeftPadding,\SM@fit@bottom)$);
			\coordinate (SMFit@r) at ($(current page.east)+(-\SMFrametitleRightPadding,0)$);
			\path let \p1 = (SMFit@top), \p2 = (SMFit@bl), \p3 = (SMFit@r) in
				\pgfextra{%
					\pgfmathsetlengthmacro\SM@fit@h{max(\y1-\SM@fit@gapexpr-\y2,1pt)}%
					\pgfmathsetlengthmacro\SM@fit@w{max(\x3-\x2,1pt)}%
					\pgfmathsetlengthmacro\SM@fit@x{\SM@fit@xexpr}%
					\pgfmathsetlengthmacro\SM@fit@y{\SM@fit@yexpr}%
					\ifdefstring{\SM@fit@valign}{top}{\pgfmathsetlengthmacro\SM@fit@y{\y1-\SM@fit@gapexpr}}{}%
					\ifdefstring{\SM@fit@valign}{center}{\pgfmathsetlengthmacro\SM@fit@y{(\y1-\SM@fit@gapexpr+\y2)/2}}{}%
					\global\let\SM@fit@H\SM@fit@h \global\let\SM@fit@W\SM@fit@w
					\global\let\SM@fit@X\SM@fit@x \global\let\SM@fit@Y\SM@fit@y
				};
			\expandafter\node\expandafter[\SM@fit@nodeopts] at (\SM@fit@X,\SM@fit@Y)
				{\includegraphics[width=\SM@fit@W, height=\SM@fit@H, keepaspectratio]{#2}};
		\end{tikzpicture}%
	}%
	\endgroup
}
\makeatother

%%%%%%%%%%%%%%%%%%%%%%%%%%%%%%%%%%%%%%%%%%%%%%%%%%%%%%%%%%%%%%%%%%%%%%%%%%%%%
% Special frame types (converted from ExampleSlides.pptx; see README.md)
%
% All five below deliberately reuse as much of standardframe -- and of each
% other -- as possible, rather than reimplementing frame chrome from
% scratch: sectionframe, figureframe, and highlightframe are literally thin
% wrappers around standardframe itself (same title bar, logo, footer, bullet
% styling; only what's documented at each one differs). titleframe and
% quoteframe can't be standardframe wrappers -- both are title-less, and
% quoteframe replaces the whole body layout -- but still reuse its palette,
% fonts, and margins (\SMFrametitleLeftPadding, the institute logo
% drawing command, etc.) rather than inventing new ones.
%%%%%%%%%%%%%%%%%%%%%%%%%%%%%%%%%%%%%%%%%%%%%%%%%%%%%%%%%%%%%%%%%%%%%%%%%%%%%

% titleframe: the deck's own title slide (ExampleSlides.pptx slide 1) --
% typically used exactly once, at the very start of a course's first deck.
% One continuous yellow box -- like standardframe's own title bar -- flush
% with the page's top and left edges, stopping short of the institute logo
% on the right (the same margin already used around the logo elsewhere:
% \SMLogoWidth + 2*\SMFrametitleRightPadding), holding the title and,
% below it, a second line (author/term/whatever) in the same black text as
% standardframe's own frametitle. Usage:
%
%   \titleframe{Financial Markets}{Stefan Palan\\Summer term 2026}
%
% A *command*, not an environment, taking the second (lower) line's content
% as its second argument -- plain text, use \\ for line breaks same as any
% normal LaTeX paragraph, not an itemize, since the source deck's own
% subtitle block is just stacked lines. This has to be a command rather
% than an environment with the body as its content: beamer's frame body is
% captured by a delimited-argument macro (\beamer@collect@@body, looked up
% by triggering the error below on purpose) hunting for a literal
% \end{frame}, and a beamercolorbox -- needed by an earlier version of this
% to wrap the subtitle text in its own colored band -- opened in a *custom
% environment's* begin code and closed in its end code confuses that hunt
% exactly like the \columns example above (same root cause, different
% environment): with
% \newenvironment{titleframe}{...\begin{beamercolorbox}{...}}{\end{beamercolorbox}...},
% compiling ends in "! File ended while scanning use of \beamer@collect@@body."
% every time, however short the body. Taking the content as a normal brace
% argument instead sidesteps the whole issue. No footer/frame number
% (matches beamer's own convention for title slides, e.g. the old
% \begin{frame}[plain,noframenumbering]\titlepage\end{frame} this replaces --
% see README.md) and no chapter-logo support: both would be meaningless on a
% slide that appears once, before any chapter exists.
%
% Unlike the frametitle hook further up, neither the logo nor the yellow box
% here need the white-patch-and-redraw trick: the box's own width already
% stops short of the logo (rather than being painted full-width and then
% partly painted back over), so there's no yellow left to paint over in the
% first place. Both still need \tikz's "remember picture" on their
% "current page"-anchored nodes, though, and so still need the usual second
% compile pass (see "Title bar and logo" in README.md) -- checked
% empirically: without it, "current page" doesn't resolve to the actual page
% box on the first pass, and the logo/yellow box end up positioned nowhere
% near the page corners they're meant to be anchored to.
%
% \SMTitleframeBoxHeight is a fixed target for the box's *bottom* edge
% (matching what was actually asked for: "reaching down to where the lower
% grey bar currently ends" in the two-band design this replaces), not the
% blank space above the text -- that space is derived from it instead (box
% height minus the title+subtitle content's own measured height), so the
% bottom edge stays fixed even if a longer title ever wraps onto a second
% line; the blank space above would just shrink to compensate. Measured
% against the old two-band design's own rendered PDF coordinates (before
% this change): its second band's bottom edge sat 187.147pt below the page's
% top edge.
\newlength{\SMTitleframeBoxWidth}
\setlength{\SMTitleframeBoxWidth}{\dimexpr\paperwidth-\SMLogoWidth-2\SMFrametitleRightPadding\relax}
% Text inset from the box's own edges: \SMFrametitleLeftPadding, the same
% length standardframe's own title text uses, plus \SMTitleframeTextSep
% (2.5ex, matching this command's own former per-band beamercolorbox "sep").
% Because the leading blank space above the text (further down) is itself
% computed as "box height minus the content's own natural height", the
% content's *bottom* padding is the only thing that determines how far the
% last line sits above the box's bottom edge -- the leading space silently
% absorbs whatever that padding is, keeping the bottom edge fixed either
% way. (The relationship is a direct 1:1 offset, checked empirically by
% comparing two rendered values, but with a nonzero intercept from beamer's
% own default top-of-frame vskip -- there's no clean closed-form for the
% right value, hence the empirical calibration here rather than a formula.)
% Calibrated (against the rendered PDF's own text coordinates, not just
% visually) to about 1.5 line-heights of the subtitle text's own \Large
% size (~20.6pt/line here) above the box's bottom edge; re-measure and
% adjust if the subtitle's font size, or beamer's own default top-of-frame
% vskip, ever change.
\newbox\SMTitleframeBox
\makeatletter
\newcommand{\titleframe}[2]{%
	% No noframenumbering (2026-08-11, direct instruction): this frame has
	% no footer at all ("plain"), so removing it has no visible effect on
	% the title slide itself -- but it does mean the title slide now
	% consumes a real "framenumber" step, same as any other frame, so the
	% deck's first real content slide is numbered 2, not 1 (matching this
	% slide's own physical position in the compiled PDF -- SlideMage's
	% own deck-editor "Render" preview relies on this: it derives a
	% previewed slide's page number from its plain list position, which
	% only equals the real compiled frame number when every \slide{}-family
	% call -- including this one -- consumes exactly one framenumber step).
	\SMOpenFrame{[plain,t]}%
	% This frame has no footer to carry the usual "(slide id: ...)" marker
	% -- see \SMInvisibleSlideIDStamp's own comment above for why that
	% breaks SlideMage's compile-time preview snapshot, and what this
	% stamps instead.
	\SMInvisibleSlideIDStamp%
	% A bare tikzpicture behaves like an hbox (same reasoning as the footer
	% box's own \vbox-to-0pt wrapper further up): sitting directly in the
	% frame's vertical list, it silently opens an implicit paragraph,
	% switching into horizontal mode. Everything placed after it -- the
	% \vspace* further down included -- would then run in horizontal mode
	% instead of vertical mode, which is what a first version of this command
	% got wrong: with two bare tikzpictures here, \vspace*'s glue ended up as
	% *horizontal* space on an implicit line together with them (visible as
	% the title/subtitle shifting sideways, not down, and the vspace amount
	% appearing to have no effect at all -- it was moving the wrong axis).
	% \vbox to 0pt{\vss ...} keeps each overlay at zero space in the normal
	% flow without ever leaving vertical mode.
	%
	% \SMArtifactBegin/\SMArtifactEnd (defined with the footer, "PDF
	% accessibility" comment): this box is purely decorative background, same
	% kind of element as the footer's own yellow box -- marked the same way,
	% for a screen reader to skip. titleframe has no footer at all to worry
	% about separately (it's a "plain" frame, no footline), so this box is
	% the only thing here that needs it; the title/subtitle text right below
	% is real content and stays unmarked.
	\SMArtifactBegin{/Type /Background}%
	\vbox to 0pt{\vss
	\begin{tikzpicture}[overlay, remember picture]
		% \SMTitleBoxColor, not a literal color token -- same setting the
		% standard frametitle bar (colorA, above) uses, so the title slide's
		% own box matches every other title bar in the deck (2026-07-16,
		% direct instruction).
		% outer sep=0pt here and on the logo below: without it, tikz's
		% default outer sep insets the fill/logo 0.2pt from the anchored
		% page corner -- a hairline white line above/left of this box, and
		% its right edge/the logo 0.2-0.4pt off every other frame's. See
		% the frametitle hook's own comment for the full story.
		\node[fill=\SMTitleBoxColor, anchor=north west, inner sep=0pt, outer sep=0pt,
		      minimum width=\SMTitleframeBoxWidth,
		      minimum height=\SMTitleframeBoxHeight]
			at (current page.north west) {};
	\end{tikzpicture}%
	}%
	\SMArtifactEnd
	\vbox to 0pt{\vss
	\begin{tikzpicture}[overlay, remember picture]
		% inner ysep derived, not \SMFrametitleRightPadding -- see the
		% frametitle hook's fixed-position logo redraw for why (matches the
		% inline logo's vertically-centered position on every other frame).
		\node[anchor=north east, inner xsep=\SMFrametitleRightPadding,
		      inner ysep=\dimexpr\SMTitleBarBottomAlign-\SMLogoHeight\relax,
		      outer sep=0pt]
			at (current page.north east) {\SMInstituteLogo};
	\end{tikzpicture}%
	}%
	\setbox\SMTitleframeBox=\vbox{%
		\hsize=\SMTitleframeBoxWidth
		\parindent=0pt
		\leftskip=\SMFrametitleLeftPadding
		\advance\leftskip by \SMTitleframeTextSep
		\rightskip=\SMFrametitleLeftPadding
		\advance\rightskip by \SMTitleframeTextSep
		% \SMTitleTextColor, not a literal color token -- follows the
		% box color change right above (2026-07-16, direct instruction):
		% leaving this hardcoded dark would be unreadable once the box
		% itself switches to a dark blue for MCI/Innsbruck.
		\color{\SMTitleTextColor}%
		% Beamer redefines \\ inside every frame to peek for an overlay spec
		% (\\<2->, for line-by-line reveals in itemize/enumerate); that
		% redefinition assumes it's running in beamer's own frame-body
		% machinery and chokes ("Undefined control sequence", silently
		% swallowing the following word) when \\ is used for a plain line
		% break in a raw \vbox like this one. Restoring plain LaTeX's own
		% \\ (\@normalcr) for this group only fixes it without affecting \\
		% anywhere else in the deck. (The @ here needs catcode 11, which is
		% why this whole command definition -- not just this line -- is
		% wrapped in \makeatletter/\makeatother: macro bodies are tokenized
		% at *definition* time, so a \makeatletter placed inside the body
		% itself would take effect too late, after "\@normalcr" had already
		% been mis-tokenized as "\@" followed by four ordinary-text letters.)
		\let\\\@normalcr
		\vskip\SMTitleframeTextSep
		\noindent{\Huge\bfseries #1\par}%
		\vskip\SMTitleframeTextSep
		\vskip\SMTitleframeBandGap
		\vskip\SMTitleframeTextSep
		\noindent{\Large #2\par}%
		\vskip\SMTitleframeBottomPad
	}%
	\vspace*{\dimexpr\SMTitleframeBoxHeight-\ht\SMTitleframeBox-\dp\SMTitleframeBox\relax}%
	\box\SMTitleframeBox
	% This command's own \color{\SMTitleTextColor} above (set inside a
	% stored box register) can leak into the *next* frame's body text --
	% see standardframe's own opening code ("Forces this frame's own body")
	% for the real fix and the full story (a trailing reset attempted
	% here first turned out to be a silent no-op, since nothing draws after
	% it within this same frame to actually force the color operator to be
	% emitted).
	\end{frame}%
}
\makeatother

% fullbleedframe: a slide that is *only* a full-screen image (FM2_2026s_ps_deck.pptx
% slide 13, a recruitment flyer pasted in as one embedded picture on an
% otherwise blank layout -- no title, no PPTX text at all) -- no title bar,
% no footer, no logo, matching that source's own blank layout intent, unlike
% figureframe (which is standardframe underneath and always keeps the usual
% title/footer chrome). The image is scaled to fit the slide while keeping
% its own aspect ratio (\includegraphics's keepaspectratio, capped by
% \paperwidth *and* \paperheight together), centered; whatever's left over
% on the binding axis (e.g. a widescreen slide holding a taller/narrower
% image) is filled by a solid black backdrop rather than left showing the
% frame's usual white canvas -- the point of a full-*bleed* image is that
% nothing reads as "empty margin" even when the aspect ratios don't match.
%
%   \fullbleedframe{media/photo-datacamp-recruitment-flyer}
%
% \SMFullBleedBlack is a fixed HTML black, deliberately *not* a
% \colorlet of c_primary_dark: this frame type's backdrop should stay
% neutral true black regardless of whether c_primary_dark itself is ever
% restyled to something else (e.g. a softer near-black, the way
% c_custom_gray already is elsewhere in this file) for branding
% reasons -- those are two different jobs (brand color vs. letterbox
% backdrop) that only happen to want the same value today.
\definecolor{SMFullBleedBlack}{HTML}{000000}
\newcommand{\fullbleedframe}[1]{%
	% No noframenumbering -- see \titleframe's own comment above for why
	% (same reasoning, same 2026-08-11 change: this frame has no footer
	% either, so nothing here changes visibly, only later slides' numbers).
	\SMOpenFrame{[plain,t]}%
	% This frame has no footer either -- same \SMInvisibleSlideIDStamp
	% need as \titleframe above, see its own comment for why.
	\SMInvisibleSlideIDStamp%
	% \SMArtifactBegin/\SMArtifactEnd (defined with the footer, "PDF
	% accessibility" comment): the black backdrop is purely decorative
	% letterboxing, marked the same way as titleframe's/quoteframe's own
	% background boxes so a screen reader skips it. The image itself is left
	% unmarked *here* -- not because it goes undescribed, but because
	% \includegraphics's own global wrapper (this file's own
	% \RenewDocumentCommand{\includegraphics}, "PDF accessibility" section)
	% already does that job for every \includegraphics call site, this one
	% included: it records the image's role/source with
	% slidemage_alt.mark()/recordSource() and, for an informative image, emits
	% its short alt text as an invisible layer via \SMAltInvisible
	% before the image itself. Nothing extra is needed at this call site.
	\SMArtifactBegin{/Type /Background}%
	\vbox to 0pt{\vss
	\begin{tikzpicture}[overlay, remember picture]
		% outer sep=0pt: see the frametitle hook's comment -- without it the
		% backdrop sits 0.2pt in from the anchored page corner.
		\node[fill=SMFullBleedBlack, anchor=south west, inner sep=0pt, outer sep=0pt,
		      minimum width=\paperwidth, minimum height=\paperheight]
			at (current page.south west) {};
	\end{tikzpicture}%
	}%
	\SMArtifactEnd
	\SMFullBleedPick{#1}{\SMFullBleedLayers{#1}}{\SMFullBleedImage{#1}}%
	\end{frame}%
}
\makeatletter
\newcommand{\SMFullBleedPick}[3]{\in@{,}{#1}\ifin@#2\else#3\fi}
\makeatother

% \fullbleedframe{a.png,b.png,c.png} (2026-09-18, direct instruction:
% combine the Fink/Palan/Theissen excursus's build-step screenshots into
% one slide): a comma-separated list layers the images on top of each
% other, revealed one per overlay step (image k appears at step k and stays,
% each covering the previous one -- the screenshots are same-size, opaque
% and each already contains every earlier step's content, so stacking
% reproduces the source's own build animation). Instructor deck only --
% the during-course/after-course versions show just the last (final)
% image, in a single step. A single path behaves
% exactly as before. Once a layer is covered by the next one it is wrapped
% in the same /Background artifact as the backdrop above, so on step k a
% screen reader reads only image k's own alt text (each screenshot is
% described as its own "build step k of n") instead of every buried layer's.
\newcommand{\SMFullBleedImage}[1]{%
	\vbox to 0pt{\vss
	\begin{tikzpicture}[overlay, remember picture]
		\node[anchor=center] at (current page.center)
			{\includegraphics[width=\paperwidth, height=\paperheight, keepaspectratio]{#1}};
	\end{tikzpicture}%
	}%
}
\newcount\SMFBtotal
\newcommand{\SMFullBleedLayers}[1]{%
	\global\SMFBtotal=0
	\foreach \SMFBimg in {#1} {\global\advance\SMFBtotal by 1 }%
	\foreach \SMFBimg [count=\SMFBi] in {#1} {%
		\ifSMInstructorVersion
			\ifnum\SMFBi<\SMFBtotal
				\pgfmathtruncatemacro{\SMFBnext}{\SMFBi+1}%
				\only<\SMFBi>{\SMFullBleedImage{\SMFBimg}}%
				\only<\SMFBnext->{%
					\SMArtifactBegin{/Type /Background}%
					\SMFullBleedImage{\SMFBimg}%
					\SMArtifactEnd
				}%
			\else
				\only<\SMFBi>{\SMFullBleedImage{\SMFBimg}}%
			\fi
		\else
			\ifnum\SMFBi=\SMFBtotal
				\SMFullBleedImage{\SMFBimg}%
			\fi
		\fi
	}%
}

% sideimageframe: a standardframe that always carries one auto-placed
% "side filler image" (the user's own term, see
% slides/dfc-Angel-Investors-current.tex's own comment for the technique
% this formalizes into a reusable frame type) -- 2026-07-06, direct
% instruction. The image is positioned flush to three edges at once: its
% top touches the title bar's own bottom edge, its right touches the
% page's own right edge, and its bottom touches the footer box's own top
% edge -- so it's scaled by *height* to exactly
% \paperheight-\SMTitleBarBottomAlign-\SMFooterBoxHeight (the
% measured gap between those two), with `keepaspectratio` so the width
% simply follows from the image's own aspect ratio rather than being
% forced/distorted to some second, independent target -- anchoring the
% node at the page's own north-east corner then makes the top+right
% touches automatic, and the height exactly filling that gap makes the
% bottom touch automatic too, with no separate width constraint needed
% (only a *square* source image would ever visibly reach some further
% left boundary; this only guarantees top/right/bottom, matching what was
% actually asked for).
%
%   \begin{sideimageframe}{Title}{media/photo-example}
%   	\begin{itemize}
%   		\item ...
%   	\end{itemize}
%   	\source{...}
%   \end{sideimageframe}
%
% \sidefootnotes is switched on automatically (not left to the caller) --
% the whole point of this frame type is that its footnotes always end up
% rotated along the left edge of the image rather than colliding with it
% the way a full-width bottom footnote would (see
% slides/u29-Implications-of-Capital-Market-Efficiency-Prices-current.tex's
% own comment for exactly that collision, on a plain standardframe). Unlike
% every other \sidefootnotes use elsewhere in this deck, though, the strip
% can't sit at its usual fixed inset from the page's own right edge here --
% that would run it directly underneath the image, the same problem u29
% had. Fixed by *computing* the image's own rendered width (via
% \settowidth against the exact same \includegraphics call used to draw
% it) and using that to override \SMSideFootnoteRightInset for just
% this one frame -- moving the whole rotated strip to sit immediately left
% of the image instead of at the page's raw edge, entirely automatically,
% with no manual x-coordinate needed from the caller (the direct
% instruction offered a manual-x-coordinate fallback "if it is possible
% [to do automatically], that would be great" -- turned out to be
% possible). \SMSideImageFrameHeight/\SMSideImageFrameWidth are
% declared once, globally, and merely *set* (not re-declared) inside the
% environment body, so the frame type can be used more than once per deck;
% the \setlength override of \SMSideFootnoteRightInset is a plain,
% non-\global assignment, which -- like every other frame-local override in
% this file (e.g. slides/u29-...'s own \setbeamercolor{footnote}{fg=white})
% -- resets on its own once \begin{frame}...\end{frame}'s own TeX group
% closes, with no explicit reset needed.
%
% v5.30.5: 3.1pt became 2.1pt. Measured on a rendered slide whose rotated footnote
% text was finally visible (see the colour note at the strip's drawing code), 3.1pt
% left 3.9pt between the ink of the line nearest the image and the image -- the
% text's own depth strut adds ~0.8pt -- against the 2.93pt this was meant to match.
%
% The extra gap added on top of the image's own measured width (2026-07-06,
% direct instruction: "replicate ... the precise margin we have on regular
% slides between footnote and footer box") was originally an unmeasured
% 4pt guess -- replaced with 3.1pt, back-solved from the deck's own real
% reference gap: measured directly on a plain bottom-footnote slide
% (slides/xhe-Active-vs-Passive-Management-2026-First-current.tex's own
% \source{} line) that the real gap between a footnote's own text and the
% footer box's top edge is 2.93pt, while the original 4pt guess here was
% actually landing the rotated text 3.8pt from the image -- 4pt - (3.8-2.93)
% = 3.13pt, rounded to 3.1pt.
\newlength{\SMSideImageFrameHeight}
\newlength{\SMSideImageFrameWidth}
\newenvironment{sideimageframe}[3][t]{%
	\setlength{\SMSideImageFrameHeight}{\dimexpr\paperheight-\SMTitleBarBottomAlign-\SMFooterBoxHeight\relax}%
	\settowidth{\SMSideImageFrameWidth}{\includegraphics[height=\SMSideImageFrameHeight, keepaspectratio]{#3}}%
	\sidefootnotes
	\setlength{\SMSideFootnoteRightInset}{\dimexpr\SMSideImageFrameWidth+2.1pt\relax}%
	\begin{standardframe}[#1]{#2}%
	\vbox to 0pt{\vss
	\begin{tikzpicture}[overlay, remember picture]
		% outer sep=0pt: see the frametitle hook's comment -- without it the
		% image sits 0.2pt in from the page's right edge (and 0.2pt low).
		\node[anchor=north east, inner sep=0pt, outer sep=0pt] at ($(current page.north east)+(0,-\SMTitleBarBottomAlign)$)
			{\includegraphics[height=\SMSideImageFrameHeight, keepaspectratio]{#3}};
	\end{tikzpicture}%
	}%
}{%
	\end{standardframe}%
}

% sectionframe: a section-opening overview (ExampleSlides.pptx slide 3) --
% same layout as standardframe (footer, logo, bullet styling all completely
% unchanged), except the title bar itself is much taller (about a third of
% the page) and set in a bigger font, and the body text is a step bigger
% too, so it reads as "somewhat greater focus" and is unmistakable at a
% glance when flipping through slides -- while still keeping the deck's
% usual yellow-background/black-text look, unlike an earlier version of this
% that inverted to dark-background/yellow-text instead.
%
% The taller bar is NOT a separately-drawn box the way titleframe's is (see
% "Title slide" above): the gotham theme's own frametitle box already sizes
% itself from \SMFrametitleTopPadding/\SMFrametitleBottomPadding plus
% the title font's own strut height (beamerthemeSlideMage.sty, the
% "frametitle" template), so overriding those two lengths locally -- the
% same trick already used below for the font/color overrides -- grows the
% *theme's own* box instead of fighting it with a second one. This also
% means the white-patch-behind-the-logo/chapter-logo machinery (tied to the
% frametitle template via \addtobeamertemplate) keeps working unmodified,
% since it measures the box's actual rendered bottom edge rather than
% assuming any particular height.
%
% The title text is bottom-anchored within the bar -- sitting a fixed "one
% line" above the bar's bottom edge, growing *upward* for a title long
% enough to wrap onto a second line -- rather than top-anchored the way
% gotham's own frametitle template flows by default (which would instead
% leave a large, growing gap *below* a wrapped title, between it and the
% bar's bottom edge). gotham's template has no built-in way to do this: its
% only two knobs are the toppading/bottompading struts, both anchored to
% *one* end of a single top-down paragraph flow. So instead of leaving
% toppading fixed (as an earlier version of this did) the whole title text
% is measured into \SMSectionframeTitleBox first -- built with the same
% font and (approximately) the same usable width gotham's own template
% gives it, so its \ht+\dp closely track the title's real rendered size,
% including wrapping onto extra lines for a long title -- and toppading is
% then computed as "whatever's left" after that measured height and the
% fixed bottom gap are subtracted from the bar's total target height. A
% longer/wrapped title therefore automatically shrinks toppading (pushing
% the whole title block up, growing upward) rather than changing the gap
% below it. \SMSectionframeTitleWidth is a slight underestimate of
% gotham's actual per-line width (which is only narrowed by the logo on
% whichever line ends up next to it, typically just the last one) --
% applying that narrower width to *every* line is a deliberate, harmless
% simplification rather than replicating gotham's per-line logic exactly.
%
% The institute logo moves down along with the title now (it's vertically
% centered against the title's own baseline, gotham's instituteLogo
% template), rather than staying pinned where it sits on every other frame
% type, as an earlier version of this kept it -- an intentional change from
% that version: with the title itself now sitting near the bar's bottom
% edge, keeping the logo up at the bar's top would visually separate it
% from the title it's meant to sit beside.
%
% \begingroup/\endgroup (rather than the \global-set-and-reset dance
% \chapterlogo{}/\sidefootnotes need) is enough here: this only ever touches
% \setbeamercolor/\setbeamerfont/\setlength/\setbox, all of which assign
% locally under the hood (checked against beamerbasecolor.sty/
% beamerbasefont.sty directly for the first two -- their key handlers are
% plain \def, not \gdef -- and \setlength/\setbox are plain TeX assignments,
% local by default), so a plain surrounding TeX group reverts them
% correctly on its own.
%
% \SMSectionframeContentDrop adds a little extra room below the bar
% before the body content starts, on top of standardframe's own
% \SMContentTopShift pull-up (applied inside \begin{standardframe}
% already, before this code regains control). All lengths here calibrated
% against the rendered PDF's own coordinates; re-measure and adjust if the
% title/body font sizes change.
% \SMSectionframeTrueBoxBottom (2026-07-25, direct instruction: "the
% image extends into the title box" on neb/u97 -- followed by "tyn actually
% looks fine, find out what is really driving this issue/difference"):
% \SMSectionframeBoxHeight above is only the *target* the toppading
% formula below solves for -- by construction that formula makes
% toppading+title-height+title-depth+bottompading algebraically equal
% boxHeight+\SMSectionframeHeightCorrection regardless of title content,
% so the *rendered* bar height should be a fixed, title-independent
% 97.14pt... except it measurably isn't (confirmed directly, PyMuPDF vector
% fill): neb/u97 ("Market Environment"/"Household Finance") render at
% 89.19pt, while tyn/vtm ("Capital asset pricing model"/"Measuring and
% managing risk") render at 85.04pt -- a real, reproducible ~4.15pt
% difference driven entirely by the title text itself, not by anything
% neb/u97 did wrong.
%
% Root cause (traced into gotham.dtx's own
% \SM@frametitlestrut@start/@end, not just inferred): the toppading
% formula's own `-\dp\SMSectionframeTitleBox` term assumes the title's
% real descender depth gets reintroduced somewhere else in the box so it
% nets out -- it doesn't. `@start` is a zero-depth `\rule` whose height is
% `toppading+\totalheightof{X}` (a *fixed reference glyph*, not the real
% title), so the line's rendered height is toppading-driven, independent of
% the title's own real height. `@end` is a zero-height `\rule` of depth
% `\SMFrametitleBottomPadding` (19.90pt for sectionframe) -- also fixed,
% not the real title's own depth. The real title text sits between these
% two struts and contributes its own real height/depth to the same TeX
% line, but a line's rendered depth is the *max* of every contributor, not
% a sum -- and 19.90pt comfortably exceeds any real title's own descender
% depth (measured directly: 0.25pt with none, e.g. "Market
% Environment"/"Household Finance"; 4.42pt with any single descender
% present -- g/j/p/q/y, letter/count doesn't matter, e.g. "Capital asset
% pricing model"/"Measuring and managing risk"/a plain "Testing gjpqy"),
% so bottompading always wins there and the title's real depth is never
% actually given back. Net effect: every point of a title's own real
% descender depth shrinks toppading (and therefore the whole visible box)
% by that same amount, with nothing compensating -- confirmed exactly:
% 4.42pt-0.25pt=4.17pt predicted vs. 89.19pt-85.04pt=4.15pt measured, a
% match well within measurement noise. \SMSectionframeHeightCorrection
% (12.1pt) was presumably calibrated against a descender'd test title (it
% lands tyn/vtm almost exactly on the 85.04pt target), leaving every
% no-descender title's own bar ~4.15pt taller than that target instead.
% Full derivation, and why this means \SMSectionframeBoxHeight was
% never actually title-independent, in README.md's own "Section overview
% (sectionframe)" section.
%
% \SMSectionframeTrueBoxBottom below is the *measured* value for a
% no-descender title specifically (matching neb/u97, the only two slides
% using it) -- it is NOT a general replacement for
% \SMSectionframeBoxHeight: a sectionframe whose title has any
% descender (tyn/vtm, confirmed landing on the *other* value, 85.04pt) must
% keep anchoring against \SMSectionframeBoxHeight directly instead
% (already correct for them, so left untouched). Anything new needing the
% bar's own true rendered bottom should re-measure per its own actual title
% text (real vector fill, not a visual guess, and not an assumption that
% either existing constant transfers) rather than reuse this length
% unless its own title also has zero descenders.
% v5.28.0: the bar height no longer depends on the title (see
% \SMSectionframeRefDepth below), so the "true" bottom is the same 85.04 for
% every title; the old 89.19pt value only fitted titles without descenders.
\newlength{\SMSectionframeTrueBoxBottom}
\setlength{\SMSectionframeTrueBoxBottom}{\SMSectionframeBoxHeight}
% (v9.8.0: \SM@sectionbarmetrics below places the title now, from the bar's own material; the
% four profile numbers named in this and the paragraphs above and below, \SMSectionframeRefDepth,
% \SMSectionframeBottomGap, \SMSectionframeBottomPadCorrection and \SMSectionframeHeightCorrection,
% are gone, and so is the formula these paragraphs explain. They stay as the record of why it was
% done this way for ten versions.)
% \SMSectionframeRefDepth (v5.28.0): the depth the toppading formula assumes for
% the title. It used to subtract the title's own measured depth, which nothing
% adds back, so a title without descenders ("Slides", "Decks and deck
% versions") got a bar 4.15pt taller than one with them, and every other depth
% (p, y, ...) a height in between. With this fixed value (the depth of "p",
% 4.42pt) every title gets the 85.04bp bar the calibrated constants below were
% fitted to, and \SMSectionframeBarBottom (where \sectionfigure starts) is that
% same constant.
% \SMSectionframeBottomGap is the *visual* target -- one huge-font line
% height (baselineskip, measured directly from a deliberately 2-line test
% title: the gap between the two lines' own rendered top edges). It is NOT
% used directly as \SMFrametitleBottomPadding: gotham's bottom strut
% (\SM@frametitlestrut@end) consistently produces a visual gap about
% 5.58pt *smaller* than the padding value fed into it (checked empirically,
% same offset for both a one-line and a two-line test title, so it's a
% fixed per-box offset rather than something that scales with content) --
% \SMSectionframeBottomPadCorrection adds that back.
% Similarly, the box's total rendered height comes out consistently ~11.4pt
% *less* than toppading+measuredTitleHeight+bottompading would suggest --
% the measured title box (built as a plain \vbox{...\par}) doesn't reproduce
% gotham's own leading strut (\SM@frametitlestrut@start, a zero-width
% rule with an explicit height rather than natural interline glue) exactly,
% and this correction closes that gap. Checked the same way as the bottom
% one above: consistent across both the one-line and two-line test titles.
\newlength{\SMSectionframeTitleWidth}
\setlength{\SMSectionframeTitleWidth}{\dimexpr\paperwidth-\SMFrametitleLeftPadding-\SMLogoWidth-2\SMFrametitleRightPadding\relax}
\newbox\SMSectionframeTitleBox
\newlength{\SMSectionframeBarBottom}
% The title's place in the bar (v9.8.0: from the title's own material, no longer from four numbers
% measured at Fira): its last baseline lies the depth of a \strut of the title font plus
% \SMSectionframeTitleGap above the bar's bottom edge, and the top padding is what makes the bar
% \SMSectionframeBoxHeight high. The bar's content is typeset once into a box, by the same
% beamercolorbox the frame title template uses and with a top padding as large as the bar; the box's
% height is that padding plus a constant from there on (any title, any number of lines), so the
% padding that gives the bar its height follows. #1 is the title.
\makeatletter
\newbox\SM@sectionprobe
\newcommand{\SM@sectionbarmetrics}[1]{%
	\begingroup\usebeamerfont{frametitle}\setbox\z@\hbox{\strut}\setlength{\dimen0}{\dimexpr\dp\z@+\SMSectionframeTitleGap\relax}%
	\xdef\SM@sectionbottom{\the\dimen0}\endgroup
	\setlength{\SMFrametitleBottomPadding}{\SM@sectionbottom}%
	\setlength{\SMFrametitleTopPadding}{\SMSectionframeBoxHeight}%
	\setbox\SM@sectionprobe=\vbox{\usebeamerfont{frametitle}%
		\begin{beamercolorbox}[wd=\paperwidth, sep=0pt, leftskip=\SMFrametitleLeftPadding,
			rightskip=\dimexpr\SMLogoWidth+2\SMFrametitleRightPadding\relax plus 1fil]{frametitle}%
			\SM@frametitlestrut@start #1\SM@frametitlestrut@end
		\end{beamercolorbox}}%
	\setlength{\SMFrametitleTopPadding}{\dimexpr2\SMSectionframeBoxHeight-\ht\SM@sectionprobe-\dp\SM@sectionprobe\relax}}
\makeatother
\newenvironment{sectionframe}[2][t]{%
	\begingroup
	\SMManualTopMargin% own title-bar geometry: no automatic Top Margin
	\global\csname SM@topsectiontrue\endcsname% its taller bar: not counted for \SMContentTopShift
	% Disable hyphenation for the title -- both here, in the measurement
	% box below, and (since this is in effect for the rest of the group,
	% which \begin{standardframe} -- triggering the real \insertframetitle
	% -- runs inside) in the actual rendered title too, so a line break
	% never lands mid-word. Setting the penalties alone is the standard,
	% sufficient way to do this (TeX treats a penalty of 10000 at a
	% breakpoint as forbidden, not just heavily discouraged, so it's never
	% chosen even when the alternative is an overfull line); the hyphenmin
	% lengths are redundant insurance, not strictly needed on top of that.
	\hyphenpenalty=10000
	\exhyphenpenalty=10000
	\lefthyphenmin=100
	\righthyphenmin=100
	\setbeamerfont{frametitle}{size=\huge}
	\csname SM@sectionbarmetrics\endcsname{#2}%
	\global\def\SMFrametitleFullHeightYellow{1}%
	\global\def\SMFrametitleLogoFixed{1}%
	\setbeamertemplate{frametitle}[SMNoInlineLogo]%
	\setlength{\SMSectionframeBarBottom}{\dimexpr\SMSectionframeBoxHeight+0.12bp\relax}% 85.14bp for the default bar: the figure starts 0.12bp under it
	% \sectionfigure sets the side-footnote lengths \global (see there); the
	% values from before this frame come back at \end{sectionframe}.
	\edef\SMSectionframeRestoreSideFootnotes{%
		\global\SMSideFootnoteRightInset=\the\SMSideFootnoteRightInset\relax
		\global\SMSideFootnoteHsize=\the\SMSideFootnoteHsize\relax}%
	\begin{standardframe}[#1]{#2}%
	\vspace*{\SMSectionframeContentDrop}%
	\large
}{%
	\end{standardframe}%
	\SMSectionframeRestoreSideFootnotes
	\endgroup
}

% \sectionfigure[includegraphics keys]{image} -- a sectionframe's side image
% (2026-10-01, direct instruction: bij's image should "fit snugly between
% header box and footer and touch the right edge of the slide", for section
% slides with images in general, with the image's source visible). Put it
% anywhere in the sectionframe body. The image is scaled by height to fill
% the gap from the title bar's bottom edge to the footer box's top edge and
% sits flush with the page's right edge -- sideimageframe's geometry, but
% for sectionframe's taller bar. That bar's bottom edge depends on the
% title: the theme's strut gives the last line a fixed depth of ~4.4pt,
% but the padding subtracts the title's real depth -- so the bottom edge is
% 89.44bp minus that depth, but never above 85.04bp (a descender deeper
% than the strut deepens the bar again). Measured in the PDF, i.e. in bp:
% 89.19bp for "Market Environment" (depth 0.25pt), 86.01bp for "Market,
% Environment" (3.44pt), 85.04bp for "Capital asset pricing model"
% (4.42pt) and "Trends in publishing" (4.73pt). (The two
% constants \SMSectionframeTrueBoxBottom/\SMSectionframeBoxHeight
% hold the same numbers in TeX pt, 0.3pt short -- neb/u97 used to add
% that back by hand.) sectionframe computes it as
% \SMSectionframeBarBottom. (Not measured
% from the SMTitleBottom node: the title bar is typeset after the body,
% so that node still belongs to the previous frame while the body runs.)
% The frame's footnotes -- the automatic "Image source(s):" line included --
% go into the rotated side strip immediately left of the image, as on
% sideimageframe. The strip's inset is read at shipout, after the frame body,
% so it and the strip length are set \global here; \end{sectionframe} puts
% the previous values back. Replaces the per-slide \settowidth/
% \SMSideFootnoteRightInset/\SMSideFootnoteHsize set-up of neb, u97,
% tyn and vtm.
\newlength{\SMSectionFigureHeight}
\newlength{\SMSectionFigureWidth}
\newcommand{\sectionfigure}[2][]{%
	\setlength{\SMSectionFigureHeight}{\dimexpr\paperheight-\SMSectionframeBarBottom-\SMFooterBoxHeight\relax}%
	\settowidth{\SMSectionFigureWidth}{\includegraphics[height=\SMSectionFigureHeight, keepaspectratio, #1]{#2}}%
	\vbox to 0pt{\vss
	\begin{tikzpicture}[overlay, remember picture]
		\node[anchor=north east, inner sep=0pt, outer sep=0pt] at ($(current page.north east)+(0,-\SMSectionframeBarBottom)$)
			{\includegraphics[height=\SMSectionFigureHeight, keepaspectratio, #1]{#2}};
	\end{tikzpicture}%
	}%
	\global\SMSideFootnotestrue
	% 3.1pt: sideimageframe's measured gap between strip and image.
	% 2.31pt: the clearance above the footer box that \SMSideFootnoteHsize's
	% 207pt default leaves under a standard title bar.
	\global\SMSideFootnoteRightInset=\dimexpr\SMSectionFigureWidth+3.1pt\relax
	\global\SMSideFootnoteHsize=\dimexpr\SMSectionFigureHeight-\SMSideFootnoteTopSep-2.31pt\relax
}

% quoteframe: a quote slide (ExampleSlides.pptx slide 4) -- institute logo
% as usual, and the quote itself set large/italic/black on a yellow
% full-bleed block, same style as titleframe/sectionframe's own title boxes
% (flush top-left, stopping short of the logo on the right, same margin as
% elsewhere) rather than the dark block an earlier version of this used. The
% usual footer still shows (unlike titleframe, which drops it -- a title
% slide is a one-off outside any particular course content, while a quote
% is still one slide *among* a deck's others, worth being able to reference
% by its footer slide ID like any other). Opening/closing curly quotes are
% added automatically (write the quote text with none of its own); the
% second argument is the person it's from, set smaller/plain, below the
% yellow box and right-aligned to its own right edge (not the page's usual
% text margin) -- pass an empty second argument for an unattributed quote.
%
%   \quoteframe{One of the funny things about the stock market is that
%   	every time one person buys, another sells, and both think they
%   	are astute.}{William Feather, 1889-1981}
%
% \SMQuoteBoxHeight is a plain percentage of the page, unlike
% titleframe's/sectionframe's own (calibrated-to-land-on-an-exact-fraction)
% box heights, since nothing here needs to line up with another element's
% measured position the way those do. \SMTitleframeBoxWidth (defined
% above) is reused as-is for the box's width -- the same
% stop-short-of-the-logo formula, so no reason to duplicate it under a new
% name.
%
% Every piece here -- box, logo, quote text, attribution -- is drawn as a
% zero-space tikz overlay anchored to absolute page coordinates (the same
% \vbox to 0pt{\vss ...} pattern used elsewhere in this file for exactly
% this reason), rather than any of it flowing through the frame's normal
% vertical list. This sidesteps needing to know/counteract beamer's own
% default top-of-frame spacing (which standardframe explicitly measures and
% cancels via \SMContentTopShift) -- there's no such frame-relative
% positioning to get right here, since literally everything is placed by
% explicit coordinate instead. Like the other page-corner overlays in this
% file, "current page" doesn't resolve correctly on a document's very first
% compile pass, so (like those) this needs the usual second pass too.
%
% A *command*, not an environment -- see the identically-reasoned comment on
% titleframe just above for why: the same beamercolorbox-spanning-a-custom-
% environment's-begin/end problem would apply here too, for the same reason.
% \@empty (the empty-attribution check below) needs \makeatletter.
% \quotesource{text} -- opt-in \source{} for \quoteframe (call *before*
% \quoteframe{}{}, same set-then-consume pattern as \chapterlogo{}
% above), needed because \quoteframe is a *command* that auto-closes its
% own \end{frame} -- a bare \source{...} typed after \quoteframe{}{} in
% the calling slide file (as if it were still inside the frame, the way
% \source{} reads everywhere else in this deck) is actually issued
% *outside* any frame, and \source{}'s own \footnote{} then has no page to
% attach to yet. Confirmed directly (2026-07-12, direct instruction:
% traced a spurious extra blank leading overlay step on
% slides/c4c-Macbeth-Spot-Removers-Proposed-Structure-current.tex -- the
% very next multi-overlay frame after
% slides/w8w-Pie-Model-of-Capital-Structure-current.tex's own bare
% \source{} -- down to that orphaned footnote specifically, via a reduced
% two-slide test deck): the dangling footnote insert doesn't render
% anywhere visible on w8w itself, but silently corrupts beamer's own
% overlay-counting pre-pass on whichever frame happens to come right
% after -- reproduced with w8w+c4c alone (bug present) vs. c4c alone (bug
% absent), and confirmed it's specifically about the footnote being
% orphaned, not about w8w's content otherwise, by removing just its
% trailing \source{} line (bug also disappears). A plain \clearpage
% between the two frames also "fixes" c4c but only by giving the orphaned
% footnote a page of its own to attach to -- confirmed directly this
% splits w8w itself into two pages (real content, then a second page with
% nothing but the footnote/footer), not an acceptable trade. The only
% correct fix is keeping the footnote inside the frame that's actually
% citing something, i.e. before \quoteframe's own \end{frame} -- exactly
% what \quotesource{} below does.
\newcommand{\SMQuoteSource}{}
\newcommand{\quotesource}[1]{\gdef\SMQuoteSource{#1}}
\makeatletter
\newcommand{\quoteframe}[2]{%
	\SMOpenFrame{[t]}%
	% Disable hyphenation for the quote text (2026-07-26, direct
	% instruction: "do not allow words to be broken across lines" --
	% confirmed directly first, e.g. on `n3q`'s own long quote,
	% "limi-tations"/"finan-cial" both landed as real mid-word line
	% breaks). Same standard, sufficient technique already used for
	% `sectionframe`'s own title (see that environment's own comment,
	% right before \SMSectionframeTitleBox) -- setting the penalties
	% to 10000 makes TeX treat a hyphenated break as forbidden rather than
	% just discouraged, so one is never chosen even over an otherwise
	% ragged/short line; the hyphenmin lengths are redundant insurance on
	% top of that. Scoped to this frame only via the enclosing
	% \begin{frame}/\end{frame} group, same as sectionframe's own
	% \begingroup/\endgroup scoping -- doesn't leak into any other frame
	% type.
	\hyphenpenalty=10000
	\exhyphenpenalty=10000
	\lefthyphenmin=100
	\righthyphenmin=100
	% \SMArtifactBegin/\SMArtifactEnd (defined with the footer, "PDF
	% accessibility" comment): same reasoning as titleframe's own background
	% box above -- purely decorative, marked so a screen reader skips it. The
	% quote text and attribution below stay unmarked (real content); the
	% footer (this frame keeps one, unlike titleframe) is already marked via
	% the global footline hook.
	\SMArtifactBegin{/Type /Background}%
	\vbox to 0pt{\vss
	\begin{tikzpicture}[overlay, remember picture]
		% \SMTitleBoxColor, not a literal color token -- same setting
		% titleframe's own box and the standard frametitle bar (colorA) use
		% (2026-07-16, direct instruction).
		% outer sep=0pt here and on the logo/attribution below: same hairline-
		% inset-from-the-anchored-corner issue as titleframe's own box -- see
		% the frametitle hook's comment for the full story.
		\node[fill=\SMTitleBoxColor, anchor=north west, inner sep=0pt, outer sep=0pt,
		      minimum width=\SMTitleframeBoxWidth,
		      minimum height=\SMQuoteBoxHeight]
			at (current page.north west) {};
	\end{tikzpicture}%
	}%
	\SMArtifactEnd
	\vbox to 0pt{\vss
	\begin{tikzpicture}[overlay, remember picture]
		% inner ysep derived, not \SMFrametitleRightPadding -- see the
		% frametitle hook's fixed-position logo redraw for why (matches the
		% inline logo's vertically-centered position on every other frame).
		\node[anchor=north east, inner xsep=\SMFrametitleRightPadding,
		      inner ysep=\dimexpr\SMTitleBarBottomAlign-\SMLogoHeight\relax,
		      outer sep=0pt]
			at (current page.north east) {\SMInstituteLogo};
	\end{tikzpicture}%
	}%
	\vbox to 0pt{\vss
	\begin{tikzpicture}[overlay, remember picture]
		% \SMTitleTextColor, not a literal color token -- follows the box
		% color change right above (2026-07-16, direct instruction). Set
		% *inside* the node's own {...} content group (not a stored box
		% register, unlike titleframe's own vbox above), so this doesn't hit
		% the same "\color leaks out of a stored \vbox" pitfall documented on
		% \titleframe's own closing code -- confirmed directly by checking
		% the frame right after this one for the same symptom.
		% The quote text is set in a \parbox instead of the node's own
		% `text width` (2026-09-24), so #1 can hold reveal steps
		% (\conditionalpause, \uncover, \pause). Two separate failures
		% with those directly inside the node, both reproduced in plain
		% beamer and TikZ (no project code): (1) TikZ reads a `text width`
		% node's text through an \aftergroup hook that a \pause/\uncover
		% inside it breaks -- "Package tikz Error: Giving up on this path.
		% Did you forget a semicolon?" on every overlay page; (2) even
		% with (1) avoided, a text-less \pause starts a "covered" state
		% that is only ended by a *global* \beamer@endpause at the next
		% \pause or the end of the frame, so everything after the quote --
		% the footer included -- stayed covered until the last step.
		% \beamer@endpause is therefore run at the end of the quote (and
		% cleared), which confines the covering to the quote itself.
		%
		% The parbox reproduces what TikZ builds for `text width`
		% (minipage[t], \leavevmode\hbox{}, then \tikz@text@action --
		% ragged-right with its own \rightskip/\spaceskip -- and the
		% paragraph ended *after* the \Large group, i.e. with the normal
		% line spacing) so that existing quote slides are unchanged:
		% verified pixel-identical on all 18 other \quoteframe slides.
		\node[anchor=center, inner sep=0pt]
			at ($(current page.north west)+(0.5\SMTitleframeBoxWidth,-0.5\SMQuoteBoxHeight)$)
			{\parbox[t]{\dimexpr\SMTitleframeBoxWidth-2\SMFrametitleLeftPadding\relax}{%
				\leavevmode\hbox{}\tikz@textfont\tikz@text@action%
				{\color{\SMTitleTextColor}\Large\itshape\textquotedblleft #1\textquotedblright%
				\beamer@endpause\global\let\beamer@endpause\@empty}\par}};
	\end{tikzpicture}%
	}%
	\def\SMQuoteAttributionTemp{#2}%
	\ifx\SMQuoteAttributionTemp\@empty\else
		\vbox to 0pt{\vss
		\begin{tikzpicture}[overlay, remember picture]
			\node[anchor=north east, inner sep=0pt, outer sep=0pt]
				at ($(current page.north west)+(\SMTitleframeBoxWidth,-\SMQuoteBoxHeight-\SMQuoteAttributionGap)$)
				{\normalfont\itshape\color{c_custom_gray} #2};
		\end{tikzpicture}%
		}%
	\fi
	\ifx\SMQuoteSource\@empty\else
		\source{\SMQuoteSource}%
	\fi
	\end{frame}%
	% Cleared here, *after* \end{frame}, not inside the frame body above
	% (2026-09-24): beamer re-runs the collected body once per reveal step,
	% so clearing it in the body left only the first step with a source
	% line -- the quote's \quotesource vanished on every later step of a
	% quote with \conditionalpause. Code after \end{frame} runs once.
	\global\let\SMQuoteSource\@empty
}
\makeatother

% figureframe: a slide that is just a figure (ExampleSlides.pptx slide 5) --
% same standardframe layout (title, footer, \source{} for the usual
% "Source: ..."/"Image source: ..." line), but with the figure positioned
% via \figurecenter (call before the figure, inside the frame body) rather
% than left-aligned, which -- like the source deck's own figure slide -- is
% the default:
%
%   \begin{figureframe}{Correlation between DJIA stock returns over 15 years}
%   	\includegraphics[width=0.5\linewidth]{correlation-matrix}
%   	\source{Source: Own illustration, data from Yahoo! Finance, 24.02.2026.}
%   \end{figureframe}
%
%   \begin{figureframe}{Title}
%   	\figurecenter
%   	\includegraphics[width=0.6\linewidth]{diagram}
%   \end{figureframe}
%
% Vertical placement is standardframe's own [t]/[c] (its existing optional
% argument, passed straight through as figureframe's own -- so
% \begin{figureframe}[c]{Title} vertically centers exactly like
% \begin{standardframe}[c]{Title} already does); \figurecenter only ever
% affects the horizontal axis.
%
% \figurecenter is plain \centering under a more discoverable name, nothing
% more -- *not* an opt-in-before-the-frame flag like \chapterlogo{}/
% \sidefootnotes (those must run before \begin{standardframe}, since they
% configure something the frame's own opening code reads). This is called
% *inside* the body instead, exactly where it takes effect, so it can just
% take effect immediately like any other formatting declaration -- no flag,
% no state to reset afterwards, and (checked by first writing it as a flag
% read in figureframe's own opening code, mirroring \sidefootnotes) no way
% for the frame's opening code to see it anyway, since that code runs
% before the body -- and \figurecenter -- does.
\newcommand{\figurecenter}{\centering}
\newenvironment{figureframe}[2][t]{%
	\begin{standardframe}[#1]{#2}%
}{%
	\end{standardframe}%
}

% highlightframe: a standardframe with two things about it made unmistakably
% different -- the title bar's *and* the footer's background color, both
% together -- so a presenter can never flip to it without immediately
% noticing (ExampleSlides.pptx slide 6 was the inspiration for wanting a
% standardframe variant here, though the source slide itself has no such
% color cue; this is a new addition, not a conversion). Everything else --
% body, logo, bullet styling -- is exactly standardframe's own, via the same
% \begingroup/\endgroup wrapper sectionframe above uses for the title-color
% half of this. \SMHighlightColor is a \colorlet, not a color baked
% directly into the environment, so retargeting every highlightframe in a
% deck to a different color -- title bar and footer both -- is still the
% same one-line change it was when this only touched the title bar.
%
% The footer box needs the \global\def dance (\SMFooterBoxColor, defined
% above with the frametitle hook's other opt-in flags, reset in
% standardframe's own closing code) that \setbeamercolor{frametitle} doesn't:
% that assigns through beamer's own color-management system, which handles
% the frame-closing group correctly on its own (checked -- see
% \SMChapterLogoPath's comment for why plain macros need \global\def but
% beamer's own set-commands don't), but the footer box's fill is read from a
% plain macro inside a tikz key, same category as \SMChapterLogoPath.
\colorlet{SMHighlightColor}{c_primary_gray}
\newenvironment{highlightframe}[2][t]{%
	\begingroup
	\setbeamercolor{frametitle}{bg=SMHighlightColor}
	\global\def\SMFooterBoxColor{SMHighlightColor}%
	\begin{standardframe}[#1]{#2}%
}{%
	\end{standardframe}%
	\endgroup
}


% The strip is an overlay (a zero-size box: the Top Margin measurement does not see it), so
% the body below it starts where the title bar's body would -- under the strip. The frame's
% own target line is therefore the Top Margin plus the strip's height
% \SMProgressStripHeight (the strip's bottom edge minus the title bar's: measured 11.08pt in
% the PDF; the node is 11.19pt minimum height): the first line sits as far below the strip
% as it sits below the title bar on any other frame, whatever the body starts with (v5.30.5; the
% frame used to be exempt from the measurement, with a fixed -3pt, which left a paragraph
% 10pt and a list 14pt below the strip against the usual 16pt -- direct instruction, slide 1j2).
\makeatletter
\newcommand{\SM@topstrip}{%
	\ifluatex
		\directlua{slidemage_measure.top_extra(\SM@topthis, \strip@pt\dimexpr\SMProgressStripHeight\relax)}%
	\fi
}
\makeatother
% \progressstrip{Step 1,Step 2,Step 3}{k} -- a full-width strip in the
% institution's tertiary yellow, flush below the title bar, that lists a
% short chain of steps joined by arrows and underlines the k-th one to mark
% where the slide sits in that chain (added 2026-09-19 for the Publishing
% masterdeck's "Field -> Topic -> Question" strip, which the source
% PowerPoint's "Progress" layout draws directly below the title bar; revised
% the same day, direct instruction: edge-to-edge, flush under the title bar,
% smaller font, tight vertical padding, secondary yellow). Written in the
% frame body, as the very first thing. The strip itself is a zero-space
% absolute overlay (so it can span the whole page width, which in-flow text
% cannot); the \vspace after it reserves its height so the body starts just
% below it.
\newcommand{\progressstrip}[2]{%
	\csname SM@topstrip\endcsname% the text below the strip: Top Margin + the strip's height (measured, like any frame)
	\vbox to 0pt{\vss
	\begin{tikzpicture}[overlay, remember picture]
		\node[fill=c_tertiary_yellow, anchor=north west, outer sep=0pt,
		      inner xsep=\SMFrametitleLeftPadding, inner ysep=1.5pt, text height=6.2pt, text depth=2pt,
		      minimum height=11.19pt,
		      text width=\dimexpr\paperwidth-2\SMFrametitleLeftPadding\relax,
		      align=left, text=c_primary_dark, font=\fontsize{6.3}{7.5}\selectfont]
			at ($(current page.north west)+(0,-\SMTitleBarBottomAlign)$)
			{\foreach \SMStepName [count=\SMStepIndex] in {#1}{%
				\ifnum\SMStepIndex>1\ $\rightarrow$\ \fi
				\ifnum\SMStepIndex=#2\underline{\SMStepName}\else\SMStepName\fi}};
	\end{tikzpicture}%
	}%
	\vspace{-3pt}%
}

% \SMCodeFile[listings keys]{path} -- LaTeX source shown on a slide, read from
% a file (2026-10-04, for the SlideMage manual: "show the result and the
% source together"). The code lives in its own small file, e.g.
% snippets/<slide id>-1.tex at the project root, so the slide itself needs no
% [fragile] frame and the code is exactly what was compiled elsewhere. Found on
% the TeX search path like any \input (the project root is on it, see
% .latexmkrc). Optional keys go to \lstinputlisting (firstline=, lastline=,
% basicstyle=...). Fira Mono Regular (the theme's \ttfamily is Fira Mono Bold),
% control sequences, comments and the panel in the SMCode* colours below
% (set per template). Write snippets in ASCII; the few non-ASCII characters below are mapped.
% SlideMage's transcript shows the file as a code block, and the pre-render
% check reports a missing file.
\usepackage{listings}
% Every control sequence, backslash included, is a keyword: the backslash is
% made a letter and is the keyword prefix (listings' prefix keywords need a
% non-empty keyword list to be active, hence the dummy). An escaped \% and a
% \\ are mapped literally (\% must not start a comment; \\ otherwise loses a
% backslash to the \% mapping) and reset the prefix flag, which would
% otherwise colour the word after them.
%
% Code colours (2026-10-05, direct instruction: "set them to something more
% readable in the institutions' template files"): three named colours a
% template defines before \input{SlideMage_base} --
% SMCodeKeywordColor (control sequences), SMCodeCommentColor (% comments),
% SMCodeBackgroundColor (the panel), SMCodeTextColor (ordinary code). A template that doesn't keeps the
% original look: highlight colour, primary gray, tertiary gray -- which in
% the Graz/MCI/Innsbruck palettes is yellow/orange and light gray on light
% gray, about 1.1-1.3:1 contrast; their templates set readable colours.
\makeatletter
\@ifundefinedcolor{SMCodeKeywordColor}{\colorlet{SMCodeKeywordColor}{c_primary_highlight}}{}
\@ifundefinedcolor{SMCodeCommentColor}{\colorlet{SMCodeCommentColor}{c_primary_gray}}{}
\@ifundefinedcolor{SMCodeBackgroundColor}{\colorlet{SMCodeBackgroundColor}{c_tertiary_gray}}{}
\@ifundefinedcolor{SMCodeTextColor}{\colorlet{SMCodeTextColor}{c_primary_dark}}{}
\newcommand{\SMCodePercent}{\textbackslash\%\global\let\lst@prefixkeyword\relax}
\newcommand{\SMCodeNewline}{\textcolor{SMCodeKeywordColor}{\textbackslash\textbackslash}\global\let\lst@prefixkeyword\relax}
\makeatother
\lstdefinestyle{SMCode}{
	basicstyle=\SMCodeFont\SMCodeSize\color{SMCodeTextColor},
	columns=fullflexible, keepspaces=true, tabsize=2,
	keywords={SMCodeNoSuchKeyword},
	alsoletter={\\}, keywordsprefix=\\,
	keywordstyle=\color{SMCodeKeywordColor},
	morecomment=[l]\%, commentstyle=\color{SMCodeCommentColor},
	showstringspaces=false, upquote=true,
	breaklines=true, breakindent=1.5em,
	aboveskip=0pt, belowskip=0pt,
	literate={\\\%}{{\SMCodePercent}}2 {\\\\}{{\SMCodeNewline}}2 {ä}{{\"a}}1 {ö}{{\"o}}1 {ü}{{\"u}}1 {Ä}{{\"A}}1 {Ö}{{\"O}}1 {Ü}{{\"U}}1
		{ß}{{\ss}}1 {é}{{\'e}}1 {–}{{--}}1 {—}{{---}}1 {·}{{$\cdot$}}1 {…}{{...}}1
		{→}{{$\rightarrow$}}1,
}
% The panel is a \colorbox around the listing, not listings' own frame and
% background: listings draws its frame outside the box it typesets, so a
% panel stuck out above the top of its column by 4-8pt (depending on the font
% size) and missed the Top Margin (v5.14.1, found by check_layout.py in the
% SlideMage manual). The box holds everything it draws.
\newcommand{\SMCodeFile}[2][]{\par\noindent{\setlength{\fboxsep}{3pt}%
	\colorbox{SMCodeBackgroundColor}{\begin{minipage}{\dimexpr\linewidth-2\fboxsep\relax}%
		\lstinputlisting[style=SMCode,#1]{#2}%
	\end{minipage}}}\par}

% \SMCode{text} / \SMCmd{name} -- code in running text: Fira Mono Regular at
% 90% (\texttt is the theme's Fira Mono Bold, which reads heavy and wide in a
% sentence of Fira Sans Light). \SMCmd{slide} prints \slide. The argument is
% typeset, so write \_ \{ \} \# as in any text. The transcript shows both as
% `code`. Ligatures=TeXOff: code shows -- and `` as typed, not as an en dash
% and a curly quote (fontspec's default for \newfontfamily).
\newcommand{\SMInlineCodeStyle}{}
\newcommand{\SMCode}[1]{{\SMInlineCodeFont\SMInlineCodeStyle #1}}
\newcommand{\SMCmd}[1]{{\SMInlineCodeFont\SMInlineCodeStyle\textbackslash #1}}

% \SMQRCode[height]{image}{caption} -- a QR code SlideMage made (2026-10-25,
% direct instruction: the Source view's Insert QR, the assistant's
% make_qr_code; src/slidemage/qr_code.py), the code at [height] (default
% 0.6\textheight: a lecture hall's back row needs about that) with an
% optional caption centred under it, as wide as the code -- the link, or a
% text such as "QR 1" that the slide refers to; {} for none. An inline box:
% several codes go side by side or in columns, and
% \cornerfigure{\SMQRCode{...}{...}} puts one in a corner. The caption's
% look is \SMQRCodeCaptionStyle, which a template may redefine.
\newsavebox{\SMQRCodeBox}
\providecommand{\SMQRCodeCaptionStyle}{\scriptsize}
\NewDocumentCommand{\SMQRCode}{O{0.6\textheight} m m}{%
	\begingroup
	\sbox{\SMQRCodeBox}{\includegraphics[height=#1]{#2}}%
	\begin{minipage}[t]{\wd\SMQRCodeBox}%
		\centering\usebox{\SMQRCodeBox}%
		\if\relax\detokenize{#3}\relax\else
			\par{\SMQRCodeCaptionStyle #3\par}%     the code's own white border is the gap
		\fi
	\end{minipage}%
	\endgroup
}

% Navigation for readers (2026-10-04, SlideMage manual: readers must find
% their way without knowing that the logo is a hidden link). Three opt-in
% pieces, each usable in any deck:
%
% \SMChapterAgenda[width] -- the slides of the current \section as a
%   two-column list, each entry a link to that slide (its frame title, or the
%   slide's \SMAgendaTitle{...}, or else its slug with spaces). Meant for the
%   body of a chapter's sectionframe (the manual puts one on every divider).
%   [width] (default \linewidth) leaves room for e.g. a \sectionfigure. Set in
%   \normalsize, else \small, else \footnotesize, the first that fits
%   \SMChapterAgendaMaxHeight.
%   Works from the .aux file like a table of contents: every slide writes
%   \SMAgendaSlide{section}{id}{label} (in \setSlideID) and its frame title
%   (\SMAgendaEntryTitle, in the frametitle hook), and the agenda shows what
%   the previous pass recorded -- latexmk's rerun fills it in. A slide shown
%   twice (\slidestep) is listed once, at its first place; the agenda's own
%   slide is not listed.
% \SMSubsectionAgenda[width] -- the same, but one entry per \subsection of
%   the current \section (2026-10-12): the subsection's title, linked to its
%   first slide in this version. For a chapter divider that introduces the
%   subsections instead of listing every slide. Slides before the chapter's
%   first \subsection are not listed; it fills in on the rerun.
% \SMAgendaTitle{text} -- in a slide file, before the frame: the entry text
%   for a slide whose frame has no title (\quoteframe, \fullbleedframe, ...)
%   or should read differently in the agenda.
% \SMDeckAgenda -- marks the deck's contents slide (put it in that slide's
%   body, next to \tableofcontents) as the target of the navigation bar's
%   first link.
% \SMNavigationBar[deck,chapter,previews] -- in the deck preamble: the footer
%   then reads "\FooterText | Navigation: <deck> · <chapter> · <previews>",
%   three links to the \SMDeckAgenda slide, the current chapter's
%   \SMChapterAgenda (the contents slide before the first \section, or where
%   a chapter has none) and the generated appendix agenda with the slide
%   previews. A link whose target is not in this PDF is left out: the last
%   one exists only when SlideMage generates the appendix (the Generate
%   appendix TOC checkbox, or \deckappendixtoc in the deck). Default labels
%   "Contents,Chapter contents,Slide previews"; the word before them is
%   \SMNavigationLead ("Navigation:"). Off by default: other decks' footers
%   and the logo link are unchanged.
% \deckappendixtoc -- in the deck preamble, inert for LaTeX (like
%   \decktags): SlideMage then generates the appendix for this deck on every
%   whole-deck compile, whatever the toolbar checkbox says.
\newcommand{\deckappendixtoc}{}
\newlength{\SMChapterAgendaMaxHeight}
\newcommand{\SMNavigationLead}{Navigation:}
\newcommand{\SMNavigationLabelDeck}{Contents}
\newcommand{\SMNavigationLabelChapter}{Chapter contents}
\newcommand{\SMNavigationLabelPreviews}{Slide previews}
\makeatletter
% -- recording (writes to the .aux) --
\newcommand{\SM@agenda@slide}[2]{%
	\ifcsname SM@ag@written@#1\endcsname\else
		\expandafter\gdef\csname SM@ag@written@#1\endcsname{}%
		% label from the rest of the id-slug ("EN-Slide-IDs-"): language
		% segment and trailing hyphen dropped, hyphens to spaces. (No % in
		% the Lua code -- TeX would read it as a comment -- and no hash
		% sign, which is a macro parameter here.)
		\edef\SM@ag@label{\directlua{
			local s = "\luaescapestring{#2}"
			while string.sub(s, -1) == "-" do s = string.sub(s, 1, -2) end
			for _, p in ipairs({"EN-", "DE-", "mult-"}) do
				if string.sub(s, 1, string.len(p)) == p then s = string.sub(s, string.len(p) + 1) end
			end
			s = string.gsub(s, "[-]", " ")
			tex.sprint(-2, s)}}%
		\immediate\write\@auxout{\string\SMAgendaSlide{\the\c@section}{#1}{\SM@ag@label}}%
	\fi
	% the first slide of each \subsection and the subsection's title, for
	% \SMSubsectionAgenda (its own guard: a slide shown twice must not hide
	% that it also opens a later subsection)
	\ifnum\c@subsection>0
		\ifcsname SM@ag@subwritten@\the\c@section.\the\c@subsection\endcsname\else
			\expandafter\gdef\csname SM@ag@subwritten@\the\c@section.\the\c@subsection\endcsname{}%
			\immediate\write\@auxout{\string\SMAgendaSubsection{\the\c@section}{#1}{\detokenize\expandafter{\subsecname}}}%
		\fi
	\fi
}
\newcommand{\SM@agenda@title}[1]{%
	\ifcsname SM@ag@titled@\slideID\endcsname\else
		\expandafter\gdef\csname SM@ag@titled@\slideID\endcsname{}%
		\immediate\write\@auxout{\string\SMAgendaEntryTitle{\slideID}{\detokenize{#1}}}%
	\fi
}
\newcommand{\SMAgendaTitle}[1]{\SM@agenda@title{#1}}
\newcommand{\SM@agenda@frametitle}{%
	\ifx\beamer@frametitle\@empty\else\expandafter\SM@agenda@title\expandafter{\beamer@frametitle}\fi
}
\newcommand{\SMAgendaFrametitleHook}{\SM@agenda@frametitle}% called by \SMFrametitleHookCode
\newcommand{\SMDeckAgenda}{%
	\hypertarget{SMDeckAgenda}{}%
	\immediate\write\@auxout{\string\SMAgendaDeckTarget}%
}
% -- reading (the .aux, at \begin{document}; its second reading at
% \end{document} changes nothing) --
\newcommand{\SMAgendaSlide}[3]{%
	\ifcsname SM@ag@seen@#2\endcsname\else
		\expandafter\gdef\csname SM@ag@seen@#2\endcsname{}%
		\listcsgadd{SM@ag@list@#1}{#2}%
		\expandafter\gdef\csname SM@ag@label@#2\endcsname{#3}%
	\fi
}
\newcommand{\SMAgendaEntryTitle}[2]{%
	\ifcsname SM@ag@title@#1\endcsname\else
		\expandafter\gdef\csname SM@ag@title@#1\endcsname{#2}%
	\fi
}
% {section}{id of the subsection's first slide}{subsection title}
\newcommand{\SMAgendaSubsection}[3]{%
	\ifcsname SM@ag@subseen@#2\endcsname\else
		\expandafter\gdef\csname SM@ag@subseen@#2\endcsname{}%
		\listcsgadd{SM@ag@sublist@#1}{#2}%
		\expandafter\gdef\csname SM@ag@subtitle@#2\endcsname{#3}%
	\fi
}
\newcommand{\SMAgendaDeckTarget}{\gdef\SM@ag@decktarget{}}
\newcommand{\SMAgendaChapterTarget}[1]{\expandafter\gdef\csname SM@ag@chtarget@#1\endcsname{}}
% -- the chapter agenda --
\newcount\SM@ag@n
\newcount\SM@ag@i
\newbox\SM@ag@box
\newcommand{\SM@ag@entry}[1]{%
	\item\hyperlink{#1}{\ifcsname SM@ag@title@#1\endcsname\csname SM@ag@title@#1\endcsname
		\else\csname SM@ag@label@#1\endcsname\fi}%
}
% a subsection agenda's entry: #1 is the id of the subsection's first slide
\newcommand{\SM@ag@subentry}[1]{%
	\item\hyperlink{#1}{\csname SM@ag@subtitle@#1\endcsname}%
}
% (\def + \ifx, not \ifstrequal: that one would compare the *name* of
% \SM@ag@self, see \decklanguage)
\newcommand{\SM@ag@count}[1]{\def\SM@ag@this{#1}\ifx\SM@ag@this\SM@ag@self\else\advance\SM@ag@n by 1 \fi}
\newcommand{\SM@ag@split}[1]{%
	\def\SM@ag@this{#1}%
	\ifx\SM@ag@this\SM@ag@self\else
		\advance\SM@ag@i by 1
		\ifnum\SM@ag@i>\SM@ag@half
			\gappto\SM@ag@right{\SM@ag@entrycmd{#1}}%
		\else
			\gappto\SM@ag@left{\SM@ag@entrycmd{#1}}%
		\fi
	\fi
}
\newlength{\SM@ag@maxh}
\newcommand{\SM@ag@typeset}[2]{% #1 width, #2 font switch
	\setbox\SM@ag@box=\hbox{\begin{minipage}[t]{#1}%
		#2%
		\begin{columns}[t,onlytextwidth]
			\begin{column}{0.48\linewidth}
				\ifx\SM@ag@left\@empty\else\begin{itemize}\SM@ag@left\end{itemize}\fi
			\end{column}
			\begin{column}{0.48\linewidth}
				\ifx\SM@ag@right\@empty\else\begin{itemize}\SM@ag@right\end{itemize}\fi
			\end{column}
		\end{columns}%
	\end{minipage}}%
}
% \SMChapterAgenda lists the chapter's slides, \SMSubsectionAgenda its
% subsections (2026-10-12, direct instruction: "an agenda slide that lists
% all the subsections in a chapter"; each entry links to the subsection's
% first slide, titled as its \subsection{}). Same two columns, font steps and
% chapter target ("Chapter contents" in the navigation bar).
\newcommand{\SMChapterAgenda}[1][\linewidth]{%
	\SM@ag@main{SM@ag@list@}{\SM@ag@entry}{#1}%
}
\newcommand{\SMSubsectionAgenda}[1][\linewidth]{%
	\SM@ag@main{SM@ag@sublist@}{\SM@ag@subentry}{#1}%
}
% #1 the name prefix of the recorded list (+ section number), #2 the entry
% macro, #3 the width
\newcommand{\SM@ag@main}[3]{%
	\let\SM@ag@entrycmd#2%
	\edef\SM@ag@sec{\the\c@section}%
	\edef\SM@ag@self{\slideID}%
	\hypertarget{SMChapterAgenda\SM@ag@sec}{}%
	\immediate\write\@auxout{\string\SMAgendaChapterTarget{\SM@ag@sec}}%
	\SM@ag@n=0
	\ifcsdef{#1\SM@ag@sec}{\forlistcsloop{\SM@ag@count}{#1\SM@ag@sec}}{}%
	\ifnum\SM@ag@n>0
		\SM@ag@i=0
		\edef\SM@ag@half{\the\numexpr(\SM@ag@n+1)/2\relax}% ceil(n/2); \numexpr rounds
		\ifodd\SM@ag@n\else\edef\SM@ag@half{\the\numexpr\SM@ag@n/2\relax}\fi
		\gdef\SM@ag@left{}\gdef\SM@ag@right{}%
		\forlistcsloop{\SM@ag@split}{#1\SM@ag@sec}%
		\ifdim\SMChapterAgendaMaxHeight=0pt
			% the room below a sectionframe's title bar, down to the Bottom Margin
			\setlength{\SM@ag@maxh}{\dimexpr\paperheight-\SMSectionframeBarBottom
				-\SMSectionframeContentDrop-\SMContentBottomAlign-14pt\relax}%
		\else
			\setlength{\SM@ag@maxh}{\SMChapterAgendaMaxHeight}%
		\fi
		\SM@ag@typeset{#3}{\normalsize}%
		\ifdim\dimexpr\ht\SM@ag@box+\dp\SM@ag@box\relax>\SM@ag@maxh
			\SM@ag@typeset{#3}{\small}%
			\ifdim\dimexpr\ht\SM@ag@box+\dp\SM@ag@box\relax>\SM@ag@maxh
				\SM@ag@typeset{#3}{\footnotesize}%
			\fi
		\fi
		\noindent\box\SM@ag@box\par
	\fi
}
% -- the navigation bar (footer) --
\newcommand{\SMNavigationBar}[1][\SMNavigationLabelDeck,\SMNavigationLabelChapter,\SMNavigationLabelPreviews]{%
	\global\SMNavigationBartrue
	\SM@nav@labels#1,,,\@nil
}
\def\SM@nav@labels#1,#2,#3,#4\@nil{%
	\gdef\SMNavigationLabelDeck{#1}\gdef\SMNavigationLabelChapter{#2}\gdef\SMNavigationLabelPreviews{#3}%
}
% Link text in the footer: the footer's own text colour, underlined in
% c_secondary_blue (the \href underline, see \SMHrefUnderline).
\newcommand{\SMLinkUnderline}[1]{%
	\sbox{\SMHrefTextBox}{#1}%
	\usebox{\SMHrefTextBox}%
	\hspace{-\wd\SMHrefTextBox}%
	\SMHrefUnderlineMeasureRefHeights
	\setbox\SMHrefUnderlineBox=\hbox{M}\SMHrefUnderlineNowHeight=\ht\SMHrefUnderlineBox
	\SMHrefUnderlineSetThickness
	\raisebox{-0.35ex}[0pt][0pt]{\textcolor{c_secondary_blue}{\rule{\wd\SMHrefTextBox}{\SMHrefUnderlineThickness}}}%
}
\newcommand{\SM@nav@link}[2]{{\hypersetup{linkcolor=\SMFooterTextColor}\hyperlink{#1}{\SMLinkUnderline{#2}}}}
% \SMLink{target}{text} -- a link to a slide (or any \hypertarget) in running
% text, underlined like an \href so a reader sees it is a link: a bare
% \hyperlink{target}{text} has no underline. One line only, like \href's
% underline: keep the link text short. Buttons still use \hyperlink with
% \beamergotobutton / \beamerreturnbutton.
\newcommand{\SMLink}[2]{\hyperlink{#1}{\SMLinkUnderline{#2}}}
\newcommand{\SM@nav@sep}{\ \textperiodcentered\ }
\newcommand{\SMNavigationLinks}{%
	\def\SM@nav@chapter{}%
	\ifnum\c@section>0
		\ifcsname SM@ag@chtarget@\the\c@section\endcsname
			\edef\SM@nav@chapter{SMChapterAgenda\the\c@section}%
		\fi
	\fi
	\ifx\SM@nav@chapter\@empty
		\ifdefined\SM@ag@decktarget\def\SM@nav@chapter{SMDeckAgenda}\fi
	\fi
	\def\SM@nav@any{}%
	\ \textbar\ \SMNavigationLead\ %
	\ifdefined\SM@ag@decktarget
		\SM@nav@link{SMDeckAgenda}{\SMNavigationLabelDeck}\def\SM@nav@any{1}%
	\fi
	\ifx\SM@nav@chapter\@empty\else
		\ifx\SM@nav@any\@empty\else\SM@nav@sep\fi
		\SM@nav@link{\SM@nav@chapter}{\SMNavigationLabelChapter}\def\SM@nav@any{1}%
	\fi
	\ifdefined\SMAppendixTocOverride
		\ifx\SM@nav@any\@empty\else\SM@nav@sep\fi
		\SM@nav@link{SMAppendixAgenda}{\SMNavigationLabelPreviews}%
	\fi
}
\makeatother

% ---- Bibliography (biblatex + biber, opt-in per deck) ----
% \SMUseBibliography[file.bib] -- in a deck's preamble (after the institution
% template is \input) instead of \usepackage{biblatex} + \addbibresource: loads
% biblatex with \SMBiblatexOptions, registers the .bib file (default
% SlideMage.bib), then runs \SMBibliographyStyle. The template decides the
% citation format by redefining the two macros (2026-10-07, direct instruction:
% the house format had been copied into every course deck); a deck that wants
% another format redefines them before the call, or does not use the command.
% Defaults: biblatex's stock authoryear style.
%
% The .bib lookup: \IfFileExists below is the biber equivalent of the
% \input@path/\graphicspath fallback: biber runs as its own process and reads
% the literal path recorded in the .bcf at compile time, so it cannot follow
% \input@path -- resolving the correct relative path once, here, at LaTeX time
% (when the real cwd is known) is what makes biber find the .bib file whether
% compiling from the project root or from inside decks/. \input@path is
% emptied for the check itself: \IfFileExists honors \input@path too, so
% without this it reports the file as "found" via the ../ fallback and then
% writes that unresolved literal path into the .bcf, which biber (a separate
% process that never sees \input@path) then cannot find.
\makeatletter
\newcommand{\SMBiblatexOptions}{style=authoryear, backend=biber}
\newcommand{\SMBibliographyStyle}{}
\newcommand{\SMUseBibliography}[1][SlideMage.bib]{%
	\expandafter\usepackage\expandafter[\SMBiblatexOptions]{biblatex}% (options are not expanded otherwise)
	\let\SM@SavedInputPath\input@path
	\def\input@path{}%
	\IfFileExists{#1}{\addbibresource{#1}}{\addbibresource{../#1}}%
	\let\input@path\SM@SavedInputPath
	\SMBibliographyStyle}
\makeatother

%   ... additions only: dark-mode beamer colors, citation format, project macros
%
% Rules for a template:
%   1. Before this file: only \SMDarkMode (written into the deck by SlideMage).
%   2. Between this file and SlideMage_base: \setlength, \renewcommand,
%      \definecolor and \colorlet on the names in this catalog. Every name
%      exists already, so a misspelt \renewcommand is an error instead of a
%      setting that silently does nothing.
%   3. After SlideMage_base: additions only.
%
% Decks never load the beamer theme themselves; this file does. A length given
% in ex or em is evaluated where the template sets it, in the normal body font,
% after the theme has loaded its fonts.

\newcommand{\SMSetupLoaded}{}

% The beamer theme, with the body fonts.
\usetheme{gotham}

% ---------------------------------------------------------------------------
% Palette
% ---------------------------------------------------------------------------
% The SlideMage look; a template redefines what it changes (usually all of
% them). The framework combines these pairs, so a template should keep them
% readable (WCAG 2.1): body text c_primary_dark on each c_secondary_* block
% title color >= 4.5:1, title/footer text on its box color >= 4.5:1, white text
% on c_primary_gray (highlightframe), link underline c_secondary_blue on the
% canvas >= 3:1.
%
% SlideMage's own values: wherever possible a color the app itself uses
% (src/slidemage/palette.py, qt_deck_table.py, assets/splash.png). White
% title/footer text on c_primary_dark is 16.7:1, white on c_primary_gray 5.0:1.
\definecolor{c_primary_dark}{HTML}{1e1e1e}      % body text, dark boxes; splash screen background (palette dark "bg")
\definecolor{c_primary_highlight}{HTML}{0066cc} % accent: bullets, highlights; SlideMage button blue (Create PDF, progress bar, links)
\definecolor{c_primary_gray}{HTML}{6f6f6f}      % highlightframe background; palette dark "disabled_fg", dark enough for white text
\definecolor{c_primary_white}{HTML}{FFFFFF}
\definecolor{c_secondary_violet}{HTML}{ba68c8}  % block titles (block); palette dark "tab_deck_editor"
\definecolor{c_tertiary_violet}{HTML}{dcb4e4}   % block bodies; derived (50% white-blend)
% Derived, not an app color: the button blue is too dark for the dark text
% an infoblock title puts on it, the app's dark-mode link blue (4da6ff) too
% light for a link underline on white; this one sits between them and passes
% both (4.9:1 and 3.4:1).
\definecolor{c_secondary_blue}{HTML}{3d8fe0}    % infoblock titles, link underline
\definecolor{c_tertiary_blue}{HTML}{9ec7f0}     % derived (50% white-blend)
\definecolor{c_secondary_green}{HTML}{4caf50}   % exampleblock titles; palette dark "status_pass"
\definecolor{c_tertiary_green}{HTML}{81c784}    % palette dark "tab_slide_editor"
\definecolor{c_secondary_red}{HTML}{f04848}     % alertblock titles; palette dark "status_fail"
\definecolor{c_tertiary_red}{HTML}{f8a4a4}      % derived (50% white-blend)
\definecolor{c_secondary_orange}{HTML}{ffb74d}  % reference block titles; palette dark "tab_archived"
\definecolor{c_tertiary_orange}{HTML}{ffdba6}   % derived (50% white-blend)
\definecolor{c_secondary_yellow}{HTML}{ffd500}  % logo/splash yellow
\definecolor{c_tertiary_yellow}{HTML}{ffe066}   % deck editor's highlight yellow
\definecolor{c_tertiary_gray}{HTML}{e4e4e4}     % block bodies, process boxes; palette light "panel_bg"
% Supporting colors ("custom": not part of the palette proper).
% c_custom_gray is hyperlink text, button fill and the quote attribution.
\definecolor{c_custom_gray}{HTML}{2d2d2d}       % palette dark "panel_bg"
\definecolor{c_custom_red}{HTML}{cc0000}        % palette light "status_fail"
\definecolor{c_custom_green}{HTML}{1a8a1a}      % palette light "status_pass"
\definecolor{c_custom_brown}{HTML}{996600}      % palette light "syntax_brace"

% ---------------------------------------------------------------------------
% Roles: which palette color a part of the slide uses (color names, read
% when used, so they follow a redefined palette color)
% ---------------------------------------------------------------------------
\newcommand{\SMBulletColor}{c_primary_highlight}       % itemize bullets
\newcommand{\SMTitleBoxColor}{c_primary_dark}          % frame title bar
\newcommand{\SMTitleTextColor}{c_primary_white}        % frame title text
\newcommand{\SMFooterBoxColorDefault}{c_primary_dark}  % footer box
\newcommand{\SMFooterTextColor}{c_primary_white}       % footer text, slide id

% Optional colors a template may define; SlideMage_base derives each one from
% the palette when the template leaves it undefined:
%   SMCanvasColor          slide background                  (white)
%   SMCardBackgroundColor  SMCard fill                       (white)
%   SMCardBorderColor      SMCard border                     (c_primary_dark)
%   SMProcessBorderColor   SMProcess border, dark mode       (c_tertiary_gray)
%   SMCodeKeywordColor     \SMCodeFile commands              (c_primary_highlight)
%   SMCodeCommentColor     \SMCodeFile comments              (c_primary_gray)
%   SMCodeBackgroundColor  \SMCodeFile panel                 (c_tertiary_gray)
%   SMCodeTextColor        \SMCodeFile text                  (c_primary_dark)

% ---------------------------------------------------------------------------
% Logo
% ---------------------------------------------------------------------------
% A vector PDF cropped to its ink, found on the TeX search path (a project
% path such as schemes/<name>/logo.pdf, or a file next to the template).
\newcommand{\SMInstituteLogoPath}{slidemage-logo.pdf}
\newlength{\SMLogoHeight}
\setlength{\SMLogoHeight}{3ex}                         % logo height in the title bar

% ---------------------------------------------------------------------------
% Frame title bar
% ---------------------------------------------------------------------------
% Space around the frame title (the theme's registers under their SlideMage
% names). The right padding is measured against the compiled PDF for the logo:
% see SlideMage_base.tex's comment on \gothamFrametitleRightpading.
\let\SMFrametitleTopPadding\gothamFrametitleToppading
\let\SMFrametitleBottomPadding\gothamFrametitleBottompading
\let\SMFrametitleLeftPadding\gothamFrametitleLeftpading
\let\SMFrametitleRightPadding\gothamFrametitleRightpading
\setlength{\SMFrametitleRightPadding}{0.28cm}

%   \definecolor{c_primary_highlight}{HTML}{...}
%   \renewcommand{\SMInstituteLogoPath}{schemes/my_logo.pdf}
%   \setlength{\SMLogoHeight}{3ex}
%   % SPDX-FileCopyrightText: 2026 Stefan Palan, Jürgen Fleiß
% SPDX-License-Identifier: LPPL-1.3c+
% A template \input{}s SlideMage_setup (the theme and every template setting
% with its default) before this file; a template from before SlideMage v9.0.0
% does not, and SlideMage's Project compatibility update fixes it.
\ifdefined\SMSetupLoaded\else
	\PackageError{SlideMage}{The template does not \string% SPDX-FileCopyrightText: 2026 Stefan Palan, Jürgen Fleiß
% SPDX-License-Identifier: LPPL-1.3c+
% SlideMage setup layer and hook catalog (SlideMage v9.0.0).
%
% A template \input{}s this file first, sets whatever it wants to change from
% the catalog below, then \input{}s SlideMage_base:
%
%   % SlideMage template name: My University
%   \input{SlideMage_setup}
%   \definecolor{c_primary_highlight}{HTML}{...}
%   \renewcommand{\SMInstituteLogoPath}{schemes/my_logo.pdf}
%   \setlength{\SMLogoHeight}{3ex}
%   \input{SlideMage_base}
%   ... additions only: dark-mode beamer colors, citation format, project macros
%
% Rules for a template:
%   1. Before this file: only \SMDarkMode (written into the deck by SlideMage).
%   2. Between this file and SlideMage_base: \setlength, \renewcommand,
%      \definecolor and \colorlet on the names in this catalog. Every name
%      exists already, so a misspelt \renewcommand is an error instead of a
%      setting that silently does nothing.
%   3. After SlideMage_base: additions only.
%
% Decks never load the beamer theme themselves; this file does. A length given
% in ex or em is evaluated where the template sets it, in the normal body font,
% after the theme has loaded its fonts.

\newcommand{\SMSetupLoaded}{}

% The beamer theme, with the body fonts.
\usetheme{gotham}

% ---------------------------------------------------------------------------
% Palette
% ---------------------------------------------------------------------------
% The SlideMage look; a template redefines what it changes (usually all of
% them). The framework combines these pairs, so a template should keep them
% readable (WCAG 2.1): body text c_primary_dark on each c_secondary_* block
% title color >= 4.5:1, title/footer text on its box color >= 4.5:1, white text
% on c_primary_gray (highlightframe), link underline c_secondary_blue on the
% canvas >= 3:1.
%
% SlideMage's own values: wherever possible a color the app itself uses
% (src/slidemage/palette.py, qt_deck_table.py, assets/splash.png). White
% title/footer text on c_primary_dark is 16.7:1, white on c_primary_gray 5.0:1.
\definecolor{c_primary_dark}{HTML}{1e1e1e}      % body text, dark boxes; splash screen background (palette dark "bg")
\definecolor{c_primary_highlight}{HTML}{0066cc} % accent: bullets, highlights; SlideMage button blue (Create PDF, progress bar, links)
\definecolor{c_primary_gray}{HTML}{6f6f6f}      % highlightframe background; palette dark "disabled_fg", dark enough for white text
\definecolor{c_primary_white}{HTML}{FFFFFF}
\definecolor{c_secondary_violet}{HTML}{ba68c8}  % block titles (block); palette dark "tab_deck_editor"
\definecolor{c_tertiary_violet}{HTML}{dcb4e4}   % block bodies; derived (50% white-blend)
% Derived, not an app color: the button blue is too dark for the dark text
% an infoblock title puts on it, the app's dark-mode link blue (4da6ff) too
% light for a link underline on white; this one sits between them and passes
% both (4.9:1 and 3.4:1).
\definecolor{c_secondary_blue}{HTML}{3d8fe0}    % infoblock titles, link underline
\definecolor{c_tertiary_blue}{HTML}{9ec7f0}     % derived (50% white-blend)
\definecolor{c_secondary_green}{HTML}{4caf50}   % exampleblock titles; palette dark "status_pass"
\definecolor{c_tertiary_green}{HTML}{81c784}    % palette dark "tab_slide_editor"
\definecolor{c_secondary_red}{HTML}{f04848}     % alertblock titles; palette dark "status_fail"
\definecolor{c_tertiary_red}{HTML}{f8a4a4}      % derived (50% white-blend)
\definecolor{c_secondary_orange}{HTML}{ffb74d}  % reference block titles; palette dark "tab_archived"
\definecolor{c_tertiary_orange}{HTML}{ffdba6}   % derived (50% white-blend)
\definecolor{c_secondary_yellow}{HTML}{ffd500}  % logo/splash yellow
\definecolor{c_tertiary_yellow}{HTML}{ffe066}   % deck editor's highlight yellow
\definecolor{c_tertiary_gray}{HTML}{e4e4e4}     % block bodies, process boxes; palette light "panel_bg"
% Supporting colors ("custom": not part of the palette proper).
% c_custom_gray is hyperlink text, button fill and the quote attribution.
\definecolor{c_custom_gray}{HTML}{2d2d2d}       % palette dark "panel_bg"
\definecolor{c_custom_red}{HTML}{cc0000}        % palette light "status_fail"
\definecolor{c_custom_green}{HTML}{1a8a1a}      % palette light "status_pass"
\definecolor{c_custom_brown}{HTML}{996600}      % palette light "syntax_brace"

% ---------------------------------------------------------------------------
% Roles: which palette color a part of the slide uses (color names, read
% when used, so they follow a redefined palette color)
% ---------------------------------------------------------------------------
\newcommand{\SMBulletColor}{c_primary_highlight}       % itemize bullets
\newcommand{\SMTitleBoxColor}{c_primary_dark}          % frame title bar
\newcommand{\SMTitleTextColor}{c_primary_white}        % frame title text
\newcommand{\SMFooterBoxColorDefault}{c_primary_dark}  % footer box
\newcommand{\SMFooterTextColor}{c_primary_white}       % footer text, slide id

% Optional colors a template may define; SlideMage_base derives each one from
% the palette when the template leaves it undefined:
%   SMCanvasColor          slide background                  (white)
%   SMCardBackgroundColor  SMCard fill                       (white)
%   SMCardBorderColor      SMCard border                     (c_primary_dark)
%   SMProcessBorderColor   SMProcess border, dark mode       (c_tertiary_gray)
%   SMCodeKeywordColor     \SMCodeFile commands              (c_primary_highlight)
%   SMCodeCommentColor     \SMCodeFile comments              (c_primary_gray)
%   SMCodeBackgroundColor  \SMCodeFile panel                 (c_tertiary_gray)
%   SMCodeTextColor        \SMCodeFile text                  (c_primary_dark)

% ---------------------------------------------------------------------------
% Logo
% ---------------------------------------------------------------------------
% A vector PDF cropped to its ink, found on the TeX search path (a project
% path such as schemes/<name>/logo.pdf, or a file next to the template).
\newcommand{\SMInstituteLogoPath}{slidemage-logo.pdf}
\newlength{\SMLogoHeight}
\setlength{\SMLogoHeight}{3ex}                         % logo height in the title bar

% ---------------------------------------------------------------------------
% Frame title bar
% ---------------------------------------------------------------------------
% Space around the frame title (the theme's registers under their SlideMage
% names). The right padding is measured against the compiled PDF for the logo:
% see SlideMage_base.tex's comment on \gothamFrametitleRightpading.
\let\SMFrametitleTopPadding\gothamFrametitleToppading
\let\SMFrametitleBottomPadding\gothamFrametitleBottompading
\let\SMFrametitleLeftPadding\gothamFrametitleLeftpading
\let\SMFrametitleRightPadding\gothamFrametitleRightpading
\setlength{\SMFrametitleRightPadding}{0.28cm}
}{%
		Add \string% SPDX-FileCopyrightText: 2026 Stefan Palan, Jürgen Fleiß
% SPDX-License-Identifier: LPPL-1.3c+
% SlideMage setup layer and hook catalog (SlideMage v9.0.0).
%
% A template \input{}s this file first, sets whatever it wants to change from
% the catalog below, then \input{}s SlideMage_base:
%
%   % SlideMage template name: My University
%   \input{SlideMage_setup}
%   \definecolor{c_primary_highlight}{HTML}{...}
%   \renewcommand{\SMInstituteLogoPath}{schemes/my_logo.pdf}
%   \setlength{\SMLogoHeight}{3ex}
%   \input{SlideMage_base}
%   ... additions only: dark-mode beamer colors, citation format, project macros
%
% Rules for a template:
%   1. Before this file: only \SMDarkMode (written into the deck by SlideMage).
%   2. Between this file and SlideMage_base: \setlength, \renewcommand,
%      \definecolor and \colorlet on the names in this catalog. Every name
%      exists already, so a misspelt \renewcommand is an error instead of a
%      setting that silently does nothing.
%   3. After SlideMage_base: additions only.
%
% Decks never load the beamer theme themselves; this file does. A length given
% in ex or em is evaluated where the template sets it, in the normal body font,
% after the theme has loaded its fonts.

\newcommand{\SMSetupLoaded}{}

% The beamer theme, with the body fonts.
\usetheme{gotham}

% ---------------------------------------------------------------------------
% Palette
% ---------------------------------------------------------------------------
% The SlideMage look; a template redefines what it changes (usually all of
% them). The framework combines these pairs, so a template should keep them
% readable (WCAG 2.1): body text c_primary_dark on each c_secondary_* block
% title color >= 4.5:1, title/footer text on its box color >= 4.5:1, white text
% on c_primary_gray (highlightframe), link underline c_secondary_blue on the
% canvas >= 3:1.
%
% SlideMage's own values: wherever possible a color the app itself uses
% (src/slidemage/palette.py, qt_deck_table.py, assets/splash.png). White
% title/footer text on c_primary_dark is 16.7:1, white on c_primary_gray 5.0:1.
\definecolor{c_primary_dark}{HTML}{1e1e1e}      % body text, dark boxes; splash screen background (palette dark "bg")
\definecolor{c_primary_highlight}{HTML}{0066cc} % accent: bullets, highlights; SlideMage button blue (Create PDF, progress bar, links)
\definecolor{c_primary_gray}{HTML}{6f6f6f}      % highlightframe background; palette dark "disabled_fg", dark enough for white text
\definecolor{c_primary_white}{HTML}{FFFFFF}
\definecolor{c_secondary_violet}{HTML}{ba68c8}  % block titles (block); palette dark "tab_deck_editor"
\definecolor{c_tertiary_violet}{HTML}{dcb4e4}   % block bodies; derived (50% white-blend)
% Derived, not an app color: the button blue is too dark for the dark text
% an infoblock title puts on it, the app's dark-mode link blue (4da6ff) too
% light for a link underline on white; this one sits between them and passes
% both (4.9:1 and 3.4:1).
\definecolor{c_secondary_blue}{HTML}{3d8fe0}    % infoblock titles, link underline
\definecolor{c_tertiary_blue}{HTML}{9ec7f0}     % derived (50% white-blend)
\definecolor{c_secondary_green}{HTML}{4caf50}   % exampleblock titles; palette dark "status_pass"
\definecolor{c_tertiary_green}{HTML}{81c784}    % palette dark "tab_slide_editor"
\definecolor{c_secondary_red}{HTML}{f04848}     % alertblock titles; palette dark "status_fail"
\definecolor{c_tertiary_red}{HTML}{f8a4a4}      % derived (50% white-blend)
\definecolor{c_secondary_orange}{HTML}{ffb74d}  % reference block titles; palette dark "tab_archived"
\definecolor{c_tertiary_orange}{HTML}{ffdba6}   % derived (50% white-blend)
\definecolor{c_secondary_yellow}{HTML}{ffd500}  % logo/splash yellow
\definecolor{c_tertiary_yellow}{HTML}{ffe066}   % deck editor's highlight yellow
\definecolor{c_tertiary_gray}{HTML}{e4e4e4}     % block bodies, process boxes; palette light "panel_bg"
% Supporting colors ("custom": not part of the palette proper).
% c_custom_gray is hyperlink text, button fill and the quote attribution.
\definecolor{c_custom_gray}{HTML}{2d2d2d}       % palette dark "panel_bg"
\definecolor{c_custom_red}{HTML}{cc0000}        % palette light "status_fail"
\definecolor{c_custom_green}{HTML}{1a8a1a}      % palette light "status_pass"
\definecolor{c_custom_brown}{HTML}{996600}      % palette light "syntax_brace"

% ---------------------------------------------------------------------------
% Roles: which palette color a part of the slide uses (color names, read
% when used, so they follow a redefined palette color)
% ---------------------------------------------------------------------------
\newcommand{\SMBulletColor}{c_primary_highlight}       % itemize bullets
\newcommand{\SMTitleBoxColor}{c_primary_dark}          % frame title bar
\newcommand{\SMTitleTextColor}{c_primary_white}        % frame title text
\newcommand{\SMFooterBoxColorDefault}{c_primary_dark}  % footer box
\newcommand{\SMFooterTextColor}{c_primary_white}       % footer text, slide id

% Optional colors a template may define; SlideMage_base derives each one from
% the palette when the template leaves it undefined:
%   SMCanvasColor          slide background                  (white)
%   SMCardBackgroundColor  SMCard fill                       (white)
%   SMCardBorderColor      SMCard border                     (c_primary_dark)
%   SMProcessBorderColor   SMProcess border, dark mode       (c_tertiary_gray)
%   SMCodeKeywordColor     \SMCodeFile commands              (c_primary_highlight)
%   SMCodeCommentColor     \SMCodeFile comments              (c_primary_gray)
%   SMCodeBackgroundColor  \SMCodeFile panel                 (c_tertiary_gray)
%   SMCodeTextColor        \SMCodeFile text                  (c_primary_dark)

% ---------------------------------------------------------------------------
% Logo
% ---------------------------------------------------------------------------
% A vector PDF cropped to its ink, found on the TeX search path (a project
% path such as schemes/<name>/logo.pdf, or a file next to the template).
\newcommand{\SMInstituteLogoPath}{slidemage-logo.pdf}
\newlength{\SMLogoHeight}
\setlength{\SMLogoHeight}{3ex}                         % logo height in the title bar

% ---------------------------------------------------------------------------
% Frame title bar
% ---------------------------------------------------------------------------
% Space around the frame title (the theme's registers under their SlideMage
% names). The right padding is measured against the compiled PDF for the logo:
% see SlideMage_base.tex's comment on \gothamFrametitleRightpading.
\let\SMFrametitleTopPadding\gothamFrametitleToppading
\let\SMFrametitleBottomPadding\gothamFrametitleBottompading
\let\SMFrametitleLeftPadding\gothamFrametitleLeftpading
\let\SMFrametitleRightPadding\gothamFrametitleRightpading
\setlength{\SMFrametitleRightPadding}{0.28cm}
 at the top of the template, before its colors,
		or run Info -> Project compatibility in SlideMage.}
\fi
% \SMLuaOnce{<lua>}: run Lua now AND remember it, so a precompiled format
% (SlideMage's fast preview, see preview_format.py) can replay it. A LuaTeX
% format stores TeX state but not Lua state: every global function, table and
% callback a preamble sets up is gone when the format is loaded. The replay list
% is run by the format's \everyjob; in an ordinary build it is never used.
\newtoks\SMLuaReplayList
\newcommand{\SMLuaOnce}[1]{%
	\directlua{#1}%
	\global\SMLuaReplayList\expandafter{\the\SMLuaReplayList\directlua{#1}}%
}
% The packages this file uses, loaded here explicitly (the theme loads most of them already, which makes
% each \usepackage below a no-op): iftex for \ifluatex, fontspec for \newfontfamily, xcolor for the colors.
\usepackage{iftex}
\usepackage{fontspec}
\usepackage{xcolor}
\usepackage{etoolbox} % \AtEndEnvironment, used for the itemize level-drop spacing below
\usepackage{tikz} % For painting a white box over title bar
\usetikzlibrary{calc}
\usetikzlibrary{arrows.meta} % for e.g. arr/.style={-{Stealth[...]}...} on diagram slides
\usetikzlibrary{positioning} % for `right=<dist> of <node>`-style relative node placement
\usetikzlibrary{decorations.text} % text-along-a-path, for a circular photo's curved caption
\usetikzlibrary{decorations.pathreplacing} % the `brace` decoration shape, for \SMUnderbrace below
% TikZ reveal steps that keep the picture's size (2026-09-30): `visible on=<3->`
% draws an element on steps 3+ and draws it invisibly (opacity 0) before, so it
% still counts toward the bounding box -- unlike \only<3->{\draw ...;}, whose
% element is simply missing on steps 1-2 and can make the whole picture jump
% when it appears (the "diagram shifts downward in step 7" class of fix).
% Same key names as the aobs-tikz library (overlay-beamer-styles), defined here
% so no extra package is needed.
\tikzset{
	invisible/.style={opacity=0, text opacity=0},
	alt/.code args={<#1>#2#3}{\alt<#1>{\pgfkeysalso{#2}}{\pgfkeysalso{#3}}},
	visible on/.style={alt={#1{}{invisible}}},
}
\usepackage{booktabs} % \toprule/\midrule/\bottomrule for SMTable
\usepackage{colortbl} % \arrayrulecolor{}, used by SMTable to color those rules
\usepackage{pifont} % \ding{51}/\ding{55} checkmark/X-mark glyphs, e.g. slides/fr3-Forms-of-Capital-Market-Efficiency-Strong-Semistrong-Weak-current.tex's own validity ratings
\usepackage{pgfplots} % real-data line/bar charts (e.g. slides/3vp-Metallgesellschaft-Oil-Price-Chart-current.tex) -- confirmed safe alongside this deck's own \tableofcontents, see CONVERSION_NOTES.md's "beamer/TeX gotchas" entry
\pgfplotsset{compat=1.18}

% The math fonts (Computer Modern next to the sans text font) and the \text{} wrapper that
% goes with them are font handling: SlideMage_layout_\SMLayoutProfile.tex, input by SlideMage_setup.

% Colors, institute logo path, and the two logo-calibration lengths
% (\SMLogoHeight, \SMFrametitleRightPadding = \SMFrametitleRightPadding)
% are template settings: SlideMage_setup.tex defines each with its default
% (the hook catalog), the template changes what it wants, and this file only
% reads them -- shared machinery, no institution-specific branding. See e.g.
% SlideMage_template_uni-graz.tex for the Graz values. \SMLogoWidth is the one
% logo-derived length that *does* stay computed here (via \settowidth,
% further below) rather than hardcoded per institution, since it already
% adapts automatically to whatever image \SMInstituteLogoPath points at.

% Style hyperlinks (\href{url}{text}) consistently across all slides:
% underlined, so a link is visually recognizable as such against plain
% body text -- yellow (this project's usual accent color) was tried first
% for the *text* itself, but yellow-on-white body text is hard to read,
% unlike yellow used as an accent against dark/black text elsewhere in the
% theme. Went through a few more candidates on 2026-07-06 (all direct
% instruction) before settling on c_secondary_blue for both text
% and underline (a real color reads more clearly as "this is a link" than
% any gray does, without yellow's legibility problem).
%
% 2026-07-07, direct instruction: split text color and underline color
% instead of using one color for both -- link text back to
% c_custom_gray (the near-black dark gray also used for
% quoteframe's attribution text), underline in a separate accent color.
% Plain \underline can't do this: its rule is drawn in whatever color is
% current at that point, so a single \underline{...} call always
% underlines in its own text's color, never a second, different one.
% \SMHrefUnderline below reimplements it instead: measure the linked
% text's own width into a savebox, print that box in customgray, then
% back up by the box's own width and draw a plain \rule of the same
% width, raised to sit just below the text's baseline, in the accent
% color. Same single-line limitation as plain \underline (and the reason
% \url{} is used instead of \href on slides/acw-Some-Topics-in-Testing-the-EMH-current.tex --
% see that slide's own comment): a \rule, like \underline, can't span a
% line break either, so this only ever wraps normal, comparatively short
% \href display text, never something that needs mid-word breaking.
%
% First tried c_primary_highlight for the underline (same day) -- confirmed
% via the compiled PDF that the rule *was* being drawn (right color,
% position, width), so not a compile-time bug, but yellow (very high
% luminance, close to white) at this rule's 0.4pt thickness reads as
% essentially invisible against the white slide background, worst at
% footnote size (see e.g. slides/za2-Startup-Financing-Categorization-current.tex's
% own long inline URL) where it's easy to miss entirely. Switched
% 2026-07-07 (direct instruction) to c_secondary_blue instead --
% far more contrast against white at the same thickness -- and the rule
% was also thickened at the same time (0.4pt -> 0.6pt), on the reasoning
% that it might need it for the same contrast reason. It didn't: blue
% already reads clearly at the original thickness, and 0.6pt instead
% just looked heavier than intended (direct feedback 2026-07-08, e.g. on
% slides/6qi-Yield-Spreads-Corporate-Bonds-current.tex). Reverted to
% 0.4pt -- only the color was ever actually asked to change.
%
% `citecolor` (2026-07-23, direct instruction: "make sure the year in the
% short citation is printed in black", found on
% slides/64m-APT-Choosing-the-Factors-Chen-Roll-Ross-current.tex's own
% `\textcite{}` call, "and generally" -- a deck-wide fix, not a one-off):
% with `colorlinks=true` and no explicit `citecolor`, hyperref falls back
% to its own factory default (green) for every citation hyperlink --
% confirmed directly, rendering 64m with the real biblatex/biber setup
% (this deck's own `\hypersetup` alone, with no citecolor, has no other
% source that could be coloring it): `\textcite{chen1986}` printed as
% "Chen, Roll, Ross (\textcolor{green}{1986})", the year in green, nothing
% else on the line affected (`\textcite` only wraps the year/label itself
% in the actual hyperlink, not the surrounding author names, which are
% plain text). `citecolor=black` set explicitly here, alongside
% urlcolor/linkcolor, rather than left to hyperref's own default, so this
% applies deck-wide to every citation command (`\cite`, `\textcite`,
% `\parencite`, ...) that ends up as a colored link, not just 64m's own
% call site.
\hypersetup{colorlinks=true, urlcolor=c_custom_gray, linkcolor=c_custom_gray, citecolor=black}


\let\SMOrigHref\href
\newsavebox{\SMHrefTextBox}
%
% 2026-07-08 (direct instruction): the underline's own rule thickness now
% scales with the ambient font size instead of a single fixed 0.2pt for
% every link everywhere -- 0.2pt at this deck's footnote size, 0.4pt at its
% normal body-text size, interpolated (and extrapolated, for any other
% size a link happens to appear at) linearly in between.
%
% Deriving "current font size" from \f@size (the textbook NFSS way) was
% tried first and confirmed broken specifically *inside* a beamer frame --
% \typeout{\meaning\f@size} right after this preamble loads shows a normal,
% defined macro (`macro:->10.95`), but the exact same check inside
% `\begin{standardframe}...\end{standardframe}` reports `\f@size` as
% *undefined*, presumably from beamer's own overlay-aware font-switching
% machinery not keeping classic NFSS's own current-size bookkeeping in
% sync the way plain (non-beamer) LaTeX does. Since every real use of this
% macro happens inside a frame, that path was a dead end regardless of how
% well it worked in isolation.
%
% Measuring a reference glyph's own rendered height instead sidesteps this
% entirely: building a box and reading back its own height just measures
% whatever font is *actually* active at the call site, independent of
% whether beamer's own internal "current size" bookkeeping happens to
% agree -- and within one font
% family, every glyph's metrics scale by exactly the same ratio the point
% size does, so a measured height is just as valid a stand-in for "current
% size" as \f@size would have been, had it worked. "M" (present in every
% context this is used in) is the reference glyph; \SMHrefUnderlineFnHeight/
% \SMHrefUnderlineBodyHeight are its height at this deck's footnote font
% and normal body font respectively, each measured once, here, in the
% relevant font context.
% Deliberately built from raw TeX primitives (\newdimen/\setbox/\ht),
% *not* \newlength/\setlength/\settoheight -- found, while building this,
% that loading this deck's theme silently turns \setlength and
% \settoheight into something else entirely: `\meaning\setlength` reads
% `macro:->\protect \setlength` right after `\usetheme{gotham}` (still the
% ordinary kernel definition beforehand), and the same for `\settoheight`.
% Not chased to which file does this or why (bisected only as far as
% "plain beamer alone doesn't, beamer+gotham does" -- most likely
% `translator`, which beamer loads for localized strings and is known to
% robustify commands this way, but not confirmed) -- what matters
% practically is that a *global* assignment through either one silently
% no-ops instead of erroring, confirmed directly and in isolation (a
% minimal beamer+gotham file, no other content): `\global\setlength{\len}
% {\otherlen}` leaves `\len` at 0pt after the group closes, while the
% exact same pattern with plain `\newdimen`/raw `=` assignment propagates
% correctly. Every plain (non-global) `\setlength` elsewhere in this
% project's own files is unaffected -- it's specifically the *global*
% form that breaks, which is why this is the first place it's shown up
% despite `\setlength` being used constantly throughout this preamble.
\newdimen\SMHrefUnderlineFnHeight
\newdimen\SMHrefUnderlineBodyHeight
\newdimen\SMHrefUnderlineNowHeight
\newdimen\SMHrefUnderlineThickness
% A dedicated box register, not the commonly-borrowed box0 -- several
% slides in this deck (the "Long-term bonds" tables) already use box0
% themselves (`\sbox0{...}`/`\usebox0`) for their own, unrelated
% mid-table-layout purposes, and a link can appear anywhere, including
% inside one of those tables -- reusing box0 here would silently clobber
% whatever they were holding there.
\newbox\SMHrefUnderlineBox
% The two reference heights are (re)measured on every single \href call,
% not once at preamble-load time -- tried once-at-load first and found a
% second, independent problem with it: `\usebeamerfont{footnote}` right
% here (before `\begin{document}`) silently has no effect at all (footnote-
% size "M" and normal-size "M" measured identically, confirmed directly),
% presumably because beamer's font-theme machinery isn't fully installed
% until the document starts. Since every real `\href` call happens well
% inside the document (inside some frame), measuring fresh each time
% sidesteps that timing question entirely rather than needing to find the
% right hook to measure once, early enough to be free but late enough to
% be correct -- the extra cost is two throwaway one-glyph boxes per link,
% not something worth optimizing away in a slide deck.
\newcommand{\SMHrefUnderlineMeasureRefHeights}{%
	{\usebeamerfont{footnote}\setbox\SMHrefUnderlineBox=\hbox{M}\global\SMHrefUnderlineFnHeight=\ht\SMHrefUnderlineBox}%
	{\normalsize\setbox\SMHrefUnderlineBox=\hbox{M}\global\SMHrefUnderlineBodyHeight=\ht\SMHrefUnderlineBox}%
}
% A plain e-TeX \dimexpr can multiply a dimension by another dimension and
% then divide by a third, as one combined operation, as long as the
% *middle* operand of the chain is itself wrapped in its own nested
% \dimexpr...\relax -- confirmed directly (an expl3 \dim_set:Nn attempt
% first, using \seq_.../\fp_... the way \setfootnotesymbols above does,
% kept hitting "Missing ) inserted for expression"/type-mismatch errors:
% dividing one dimension by another is a *ratio*, not a dimension, and
% \dim_set:Nn's own expression grammar doesn't implicitly convert between
% the two the way this nested-\dimexpr form does). No expl3 needed here at
% all -- plain e-TeX arithmetic, already available in any current LaTeX.
\newcommand{\SMHrefUnderlineSetThickness}{%
	\SMHrefUnderlineThickness=\dimexpr 0.2pt +
		(\SMHrefUnderlineNowHeight - \SMHrefUnderlineFnHeight)
		* \dimexpr0.2pt\relax
		/ (\SMHrefUnderlineBodyHeight - \SMHrefUnderlineFnHeight)\relax
}
\newcommand{\SMHrefUnderline}[1]{%
	\sbox{\SMHrefTextBox}{#1}%
	\textcolor{c_custom_gray}{\usebox{\SMHrefTextBox}}%
	\hspace{-\wd\SMHrefTextBox}%
	\SMHrefUnderlineMeasureRefHeights
	\setbox\SMHrefUnderlineBox=\hbox{M}\SMHrefUnderlineNowHeight=\ht\SMHrefUnderlineBox
	\SMHrefUnderlineSetThickness
	\raisebox{-0.35ex}[0pt][0pt]{\textcolor{c_secondary_blue}{\rule{\wd\SMHrefTextBox}{\SMHrefUnderlineThickness}}}%
}
\renewcommand{\href}[2]{\SMOrigHref{#1}{\SMHrefUnderline{#2}}}
% \url{} (the `url` package hyperref depends on/patches) defaults to a
% monospace-ish "tt" style regardless of the above -- used deliberately in
% only one place in this deck,
% slides/acw-Some-Topics-in-Testing-the-EMH-current.tex, specifically
% *for* its line-breaking-at-slashes behavior (a plain \href with the
% URL as its own display text has no break points at all -- see that
% slide's own comment). Every other same-text-as-target URL in the deck
% (including slides/bb6-A-Perpetual-Bond-current.tex's own video link,
% which used to be a \url{} too) uses plain \href{url}{url} instead, since
% \url{} doesn't pick up this deck's custom colored underline below (its
% breakable content can't be measured into the savebox the underline
% needs -- see that mechanism's own comment) and most of these URLs are
% short enough to not need the breaking \url{} exists for anyway. 2026-07-07,
% direct instruction: that one genuinely-needed \url{} (on s-acw) should
% still look like this deck's other hyperlinks, i.e. the surrounding body
% font, not \url{}'s own default monospace.
% `\urlstyle{same}` (a standard url-package switch) does exactly that,
% with no effect on \url{}'s own breaking behavior. Set deck-wide, not
% only on that one slide, so any future \url{} use matches by default
% too, rather than needing to remember this each time. Coloring already
% matches automatically (hyperref's own \url respects
% \hypersetup{urlcolor=...} the same way \href does); the underline
% \href gets above doesn't carry over, since \underline's own argument
% can't contain line breaks -- wrapping \url{}'s necessarily-breakable
% text in \underline would silently undo the breaking this exists for,
% so that one difference is accepted rather than "fixed."
\urlstyle{same}



% The colour roles of the theme follow the template's palette: SMColorDark and SMColorA, the canvas
% default and the roles (beamerthemeSlideMage.sty). Needs to go after the template's color redefinitions.
\SMApplyColorRoles
% The hooks that go into registers of beamer's and LaTeX's own (text margins, footnote font
% and spacing): needs to go after the template as well.
\SMApplyLayoutHooks

% Derived from the template's hooks and the font (v9.7.0; they were hand-kept numbers of the catalog
% and the layout profile before). The setup layer holds a macro under each of these four names that
% stops with an error when a template or a profile sets one; a name still holding that macro is
% replaced by the length it stands for, and a name that was redefined is reported.
\makeatletter
\def\SM@derivedlength#1{%
	\expandafter\ifx\csname #1\expandafter\endcsname\csname SM@guard@#1\endcsname\else
		\PackageError{SlideMage}{\expandafter\string\csname #1\endcsname\space is derived, not set}{%
			A template or a layout profile redefined it. SlideMage_base derives it from the hooks and
			the font; remove that line from the template or the profile.}%
	\fi
	\expandafter\let\csname #1\endcsname\relax
	\expandafter\newlength\csname #1\endcsname}
\SM@derivedlength{SMContentTopAlign}
\SM@derivedlength{SMTitleBarBottomAlign}
\SM@derivedlength{SMFootnoteRightShift}
\SM@derivedlength{SMFooterRightShift}
% v9.8.0: the lengths below were numbers of the layout profile (measured once at Fira) or of the
% catalog; each is evaluated in the font where it is used. The footer's, the lists' and the section
% frame's follow with the code that uses them.
\SM@derivedlength{SMFooterTextDrop}
\SM@derivedlength{SMFootlineVOffset}
\SM@derivedlength{SMLabelItemWidth}
\SM@derivedlength{SMLabelItemSep}
\SM@derivedlength{SMEnumerateNestedMinNumber}
\SM@derivedlength{SMEnumerateNestedLabelSep}
\SM@derivedlength{SMEnumerateNestedLeftMargin}
\SM@derivedlength{SMEnumerateLabelBox}
\SM@derivedlength{SMSectionframeRefDepth}
\SM@derivedlength{SMSectionframeBottomGap}
\SM@derivedlength{SMSectionframeBottomPadCorrection}
\SM@derivedlength{SMSectionframeHeightCorrection}
% v9.9.0: the pull-ups of the frames the automatic Top Margin does not correct by their content were
% the last three numbers of the layout profiles. \SMContentTopShift is measured (the frame's origin,
% see AUTOMATIC TOP MARGIN; in a run without \jobname.top, from the first frame that counts on), the two columns
% corrections are measured where a columns environment starts such a frame.
\SM@derivedlength{SMContentTopShift}
\SM@derivedlength{SMColumnsTopCorrection}
\SM@derivedlength{SMColumnsTTopCompensation}
\makeatother
% - \SMContentTopAlign, the Top Margin as an overlay coordinate: the same line the automatic Top Margin
%   puts the first text line's box on (the Top Margin is in bp, the coordinate in pt: 41.58bp is 41.7359pt).
\setlength{\SMContentTopAlign}{\SMTopMargin}
% - \SMTitleBarBottomAlign, the visible bottom of the title bar: where the logo ends. The frame title
%   template centres the logo on the capital height of the title font (\raisebox{.5\fontcharht\font`E-0.5\height}
%   at the baseline, which is the top padding plus the height of an X below the bar's top), so the
%   bottom is the top padding, plus that X, minus half an E, plus half the logo.
\begingroup\usebeamerfont{frametitle}\settototalheight{\dimen0}{X}\dimen1=\fontcharht\font`E
\global\SMTitleBarBottomAlign=\dimexpr\SMFrametitleTopPadding+\dimen0-0.5\dimen1+0.5\SMLogoHeight\relax\endgroup
% - \SMFootnoteRightShift: a footnote ends where the text margin does; the logo ends its right padding
%   from the page's edge, so the footnotes are pushed on by the difference.
\setlength{\SMFootnoteRightShift}{\dimexpr\SMTextMarginRight-\SMFrametitleRightPadding\relax}
% (\SMFooterRightShift is derived with the footer, below.)
% - \SMFootlineVOffset (v9.8.0; the catalog's -0.5ex-0.77pt until then), the space the footline
%   template puts before the footline box. Beamer sets the frame text area to end where the footline
%   box begins (its height plus depth, and the offset, are its \footheight) plus a fixed 4pt
%   (beamerbaseframecomponents.sty, \beamer@calculateheadfoot), and the last footnote line's baseline
%   is the bottom of that area. So the baseline lies
%       \SMFootlineHeight + \SMFootlineDepth + offset + 4pt
%   above the page's bottom edge, and the offset that puts the descenders of that line (the depth of
%   "gjpqy" in the footnote font) \SMFootnoteFooterGap above the footer box follows.
\makeatletter
\begingroup\usebeamerfont{footnote}\settodepth{\dimen0}{gjpqy}
\global\SMFootlineVOffset=\dimexpr\SMFooterBoxHeight+\SMFootnoteFooterGap+\dimen0-4pt-\SMFootlineHeight-\SMFootlineDepth\relax
\xdef\SMPDFInfoFootnoteDescender{\strip@pt\dimexpr\dimen0*7200/7227\relax}\endgroup
\makeatother

% The effective margins, for tools/check_layout.py (SlideMage v9.5.0): custom keys of the PDF's
% info dictionary, in bp (PDF units, what the tool measures in). SMTopMargin is the Top Margin
% (\SMTopMargin), SMLeftMargin the frame title text's left edge (\SMFrametitleLeftPadding). A template
% has set its hooks by now; a slide cannot change them. SMTitleBarBottom (v9.6.0) is the bottom edge of
% the title bar (\SMTitleBarBottomAlign): the tool tells a frame with the ordinary title bar by it.
% SMFontAscender, the text font's ascender, follows once it is known (below, with the Lua file).
% v9.8.0: SMFootnoteFooterGap, SMFootnoteDescender (the depth of "gjpqy" in the footnote font) and
% SMFooterBoxHeight, which tools/calibrate_font_profile.py checks the footnotes' position against.
\ExplSyntaxOn
\edef\SMPDFInfoTopMargin{\dim_to_decimal_in_bp:n{\SMTopMargin}}
\edef\SMPDFInfoLeftMargin{\dim_to_decimal_in_bp:n{\SMFrametitleLeftPadding}}
\edef\SMPDFInfoTitleBarBottom{\dim_to_decimal_in_bp:n{\SMTitleBarBottomAlign}}
\edef\SMPDFInfoFootnoteGap{\dim_to_decimal_in_bp:n{\SMFootnoteFooterGap}}
\edef\SMPDFInfoFooterBoxHeight{\dim_to_decimal_in_bp:n{\SMFooterBoxHeight}}
\ExplSyntaxOff
\hypersetup{pdfinfo={SMTopMargin={\SMPDFInfoTopMargin}, SMLeftMargin={\SMPDFInfoLeftMargin},
	SMTitleBarBottom={\SMPDFInfoTitleBarBottom},
	SMFootnoteFooterGap={\SMPDFInfoFootnoteGap}, SMFootnoteDescender={\SMPDFInfoFootnoteDescender},
	SMFooterBoxHeight={\SMPDFInfoFooterBoxHeight}}}

% Shared diagram roles: templates set the colors once; slides use these styles.
\makeatletter
\@ifundefinedcolor{SMCardBackgroundColor}{\colorlet{SMCardBackgroundColor}{white}}{}
\@ifundefinedcolor{SMCardBorderColor}{\colorlet{SMCardBorderColor}{c_primary_dark}}{}
\@ifundefinedcolor{SMProcessBorderColor}{\colorlet{SMProcessBorderColor}{c_tertiary_gray}}{}
\makeatother
\tikzset{
	SMCard/.style={draw=SMCardBorderColor, fill=SMCardBackgroundColor, text=c_primary_dark},
	SMProcess/.style={draw=none, fill=c_tertiary_gray, text=c_primary_dark},
}

% The theme's divider frames for \part, \section, \subsection and \subsubsection
% (\SMDividerFrames in beamerthemeSlideMage.sty) are off, so \section{} (used
% for a deck's chapter grouping/table of contents) does not insert its own
% plain, unstyled divider frame plus a small "current section highlighted in
% the outline" frame -- redundant with (and visually inconsistent next to)
% this project's own hand-styled sectionframe (see README.md's "Special frame
% types" section), which every chapter in this deck already opens with as an
% explicit \slide{}. \section{} only feeds \tableofcontents (see
% slides/nkj-Agenda-current.tex).
%
% Same for \subsection{} (2026-07-15, direct instruction): \subsection{} calls
% still register in the PDF's own bookmark/outline structure (Acrobat's
% "chapters" pane -- a plain PDF sectioning feature, independent of whether
% beamer auto-inserts a frame for it) and still feed 	ableofcontents (see
% slides/nkj-Agenda-current.tex's own `	ableofcontents[hideallsubsections]`,
% which then hides them from that one slide specifically), without also
% spawning their own unstyled divider slide in the middle of a chapter every
% time one is added.


% A color leak from \titleframe/\quoteframe into whichever frame follows
% them is fixed in standardframe's own opening code below (search
% "Forces this frame's own body") rather than here -- an
% \AtBeginEnvironment{frame}{\color{c_primary_dark}} hook was tried first
% and confirmed, by isolated reproduction, to have *no effect at all*:
% beamer's own "frame" isn't a plain \newenvironment etoolbox's
% \AtBeginEnvironment can patch that way, so the hook silently never fired.


%\colorlet{colorPale}{gPaleYell} % BG in light/normal mode
%\colorlet{colorDark}{gDarkBlack} % FG in light/normal mode
%\colorlet{colorA}{gAnthracite} % frametitle, standin.out,
%\colorlet{colorAreversed}{gLightTeal} % frametitle, standin.in,
%\colorlet{colorB}{gMidGrey} % gray BG : progress bar, blocks
%\colorlet{colorC}{gDeepYellOr} % progress bar
%\colorlet{colorD}{gLightOrange} % alert
%\colorlet{colorE}{gLightGreen} % example

% Opt-in yellow border for an image placed via a tikz node (e.g.
% \node[SMImageBorder] {\includegraphics{...}};) -- matches the
% 3pt/accent1-yellow outline PowerPoint draws directly on a picture shape's
% own border (<a:ln w="38100"><a:solidFill><a:schemeClr val="accent1"/>,
% where accent1 turned out to already be this project's own
% c_primary_highlight, not a new color to map). A tikz *style* (a named
% bundle of node keys), not a wrapper command: adding or removing it is
% just adding or deleting one word from that one \node's own option list,
% no separate macro call/argument needed -- add `SMImageBorder` to draw
% the border, leave it off for a plain image, same as any other tikz key.
%
% Thickness is a single named length, \SMImageBorderWidth, not a bare
% "3pt" baked directly into the style (2026-07-15, direct instruction: "make
% sure the yellow image borders have consistent thickness throughout the
% project by defining a single variable") -- confirmed directly, by
% measuring the actual PDF content-stream "w" operator, that every
% SMImageBorder use already rendered at the exact same 3pt regardless
% (a genuinely shared style was already in effect, no numeric drift found),
% so this is a refactor for future adjustability/discoverability, not a fix
% for an actual inconsistency -- any perceived thickness difference between
% slides traces to something else instead (e.g. a border straddling flush
% image content on one slide and not another, which reads as "heavier" even
% at an identical stroke width -- see the SMImageBorder padding gotcha
% below).
% (\SMImageBorderWidth, 2.5pt, is a hook of SlideMage_setup.tex.)
% outer sep = half the stroke (2026-10-01, direct instruction: "if a figure
% has a border, that border should determine its position"): a bordered
% node's anchors (north east, west, ...) sit on the border's *outer* edge, so
% anchoring it at a margin puts the border on the margin. That is TikZ's own
% default (0.5\pgflinewidth, evaluated with the node's final line width, so a
% slide's line width=1.2pt override still gets 0.6pt); the style sets it
% explicitly so that an "outer sep=0pt" earlier in the option list
% (\cornerfigure's node, for an unbordered image) doesn't take it away. A
% slide that writes outer sep=0pt *after* the style opts out again -- ts8 did,
% and its border stuck out 1.25pt past the Right Margin.
\tikzset{SMImageBorder/.style={draw=c_primary_highlight, line width=\SMImageBorderWidth, outer sep=0.5\pgflinewidth}}

% Named styles for the drawing conventions slides repeat by hand
% (2026-09-30): a restyle -- "make the arrows a bit bolder", "all
% highlights dashed" -- is then one edit here instead of one per slide.
% Each is exactly the option list the slides used, so switching a slide to
% the name changes nothing on the page.
\tikzset{
	SMHighlight/.style={c_custom_red, very thick, rounded corners},   % red box around part of a figure/table (18x)
	SMDottedHighlight/.style={densely dotted, thick, c_custom_red},  % red dotted guide line (19x)
	SMArrow/.style={-{Stealth[length=2mm]}},                         % the usual arrow tip (54x)
}

% \SMTikzmark{name} -- a zero-size remember-picture node at the current
% text position (on the baseline, 2pt up), for arrows/braces that point at
% a word. Three slides defined their own identical copy (\tikzmark on 33r
% and ha9 -- two \newcommand's of one name, an error the moment both land in
% one deck -- and ebq's \ebqTikzmark).
\newcommand{\SMTikzmark}[1]{\tikz[remember picture, overlay, baseline] \node[inner sep=0pt, anchor=base, yshift=2pt] (#1) {};}

% \SMTikzword{name}{text} -- typesets text as usual, but as a named
% remember-picture node, so an overlay arrow can point at the word itself
% ((name.south), (name.north east), ...) and follows it when the layout
% moves (2026-10-01, direct instruction: 3ac's arrow at measured page
% coordinates had drifted off its words). Transcripts keep the text.
\newcommand{\SMTikzword}[2]{\tikz[remember picture, baseline=(#1.base)] \node[inner sep=0pt, outer sep=0pt] (#1) {#2};}

% \SMLineTop -- a zero-width strut as tall as a text line's *line box*
% (\SMFontAscender em above the baseline, the text font's ascender, 0.935 for Fira Sans: the box the Top Margin is
% measured on). Put at the start of the first line of a text column that sits
% next to a figure in top-aligned minipages,
%   \begin{minipage}[t]{0.4\linewidth}\vspace{0pt}
%   	\begin{itemize}\item \SMLineTop First item ...
%   \end{minipage}\hfill
%   \begin{minipage}[t]{0.56\linewidth}\vspace{0pt} <figure> \end{minipage}
% and the text's line box starts exactly where the figure starts -- so the
% automatic Top Margin puts both on it. Without it the first line's box is
% only as tall as its glyphs and the figure lands ~1.8pt below the margin
% (2026-10-01, m2k: its two minipages were vertically *centred*, LaTeX's
% default, and the text column started 12pt below the Top Margin;
% check_layout.py's `coltop` check now reports that pattern).
\newcommand{\SMLineTop}{\rule{0pt}{\SMFontAscender em}}

% \EUR -- the euro sign in Fira Sans, in text and in math alike. In math,
% \text{\texteuro} came out in the math font's style ("Euro symbols ...
% proper", 2026-07-25); slides wrote \text{\sffamily€} by hand (34x).
\newcommand{\EUR}{\ifmmode\text{\sffamily€}\else\texteuro\fi}

% \shrinklists{n} -- all list levels n font steps smaller, for a slide with
% more text than fits (put it inside the frame, like the \setbeamerfont
% lines it replaces; the frame's group resets it). The beamer defaults are
% normal/\small/\footnotesize for levels 1-3; n=1 gives
% \small/\footnotesize/\scriptsize, n=2 \footnotesize/\scriptsize/\scriptsize,
% n=3 \scriptsize on all three levels (the floor; v5.30.0 -- before, every
% n>=2 meant n=2). 13 slides wrote 2-3 \setbeamerfont lines for this, in five
% different combinations.
% The item gaps shrink with the font since 2026-10-01 -- see \SMNoListParsep
% next to the itemize code.
\newcommand{\shrinklists}[1]{%
	\ifcase#1\relax
	\or
		\setbeamerfont{itemize/enumerate body}{size=\small}%
		\setbeamerfont{itemize/enumerate subbody}{size=\footnotesize}%
		\setbeamerfont{itemize/enumerate subsubbody}{size=\scriptsize}%
	\or
		\setbeamerfont{itemize/enumerate body}{size=\footnotesize}%
		\setbeamerfont{itemize/enumerate subbody}{size=\scriptsize}%
		\setbeamerfont{itemize/enumerate subsubbody}{size=\scriptsize}%
	\else
		\setbeamerfont{itemize/enumerate body}{size=\scriptsize}%
		\setbeamerfont{itemize/enumerate subbody}{size=\scriptsize}%
		\setbeamerfont{itemize/enumerate subsubbody}{size=\scriptsize}%
	\fi
}

% The "padding gotcha" promised above, finally root-caused (2026-07-23,
% background-task follow-up to the IA2 fine-tuning batch #2 anomaly on
% slides/m2k-Exercise-CAPM-2-current.tex/
% slides/n8e-Tail-Risk-Variance-Clustering-current.tex/
% slides/fm6-Tail-Risk-Market-vs-Simulated-current.tex/
% slides/xqq-Fat-Tails-current.tex, all fixed there with a pragmatic
% `line width=1.2pt` override and an explicitly *unresolved* root cause at
% the time -- "declares an identical stroke width and CTM in the compiled
% PDF's own raw content stream, yet renders at a genuinely different pixel
% thickness", with CTM scale/clips/dash/cap/join all directly checked and
% ruled out). Confirmed via a controlled reproducer (four `SMImageBorder`
% nodes -- plain opaque image, nested pgfplots axis, plain image at
% `inner sep=3pt`, nested plain tikzpicture at `inner sep=3pt` -- rendered
% at 600dpi and measured by a half-max cross-section profile, not just
% eyeballed): the one variable that actually predicts the rendered width is
% `inner sep`, not content type at all. At `inner sep=0pt` a *fully opaque*
% image measured 1.2pt; the exact same image at `inner sep=3pt` measured
% 2.52pt (matching the declared 2.5pt); a nested pgfplots axis and a plain
% nested tikzpicture, both at `inner sep=3pt`, also measured 2.52pt. Two
% more points confirmed it's a smooth transition, not a binary switch:
% `inner sep=0.6pt` measured 1.8pt (= the 1.2pt floor plus that 0.6pt gap,
% exactly); `inner sep=1.25pt` (= \SMImageBorderWidth/2, the stroke's
% own half-width) already measured the full 2.52pt.
%
% The mechanism: tikz draws a node's border stroke centered on the node's
% own boundary path, then paints the node's *content* on top of it,
% afterwards. At `inner sep=0`, the content's own bounding box coincides
% exactly with that boundary path, so if the content is opaque all the way
% to its own edge, it paints directly over the stroke's inner half,
% visually leaving only the outer half (`\SMImageBorderWidth`/2,
% ~1.2pt) -- not a rendering bug, just tikz's normal draw order interacting
% with zero padding. A gap of at least `\SMImageBorderWidth`/2 (or
% content that's transparent right at its own edge -- see below) leaves the
% border with nothing opaque to cover it, so it renders at its true
% declared width instead.
%
% This also explains why the fix landed on exactly the 4 slides it did,
% each for a different reason: slides/m2k-Exercise-CAPM-2-current.tex
% used `inner sep=3pt` for unrelated layout reasons (room for the axis's
% own tick labels), well past the 1.25pt threshold, so nothing ever covered
% its border. slides/n8e-Tail-Risk-Variance-Clustering-current.tex/
% slides/fm6-Tail-Risk-Market-vs-Simulated-current.tex/
% slides/xqq-Fat-Tails-current.tex all use `inner sep=0pt` -- matching the
% deck's own overwhelming convention -- but their own chart PNGs
% (`chart-index-daily-returns-*.png`/`chart-index-fat-tails-histogram*.png`,
% matplotlib exports with a transparent background) measured alpha=0 at
% every edge pixel (confirmed directly, Pillow), so there was no opaque
% content to cover their own border either. Every other `SMImageBorder`
% use in this project that was never flagged uses an opaque image (a JPG,
% or a PNG with no transparent margin) at `inner sep=0pt` -- which
% *accidentally* renders at the ~1.2pt covered-half width instead of the
% full 2.5pt this constant nominally declares, and that accidental width is
% what reads as "the deck's normal border thickness" throughout, since it's
% by far the most common case.
%
% Practical upshot for a *new* slide: a `SMImageBorder` node with
% `inner sep=0pt` wrapping opaque content (the common case) needs no
% special handling -- it will match the rest of the deck automatically. A
% node that instead uses a nonzero `inner sep`, or wraps content with a
% transparent edge (a matplotlib export, an SVG-derived PDF with padding,
% ...), will render at the *full* \SMImageBorderWidth and look
% noticeably heavier next to the deck's other borders -- either add an
% explicit `line width=\SMImageBorderWidth/2` (matching the 4 fixed
% slides' own literal 1.2pt) to match the accidental deck-wide look, or
% reduce `inner sep` below \SMImageBorderWidth/2 so the same coverage
% effect kicks in naturally. The 4 existing per-slide overrides above are
% left as-is -- they're already the correct fix for what's now a fully
% understood cause, not a mystery workaround needing revisiting.

% Reusable "mugshot" -- a circular-cropped photo, bordered like any other
% SMImageBorder image, with a caption curving along its own bottom arc
% (2026-07-06, direct instruction, generalized from
% slides/4uq-Definition-of-Capital-Market-Efficiency-current.tex's own
% Eugene Fama photo -- see that slide's earlier comment for the source
% PPTX effect this reproduces, `<a:prstGeom prst="ellipse">` +
% `<a:prstTxWarp prst="textArchDown">`, and CONVERSION_NOTES.md's own
% "beamer/TeX gotchas" entry, search "circular photo crop", for the
% clip/border/decorations.text mechanics themselves -- unchanged here,
% just factored out so a second photo doesn't need to repeat them).
%
%   \coordinate (myMugshotCenter) at (...);
%   \SMMugshot{myMugshotCenter}{media/photo-example.jpeg}{2.6cm}{Some Person, NP 2013}
%
% #1 the *name* of a coordinate the caller has already placed (not a raw
% point expression -- needed as a plain name, not just a point, since the
% caption's own arc path re-references it by name via `++(...)`, which
% only works with a named coordinate/node, not an arbitrary `($...$)`
% calc expression); #2 the image path; #3 the desired final diameter;
% #4 the caption text.
%
% Scales the source photo so it fully covers the circle with no gap
% anywhere -- both dimensions end up at least as big as the diameter --
% while using as much of the source photo as possible (2026-07-07,
% direct instruction: "should not include whitespace... scale the photo
% such that its edges just touch the circle's inner border"). This
% briefly (same day) scaled the *other* way -- by the photo's own
% diagonal, the largest scale at which the *entire* rectangle fits
% inside the circle with nothing clipped at all -- which does avoid
% ever clipping real content, but leaves the circle's own four corners
% uncovered (visible as plain background there, since a non-circular
% photo can only touch an inscribing circle at its corners, not along
% the whole rim) -- exactly the "whitespace" this round's feedback
% ruled out in favor of full coverage. Covering scale is
% `diameter / min(w,h)`: whichever of the photo's own two dimensions is
% *smaller* becomes exactly the diameter (that edge pair touches the
% circle's own inner border, per the direct instruction), and the other,
% larger dimension overhangs the diameter and gets clipped by the
% circular `\clip` -- for a portrait photo (every one used here so far),
% that means width becomes the diameter and a sliver clips off the top
% and bottom, which is the origin of the clipping problem the diagonal
% approach above existed to solve in the first place. Reintroduced
% deliberately: showing as much of the photo as the "no whitespace"
% constraint allows is now the stated priority, over never clipping
% anything.
%
% Measuring the photo's own native size needs `\savebox` (into
% \SMMugshotBox, declared once below, not a fresh box register per
% call) -- `\includegraphics` with no size option at all renders at the
% image's native resolution, and `\wd`/`\ht` then read that box's actual
% dimensions. pgfmath (already loaded for tikz) converts any TeX
% dimension fed into it to a plain pt-count number automatically, so
% arithmetic on the two measured lengths just works. `min(1,ratio)`
% (ratio = height/width) picks out whichever of the two is smaller
% *relative to the other*, without needing a `sqrt`/`^` at all this time
% -- so this version doesn't hit the "! Dimension too large" pgfmath
% overflow the diagonal version's own arithmetic did (see this file's
% own comment history / CONVERSION_NOTES.md's "beamer/TeX gotchas" for
% that point (easy to misread as a harmless warning, since compilation
% continues and produces *a* number regardless). Reformulating via the
% aspect ratio sidesteps the whole problem: `scale = diameter /
% sqrt(1+ratio^2)` is the same formula algebraically (factor `w` out of
% `sqrt(w^2+h^2)` and it becomes `w*sqrt(1+(h/w)^2)`), but every
% the full story of that overflow, from when this macro's own arithmetic
% did use a `sqrt`).
%
% Caption font reverted to this deck's normal body font (2026-07-07,
% direct instruction: "change the mugshot template back to using regular
% font instead of monospace") -- `\ttfamily` (Fira Mono) was a direct
% instruction on 2026-07-06, one day earlier; \normalfont resets
% explicitly rather than just omitting \ttfamily, so this doesn't
% silently inherit whatever font happens to be active at the call site.
\newsavebox{\SMMugshotBox}
\newcommand{\SMMugshot}[4]{%
	\savebox{\SMMugshotBox}{\includegraphics{#2}}%
	\pgfmathsetmacro{\SMMugshotImgW}{\wd\SMMugshotBox}%
	\pgfmathsetmacro{\SMMugshotImgH}{\ht\SMMugshotBox}%
	\pgfmathsetmacro{\SMMugshotRatio}{\SMMugshotImgH/\SMMugshotImgW}%
	\pgfmathsetmacro{\SMMugshotScaledW}{(#3)/min(1,\SMMugshotRatio)}%
	\begin{scope}
		\clip (#1) circle (0.5*#3);
		\node[anchor=center, inner sep=0pt] at (#1)
			{\includegraphics[width=\SMMugshotScaledW pt]{#2}};
	\end{scope}
	\draw[SMImageBorder] (#1) circle (0.5*#3);
	{\normalfont\tiny
	\path[decoration={text along path, text={#4}, text align=center}, decorate]
		(#1) ++(200:0.5*#3+0.25cm) arc (200:340:0.5*#3+0.25cm);
	}
}

% Reusable "warning sign" -- a rotated, no-fill, thick-bordered box with
% centered bold text, generalized (2026-07-24, direct instruction: "build a
% GENERAL reusable command for placing such signs") from the red
% "Extremely skewed outcomes!" callout on slides/awe-Example-Investment-
% Return-Simulation-current.tex's own source PPTX slide (233) -- read
% directly from that slide's `ppt/slides/slide233.xml`: a rotated
% no-fill rectangle (`rot="20410887"`, i.e. 340.18 degrees clockwise =
% -19.82 degrees), a thick solid border in the same color as the text
% (`<a:ln w="57150">` = 4.5pt, color C00000 = this project's own
% `c_custom_red`), text centered inside.
%
% Usage: \SMWarningSign[color]{x}{y}{text}{aspect ratio}{rotation}
% #1 (optional) the border/text color, defaults to \c_custom_red; #2/#3
% the sign's own center, as a page-absolute (x,-y) offset from
% "current page.north west" in pt -- the same coordinate convention as
% this deck's other absolute-overlay figures (e.g.
% slides/ivw-Divestment-of-Polluting-Plants-Overview-current.tex's own
% photo nodes); #4 the sign's text; #5 the desired width:height aspect
% ratio of the box; #6 the rotation angle in degrees, tikz convention
% (counterclockwise positive) -- note OOXML's own `rot` attribute on a
% source PPTX shape is clockwise positive instead, so a value read
% directly off a source slide needs negating first.
%
% The box is sized to fit the text *wrapped* at the requested aspect
% ratio, not a caller-supplied width -- matching how the source PPTX's
% own sign isn't a single unwrapped line either (rendering the actual
% source deck's own slide 233 to a flat image for comparison confirmed
% its text wraps across 3 lines inside a squarish box, not one long
% strip). #5 drives that wrap width via a closed-form estimate (start
% from the unwrapped text's own natural width #4 and single-line height
% #4, i.e. `sqrt(ratio*width*height)`, the width at which a box of
% roughly that many lines would hit the target ratio if text-ink area
% stayed constant under rewrapping) plus one correction pass (rewrap at
% that estimate, remeasure the real height, then solve
% `sqrt(width*height*ratio)` again treating the *measured* width*height
% as the area instead of the rough estimate) -- not exact, since
% rewrapping to a handful of lines isn't a smooth function of width, but
% close enough for the short warning-sign text this is meant for; a
% caller can always nudge #5 if a specific case lands visibly off.
% (\SMWarningSignPadding, 6pt, and \SMWarningSignBorderWidth, 4.5pt, are hooks of SlideMage_setup.tex.)
\newsavebox{\SMWarningSignBox}
\newlength{\SMWarningSignW}
\newlength{\SMWarningSignH}
\newcommand{\SMWarningSign}[6][c_custom_red]{%
	\savebox{\SMWarningSignBox}{\bfseries\color{#1}#4}%
	\pgfmathsetlength{\SMWarningSignW}{sqrt((#5)*\wd\SMWarningSignBox*(\ht\SMWarningSignBox+\dp\SMWarningSignBox))}%
	\savebox{\SMWarningSignBox}{\parbox{\SMWarningSignW}{\centering\bfseries\color{#1}#4}}%
	\pgfmathsetlength{\SMWarningSignW}{sqrt((#5)*\SMWarningSignW*(\ht\SMWarningSignBox+\dp\SMWarningSignBox))}%
	\savebox{\SMWarningSignBox}{\parbox{\SMWarningSignW}{\centering\bfseries\color{#1}#4}}%
	\setlength{\SMWarningSignH}{\ht\SMWarningSignBox+\dp\SMWarningSignBox}%
	% \vbox to 0pt{\vss ...} keeps this at zero space in the normal vertical
	% list, the same reason every other absolute-page-coordinate tikz overlay
	% in this file uses it (footer/logo/title-bar boxes, y5z's/22s's own
	% corner images, ...): a bare tikzpicture here is otherwise sized from
	% its own local origin (wherever this macro happens to be called in the
	% flow) out to the far-away absolute page coordinate the node is drawn
	% at, so without this wrapper the "sign" silently pushes real, sizeable
	% blank space into the frame's own body -- confirmed directly on
	% slides/awe-Example-Investment-Return-Simulation-current.tex: with this
	% call gated behind `\deckreveal[3]{...}`, the reveal step where it first
	% turns on pushed the *unrelated* left-column text down by about 100pt
	% relative to the steps before it, a "content moves between reveal
	% steps" bug rather than anything wrong with that text itself.
	\vbox to 0pt{\vss
	\begin{tikzpicture}[overlay, remember picture]
		\node[rotate=#6, draw=#1, line width=\SMWarningSignBorderWidth,
		      minimum width=\dimexpr\SMWarningSignW+2\SMWarningSignPadding\relax,
		      minimum height=\dimexpr\SMWarningSignH+2\SMWarningSignPadding\relax,
		      inner sep=0pt] at ($(current page.north west)+(#2,-#3)$)
			{\usebox{\SMWarningSignBox}};
	\end{tikzpicture}}%
}

% A \frac/\dfrac-stacked formula reads as cramped against a plain-text
% neighbor in an itemize -- it's visibly taller than a normal line (both
% above *and* below the line, since the fraction bar sits on the text
% baseline with numerator/denominator stacked around it), so both the
% fraction's own item *and* the item that follows it need a bit of extra
% headroom that a plain-text \item doesn't. Started as a manually-repeated
% \vspace{3pt} before each such \item (first on
% slides/u74-Example-Dynamic-Mattress-Company-Periods-current.tex);
% \tallitem below is that same convention factored into a single drop-in
% \item replacement, so a future fraction-heavy slide (or a change to the
% spacing amount) doesn't mean re-editing several slide files by hand.
% Takes an optional leading multiplier (default 1, i.e. 3pt): \tallitem[1.5]
% for 4.5pt, for a formula tall enough that the default doesn't read as
% enough room. Usage: replace \item with \tallitem both on any itemize line
% containing a stacked fraction *and* on the item immediately following it
% (same nesting level; skip this second replacement if that next item is
% itself already a \tallitem for its own fraction, or if the fraction item
% is the last one in its list) -- it's a plain substitute, not an
% environment, so it doesn't change how the surrounding list itself is
% structured. (Named \SMFracItem originally; renamed 2026-07-21 to the
% more discoverable \tallitem, and given the optional multiplier and the
% following-item rule above at the same time -- no behavior change for any
% existing default-strength call site.)
% Ratio, not a frozen length (2026-07-22, direct instruction: "modify
% \tallitem and \tallitemnobullet to also use scaling relative to font
% size") -- the original \newlength+\setlength{...}{3pt} pair freezes to
% one fixed pt value the moment \setlength runs, the same issue (and same
% fix) as \SMItemsepRatioI/II/III/IV further below: converted to a
% bare ratio macro, with the "em" unit written back on inside \tallitem's
% and \tallitemnobullet's own bodies (not here), so it's recomputed fresh
% against whatever font is actually live at each call site -- correct
% regardless of nesting level or a \footnotesize/\scriptsize override,
% unlike the old fixed 3pt (calibrated against one specific level's font,
% silently wrong at any other). 0.3 reproduces the original 3pt exactly
% at level 2's own font (10pt) -- the level \tallitem was first used at,
% on slides/u74-Example-Dynamic-Mattress-Company-Periods-current.tex (see
% this macro's own comment above) -- so no visual change there, while now
% correctly scaling up/down at every other level and font override.
% (\SMTallItemBaseRatio, 0.3, is a hook of SlideMage_setup.tex.)
\newlength{\SMTallItemBaseUnit}
\newlength{\SMTallItemSpace}
% The optional argument is either a number (the multiplier above -- any
% argument made only of digits and dots) or anything else, which is taken as
% a bullet label and passed on as \item[<label>] with the default spacing:
% \tallitem[\SMBulletEllipsis] gives an ellipsis-labelled item with the
% usual headroom (2026-09-26, direct instruction), \tallitem[] a label-less
% one. (For a label-less item with a custom multiplier, use
% \tallitemnobullet[<multiplier>].)
\ExplSyntaxOn
\cs_new_protected:Npn \SMTallItemIfNumberTF #1
	{ \regex_match:nnTF { \A [0-9.]+ \Z } {#1} }
\ExplSyntaxOff
\newcommand{\tallitem}[1][1]{%
	\setlength{\SMTallItemBaseUnit}{\SMTallItemBaseRatio em}%
	\SMTallItemIfNumberTF{#1}%
		{\setlength{\SMTallItemSpace}{#1\SMTallItemBaseUnit}%
		 \vspace{\SMTallItemSpace}\item}%
		{\setlength{\SMTallItemSpace}{\SMTallItemBaseUnit}%
		 \vspace{\SMTallItemSpace}\item[#1]}}

% For a line that needs \tallitem's own extra headroom but no bullet at
% all -- a display equation continuing the previous \item's own sentence,
% rather than a new bulleted point of its own (see README.md's own
% \tallitem section, which already documents this exact combination as a
% rare, previously hand-written case: "\vspace{\SMTallItemBaseUnit}\item[]").
% \tallitemnobullet (2026-07-21, direct instruction) factors that hand-
% written idiom into the same kind of drop-in replacement \tallitem
% already is for its own case, same optional multiplier: \tallitemnobullet[1.5].
\newcommand{\tallitemnobullet}[1][1]{%
	% Recomputes \SMTallItemBaseUnit itself, fresh, rather than reusing
	% whatever \tallitem's own body last left it at -- since it's now a
	% per-call, live "em"-based value (see \tallitem's own comment above),
	% not the one-time preamble-load-time constant it used to be, relying
	% on call order/a preceding \tallitem to have set it correctly would
	% silently break the moment \tallitemnobullet is used on its own, or
	% at a different font size than whichever \tallitem last ran.
	\setlength{\SMTallItemBaseUnit}{\SMTallItemBaseRatio em}%
	\setlength{\SMTallItemSpace}{#1\SMTallItemBaseUnit}%
	\vspace{\SMTallItemSpace}\item[]}

% flalign* (used to left-align a display equation inside an itemize --
% see CONVERSION_NOTES.md's "flalign* for left-aligning..." entry) inherits
% plain LaTeX's \abovedisplayskip/\belowdisplayskip, sized for a full
% document page, not a compact slide itemize -- reads as far too much blank
% space above/below the equation next to the surrounding item text (found
% 2026-07-21, direct instruction; most visible on
% slides/bwi-Pearson-Correlation-Definition-current.tex's own single-line
% equation, but the same oversized gap is present everywhere flalign* is
% used this way, e.g. slides/7se-Portfolio-Risk-Three-Assets-current.tex
% and slides/rb3-Asset-Allocation-Example-CAL-current.tex). \SMTightDisplaySkips
% tightens both skips -- and their "short" variants, since amsmath picks
% between the two depending on how much text precedes the display, so
% setting only one pair left the gap unfixed whenever the preceding line
% happened to be short -- down to something that reads as part of the
% same list rather than a separate floating display. Usage: wrap the
% *real*, unmodified flalign* in a plain TeX group together with this
% command, `{\SMTightDisplaySkips\begin{flalign*} ... \end{flalign*}}`,
% rather than a new same-named environment -- tried that first
% (`\newenvironment{SMTightFlalign*}{...\begin{flalign*}}{\end{flalign*}}`)
% and it broke LaTeX's own environment-nesting tracking (`\@currenvir`),
% surfacing as "\begin{flalign*} ... ended by \end{SMTightFlalign*}"
% and a fatal, PDF-less compile on the very first slide that used it. A
% plain group sidesteps that class of bug entirely since it never
% introduces a same-named wrapper environment.
% (\SMTightDisplaySkip, 3pt, is a hook of SlideMage_setup.tex.)
\newcommand{\SMTightDisplaySkips}{%
	\abovedisplayskip=\SMTightDisplaySkip
	\abovedisplayshortskip=\SMTightDisplaySkip
	\belowdisplayskip=\SMTightDisplaySkip
	\belowdisplayshortskip=\SMTightDisplaySkip}


% \SMFrametitleRightPadding (left/right padding around the logo, chosen
% to match its top/bottom margins -- measured against the title colorbox's
% own top and bottom edges -- rather than the theme's 2.2ex default) is
% also set by the per-university preamble, not here: the top/bottom margins
% are a side effect of vertically centering a \SMLogoHeight-tall logo in
% the title bar, not something set directly, so this value has to be
% re-measured and set per institution whenever its own logo/height differs
% from another's.

% The logo macro \SMInstituteLogo, the frame title hook that paints the white patch behind the
% logo and redraws it (\SMFrametitleHookCode), \chapterlogo and the template without the inline
% logo (SMNoInlineLogo) are the theme's (beamerthemeSlideMage.sty). The logo width, which the
% hook uses to size the white patch, is measured here from the template's logo path and height.
\settowidth{\SMLogoWidth}{\includegraphics[height=\SMLogoHeight]{\SMInstituteLogoPath}}



% Define slideID as a simple value holder (initially empty)
\newcommand{\slideID}{}%

% Set (locally, by \slidestep) when this inclusion of a slide id must not define
% the id's \hypertarget again, because an earlier \slidestep of the same id
% already owns it (see \slidestep's own comment).
\newif\ifSMSuppressSlideTarget

% Define slideID as a simple value holder (not a command with arguments)
\makeatletter
\newcommand{\setSlideID}[1]{%
	% Grab everything up to the first '-' and store it in \slideID
	\edef\tempa{#1}%
	\expandafter\setSlideID@aux\tempa-\relax
}
\def\setSlideID@aux#1-#2\relax{%
	\renewcommand{\slideID}{#1}%
	% Non-linear navigation (direct instruction, 2026-09-15): every real
	% slide inclusion -- \slide{}, \slidefinal{} (delegates to \slide{}),
	% \slideversion{}, \omitslide{}, \omitslideversion{}, \conditionalslide{}
	% -- calls \setSlideID{} immediately before \input{}ing the slide file,
	% regardless of which frame type (standardframe, sideimageframe,
	% figureframe, sectionframe, highlightframe, ...) that file happens to
	% use internally. This is the one place all of them funnel through, so
	% it is the correct place for a single, frame-type-independent
	% \hypertarget{} -- giving every slide a free, permanent jump target
	% keyed to its own existing id, with no per-slide authoring and no new
	% bookkeeping. Deliberately \hypertarget (hyperref's own mechanism, not
	% plain LaTeX \label{}: slides/ZZM-EN-Slidemaster-colors-current.tex's
	% own comment documents a real, already-hit collision between \label{}
	% and pgffor's \foreach inside a tikz loop, which this sidesteps
	% entirely. Used by \hyperlink{<bare id>}{...} anywhere in the deck
	% (README.md's own "Non-linear navigation" section) to jump straight to
	% this slide -- e.g. an excursus link's own "Back" button, or a
	% generated appendix's thumbnail grid (SlideMage's "Generate
	% appendix TOC" checkbox, appendix_toc.py).
	%
	% Known, pre-existing limitation, not introduced by this: if the same
	% id is genuinely \slide{}'d twice in one deck (it happens -- e.g. a
	% \slidefinal{} recap of an earlier slide), hyperref just warns
	% ("multiply defined") and the *later* occurrence's target wins, so a
	% link to that id lands on whichever instance is last in the document
	% -- inherent to any permanent-id-keyed target scheme, not specific to
	% this one.
	\ifSMSuppressSlideTarget\else\hypertarget{\slideID}{}\SM@agenda@slide{#1}{#2}\fi%
}
\makeatother

% \setSlideVersion{X-Y-Z} -- footer's own version display (added
% 2026-08-20, direct instruction: show the slide's version next to its id
% in every compiled deck). Takes the same filename-safe hyphenated form a
% version file's own "-vX-Y-Z.tex" suffix and \slideversion{}'s own second
% argument already use, and reformats it with dots ("1.0.0") for display,
% the conventional way semver is written -- \setSlideVersion@aux splits it
% the same way \setSlideID@aux above splits an id-slug, just at both
% hyphens instead of the first one only. \ifSMSlideVersionSet tracks
% whether *this* frame actually has a tracked version to show at all --
% reset false at the top of every \slide{}/\omitslide{}/\conditionalslide{}
% call (current.tex may or may not embed a \setSlideVersion{} of its own,
% see SlideMage's own automatic-versioning save flow) so an untracked
% slide's footer never inherits whatever the immediately preceding frame
% last set. \slideversion{}/\omitslideversion{} instead set it directly
% from their own second argument -- already an explicit version number at
% the call site, so there is nothing to read from the pinned file itself.
\newcommand{\slideVersion}{}%
\newif\ifSMSlideVersionSet
\makeatletter
\newcommand{\setSlideVersion}[1]{%
	\edef\SMTempSlideVersion{#1}%
	% No trailing "-" appended before \relax here, unlike \setSlideID@aux's
	% own \tempa-\relax above -- that extra hyphen is harmless there since
	% \setSlideID@aux's own #2 (everything after the first "-") is thrown
	% away entirely, but this pattern's #3 is a *used* greedy match
	% ("everything up to \relax"), so an appended "-" would land inside it
	% verbatim (confirmed directly: produced "v2.3.1-" instead of
	% "v2.3.1" the first time this was compiled).
	\expandafter\setSlideVersion@aux\SMTempSlideVersion\relax
	\global\SMSlideVersionSettrue
}
\def\setSlideVersion@aux#1-#2-#3\relax{%
	\renewcommand{\slideVersion}{#1.#2.#3}%
}
\makeatother

% \setSlideLanguageFromSlug{ID-slug} -- multi-language support (added
% 2026-08-24). Language lives in the id-slug's own *second* hyphen segment
% rather than a separate in-body macro call, so it's visible directly in
% Explorer/grep and in every \slide{}/\slideversion{}/etc. call site, not
% just inside the file it names: slides/<id>-EN-<slug>-current.tex,
% slides/<id>-DE-<slug>-current.tex, slides/<id>-mult-<slug>-current.tex
% (mult = multilingual, i.e. a single file holding both languages via
% \multilingualtext{}{} below, picked at compile time by \ifSMGerman).
% A slide that predates this feature simply has no such segment (its own
% second hyphen segment is the first word of an ordinary slug, e.g.
% "neb-Market-Environment-...") -- \ifSMSlideLanguageSet stays false
% for it, exactly like an untracked \setSlideVersion{} today shows no
% "v..." footer segment, so the other ~665 files need no migration.
% Same three-way split idiom as \setSlideVersion@aux above, but the
% trailing padding ("-zzz-zzz") is load-bearing here in a way it isn't
% there: \setSlideVersion@aux's own caller (X-Y-Z) always has exactly two
% hyphens, but an id-slug can legally have as few as *one* (e.g. a
% single-word slug like "abc-Word"), which would leave #2-#3\relax with no
% closing delimiter to match at all ("Use of ... doesn't match its
% definition") without this padding. #3 (and the padding itself) is never
% read, so its exact content doesn't matter -- only that the two extra
% hyphens are guaranteed to exist.
\newcommand{\slideLanguage}{}%
\newif\ifSMSlideLanguageSet
% \ifSMSlideLanguageMult -- set alongside \slideLanguage whenever it's
% specifically "mult", so the footer tag below never has to run
% \ifstrequal against \slideLanguage itself: etoolbox's \ifstrequal
% compares via \detokenize, which does not expand its argument -- fed a
% bare control sequence like \slideLanguage it would stringify the
% *macro's name*, not its meaning, and always take the false branch (see
% \decklanguage's own comment above, where this was caught the same way).
% A dedicated boolean, set right here where the value is still a plain
% matched string (#2, not a macro), sidesteps the trap entirely.
\newif\ifSMSlideLanguageMult
\makeatletter
\newcommand{\setSlideLanguageFromSlug}[1]{%
	\global\SMSlideLanguageSetfalse
	\global\SMSlideLanguageMultfalse
	\edef\SMTempSlideLangSlug{#1}%
	\expandafter\setSlideLanguageFromSlug@aux\SMTempSlideLangSlug-zzz-zzz\relax
}
\def\setSlideLanguageFromSlug@aux#1-#2-#3\relax{%
	\ifstrequal{#2}{EN}{\renewcommand{\slideLanguage}{EN}\global\SMSlideLanguageSettrue}{%
	\ifstrequal{#2}{DE}{\renewcommand{\slideLanguage}{DE}\global\SMSlideLanguageSettrue}{%
	\ifstrequal{#2}{mult}{\renewcommand{\slideLanguage}{mult}\global\SMSlideLanguageSettrue\global\SMSlideLanguageMulttrue}{}}}%
}
\makeatother

% Define a new \slide command to automatically pick the current version of the slide
% 2026-07-08 (direct instruction): slide filenames dropped their leading
% "s-" (e.g. slides/ab3-Some-Slide-current.tex, not the old
% slides/ab3-Some-Slide-current.tex) so slide IDs are directly
% searchable in Explorer -- this macro (and the other three below) is the
% *only* place that ever needed to change for that: every deck's own
% \slide{}/\slideversion{}/\omitslide{}/\omitslideversion{} calls pass just
% "ID-slug", never the filename itself, so renaming every file on disk to
% match needed no changes anywhere else (not the deck files, not
% slide_ids.csv, which only ever stored bare IDs to begin with).
\newcommand{\slide}[1]{%
	\global\SMSlideVersionSetfalse
	\setSlideID{#1}%
	\setSlideLanguageFromSlug{#1}%
	\input{slides/#1-current.tex}%
}

% \slidefinal{ID-slug} -- like \slide{}, but renders the slide at its final,
% fully-revealed overlay state in every deck version, instructor included
% (2026-07-12, direct instruction: show jww's finished build in place of a
% separate recap slide, "irrespective of which deck version is being
% compiled" -- that recap slide has since been deleted entirely, its own ID
% no longer in use). A slide
% built with \conditionalpause[instructor] already renders fully built in
% the during-course/after-course compiles -- that pause is a no-op for
% them, so nothing there ever waits on a reveal step -- the only case
% actually needing help is the instructor deck, which would otherwise still
% step through the build. Locally shadowing \conditionalpause to a no-op
% (same optional-argument signature as the real one, so it still accepts
% e.g. \conditionalpause[instructor] without erroring) for the duration of
% one \slide{} call makes every \conditionalpause in that one inclusion
% inert, regardless of the actual compiled \SMDeckVersion -- scoped
% inside \begingroup/\endgroup, so it reverts immediately after, and has no
% effect on any *other* \slide{}/\slidefinal{} call to the same ID
% elsewhere in the same deck (e.g. an earlier, normally-building instance of
% the same slide). \deckonly/\deckuncover/\deckvisible are untouched, so
% deck-version-tagged content (as opposed to reveal-step-tagged content)
% still follows its usual per-version rules.
\newcommand{\slidefinal}[1]{%
	\begingroup
		\renewcommand{\conditionalpause}[1][instructor]{}%
		\slide{#1}%
	\endgroup
}

% Load a specific, pinned version of a slide instead of whichever is "current"
% Usage: \slideversion{ID-slug}{M-m-p} loads slides/ID-slug-vM-m-p.tex
\newcommand{\slideversion}[2]{%
	\setSlideID{#1}%
	\setSlideLanguageFromSlug{#1}%
	\setSlideVersion{#2}%
	\input{slides/#1-v#2.tex}%
}

% Student-version decks: \studentversion (called anywhere in a deck, e.g.
% right after \input{SlideMage_template_uni-graz}, similar to \sidefootnotes or
% \setfootnotesymbols) switches that compile into the "handout" variant --
% including, as of 2026-07-03, actually forcing beamer's own "handout" mode
% (\gdef\beamer@currentmode{handout}, normally only reachable via the
% \documentclass[handout]{beamer} *class option*, set once at class-loading
% time -- too early for a command like \studentversion, called later from a
% deck's own preamble, to use that route). This was checked, not assumed:
% until this point \studentversion only ever touched \omitslide{}, so
% nothing suppressed \pause/overlay specs -- true only by coincidence, since
% no slide converted so far happens to use them. Late-redefining
% \beamer@currentmode was verified (isolated minimal-beamer tests, plain
% \pause *and* inline specs like \item<2->) to reproduce the same collapse
% -to-one-page, show-the-final-revealed-state behavior as the class option,
% even set well after \documentclass{beamer} -- beamer's overlay decoder
% (beamerbasedecode.sty) reads \beamer@currentmode at frame-processing time,
% not just once at class-option time, so this is safe, not a hack that
% happens to work by luck of file ordering.
% \omitslide{}/\omitslideversion{} are drop-in replacements for
% \slide{}/\slideversion{} that tag a slide as excluded from student deck versions:
%
%   \omitslide{xyz-Answer-key}                 % instead of \slide{...}
%   \omitslideversion{xyz-Answer-key}{1-0-0}    % instead of \slideversion{...}{...}
%
% In a \studentversion compile or a deck version other than instructor,
% a tagged slide's \input{} is skipped entirely and its frame number still
% advances, keeping numbering aligned across deck versions. In an
% instructor compile (the default), a tagged slide loads exactly as it
% would via \slide{}/\slideversion{}, with
% one difference: \SMOmittedSlide is set for the duration of that one
% \input{}, so the footer template below can print a marker next to that
% frame's slide ID -- a deliberately visible reminder, while reviewing the
% instructor deck, of which slides won't make it to students. The same tag
% works for all deck versions because it lives at the call site in the *deck*
% file (not inside the slide's own content file): the identical slide list
% can serve the instructor deck and both student deck versions, e.g.
%
%   \input{SlideMage_template_uni-graz}
%   \FooterText{...}
%   \makeatletter\ifdefined\HANDOUT\studentversion\fi\makeatother  % or similar
%   \input{course-content.tex}   % the shared \slide{}/\omitslide{} list
%
% with two thin wrapper decks (one plain, one calling \studentversion first)
% both \input-ing the same shared content file -- see README.md.
\newif\ifSMStudentVersion
% \makeatletter must wrap the \newcommand itself, not sit *inside* the macro
% body: catcodes apply at tokenization time, which for a macro body happens
% once, immediately, when \newcommand reads it -- not deferred until the
% macro is later called. \makeatletter placed inside the body (tried first)
% was too late: \beamer@currentmode had already been mis-tokenized as
% "\beamer" + the six-way-otherwise-ordinary character "@" + eleven plain
% letters "currentmode", by the time that inner \makeatletter could act on
% it, so the \gdef silently created/redefined the wrong thing and never
% touched beamer's actual mode variable. Confirmed via \meaning\beamer@
% currentmode before/after \studentversion -- unchanged with the inside-the-
% body version, fixed once \makeatletter moved out here.
\makeatletter
\newcommand{\studentversion}{%
	\global\SMStudentVersiontrue
	\gdef\beamer@currentmode{handout}%
}
\makeatother

% Deck version system - new boolean flags and string for current version
\newif\ifSMInstructorVersion
\newif\ifSMDuringCourseVersion
\newif\ifSMAfterCourseVersion
\gdef\SMDeckVersion{instructor}
\gdef\SMDeckVersionText{}

% \deckversion[] - set the deck version(s)
% Usage: \deckversion[instructor], \deckversion[during-course], etc.
% Multiple versions: \deckversion[instructor,during-course,after-course]
% Parses #1 as a comma-separated list via etoolbox's \forcsvlist (rather
% than the single \ifstrequal{whole-list}{"instructor"} check this used to
% do, which only ever matched a single-version list -- a multi-version
% \deckversion[instructor,during-course] silently set none of the three
% booleans, since the compared string was the literal, unsplit
% "instructor,during-course"). Note \forcsvlist, not etoolbox's similarly-
% named \forlistloop: \forlistloop iterates etoolbox's own internal
% pipe-delimited list format (built via \listadd), not a plain comma string
% like this command's own argument -- confirmed directly (isolated test),
% passing a raw comma string to it doesn't error or split on the commas, it
% just treats the whole string as a single one-item "list".
\makeatletter
\newcommand{\deckversion}[1][instructor]{%
	% \xdef (not \gdef): #1 needs to be fully expanded here, not just
	% stored as whatever token(s) were passed -- \deckversion[\myvar]
	% (compile-all-versions.bat's own override mechanism, via a wrapper
	% file's \def\deckversionoverride{...}) passes a macro, not a literal
	% string, and everywhere else in this deck-version system that reads
	% \SMDeckVersion only ever applies one \expandafter to it (see
	% \ifdeckversioninlist's own comment) -- enough to unwrap \SMDeckVersion
	% itself, but not a further macro still sitting inside it. Confirmed
	% directly: with the old \gdef, \deckversion[\deckversionoverride]
	% left \SMDeckVersionText empty and every \ifdeckversioninlist
	% check false, since \SMDeckVersion's own stored content was the
	% single unexpanded token \deckversionoverride, not "during-course".
	\xdef\SMDeckVersion{#1}%
	\global\SMInstructorVersionfalse
	\global\SMDuringCourseVersionfalse
	\global\SMAfterCourseVersionfalse
	\gdef\SMDeckVersionText{}%
	\def\do@deckversions##1{%
		\ifstrequal{##1}{instructor}{\global\SMInstructorVersiontrue\gdef\SMDeckVersionText{instructor deck}}{}%
		\ifstrequal{##1}{during-course}{\global\SMDuringCourseVersiontrue\gdef\SMDeckVersionText{student course deck}}{}%
		\ifstrequal{##1}{after-course}{\global\SMAfterCourseVersiontrue\gdef\SMDeckVersionText{student final deck}}{}%
	}%
	% \forcsvlist itself splits on literal commas in its own raw,
	% unexpanded argument (confirmed directly against etoolbox's own
	% source, \etb@defparser@do@aux) -- passing #1 straight through here
	% has the exact same problem as the \gdef above did, one call site
	% later. Building the whole \forcsvlist{...}{...} call inside an
	% \edef (protecting \forcsvlist/\do@deckversions themselves with
	% \noexpand so they're invoked *later*, once this constructed call is
	% run, not expanded away here) expands only \SMDeckVersion --
	% already fully resolved by the \xdef above -- into the list argument.
	\edef\SM@deckversion@call{\noexpand\forcsvlist{\noexpand\do@deckversions}{\SMDeckVersion}}%
	\SM@deckversion@call
}
\makeatother

% \deckoutputfolder[version]{path} - declares where SlideMage should
% place this deck's own compiled PDF, read directly out of this file's own
% source text by SlideMage/deck_scan.py (OUTPUT_FOLDER_RE) at deck-scan
% time -- never expanded by LaTeX itself. A silent no-op here (2026-09-16,
% direct instruction) purely so a deck carrying this directive still
% compiles cleanly under a plain lualatex/TeXstudio run, the same reason
% \deckversionoverride's own \ifdefined door (used elsewhere in this file)
% never needs a real definition either. #1 (the optional [version]) is
% intentionally unused: unlike \deckversion above, this macro's own
% argument is never inspected by LaTeX at all, only by SlideMage's own
% Python-side regex -- see README.md's "Custom output folder" section for
% the actual directive syntax.
\newcommand{\deckoutputfolder}[2][]{}

% \decktags{tag1,tag2} - declares this deck's tags (2026-09-19), read
% directly out of this file's own source text by SlideMage/deck_scan.py
% (DECKTAGS_RE) and never expanded by LaTeX itself -- a silent no-op here
% for exactly the same reason \deckoutputfolder above is. SlideMage's deck
% tables group decks by their tag: work (the pinned workdeck), course,
% research, test; a deck with none of those is a course deck. Written and
% edited from SlideMage itself (the "Tags" button, and the Duplicate
% dialog), placed just before \begin{document}.
\newcommand{\decktags}[1]{}

% Deck language system (added 2026-08-24) -- fixed per deck, like
% \FooterText or the institution \input{} choice, not a per-compile-run
% toggle like \deckversion above: a German course is its own deck file
% with mostly \setSlideLanguageFromSlug-tagged German/multilingual slide
% IDs, not the same deck file recompiled twice. \decklanguage[en]/
% \SMDeckLanguage/\ifSMGerman mirror \deckversion/\SMDeckVersion
% mechanically (same \xdef for the same \deckversion[\deckversionoverride]-
% style macro-argument reason documented above), so every deck gets the
% same \ifdefined\decklanguageoverride ... \fi override door
% \deckversionoverride already has, for SlideMage's own one-off render
% preview (slide_preview.py) and for ad-hoc hand-testing via
% decks/#workdeck.tex -- never wired into the Active-decks table itself.
\newif\ifSMGerman
\gdef\SMDeckLanguage{en}
\makeatletter
\newcommand{\decklanguage}[1][en]{%
	\xdef\SMDeckLanguage{#1}%
	\global\SMGermanfalse
	% \ifstrequal{<macro>}{...} is a trap, not just here but anywhere in this
	% file: its etoolbox implementation compares via
	% \edef\tempX{\detokenize{<arg>}} -- \detokenize does NOT expand its
	% argument first, so handed a bare control sequence it stringifies the
	% *macro's name* ("\SMDeckLanguage"), never its meaning ("de"),
	% and the comparison silently always takes the false branch (confirmed
	% directly: \ifSMGerman never went true this way, even under
	% \decklanguage[de]). \deckversion's own \do@deckversions above sidesteps
	% this by building the comparison inside an \edef first, so \ifstrequal
	% only ever sees an already-expanded literal string -- same fix here.
	\def\do@decklanguage##1{%
		\gdef\SMPDFLanguage{##1}% any other language: its own code (fr, es, ...)
		\ifstrequal{##1}{en}{\gdef\SMPDFLanguage{en-US}}{}%
		\ifstrequal{##1}{de}{\global\SMGermantrue\gdef\SMPDFLanguage{de-DE}}{}}%
	\edef\SM@decklanguage@call{\noexpand\do@decklanguage{\SMDeckLanguage}}%
	\SM@decklanguage@call
	% Document-level PDF /Lang, independent of (and not blocked by) the
	% structural PDF/UA tagging this project can't use under beamer -- see
	% \SMArtifactBegin's own comment above for why full tagging is out;
	% a screen reader still uses this Catalog-level entry to pick
	% pronunciation for the whole document.
	\hypersetup{pdflang=\SMPDFLanguage}%
}
\makeatother

% \multilingualtext{EN text}{DE text} -- for the handful of light,
% image-only/title-only slides (id-slug's own language segment "mult",
% see \setSlideLanguageFromSlug above) that hold both languages in one
% file rather than being split into two separate slide IDs, switching on
% whichever language the *deck* is compiling as. The two-argument form is
% EN then DE, as it always was.
%
% \multilingualtext*{EN={Hello}, DE={Hallo}, FR={Bonjour}} -- the keyed form
% (2026-10-05, direct instruction), for any number of languages: one
% key=value pair per language, the key being the deck's \decklanguage code
% (upper or lower case: EN = en). The value is picked whose key is the
% deck's language; a deck in a language the list does not have gets the EN
% value, or the first one if there is no EN. Commas and equals signs inside a
% value need the braces shown. A deck in another language than en/de just says
% \decklanguage[fr] (fr, es, ...): \ifSMGerman stays false for it, and the
% PDF's /Lang is the code itself.
\ExplSyntaxOn
\tl_new:N \l__sm_mlt_pick_tl
\tl_new:N \l__sm_mlt_en_tl
\tl_new:N \l__sm_mlt_first_tl
\bool_new:N \l__sm_mlt_pick_bool
\bool_new:N \l__sm_mlt_en_bool
\bool_new:N \l__sm_mlt_first_bool
\cs_new:Npn \__sm_mlt_pair:nn #1#2 { \ifSMGerman #2 \else #1 \fi }
\cs_new_protected:Npn \__sm_mlt_item:nn #1#2
  {
    \bool_if:NF \l__sm_mlt_first_bool
      { \tl_set:Nn \l__sm_mlt_first_tl {#2} \bool_set_true:N \l__sm_mlt_first_bool }
    \str_if_eq:eeT { \str_casefold:n {#1} } { en }
      { \tl_set:Nn \l__sm_mlt_en_tl {#2} \bool_set_true:N \l__sm_mlt_en_bool }
    \str_if_eq:eeT { \str_casefold:n {#1} } { \str_casefold:V \SMDeckLanguage }
      { \tl_set:Nn \l__sm_mlt_pick_tl {#2} \bool_set_true:N \l__sm_mlt_pick_bool }
  }
\cs_new_protected:Npn \__sm_mlt_item_bad:n #1
  { \msg_warning:nnn { slidemage } { multilingual-key-without-value } {#1} }
\msg_new:nnn { slidemage } { multilingual-key-without-value }
  { \string\multilingualtext*:~"#1"~has~no~=~value;~ignored. }
\cs_new_protected:Npn \__sm_mlt_keyed:n #1
  {
    \group_begin:
      \bool_set_false:N \l__sm_mlt_pick_bool
      \bool_set_false:N \l__sm_mlt_en_bool
      \bool_set_false:N \l__sm_mlt_first_bool
      \keyval_parse:NNn \__sm_mlt_item_bad:n \__sm_mlt_item:nn {#1}
      \bool_if:NTF \l__sm_mlt_pick_bool
        { \exp_last_unbraced:NV \group_end: \l__sm_mlt_pick_tl }
        {
          \bool_if:NTF \l__sm_mlt_en_bool
            { \exp_last_unbraced:NV \group_end: \l__sm_mlt_en_tl }
            { \exp_last_unbraced:NV \group_end: \l__sm_mlt_first_tl }
        }
  }
\NewDocumentCommand{\multilingualtext}{s m}
  { \IfBooleanTF {#1} { \__sm_mlt_keyed:n {#2} } { \__sm_mlt_pair:nn {#2} } }
\ExplSyntaxOff

\newif\ifSMOmittedSlide
\newcommand{\SMOmittedSlideMarker}{ | omitted from student deck}
% \stepcounter{framenumber} in the skipped branches below keeps the student
% deck's slide numbers matching the instructor deck's, gap and all, so "see
% slide 12" means the same physical slide to both audiences: beamer steps
% "framenumber" exactly once per \begin{frame} (checked in
% beamerbaseframe.sty -- once per *frame*, not per overlay page, so this is
% right regardless of how many overlay steps the omitted slide would have
% had), and this is the one place that increment is skipped along with the
% frame itself, so it has to be replayed by hand.
% Both forms are instructor-only; \studentversion also skips them.
\newcommand{\omitslide}[1]{%
	\ifSMStudentVersion
		\stepcounter{framenumber}%
	\else
		\ifdeckversioninlist{\SMDeckVersion}{instructor}%
			{\global\SMOmittedSlidetrue
			 \global\SMSlideVersionSetfalse
			 \setSlideID{#1}%
			 \setSlideLanguageFromSlug{#1}%
			 \input{slides/#1-current.tex}%
			 \global\SMOmittedSlidefalse}%
			{\stepcounter{framenumber}}%
	\fi
}
\newcommand{\omitslideversion}[2]{%
	\ifSMStudentVersion
		\stepcounter{framenumber}%
	\else
		\ifdeckversioninlist{\SMDeckVersion}{instructor}%
			{\global\SMOmittedSlidetrue
			 \setSlideID{#1}%
			 \setSlideLanguageFromSlug{#1}%
			 \setSlideVersion{#2}%
			 \input{slides/#1-v#2.tex}%
			 \global\SMOmittedSlidefalse}%
			{\stepcounter{framenumber}}%
	\fi
}

% Deck version-aware overlay commands
% Syntax: \deckonly<overlay>[versions]{content}
% If [versions] omitted, defaults to all decks (instructor,during-course,after-course)
%
% Built on etoolbox's own \forcsvlist for the comma-splitting (rather than a
% hand-rolled parser: an earlier version tried to split a list at each comma
% using \newcommand's plain, non-delimited argument grabbing
% (\newcommand{\ifdeckversioninlist@split}[3]{...} called as
% \ifdeckversioninlist@split#1\@nil{#1}) -- but a plain \newcommand doesn't do
% delimited-parameter matching against \@nil the way a \def with a matching
% pattern would; it just grabs the next 3 tokens/groups one at a time. So
% e.g. "instructor" silently split into single characters ("i", "n", "s",
% ...) instead of being read as one list item, which eventually ran the
% parser off the end of its own input and surfaced as "Undefined control
% sequence" wherever the first \conditionalpause[...]/\deckonly[...] in a
% deck happened to be used -- confirmed directly against the compile log,
% traced to slide y5z-CoCo-Bonds's own \conditionalpause[instructor].
%
% NOT etoolbox's own \ifinlist{item}{list}, despite the name similarity: that
% one tests membership in etoolbox's own internal list format, built only via
% \listadd (items joined with a literal "|", not ","); fed a plain comma
% string instead, confirmed directly (isolated test) to silently return
% false for every item, including an exact single-item match.
%
% The membership test itself needs one more subtlety: #1 here is the deck
% version to search for -- always \SMDeckVersion at every call site
% below, i.e. a macro, not a literal string -- so comparing it via a plain
% \ifstrequal{##1}{#1} does not work either: etoolbox's \ifstrequal compares
% via \detokenize, which (confirmed directly, isolated test) stringifies an
% unexpanded macro token to its own control-sequence *name* ("SMDeckVersion"),
% not the current value that macro holds. \expandafter forces that expansion
% before \ifstrequal ever sees it.
\makeatletter

\let\SM@original@only\only
\let\SM@original@uncover\uncover
\let\SM@original@visible\visible

\newif\ifSM@deckversion@found

% Usage: \ifdeckversioninlist{version}{list}{true-branch}{false-branch}
\newcommand{\ifdeckversioninlist}[2]{%
	\global\SM@deckversion@foundfalse
	\def\SM@deckversion@check##1{%
		\expandafter\ifstrequal\expandafter{#1}{##1}{\global\SM@deckversion@foundtrue}{}%
	}%
	\forcsvlist{\SM@deckversion@check}{#2}%
	\ifSM@deckversion@found
		\expandafter\@firstoftwo
	\else
		\expandafter\@secondoftwo
	\fi
}

% \@ifnextchar-based parsing (the same idiom LaTeX2e's own kernel commands
% use for an optional argument that isn't a plain trailing [...]) so the
% overlay spec reads exactly like beamer's own \only<...>, with the deck-
% version list as a further, genuinely optional [...] afterwards:
%   \deckonly<2>{content}                              visible in all versions
%   \deckonly<2>[instructor]{content}                  instructor only
%   \deckonly<2>[during-course,after-course]{content}  during + after
\newcommand{\deckonly}{\@ifnextchar<{\SM@deckonly@overlay}{\SM@deckonly@overlay<1->}}
\def\SM@deckonly@overlay<#1>{\@ifnextchar[{\SM@deckonly@versions<#1>}{\SM@deckonly@versions<#1>[instructor,during-course,after-course]}}
\def\SM@deckonly@versions<#1>[#2]#3{%
	\ifSMInstructorVersion
		\SM@original@only<#1>{#3}%
	\else
		\ifdeckversioninlist{\SMDeckVersion}{#2}{\SM@original@only<#1>{#3}}{}%
	\fi
}

% Same interface as \deckonly, wrapping \uncover and \visible respectively
\newcommand{\deckuncover}{\@ifnextchar<{\SM@deckuncover@overlay}{\SM@deckuncover@overlay<1->}}
\def\SM@deckuncover@overlay<#1>{\@ifnextchar[{\SM@deckuncover@versions<#1>}{\SM@deckuncover@versions<#1>[instructor,during-course,after-course]}}
\def\SM@deckuncover@versions<#1>[#2]#3{%
	\ifSMInstructorVersion
		\SM@original@uncover<#1>{#3}%
	\else
		\ifdeckversioninlist{\SMDeckVersion}{#2}{\SM@original@uncover<#1>{#3}}{}%
	\fi
}

\newcommand{\deckvisible}{\@ifnextchar<{\SM@deckvisible@overlay}{\SM@deckvisible@overlay<1->}}
\def\SM@deckvisible@overlay<#1>{\@ifnextchar[{\SM@deckvisible@versions<#1>}{\SM@deckvisible@versions<#1>[instructor,during-course,after-course]}}
\def\SM@deckvisible@versions<#1>[#2]#3{%
	\ifSMInstructorVersion
		\SM@original@visible<#1>{#3}%
	\else
		\ifdeckversioninlist{\SMDeckVersion}{#2}{\SM@original@visible<#1>{#3}}{}%
	\fi
}

% \conditionalpause[versions] - \pause, but only takes effect in the listed
% deck versions (default: instructor only); a no-op, not even a skipped
% overlay step, in any other version
% Usage: \conditionalpause[instructor,during-course] or \conditionalpause[after-course]
\newcommand{\conditionalpause}[1][instructor]{%
	\ifdeckversioninlist{\SMDeckVersion}{#1}{\pause}{}%
}

% \deckreveal[step]{content} - the "hidden answer, click-revealed for the
% instructor, absent for during-course, pre-shown for after-course" pattern
% already documented by hand on several exercise slides (e.g. mbz/x7q/37s/
% 3cv's own header comments: "a slide/reveal-step that is hidden in the
% source PowerPoint should, by default, only be made visible on the
% instructor version"). A bare, unconditional \only<N->{content} looks like
% it does this, but doesn't: \only's overlay range registers with beamer's
% own overlay-counting pass regardless of \ifSMInstructorVersion/
% \conditionalpause -- neither of those touch \only itself, they only guard
% *other* content around it -- so a frame with nothing else establishing a
% second overlay step still gets a real second slide out of a bare
% \only<2->, in *every* deck version, not just the instructor one. Confirmed
% directly (2026-07-21, direct instruction) on 37s/3cv: both a) leaked the
% "hidden" content to during-course as a genuine second slide, the exact
% thing this pattern exists to prevent, and b) showed a blank second bullet/
% cell first on that leaked slide's own step 1 -- \only<2-> hides the
% *content*, not the space it would occupy, so during-course's own step 1
% still had nothing else deciding where the bullet/cell landed.
% \deckreveal fixes both by picking a genuinely different rendering per
% version rather than reusing the same overlay-restricted \only call
% everywhere: instructor gets the real overlay build (\only<step->,
% default step 2, matching \tallitem's own default-multiplier convention);
% during-course gets nothing at all (no token emitted, so beamer's overlay
% pass never sees an overlay spec to register); after-course gets the
% content unconditionally, with no overlay spec of its own either, so it
% doesn't silently force a second slide there either. Content-level only
% (wraps a value/expression, like \only itself) -- doesn't cover an
% itemize \item's own bullet marker, which needs its own overlay spec
% (`\item<step->`) directly on the \item, not wrapped from outside; see
% 37s's own comment for that variant, hand-written rather than factored
% into a second macro since only one call site needs it so far.
% (2026-09-30: that variant is now \deckitem<step->, below.)
\NewDocumentCommand{\deckreveal}{s O{2} m}{%
	\ifSMDuringCourseVersion
	\else
		\ifSMInstructorVersion
			\IfBooleanTF{#1}{\visible<#2->{#3}}{\only<#2->{#3}}%
		\else
			#3%
		\fi
	\fi
}

% \deckrevealfinal[step]{content} - reveal step for the instructor deck only,
% everything shown at once in every other deck version (added 2026-09-19,
% direct instruction: "make the student versions show the final slide
% version immediately"). Unlike \deckreveal, the during-course deck gets the
% content too (not omitted), and like it no overlay spec is emitted outside
% the instructor version, so a non-instructor compile has no extra pages.
% For plain \pause steps use \conditionalpause[instructor] instead.
\NewDocumentCommand{\deckrevealfinal}{s O{2} m}{%
	\ifSMInstructorVersion
		\IfBooleanTF{#1}{\visible<#2->{#3}}{\only<#2->{#3}}%
	\else
		#3%
	\fi
}

% Starred forms, \deckreveal*[step]{content} and \deckrevealfinal*[step]{content}
% (2026-09-30, direct instruction: systematic replacements for per-slide
% fixes): the instructor steps before `step` keep the content's space
% (\visible instead of \only), so nothing around it moves when it appears --
% the "the diagram/equation/text jumps between reveal steps" class of
% complaint. Replaces the per-slide helpers that did the same by hand
% (\only<N->{x}\only<1..N-1>{\phantom{x}}; vxd's \vxdRevealB/C/D, c4c's
% \cRevealVisible). Other versions behave exactly as the unstarred forms.

% \deckitem<spec> - a list item whose overlay spec only applies in the
% instructor version; every other version gets a plain \item (the
% "student versions show the final state" rule). Replaces the pattern
%   \ifSMInstructorVersion \item<3-> text \else \item text \fi
% which repeated each item's text twice (42 items in 12 slides by
% 2026-09-30). An item label may follow: \deckitem<2->[\SMBulletEllipsis].
% Without a <spec> it is a plain \item.
\newcommand{\deckitem}{\@ifnextchar<{\SM@deckitem}{\item}}
\def\SM@deckitem<#1>{%
	\ifSMInstructorVersion
		\def\SM@next{\item<#1>}%
	\else
		\let\SM@next\item
	\fi
	\SM@next
}

% \conditionalslide[versions]{ID-slug} - \slide{}, but only included in the
% listed deck versions; keeps the frame counter in step with the omitted
% version otherwise (same \stepcounter{framenumber} technique \omitslide{}
% already uses for the student deck), so "see slide N" still points at the
% same physical slide across every deck version.
% Usage: \conditionalslide[instructor,after-course]{xyz-Answer-key}
\newcommand{\conditionalslide}[2][instructor,during-course,after-course]{%
	\ifdeckversioninlist{\SMDeckVersion}{#1}%
		{\global\SMSlideVersionSetfalse
		 \setSlideID{#2}%
		 \setSlideLanguageFromSlug{#2}%
		 \input{slides/#2-current.tex}}%
		{\stepcounter{framenumber}}%
}

% \slidestep{ID-slug}[versions]{steps} - \conditionalslide{} that also restricts
% which reveal (overlay) steps of the slide are typeset (2026-09-20, direct
% instruction). One slide file with several reveal steps can then appear at
% several places of a deck, each time showing only the steps that belong
% there, and a different set of steps per deck version:
%   \slidestep{egk-EN-Structure-of-a-Scientific-Paper-Overview}[instructor]{1,2}
%   \slidestep{egk-EN-Structure-of-a-Scientific-Paper-Overview}[during-course,after-course]{1}
% [versions] defaults to instructor,during-course,after-course (same as
% \conditionalslide; a presenter compile runs as version "instructor", so it
% always shows what the instructor deck shows); {steps} is a beamer overlay
% specification without the angle brackets ("1,2", "3", "2-", "1-2,4").
% The slide file itself is written once, with plain overlay specifications
% (\textcolor<2>{...}{...}, \only<3>{...}) and no version conditionals: which
% steps end up in the PDF is decided here, by handing the specification to
% beamer's own frame-level overlay restriction, \begin{frame}<steps>. Every
% frame environment opens its frame through \SMOpenFrame below, which
% inserts it.
%
% - Versions not in the list: the call does nothing at all -- not even a
%   \stepcounter{framenumber}, unlike \conditionalslide{}. Two \slidestep calls
%   at one position (say [instructor]{1,2} and [during-course,after-course]{1})
%   are alternatives, so every deck version has exactly one frame there;
%   stepping the counter for the omitted one would shift every later slide
%   number by one.
% - Jump target: the first \slidestep{} of a given id-slug in the document
%   defines the id's \hypertarget (see \setSlideID); later ones do not, so
%   there is no "multiply defined" warning and \hyperlink{id}{...} lands on the
%   first place the slide is shown.
% - \SMSlideSteps is cleared again right after the slide file has been
%   read, so it can never leak into a following plain \slide{}.
\gdef\SMSlideSteps{}
% \begin{frame}<steps>[options]{title}: the <...> has to be the very next
% token after "\begin{frame}" (beamer looks for it with \@ifnextchar), so the
% whole opening is built by \edef and only then executed. #1 is everything that
% follows "\begin{frame}" in the caller ("[#1]{#2}", "[plain,t]", ...).
% \SMFinalStateForStudents -- put it in a slide file, before the frame environment: in the
% student deck versions (everything but "instructor") that frame then typesets only its last
% overlay page, i.e. the final state of all its reveal steps (2026-09-21, direct instruction:
% "student decks should only show the final state"). Instructor and presenter compiles are
% unaffected. Unlike \conditionalpause/\deckrevealfinal, which each turn one reveal
% construct off, this needs no change to the slide body, so it also covers overlay specs
% like \only<2>, \item<3->, \uncover, \textcolor<2>, \alt and reveal steps that hide things
% again (the final state is what is left at the last step).
% How: beamer's frame-level overlay restriction \begin{frame}<N> (see \slidestep) with N =
% the frame's overlay count, which a frame only knows after it was typeset once. Every
% student compile therefore records N per frame (keyed by beamer's frame number, which is
% stable) in \jobname.ovl and reads it back at the start of the next run; the first run of
% a fresh job still typesets all pages, latexmk's rerun then restricts them. Beamer reports
% the same N for a restricted frame, so the value is stable.
\newif\ifSMFinalOnly
\newcommand{\SMFinalStateForStudents}{\global\SMFinalOnlytrue}
\makeatletter
\newwrite\SMOverlayFile
\newcommand{\SMSetOverlays}[2]{\global\@namedef{SMOverlays@#1}{#2}}
% Build files written before v5.3.0 (2026-10-01: this framework's \gotham...
% names became \SM...) still use the old names: .ovl (\gotham...SetOverlays),
% .top (...SetTopCorrection). The old names are assembled from \SMOldPrefix so a
% rerun of slidemage.sm_rename leaves them alone. The next compile rewrites
% the files with the new names.
\def\SMOldPrefix{gotham}
\expandafter\def\csname\SMOldPrefix SetOverlays\endcsname{\SMSetOverlays}
\expandafter\def\csname\SMOldPrefix SetTopCorrection\endcsname{\SMSetTopCorrection}
% A .aux written before v9.1.0 holds each subsection title (\SMAgendaSubsection)
% as `\gotham@subsectiontitleformat {Title}`: the Gotham theme wrapped titles
% in that no-op switch, which the SlideMage theme no longer has.
\expandafter\let\csname\SMOldPrefix @subsectiontitleformat\endcsname\@empty
\newcommand{\SMRecordOverlays}{%
	\ifSMInstructorVersion\else
		\immediate\write\SMOverlayFile{\string\SMSetOverlays{\the\c@framenumber}{\the\beamer@slideinframe}}%
	\fi
}
\apptocmd{\beamer@writeslidentry}{\SMRecordOverlays}{}{}
% "No file" is the line LaTeX's own \@input prints for a missing .aux or .nav: latexmk reads it and then knows the
% file as an input from the first run, so a second run that reads an unchanged one asks for no third.
\AtBeginDocument{%
	\InputIfFileExists{\jobname.ovl}{}{\typeout{No file \jobname.ovl.}}%
	\immediate\openout\SMOverlayFile=\jobname.ovl\relax
}
% Expands to <N> when the frame that is about to open is flagged and N is known.
\newcommand{\SMFinalSpec}{%
	\ifSMFinalOnly
		\ifSMInstructorVersion\else
			\ifcsname SMOverlays@\the\numexpr\c@framenumber+1\relax\endcsname
				<\csname SMOverlays@\the\numexpr\c@framenumber+1\relax\endcsname>%
			\fi
		\fi
	\fi
}
\newcommand{\SMOpenFrame}[1]{%
	\edef\SMOpenFrame@cmd{%
		\noexpand\begin{frame}%
		\ifx\SMSlideSteps\empty\SMFinalSpec\else<\SMSlideSteps>\fi
		\unexpanded{#1}}%
	\global\SMFinalOnlyfalse
	\SMOpenFrame@cmd
}
\NewDocumentCommand{\slidestep}{m O{instructor,during-course,after-course} m}{%
	\ifdeckversioninlist{\SMDeckVersion}{#2}%
		{\ifcsname SMSlideStepSeen@#1\endcsname
			\global\SMSuppressSlideTargettrue
		 \else
			\global\SMSuppressSlideTargetfalse
			\expandafter\gdef\csname SMSlideStepSeen@#1\endcsname{}%
		 \fi
		 \gdef\SMSlideSteps{#3}%
		 \global\SMSlideVersionSetfalse
		 \setSlideID{#1}%
		 \global\SMSuppressSlideTargetfalse
		 \setSlideLanguageFromSlug{#1}%
		 \input{slides/#1-current.tex}%
		 \gdef\SMSlideSteps{}}%
		{}%
}

\makeatother

% \SMUnderbrace{math}{label} - seamless TikZ-drawn brace, replacing plain
% amsmath \underbrace (2026-07-23, direct instruction, after presenting
% options on slide ebq -- see SESSION_HANDOFF.md's own entry for that task):
% amsmath assembles \underbrace's glyph from separate piecewise brace-half
% characters, which visibly shows a seam/notch at narrow widths (ebq's own
% "risk-free rate" term is narrow enough to show it plainly); this instead
% draws one continuous vector brace path via TikZ's `decorate` +
% `decoration={brace}` (decorations.pathreplacing, loaded above), which has
% no such seam at any width. #1 is typeset in math mode automatically (do
% not wrap the call itself in $...$); #2 is the label, set in upright text
% (font=\small, align=center so a multi-line label via \\ stays centered
% under the brace, e.g. two-line labels converted from \substack{}). Each
% call gets its own auto-numbered internal node name (SMUnderbraceCounter)
% so multiple calls can coexist on one slide without a caller-supplied name
% argument colliding (unlike the original \tikzbrace{}{}{} demo macro this
% replaces, which took the node name as a third argument).
\newcounter{SMUnderbraceCounter}
\newcommand{\SMUnderbrace}[2]{%
	\stepcounter{SMUnderbraceCounter}%
	\tikz[baseline=(SMUnderbraceNode\theSMUnderbraceCounter.base)]{
		\node[inner sep=0pt] (SMUnderbraceNode\theSMUnderbraceCounter) {$#1$};
		\draw[decorate,decoration={brace,amplitude=4pt,mirror,raise=1pt}]
			(SMUnderbraceNode\theSMUnderbraceCounter.south west) -- (SMUnderbraceNode\theSMUnderbraceCounter.south east)
			node[midway,below=6pt,font=\small,align=center] {#2};
	}
}

% PDF accessibility: mark purely decorative/running content (footer, below;
% also reused on titleframe's/quoteframe's own decorative background boxes,
% and -- since 2026-08-14, via the \includegraphics wrapper further down --
% on every image the alt-text manifest classifies as decorative, the
% institute logo among them) as a PDF "Artifact", so screen readers skip it
% -- via
% LuaTeX's raw \pdfextension literal
% primitive. \ifluatex-guarded (iftex is already loaded, via beamer.cls
% itself) because \pdfextension is LuaTeX-only -- this project's README
% documents xelatex as a supported alternative engine, and compiling with it
% hit exactly this: "Undefined control sequence: \pdfextension", hard-failing
% every frame, since the first version of this didn't guard the primitive at
% all. Under any non-LuaTeX engine these two commands now silently do
% nothing instead: decorative content just isn't marked as an Artifact under
% xelatex/pdflatex, rather than breaking the compile -- a real loss of
% accessibility under those engines specifically, not a full equivalent fix,
% but not something to solve today (see also: compiling with plain pdflatex
% already silently mis-renders the whole deck's fonts, so xelatex is really
% the only *other* engine this needs to not break under; pdflatex was never
% viable here to begin with, per "Compiling" in README.md).
% Deliberately NOT tagpdf/kernel tagging (\DocumentMetadata):
% beamer 3.77 refuses to compile at all when that's active ("The beamer
% class is not compatible with \DocumentMetadata ... switch to the new
% ltx-talk class if you require tagged output"), and ltx-talk is a from-
% scratch class (not beamer-based, no theme/template system) that this
% framework can't realistically migrate to. Bypassing beamer's check by
% force-loading the underlying package before \documentclass also failed --
% it gets past hyperref but then hits a genuine internal error in tagpdf's
% LuaTeX backend from the non-standard load order, not something to build
% real coursework on.
%
% /Artifact is a standalone PDF marked-content tag (PDF spec 14.8.2.2):
% valid in *any* PDF, with or without a structure tree, no /Properties
% resource entry or MCID required -- unlike structure-tree tagging, this
% doesn't require kernel tagging support at all. It's a real, working
% partial fix, not full PDF/UA conformance: there's no structure tree, so
% no true semantic headings, and reading order elsewhere on the slide still
% depends on plain content-stream order (verified, not enforced, for
% standardframe -- see its own comment).
% #1 is the artifact's inline PDF dictionary content, e.g.
% "/Type /Pagination /Subtype /Footer" for a running footer, or
% "/Type /Background" for a decorative fill -- not hardcoded to one /Type,
% since not everything decorative here is "Pagination" in the PDF spec's own
% sense of that word (running headers/footers/page numbers/watermarks);
% titleframe's/quoteframe's plain background color fills are semantically
% "Background" artifacts instead, a distinct standard type.
\newcommand{\SMArtifactBegin}[1]{%
	\ifluatex\pdfextension literal direct{/Artifact <<#1>> BDC}\fi
}
\newcommand{\SMArtifactEnd}{\ifluatex\pdfextension literal direct{EMC}\fi}

% Per-frame "Image source(s):" auto-footnote (ALT_TEXT_NOTES.md phase 6 step
% 3). The \includegraphics wrapper further down calls
% slidemage_alt.recordSource() for every image (decorative or informative --
% source is independent of role, same reasoning as media/alt-text.json's own
% ImageEntry.source), accumulating each frame's distinct, order-of-first-
% -appearance source strings in Lua. These two commands are standardframe's
% own hooks into that: reset at frame open (same per-overlay-replay reset
% \SMClearSideFootnotesBox already does), emit at frame close.
%
% Defined here (outside/before the "Alt-text manifest" \ifluatex block
% further down that actually loads SlideMage_alt.lua and redefines
% \includegraphics) rather than inside it, matching \SMArtifactBegin/End
% just above: standardframe unconditionally calls both, so they need to
% exist regardless of engine, each independently \ifluatex-guarded inside
% -- under any non-LuaTeX engine \SMResetImageSources/
% \SMEmitImageSourceFootnote are both no-ops, same degradation as the
% rest of this accessibility system under a non-LuaTeX engine. Safe to
% define before SlideMage_alt.lua is actually loaded (\AtBeginDocument,
% further down): these are only ever *called* from inside a real frame,
% well after \begin{document} has fired.
\newif\ifSMImageSourcesHasText
\newcommand{\SMImageSourcesText}{}
\newcommand{\SMResetImageSources}{%
	\ifluatex\directlua{slidemage_alt.resetFrameSources()}\fi
}
% \source{} is itself implemented as an empty-marker \footnote{} (see its
% own definition further below) -- calling it here gives the auto-footnote
% the exact same rendering (font, spacing, right-alignment, stretch-to-
% -sit-above-the-footer, side-footnote redirection under \sidefootnotes) as
% a hand-authored \source{} line, for free, rather than reimplementing any
% of that. Coexists with a real hand-authored \source{} call elsewhere in
% the same frame -- both are just separate unmarked footnote-styled lines,
% stacking the same way multiple real \footnote{}s already do.
% \SMImageSourceLabel -- the "Image source(s):"/"Bildquelle(n):" label,
% driven by \ifSMGerman (set from the compiling deck's own
% \decklanguage) rather than hardcoded, so a German deck's auto-footnote
% reads naturally instead of showing an English label on German content.
\newcommand{\SMImageSourceLabel}{\ifSMGerman Bildquelle(n)\else Image source(s)\fi}
\newcommand{\SMEmitImageSourceFootnote}{%
	\ifluatex
	\SMImageSourcesHasTextfalse
	\directlua{slidemage_alt.emitSourceFootnote()}%
	\ifSMImageSourcesHasText
		\source{\SMImageSourceLabel: \SMImageSourcesText}%
	\fi
	\fi
}

% Invisible slide-id stamp for footer-less frame types. \titleframe and
% \fullbleedframe (both defined further down) render via beamer's "plain"
% frame option, which suppresses the footline entirely -- so the visible
% footer's own "(slide id: ...)" text (see the footer template below) never
% reaches those pages. SlideMage's own compile-time preview snapshot
% (slide_snapshot.py) identifies which slide a compiled page belongs to
% purely by scanning its extracted text for that "slide id: XXX" string, so
% a plain frame with none of that text never gets a snapshot captured at
% all (caught 2026-08-15: titleframe's own title slide and 25 straight
% fullbleedframe "Excursus" slides in the IA2 deck all had no preview).
% This stamps the same marker, invisibly (PDF text-render mode 3, the same
% technique \SMAltInvisible below uses for image alt text) and wrapped
% as a /Pagination/Footer artifact, same as the real footer's own slide-id
% text -- so it reaches neither a sighted reader nor a screen reader,
% purely a technical marker for the compile pipeline. Only \slideID itself
% is stamped, not the omitted-slide/version suffix the visible
% footer also carries: nothing here is ever read by a human, and
% slide_snapshot.py's own regex only looks for the id.
%
% \vbox to 0pt{\vss ...}: same reasoning as every other absolute/zero-space
% overlay in this file (see \titleframe's own comment on this) -- \sbox
% builds an hbox, and a bare hbox placed directly in a frame's vertical
% list silently opens an implicit paragraph, switching following material
% into horizontal mode. Wrapping in a vbox contains that entirely within
% its own internal vertical mode instead of leaking out to the frame body.
\newsavebox{\SMSlideIDStampBox}
\newcommand{\SMInvisibleSlideIDStamp}{%
	\ifluatex
	\vbox to 0pt{\vss
	\begingroup
	\hbadness=10000 \hfuzz=\maxdimen
	\sbox{\SMSlideIDStampBox}{%
		\parbox[t]{\textwidth}{%
			\tiny\raggedright
			\SMArtifactBegin{/Type /Pagination /Subtype /Footer}%
			\pdfextension literal direct{3 Tr}%
			% Trailing \space before the closing paren, not just "\slideID)":
			% slide_snapshot.py's own regex ("slide id:\s*(\S+)") reads
			% everything up to the next whitespace as the id, and without a
			% space there it would capture "gcs)" instead of "gcs" (caught
			% by testing the actual compiled+extracted text, not assumed --
			% the real footer's own "... \slideID ... | ..." format always
			% has a space after \slideID for the same reason).
			(slide id: \slideID\space)%
			\pdfextension literal direct{0 Tr}%
			\SMArtifactEnd
		}%
	}%
	\wd\SMSlideIDStampBox=0pt
	\ht\SMSlideIDStampBox=0pt
	\dp\SMSlideIDStampBox=0pt
	\usebox{\SMSlideIDStampBox}%
	\endgroup
	}%
	\fi
}

% ---------------------------------------------------------------------------
% Alt-text manifest: automatic /Artifact marking of decorative images
% ---------------------------------------------------------------------------
% See ALT_TEXT_NOTES.md for the whole design; this is phase 1 of it. Every
% image in the project has a row in media/alt-text.csv saying whether it is
% "decorative" (carries no information beyond the slide text -- Midjourney
% illustrations, icons, chapter logos, institute logos) or "informative" (a
% chart/screenshot/photo a reader actually needs described). Decorative ones
% get wrapped in /Artifact here, so a screen reader skips them instead of
% announcing a bare filename. Informative ones are left alone for now; their
% descriptions are authored per slide and emitted later (phases 4 and 5).
%
% Why wrap \includegraphics rather than change the 265 call sites: the call
% sites don't need to change at all this way, including the ones inside tikz
% \node{}s and the ones that pass the filename through a macro argument
% (\SMInstituteLogoPath, or a #1 inside another macro's definition) --
% those only know their filename at expansion time, which a source-level
% rewrite couldn't have handled anyway. An image with no manifest row simply
% renders exactly as before, so a missing or unreadable CSV degrades to
% today's behavior rather than breaking every image in the deck.
%
% LuaTeX-only, same \ifluatex reasoning as \SMArtifactBegin above -- and
% here the guard covers the whole redefinition, not just the marking, since
% the lookup itself is Lua. Under any other engine \includegraphics is left
% completely untouched.
\ifluatex
	% The Lua itself lives in SlideMage_alt.lua rather than inline here --
	% see that file's own header for why (short version: TeX tokenizes a
	% \directlua argument first, so every "%" in a Lua string pattern
	% silently comments out the rest of its line, and the first version of
	% this failed exactly that way). This loader snippet is deliberately
	% free of "%" and "#" so it is safe to keep inline. Found via kpathsea
	% (TEXINPUTS), since it lives in latex/, outside the project.
	\SMLuaOnce{
		local path = kpse.find_file("SlideMage_alt.lua", "lua")
			or kpse.find_file("SlideMage_alt.lua", "tex")
		if path then
			dofile(path)
		end
	}

	\newif\ifSMAltDecorative
	\newif\ifSMAltHasText
	\newsavebox{\SMAltTextBox}
	\newcommand{\SMAltShort}{}

	% Invisible alt text (ALT_TEXT_NOTES.md phase 5). An untagged PDF is read
	% by *text extraction*, so text placed in the content stream at the image's
	% position is what a screen reader encounters there -- text rendering mode
	% 3 ("neither fill nor stroke", PDF spec 9.3.6) makes it invisible on the
	% page while leaving it fully extractable.
	%
	% Two deliberate properties:
	%
	% - The box's three dimensions are zeroed after typesetting, so the text
	%   contributes nothing whatsoever to layout. Setting \wd/\ht/\dp of an
	%   already-built box is the only way to be sure of that: \rlap/\smash
	%   each zero only some of them, and a stray descender would shift a
	%   hand-positioned slide.
	% - \SMAltImagePrefix ("Image: "/"Bild: ") prefixes the description so a listener can tell it apart
	%   from the slide's own prose, which they otherwise cannot.
	%
	% Known and accepted side effect: this text is real text, so it appears in
	% copy-paste and in Ctrl-F. That was the explicit tradeoff for reaching a
	% screen reader at all without a structure tree, which beamer cannot give
	% us (see README.md's accessibility section).
	% \SMAltImagePrefix -- language-conditional counterpart to
	% \SMImageSourceLabel above, for this invisible layer's own "Image: "/
	% "Bild: " prefix.
	\newcommand{\SMAltImagePrefix}{\ifSMGerman Bild: \else Image: \fi}
	\newcommand{\SMAltInvisible}[1]{%
		\begingroup
		% A long alt string in an unbreakable hbox would otherwise flood the
		% log with overfull-box warnings on every image.
		\hbadness=10000 \hfuzz=\maxdimen
		% Wrapped in a \parbox, and set \tiny, so the text stays *inside the
		% page*. Neither is cosmetic: an unbreakable line starting at the
		% image's own x-position runs past the right page edge after about 80
		% characters, and every PDF text extractor tested (pdftotext with and
		% without -layout, PyMuPDF) silently drops whatever falls outside the
		% page rectangle -- so a longer alt was being truncated mid-sentence
		% with no error anywhere. \tiny is free here (the text is invisible)
		% and fits several times more characters per line, which keeps even a
		% long description to one or two lines.
		\sbox{\SMAltTextBox}{%
			\parbox[t]{\textwidth}{%
				\tiny\raggedright
				\pdfextension literal direct{3 Tr}%
				\SMAltImagePrefix#1%
				\pdfextension literal direct{0 Tr}%
			}%
		}%
		\wd\SMAltTextBox=0pt
		\ht\SMAltTextBox=0pt
		\dp\SMAltTextBox=0pt
		\usebox{\SMAltTextBox}%
		\endgroup
	}

	% Deferred to \AtBeginDocument so it cannot possibly run before beamer has
	% loaded graphicx, and so a deck that loads further graphics packages of
	% its own still ends up with this wrapper outermost.
	\AtBeginDocument{%
		\directlua{
			local loaded = slidemage_alt.load()
			if loaded == nil then
				texio.write_nl("slidemage_alt: media/alt-text.csv not found -- "
					.. "decorative images will not be marked as artifacts.")
			end
		}%
		% \NewCommandCopy (LaTeX 2020-10-01 and later) rather than \let: it
		% copies a robust command's full definition correctly, which \let
		% does not for commands whose top level is a \protected wrapper.
		\NewCommandCopy{\SMOrigIncludegraphics}{\includegraphics}%
		% Signature "s o o m": the starred form and the legacy two-optional
		% bounding-box form aren't used anywhere in this project today
		% (checked, 2026-08-14: zero occurrences of either), but graphicx
		% accepts them, so the wrapper accepts and forwards them rather than
		% silently breaking the day one gets used. "o" (not "O{}") because
		% only -NoValue- distinguishes an absent optional argument from an
		% empty one, and graphicx reads [a][b] and [a] as different things.
		\RenewDocumentCommand{\includegraphics}{s o o m}{%
			\begingroup
			% \edef so a filename reaching this through a macro
			% (\SMInstituteLogoPath, or a #1 in another macro's body) is
			% resolved to the actual path before the lookup. Every path-valued
			% macro in this project is a plain expandable string; a fragile one
			% would need \protected@edef here instead.
			\edef\SMAltReference{#4}%
			\SMAltDecorativefalse
			\directlua{slidemage_alt.mark("\luaescapestring{\SMAltReference}")}%
			% Unconditional -- unlike SMAltInvisible below (informative
			% images only), a source is recorded regardless of role.
			\directlua{slidemage_alt.recordSource("\luaescapestring{\SMAltReference}", "\luaescapestring{\SMDeckLanguage}")}%
			% /Type /Layout is the PDF spec's artifact type for cosmetic design
			% elements, which is what a decorative illustration is. Not
			% /Background: that one means a fill or image behind the whole
			% page, and is already used for exactly that by titleframe and
			% quoteframe above.
			\ifSMAltDecorative\SMArtifactBegin{/Type /Layout}\fi
			% Emitted *before* the image, so the description precedes it in the
			% content stream and therefore in extraction order -- a listener
			% hears what the image is before reaching it, rather than after.
			% slidemage_alt.altText prints nothing for a decorative image or one
			% with no short alt yet, so this line is inert until the manifest
			% is filled in.
			\SMAltHasTextfalse
			\ifSMAltDecorative\else
				\directlua{slidemage_alt.altText("\luaescapestring{\SMAltReference}", "\luaescapestring{\SMDeckLanguage}")}%
				\ifSMAltHasText\SMAltInvisible{\SMAltShort}\fi
			\fi
			\IfBooleanTF{#1}%
				{\IfNoValueTF{#2}%
					{\SMOrigIncludegraphics*{#4}}%
					{\IfNoValueTF{#3}%
						{\SMOrigIncludegraphics*[#2]{#4}}%
						{\SMOrigIncludegraphics*[#2][#3]{#4}}}}%
				{\IfNoValueTF{#2}%
					{\SMOrigIncludegraphics{#4}}%
					{\IfNoValueTF{#3}%
						{\SMOrigIncludegraphics[#2]{#4}}%
						{\SMOrigIncludegraphics[#2][#3]{#4}}}}%
			\ifSMAltDecorative\SMArtifactEnd\fi
			\endgroup
		}%
	}
\fi

% \tikzillustration{ID}: extends the alt-text system above -- which only
% ever sees a real \includegraphics file -- to a hand-drawn tikz diagram
% with no file behind it at all (ALT_TEXT_NOTES.md). #2 is the
% illustration's own opaque, SlideMage-assigned ID (never hand-invented
% -- same "permanent ID" convention slides themselves use, so a description
% survives edits to the drawing itself); #1 (optional) is passed straight
% through as the inner tikzpicture's own options. The body between
% \begin/\end is the tikzpicture's content, completely unmodified.
%
% Role (decorative/informative) is looked up exactly the way an image's is
% -- slidemage_alt.mark(), the same Lua call \includegraphics's own wrapper
% above makes, against the same media/alt-text.csv -- rather than a second
% bracket argument SlideMage would otherwise have to text-edit into this
% call in place every time the role changes. SlideMage writes this
% illustration's own row into media/alt-text.json (and its compile-time
% media/alt-text.csv projection) keyed "tikz:ID", never a real media/ path;
% slidemage_alt.mark() itself needed no changes, since it already keys purely
% on whatever string it is given (SlideMage/alt_images.py).
%
% A decorative illustration is wrapped in the same \SMArtifactBegin{/Type
% /Background}/\End pair sideimageframe's/fullbleedframe's own decorative
% filler content already uses, so a screen reader skips it entirely; an
% informative one is left completely unmarked, same as real content in
% figureframe/fullbleedframe. Unlike an informative *image*, it gets no
% invisible-text alt layer (phase 5) of its own -- that layer's positioning
% is built around \includegraphics's own bounding box, which a freeform
% tikzpicture doesn't have the same way. The per-deck accessible transcript
% (README.md) is this illustration's own route to a screen reader instead,
% the same status phase 5 already carries for every image.
%
% \ifluatex-guarded exactly like \SMArtifactBegin/\SMArtifactEnd
% above (both \SMAltDecorativefalse/\ifSMAltDecorative only exist
% inside that same guard) -- under any other engine this degrades to a
% plain, unmarked tikzpicture, the same "no accessibility marking, but
% still compiles and renders" degradation the rest of this system already
% has. Safe even though those two control sequences are otherwise undefined
% under a non-LuaTeX engine: TeX's own \iffalse-branch skipping never looks
% up a token's meaning, so an unrecognized name inside the untaken
% \ifluatex branch is simply skipped over, never executed.
\newenvironment{tikzillustration}[2][]{%
	\ifluatex
		\SMAltDecorativefalse
		\directlua{slidemage_alt.mark("\luaescapestring{tikz:#2}")}%
		\ifSMAltDecorative\SMArtifactBegin{/Type /Background}\fi
	\fi
	\begin{tikzpicture}[#1]
}{%
	\end{tikzpicture}%
	\ifluatex\ifSMAltDecorative\SMArtifactEnd\fi\fi
}

% Footer
% The footer box on the footline, the footnote template with its symbols, \sidefootnotes and the
% strip they are drawn in are the theme's (beamerthemeSlideMage.sty); the footer text, which
% depends on the deck version, is the "frame footer" template below.
% Smaller footer + a bit of padding under it
\setbeamerfont{frame footer}{size=\SMFooterTextSize} % the catalog: {font size}{baseline skip} in pt

% Shrinks the gap between footnotes and the footer box below them. The
% theme positions the footer's content via \vspace*{\SMFootlineVOffset}
% (a fixed, negative length) before it starts; making that more negative
% turns out to give footnotes (which are placed just before the footer, and
% stretch to fill whatever room remains via beamer's own
% \beamer@framebottomskip) more room to extend into, pulling them down
% towards the footer -- checked empirically against the PDF's own
% coordinates, this does *not* move the footer's own text at all (that
% still sits exactly where \SMFooterBoxHeight below was calibrated for),
% so the footer box itself needs no compensating adjustment.
%
% Partly given back (+1.4pt) 2026-07-05: a footnote/source containing an
% underlined \href (see the hyperlink-styling comment above) sits close
% enough to the footer at the original -1.3pt that the underline visually
% touches the footer box (checked directly against the compiled PDF's own
% pixel coordinates: a plain-text footnote's own gap to the footer measured
% ~1.9pt, but an underlined-link footnote's gap measured only ~0.6pt at
% the same setting, since the strut rebuilt for footnotes just below --
% correctly sized for plain text's own natural depth -- doesn't account
% for the underline rule extending further down). +1.4pt splits the
% difference: a linked footnote's gap now measures close to the plain-
% footnote reference (~1.9pt) instead of colliding with the footer, and a
% plain footnote's own gap just becomes a little more relaxed (~3.3pt)
% rather than stretched uncomfortably tight -- applied globally rather
% than only to footnotes with links, since there's no clean per-footnote
% hook for it and the plain-footnote side effect is minor.
%
% Nudged back down by 0.4pt the next day (direct instruction: move
% footnotes down "by the width of the underline line") -- the +1.4pt above
% left a bit more headroom than necessary; 0.4pt is the underline rule's
% own measured stroke thickness (a solid ~1px core at 600dpi, i.e.
% 1/600in \approx 0.12pt, plus antialiasing bleed either side consistent
% with the common 0.4pt LaTeX default rule width -- measured directly
% against the compiled PDF's own pixel values, not assumed).
%
% Reverted the very next round (direct instruction: move back *up* by that
% same underline thickness) -- net effect is back to the plain +1.4pt this
% section started with; kept as two separate historical adjustments below
% rather than collapsing them into one, so this back-and-forth stays
% visible rather than silently disappearing.
% (v9.4.0: the adjustments are one value now, the default of the catalog's
% \SMFootlineVOffset in SlideMage_setup.tex: -0.5ex-1.3pt+1.4pt-0.4pt+0.4pt-0.87pt
% = -0.5ex-0.77pt.)

% 2026-07-06 (direct instruction): move footnotes down by half the
% measured distance between a linked footnote's own underline and the
% footer box, checked on slides/3ic-Corporate-Venture-Capital-
% current.tex specifically (the underline sits at y=245.95pt, the footer
% box's own top edge at y=247.69pt, a 1.74pt gap) -- half of that is
% 0.87pt, applied as a further negative adjustment (more negative =
% footnotes get pulled further down, per this section's own comment
% above).

% \SMFooterTextDrop (v9.8.0: derived; it was a number of the layout profile since v9.4.0, -0.49pt for
% Fira Sans, -0.37pt for Latin Modern Sans) centres the ink of the frame number in the footer box. The
% template below starts the footer line with \vspace*{drop}, which in the footline box's horizontal
% mode is a \vadjust below the line, as is the \vspace*{0.4ex} that ends it; the box has the declared
% depth \SMFootlineDepth, so the number's baseline lies
%     \SMFootlineDepth + (the depth of the line) + 0.4ex + drop
% above the page's bottom edge. The depth of the line is that of the slide id's strut
% (\SMFooterIDStrut makes it the same on every slide). The ink of the digits is what luaotfload
% keeps as the height and depth of a digit's box (the glyph's bounding box: 676 and 19 units in Fira
% Sans Light), so centring it in the footer box of height \SMFooterBoxHeight gives the drop.
\makeatletter
\newlength{\SM@footerlinedepth}
\newcommand{\SM@footertextdrop}{%
	\setbox\z@\hbox{\SMFooterIDSize\SMFooterIDStrut}\setlength{\SM@footerlinedepth}{\dp\z@}%
	\setbox\z@\hbox{0123456789}%
	\setlength{\SMFooterTextDrop}{\dimexpr(\SMFooterBoxHeight-\ht\z@+\dp\z@)/2-\SMFootlineDepth-\SM@footerlinedepth-0.4ex\relax}}
\makeatother
%
% \SMFooterRightShift (v9.7.0: derived; it was a number of the layout profile, 4.6pt for Fira Sans,
% -0.33pt for Latin Modern Sans) puts the frame number's right edge on the logo's right edge above
% it, as the footnotes' lines end there. The footer line is typeset in the footline box, which is
% \hsize wide and ends its lines \rightskip before that; after the number come, in this order,
% \hspace*{-4em} and the shift (set here), then the box's \SMRightFooterPadding and its
% \SMFootlineHRightOffset. So the number's right edge is
%     \hsize - \rightskip - \SMRightFooterPadding - \SMFootlineHRightOffset + 4em + shift
% and the shift that puts it at \paperwidth - \SMFrametitleRightPadding follows. Every term is the
% one in force where the footer is typeset (the \rightskip is 3ex of the footline font, the 4em is
% the footer's), which is why the shift is computed in the frame footer template and not once: the
% two fonts differ in size, and a profile with another x-height has other paddings in ex.

% The footer no longer shows a compile date/time (removed in v8.25.0: a
% changing stamp made every wrapper and PDF differ between compiles). The
% two names below remain as harmless no-ops so decks that still call
% \SMShowCompileTimestampfalse/true or print \SMCompileTimestamp compile.
\newif\ifSMShowCompileTimestamp
\newcommand{\SMCompileTimestamp}{}

% \SMPresenterModeOverride: SlideMage's own hook for a "Presenter"
% compile -- a separate output alongside the plain instructor PDF, never a
% toggle on it (see SlideMage/compiler.py's own CompileJob.presenter
% field and SlideMage/qt_deck_table.py's own "Presenter" checkbox,
% independent of the three version checkboxes the same way "After-course,
% per chapter" already is). Same \ifdefined-hook pattern as
% \deckversionoverride: compiler.py's wrapper file
% \def's it, before \input{}ing the real deck, whenever that checkbox is
% ticked for a job.
%
% Turns on beamer's own built-in dual-screen presenter mechanism:
% \note{...} content (see "Speaker notes" below -- already authored
% throughout this project's own slides, previously never actually emitted
% anywhere at all) becomes a facing "notes page" after every real slide,
% which presenter software (pympress is this project's own recommended
% choice -- see README.md's "Presenting" section) auto-detects and splits
% across the projector (plain slides only) and the presenter's own screen
% (current slide, next slide, and the note text). \usepackage{pgfpages} is
% loaded explicitly rather than trusted to \setbeameroption's own
% auto-load, since which beamer versions actually do that on their own
% isn't worth depending on.
%
% A presenter compile's own CompileJob always carries version="instructor"
% (or the deck's own hardcoded content, for a deck with no
% \deckversionoverride at all) regardless of whether the plain "Instructor"
% checkbox is *also* ticked in the same batch -- notes and full \pause
% reveals only make sense against the instructor's own content, so
% Presenter never needs (or offers) During-course/After-course versions of
% itself. slide_snapshot.py's own compile-time preview refresh is
% deliberately skipped for a presenter job (compiler.py's own
% _relocate_pdf_and_finish()) -- its page count and per-page layout are
% nothing like a normal compile's, so reading it the normal way would
% silently corrupt the real preview cache with wrong or missing pages.
\ifdefined\SMPresenterModeOverride
	\usepackage{pgfpages}
	\setbeameroption{show notes on second screen=right}
	% Presenter-mode note page (direct instruction, 2026-09-10): beamer's
	% own default "note page" template (beamerouterthemedefault.sty,
	% \defbeamertemplate*{note page}{default}) reserves its top quarter for
	% a thumbnail of the current slide (\insertslideintonotes) plus a
	% rotated compile-date printed beside it, then gives the title band
	% (short title + section/subsection/frametitle breadcrumb) the same
	% 0.25\paperheight, leaving only the bottom 0.75 for \insertnote. Both
	% the thumbnail and date are redundant here -- pympress (this project's
	% own recommended presenter software, see README.md) already shows the
	% current and next slide on its own -- so this redefinition drops that
	% whole top block (thumbnail + date) outright and shrinks the title
	% band from 0.25\paperheight to 0.08\paperheight (still one line for
	% the short title plus room for the breadcrumb, per beamer's own
	% layout below, copied unchanged apart from widening each line from
	% beamer's own 0.71-0.75\textwidth, sized to leave room for the now-
	% removed thumbnail, to 0.85-0.95\textwidth), giving \insertnote the
	% remaining ~0.92\paperheight instead of 0.75. No color-theme "note
	% title"/"note page" background is set anywhere in this project (a
	% plain grep confirms it), so beamer's own \ifbeamercolorempty guards
	% around both of those bands are dropped too -- they were already
	% pure no-ops.
	%
	% Registered as a new named variant ("SMPresenter"), not a
	% redefinition of beamer's own "default" -- confirmed directly that
	% \defbeamertemplate*{note page}{default}{...} here errors ("Command
	% \beamer@@tmpop@note page@default already defined"): both the
	% starred and unstarred \defbeamertemplate forms register a variant
	% name via a plain \newcommand internally (beamerbasetemplates.sty),
	% which always errors on a name beamer's own outer theme already
	% registered, star or not. Same pattern this file's own footnote
	% template already uses just above (\defbeamertemplate{footnote}
	% {SMFootnote}{...} then \setbeamertemplate{footnote}[SMFootnote]) --
	% define-then-select instead of overwrite-in-place.
	%
	% \makeatletter/\makeatother wraps this specific block (not the whole
	% \ifdefined) because \beamer@tempbox below is an @-command: this
	% point in the preamble sits between two of this file's own
	% \makeatletter/\makeatother pairs (the previous one closes well
	% above, the next opens well below), so @ is still catcode-other here
	% otherwise -- confirmed directly rather than assumed.
	\makeatletter
	\defbeamertemplate{note page}{SMPresenter}{%
		{%
			\scriptsize
			% The breadcrumb typesets the frame title (twice: once to measure, once
			% to show it), and a \footnote{}/\source{} in a frame title would run each
			% time: the slide's last reveal step then carried its footnote three
			% times (direct report 2026-10-13, mbz). Footnotes are switched off in
			% this group; the slide itself keeps its own, single one.
			\renewcommand*{\footnote}[2][]{}%
			\renewcommand*{\footnotemark}[1][]{}%
			\usebeamerfont{note title}\usebeamercolor[fg]{note title}%
			% Natural-height \vbox (no "to <len>"/\vfil): this project's own
			% real decks call \section{}/subsection ones sometimes not (a
			% plain grep of decks/*.tex confirms both), so a fixed-height
			% band sized for the *shortest* case overlaps \insertnote below
			% on any slide needing all three breadcrumb lines -- confirmed
			% directly: a first attempt at a fixed 0.08\paperheight did
			% exactly that against a real slide's own \note{}. A natural-
			% height box is exactly as tall as whichever of these three
			% lines are non-empty for *this* slide, never more, so it can
			% never collide with the notes below it, and hands \insertnote
			% every bit of space a section/subsection-free slide doesn't
			% need for a breadcrumb. \insertshorttitle (beamer's own first
			% line here, normally the deck's own \title{}) is dropped
			% outright rather than shrunk -- confirmed via grep that
			% \title{} is never called anywhere in this project, so that
			% line has always been silently empty.
			\vbox{\vskip0.3em
				\setbox\beamer@tempbox=\hbox{\insertsection}%
				\hbox{\ifdim\wd\beamer@tempbox>1pt{\hskip4pt\raise3pt\hbox{\vrule
							width0.4pt height7pt\vrule width 9pt
							height0.4pt}}\hskip1pt\hbox{\begin{minipage}[t]{0.95\textwidth}\def\breakhere{}\insertsection\end{minipage}}\fi%
				}%
				\setbox\beamer@tempbox=\hbox{\insertsubsection}%
				\hbox{\ifdim\wd\beamer@tempbox>1pt{\hskip17.4pt\raise3pt\hbox{\vrule
							width0.4pt height7pt\vrule width 9pt
							height0.4pt}}\hskip1pt\hbox{\begin{minipage}[t]{0.95\textwidth}\def\breakhere{}\insertsubsection\end{minipage}}\fi%
				}%
				\setbox\beamer@tempbox=\hbox{\insertshortframetitle}%
				\hbox{\ifdim\wd\beamer@tempbox>1pt{\hskip30.8pt\raise3pt\hbox{\vrule
							width0.4pt height7pt\vrule width 9pt
							height0.4pt}}\hskip1pt\hbox{\insertshortframetitle[width=0.9\textwidth]}\fi%
				}%
				\vskip0.3em}%
		}%
		\nointerlineskip
		\insertnote
	}
	\setbeamertemplate{note page}[SMPresenter]
	\makeatother
\fi
% \SMHideFooterText: opt-in flag suppressing just \FooterText (the left
% part of the footer) -- default false. Set true by the standout template's
% own hook below, so `\begin{frame}[standout]` hides only that left part,
% keeping the "(slide id: ...)" block and frame number visible (2026-07-16,
% direct instruction). Reset immediately after being read, inside the
% footer template itself, rather than relying on a per-environment closing
% hook (standout frames are bare `\begin{frame}[standout]...\end{frame}`,
% not one of this project's own environments like standardframe, so there's
% no custom closing code to reset it in) -- safe because \SMHideFooterText
% is only ever consumed here, once, immediately before being cleared.
\newif\ifSMHideFooterText
% The theme's own \defbeamertemplate*{standout}{SMDefault}
% (beamerthemeSlideMage.sty) is what `\begin{frame}[standout]` actually
% activates; \addtobeamertemplate's "after" hook appends to it without
% needing to touch the theme.
\addtobeamertemplate{standout}{}{\global\SMHideFooterTexttrue}

% \SMSlideLanguageFooterTag -- resolved once here rather than inlined
% into the footer template below, same reason \SMDeckVersionText is
% precomputed rather than built inline in \deckversion. Only ever appears
% appended directly onto the version digits ("v1.1.0EN"), so it's a no-op
% whenever \ifSMSlideVersionSet is false -- ties the language tag's
% rollout to the same version-bump touchpoint slides already go through,
% instead of needing a one-time pass over every slide. For an "mult"
% slide (\multilingualtext above) this shows whichever language the deck
% itself is compiling as, not a fixed value -- \MakeUppercase{en} = "EN".
\newcommand{\SMSlideLanguageFooterTag}{%
	\ifSMSlideLanguageSet
		\ifSMSlideLanguageMult\MakeUppercase{\SMDeckLanguage}\else\slideLanguage\fi
	\fi
}

% \SMFooterIDStrut -- fixed-size strut for the footer's "(slide id: ...)"
% block (2026-09-19, direct instruction: footer text moved by ~0.27pt between
% slides). Cause: the footline's beamercolorbox only has a *minimum* height and
% depth, and the id block is the only footer text that differs from slide to
% slide -- an id containing a descender letter (g, j, p, q, y: "egk", "y89")
% is deeper than the box's minimum depth and lifts the whole footer, an id
% without one ("kaz", "c2t") leaves it at the minimum. The strut spans the
% largest ascender and descender the block can ever contain (measured with the
% block's own font size, at use time, so it follows a font change), which makes
% the footline's height and depth identical for every slide id. Must be
% called after the block's \fontsize.
\newsavebox{\SMFooterIDMeasureBox}
\newcommand{\SMFooterIDStrut}{%
	\sbox{\SMFooterIDMeasureBox}{gjpqy()[]|bdfhklt0123456789}%
	\rule[-\dp\SMFooterIDMeasureBox]{0pt}{\ht\SMFooterIDMeasureBox}%
}

% \SMFooterHide: which footer elements a deck leaves out, a comma list of
% text (the left \FooterText), id ("slide id: ... | v..."), deck ("instructor
% deck" etc.). Empty by default = everything shown. SlideMage's deck
% settings write it as \def\SMFooterHide{id,deck} in the deck's preamble;
% it is read once at \begin{document}, so the deck may set it before or
% after \input{SlideMage_template_...}. The frame number is always shown.
\providecommand{\SMFooterHide}{}
\newif\ifSMFooterShowText
\newif\ifSMFooterShowID
\newif\ifSMFooterShowDeck
\newcommand{\SMFooterHideItem}[1]{%
	\ifstrequal{#1}{text}{\SMFooterShowTextfalse}{}%
	\ifstrequal{#1}{id}{\SMFooterShowIDfalse}{}%
	\ifstrequal{#1}{deck}{\SMFooterShowDeckfalse}{}%
}
\AtBeginDocument{%
	\SMFooterShowTexttrue\SMFooterShowIDtrue\SMFooterShowDecktrue
	\expandafter\forcsvlist\expandafter{\expandafter\SMFooterHideItem\expandafter}\expandafter{\SMFooterHide}%
}

\newif\ifSMNavigationBar % \SMNavigationBar, see "Navigation for readers" at the end
\setbeamertemplate{frame footer}{%
	\csname SM@footertextdrop\endcsname%
	\vspace*{\SMFooterTextDrop}%
	\ifSMHideFooterText
		\global\SMHideFooterTextfalse
	\else
		\ifSMFooterShowText
			\FooterText\ifSMNavigationBar\SMNavigationLinks\fi
		\fi
	\fi
	\hfill%
	{\SMFooterIDSize \SMFooterIDStrut
		\ifSMFooterShowID(slide id: \slideID\ifSMOmittedSlide\SMOmittedSlideMarker\fi\ifSMSlideVersionSet\ | v\slideVersion\SMSlideLanguageFooterTag\fi\ifSMFooterShowDeck\ | \SMDeckVersionText\fi)%
		\else\ifSMFooterShowDeck(\SMDeckVersionText)\fi\fi}\hspace*{2em}\insertframenumber%
	\hspace*{-4em}%
	\setlength{\SMFooterRightShift}{\dimexpr\paperwidth-\SMFrametitleRightPadding-\hsize+\rightskip+\SMRightFooterPadding+\SMFootlineHRightOffset-4em\relax}%
	\hspace*{-\SMFooterRightShift}%
	\vspace*{0.4ex}%
}




% \source{text} -- for a "Source: ..."/"Image source: ..."/etc. citation line
% (common at the bottom/edge of a converted PowerPoint slide) that should look
% and sit *exactly* like the footnotes above -- same font, spacing,
% right-alignment, and the same stretch-to-sit-just-above-the-footer
% behaviour beamer already gives real footnotes -- but with no marker, neither
% inline in the body text nor next to the note itself, since a source line
% isn't anchored to one specific word. Implemented as a real \footnote{} (so
% it gets that positioning for free, rather than reimplementing it) with
% \thefootnote emptied locally; since \footnote{} still advances the footnote
% counter even with an empty marker, \addtocounter reverses that so a real
% \footnote{} elsewhere on the same slide doesn't skip a symbol in the
% \fnsymbol sequence. Unlike a real footnote, no label is added automatically
% here -- the label varies ("Source:", "Image source:", ...), so write it as
% part of #1.
\newcommand{\source}[1]{%
	\begingroup
	\renewcommand{\thefootnote}{}%
	\footnote{#1}%
	\endgroup
	\addtocounter{footnote}{-1}%
}

% A plain tabular with yellow \toprule/\midrule/\bottomrule, matching the
% frametitle/footer accent color. Takes the same column-spec argument as
% tabular itself, e.g. \begin{SMTable}{lccc} ... \end{SMTable}. Only
% booktabs rules are colored (via \arrayrulecolor, scoped to this table by
% \begingroup/\endgroup) -- if a table ever needs a plain black \hline too,
% that's unaffected since \arrayrulecolor only recolors rules drawn *after*
% it takes effect, and booktabs' rules read it fresh each time.
%
% RULE-BLEED WARNING (2026-07-12, direct instruction: "make sure the table
% lines do not extend beyond the text"), read this before writing a new
% SMTable: a plain column spec like `{lccc}` above leaves the default
% \tabcolsep (6pt, or whatever a slide-local \setlength last set it to) as
% *implicit* padding on the table's own outermost edges -- invisible to the
% table's own text (which starts/ends flush with each column's actual
% content), but very visible to booktabs' \toprule/\midrule/\bottomrule and
% to colortbl's \rowcolor/\cellcolor, both of which span the *declared*
% column width including that implicit padding. The result, confirmed
% directly and independently on at least three slides so far
% (slides/u5h-Long-term-Bonds-Provisions-current.tex's own v1-1-0/v1-4-0,
% slides/c4c-Macbeth-Spot-Removers-Proposed-Structure-current.tex's own
% v1-2-2, slides/tw5-Exercise-Leverage-and-ROE-Solution-current.tex's own
% v1-4-0): every rule/highlight bleeds \tabcolsep past the visible text on
% both the left and right edges, most noticeable on a left-aligned
% (non-\centering) table where the mismatch is easy to see against the
% title's own left margin.
%
% (2026-09-30: write new tables as \begin{SMTable}[gap=12pt]{...}
% instead -- see "FLUSH EDGES BUILT IN" below; the rest of this comment
% describes the hand-written form that option now produces.)
% Not fixed here in the environment definition itself: doing so would
% change every existing SMTable's own column spacing out from under it
% (this environment has no way to tell "plain default spacing was intended"
% apart from "the bleed was never noticed"). Fix per table instead, the
% same three-part pattern every slide above uses -- worth reaching for by
% default on any *new* SMTable, not just when the bleed is spotted:
%   {\setlength{\tabcolsep}{0pt}
%   \begin{SMTable}{@{}l<{\hspace{<gap>pt}}c<{\hspace{<gap>pt}}c@{}}
%   ...
%   \end{SMTable}}
% -- \tabcolsep zeroed (scoped locally, e.g. by wrapping the whole call in
% a `{...}` group, so it doesn't leak into later tables), explicit `@{}` at
% both outer edges of the column spec (so booktabs/colortbl have no
% implicit padding left to bleed into), and array's `<{\hspace{<gap>pt}}>`
% column modifier appended after every column but the last, to put back an
% explicit, *content-column* gap (matching whatever `2*\tabcolsep` the
% table used to have, if preserving its exact look, e.g. `<{\hspace{12pt}}>`
% for the ordinary 6pt-\tabcolsep default) -- `<{...}>` puts the gap
% *inside* the preceding column's own content, which `\rowcolor` (if used)
% colors along with everything else, unlike a plain `@{...}` separator,
% which `\rowcolor` deliberately leaves untinted (see u5h's own v1-6-0
% comment for that exact failure mode if a highlighted table ever needs
% this too).
%
% RE-CONFIRMED 2026-07-16, direct instruction, after this exact bleed
% reappeared on essentially every SMTable in a freshly-converted
% chapter (Capital Budgeting Techniques) despite the warning above already
% existing at the time -- the rule above is not optional/only-if-noticed,
% it's the default way to write *any* new SMTable, full stop. Check
% every new SMTable against it before considering a slide done, not
% just when a rendered page happens to make the bleed visually obvious.
%
% BLOCK-RULE-COLOR CONVENTION (2026-07-16, direct instruction): a
% SMTable placed *inside* any block environment (`infoblock`,
% `exampleblock`, `referenceblock`, `definitionblock`, plain `block`, or
% their `SMFloating*Block` wrappers) should use black rules, not this
% environment's own default yellow -- yellow-on-the-block's-own-tinted-
% background reads poorly, unlike yellow-on-plain-white outside a block.
% Override locally, *inside* `\begin{SMTable}{...}`, before `\toprule`
% (not before `\begin{SMTable}` itself -- the environment's own
% opening code sets yellow *after* its `\begingroup`, so a color set
% before that point is simply overwritten; confirmed directly, first
% attempt this same day):
%   \begin{SMTable}{...}
%   	\arrayrulecolor{c_primary_dark}
%   	\toprule
%   	...
%   \end{SMTable}
% Booktabs rules read `\arrayrulecolor` fresh at each rule, and
% `\endgroup` (SMTable's own closing code) resets it back to yellow
% afterward, so this never leaks into a later table. First applied on
% slides/6ux-Example-51-BMA-Payback-Table-current.tex's own
% `SMFloatingInfoBlock`.
%
% FLUSH EDGES BUILT IN (2026-09-30, direct instruction: "propose a more
% systematic solution" for per-slide fixes -- 53 tables carried the
% three-part pattern above by hand): give the column gap as an option and
% write the plain column spec,
%   \begin{SMTable}[gap=12pt]{lccc} ... \end{SMTable}
% which produces exactly the hand-written form -- \tabcolsep 0pt, @{} at
% both edges, <{\hspace{12pt}} after every column that is followed by
% another one (not before an explicit @{...}/!{...} separator), *{n}{...}
% expanded -- so a table rewritten this way renders identically. After a
% fixed-width p/m/b/w column the gap is an @{\hspace{...}} separator
% instead: array places a <{...} suffix inside such a column's box, where
% it only pads the end of the cell text and leaves the columns touching
% (the hand-written pattern had exactly that flaw for p columns). The
% expansion is done by SlideMage_measure.lua (slidemage_measure.gap_spec).
% gap=12pt is the ordinary look (2 x beamer's 6pt \tabcolsep). Without the
% option the environment behaves as it always has (plain tabular with
% yellow rules), so every existing table is untouched.
\makeatletter
\pgfqkeys{/SMTable}{
	gap/.store in=\SMTable@gap,
}
\ifluatex
	\SMLuaOnce{
		local path = kpse.find_file("SlideMage_measure.lua", "lua")
			or kpse.find_file("SlideMage_measure.lua", "tex")
		if path then
			dofile(path)
		end
	}
	% The two numbers the Top Margin measurement needs. The ascender is read from the text font in
	% force (luaotfload's OpenType metrics) and replaces the layout profile's \SMFontAscender where
	% the font gives one; the Top Margin is the template's (SlideMage_setup.tex). Not replayed by
	% a precompiled format: the macro it sets lives in the format.
	\begingroup\normalfont\selectfont\directlua{slidemage_measure.derive_ascender()}\endgroup
	\SMLuaOnce{
		slidemage_measure.configure{top_margin = "\SMTopMargin", ascender = \SMFontAscender}
	}
\fi
\hypersetup{pdfinfo={SMFontAscender={\SMFontAscender}}}
\def\SMTable@begin#1{\begin{tabular}{#1}}
% BLOCK-RULE-COLOR, automatic (2026-09-30): inside any block (the
% convention above) the rules default to c_primary_dark -- all 11 tables
% placed in a block so far set that by hand. An \arrayrulecolor inside the
% table still overrides it, as before.
\newif\ifSM@inblock
\AtBeginEnvironment{block}{\SM@inblocktrue}
\AtBeginEnvironment{exampleblock}{\SM@inblocktrue}
\AtBeginEnvironment{alertblock}{\SM@inblocktrue}
\AtBeginEnvironment{infoblock}{\SM@inblocktrue}
\AtBeginEnvironment{referenceblock}{\SM@inblocktrue}
\AtBeginEnvironment{definitionblock}{\SM@inblocktrue}
\NewDocumentEnvironment{SMTable}{O{}m}{%
	\begingroup
	\ifSM@inblock
		\arrayrulecolor{c_primary_dark}%
	\else
		\arrayrulecolor{c_primary_highlight}%
	\fi
	\let\SMTable@gap\@empty
	\pgfqkeys{/SMTable}{#1}%
	\ifx\SMTable@gap\@empty
		\def\SMTable@spec{#2}%
	\else
		\setlength{\tabcolsep}{0pt}%
		\directlua{slidemage_measure.set_gap_spec(
			"\luaescapestring{\detokenize{#2}}",
			"\luaescapestring{\detokenize\expandafter{\SMTable@gap}}")}%
	\fi
	\expandafter\SMTable@begin\expandafter{\SMTable@spec}%
}{%
	\end{tabular}%
	\endgroup
}
\makeatother

% For a side-corner image overlay (the usual \vbox to 0pt{...} tikz overlay
% pattern used throughout this deck's converted slides) whose *top edge*
% should line up with the content text's own top edge -- distance from the
% page's own top edge to the cap-height top of a standardframe's first body
% line, measured directly against the compiled PDF's own text coordinates,
% not derived from \SMContentTopShift or guessed. Previous per-slide
% image placements each guessed their own downward shift (e.g. a bare
% "-1.3cm") independently -- close, but not exact, landing the image's top
% a couple mm above the text's own top; use this constant instead of a
% fresh guess for any new "text top edge"-aligned image.
%
% 41.58pt (2026-07-22, re-measured): the previous value, 44.0pt, went stale
% the moment \SMContentTopShift itself was increased (7.33pt -> 11.5964pt,
% "reduce the Top Margin", see that length's own comment above) without a
% matching update here -- that change pulls plain body text up, but this
% length is a separate, hand-measured constant with no formula tying it to
% \SMContentTopShift, so it kept anchoring images at the *old*, lower
% position. Confirmed directly (pymupdf against this same freshly-compiled
% preview deck, not assumed from the \SMContentTopShift delta alone):
% three independent plain-first-line standardframes with no per-slide
% top-shift correction (slides/5v4-CML-vs-SML-current.tex,
% slides/u5j-Exercise-CAPM-4-current.tex,
% slides/nt4-Exercise-CAPM-5-current.tex) all land their first line's own
% top edge at y=41.58pt post-change, not 44.0pt. Re-measure and adjust
% again if \SMContentTopShift ever changes in the future.
% (v9.7.0: \SMContentTopAlign is derived from \SMTopMargin, further up; it was a hook of SlideMage_setup.tex, 41.58pt, in v9.4.0-v9.6.0.)

% For a "side filler image" -- a decorative image stuck flush to the page's
% own right edge (zero margin) with its top edge flush against the title
% bar's own bottom edge (zero gap), rather than sitting in the usual
% content-aligned corner inset -- distance from the page's own top edge
% down to the standardframe title bar's bottom edge, measured directly
% against the compiled PDF (a pixel scan for the yellow-to-white
% transition, at this deck's 16cm-wide/9cm-tall page size). First used on
% slides/dfc-Angel-Investors-current.tex.
%
% 25.43pt -- the raw measured value (originally 25.44pt; re-measured
% 2026-07-07 against the current compiled PDF via the title bar's own
% white-patch rectangle, which came back 25.427pt, close enough to the
% original measurement to be the same underlying edge, not a drift).
% Briefly set to 24.94pt instead (2026-07-06-ish), a deliberate 0.5pt
% overlap *into* the title bar, after a report of a visible hairline
% white gap right at this seam in one render -- not reproduced in this
% project's own pdftoppm-based rendering checks at the time (every pixel
% scan at the raw value showed a clean cut, no gap), so assumed to be a
% renderer-specific sub-pixel rounding artifact rather than a real error
% in the constant, and "fixed" by deliberately overlapping instead of
% chasing the specific viewer. That overlap turned out to be the wrong
% call: direct feedback 2026-07-07 reported the *opposite* symptom on
% `sideimageframe` (the image's own top edge visibly overlapping the
% title box) -- exactly what a real 0.5pt overlap looks like once it's
% actually noticed, rather than a rendering artifact. Reverted back to
% the precise raw measurement: a real, small, deliberate overlap chosen
% to paper over an unreproduced gap in an unspecified other renderer
% turns out to be a worse trade than just using the measured value and
% letting a genuine future gap report (if one ever recurs, in a
% specific, identified renderer) be diagnosed on its own terms.
%
% Re-measured 2026-07-16, 25.43pt -> 25.32pt: the frametitle hook's
% outer sep fix (see its own comment) moved the visible yellow bottom
% edge up by the 0.2pt the trim patch's fill had previously been inset
% below its SMLogoAnchor.south anchor -- the yellow now genuinely
% ends flush with the logo's bottom edge, as always intended. New raw
% measurement against the rebuilt PDF: the trim patch's fill top edge
% (= the visible yellow bottom) sits 25.228 PDF units below the page
% top, i.e. 25.3226pt (earlier measurements in this comment quoted PDF
% coordinate values directly as "pt"; PDF units are big points, 1bp =
% 1.00375pt -- a systematic ~0.1pt understatement at this magnitude,
% harmless then, but corrected now that the constant is being reset
% anyway). Note this value is also what the overlay-logo nodes
% (titleframe/quoteframe/the hook's fixed-position redraw) derive their
% vertical inset from: visible-bar-bottom minus \SMLogoHeight is
% exactly the inline logo's own vertically-centered top margin.
% (v9.7.0: \SMTitleBarBottomAlign is derived from the paddings, the frame title font and the logo height, further up; it was a hook of SlideMage_setup.tex, 25.32pt, in v9.4.0-v9.6.0.)

% Note on the [t] default below: it top-aligns content within the frame
% (beamer's own default is vertical centering, "c"), so slides don't have
% their content floating in the middle. It also sidesteps a real beamer bug:
% \begin{frame}[]{Title} (a literal empty options bracket) fails outright,
% even in plain beamer with no theme at all -- so the optional argument must
% never default to an empty string. The title is always passed as an
% explicit argument (rather than left for beamer to pick up from a trailing
% brace group), since anything placed between \begin{frame} and its title --
% like our \setbeamertemplate calls -- breaks that trailing-brace mechanism.

% Pulls the frame's content up closer to the title bar: measured (via the
% PDF's own coordinates) at about 22pt of gap by default between the title
% bar's bottom edge and the first line of body content: this is a third of
% that, shrinking the gap to about two thirds of its original size.
%
% Called "Top Margin" in conversation from 2026-07-12 onward (direct
% instruction) -- the deck-wide name for the resulting gap between the
% title bar and a standardframe's first line of body content, i.e. the
% vertical position a plain-text frame's first baseline naturally lands at
% (measured directly against a reference slide's own compiled PDF, e.g.
% slides/aii-Leverage-Effects-current.tex's first line). Any first content
% that isn't plain text (a table, a minipage, an itemize with an unusual
% first item, ...) needs its own per-slide \SMIncreaseContentTopShift/
% \SMReduceContentTopShift correction to land on this same baseline --
% not portable between slides even for superficially similar content (see
% slides/zgz-EBIT-Definition-current.tex's own v1-2-0 comment for a case
% where an earlier, wrong correction was carried over uncritically).
% (\SMContentTopShift was a number of the layout profile from v9.4.0 to v9.8.0 and is measured
% since v9.9.0, see AUTOMATIC TOP MARGIN; its history is below.)
% 11.5964pt (2026-07-22, direct instruction: "reduce the Top Margin"
% generally, tested at 20%/30%/40% smaller against
% slides/8s8-Unit-Investment-Trusts-current.tex before settling on 40% --
% up from the original 7.33pt so the *resulting* gap (this length pulls
% content up by, so a bigger value means a smaller visible margin) comes
% out to 60% of its old size: measured directly on 8s8 at 10.67pt before,
% 6.42pt after (target was 6.40pt, within 0.02pt). Since every
% standardframe applies this same pull-up before any per-slide
% \SMIncreaseContentTopShift/\SMReduceContentTopShift correction,
% changing it here shifts the baseline project-wide, not just on 8s8.

% A `columns` environment that is the first content of a frame starts its first line
% ~10.17pt lower than a plain list does (measured on 5s6/vbb/yxv/c2t/n83/tj4/eey/yvw: first
% line box at 51.7pt against the 41.58pt Top Margin). It used to be corrected slide by
% slide with \SMIncreaseContentTopShift{10.13..10.17pt}, each value re-measured (and
% re-tuned when \SMContentTopShift changed); now the environment corrects itself
% (2026-09-21, direct instruction: "if this is a general problem for two-column slides, fix
% it generally"). "First content" = the frame's vertical list holds no box taller than
% 0.5pt yet -- the frame's own zero-size stamp boxes, glue and a zero-height overlay picture
% (\vbox to 0pt{...}) in front of the columns do not count, so a slide that puts such a
% picture first needs no special treatment either. (\prevdepth cannot tell: those zero-size
% boxes already set it.) Checked in LuaTeX by walking the current list.
% (\SMColumnsTopCorrection was a number of the layout profile from v9.4.0 to v9.8.0.)
\SMLuaOnce{
	slidemage_list_has_no_visible_box = function()
		local head = tex.nest[tex.nest.ptr].head
		for n in node.traverse(head) do
			if (n.id == node.id("hlist") or n.id == node.id("vlist")) and n.height > 32768 then
				return false
			end
		end
		return true
	end
}
% Since 2026-09-30 only for frames the automatic Top Margin doesn't manage
% (see AUTOMATIC TOP MARGIN below: [c] frames, sectionframe,
% \SMManualTopMargin, a deck with \SMAutoTopMarginfalse) -- in every
% other frame the measurement puts a columns-first slide on the margin by
% itself, so this correction (and the [T] compensation below) stays out.
% v9.9.0: the two corrections are measured where they are used, not numbers of the layout
% profile (10.17pt and 13.76pt, fitted on Fira; with Latin Modern Sans under another title bar a
% [T] columns landed 5.2pt high). Three reference lists are typeset in boxes at this point -- a
% top-level itemize, the same itemize in the first column of a columns[t], and of a columns[T]
% -- and \SMColumnsTopCorrection is how much lower the [t] one starts, \SMColumnsTTopCompensation
% how much of that the [T] one has taken back; a list in a column of either kind then starts
% where the list would start without the columns. Both names stay readable (the last values
% measured, 0pt before the first columns of the deck).
\makeatletter
\newcommand{\SM@columnsprobe}{%
	\SM@toprefprobe{\begin{itemize}\item\SM@topreftext\end{itemize}}\let\SM@columnsref\SM@topref
	\SM@toprefprobe{\begin{columns}[t]\begin{column}{0.46\linewidth}\begin{itemize}\item\SM@topreftext\end{itemize}\end{column}\end{columns}}%
	\global\SMColumnsTopCorrection=\dimexpr\SM@topref sp-\SM@columnsref sp\relax
	\SM@toprefprobe{\begin{columns}[T]\begin{column}{0.46\linewidth}\begin{itemize}\item\SM@topreftext\end{itemize}\end{column}\end{columns}}%
	\global\SMColumnsTTopCompensation=\dimexpr\SMColumnsTopCorrection+\SM@columnsref sp-\SM@topref sp\relax
}
\BeforeBeginEnvironment{columns}{%
	\ifSM@topprobing\else
	\ifnum0\ifdefined\SMTopAutoActive\SMTopAutoActive\fi=1
		\def\SMColumnsAtTop{0}%
	\else
		\edef\SMColumnsAtTop{\directlua{tex.sprint(slidemage_list_has_no_visible_box() and "1" or "0")}}%
	\fi
	\ifnum\SMColumnsAtTop=1
		\ifluatex\ifvmode\SM@columnsprobe\fi\fi
		\vspace*{-\SMColumnsTopCorrection}%
	\fi
	\fi
}
\makeatother
% The correction above is measured for `columns[t]`. `columns[T]` (top-aligned, used on
% most text-next-to-image slides) already starts its columns higher by itself -- beamer's
% own T key adds \vskip-1ex\nointerlineskip at every column head -- so the same pull-up left
% its first line 13.70pt (itemize) / 13.82pt (plain text) above the Top Margin, flush
% against the title bar (2026-09-30, direct instruction, seen on c58-DE-Unikat-Tipps).
% Measured against a plain standardframe: itemize 41.58pt, [t] 41.58pt, [T] 27.89pt;
% plain text 38.60pt, [T] 24.78pt. The T key gives the difference back while the columns
% environment is still in vertical mode, i.e. only for the columns[T] call itself (the
% per-column \setkeys calls run inside the columns' hbox, where \ifvmode is false) and only
% when the hook above applied its pull-up. A block as first content of a [T] column lands
% ~7pt below its plain-frame position with this -- extra clearance, not a collision.
% The column head also sets LaTeX's "top of a minipage" flag again (v5.14.1, found by
% check_layout.py's `coltop` in the SlideMage manual): beamer opens each column with
% \leavevmode, which clears it, so a list at the top of a [T] column kept its \topsep --
% zero at normal size, but a smaller list (\shrinklists, \small) loads \@listi's 4pt and
% started that much below a figure beside it that sat on the Top Margin.
% (\SMColumnsTTopCompensation was a number of the layout profile from v9.4.0 to v9.8.0; it is
% measured since v9.9.0, see above, and the figures in this paragraph are history.)
\makeatletter
\define@key{beamer@col}{T}[true]{%
	\def\beamer@colmode{T}\def\beamer@colalign{t}\def\beamer@colheadskip{\vskip-1ex\nointerlineskip\@minipagetrue}%
	\ifvmode
		\ifnum0\ifdefined\SMColumnsAtTop\SMColumnsAtTop\fi=1
			\vspace*{\SMColumnsTTopCompensation}%
		\fi
	\fi
}
\makeatother


% "Top Margin" (above) is one of four deck-wide margin names introduced in
% conversation from 2026-07-12 onward (direct instruction) for the plain
% page-edge/title-bar/footer gaps every slide is expected to obey by
% default. The other three, for reference (each maps to a length already
% declared elsewhere, not a new one introduced here):
% - "Left Margin": the default gap from the page's own left edge to the
%   title text and body content -- \SMFrametitleLeftPadding
%   (beamerthemeSlideMage.sty, 2.2ex, measured
%   directly against the compiled PDF at 12.552pt in this deck's usual
%   fonts -- see e.g. slides/tw5-Exercise-Leverage-and-ROE-Solution-current.tex's
%   own v1-1-0 comment for that same measured value).
% - "Right Margin": the default gap from the page's own right edge to the
%   institute logo (and, by the same convention, to anything else
%   deliberately flush-right, e.g. a corner-anchored image) --
%   \SMFrametitleRightPadding (0.28cm/7.93pt, this file's own
%   \SMInstituteLogo placement above uses it directly as the
%   logo's own inner sep).
% - "Bottom Margin": the default gap between a frame's own body content
%   and the \source{}/footnote line below it. Unlike the other three, this
%   one has no single settable length behind it -- \source{}'s own
%   vertical position isn't fixed the way the footer box is; it tracks the
%   frame's actual content height once that content gets tall enough (see
%   slides/u5h-Long-term-Bonds-Provisions-current.tex's own v1-3-0 comment,
%   "past a certain table height it starts tracking the table's own growth
%   roughly 1:1"). "Respecting Bottom Margin" in practice means leaving
%   comparable clearance to whatever a reference slide with unremarkable
%   content already shows, the same empirical, measure-against-a-reference-
%   slide approach "Top Margin" itself uses, rather than a fixed constant.

% Compensates a gap that opens up specifically when a standardframe's body
% starts with a \begin{minipage}{...} (used to width-limit text next to an
% image/diagram overlay, e.g. slides/igr-Budgeting-current.tex,
% slides/bdh-Working-Capital-Example-current.tex -- see README.md's
% "Available layouts" section for why minipage instead of columns there) --
% a minipage is a *box*, and TeX's vertical glue insertion (\topskip
% specifically) treats the first box in a vertical list differently from
% the first line of plain paragraph text, so the same \SMContentTopShift
% pull-up above lands closer to the title for minipage-wrapped content than
% for plain content. Cannot be fixed by \vspace(*) right after
% \begin{minipage}{...}: both plain and starred forms were tried and
% measurably had zero effect (the star only protects glue from page-break
% discard, not from \topskip's own separate glue-insertion calculation,
% which absorbs leading glue at the start of a vertical list regardless).
% Fixed instead by shrinking standardframe's own pull-up for the affected
% frame -- see \SMContentTopShiftReduction below. The value below is
% calibrated empirically, not derived from first principles: an isolated
% two-frame test (single-level itemize, one plain, one minipage-wrapped)
% measured only a 4.4pt gap, but real slides (nested itemize inside the
% minipage) showed a larger ~7.5-7.6pt gap against a real plain-itemize
% slide (slides/jyn-Categorization-current.tex) -- the extra few points
% likely come from itemize's own topsep/partopsep being computed differently
% depending on whether the list opens a box or a paragraph, on top of the
% topskip/box effect alone. Tuned so slides/igr-Budgeting-current.tex and
% slides/bdh-Working-Capital-Example-current.tex's first content line
% lands within 0.1pt of slides/jyn-Categorization-current.tex's own
% (200.917pt vs. 200.871pt/200.980pt in the compiled PDF).
% v9.4.0: the four corrections below (SMMinipageTopCorrection, SMMinipageWithCornerImageExtraCorrection,
% SMSideImageFrameMinipageCorrection, SMFlatItemizeTopCorrection) are 0pt. The automatic Top Margin (below) moves
% every frame it manages onto the Top Margin, so the slides that pass them to \SMReduceContentTopShift or
% \SMIncreaseContentTopShift land where they did: a copy of a real project's 16 decks (5274 pages, 122 slides
% that read SMMinipageTopCorrection, 123 SMFlatItemizeTopCorrection) rendered pixel-identically with the
% values at 0. The names stay defined because slides read them. The numbers in the comments are history.
\newlength{\SMMinipageTopCorrection}
% 9.38pt (2026-07-22, re-measured): the previous 7.6pt value went stale
% the moment \SMContentTopShift itself was increased (7.33pt ->
% 11.5964pt, "reduce the Top Margin", see that length's own comment above)
% without a matching update here -- confirmed directly (pymupdf against a
% freshly-compiled preview deck, not assumed from the \SMContentTopShift
% delta alone, which doesn't translate 1:1 into a minipage's own \topskip
% behavior): four independent slides all using this constant completely
% unmodified (slides/fr3-Forms-of-Capital-Market-Efficiency-Strong-Semistrong-Weak-current.tex,
% slides/za2-Startup-Financing-Categorization-current.tex,
% slides/yke-Implications-of-Homemade-Dividends-current.tex,
% slides/6j4-Modigliani-Miller-Assumptions-current.tex) all land their
% own first line at an identical y=39.80pt post-change, 1.78pt short of
% this deck's current 41.58pt Top Margin reference -- not slide-specific
% noise, the same deterministic amount on every one of them. Re-measure
% and adjust again if \SMContentTopShift ever changes in the future.
\setlength{\SMMinipageTopCorrection}{0pt}% obsolete since v9.4.0 (9.38pt before): the automatic Top Margin corrects it; slides still read the name

% A *third* minipage-top correction, on top of \SMMinipageTopCorrection
% above rather than replacing it, for a minipage whose own first internal
% content is a zero-height `\vbox to 0pt{...}` tikz overlay (a
% corner-anchored decorative image, placed there specifically so it stays
% unconditionally visible even when the minipage's own later content
% includes a `\pause` -- see e.g.
% slides/y5z-CoCo-Bonds-current.tex's own "v1-2-0" note for the fuller
% pause-inheritance story behind why the image has to sit there at all).
% 2026-07-08, found while chasing "the top margin of the text on y5z/22s
% is too low" -- measured directly against the compiled PDF, y5h/22s both
% landed 16.5pt too far down despite already using
% \SMMinipageTopCorrection, with the zero-height image overlay sitting
% *before* the minipage entirely as the actual root cause (that ordering
% makes the image, not the minipage, the frame body's true first box,
% which \topskip's own special first-box glue handling reacts to
% differently than a plain first minipage). Moving the image inside the
% minipage as its own first content closed most of the gap on its own
% (restoring the minipage as the frame's true first box) but left a
% further 10.35pt uncorrected -- confirmed the same on two independent
% slides (y5z, 22s) with identical structure, not assumed transferable
% from one to the other. Call via
% `\SMReduceContentTopShift{\dimexpr\SMMinipageTopCorrection-\SMMinipageWithCornerImageExtraCorrection\relax}`
% for any slide combining this exact pattern (corner image as a minipage's
% own first content, ahead of a `\pause`) -- confirmed transferable as-is
% to two more slides with the same structure
% (slides/qbj-Bank-Loans-Credit-Lines-current.tex,
% slides/jkf-Bank-Loans-Term-Loan-current.tex), neither of which
% actually uses `\pause` itself, so this isn't specific to that -- the
% image-inside-minipage-as-first-content shape alone is what triggers it.
\newlength{\SMMinipageWithCornerImageExtraCorrection}
\setlength{\SMMinipageWithCornerImageExtraCorrection}{0pt}% obsolete since v9.4.0 (10.35pt before), see above

% A *different* minipage-top correction, specifically for a minipage+
% itemize sitting as the first content of a `sideimageframe` (2026-07-06,
% found while chasing "the top margin of the text on ezz/acw is again
% off" -- see CONVERSION_NOTES.md's own "beamer/TeX gotchas" entry for the
% full story). `\SMMinipageTopCorrection` above is calibrated for a
% *plain* standardframe, and simply does not transfer to `sideimageframe`:
% measured directly on three real sideimageframe slides
% (slides/xmp-Implications-of-Capital-Market-Efficiency-Considerations-current.tex,
% slides/dfc-Angel-Investors-current.tex,
% slides/ezz-Examples-of-Tests-for-Market-Efficiency-current.tex --
% one flat itemize, one flat itemize, one nested-first-item itemize
% respectively, so this isn't a flat-vs-nested effect either), all three
% landed within 0.6pt of each other at their own *uncorrected* natural
% position (52.03-52.57pt against this deck's usual 44.0pt first-line
% reference) regardless of which of the two *wrong* corrections
% (`\SMReduceContentTopShift{\SMMinipageTopCorrection}` or
% `\SMIncreaseContentTopShift{\SMFlatItemizeTopCorrection}`) had
% been applied to each beforehand -- i.e. `sideimageframe` itself pushes
% its own first-content minipage down by a consistent ~8.2pt relative to
% a plain standardframe, on top of whatever generic itemize-shape
% correction would otherwise apply, most likely from the extra
% \sidefootnotes/\settowidth setup work `sideimageframe`'s own opening
% code does before `\begin{standardframe}` ever runs. Use
% `\SMIncreaseContentTopShift{\SMSideImageFrameMinipageCorrection}`
% (note: *increase*, not reduce) as the *only* topshift call on any
% future sideimageframe slide whose first content is a minipage+itemize --
% not the generic `\SMMinipageTopCorrection`/
% `\SMFlatItemizeTopCorrection` this deck uses everywhere else.
\newlength{\SMSideImageFrameMinipageCorrection}

% 6.26pt (2026-07-22, re-measured): stale for the same reason
% \SMMinipageTopCorrection above went stale -- \SMContentTopShift
% itself changed (7.33pt -> 11.5964pt) without a matching update here.
% Re-measured directly against a freshly-compiled preview deck across six
% independent sideimageframe slides using this constant unmodified: four
% (ezz, xmp, hzx, u6n) land their own first line at an identical
% y=39.64pt, 1.94pt short of this deck's current 41.58pt Top Margin
% reference (i.e. already pulled up *too far*); the other two (c7w, dfc)
% are close but not identical (39.53pt/-2.05, 40.18pt/-1.4) -- no single
% value fixes all six exactly, so this recalibrates to the four-way
% cluster (the majority case) and accepts a small (<0.6pt) residual on the
% two outliers. Since this constant feeds
% `\SMIncreaseContentTopShift{}` (more pull-up, i.e. *smaller* y), and
% the observed y was already short of target, the fix is a *decrease*
% (8.2pt - 1.94pt), not an increase -- caught by re-measuring after a
% first, wrong-direction attempt at 10.14pt drove the four-slide cluster
% to y=37.70pt, further from target instead of closer.
\setlength{\SMSideImageFrameMinipageCorrection}{0pt}% obsolete since v9.4.0 (6.26pt before, 8.2pt originally), see above

% Opt-in override: reduces standardframe's \SMContentTopShift pull-up by
% this amount for the frame currently being opened. Empty-default-style flag
% (mirrors \SMChapterLogoPath etc.): 0pt by default (no effect), set via
% \SMReduceContentTopShift{<length>} called *before* \begin{standardframe},
% reset back to 0pt in standardframe's own closing code so it never leaks
% into the next frame. Used by minipage-based slides to counteract
% \SMMinipageTopCorrection above.
\newcommand{\SMContentTopShiftReduction}{0pt}
\newcommand{\SMReduceContentTopShift}[1]{\global\def\SMContentTopShiftReduction{#1}}

% A *second*, separate first-line-position quirk, found and fixed
% 2026-07-05 after a run of slides (batch 4 plus several earlier ones) were
% found to land at two visibly different heights depending purely on
% content shape, with no minipage involved at all: a plain (non-minipage)
% standardframe whose top-level itemize's very *first* \item is immediately
% followed by its own nested \begin{itemize}/\begin{enumerate} lands 1.834pt
% higher than one whose first \item has no such immediate child (checked
% against every content slide converted so far via each compiled page's own
% first-line PDF baseline coordinate, not eyeballed -- 13 "nested-first-item"
% slides all landed within 0.05pt of each other, and 9 "flat-first-item"
% slides all landed within 0.05pt of each other, at a clean +1.834pt offset
% from the first group; isolated two-frame test cases, differing only in
% presence/absence of that immediate nested list, reproduced the exact same
% 1.834pt gap independent of which text/list type/item count was used).
% Not chased to one named TeX primitive (most likely some interaction of
% \addvspace's glue-coalescing behaviour with \partopsep/\topsep at the
% \begin{itemize}/\begin{frame} boundary, given the existing
% \SMMinipageTopCorrection comment above already flags topsep/partopsep
% as varying with box vs. paragraph context) -- calibrated empirically like
% \SMMinipageTopCorrection, not derived from first principles.
% Interacts with minipage differently, not additively: inside a minipage,
% this nested-vs-flat distinction disappears entirely (both land within
% 0.05pt of each other) -- so this correction is for *plain* itemize-first
% standardframes only, never combined with \SMMinipageTopCorrection.
\newlength{\SMFlatItemizeTopCorrection}

% 0.01pt (2026-07-22, re-measured): stale for the same reason
% \SMMinipageTopCorrection above went stale -- \SMContentTopShift
% itself changed (7.33pt -> 11.5964pt) without a matching update here.
% Re-measured directly against three independent flat-itemize slides using
% this constant unmodified (3ce, 5cr, c8e): all three landed their own
% first line at an identical y=39.76pt, 1.82pt short of this deck's
% current 41.58pt Top Margin reference -- coincidentally, the recalibrated
% \SMContentTopAlign reference now already accounts for almost all of
% what this correction used to supply on its own, leaving next to nothing
% left to add. Effectively zeroed out rather than removing the call
% itself, so every flat-itemize slide's own
% `\SMIncreaseContentTopShift{\SMFlatItemizeTopCorrection}` line
% keeps working unchanged.
\setlength{\SMFlatItemizeTopCorrection}{0pt}% obsolete since v9.4.0 (0.01pt before, 1.834pt originally), see above

% Opposite-direction counterpart to \SMReduceContentTopShift: *increases*
% the pull-up instead of reducing it (needed for \SMFlatItemizeTopCorrection
% above, since a flat first item needs *more* pull-up, not less). Same
% storage variable, just negated -- \SMContentTopShiftReduction is
% allowed to go negative, since the frame-opening code only ever computes
% \SMContentTopShift-\SMContentTopShiftReduction, never checks its
% sign.
\newcommand{\SMIncreaseContentTopShift}[1]{\global\def\SMContentTopShiftReduction{-#1}}

% AUTOMATIC TOP MARGIN (2026-09-30, direct instruction: "propose a more
% systematic solution" -- 195 slides carried a hand-measured
% \SMIncrease/ReduceContentTopShift, 89 of them a raw pt value, every one
% of them re-measured by hand whenever \SMContentTopShift or a frame type
% changed, and a stale one survived a central change unnoticed more than
% once, e.g. six slides double-corrected after the 2026-09-21 columns fix).
% Where a frame's first line lands depends on what comes first -- a plain
% paragraph sits 3pt higher than a list, an image 3.5pt higher, a minipage,
% columns, a table each differently -- so no constant pull-up can be right
% for all. Instead every top-aligned standardframe (and figureframe,
% highlightframe, sideimageframe) now measures where its first visible
% content actually landed (at shipout, SlideMage_measure.lua; overlays at
% absolute positions don't count, they are placed on the margin already) and
% moves its body onto the Top Margin line on the next run, via \jobname.top
% -- the same record-then-apply-on-rerun pattern as \jobname.ovl above;
% latexmk reruns while the file changes, and one rerun converges.
% Any older \SMIncrease/ReduceContentTopShift call still works but no
% longer matters: the measurement corrects whatever is left.
% Exceptions: [c] frames (centred on purpose: no pull-up at all), sectionframe
% (own geometry), and any frame opened after
%   \SMManualTopMargin
% which keeps the frame's default position: no correction that depends on what
% the frame starts with. (A \progressstrip frame is measured since v5.30.5.)
% The default position (v9.9.0) is \SMContentTopShift, the pull-up that puts the
% first line of a top-level list on the Top Margin in a standard frame. It used to
% be a number of the layout profile (11.5964pt, fitted on Fira under the default
% title bar: with another title bar or font the frames above landed up to 2.9pt
% off); now every top-aligned frame records where its body begins (the origin
% marker) and how far below that a top-level itemize's first line would start
% (typeset in a box there), and the measurement writes the pull-up most frames
% share into \jobname.top. The corrections of the measured frames are relative
% to it. In a run without \jobname.top the first titled frame's own value is used
% from the next frame on, so the first run already places the deck as the next one
% will (no deck needs a run more than with the old stored number); a frame that
% started that run elsewhere (the first one) is corrected from its default position. \SMAutoTopMarginfalse in a deck's preamble places every frame like
% \SMManualTopMargin does (the default position, measured; no content-dependent
% correction anywhere in the deck); it does not switch the measurement off.
\newif\ifSMAutoTopMargin
\SMAutoTopMargintrue
\makeatletter
\newif\ifSM@topmanual
\newif\ifSM@topsection
\newcommand{\SMManualTopMargin}{\global\SM@topmanualtrue}
\newcount\SM@topindex
\newwrite\SM@topfile
\newcommand{\SMSetTopCorrection}[2]{\global\@namedef{SM@topcorr@#1}{#2}}
\newcommand{\SMSetContentTopShift}[1]{\global\SMContentTopShift=#1\relax
	\ifluatex\directlua{slidemage_measure.top_shift_read = \strip@pt\SMContentTopShift}\fi}
\ifluatex
	\SMLuaOnce{
		luatexbase.add_to_callback("pre_shipout_filter", function(head)
			if slidemage_measure and slidemage_measure.top_measure_page then
				slidemage_measure.top_measure_page(head)
			end
			return true
		end, "slidemage top margin")
	}
	\AtBeginDocument{%
		\InputIfFileExists{\jobname.top}{}{\typeout{No file \jobname.top.}}%     % "No file": see \jobname.ovl
		\immediate\openout\SM@topfile=\jobname.top\relax
	}
	\AtEndDocument{\directlua{slidemage_measure.top_emit()}\immediate\closeout\SM@topfile}
\fi
% Before \begin{frame} (runs once per frame, not once per overlay page):
% number the frame and fetch its correction from the previous run.
\newcommand{\SM@topprepare}{%
	\ifdefined\SM@topshiftrun% a first run: measured on the first frame that counts
		\global\SMContentTopShift=\SM@topshiftrun\relax
		\global\let\SM@topshiftrun\undefined
	\fi
	\global\advance\SM@topindex by 1
	\xdef\SM@topthis{\the\SM@topindex}%
	\ifcsname SM@topcorr@\SM@topthis\endcsname
		\xdef\SM@topcorr{\csname SM@topcorr@\SM@topthis\endcsname}%
	\else
		\gdef\SM@topcorr{0pt}%
	\fi
	\ifSMAutoTopMargin\ifSM@topmanual\gdef\SM@topcorr{0pt}\fi\else\gdef\SM@topcorr{0pt}\fi
}
% In the frame body, right after the pull-up: the start marker, then the origin marker.
\newcommand{\SM@topstart}{%
	\ifluatex
		\directlua{slidemage_measure.top_begin(\SM@topthis, \strip@pt\dimexpr\SM@topcorr\relax,
			\number\dimexpr\SMContentTopShift-\SMContentTopShiftReduction+\SM@topcorr\relax,
			\strip@pt\SMContentTopShift, \ifSMAutoTopMargin\ifSM@topmanual false\else true\fi\else false\fi)}%
		\ifSMAutoTopMargin\ifSM@topmanual\else
			\def\SMTopAutoActive{1}% this frame is measured: the columns hook stays out
		\fi\fi
		\SM@toprefprobe{\begin{itemize}\item\SM@topreftext\end{itemize}}%
		\directlua{slidemage_measure.top_origin(\SM@topthis, \SM@topref,
			\ifSM@topsection false\else\ifSM@toptitled true\else false\fi\fi)}%
	\fi
}
% The reference list's first line box, in sp below this point (\SM@topref): #1 typeset in
% a box that starts as the body does here (the same \prevdepth, font, \parskip and list
% settings), measured by SlideMage_measure.lua like a frame's first content. Nothing of it
% reaches the page; a deck's default overlay specification stays out (an \item<+-> would
% count a step).
% The reference line's text: a capital and an ascender, as tall as a line of running text (a
% [T] column aligns its first line by the line's height, not by \baselineskip).
\newcommand{\SM@topreftext}{Xh}
\newbox\SM@topprobebox
\newif\ifSM@topprobing
\newif\ifSM@toptitled
\newcommand{\SM@toprefprobe}[1]{%
	\begingroup
	\SM@topprobingtrue
	\def\beamer@defaultospec{<*>}% beamer's own default: every step
	\def\SMColumnsAtTop{0}%
	\ifvmode\dimen@=\prevdepth\else\dimen@=\z@\fi
	\setbox\SM@topprobebox=\vbox{\prevdepth=\dimen@#1\par}%
	\xdef\SM@topref{\directlua{tex.sprint(slidemage_measure.box_content_top(\number\SM@topprobebox))}}%
	\endgroup
}
% End of the body, or anywhere a frame wants to stop the measurement
% (\progressstrip): nothing after this marker counts as first content.
\newcommand{\SM@topstop}{%
	\ifluatex
		\directlua{slidemage_measure.top_end(\SM@topthis)}%
	\fi
}
\makeatother

% Define custom frame type with en-dash bullets
\newenvironment{standardframe}[2][t]{%
	% Restart footnote numbering (the *, dagger, ... symbols) fresh on every
	% slide, rather than letting it run on cumulatively across the whole deck
	% the way beamer does by default. Two reasons: \fnsymbol only defines 9
	% symbols before erroring out, which a real multi-slide deck would run
	% through in no time if every \footnote{}/\source{} call across the whole
	% talk drew from one shared counter; and each slide's footnotes are read
	% and referenced only on that one slide anyway, so there's no reason for
	% their numbering to depend on how many footnotes earlier slides happened
	% to have. This must run *before* \begin{frame}: beamer saves the
	% counter's current value as that frame's baseline as part of
	% \begin{frame}'s own internal setup (restoring it before every overlay
	% replay, so overlays within one frame don't double-count) -- resetting
	% after \begin{frame} would be reset too late, after that baseline was
	% already captured.
	\setcounter{footnote}{0}%
	\csname SM@topprepare\endcsname
	\ifstrempty{#2}{\global\csname SM@toptitledfalse\endcsname}{\global\csname SM@toptitledtrue\endcsname}%
	\SMOpenFrame{[#1]{#2}}%
	% Forces this frame's own body to start with the correct default text
	% color, regardless of what any *previous* frame left active
	% (2026-07-16, direct instruction, root-caused via an isolated minimal
	% reproduction). \titleframe's/\quoteframe's own
	% \color{\SMTitleTextColor} calls -- set inside content built via a
	% stored box register or a zero-space overlay, used/shipped out *later*
	% rather than immediately -- don't reliably get "undone": xcolor's
	% \color doesn't emit a PDF color operator immediately when called, only
	% when the *next* real glyph/path actually needs it, so whatever color
	% one of those frame types last set can end up as the active PDF color
	% well into the *next* frame's own first draw call. Invisible for Graz
	% (the leaked color, c_primary_dark, happens to match ordinary body text
	% anyway) but catastrophic for MCI/Innsbruck (c_primary_white, rendering
	% plain non-itemize body text invisible against its own white
	% background -- itemize items escape it, since beamer's own itemize
	% template sets color fresh per bullet). An \AtBeginEnvironment{frame}
	% hook (etoolbox) was tried first and confirmed *not* to fire at all --
	% beamer's own "frame" isn't a plain \newenvironment etoolbox can patch
	% that way -- so this resets directly in standardframe's own opening
	% code instead, exactly where this project already has full control.
	\color{c_primary_dark}%
	% Defensive: the tokens after \begin{frame} are collected into the frame
	% body and replayed once per overlay slide, so this clears the side-
	% footnote collector at the start of every slide's pass -- normally a
	% no-op (the previous slide's footline already consumed the box), but it
	% stops stale footnotes leaking in if a previous frame ever collected
	% without shipping a footline.
	\SMClearSideFootnotesBox%
	\SMResetImageSources%
	% The itemize bullets are the theme's (\SMBulletDot, beamerthemeSlideMage.sty).
	% Only for top-aligned frames (the default, #1=t): a [c]-centered frame
	% (2026-07-22, direct instruction, found while centering
	% slides/bw9-Pearson-Correlation-Illustration-current.tex/agc/crj) was
	% landing its content almost flush against the title bar instead of
	% truly centered -- root cause was this same pull-up firing
	% unconditionally regardless of #1. beamer's own [c] option already
	% centers the frame body within the available space; stacking this
	% deck's *own* top-alignment pull-up on top of that just drags the
	% already-centered box further upward by the same fixed amount, which
	% is exactly the asymmetry that was measured (a [c] figureframe's image
	% landing ~0pt from the title bar and ~10pt from the footer instead of
	% symmetric). Guarded with \ifstrequal (etoolbox, already loaded above)
	% rather than \ifx, matching this file's own \ifstrequal{##1}{...}
	% pattern elsewhere (\ifx would need #1 fully expanded first).
	\ifstrequal{#1}{c}{}{%
		\vspace*{-\dimexpr\SMContentTopShift-\SMContentTopShiftReduction+\csname SM@topcorr\endcsname\relax}%
		\csname SM@topstart\endcsname}%
}{%
	\csname SM@topstop\endcsname
	% Must run *before* \end{frame}: like any other frame content, a
	% \footnote{}/\source{} call has to happen while the frame body is
	% still open -- everything in this closing block from \end{frame}
	% onward runs *after* the frame has already closed (see README.md's
	% own \titleframe comment on beamer's \beamer@collect@@body for the
	% general shape of this constraint).
	\SMEmitImageSourceFootnote%
	\end{frame}%
	\global\def\SMChapterLogoPath{}% opt-in only for the frame just closed
	\global\def\SMContentTopShiftReduction{0pt}% opt-in only for the frame just closed
	\global\csname SM@topmanualfalse\endcsname% opt-in only for the frame just closed
	\global\csname SM@topsectionfalse\endcsname% set by sectionframe for the frame just closed
	\global\SMSideFootnotesfalse% opt-in only for the frame just closed
	\global\def\SMFrametitleFullHeightYellow{}% opt-in only for the frame just closed
	\global\def\SMFrametitleLogoFixed{}% opt-in only for the frame just closed
	\global\def\SMFooterBoxColor{\SMFooterBoxColorDefault}% opt-in only for the frame just closed
}


% "a)"/"b)"/"c)"-labeled list, for exercise sub-tasks that read better
% lettered than bulleted (2026-07-20, direct instruction, first used on
% slides/y3a-Exercise-Diversification-2-current.tex's own bottom three
% items).
%
% First attempt (superseded, kept as a documented wrong turn -- see
% README.md's own "lualatex vs latex" section for why this project keeps
% wrong write-ups rather than deleting them): reuse itemize's own top-level
% `\llap`-based label positioning by locally overriding
% `\setbeamertemplate{itemize item}` inside a plain `\begin{itemize}` --
% the reasoning being that `\llap` draws its content flush against the
% text's own left edge "regardless of the label's width." That reasoning
% turned out to be wrong in a way that only showed up once actually
% measured against the compiled PDF (not just visually eyeballed): a plain
% bullet's own `\llap`-drawn footprint (title/Left-Margin-relative-left-
% edge to text-start) measured only ~9.05pt end-to-end on this deck's real
% slides, matching the Left-Margin-to-text-start gap exactly -- but "a)"
% rendered through the *exact same* mechanism measured a ~13.97pt
% footprint for the same gap, not the ~8.5-9pt `\settowidth` reported for
% the bare glyphs in isolation. The extra ~5pt was never fully root-caused
% (this custom itemize's own `\list` call, further below, only overrides
% `\makelabel` -- it never sets `\leftmargin`/`\labelwidth` itself, so
% those are left at whatever this deck's global itemize defaults happen to
% be, and `\llap`'s actual box-width behavior inside that inherited,
% unexamined list-parameter state is where the extra width most likely
% comes from) -- rather than keep guessing at *why*, measured the actual
% failure directly (the compiled "a)" sat ~4.9pt left of the Left Margin,
% confirmed via exact PDF text-position extraction, not a visual guess)
% and abandoned the approach once even `\tiny` plus manual per-letter
% kerning only just barely closed the gap while reading as illegibly small
% -- a real dead end, not a tuning problem.
%
% Working fix: don't reuse itemize's `\llap` label mechanism at all -- use
% a real, width-reserving `\makebox[W][r]{...}` label instead (the same
% *kind* of mechanism `\SMEnumerateLabelWidth` above already uses for
% a nested enumerate, just with an explicit `\leftmargin`/`\labelwidth`/
% `\labelsep` this environment sets itself rather than inheriting), inside
% its own plain `\list`, not a reused `\itemize`. `\leftmargin` (label box
% + gap) is set to 11.1pt -- 9.1pt label box (fits "a)"/"b)"/"c)" at this
% deck's normal `\small` itemize size with a little clearance, confirmed
% directly: measured label-left-edge stays right of the 12.552pt Left
% Margin for all three letters, closest case (widest letter, "c)") 1.37pt
% clear) plus 2pt `\labelsep` for a visible gap before the text. Trade-off,
% accepted deliberately rather than found by accident: item text lands
% ~2.2pt right of where a plain bulleted item's text would (23.8pt vs.
% 21.6pt) -- fitting a legible "a)"/"b)"/"c)" at normal reading size plus
% a real gap simply doesn't fit in the ~9pt a plain bullet's tiny dot
% leaves, so *exact* text-position parity with plain bullets and a legible
% label are mutually exclusive here; legibility won given the two-point
% shift is visually negligible and there's no plain-bulleted item left on
% the same slide to compare against directly anyway (y3a's own first item
% moved into its infoblock). Re-derive `\leftmargin`/`\labelwidth` the same
% way (render, measure with exact text-position extraction, adjust) if
% this is ever used at a different font size than this deck's normal
% itemize `\small`.
% The two lengths of every list with "a)" / "1." / "1." / "a." labels (v9.8.0: derived; they were
% 9.1pt and 4.75pt in the Fira profile and 9.16pt and 5.20pt in the lmsans one). \SMLabelItemWidth
% is the box that holds the widest label of the first items, a)-c) and 1)-9) with a bracket or a
% period, in the font in force where the list begins, plus 0.1pt of clearance; \SMLabelItemSep is the
% hook \SMLabelItemGap (in em of that font). Called first in each list's setup, so that
% the lists follow \small, \footnotesize and \shrinklists, and once at \begin{document}
% (in \small) for a slide that reads the two lengths itself.
\makeatletter
\newcommand{\SM@labelitemmetrics}{%
	\setlength{\SMLabelItemWidth}{0pt}%
	\@for\SM@labelprobe:=a),b),c),a.,b.,c.,1),2),3),4),5),6),7),8),9),1.,2.,3.,4.,5.,6.,7.,8.,9.\do{%
		\settowidth{\dimen@}{\SM@labelprobe}%
		\ifdim\dimen@>\SMLabelItemWidth\SMLabelItemWidth=\dimen@\fi}%
	\addtolength{\SMLabelItemWidth}{0.1pt}%
	\setlength{\SMLabelItemSep}{\SMLabelItemGap}}
\makeatother
\newcounter{SMAlphaItemCounter}
\newcommand{\SMAlphaItem}{\stepcounter{SMAlphaItemCounter}\alph{SMAlphaItemCounter})}
\newenvironment{SMAlphaItemize}{%
	\setcounter{SMAlphaItemCounter}{0}%
	\list{\SMAlphaItem}{%
		% 4.75pt (2026-07-20, direct instruction: "half the width of the
		% 'C' in 'Calculate'" -- \settowidth{C} at this deck's normal
		% itemize \small measured 5.5pt, so 2pt original + 2.75pt half-C).
		% Only \labelsep changes the text's own start position -- the
		% label's own box stays exactly where it was (its left edge is
		% pinned to this list's origin regardless of \labelsep, since
		% \leftmargin=\labelwidth+\labelsep always cancels \labelsep back
		% out of the label box's own position), so this widens the gap
		% without reopening the Left-Margin question already solved above.
		\csname SM@labelitemmetrics\endcsname\setlength{\labelsep}{\SMLabelItemSep}%
		\setlength{\labelwidth}{\SMLabelItemWidth}%
		\setlength{\leftmargin}{\labelwidth}%
		\addtolength{\leftmargin}{\labelsep}%
		\setlength{\itemindent}{0pt}%
		% \itemsep/\parsep left at plain \list's own defaults otherwise --
		% since this environment is built directly on \list rather than
		% reusing beamer's own itemize template, it never got beamer's own
		% itemize spacing hooks at all, only whatever \list falls back to.
		% Measured directly (2026-07-21, direct instruction, found on
		% slides/nt4-Exercise-CAPM-5-current.tex): that default reads as
		% *less* space between two items (13.7pt, item-to-item) than
		% between two wrapped lines within one *single* item (15.6pt,
		% pure \baselineskip) -- backwards from how a list should read
		% (item boundaries should be at least as open as a same-item line
		% wrap, not tighter). \itemsep set explicitly so the total gap at
		% an item boundary (baselineskip + itemsep) reads as clearly more
		% open than a plain wrapped line, not accidentally tighter.
		% Reduced from the original 6pt to 2.5pt (2026-07-22, found during
		% a full-project check after the itemize-rules work, not part of
		% that work itself): 6pt overshot by more than needed to clear the
		% 1.9pt gap (15.6 vs. 13.7), and slides/kmf-Exercise-CAPM-6-current.tex's
		% own 5-item list -- longer than nt4's own -- compounded that
		% excess across 4 item boundaries into a ~25pt overfull frame, its
		% own last item cut off by the footer (confirmed directly,
		% rendered and compared against a pre-regression PDF). 2.5pt
		% keeps a real, visible margin over the 1.9pt needed (not just
		% barely crossing it) while re-measuring clean (no overfull) on
		% kmf; re-verify against both nt4 and kmf directly if this is
		% ever changed again -- they're this environment's own tightest
		% (most items) and originally-tuned-against cases respectively.
		\setlength{\itemsep}{\SMLabelItemItemsep}%
	}
	\raggedright% plain \list defaults to justified+hyphenated text, unlike
	% this deck's own itemize (which calls this too, further below) --
	% missed on the first pass, only caught by comparing a rendered a)/b)/c)
	% item side by side against a plain bulleted one ("stan-\ndard deviation"
	% hyphenating where a plain-bulleted version of the same sentence
	% elsewhere in this deck doesn't).
}{%
	\endlist
}

% Three siblings of \SMAlphaItemize above, same measured \list metrics
% (labelsep/labelwidth/leftmargin/itemindent/itemsep), only the label text
% itself differs -- added 2026-08-09 for slides whose numbering has to
% match a real external scheme (statute paragraph numbers, PowerPoint's
% own "1./2." or "a./b." auto-numbering) rather than plain letters-with-
% parenthesis. Not merged into one parameterized environment: each name is
% called directly from a slide file, and duplicating six short \setlength
% lines three times is cheaper to read at the call site than a macro
% layer over \newenvironment for four near-identical bodies.
\newcounter{SMNumItemCounter}
\newcommand{\SMNumItem}{\stepcounter{SMNumItemCounter}\arabic{SMNumItemCounter})}
\newenvironment{SMNumItemize}{%
	\setcounter{SMNumItemCounter}{0}%
	\list{\SMNumItem}{%
		\csname SM@labelitemmetrics\endcsname\setlength{\labelsep}{\SMLabelItemSep}%
		\setlength{\labelwidth}{\SMLabelItemWidth}%
		\setlength{\leftmargin}{\labelwidth}%
		\addtolength{\leftmargin}{\labelsep}%
		\setlength{\itemindent}{0pt}%
		\setlength{\itemsep}{\SMLabelItemItemsep}%
	}
	\raggedright%
}{%
	\endlist
}

\newcounter{SMAlphaPeriodItemCounter}
\newcommand{\SMAlphaPeriodItem}{\stepcounter{SMAlphaPeriodItemCounter}\alph{SMAlphaPeriodItemCounter}.}
\newenvironment{SMAlphaPeriodItemize}{%
	\setcounter{SMAlphaPeriodItemCounter}{0}%
	\list{\SMAlphaPeriodItem}{%
		\csname SM@labelitemmetrics\endcsname\setlength{\labelsep}{\SMLabelItemSep}%
		\setlength{\labelwidth}{\SMLabelItemWidth}%
		\setlength{\leftmargin}{\labelwidth}%
		\addtolength{\leftmargin}{\labelsep}%
		\setlength{\itemindent}{0pt}%
		\setlength{\itemsep}{\SMLabelItemItemsep}%
	}
	\raggedright%
}{%
	\endlist
}

\newcounter{SMNumPeriodItemCounter}
\newcommand{\SMNumPeriodItem}{\stepcounter{SMNumPeriodItemCounter}\arabic{SMNumPeriodItemCounter}.}
\newenvironment{SMNumPeriodItemize}{%
	\setcounter{SMNumPeriodItemCounter}{0}%
	\list{\SMNumPeriodItem}{%
		\csname SM@labelitemmetrics\endcsname\setlength{\labelsep}{\SMLabelItemSep}%
		\setlength{\labelwidth}{\SMLabelItemWidth}%
		\setlength{\leftmargin}{\labelwidth}%
		\addtolength{\leftmargin}{\labelsep}%
		\setlength{\itemindent}{0pt}%
		\setlength{\itemsep}{\SMLabelItemItemsep}%
	}
	\raggedright%
}{%
	\endlist
}

% Itemize indentation: level 1 bullets should sit flush with the frame title
% (no extra left margin), and each deeper level should be indented just
% enough for its bullet to line up under the *text* (not the bullet) of the
% level above. A nested list's \leftmarginN is added on top of its parent's
% (indentation is cumulative), so giving every level the same "step" --
% label width + labelsep, the width a bullet plus its gap to text takes up --
% makes level N's text sit at N*step and level N's bullet at (N-1)*step, i.e.
% exactly at level (N-1)'s text (0 for level 1, flush with the title).
\makeatletter
\newlength{\SMBulletWidth}
\settowidth{\SMBulletWidth}{\SMBulletDot}
\newlength{\SMItemStep}
\setlength{\SMItemStep}{\labelsep}
\addtolength{\SMItemStep}{\SMBulletWidth}
\setlength{\leftmargini}{\SMItemStep}
\setlength{\leftmarginii}{\SMItemStep}
\setlength{\leftmarginiii}{\SMItemStep}
\setlength{\leftmarginiv}{\SMItemStep}
\makeatother

% \item[\SMBulletEllipsis] -- an ellipsis used as the item label (the source's
% "...how you contribute" continuation items). A plain \item[\ldots] right-aligns the
% label to the bullet dot's right edge, so the wider ellipsis pokes out to the left
% of where a bullet starts; this one starts at the dot's left edge instead
% (2026-09-20, direct instruction: the ellipses were too far left).
\newcommand{\SMBulletEllipsis}{\hbox to \SMBulletWidth{\rlap{\ldots}\hss}}

% \setbeamertemplate{itemize subitem}{\SMBulletEllipsisScript} -- the same for a list whose *bullet itself* is a small ellipsis (the source's
% "..." sub-points; slides quy, 3a8, unt). The bare \scriptsize… bullet used to sit 1.6pt left of where a dot starts (2026-09-20).
\newcommand{\SMBulletEllipsisScript}{\hbox to \SMBulletWidth{\rlap{\scriptsize…}\hss}}

% A 4th itemize level -- beamer's own \itemize (beamerbaselocalstructure.sty)
% hard-caps nesting at 3 levels (\ifnum \@itemdepth >2\relax\@toodeep),
% independently of anything this deck sets: a 4th \begin{itemize} always
% raised "! LaTeX Error: Too deeply nested." even though plain LaTeX's own
% kernel itemize supports 4. Found 2026-07-15, direct instruction, on
% slides/hrg-Exchange-Rate-Quotation-Forward-current.tex -- genuinely
% 4 levels deep in the source pptx itself (confirmed via a PowerPoint COM
% render, not assumed), not a conversion artifact to design around.
% Since this deck already styles every existing level identically (same
% \SMBulletDot, same \SMItemStep margin -- see the comments right
% above/below), extending to a 4th level needs no new visual design, just
% patching beamer's own depth limit up by one and telling it what to call
% the new level ("subsubsub", the next name in beamer's own
% sub/subsub/... naming scheme) so its template/color lookups
% (\usebeamertemplate{itemize subsubsubitem}, \usebeamercolor{itemize/
% enumerate subsubsub body}, etc.) resolve to *something* instead of
% erroring -- undefined beamer templates/colors are expected to fall back
% to sensible defaults (no extra spacing hook, plain text color) rather
% than raising an error, same as any other beamer template/color that a
% given theme never bothers to set; the actual bullet symbol still needs
% its own \SMBulletDot template (in the theme), added explicitly the same way
% as the existing 3 levels, not left to whatever fallback beamer has for
% an unset "itemize subsubsubitem" template.
% \beamer@computepref is patched wholesale (not via \patchcmd) since its
% body is one single \ifcase chain, not a stable string to match against.
% \itemize itself was *first* tried via \patchcmd{\itemize}{...}
% (matching this deck's own established practice for beamer/theme
% internals, e.g. beamerthemeSlideMage.sty's own
% \patchcmd{\beamer@title}{...}) -- this failed outright, confirmed
% directly (\patchcmd's own failure branch fired, raising a clean
% "Patching failed" error rather than silently doing nothing): recent
% LaTeX2e kernels implement an optional-argument environment like
% beamer's own `\renewcommand{\itemize}[1][]{...}` via
% `\@protected@testopt`, which stores the *real* body not under `\itemize`
% itself but under a differently-named internal control sequence (`\show
% \itemize` confirms this directly: it prints as
% `\@protected@testopt \itemize \\itemize {}`, i.e. plain `\itemize` is
% only a thin dispatcher, and the actual `\ifnum \@itemdepth >2...` body
% \patchcmd was trying to match against was never there to find). Rather
% than chase that internal storage name (liable to differ across LaTeX
% kernel versions, which \patchcmd's own documented usage doesn't need to
% care about), \itemize is fully re-declared here instead, copied
% verbatim from beamer's own beamerbaselocalstructure.sty (this project's
% installed version) with only the depth bound itself changed from 2 to
% 3 -- the same declaration form beamer itself used
% (\renewcommand{\itemize}[1][]{...}), which regenerates both the
% dispatcher and the real body correctly under whatever internal name
% the current kernel expects, sidestepping the naming question entirely.
% Trade-off, stated plainly: if a future beamer release changes this
% macro's own body for an unrelated reason, this override will silently
% keep using today's (2026-07) copy instead of picking up that change --
% acceptable here since the override is one line different from upstream
% and clearly marked, but re-diff against beamerbaselocalstructure.sty's
% own \itemize if beamer is ever upgraded and itemize behavior looks off.
% \SMNoListParsep (2026-10-01, direct instruction: "Implement a general
% fix"): a level-1 list typeset at \small or \footnotesize -- an ambient size
% switch before it, \setbeamerfont, \shrinklists -- loads LaTeX's
% size-specific \@listi, which brings \parsep 3pt plus 2pt (\footnotesize:
% similar) and \parskip with it; a normal-size list has neither. So a smaller
% list's items sat *further* apart than at normal size (6j4: 19.9pt per item
% at normal size, 20.8pt at \small), against this deck's rule that list
% spacing is em-based and shrinks with the font (\SMItemsepRatioI etc.).
% Removed wherever \parsep has a natural (non-stretch) part, i.e. exactly in
% those lists; deeper levels (0pt plus 1pt) are left as they were.
\makeatletter
\newcommand{\SMNoListParsep}{%
	\ifdim\parsep>\z@
		\setlength{\parsep}{\z@}\setlength{\parskip}{\z@}%
	\fi
}
\long\def\beamer@computepref#1{%
	\let\beameritemnestingprefix\@empty%
	\ifcase#1\or\or\def\beameritemnestingprefix{sub}\or\def\beameritemnestingprefix{subsub}\or\def\beameritemnestingprefix{subsubsub}\or\@toodeep\fi%
}
\renewcommand{\itemize}[1][]{%
	\ifblank{#1}{}{\def\beamer@defaultospec{#1}}%
	\ifnum \@itemdepth >3\relax\@toodeep\else
		\advance\@itemdepth\@ne
		\beamer@computepref\@itemdepth% sets \beameritemnestingprefix
		\usebeamerfont{itemize/enumerate \beameritemnestingprefix body}%
		\usebeamercolor[fg]{itemize/enumerate \beameritemnestingprefix body}%
		\usebeamertemplate{itemize/enumerate \beameritemnestingprefix body begin}%
		\list
			{\usebeamertemplate{itemize \beameritemnestingprefix item}}
			{\def\makelabel##1{%
					{%
						\hss\llap{{%
								\usebeamerfont*{itemize \beameritemnestingprefix item}%
								\usebeamercolor[fg]{itemize \beameritemnestingprefix item}##1}}%
					}%
				}%
			}
	\fi%
	\beamer@cramped%
	% \beamer@cramped (beamerbaselocalstructure.sty, unmodified upstream
	% code) unconditionally does `\itemsep0pt\topsep0.5pt`, every single
	% time -- called right here, in this project's own copy of beamer's
	% own \itemize body, same position as upstream. Setting \itemsep any
	% earlier (first tried as a separate \AtBeginEnvironment{itemize} hook,
	% run *before* this whole macro body even starts) gets silently wiped
	% out by this same line a moment later -- confirmed directly via
	% \typeout debugging (not assumed): the hook's own \setlength really
	% did run with the right values, \itemsep still read back as 0pt
	% immediately after. Setting it *here*, right after \beamer@cramped
	% already ran (not before), is what actually survives.
	\ifcase\@itemdepth
	\or\setlength{\itemsep}{\SMItemsepRatioI em}% level 1
	\or\setlength{\itemsep}{\SMItemsepRatioII em}% level 2
	\or\setlength{\itemsep}{\SMItemsepRatioIII em}% level 3
	\or\setlength{\itemsep}{\SMItemsepRatioIV em}% level 4
	\fi%
	\SMNoListParsep
	\raggedright%
	\beamer@firstlineitemizeunskip%
}
\makeatother

% A nested \begin{enumerate} inherits the same \leftmarginN step as itemize
% above for *where its text starts* (matching itemize's own "one step past
% the parent level's text" rule -- confirmed directly against the compiled
% PDF: an enumerate's own body text, e.g. "Senior secured debt", already
% lands exactly one \SMItemStep past its parent item's text), but its
% *label* ("1.", "2.", ...) doesn't follow \labelwidth the way plain LaTeX's
% list mechanism would -- beamer draws enumerate labels through its own
% template machinery, which turned out not to respect \labelwidth at all:
% \setlength{\labelwidth}{...} right after \begin{enumerate} measurably took
% effect (checked with \typeout{\the\labelwidth} immediately after) but the
% compiled numeral's own position never moved. What *does* move it: directly
% overriding \makelabel with a fixed-width \makebox -- empirically
% calibrated (render, measure the numeral's own x-position in the compiled
% PDF, adjust, repeat) since the box's own left edge turned out to sit at a
% fixed position unrelated to \labelwidth or \leftmarginN, with the
% right-aligned numeral inside it moving in lockstep with the box's width
% (confirmed linear: +1pt width -> +~1pt numeral position, across several
% test widths). 19.23pt (at this deck's default itemize font size) lands a
% single-digit numeral flush with the parent item's own text -- only valid
% at that font size and for single-digit lists (1-9); re-calibrate the same
% way (not by recomputing from first principles) for a differently-sized or
% 10+-item nested enumerate. Call right after \begin{enumerate}, before the
% first \item.
%
% 20.33pt (2026-07-08, direct instruction: "the numbers still extend a
% little farther to the left than the text above them does" --
% slides/f6m-Recovery-Rates-current.tex's own "Pecking order:" list).
% 19.23pt landed *close* but not quite flush: measured directly against
% the compiled PDF, every numeral's own left edge (20.55pt-20.95pt,
% varying slightly by digit width even though all are right-aligned to
% the same box) still sat up to ~1.05pt left of "Pecking order:"'s own
% text start (21.6pt) -- the widest digit tried, "6.", was the binding
% case. +1.1pt of box width (confirmed +1pt width -> +1pt numeral
% position, per this comment's own note above) covers that with a small
% margin, re-confirmed against the recompiled PDF.
\newcommand{\SMEnumerateLabelWidth}{\renewcommand{\makelabel}[1]{\makebox[\SMEnumerateLabelBox][r]{##1}}}

% Automatic version of the above (2026-07-23, background-task follow-up
% flagged in slides/kfh-Asset-Allocation-Process-current.tex's own v1-1-0
% comment: "fix the enumerate behavior generally, if possible", task
% `task_bb3fbb66`). That comment's first attempt, an
% `\AtBeginEnvironment{enumerate}{...}` hook, failed because the hook fires
% *before* `\enumerate`'s own body runs, and that body's own internal
% `\list{...}{\def\makelabel...}` setup (the beamer copy directly below the
% `\itemize` override above) clobbers whatever a hook that early leaves
% behind -- the exact same "template runs its own setting, external pre-set
% gets overwritten" problem `\itemize`'s own `\itemsep` comment above
% already documents. The real fix, same idea `\itemize` itself used above:
% override beamer's own low-level piece directly rather than hooking
% \AtBeginEnvironment. Unlike `\itemize` (declared via
% `\@protected@testopt`, needing a full re-declaration to reach -- see the
% `\itemize` override's own comment for the full story), `\enumerate`'s
% relevant piece, `\beamer@enum@` (beamerbaselocalstructure.sty), is a
% plain `\def` with no optional-argument indirection, so it can be
% `\renewcommand`'d directly, verbatim from upstream except for the
% `\ifnum\@itemdepth>1` branch added below.
%
% Only fires for a *nested* enumerate (`\@itemdepth>1`, i.e. already inside
% at least one `\itemize`/`\enumerate` when this one opens) -- a top-level
% enumerate (`\@itemdepth=1`, e.g.
% slides/t9z-Outland-Steel-Hedging-Borrowing-Investing-current.tex/
% slides/g4c-Outland-Steel-Hedging-Forward-Sale-current.tex/
% slides/j84-Outland-Steel-No-Hedging-current.tex's own outermost numbered
% list, or slides/nkj-Agenda-current.tex's own \tableofcontents-driven one)
% is left at beamer's stock label position entirely untouched, matching
% this project's own established behavior for that shape (never reported
% as misaligned, unlike the nested case) -- and, for nkj specifically,
% doesn't matter either way, since that slide's own
% `\SMTocLabelFlushLeft` call (right after its own `\begin{enumerate}`)
% overrides `\makelabel` again regardless of whatever this default sets.
%
% Reuses the exact same 20.33pt calibration as the manual
% `\SMEnumerateLabelWidth` above, confirmed (pymupdf, compiled PDF) to
% land correctly not just at `\@itemdepth=2` (one `\itemize` above the
% enumerate -- slides/kfh-Asset-Allocation-Process-current.tex/
% slides/f6m-Recovery-Rates-current.tex/
% slides/6pz-Outland-Steel-Hedging-Strategies-Overview-current.tex, all
% previously needing the manual call) but also at `\@itemdepth=3` (two
% `\itemize`s above it -- slides/bdh-Working-Capital-Example-current.tex's
% own "Development of balance sheet:" list, also previously needing the
% manual call, with the identical 20.33pt value) -- this deck's itemize/
% enumerate templates don't shrink font size by nesting depth (only the
% bullet glyph/margin step does), so the numeral's own glyph width -- the
% only thing \SMEnumerateLabelWidth's own box width actually
% compensates for -- stays the same at every depth, and one constant covers
% both cases without needing a depth-dependent formula. Not verified beyond
% depth 3 (no slide currently nests an enumerate that deep); if one ever
% does, re-measure before trusting this constant that far down, same as any
% other empirical calibration in this file.
% (slides/scs-Political-Risk-Stability-Index-current.tex also matches
% "`\SMEnumerateLabelWidth`" in a grep -- but only inside its own prose
% comments describing the general technique; its actual call is the
% bespoke `\SMEnumerateLabelFlushLeftShifted{16pt}{10.2pt}` for its own
% alphabetic, left-aligned, narrower-box labels, unaffected by this
% override either way since that call runs after `\begin{enumerate}` opens,
% same ordering as every other post-open `\makelabel` override here.)
%
% Former per-slide manual `\SMEnumerateLabelWidth` calls removed from
% the 4 slides above now that this covers them automatically; the macro
% itself is kept defined, both for this override to reuse and as a manual
% escape hatch (e.g. a differently-calibrated width some future slide might
% need).
% (v9.8.0: \SMEnumerateTextExtra, 6.3pt, is a hook of the catalog; the four lengths below were numbers of the
% layout profile: 7.85pt, 13.1pt, 5.25pt and 20.33pt for Fira Sans, 8.18pt, 13.43pt, 5.25pt and 20.33pt for
% Latin Modern Sans.)
% The numeral column of a nested enumerate, from the widest numeral of the list's own font (the digits
% 1-9 with the dot, which is what the item template sets). The slot a numeral sits in is that width plus
% \SMEnumerateNestedSlotExtra, the text starts \SMEnumerateNestedLabelGap after it
% (\SMEnumerateNestedLeftMargin is both), and the box the numeral is right-aligned in only has to hold a
% numeral of three digits ("999."), so it takes the width of that. Called first in the list's setup of
% every enumerate, and once at \begin{document} (in \small) for a slide that reads the lengths itself.
\makeatletter
\newcommand{\SM@enumeratemetrics}{%
	\setlength{\SMEnumerateNestedMinNumber}{0pt}%
	\setlength{\SMEnumerateLabelBox}{0pt}%
	\@for\SM@numprobe:=1,2,3,4,5,6,7,8,9\do{%
		\settowidth{\dimen@}{\SM@numprobe.}%
		\ifdim\dimen@>\SMEnumerateNestedMinNumber\SMEnumerateNestedMinNumber=\dimen@\fi
		\settowidth{\dimen@}{\SM@numprobe\SM@numprobe\SM@numprobe.}%
		\ifdim\dimen@>\SMEnumerateLabelBox\SMEnumerateLabelBox=\dimen@\fi}%
	\addtolength{\SMEnumerateNestedMinNumber}{\SMEnumerateNestedSlotExtra}%
	\setlength{\SMEnumerateNestedLabelSep}{\SMEnumerateNestedLabelGap}%
	\setlength{\SMEnumerateNestedLeftMargin}{\dimexpr\SMEnumerateNestedMinNumber+\SMEnumerateNestedLabelSep\relax}}
\makeatother
\makeatletter
\renewcommand{\beamer@enum@}{%
	\beamer@computepref\@itemdepth% sets \beameritemnestingprefix
	\usebeamerfont{itemize/enumerate \beameritemnestingprefix body}%
	\usebeamercolor[fg]{itemize/enumerate \beameritemnestingprefix body}%
	\usebeamertemplate{itemize/enumerate \beameritemnestingprefix body begin}%
	% Level-1 enumerate text sits further right than bullet text, as in the source
	% (numeral at the bullet dot's left edge, item text 27.9pt from the page's left edge
	% against 21.6pt for a bullet's text; 2026-09-20). Nested lists follow it, as they do for
	% plain bullets: a level-2 bullet starts under the level-1 text (2026-09-21, direct
	% instruction -- an earlier version gave the extra back to the nested list, which left
	% level-2 bullets 6.3pt left of the level-1 text above them). Local to this list.
	\ifnum\@itemdepth=1
		\addtolength{\leftmargini}{\SMEnumerateTextExtra}%
	\fi
	\expandafter
		\list
			{\usebeamertemplate{\beamer@enumtempl}}
			{\usecounter\@enumctr%
				\SM@enumeratemetrics% the numeral column from this list's font (v9.8.0)
				\ifnum\@itemdepth>1
					% Nested numerals: same rule as level 1 -- left edge at the left
					% edge of this level's bullet dot (2026-09-20). The box below is
					% the old 20.33pt right-aligned one, whose right end the numerals
					% used to fill flush; giving the numeral a 7.85pt left-aligned
					% slot at that end puts every numeral's left edge exactly at the
					% dot's, instead of up to 1pt right of it for narrow digits
					% ("1." vs "4."). A wider numeral keeps the old right alignment.
					% Hanging indent (2026-09-21, direct instruction: the second line of
					% an item did not start under its first line's text). The 20.33pt label
					% box is wider than beamer's \labelwidth, so it pushed the first
					% line's text 4.3pt right of \leftmargin, where every later line
					% starts. Sizing \labelwidth to the box and \leftmargin to
					% "numeral slot + \labelsep" puts the numeral's left edge at the dot's
					% and the text of every line at the same x (13.1pt from the parent's
					% text, i.e. where the first line always was).
					\setlength{\labelwidth}{\SMEnumerateLabelBox}%
					\setlength{\labelsep}{\SMEnumerateNestedLabelSep}%
					\setlength{\leftmargin}{\SMEnumerateNestedLeftMargin}%
					\def\makelabel##1{\makebox[\SMEnumerateLabelBox][r]{{%
							\setbox\@tempboxa\hbox{##1}%
							\ifdim\wd\@tempboxa>\SMEnumerateNestedMinNumber
								\box\@tempboxa
							\else
								\hbox to \SMEnumerateNestedMinNumber{\unhbox\@tempboxa\hss}%
							\fi}}}%
				\else
					% Level-1 numerals start at the left edge of the level-1 bullet
					% dot, like the source's numerals (2026-09-20, direct
					% instruction: the numbers were too far left). They used to
					% be right-aligned to the dot's right edge, which for a
					% numeral wider than the dot (every "1.") puts its left edge
					% left of where a bullet would start. A numeral too wide to
					% keep 1pt clear of the item text ("10.") keeps the old
					% right alignment instead of running into the text.
					\def\makelabel##1{{\hss\llap{{%
							\usebeamerfont*{enumerate \beameritemnestingprefix item}%
							\usebeamercolor[fg]{enumerate \beameritemnestingprefix item}%
							\setbox\@tempboxa\hbox{##1}%
							\ifdim\dimexpr\wd\@tempboxa+1pt\relax>\dimexpr\SMItemStep+\SMEnumerateTextExtra\relax
								\box\@tempboxa
							\else
								\hbox to \dimexpr\SMBulletWidth+\SMEnumerateTextExtra\relax{\rlap{\unhbox\@tempboxa}\hss}%
							\fi}}}}%
				\fi}%
	\beamer@cramped%
	\SMNoListParsep
	\raggedright%
	\beamer@firstlineitemizeunskip%
}
% The lists' lengths as a slide that reads them outside a list sees them: those of the \small font.
\AtBeginDocument{%
	\begingroup\small\SM@labelitemmetrics\SM@enumeratemetrics
	\xdef\SM@listmetrics{%
		\global\SMLabelItemWidth=\the\SMLabelItemWidth\relax
		\global\SMLabelItemSep=\the\SMLabelItemSep\relax
		\global\SMEnumerateNestedMinNumber=\the\SMEnumerateNestedMinNumber\relax
		\global\SMEnumerateNestedLabelSep=\the\SMEnumerateNestedLabelSep\relax
		\global\SMEnumerateNestedLeftMargin=\the\SMEnumerateNestedLeftMargin\relax
		\global\SMEnumerateLabelBox=\the\SMEnumerateLabelBox\relax}%
	\expandafter\endgroup\SM@listmetrics}
\makeatother

% Extra breathing room when a nested itemize closes and control returns to
% its *parent* list, right before that parent's next \item -- e.g. sub-items
% under "Mid-Term Exam" followed by the next top-level "Final Exam" -- so a
% new top-level group reads as visually distinct from the one before it,
% rather than the same uniform gap used between every other item regardless
% of level. \@itemdepth (LaTeX's own nesting-level counter, already
% decremented back to the parent's level by the time \end{itemize}'s normal
% closing has run) says which level control returned to -- 0 means the
% outermost list itself just closed (nothing left to add space before,
% deliberately excluded), 1 means back to a level-1 item, 2 means back to a
% level-2 item.
%
% Level-dependent (2026-07-06, direct instruction) -- a flat, single
% constant regardless of *which* level control returned to (the original
% version of this hook) reads as "extra"/inconsistent precisely because a
% return-to-level-1 and a return-to-level-2 transition both got the exact
% same gap, rather than level 1 getting a visibly bigger gap than level 2 --
% found on slides/57r-Reasons-for-Leasing-current.tex, where a level-2
% sub-list closing back to a level-1 item looked like a random extra gap
% rather than "level 1 items are consistently more separated than level 2
% ones." \addvspace (not plain \vspace) matters too: it coalesces with
% whatever glue itemize's own \itemsep/\topsep already contributes at the
% closing boundary (taking whichever is larger) rather than *adding* on top
% of it unconditionally -- \vspace's own stacking behavior was the other
% half of why the old flat version could look like too much space.
\makeatletter
% Ratios, not lengths (2026-07-21, direct instruction: "all of these
% spacings should scale with font size") -- a \newlength+\setlength{...}{0.5em}
% pair, the way this used to be written, freezes to a fixed *pt* value the
% one time \setlength runs (at preamble-load time, using whatever font is
% active then), despite the "em" in its own definition -- it does NOT
% re-track whichever font size is actually active on a given slide/level
% later. Confirmed directly: a \typeout{\f@size} placed at the point each
% of these actually gets used (inside \AtEndEnvironment{itemize}'s own
% hook body, further below) showed the *closing* list's own font still
% active there (9pt/10pt/10.95pt for a level-4/3/2 list closing,
% respectively -- not yet restored to the parent's own font), confirming
% this hook fires early enough in `\end{itemize}`'s own processing that a
% live "0.5em" computed *there* would correctly track whichever level is
% actually involved. Fixed by storing the bare ratio as a plain macro
% instead, and writing the literal "em" back on at the point of use
% (`\addvspace{\SMItemizeLevelDropSpaceRatioI em}`), so it's
% recomputed fresh, against whatever font is live, every time.
% Level 3's own drop-space, for when a closed *4th*-level list returns
% control to a level-3 item -- added 2026-07-15 alongside the 4th itemize
% level itself (search "subsubsub" further above). Interpolated from the
% existing I/II progression (0.5, 0.3 -- each step roughly 60% of the one
% before) rather than empirically re-measured against a compiled PDF the
% way I/II themselves were: no slide used a 4th level before
% slides/hrg-Exchange-Rate-Quotation-Forward-current.tex to measure
% against. Re-derive properly (render, measure, adjust) if this ever
% looks visibly off on a real slide.
% Level 4's own drop-space, for when a closed level-4 list returns control
% to a level-3 item -- added 2026-07-21 alongside \SMItemsepRatioIV.
% \@itemdepth at this hook's fire time is the depth of the list that *just
% closed*, confirmed directly via a temporary \typeout{\the\@itemdepth}
% (not assumed): closing the level-4 list on
% slides/ZZO-Itemize-Spacing-Reference-current.tex measured itemdepth=4
% there, closing level-3 measured 3, level-2 measured 2, and the outermost
% level-1 list itself closing measured 1 (never 0 -- \@itemdepth's own
% restore to the parent's depth happens *after* this hook runs, so the "0:
% outermost list itself closed" case below is dead code, kept only because
% removing a genuinely unreachable \ifcase branch isn't worth the risk of
% getting the numbering wrong a second time). So this is very much a real,
% reachable case -- not decorative symmetry -- and needed its own
% \addvspace branch (previously missing entirely, silently leaving the
% level-4-to-level-3 closing transition with no drop-space at all).
% Interpolated from the existing I/II/III progression as a starting guess,
% then re-measured directly against the compiled ZZO reference and bumped:
% 0.13 measured at 12.6pt, still *less* than level 4's own same-level gap
% of 13.4pt (violates rule III -- a level's own same-level gap must be
% smaller than that level's own closing-transition gap); raised until
% re-measurement cleared it with a comfortable margin again, matching the
% margin every other level's own pair already has.
\AtEndEnvironment{itemize}{%
	% Force vertical mode before \addvspace below (2026-09-22, found while
	% isolating a "Something's wrong--perhaps a missing \item" error that
	% fired at the *frame's* own \end{...}, not the itemize's, on every
	% slide using a plain, single-level \begin{itemize}...\end{itemize}
	% with nothing else in the frame -- confirmed via a from-scratch minimal
	% reproduction outside this project's own slides entirely, then
	% bisected to this exact hook by diffing against the same reproduction
	% with this \AtEndEnvironment call disabled). \AtEndEnvironment inserts
	% this code at the very *start* of \enditemize, before beamer's own
	% `\ifhmode\unskip\fi\endlist` -- so without this guard, \addvspace
	% below runs while TeX is often still in horizontal mode (right after
	% the last \item's own text, mid-paragraph). \addvspace's own hmode
	% branch defers via `\vadjust{\addvspace{...}}` instead of inserting
	% the glue directly, and that deferred insertion colliding with
	% beamer's own frame-body/overlay-replay collection is what actually
	% produced the corruption -- not reproducible as a plain "wrong
	% \item" typo anywhere in the affected slides' own source, only via
	% this exact hmode-vs-vmode timing. \ifhmode\unskip\par\fi here forces
	% vertical mode *before* \addvspace ever runs, matching the standard,
	% documented way to call \addvspace (kernel doc: "should only be used
	% between paragraphs") -- same idiom beamer's own \enditemize already
	% uses \unskip for, just moved earlier so it also covers this hook.
	% Doesn't disturb the \@itemdepth-based \ifcase dispatch just below:
	% \par doesn't touch \@itemdepth, so the closing list's own depth is
	% still read correctly afterward, same as before this guard existed.
	\ifhmode\unskip\par\fi
	\ifcase\@itemdepth
		% 0: dead code -- @itemdepth is never actually 0 when this hook
		% fires (see \SMItemizeLevelDropSpaceRatioIV's own comment
		% just above); kept only for alignment with the \ifcase branches
		% below, each keyed on the depth of the list that just closed.
	\or\addvspace{\SMItemizeLevelDropSpaceRatioI em}% level-1 (outermost) list itself closed -- a pre-existing quirk (not introduced by this comment's own edit): there's no parent item above level 1 for this to visually separate from, so this only actually matters if something follows the outermost list directly on the same frame; harmless on every slide seen so far since nothing does. Out of scope to redesign here.
	\or\addvspace{\SMItemizeLevelDropSpaceRatioII em}% level-2 list closed, back to a level-1 item
	\or\addvspace{\SMItemizeLevelDropSpaceRatioIII em}% level-3 list closed, back to a level-2 item
	\or\addvspace{\SMItemizeLevelDropSpaceRatioIV em}% level-4 list closed, back to a level-3 item
	\fi
}

% Per-level \itemsep -- the item-to-item gap *within* one list, at one
% fixed nesting depth (as opposed to the drop-space mechanism just above,
% which only fires when a nested list *closes* back up to its parent).
% Ratios, not lengths -- same "scale with font size" reasoning and fix as
% \SMItemizeLevelDropSpaceRatioI/II/III above (that comment has the
% full story on why a plain \newlength+\setlength pair here would freeze
% to one fixed pt value instead of actually tracking each level's own
% font size). Values only, defined here for discoverability alongside the
% drop-space ones above -- actually *applied* inside the `\itemize`
% redefinition itself (search "\beamer@cramped (beamerbaselocalstructure.sty"
% further above), not via a separate hook here; see that comment for why.
% Found missing entirely (2026-07-21, direct instruction, see
% slides/ZZO-Itemize-Spacing-Reference-current.tex's own header comment
% for the rules this and the drop-space mechanism together need to
% satisfy): this deck's own itemize never set \itemsep itself at any
% level -- confirmed directly, grepping this whole file -- so every depth
% shared whatever bare default beamer's own theme happens to leave
% \itemsep at, uniformly. That produced two visible problems, both
% measured directly (pymupdf against slides/2c2-Topics-Around-the-CAPM-Empirical-Validity-current.tex
% and an isolated bisection test, not guessed): (a) level-2 and level-3
% same-level item-to-item gaps came out *identical* (12.6pt each,
% confirmed against an isolated reproduction), since nothing distinguished
% them by depth at all -- only a closing transition ever got a
% level-dependent gap, never a same-level one; (b) that shared default
% \itemsep read as *smaller* than a single item's own wrapped-line
% \baselineskip (13.7-13.8pt within a wrap vs. 12.6pt between items),
% backwards from how a list should read -- an item boundary should never
% be tighter than a same-item line wrap.
% Ratios themselves are not arbitrary -- each chosen, then re-measured
% against the compiled ZZO reference slide and adjusted, specifically to
% satisfy all three of that slide's own header-comment rules together
% (extended to level 4 the same way, 2026-07-21), not just "level 1
% biggest" in isolation:
% - Rule I (level-1 same-level gap > level-2's > level-3's > level-4's)
%   falls straight out of \SMItemsepRatioI > ...II > ...III > ...IV,
%   added on top of the same \baselineskip every level shares (itself
%   already smaller per level here, since this theme's own itemize/
%   enumerate body font shrinks with depth too -- confirmed directly,
%   10.95pt/10pt/9pt/9pt for levels 1-4 respectively, level 4 sharing
%   level 3's own font since this theme has no distinct "subsubsub" size
%   of its own).
% - Rule II (every level's own same-level gap > that level's own wrapped-
%   line gap) needs \SMItemsepRatioIV alone (the smallest, hence the
%   binding case) to still be positive and large enough that it clearly
%   reads as more open than a plain wrapped line once rendered, not just
%   bigger on paper.
% - Rule III (a level's own same-level gap < that same level's own
%   *closing*-transition gap) is *not* automatic just because
%   \SMItemizeLevelDropSpaceRatio* was already bigger than the old
%   shared default -- \addvspace's own coalesce-to-the-larger-value
%   behavior (documented above) means the closing gap is
%   max(natural glue, that *drop-space* ratio), not a sum of the two. Each
%   drop-space ratio's own \ifcase branch is keyed on the depth of the
%   list that *just closed* (see \SMItemizeLevelDropSpaceRatioIV's own
%   comment for the full story on why, confirmed via \typeout, not
%   assumed), so it lines up with the same-numbered itemsep ratio despite
%   describing a different transition -- \SMItemizeLevelDropSpaceRatioIV
%   governs the level-4-list-closing-back-to-level-3 transition and is the
%   one that has to individually stay bigger than \SMItemsepRatioIV.
%   This branch was missing entirely at first (only I/II/III existed, a
%   leftover from before level 4 existed) -- confirmed directly on ZZO: the
%   level-4-to-level-3 closing transition measured 12.6pt with no drop-space
%   branch at all, less than level 4's own same-level gap of 13.4pt, a real
%   rule III violation. Adding the branch (\SMItemizeLevelDropSpaceRatioIV,
%   starting at 0.5) fixed it. Re-verify rule III directly against ZZO's own
%   compiled PDF (don't assume it still holds on paper) if any of these ratios is
%   ever changed again.

% Per-level "opening" space -- the gap between a level-n item and the
% first item of a level-(n+1) list that immediately follows it, as
% opposed to \SMItemizeLevelDropSpaceRatio* (a nested list *closing*
% back up) or \SMItemsepRatio* (two items at the *same* depth).
% Added 2026-07-21, direct instruction, replacing the original three
% separate rules on slides/ZZO-Itemize-Spacing-Reference-current.tex's
% own header comment with one multi-part inequality stated per level n:
%   level-n -> level-(n-1) close  >  level-n same-level  >
%   level-n -> level-(n+1) open  >  level-n's own wrap gap
% The "open" term was entirely unhandled before this -- confirmed
% directly on ZZO: every opening transition (level 1 into level 2, level
% 2 into level 3, level 3 into level 4) measured *smaller* than that same
% level's own wrap gap, backwards from what the rule needs, since nothing
% but beamer's own bare default \topsep governed it.
% \AtBeginEnvironment{itemize}, not \topsep -- \topsep was tried first
% (set via the same \ifcase\@itemdepth block, right after \itemsep,
% inside the `\itemize` redefinition above) and measured *unreliable*:
% two of the three levels barely moved despite a sizeable ratio, while
% the third moved by roughly the right amount -- \topsep's own value is
% evidently consulted by `\list` at some point this override comes too
% late for, at least for some levels, not evenly. \addvspace via
% \AtBeginEnvironment instead measured consistently and predictably
% (each level's own change matched its own ratio's calculated pt value
% directly), the same mechanism \SMItemizeLevelDropSpaceRatio* already
% relies on for the closing side. This hook fires with \@itemdepth
% reflecting the *parent* item's own depth (confirmed via \typeout, not
% assumed: opening the outermost level-1 list itself measured 0, opening
% a level-2 list from within a level-1 item measured 1, level-3-from-
% level-2 measured 2, level-4-from-level-3 measured 3) -- the opposite of
% \SMItemizeLevelDropSpaceRatio*'s own hook, which keys on the depth
% of the list that *just closed*. That makes the \ifcase numbering here
% match the "n" in the rule directly: case 1 is level-1's own open-to-
% level-2 ratio, case 2 is level-2's own open-to-level-3, case 3 is
% level-3's own open-to-level-4; case 0 (opening the outermost list
% itself) is a genuine no-op, correctly -- there's no parent item above
% level 1 for that transition to separate from.
% Ratios chosen, then re-measured against the compiled ZZO reference and
% adjusted, the same way as every other ratio here: each has to
% individually land strictly between that level's own \SMItemsepRatio*
% (the ceiling) and its own wrap gap (the floor). Re-verify directly
% against ZZO's own compiled PDF (don't assume it still holds on paper)
% if any of these, or the itemsep/drop-space ratios they're measured
% against, are ever changed again.
% \SMItemizeLevelOpenSpaceRatioI trimmed from 0.28 to 0.02 (2026-07-21,
% direct instruction: "the level-1 to level-2 distance...seems a bit
% much"). Confirmed directly: even at 0, this transition still measured
% ~16.3pt -- \addvspace's own coalesce-to-the-larger-value behavior means
% the ratio can only ever push this *above* whatever glue the list
% mechanism inserts here regardless (that natural floor, not the ratio),
% so 0.28 down to 0 only ever moved the actual gap by ~2pt (18.5->16.3);
% 0.02 splits the difference, keeping the mechanism nominally "in effect"
% rather than reading as accidentally disabled. \SMItemizeLevelOpenSpaceRatioIII
% (level-3-to-level-4) had to come down too, from 0.08 to 0.02, purely to
% restore its own margin -- changing RatioI shifted that unrelated
% transition's own measured gap closer to level 3's own same-level gap
% (12.1 vs. 12.0, too thin), confirmed directly on ZZO, not assumed;
% ratios elsewhere in this file have shown the same small-ripple behavior
% before (see \SMItemizeLevelDropSpaceRatioIV's own comment) --
% re-measure the *whole* table after touching any single one of these,
% not just the transition being changed.
\AtBeginEnvironment{itemize}{%
	% Same \ifhmode\unskip\par\fi guard as \AtEndEnvironment{itemize} above,
	% and for the identical reason -- see that hook's own comment for the
	% full story. This hook fires when a *nested* itemize opens (\@itemdepth
	% here is the *parent* item's depth, per the comment below), i.e. right
	% after "\item ...text..." mid-paragraph, immediately before the
	% \addvspace calls further down -- the same hmode-vs-vmode timing that
	% broke the closing hook, just on the opening side. Case 0 (opening the
	% outermost level-1 list itself) never reaches an \addvspace call below
	% regardless, so a plain single-level itemize was never affected by
	% this half of the bug -- only a *nested* \begin{itemize} (depth >=1 at
	% open time) actually needed this guard, which is why the "missing
	% \item" error survived this fix's own first pass specifically on the
	% two demo slides with a nested itemize (sectionframe's own ZZR,
	% standardframe's own ZZZ) while every plain single-level-itemize slide
	% was already fixed by the closing-hook guard alone.
	\ifhmode\unskip\par\fi
	\ifcase\@itemdepth
		% 0: opening the outermost (level-1) list itself -- no parent item
		% precedes it, nothing to add.
	\or\addvspace{\SMItemizeLevelOpenSpaceRatioI em}% opening a level-2 list from a level-1 item
	\or\addvspace{\SMItemizeLevelOpenSpaceRatioII em}% opening a level-3 list from a level-2 item
	\or\addvspace{\SMItemizeLevelOpenSpaceRatioIII em}% opening a level-4 list from a level-3 item
	\fi
}
\makeatother


% Even with the plain-enumerate template (SMPlainEnumerate, defined in the theme), \tableofcontents still
% doesn't match this deck's own normal itemize spacing -- beamer's own
% \beamer@sectionintoc (beamerbasetoc.sty), which wraps *every* section
% entry regardless of which "section in toc" template is active,
% unconditionally inserts a fixed \vskip1.5em before each entry (an
% unbounded \vfill, even worse, if subsections aren't hidden -- neither is
% reachable via a public \tableofcontents option, just this show/hide
% binary), *and* wraps the entry itself in \hbox{\vbox{...}} -- box
% wrapping that (like \SMMinipageTopCorrection's own minipage case
% elsewhere in this file) changes how surrounding baselineskip/topskip
% glue gets inserted compared to plain paragraph text. Both together added
% roughly 13pt of extra gap
% between entries beyond this deck's own normal ~18.6pt level-1 item
% spacing (measured directly against
% slides/57r-Reasons-for-Leasing-current.tex's own compiled PDF).
% Neither is controllable via a public \tableofcontents option, so this
% redefines \beamer@sectionintoc itself -- scoped safely since
% \tableofcontents is only ever called on the one Agenda slide in this
% deck -- keeping its numbering/hyperlink/part-number/hide-checking logic
% intact but dropping the vskip/vfill and the hbox/vbox wrapping, so an
% entry behaves exactly like a plain \item once printed.
\makeatletter
\def\beamer@sectionintoc#1#2#3#4#5{%
	\ifnum\c@tocdepth>0%
	\ifnum#4=\beamer@showpartnumber%
	{%
	\beamer@saveanother%
	\gdef\beamer@todo{}%
	\beamer@slideinframe=#1\relax%
	\expandafter\only\beamer@tocsections{\gdef\beamer@todo{%
			\beamer@tempcount=#5\relax%
			\advance\beamer@tempcount by\beamer@sectionadjust%
			\ifnum\beamer@tempcount>0
				\ifnum\beamer@tempcount>\beamer@toclastsection
					\def\inserttocsectionnumber{}%
				\else
					\edef\inserttocsectionnumber{\the\beamer@tempcount}%
				\fi
			\else
				\def\inserttocsectionnumber{}%
			\fi%
			\def\inserttocsection{\hyperlink{Navigation#3}{#2}}%
			\beamer@tocifnothide{\ifnum\c@section=#1\beamer@toc@cs\else\beamer@toc@os\fi}%
			{%
				\beamer@tocact{\ifnum\c@section=#1\beamer@toc@cs\else\beamer@toc@os\fi}{section in toc}%
			}%
		}%
	}%
	\beamer@restoreanother%
	}%
	\beamer@todo%
	\fi\fi%
}
\makeatother

% Left-aligns an enumerate's labels within a caller-chosen \labelwidth --
% used by the \tableofcontents Agenda slide (see
% slides/nkj-Agenda-current.tex's own comment for why its labels need
% this at 3.6pt). Defined here as a named helper rather than written as a
% literal \renewcommand{\makelabel}[1]{...#1...} inside the slide's own
% frame body, because that literal form -- specifically its raw `#1`
% parameter token -- breaks inside any frame that \tableofcontents forces
% through beamer's overlay reprocessing: beamer replays such a frame's
% collected body via \def-based machinery (\loop's own \iterate, and
% \beamer@doifinframe), and a `#` in the material of any \def must be
% doubled -- so the compile hits "! Illegal parameter number in definition
% of \iterate." (twice) plus the same for \beamer@doifinframe, once per
% compile, at that frame's own \end{frame}. TeX's error recovery ("I'll
% assume you meant ##") happens to produce the intended definition anyway,
% which is why the rendered PDF was always correct despite the errors --
% but lualatex still exits nonzero, which is exactly what kept breaking
% latexmk's automated lualatex/biber/rerun sequencing on the main deck
% (the "Known open issue" long tracked in CONVERSION_NOTES.md, originally
% misattributed to a biblatex interaction -- see its "Resolved" entry).
% Hoisting the definition out here means the frame body carries only this
% macro's *call* token, with no `#` in it, which sidesteps the whole
% mechanism. The inner \makelabel parameter is written `##1` below since
% it sits inside another \newcommand's body -- standard nesting, and also,
% incidentally, exactly the doubling rule the inline version fell foul of.
\newcommand{\SMTocLabelFlushLeft}[1]{%
	\setlength{\labelwidth}{#1}%
	\renewcommand{\makelabel}[1]{\makebox[\labelwidth][l]{##1}}%
}

% Same idea as \SMTocLabelFlushLeft above, plus an extra fixed
% \hspace shift before the label (for a list that needs to start further
% right than the plain nested-list indent, e.g. to line up with a
% sibling item's own text rather than its bullet -- see
% slides/scs-Political-Risk-Stability-Index-current.tex's own v1-5-0/
% v1-6-0 comments for the worked example this was hoisted out of).
% Found 2026-07-20 that the raw-`#1`-in-frame-body bug documented above
% is *not* actually \tableofcontents-specific, contrary to that comment's
% own framing -- slides/scs-Political-Risk-Stability-Index-current.tex has
% no \tableofcontents anywhere in the deck-wide overlay chain leading to
% it, yet hit the identical "Illegal parameter number in definition of
% \iterate" at its own \end{standardframe} from the exact same literal-
% \renewcommand{\makelabel}[1]{...#1...}-in-frame-body pattern. Whatever
% the precise trigger condition actually is, the same hoist-to-preamble
% fix (a bare macro *call*, no `#` token, left in the frame body) sidesteps
% it either way.
\newcommand{\SMEnumerateLabelFlushLeftShifted}[2]{%
	\setlength{\labelwidth}{#2}%
	\renewcommand{\makelabel}[1]{\hspace{#1}\makebox[\labelwidth][l]{##1}}%
}

% Width of a single digit in this deck's normal content font, measured once
% here (safe, plain preamble-load-time \settowidth) rather than inline in
% the Agenda slide's own frame body -- \tableofcontents there puts the whole
% frame through beamer's overlay-replay machinery, which the comment above
% \SMTocLabelFlushLeft already found chokes on raw tokens (like a bare
% \settowidth call) sitting directly in a replayed frame body, the same
% "Missing number"/"Illegal unit of measure" class of error as that `#`
% problem, confirmed directly the same way. Used by
% slides/nkj-Agenda-current.tex's own "widen the number-to-title gap"
% fix, which only ever needs to *reference* this precomputed value inside
% the frame, not measure anything itself there.
\newlength{\SMTocLabelExtraGap}
\settowidth{\SMTocLabelExtraGap}{0}

% For a two-column layout (e.g. bullets next to a chart/formula), nest
% beamer's own "columns" directly inside standardframe -- do NOT bake
% \begin{columns}/\end{columns} into a custom environment's begin/end code:
% beamer collects each frame's body as a token list to replay it once per
% overlay, and a columns environment that is opened in the "before" half and
% closed in the "after" half of a *different* environment confuses that
% collection (it never finds a properly matched \end{frame}). Nesting it in
% the document body, as below, avoids the issue entirely:
%
%   \begin{standardframe}{Title}
%   	\begin{columns}[t,onlytextwidth]
%   		\begin{column}{0.48\linewidth} ... \end{column}
%   		\begin{column}{0.48\linewidth} ... \end{column}
%   	\end{columns}
%   \end{standardframe}

% For automatic section dividers and single-statement "standout" slides, use
% what the gotham theme already provides instead of a custom environment:
%   \section{...}                                                (section divider, automatic)
%   \begin{frame}[standout]{Key takeaway}\end{frame}             (big single statement)
%   \begin{block}{...}/\begin{alertblock}{...}/\begin{exampleblock}{...}/\begin{referenceblock}{...}  (definitions, formulas, examples, references)
% For a course's own title slide, section-topic overview, quote, figure-only,
% and presenter-attention slides, see the five environments below.

% The block templates (block, alerted, example, reference, info and definition begin and end), their
% colors and the helpers they use (\SMBlockSideMargin, \SMBlockTitleWithMarker, \SMSetBlockBullets,
% \SMAlertMarker) are the theme's (beamerthemeSlideMage.sty); the three block types the beamer
% environments do not have are defined here.
\newenvironment<>{referenceblock}[1]{%
	\begin{actionenv}#2%
		\def\insertblocktitle{#1}%
		\par%
		\usebeamertemplate{block reference begin}}
	{\par%
		\usebeamertemplate{block reference end}%
	\end{actionenv}}

\newenvironment<>{infoblock}[1]{%
	\begin{actionenv}#2%
		\def\insertblocktitle{#1}%
		\par%
		\usebeamertemplate{block info begin}}
	{\par%
		\usebeamertemplate{block info end}%
	\end{actionenv}}

\newenvironment<>{definitionblock}[1]{%
	\begin{actionenv}#2%
		\def\insertblocktitle{#1}%
		\par%
		\usebeamertemplate{block definition begin}}
	{\par%
		\usebeamertemplate{block definition end}%
	\end{actionenv}}


% A narrower, freely-positionable version of \begin{block}{...} (2026-07-06,
% direct instruction) -- a plain block always spans the *current*
% \linewidth (that's what block begin/end's own \hspace+minipage wrapper
% above is measured against), so it can't be dropped into a fixed-width
% spot (e.g. beside a table, inside a tikz overlay \node) without first
% narrowing \linewidth some other way. \SMFloatingBlock just wraps a
% plain block in one more, *outer* minipage of the caller's own chosen
% width -- block begin/end's own inner minipage then measures itself
% against that narrower width instead of the frame's full one, nesting
% correctly with no changes needed to the block templates themselves.
% Title/body colors, margins, and font are exactly \block's own (same
% "block title"/"block body" beamer color names) -- this only changes the
% overall width and (by taking it out of the normal document flow, left to
% the caller) the position, not the look. Font size is also left to the
% caller (wrap the whole environment in a size declaration, e.g.
% {\footnotesize \begin{SMFloatingBlock}{5cm}{Excel} ... \end{SMFloatingBlock}})
% rather than fixed here, since the right size depends entirely on how
% much room the block ends up placed in.
%
% INNER PADDING (2026-09-30, direct instruction: systematic replacements for
% per-slide fixes): 20 of 27 floating blocks opened with a hand-written
% \vspace{2mm} and 17 closed with \par\vspace{1.2mm} (or 1mm/3.4pt), to give
% the body text room from the title bar and the block's bottom edge. All
% three SMFloating*Block environments now do that themselves, via an
% optional key list:
%   \begin{SMFloatingInfoBlock}[top=2mm, bottom=1.2mm]{6cm}{Title}
% Defaults top=2mm, bottom=1.2mm (the common values); 0pt inserts nothing
% at all. Blocks that had no or other padding carry their own values
% explicitly (e.g. [top=0pt]), so every block kept its exact look.
\makeatletter
\pgfqkeys{/SMFloatingBlock}{
	top/.store in=\SM@fb@top,
	bottom/.store in=\SM@fb@bottom,
}
\newcommand{\SM@fb@begin}[1]{%
	\pgfqkeys{/SMFloatingBlock}{top=2mm, bottom=1.2mm, #1}%
	\ifdim\SM@fb@top=0pt \else\vspace{\SM@fb@top}\fi
}
\newcommand{\SM@fb@end}{%
	\ifdim\SM@fb@bottom=0pt \else\par\vspace{\SM@fb@bottom}\fi
}
\NewDocumentEnvironment{SMFloatingBlock}{O{}mm}{%
	\begin{minipage}{#2}
	\begin{block}{#3}
	\SM@fb@begin{#1}%
}{%
	\SM@fb@end
	\end{block}
	\end{minipage}
}
\makeatother

% Same idea, but wrapping \exampleblock instead of \block (2026-07-06,
% direct instruction, first used on
% slides/fr3-Forms-of-Capital-Market-Efficiency-Strong-Semistrong-Weak-current.tex
% and slides/ha9-Monacorp-Truck-Lease-Breakeven-Rent-current.tex, both
% switching an existing \SMFloatingBlock to the "example" green header
% instead of the plain violet one) -- not worth generalizing into a single
% macro parameterized over the block name: \SMFloatingBlock's own body
% is only two lines, so a third variant is cheaper to just write out than
% to build a block-name-dispatch layer for.
\makeatletter
\NewDocumentEnvironment{SMFloatingExampleBlock}{O{}mm}{%
	\begin{minipage}{#2}
	\begin{exampleblock}{#3}
	\SM@fb@begin{#1}%
}{%
	\SM@fb@end
	\end{exampleblock}
	\end{minipage}
}
\makeatother

% Same idea again, wrapping \infoblock (2026-07-06, direct instruction,
% first used on
% slides/fr3-Forms-of-Capital-Market-Efficiency-Strong-Semistrong-Weak-current.tex's
% own "Random walk" formula box, switched from the green
% \SMFloatingExampleBlock to this blue one the same day).
\makeatletter
\NewDocumentEnvironment{SMFloatingInfoBlock}{O{}mm}{%
	\begin{minipage}{#2}
	\begin{infoblock}{#3}
	\SM@fb@begin{#1}%
}{%
	\SM@fb@end
	\end{infoblock}
	\end{minipage}
}
\makeatother

% Every SMFloatingInfoBlock anchored top-right against the institute
% logo (the deck's standard
% "(current page.north east)+(-\SMFrametitleRightPadding,-\SMContentTopAlign)"
% corner pattern) lands its own visible right edge short of the logo's,
% by a fixed amount independent of the block's own chosen width (#1) --
% \infoblock's underlying template insets its colored fill by
% \SMBlockSideMargin from *both* sides of whatever width it's given
% (that margin exists to match a first-level bullet's own left-hand
% position, see its own comment above; the fact that it also narrows the
% *right* side is just a side effect of reusing one symmetric
% hspace+minipage trick for both, not a deliberate right-margin choice),
% partly offset by beamer's own beamercolorbox "colsep*" background bleed
% (colsep* pads the colored fill outward past its own nominal box edge
% without also insetting the text, unlike unstarred "colsep" -- see
% beamerbasecolor.sty's own \beamercolorbox definition). Root-caused and
% measured directly (pymupdf against a compiled probe deck, not derived
% purely on paper) 2026-07-21, after
% slides/j8y-Asset-Allocation-Example-Weights-current.tex's own one-off
% "+2.2pt" nudge (measured that same day against just the institute
% logo's yellow chip) turned out to disagree with
% slides/fr3-Forms-of-Capital-Market-Efficiency-Strong-Semistrong-Weak-current.tex's
% own older "+3.165pt" nudge (2026-07-06, measured against the logo's
% *full* gray-bordered footprint, shadow/card included) -- both slides'
% own nudges were each individually correct for their own chosen
% reference edge, they just didn't agree on what "the logo's own edge"
% means. Direct instruction, 2026-07-21: the *full* gray-bordered
% footprint is the deck's intended reference, so 3.165pt is the value
% kept here; j8y's own newer 2.2pt nudge was the one revised to match,
% not this one.
%
% Use "-\SMFloatingInfoblockRightAlign" (not the bare
% "-\SMFrametitleRightPadding") for a SMFloatingInfoBlock's own
% anchor node whenever its right edge is meant to align with the logo --
% every current use of this pattern across the deck has been updated to
% match; new slides using this same corner pattern should use this
% length from the start rather than re-measuring their own one-off nudge.
\newlength{\SMFloatingInfoblockRightAlign}
\setlength{\SMFloatingInfoblockRightAlign}{\dimexpr\SMFrametitleRightPadding-3.165pt\relax}

% An infoblock's colored title bar starts 6.02pt below the top of its own
% box (the block template's top padding); a floating infoblock anchored at
% the Top Margin therefore sits 6.02pt low unless its anchor is lifted by
% this much. Measured 2026-07-23 (y3a/gf2: title-bar rect at 47.60pt vs.
% the 41.58pt Top Margin), written as a literal "+6.02pt" on 15 anchors
% until 2026-09-30.
\newlength{\SMFloatingInfoblockTopAlign}
\setlength{\SMFloatingInfoblockTopAlign}{6.02pt}

% "Bottom Margin" as a length (2026-09-30): the gap between the footer box
% and content anchored at the bottom of a slide, the mirror image of the Top
% Margin gap between the title bar and the first line (direct instruction on
% bwi, 2026-07-2x: "the same clearance from the footer that the text has from
% the title bar"). Measured from the page's bottom edge, like
% \SMContentTopAlign is from its top edge.
\newlength{\SMContentBottomAlign}
\setlength{\SMContentBottomAlign}{\dimexpr\SMFooterBoxHeight+\SMContentTopAlign-\SMTitleBarBottomAlign\relax}

% \cornerfigure[keys]{image file or content} (2026-09-30, direct instruction:
% systematic replacements for per-slide fixes) -- a figure in a corner of the
% content area, as a zero-space overlay: the construct 105 slides wrote out
% by hand,
%   \vbox to 0pt{\vss\begin{tikzpicture}[overlay, remember picture]
%     \node[anchor=north east, inner sep=0pt] at ($(current page.north east)
%       +(-\SMFrametitleRightPadding,-\SMContentTopAlign)$) {...};
%   \end{tikzpicture}}
% The usual call (2026-10-01, direct instruction: "the width or height of
% the figure has to be specified in the code, that's all that should be
% needed") names the image and one size; the other follows from the aspect
% ratio:
%   \cornerfigure[height=5.6cm]{media/chart-x.png}
%   \cornerfigure[width=4cm, border]{media/photo-x.jpg}
% Keys:
%   width=, height=      the image's size; with either one the argument is an
%                        image file (both: fitted inside, aspect ratio kept).
%                        Without them the argument is arbitrary content (a
%                        block, a \parbox, an \includegraphics with options).
%   border               the yellow image border (SMImageBorder); the
%                        border's outer edge sits on the margins
%   corner=ne|nw|se|sw   which corner (default ne): Top/Bottom Margin
%                        (\SMContentTopAlign/\SMContentBottomAlign)
%                        and Right/Left Margin (\SMFrametitleRightPadding,
%                        \SMFrametitleLeftPadding)
%   corner=e|w           on the Right/Left Margin, centred vertically between
%                        Top and Bottom Margin
%   block                the content is a SMFloatingInfoBlock: lifts it by
%                        \SMFloatingInfoblockTopAlign and aligns its
%                        visible right edge with the logo
%                        (\SMFloatingInfoblockRightAlign)
%   xshift, yshift       a deliberate offset from that corner (right/up +)
%   node={...}           extra TikZ node options
% The node has no inner or outer sep, so the figure's own edge -- or, with a
% border, the border's outer edge (see SMImageBorder) -- is what lands on
% the margins. Like any remember-picture overlay it needs the usual second
% compile. Text that must not run under the figure: besidefigure below.
\makeatletter
\newif\ifSM@cf@block
\newif\ifSM@cf@border
\pgfqkeys{/SMCorner}{
	corner/.store in=\SM@cf@corner,
	xshift/.store in=\SM@cf@x,
	yshift/.store in=\SM@cf@y,
	node/.store in=\SM@cf@node,
	width/.store in=\SM@cf@w,
	height/.store in=\SM@cf@h,
	gap/.store in=\SM@bf@gap,
	block/.code={\SM@cf@blocktrue},
	border/.code={\SM@cf@bordertrue},
}
% all keys back to their defaults, then the caller's
\def\SM@cf@setkeys#1{%
	\SM@cf@blockfalse\SM@cf@borderfalse
	\pgfqkeys{/SMCorner}{corner=ne, xshift=0pt, yshift=0pt, node={},
		width={}, height={}, gap=12pt, #1}%
}
% \SM@cf@content: the figure itself -- \includegraphics[width/height]{#1}
% when a size key was given, #1 unchanged otherwise (\long: content such as
% a floating block may hold paragraphs)
\long\def\SM@cf@makecontent#1{%
	\def\SM@cf@gkeys{}%
	\ifx\SM@cf@w\@empty\else
		\edef\SM@cf@gkeys{width=\unexpanded\expandafter{\SM@cf@w}}%
	\fi
	\ifx\SM@cf@h\@empty\else
		\edef\SM@cf@gkeys{\ifx\SM@cf@gkeys\@empty\else\unexpanded\expandafter{\SM@cf@gkeys},keepaspectratio,\fi
			height=\unexpanded\expandafter{\SM@cf@h}}%
	\fi
	\ifx\SM@cf@gkeys\@empty
		\def\SM@cf@content{#1}%
	\else
		\edef\SM@cf@content{\noexpand\includegraphics[\unexpanded\expandafter{\SM@cf@gkeys}]{\unexpanded{#1}}}%
	\fi
}
\NewDocumentCommand{\cornerfigure}{O{} +m}{%
	\begingroup
	\SM@cf@setkeys{#1}%
	\SM@cf@makecontent{#2}%
	\ifSM@cf@block
		\edef\SM@cf@right{\SMFloatingInfoblockRightAlign}%
		\edef\SM@cf@top{\dimexpr\SMContentTopAlign-\SMFloatingInfoblockTopAlign\relax}%
	\else
		\edef\SM@cf@right{\SMFrametitleRightPadding}%
		\edef\SM@cf@top{\SMContentTopAlign}%
	\fi
	\def\SM@cf@anchor{north east}%
	\edef\SM@cf@at{(current page.north east)+(-\SM@cf@right+\SM@cf@x,-\SM@cf@top+\SM@cf@y)}%
	\ifdefstring{\SM@cf@corner}{nw}{%
		\def\SM@cf@anchor{north west}%
		\edef\SM@cf@at{(current page.north west)+(\SMFrametitleLeftPadding+\SM@cf@x,-\SM@cf@top+\SM@cf@y)}}{}%
	\ifdefstring{\SM@cf@corner}{se}{%
		\def\SM@cf@anchor{south east}%
		\edef\SM@cf@at{(current page.south east)+(-\SM@cf@right+\SM@cf@x,\SMContentBottomAlign+\SM@cf@y)}}{}%
	\ifdefstring{\SM@cf@corner}{sw}{%
		\def\SM@cf@anchor{south west}%
		\edef\SM@cf@at{(current page.south west)+(\SMFrametitleLeftPadding+\SM@cf@x,\SMContentBottomAlign+\SM@cf@y)}}{}%
	% e/w: on the Right/Left Margin, centred vertically between Top and Bottom Margin
	\ifdefstring{\SM@cf@corner}{e}{%
		\def\SM@cf@anchor{east}%
		\edef\SM@cf@at{(current page.north east)+(-\SM@cf@right+\SM@cf@x,-0.5\paperheight-0.5\SMContentTopAlign+0.5\SMContentBottomAlign+\SM@cf@y)}}{}%
	\ifdefstring{\SM@cf@corner}{w}{%
		\def\SM@cf@anchor{west}%
		\edef\SM@cf@at{(current page.north west)+(\SMFrametitleLeftPadding+\SM@cf@x,-0.5\paperheight-0.5\SMContentTopAlign+0.5\SMContentBottomAlign+\SM@cf@y)}}{}%
	% TikZ's calc parser must see the coordinate itself, not a macro: build
	% the \node call with everything but the content already expanded
	\edef\SM@cf@cmd{\noexpand\node[anchor=\SM@cf@anchor, inner sep=0pt, outer sep=0pt%
		\ifSM@cf@border,SMImageBorder\fi
		\ifx\SM@cf@node\@empty\else,\unexpanded\expandafter{\SM@cf@node}\fi]
		at ($\SM@cf@at$)}%
	\vbox to 0pt{\vss
		\begin{tikzpicture}[overlay, remember picture]
			\SM@cf@cmd {\SM@cf@content};
		\end{tikzpicture}%
	}%
	\endgroup
}

% \begin{besidefigure}[keys]{image file or content} text \end{besidefigure}
% (2026-09-30; 2026-10-01 the figure's width is measured, not passed) -- text
% next to a figure in the top-right corner: the figure goes into \cornerfigure
% (same keys: height=/width=, border, ...), the text into a top-aligned
% minipage exactly as wide as the room left of the figure minus `gap`
% (default 12pt). 84 slides built this by hand, each with its own guessed
% minipage width (9.3cm, 0.58\linewidth, ...).
%   \begin{besidefigure}[height=5.6cm, border]{media/chart-x.png}
%   	\begin{itemize} ... \end{itemize}
%   \end{besidefigure}
% Optical top alignment (v5.27.0): the text column's first line sits on the Top
% Margin by its line box (the ascender height, 0.935em), so its visible letters
% start about 0.2em lower; a figure placed on the margin itself looked that
% much too high next to it. besidefigure therefore lowers its figure by this
% much (\SMBesideFigureTopGap, default 0.2em of the surrounding text), so the
% figure's top edge meets the top of the first line's letters. Give
% yshift=0pt among the keys to put the figure on the margin again.
\newlength{\SM@bf@width}
\NewDocumentEnvironment{besidefigure}{O{} +m}{%
	\SM@cf@setkeys{#1}%
	\SM@cf@makecontent{#2}%
	\settowidth{\SM@bf@width}{\SM@cf@content}%
	\ifSM@cf@border\addtolength{\SM@bf@width}{\SMImageBorderWidth}\fi
	\cornerfigure[yshift=-\SMBesideFigureTopGap,#1]{#2}%
	% text column: from the text's left edge to the figure's left edge minus gap
	\setlength{\SM@bf@width}{\dimexpr\linewidth+\beamer@rightmargin-\SMFrametitleRightPadding-\SM@bf@width-\SM@bf@gap\relax}%
	\begin{minipage}[t]{\SM@bf@width}%
}{%
	\end{minipage}%
}

% \fitfigure[keys]{image file} (2026-09-30) -- an image as large as it fits,
% keeping its aspect ratio, in the box from a top line down to the Bottom
% Margin (\SMContentBottomAlign) and from the Left to the Right Margin,
% placed as a zero-space overlay. Replaces the per-slide
% "scale up ... as much as possible" arithmetic (hand-computed widths and
% \vspace offsets, often converged in two passes).
%   top=here|margin   where the box starts: 'here' (default) = the current
%                     position in the slide (below the text so far), plus
%                     `gap`; 'margin' = the Top Margin
%   gap=<length>      space below the text for top=here (default 6pt)
%   valign=top|center|bottom, halign=left|center|right  (default center, center)
%   bottom=<length>   the box's bottom, from the page's bottom edge (default
%                     \SMContentBottomAlign)
%   node={...}        extra TikZ node options (e.g. SMImageBorder)
% The position is a remember-picture coordinate, so like every overlay here
% it needs the second compile.
\newcount\SM@fit@count
\def\SM@fit@settop#1{\coordinate (SMFit@top) at #1;}
\pgfqkeys{/SMFit}{
	top/.store in=\SM@fit@top,
	gap/.store in=\SM@fit@gap,
	valign/.store in=\SM@fit@valign,
	halign/.store in=\SM@fit@halign,
	bottom/.store in=\SM@fit@bottom,
	node/.store in=\SM@fit@node,
}
\NewDocumentCommand{\fitfigure}{O{} m}{%
	\begingroup
	\pgfqkeys{/SMFit}{top=here, gap=6pt, valign=center, halign=center,
		bottom=\SMContentBottomAlign, node={}, #1}%
	\global\advance\SM@fit@count by 1
	\edef\SM@fit@name{SMFit\the\SM@fit@count}%
	\ifdefstring{\SM@fit@top}{here}{%
		\vbox to 0pt{\hbox{\tikz[remember picture, overlay]\coordinate (\SM@fit@name) at (0,0);}\vss}%
	}{}%
	% alignment choices, decided here (pgfmath can't evaluate \ifstrequal)
	\def\SM@fit@xexpr{(\x2+\x3)/2}\def\SM@fit@h@anchor{}%
	\ifdefstring{\SM@fit@halign}{left}{\def\SM@fit@xexpr{\x2}\def\SM@fit@h@anchor{west}}{}%
	\ifdefstring{\SM@fit@halign}{right}{\def\SM@fit@xexpr{\x3}\def\SM@fit@h@anchor{east}}{}%
	\def\SM@fit@yexpr{(\y1+\y2)/2}\def\SM@fit@v@anchor{}%
	\ifdefstring{\SM@fit@valign}{top}{\def\SM@fit@yexpr{\y1}\def\SM@fit@v@anchor{north}}{}%
	\ifdefstring{\SM@fit@valign}{bottom}{\def\SM@fit@yexpr{\y2}\def\SM@fit@v@anchor{south}}{}%
	\edef\SM@fit@anchor{\SM@fit@v@anchor
		\ifx\SM@fit@v@anchor\@empty\else\ifx\SM@fit@h@anchor\@empty\else\space\fi\fi
		\SM@fit@h@anchor}%
	\ifx\SM@fit@anchor\@empty\def\SM@fit@anchor{center}\fi
	\ifdefstring{\SM@fit@top}{here}{%
		\edef\SM@fit@topcoord{(\SM@fit@name)}%
		\edef\SM@fit@gapexpr{\SM@fit@gap}%
	}{%
		\def\SM@fit@topcoord{($(current page.north)+(0,-\SMContentTopAlign)$)}%
		\def\SM@fit@gapexpr{0pt}%
	}%
	\edef\SM@fit@nodeopts{anchor=\SM@fit@anchor, inner sep=0pt%
		\ifx\SM@fit@node\@empty\else,\unexpanded\expandafter{\SM@fit@node}\fi}%
	\vbox to 0pt{\vss
		\begin{tikzpicture}[overlay, remember picture]
			\expandafter\SM@fit@settop\expandafter{\SM@fit@topcoord}%
			\coordinate (SMFit@bl) at ($(current page.south west)+(\SMFrametitleLeftPadding,\SM@fit@bottom)$);
			\coordinate (SMFit@r) at ($(current page.east)+(-\SMFrametitleRightPadding,0)$);
			\path let \p1 = (SMFit@top), \p2 = (SMFit@bl), \p3 = (SMFit@r) in
				\pgfextra{%
					\pgfmathsetlengthmacro\SM@fit@h{max(\y1-\SM@fit@gapexpr-\y2,1pt)}%
					\pgfmathsetlengthmacro\SM@fit@w{max(\x3-\x2,1pt)}%
					\pgfmathsetlengthmacro\SM@fit@x{\SM@fit@xexpr}%
					\pgfmathsetlengthmacro\SM@fit@y{\SM@fit@yexpr}%
					\ifdefstring{\SM@fit@valign}{top}{\pgfmathsetlengthmacro\SM@fit@y{\y1-\SM@fit@gapexpr}}{}%
					\ifdefstring{\SM@fit@valign}{center}{\pgfmathsetlengthmacro\SM@fit@y{(\y1-\SM@fit@gapexpr+\y2)/2}}{}%
					\global\let\SM@fit@H\SM@fit@h \global\let\SM@fit@W\SM@fit@w
					\global\let\SM@fit@X\SM@fit@x \global\let\SM@fit@Y\SM@fit@y
				};
			\expandafter\node\expandafter[\SM@fit@nodeopts] at (\SM@fit@X,\SM@fit@Y)
				{\includegraphics[width=\SM@fit@W, height=\SM@fit@H, keepaspectratio]{#2}};
		\end{tikzpicture}%
	}%
	\endgroup
}
\makeatother

%%%%%%%%%%%%%%%%%%%%%%%%%%%%%%%%%%%%%%%%%%%%%%%%%%%%%%%%%%%%%%%%%%%%%%%%%%%%%
% Special frame types (converted from ExampleSlides.pptx; see README.md)
%
% All five below deliberately reuse as much of standardframe -- and of each
% other -- as possible, rather than reimplementing frame chrome from
% scratch: sectionframe, figureframe, and highlightframe are literally thin
% wrappers around standardframe itself (same title bar, logo, footer, bullet
% styling; only what's documented at each one differs). titleframe and
% quoteframe can't be standardframe wrappers -- both are title-less, and
% quoteframe replaces the whole body layout -- but still reuse its palette,
% fonts, and margins (\SMFrametitleLeftPadding, the institute logo
% drawing command, etc.) rather than inventing new ones.
%%%%%%%%%%%%%%%%%%%%%%%%%%%%%%%%%%%%%%%%%%%%%%%%%%%%%%%%%%%%%%%%%%%%%%%%%%%%%

% titleframe: the deck's own title slide (ExampleSlides.pptx slide 1) --
% typically used exactly once, at the very start of a course's first deck.
% One continuous yellow box -- like standardframe's own title bar -- flush
% with the page's top and left edges, stopping short of the institute logo
% on the right (the same margin already used around the logo elsewhere:
% \SMLogoWidth + 2*\SMFrametitleRightPadding), holding the title and,
% below it, a second line (author/term/whatever) in the same black text as
% standardframe's own frametitle. Usage:
%
%   \titleframe{Financial Markets}{Stefan Palan\\Summer term 2026}
%
% A *command*, not an environment, taking the second (lower) line's content
% as its second argument -- plain text, use \\ for line breaks same as any
% normal LaTeX paragraph, not an itemize, since the source deck's own
% subtitle block is just stacked lines. This has to be a command rather
% than an environment with the body as its content: beamer's frame body is
% captured by a delimited-argument macro (\beamer@collect@@body, looked up
% by triggering the error below on purpose) hunting for a literal
% \end{frame}, and a beamercolorbox -- needed by an earlier version of this
% to wrap the subtitle text in its own colored band -- opened in a *custom
% environment's* begin code and closed in its end code confuses that hunt
% exactly like the \columns example above (same root cause, different
% environment): with
% \newenvironment{titleframe}{...\begin{beamercolorbox}{...}}{\end{beamercolorbox}...},
% compiling ends in "! File ended while scanning use of \beamer@collect@@body."
% every time, however short the body. Taking the content as a normal brace
% argument instead sidesteps the whole issue. No footer/frame number
% (matches beamer's own convention for title slides, e.g. the old
% \begin{frame}[plain,noframenumbering]\titlepage\end{frame} this replaces --
% see README.md) and no chapter-logo support: both would be meaningless on a
% slide that appears once, before any chapter exists.
%
% Unlike the frametitle hook further up, neither the logo nor the yellow box
% here need the white-patch-and-redraw trick: the box's own width already
% stops short of the logo (rather than being painted full-width and then
% partly painted back over), so there's no yellow left to paint over in the
% first place. Both still need \tikz's "remember picture" on their
% "current page"-anchored nodes, though, and so still need the usual second
% compile pass (see "Title bar and logo" in README.md) -- checked
% empirically: without it, "current page" doesn't resolve to the actual page
% box on the first pass, and the logo/yellow box end up positioned nowhere
% near the page corners they're meant to be anchored to.
%
% \SMTitleframeBoxHeight is a fixed target for the box's *bottom* edge
% (matching what was actually asked for: "reaching down to where the lower
% grey bar currently ends" in the two-band design this replaces), not the
% blank space above the text -- that space is derived from it instead (box
% height minus the title+subtitle content's own measured height), so the
% bottom edge stays fixed even if a longer title ever wraps onto a second
% line; the blank space above would just shrink to compensate. Measured
% against the old two-band design's own rendered PDF coordinates (before
% this change): its second band's bottom edge sat 187.147pt below the page's
% top edge.
\newlength{\SMTitleframeBoxWidth}
\setlength{\SMTitleframeBoxWidth}{\dimexpr\paperwidth-\SMLogoWidth-2\SMFrametitleRightPadding\relax}
% Text inset from the box's own edges: \SMFrametitleLeftPadding, the same
% length standardframe's own title text uses, plus \SMTitleframeTextSep
% (2.5ex, matching this command's own former per-band beamercolorbox "sep").
% Because the leading blank space above the text (further down) is itself
% computed as "box height minus the content's own natural height", the
% content's *bottom* padding is the only thing that determines how far the
% last line sits above the box's bottom edge -- the leading space silently
% absorbs whatever that padding is, keeping the bottom edge fixed either
% way. (The relationship is a direct 1:1 offset, checked empirically by
% comparing two rendered values, but with a nonzero intercept from beamer's
% own default top-of-frame vskip -- there's no clean closed-form for the
% right value, hence the empirical calibration here rather than a formula.)
% Calibrated (against the rendered PDF's own text coordinates, not just
% visually) to about 1.5 line-heights of the subtitle text's own \Large
% size (~20.6pt/line here) above the box's bottom edge; re-measure and
% adjust if the subtitle's font size, or beamer's own default top-of-frame
% vskip, ever change.
\newbox\SMTitleframeBox
\makeatletter
\newcommand{\titleframe}[2]{%
	% No noframenumbering (2026-08-11, direct instruction): this frame has
	% no footer at all ("plain"), so removing it has no visible effect on
	% the title slide itself -- but it does mean the title slide now
	% consumes a real "framenumber" step, same as any other frame, so the
	% deck's first real content slide is numbered 2, not 1 (matching this
	% slide's own physical position in the compiled PDF -- SlideMage's
	% own deck-editor "Render" preview relies on this: it derives a
	% previewed slide's page number from its plain list position, which
	% only equals the real compiled frame number when every \slide{}-family
	% call -- including this one -- consumes exactly one framenumber step).
	\SMOpenFrame{[plain,t]}%
	% This frame has no footer to carry the usual "(slide id: ...)" marker
	% -- see \SMInvisibleSlideIDStamp's own comment above for why that
	% breaks SlideMage's compile-time preview snapshot, and what this
	% stamps instead.
	\SMInvisibleSlideIDStamp%
	% A bare tikzpicture behaves like an hbox (same reasoning as the footer
	% box's own \vbox-to-0pt wrapper further up): sitting directly in the
	% frame's vertical list, it silently opens an implicit paragraph,
	% switching into horizontal mode. Everything placed after it -- the
	% \vspace* further down included -- would then run in horizontal mode
	% instead of vertical mode, which is what a first version of this command
	% got wrong: with two bare tikzpictures here, \vspace*'s glue ended up as
	% *horizontal* space on an implicit line together with them (visible as
	% the title/subtitle shifting sideways, not down, and the vspace amount
	% appearing to have no effect at all -- it was moving the wrong axis).
	% \vbox to 0pt{\vss ...} keeps each overlay at zero space in the normal
	% flow without ever leaving vertical mode.
	%
	% \SMArtifactBegin/\SMArtifactEnd (defined with the footer, "PDF
	% accessibility" comment): this box is purely decorative background, same
	% kind of element as the footer's own yellow box -- marked the same way,
	% for a screen reader to skip. titleframe has no footer at all to worry
	% about separately (it's a "plain" frame, no footline), so this box is
	% the only thing here that needs it; the title/subtitle text right below
	% is real content and stays unmarked.
	\SMArtifactBegin{/Type /Background}%
	\vbox to 0pt{\vss
	\begin{tikzpicture}[overlay, remember picture]
		% \SMTitleBoxColor, not a literal color token -- same setting the
		% standard frametitle bar (colorA, above) uses, so the title slide's
		% own box matches every other title bar in the deck (2026-07-16,
		% direct instruction).
		% outer sep=0pt here and on the logo below: without it, tikz's
		% default outer sep insets the fill/logo 0.2pt from the anchored
		% page corner -- a hairline white line above/left of this box, and
		% its right edge/the logo 0.2-0.4pt off every other frame's. See
		% the frametitle hook's own comment for the full story.
		\node[fill=\SMTitleBoxColor, anchor=north west, inner sep=0pt, outer sep=0pt,
		      minimum width=\SMTitleframeBoxWidth,
		      minimum height=\SMTitleframeBoxHeight]
			at (current page.north west) {};
	\end{tikzpicture}%
	}%
	\SMArtifactEnd
	\vbox to 0pt{\vss
	\begin{tikzpicture}[overlay, remember picture]
		% inner ysep derived, not \SMFrametitleRightPadding -- see the
		% frametitle hook's fixed-position logo redraw for why (matches the
		% inline logo's vertically-centered position on every other frame).
		\node[anchor=north east, inner xsep=\SMFrametitleRightPadding,
		      inner ysep=\dimexpr\SMTitleBarBottomAlign-\SMLogoHeight\relax,
		      outer sep=0pt]
			at (current page.north east) {\SMInstituteLogo};
	\end{tikzpicture}%
	}%
	\setbox\SMTitleframeBox=\vbox{%
		\hsize=\SMTitleframeBoxWidth
		\parindent=0pt
		\leftskip=\SMFrametitleLeftPadding
		\advance\leftskip by \SMTitleframeTextSep
		\rightskip=\SMFrametitleLeftPadding
		\advance\rightskip by \SMTitleframeTextSep
		% \SMTitleTextColor, not a literal color token -- follows the
		% box color change right above (2026-07-16, direct instruction):
		% leaving this hardcoded dark would be unreadable once the box
		% itself switches to a dark blue for MCI/Innsbruck.
		\color{\SMTitleTextColor}%
		% Beamer redefines \\ inside every frame to peek for an overlay spec
		% (\\<2->, for line-by-line reveals in itemize/enumerate); that
		% redefinition assumes it's running in beamer's own frame-body
		% machinery and chokes ("Undefined control sequence", silently
		% swallowing the following word) when \\ is used for a plain line
		% break in a raw \vbox like this one. Restoring plain LaTeX's own
		% \\ (\@normalcr) for this group only fixes it without affecting \\
		% anywhere else in the deck. (The @ here needs catcode 11, which is
		% why this whole command definition -- not just this line -- is
		% wrapped in \makeatletter/\makeatother: macro bodies are tokenized
		% at *definition* time, so a \makeatletter placed inside the body
		% itself would take effect too late, after "\@normalcr" had already
		% been mis-tokenized as "\@" followed by four ordinary-text letters.)
		\let\\\@normalcr
		\vskip\SMTitleframeTextSep
		\noindent{\Huge\bfseries #1\par}%
		\vskip\SMTitleframeTextSep
		\vskip\SMTitleframeBandGap
		\vskip\SMTitleframeTextSep
		\noindent{\Large #2\par}%
		\vskip\SMTitleframeBottomPad
	}%
	\vspace*{\dimexpr\SMTitleframeBoxHeight-\ht\SMTitleframeBox-\dp\SMTitleframeBox\relax}%
	\box\SMTitleframeBox
	% This command's own \color{\SMTitleTextColor} above (set inside a
	% stored box register) can leak into the *next* frame's body text --
	% see standardframe's own opening code ("Forces this frame's own body")
	% for the real fix and the full story (a trailing reset attempted
	% here first turned out to be a silent no-op, since nothing draws after
	% it within this same frame to actually force the color operator to be
	% emitted).
	\end{frame}%
}
\makeatother

% fullbleedframe: a slide that is *only* a full-screen image (FM2_2026s_ps_deck.pptx
% slide 13, a recruitment flyer pasted in as one embedded picture on an
% otherwise blank layout -- no title, no PPTX text at all) -- no title bar,
% no footer, no logo, matching that source's own blank layout intent, unlike
% figureframe (which is standardframe underneath and always keeps the usual
% title/footer chrome). The image is scaled to fit the slide while keeping
% its own aspect ratio (\includegraphics's keepaspectratio, capped by
% \paperwidth *and* \paperheight together), centered; whatever's left over
% on the binding axis (e.g. a widescreen slide holding a taller/narrower
% image) is filled by a solid black backdrop rather than left showing the
% frame's usual white canvas -- the point of a full-*bleed* image is that
% nothing reads as "empty margin" even when the aspect ratios don't match.
%
%   \fullbleedframe{media/photo-datacamp-recruitment-flyer}
%
% \SMFullBleedBlack is a fixed HTML black, deliberately *not* a
% \colorlet of c_primary_dark: this frame type's backdrop should stay
% neutral true black regardless of whether c_primary_dark itself is ever
% restyled to something else (e.g. a softer near-black, the way
% c_custom_gray already is elsewhere in this file) for branding
% reasons -- those are two different jobs (brand color vs. letterbox
% backdrop) that only happen to want the same value today.
\definecolor{SMFullBleedBlack}{HTML}{000000}
\newcommand{\fullbleedframe}[1]{%
	% No noframenumbering -- see \titleframe's own comment above for why
	% (same reasoning, same 2026-08-11 change: this frame has no footer
	% either, so nothing here changes visibly, only later slides' numbers).
	\SMOpenFrame{[plain,t]}%
	% This frame has no footer either -- same \SMInvisibleSlideIDStamp
	% need as \titleframe above, see its own comment for why.
	\SMInvisibleSlideIDStamp%
	% \SMArtifactBegin/\SMArtifactEnd (defined with the footer, "PDF
	% accessibility" comment): the black backdrop is purely decorative
	% letterboxing, marked the same way as titleframe's/quoteframe's own
	% background boxes so a screen reader skips it. The image itself is left
	% unmarked *here* -- not because it goes undescribed, but because
	% \includegraphics's own global wrapper (this file's own
	% \RenewDocumentCommand{\includegraphics}, "PDF accessibility" section)
	% already does that job for every \includegraphics call site, this one
	% included: it records the image's role/source with
	% slidemage_alt.mark()/recordSource() and, for an informative image, emits
	% its short alt text as an invisible layer via \SMAltInvisible
	% before the image itself. Nothing extra is needed at this call site.
	\SMArtifactBegin{/Type /Background}%
	\vbox to 0pt{\vss
	\begin{tikzpicture}[overlay, remember picture]
		% outer sep=0pt: see the frametitle hook's comment -- without it the
		% backdrop sits 0.2pt in from the anchored page corner.
		\node[fill=SMFullBleedBlack, anchor=south west, inner sep=0pt, outer sep=0pt,
		      minimum width=\paperwidth, minimum height=\paperheight]
			at (current page.south west) {};
	\end{tikzpicture}%
	}%
	\SMArtifactEnd
	\SMFullBleedPick{#1}{\SMFullBleedLayers{#1}}{\SMFullBleedImage{#1}}%
	\end{frame}%
}
\makeatletter
\newcommand{\SMFullBleedPick}[3]{\in@{,}{#1}\ifin@#2\else#3\fi}
\makeatother

% \fullbleedframe{a.png,b.png,c.png} (2026-09-18, direct instruction:
% combine the Fink/Palan/Theissen excursus's build-step screenshots into
% one slide): a comma-separated list layers the images on top of each
% other, revealed one per overlay step (image k appears at step k and stays,
% each covering the previous one -- the screenshots are same-size, opaque
% and each already contains every earlier step's content, so stacking
% reproduces the source's own build animation). Instructor deck only --
% the during-course/after-course versions show just the last (final)
% image, in a single step. A single path behaves
% exactly as before. Once a layer is covered by the next one it is wrapped
% in the same /Background artifact as the backdrop above, so on step k a
% screen reader reads only image k's own alt text (each screenshot is
% described as its own "build step k of n") instead of every buried layer's.
\newcommand{\SMFullBleedImage}[1]{%
	\vbox to 0pt{\vss
	\begin{tikzpicture}[overlay, remember picture]
		\node[anchor=center] at (current page.center)
			{\includegraphics[width=\paperwidth, height=\paperheight, keepaspectratio]{#1}};
	\end{tikzpicture}%
	}%
}
\newcount\SMFBtotal
\newcommand{\SMFullBleedLayers}[1]{%
	\global\SMFBtotal=0
	\foreach \SMFBimg in {#1} {\global\advance\SMFBtotal by 1 }%
	\foreach \SMFBimg [count=\SMFBi] in {#1} {%
		\ifSMInstructorVersion
			\ifnum\SMFBi<\SMFBtotal
				\pgfmathtruncatemacro{\SMFBnext}{\SMFBi+1}%
				\only<\SMFBi>{\SMFullBleedImage{\SMFBimg}}%
				\only<\SMFBnext->{%
					\SMArtifactBegin{/Type /Background}%
					\SMFullBleedImage{\SMFBimg}%
					\SMArtifactEnd
				}%
			\else
				\only<\SMFBi>{\SMFullBleedImage{\SMFBimg}}%
			\fi
		\else
			\ifnum\SMFBi=\SMFBtotal
				\SMFullBleedImage{\SMFBimg}%
			\fi
		\fi
	}%
}

% sideimageframe: a standardframe that always carries one auto-placed
% "side filler image" (the user's own term, see
% slides/dfc-Angel-Investors-current.tex's own comment for the technique
% this formalizes into a reusable frame type) -- 2026-07-06, direct
% instruction. The image is positioned flush to three edges at once: its
% top touches the title bar's own bottom edge, its right touches the
% page's own right edge, and its bottom touches the footer box's own top
% edge -- so it's scaled by *height* to exactly
% \paperheight-\SMTitleBarBottomAlign-\SMFooterBoxHeight (the
% measured gap between those two), with `keepaspectratio` so the width
% simply follows from the image's own aspect ratio rather than being
% forced/distorted to some second, independent target -- anchoring the
% node at the page's own north-east corner then makes the top+right
% touches automatic, and the height exactly filling that gap makes the
% bottom touch automatic too, with no separate width constraint needed
% (only a *square* source image would ever visibly reach some further
% left boundary; this only guarantees top/right/bottom, matching what was
% actually asked for).
%
%   \begin{sideimageframe}{Title}{media/photo-example}
%   	\begin{itemize}
%   		\item ...
%   	\end{itemize}
%   	\source{...}
%   \end{sideimageframe}
%
% \sidefootnotes is switched on automatically (not left to the caller) --
% the whole point of this frame type is that its footnotes always end up
% rotated along the left edge of the image rather than colliding with it
% the way a full-width bottom footnote would (see
% slides/u29-Implications-of-Capital-Market-Efficiency-Prices-current.tex's
% own comment for exactly that collision, on a plain standardframe). Unlike
% every other \sidefootnotes use elsewhere in this deck, though, the strip
% can't sit at its usual fixed inset from the page's own right edge here --
% that would run it directly underneath the image, the same problem u29
% had. Fixed by *computing* the image's own rendered width (via
% \settowidth against the exact same \includegraphics call used to draw
% it) and using that to override \SMSideFootnoteRightInset for just
% this one frame -- moving the whole rotated strip to sit immediately left
% of the image instead of at the page's raw edge, entirely automatically,
% with no manual x-coordinate needed from the caller (the direct
% instruction offered a manual-x-coordinate fallback "if it is possible
% [to do automatically], that would be great" -- turned out to be
% possible). \SMSideImageFrameHeight/\SMSideImageFrameWidth are
% declared once, globally, and merely *set* (not re-declared) inside the
% environment body, so the frame type can be used more than once per deck;
% the \setlength override of \SMSideFootnoteRightInset is a plain,
% non-\global assignment, which -- like every other frame-local override in
% this file (e.g. slides/u29-...'s own \setbeamercolor{footnote}{fg=white})
% -- resets on its own once \begin{frame}...\end{frame}'s own TeX group
% closes, with no explicit reset needed.
%
% v5.30.5: 3.1pt became 2.1pt. Measured on a rendered slide whose rotated footnote
% text was finally visible (see the colour note at the strip's drawing code), 3.1pt
% left 3.9pt between the ink of the line nearest the image and the image -- the
% text's own depth strut adds ~0.8pt -- against the 2.93pt this was meant to match.
%
% The extra gap added on top of the image's own measured width (2026-07-06,
% direct instruction: "replicate ... the precise margin we have on regular
% slides between footnote and footer box") was originally an unmeasured
% 4pt guess -- replaced with 3.1pt, back-solved from the deck's own real
% reference gap: measured directly on a plain bottom-footnote slide
% (slides/xhe-Active-vs-Passive-Management-2026-First-current.tex's own
% \source{} line) that the real gap between a footnote's own text and the
% footer box's top edge is 2.93pt, while the original 4pt guess here was
% actually landing the rotated text 3.8pt from the image -- 4pt - (3.8-2.93)
% = 3.13pt, rounded to 3.1pt.
\newlength{\SMSideImageFrameHeight}
\newlength{\SMSideImageFrameWidth}
\newenvironment{sideimageframe}[3][t]{%
	\setlength{\SMSideImageFrameHeight}{\dimexpr\paperheight-\SMTitleBarBottomAlign-\SMFooterBoxHeight\relax}%
	\settowidth{\SMSideImageFrameWidth}{\includegraphics[height=\SMSideImageFrameHeight, keepaspectratio]{#3}}%
	\sidefootnotes
	\setlength{\SMSideFootnoteRightInset}{\dimexpr\SMSideImageFrameWidth+2.1pt\relax}%
	\begin{standardframe}[#1]{#2}%
	\vbox to 0pt{\vss
	\begin{tikzpicture}[overlay, remember picture]
		% outer sep=0pt: see the frametitle hook's comment -- without it the
		% image sits 0.2pt in from the page's right edge (and 0.2pt low).
		\node[anchor=north east, inner sep=0pt, outer sep=0pt] at ($(current page.north east)+(0,-\SMTitleBarBottomAlign)$)
			{\includegraphics[height=\SMSideImageFrameHeight, keepaspectratio]{#3}};
	\end{tikzpicture}%
	}%
}{%
	\end{standardframe}%
}

% sectionframe: a section-opening overview (ExampleSlides.pptx slide 3) --
% same layout as standardframe (footer, logo, bullet styling all completely
% unchanged), except the title bar itself is much taller (about a third of
% the page) and set in a bigger font, and the body text is a step bigger
% too, so it reads as "somewhat greater focus" and is unmistakable at a
% glance when flipping through slides -- while still keeping the deck's
% usual yellow-background/black-text look, unlike an earlier version of this
% that inverted to dark-background/yellow-text instead.
%
% The taller bar is NOT a separately-drawn box the way titleframe's is (see
% "Title slide" above): the gotham theme's own frametitle box already sizes
% itself from \SMFrametitleTopPadding/\SMFrametitleBottomPadding plus
% the title font's own strut height (beamerthemeSlideMage.sty, the
% "frametitle" template), so overriding those two lengths locally -- the
% same trick already used below for the font/color overrides -- grows the
% *theme's own* box instead of fighting it with a second one. This also
% means the white-patch-behind-the-logo/chapter-logo machinery (tied to the
% frametitle template via \addtobeamertemplate) keeps working unmodified,
% since it measures the box's actual rendered bottom edge rather than
% assuming any particular height.
%
% The title text is bottom-anchored within the bar -- sitting a fixed "one
% line" above the bar's bottom edge, growing *upward* for a title long
% enough to wrap onto a second line -- rather than top-anchored the way
% gotham's own frametitle template flows by default (which would instead
% leave a large, growing gap *below* a wrapped title, between it and the
% bar's bottom edge). gotham's template has no built-in way to do this: its
% only two knobs are the toppading/bottompading struts, both anchored to
% *one* end of a single top-down paragraph flow. So instead of leaving
% toppading fixed (as an earlier version of this did) the whole title text
% is measured into \SMSectionframeTitleBox first -- built with the same
% font and (approximately) the same usable width gotham's own template
% gives it, so its \ht+\dp closely track the title's real rendered size,
% including wrapping onto extra lines for a long title -- and toppading is
% then computed as "whatever's left" after that measured height and the
% fixed bottom gap are subtracted from the bar's total target height. A
% longer/wrapped title therefore automatically shrinks toppading (pushing
% the whole title block up, growing upward) rather than changing the gap
% below it. \SMSectionframeTitleWidth is a slight underestimate of
% gotham's actual per-line width (which is only narrowed by the logo on
% whichever line ends up next to it, typically just the last one) --
% applying that narrower width to *every* line is a deliberate, harmless
% simplification rather than replicating gotham's per-line logic exactly.
%
% The institute logo moves down along with the title now (it's vertically
% centered against the title's own baseline, gotham's instituteLogo
% template), rather than staying pinned where it sits on every other frame
% type, as an earlier version of this kept it -- an intentional change from
% that version: with the title itself now sitting near the bar's bottom
% edge, keeping the logo up at the bar's top would visually separate it
% from the title it's meant to sit beside.
%
% \begingroup/\endgroup (rather than the \global-set-and-reset dance
% \chapterlogo{}/\sidefootnotes need) is enough here: this only ever touches
% \setbeamercolor/\setbeamerfont/\setlength/\setbox, all of which assign
% locally under the hood (checked against beamerbasecolor.sty/
% beamerbasefont.sty directly for the first two -- their key handlers are
% plain \def, not \gdef -- and \setlength/\setbox are plain TeX assignments,
% local by default), so a plain surrounding TeX group reverts them
% correctly on its own.
%
% \SMSectionframeContentDrop adds a little extra room below the bar
% before the body content starts, on top of standardframe's own
% \SMContentTopShift pull-up (applied inside \begin{standardframe}
% already, before this code regains control). All lengths here calibrated
% against the rendered PDF's own coordinates; re-measure and adjust if the
% title/body font sizes change.
% \SMSectionframeTrueBoxBottom (2026-07-25, direct instruction: "the
% image extends into the title box" on neb/u97 -- followed by "tyn actually
% looks fine, find out what is really driving this issue/difference"):
% \SMSectionframeBoxHeight above is only the *target* the toppading
% formula below solves for -- by construction that formula makes
% toppading+title-height+title-depth+bottompading algebraically equal
% boxHeight+\SMSectionframeHeightCorrection regardless of title content,
% so the *rendered* bar height should be a fixed, title-independent
% 97.14pt... except it measurably isn't (confirmed directly, PyMuPDF vector
% fill): neb/u97 ("Market Environment"/"Household Finance") render at
% 89.19pt, while tyn/vtm ("Capital asset pricing model"/"Measuring and
% managing risk") render at 85.04pt -- a real, reproducible ~4.15pt
% difference driven entirely by the title text itself, not by anything
% neb/u97 did wrong.
%
% Root cause (traced into gotham.dtx's own
% \SM@frametitlestrut@start/@end, not just inferred): the toppading
% formula's own `-\dp\SMSectionframeTitleBox` term assumes the title's
% real descender depth gets reintroduced somewhere else in the box so it
% nets out -- it doesn't. `@start` is a zero-depth `\rule` whose height is
% `toppading+\totalheightof{X}` (a *fixed reference glyph*, not the real
% title), so the line's rendered height is toppading-driven, independent of
% the title's own real height. `@end` is a zero-height `\rule` of depth
% `\SMFrametitleBottomPadding` (19.90pt for sectionframe) -- also fixed,
% not the real title's own depth. The real title text sits between these
% two struts and contributes its own real height/depth to the same TeX
% line, but a line's rendered depth is the *max* of every contributor, not
% a sum -- and 19.90pt comfortably exceeds any real title's own descender
% depth (measured directly: 0.25pt with none, e.g. "Market
% Environment"/"Household Finance"; 4.42pt with any single descender
% present -- g/j/p/q/y, letter/count doesn't matter, e.g. "Capital asset
% pricing model"/"Measuring and managing risk"/a plain "Testing gjpqy"),
% so bottompading always wins there and the title's real depth is never
% actually given back. Net effect: every point of a title's own real
% descender depth shrinks toppading (and therefore the whole visible box)
% by that same amount, with nothing compensating -- confirmed exactly:
% 4.42pt-0.25pt=4.17pt predicted vs. 89.19pt-85.04pt=4.15pt measured, a
% match well within measurement noise. \SMSectionframeHeightCorrection
% (12.1pt) was presumably calibrated against a descender'd test title (it
% lands tyn/vtm almost exactly on the 85.04pt target), leaving every
% no-descender title's own bar ~4.15pt taller than that target instead.
% Full derivation, and why this means \SMSectionframeBoxHeight was
% never actually title-independent, in README.md's own "Section overview
% (sectionframe)" section.
%
% \SMSectionframeTrueBoxBottom below is the *measured* value for a
% no-descender title specifically (matching neb/u97, the only two slides
% using it) -- it is NOT a general replacement for
% \SMSectionframeBoxHeight: a sectionframe whose title has any
% descender (tyn/vtm, confirmed landing on the *other* value, 85.04pt) must
% keep anchoring against \SMSectionframeBoxHeight directly instead
% (already correct for them, so left untouched). Anything new needing the
% bar's own true rendered bottom should re-measure per its own actual title
% text (real vector fill, not a visual guess, and not an assumption that
% either existing constant transfers) rather than reuse this length
% unless its own title also has zero descenders.
% v5.28.0: the bar height no longer depends on the title (see
% \SMSectionframeRefDepth below), so the "true" bottom is the same 85.04 for
% every title; the old 89.19pt value only fitted titles without descenders.
\newlength{\SMSectionframeTrueBoxBottom}
\setlength{\SMSectionframeTrueBoxBottom}{\SMSectionframeBoxHeight}
% (v9.8.0: \SM@sectionbarmetrics below places the title now, from the bar's own material; the
% four profile numbers named in this and the paragraphs above and below, \SMSectionframeRefDepth,
% \SMSectionframeBottomGap, \SMSectionframeBottomPadCorrection and \SMSectionframeHeightCorrection,
% are gone, and so is the formula these paragraphs explain. They stay as the record of why it was
% done this way for ten versions.)
% \SMSectionframeRefDepth (v5.28.0): the depth the toppading formula assumes for
% the title. It used to subtract the title's own measured depth, which nothing
% adds back, so a title without descenders ("Slides", "Decks and deck
% versions") got a bar 4.15pt taller than one with them, and every other depth
% (p, y, ...) a height in between. With this fixed value (the depth of "p",
% 4.42pt) every title gets the 85.04bp bar the calibrated constants below were
% fitted to, and \SMSectionframeBarBottom (where \sectionfigure starts) is that
% same constant.
% \SMSectionframeBottomGap is the *visual* target -- one huge-font line
% height (baselineskip, measured directly from a deliberately 2-line test
% title: the gap between the two lines' own rendered top edges). It is NOT
% used directly as \SMFrametitleBottomPadding: gotham's bottom strut
% (\SM@frametitlestrut@end) consistently produces a visual gap about
% 5.58pt *smaller* than the padding value fed into it (checked empirically,
% same offset for both a one-line and a two-line test title, so it's a
% fixed per-box offset rather than something that scales with content) --
% \SMSectionframeBottomPadCorrection adds that back.
% Similarly, the box's total rendered height comes out consistently ~11.4pt
% *less* than toppading+measuredTitleHeight+bottompading would suggest --
% the measured title box (built as a plain \vbox{...\par}) doesn't reproduce
% gotham's own leading strut (\SM@frametitlestrut@start, a zero-width
% rule with an explicit height rather than natural interline glue) exactly,
% and this correction closes that gap. Checked the same way as the bottom
% one above: consistent across both the one-line and two-line test titles.
\newlength{\SMSectionframeTitleWidth}
\setlength{\SMSectionframeTitleWidth}{\dimexpr\paperwidth-\SMFrametitleLeftPadding-\SMLogoWidth-2\SMFrametitleRightPadding\relax}
\newbox\SMSectionframeTitleBox
\newlength{\SMSectionframeBarBottom}
% The title's place in the bar (v9.8.0: from the title's own material, no longer from four numbers
% measured at Fira): its last baseline lies the depth of a \strut of the title font plus
% \SMSectionframeTitleGap above the bar's bottom edge, and the top padding is what makes the bar
% \SMSectionframeBoxHeight high. The bar's content is typeset once into a box, by the same
% beamercolorbox the frame title template uses and with a top padding as large as the bar; the box's
% height is that padding plus a constant from there on (any title, any number of lines), so the
% padding that gives the bar its height follows. #1 is the title.
\makeatletter
\newbox\SM@sectionprobe
\newcommand{\SM@sectionbarmetrics}[1]{%
	\begingroup\usebeamerfont{frametitle}\setbox\z@\hbox{\strut}\setlength{\dimen0}{\dimexpr\dp\z@+\SMSectionframeTitleGap\relax}%
	\xdef\SM@sectionbottom{\the\dimen0}\endgroup
	\setlength{\SMFrametitleBottomPadding}{\SM@sectionbottom}%
	\setlength{\SMFrametitleTopPadding}{\SMSectionframeBoxHeight}%
	\setbox\SM@sectionprobe=\vbox{\usebeamerfont{frametitle}%
		\begin{beamercolorbox}[wd=\paperwidth, sep=0pt, leftskip=\SMFrametitleLeftPadding,
			rightskip=\dimexpr\SMLogoWidth+2\SMFrametitleRightPadding\relax plus 1fil]{frametitle}%
			\SM@frametitlestrut@start #1\SM@frametitlestrut@end
		\end{beamercolorbox}}%
	\setlength{\SMFrametitleTopPadding}{\dimexpr2\SMSectionframeBoxHeight-\ht\SM@sectionprobe-\dp\SM@sectionprobe\relax}}
\makeatother
\newenvironment{sectionframe}[2][t]{%
	\begingroup
	\SMManualTopMargin% own title-bar geometry: no automatic Top Margin
	\global\csname SM@topsectiontrue\endcsname% its taller bar: not counted for \SMContentTopShift
	% Disable hyphenation for the title -- both here, in the measurement
	% box below, and (since this is in effect for the rest of the group,
	% which \begin{standardframe} -- triggering the real \insertframetitle
	% -- runs inside) in the actual rendered title too, so a line break
	% never lands mid-word. Setting the penalties alone is the standard,
	% sufficient way to do this (TeX treats a penalty of 10000 at a
	% breakpoint as forbidden, not just heavily discouraged, so it's never
	% chosen even when the alternative is an overfull line); the hyphenmin
	% lengths are redundant insurance, not strictly needed on top of that.
	\hyphenpenalty=10000
	\exhyphenpenalty=10000
	\lefthyphenmin=100
	\righthyphenmin=100
	\setbeamerfont{frametitle}{size=\huge}
	\csname SM@sectionbarmetrics\endcsname{#2}%
	\global\def\SMFrametitleFullHeightYellow{1}%
	\global\def\SMFrametitleLogoFixed{1}%
	\setbeamertemplate{frametitle}[SMNoInlineLogo]%
	\setlength{\SMSectionframeBarBottom}{\dimexpr\SMSectionframeBoxHeight+0.12bp\relax}% 85.14bp for the default bar: the figure starts 0.12bp under it
	% \sectionfigure sets the side-footnote lengths \global (see there); the
	% values from before this frame come back at \end{sectionframe}.
	\edef\SMSectionframeRestoreSideFootnotes{%
		\global\SMSideFootnoteRightInset=\the\SMSideFootnoteRightInset\relax
		\global\SMSideFootnoteHsize=\the\SMSideFootnoteHsize\relax}%
	\begin{standardframe}[#1]{#2}%
	\vspace*{\SMSectionframeContentDrop}%
	\large
}{%
	\end{standardframe}%
	\SMSectionframeRestoreSideFootnotes
	\endgroup
}

% \sectionfigure[includegraphics keys]{image} -- a sectionframe's side image
% (2026-10-01, direct instruction: bij's image should "fit snugly between
% header box and footer and touch the right edge of the slide", for section
% slides with images in general, with the image's source visible). Put it
% anywhere in the sectionframe body. The image is scaled by height to fill
% the gap from the title bar's bottom edge to the footer box's top edge and
% sits flush with the page's right edge -- sideimageframe's geometry, but
% for sectionframe's taller bar. That bar's bottom edge depends on the
% title: the theme's strut gives the last line a fixed depth of ~4.4pt,
% but the padding subtracts the title's real depth -- so the bottom edge is
% 89.44bp minus that depth, but never above 85.04bp (a descender deeper
% than the strut deepens the bar again). Measured in the PDF, i.e. in bp:
% 89.19bp for "Market Environment" (depth 0.25pt), 86.01bp for "Market,
% Environment" (3.44pt), 85.04bp for "Capital asset pricing model"
% (4.42pt) and "Trends in publishing" (4.73pt). (The two
% constants \SMSectionframeTrueBoxBottom/\SMSectionframeBoxHeight
% hold the same numbers in TeX pt, 0.3pt short -- neb/u97 used to add
% that back by hand.) sectionframe computes it as
% \SMSectionframeBarBottom. (Not measured
% from the SMTitleBottom node: the title bar is typeset after the body,
% so that node still belongs to the previous frame while the body runs.)
% The frame's footnotes -- the automatic "Image source(s):" line included --
% go into the rotated side strip immediately left of the image, as on
% sideimageframe. The strip's inset is read at shipout, after the frame body,
% so it and the strip length are set \global here; \end{sectionframe} puts
% the previous values back. Replaces the per-slide \settowidth/
% \SMSideFootnoteRightInset/\SMSideFootnoteHsize set-up of neb, u97,
% tyn and vtm.
\newlength{\SMSectionFigureHeight}
\newlength{\SMSectionFigureWidth}
\newcommand{\sectionfigure}[2][]{%
	\setlength{\SMSectionFigureHeight}{\dimexpr\paperheight-\SMSectionframeBarBottom-\SMFooterBoxHeight\relax}%
	\settowidth{\SMSectionFigureWidth}{\includegraphics[height=\SMSectionFigureHeight, keepaspectratio, #1]{#2}}%
	\vbox to 0pt{\vss
	\begin{tikzpicture}[overlay, remember picture]
		\node[anchor=north east, inner sep=0pt, outer sep=0pt] at ($(current page.north east)+(0,-\SMSectionframeBarBottom)$)
			{\includegraphics[height=\SMSectionFigureHeight, keepaspectratio, #1]{#2}};
	\end{tikzpicture}%
	}%
	\global\SMSideFootnotestrue
	% 3.1pt: sideimageframe's measured gap between strip and image.
	% 2.31pt: the clearance above the footer box that \SMSideFootnoteHsize's
	% 207pt default leaves under a standard title bar.
	\global\SMSideFootnoteRightInset=\dimexpr\SMSectionFigureWidth+3.1pt\relax
	\global\SMSideFootnoteHsize=\dimexpr\SMSectionFigureHeight-\SMSideFootnoteTopSep-2.31pt\relax
}

% quoteframe: a quote slide (ExampleSlides.pptx slide 4) -- institute logo
% as usual, and the quote itself set large/italic/black on a yellow
% full-bleed block, same style as titleframe/sectionframe's own title boxes
% (flush top-left, stopping short of the logo on the right, same margin as
% elsewhere) rather than the dark block an earlier version of this used. The
% usual footer still shows (unlike titleframe, which drops it -- a title
% slide is a one-off outside any particular course content, while a quote
% is still one slide *among* a deck's others, worth being able to reference
% by its footer slide ID like any other). Opening/closing curly quotes are
% added automatically (write the quote text with none of its own); the
% second argument is the person it's from, set smaller/plain, below the
% yellow box and right-aligned to its own right edge (not the page's usual
% text margin) -- pass an empty second argument for an unattributed quote.
%
%   \quoteframe{One of the funny things about the stock market is that
%   	every time one person buys, another sells, and both think they
%   	are astute.}{William Feather, 1889-1981}
%
% \SMQuoteBoxHeight is a plain percentage of the page, unlike
% titleframe's/sectionframe's own (calibrated-to-land-on-an-exact-fraction)
% box heights, since nothing here needs to line up with another element's
% measured position the way those do. \SMTitleframeBoxWidth (defined
% above) is reused as-is for the box's width -- the same
% stop-short-of-the-logo formula, so no reason to duplicate it under a new
% name.
%
% Every piece here -- box, logo, quote text, attribution -- is drawn as a
% zero-space tikz overlay anchored to absolute page coordinates (the same
% \vbox to 0pt{\vss ...} pattern used elsewhere in this file for exactly
% this reason), rather than any of it flowing through the frame's normal
% vertical list. This sidesteps needing to know/counteract beamer's own
% default top-of-frame spacing (which standardframe explicitly measures and
% cancels via \SMContentTopShift) -- there's no such frame-relative
% positioning to get right here, since literally everything is placed by
% explicit coordinate instead. Like the other page-corner overlays in this
% file, "current page" doesn't resolve correctly on a document's very first
% compile pass, so (like those) this needs the usual second pass too.
%
% A *command*, not an environment -- see the identically-reasoned comment on
% titleframe just above for why: the same beamercolorbox-spanning-a-custom-
% environment's-begin/end problem would apply here too, for the same reason.
% \@empty (the empty-attribution check below) needs \makeatletter.
% \quotesource{text} -- opt-in \source{} for \quoteframe (call *before*
% \quoteframe{}{}, same set-then-consume pattern as \chapterlogo{}
% above), needed because \quoteframe is a *command* that auto-closes its
% own \end{frame} -- a bare \source{...} typed after \quoteframe{}{} in
% the calling slide file (as if it were still inside the frame, the way
% \source{} reads everywhere else in this deck) is actually issued
% *outside* any frame, and \source{}'s own \footnote{} then has no page to
% attach to yet. Confirmed directly (2026-07-12, direct instruction:
% traced a spurious extra blank leading overlay step on
% slides/c4c-Macbeth-Spot-Removers-Proposed-Structure-current.tex -- the
% very next multi-overlay frame after
% slides/w8w-Pie-Model-of-Capital-Structure-current.tex's own bare
% \source{} -- down to that orphaned footnote specifically, via a reduced
% two-slide test deck): the dangling footnote insert doesn't render
% anywhere visible on w8w itself, but silently corrupts beamer's own
% overlay-counting pre-pass on whichever frame happens to come right
% after -- reproduced with w8w+c4c alone (bug present) vs. c4c alone (bug
% absent), and confirmed it's specifically about the footnote being
% orphaned, not about w8w's content otherwise, by removing just its
% trailing \source{} line (bug also disappears). A plain \clearpage
% between the two frames also "fixes" c4c but only by giving the orphaned
% footnote a page of its own to attach to -- confirmed directly this
% splits w8w itself into two pages (real content, then a second page with
% nothing but the footnote/footer), not an acceptable trade. The only
% correct fix is keeping the footnote inside the frame that's actually
% citing something, i.e. before \quoteframe's own \end{frame} -- exactly
% what \quotesource{} below does.
\newcommand{\SMQuoteSource}{}
\newcommand{\quotesource}[1]{\gdef\SMQuoteSource{#1}}
\makeatletter
\newcommand{\quoteframe}[2]{%
	\SMOpenFrame{[t]}%
	% Disable hyphenation for the quote text (2026-07-26, direct
	% instruction: "do not allow words to be broken across lines" --
	% confirmed directly first, e.g. on `n3q`'s own long quote,
	% "limi-tations"/"finan-cial" both landed as real mid-word line
	% breaks). Same standard, sufficient technique already used for
	% `sectionframe`'s own title (see that environment's own comment,
	% right before \SMSectionframeTitleBox) -- setting the penalties
	% to 10000 makes TeX treat a hyphenated break as forbidden rather than
	% just discouraged, so one is never chosen even over an otherwise
	% ragged/short line; the hyphenmin lengths are redundant insurance on
	% top of that. Scoped to this frame only via the enclosing
	% \begin{frame}/\end{frame} group, same as sectionframe's own
	% \begingroup/\endgroup scoping -- doesn't leak into any other frame
	% type.
	\hyphenpenalty=10000
	\exhyphenpenalty=10000
	\lefthyphenmin=100
	\righthyphenmin=100
	% \SMArtifactBegin/\SMArtifactEnd (defined with the footer, "PDF
	% accessibility" comment): same reasoning as titleframe's own background
	% box above -- purely decorative, marked so a screen reader skips it. The
	% quote text and attribution below stay unmarked (real content); the
	% footer (this frame keeps one, unlike titleframe) is already marked via
	% the global footline hook.
	\SMArtifactBegin{/Type /Background}%
	\vbox to 0pt{\vss
	\begin{tikzpicture}[overlay, remember picture]
		% \SMTitleBoxColor, not a literal color token -- same setting
		% titleframe's own box and the standard frametitle bar (colorA) use
		% (2026-07-16, direct instruction).
		% outer sep=0pt here and on the logo/attribution below: same hairline-
		% inset-from-the-anchored-corner issue as titleframe's own box -- see
		% the frametitle hook's comment for the full story.
		\node[fill=\SMTitleBoxColor, anchor=north west, inner sep=0pt, outer sep=0pt,
		      minimum width=\SMTitleframeBoxWidth,
		      minimum height=\SMQuoteBoxHeight]
			at (current page.north west) {};
	\end{tikzpicture}%
	}%
	\SMArtifactEnd
	\vbox to 0pt{\vss
	\begin{tikzpicture}[overlay, remember picture]
		% inner ysep derived, not \SMFrametitleRightPadding -- see the
		% frametitle hook's fixed-position logo redraw for why (matches the
		% inline logo's vertically-centered position on every other frame).
		\node[anchor=north east, inner xsep=\SMFrametitleRightPadding,
		      inner ysep=\dimexpr\SMTitleBarBottomAlign-\SMLogoHeight\relax,
		      outer sep=0pt]
			at (current page.north east) {\SMInstituteLogo};
	\end{tikzpicture}%
	}%
	\vbox to 0pt{\vss
	\begin{tikzpicture}[overlay, remember picture]
		% \SMTitleTextColor, not a literal color token -- follows the box
		% color change right above (2026-07-16, direct instruction). Set
		% *inside* the node's own {...} content group (not a stored box
		% register, unlike titleframe's own vbox above), so this doesn't hit
		% the same "\color leaks out of a stored \vbox" pitfall documented on
		% \titleframe's own closing code -- confirmed directly by checking
		% the frame right after this one for the same symptom.
		% The quote text is set in a \parbox instead of the node's own
		% `text width` (2026-09-24), so #1 can hold reveal steps
		% (\conditionalpause, \uncover, \pause). Two separate failures
		% with those directly inside the node, both reproduced in plain
		% beamer and TikZ (no project code): (1) TikZ reads a `text width`
		% node's text through an \aftergroup hook that a \pause/\uncover
		% inside it breaks -- "Package tikz Error: Giving up on this path.
		% Did you forget a semicolon?" on every overlay page; (2) even
		% with (1) avoided, a text-less \pause starts a "covered" state
		% that is only ended by a *global* \beamer@endpause at the next
		% \pause or the end of the frame, so everything after the quote --
		% the footer included -- stayed covered until the last step.
		% \beamer@endpause is therefore run at the end of the quote (and
		% cleared), which confines the covering to the quote itself.
		%
		% The parbox reproduces what TikZ builds for `text width`
		% (minipage[t], \leavevmode\hbox{}, then \tikz@text@action --
		% ragged-right with its own \rightskip/\spaceskip -- and the
		% paragraph ended *after* the \Large group, i.e. with the normal
		% line spacing) so that existing quote slides are unchanged:
		% verified pixel-identical on all 18 other \quoteframe slides.
		\node[anchor=center, inner sep=0pt]
			at ($(current page.north west)+(0.5\SMTitleframeBoxWidth,-0.5\SMQuoteBoxHeight)$)
			{\parbox[t]{\dimexpr\SMTitleframeBoxWidth-2\SMFrametitleLeftPadding\relax}{%
				\leavevmode\hbox{}\tikz@textfont\tikz@text@action%
				{\color{\SMTitleTextColor}\Large\itshape\textquotedblleft #1\textquotedblright%
				\beamer@endpause\global\let\beamer@endpause\@empty}\par}};
	\end{tikzpicture}%
	}%
	\def\SMQuoteAttributionTemp{#2}%
	\ifx\SMQuoteAttributionTemp\@empty\else
		\vbox to 0pt{\vss
		\begin{tikzpicture}[overlay, remember picture]
			\node[anchor=north east, inner sep=0pt, outer sep=0pt]
				at ($(current page.north west)+(\SMTitleframeBoxWidth,-\SMQuoteBoxHeight-\SMQuoteAttributionGap)$)
				{\normalfont\itshape\color{c_custom_gray} #2};
		\end{tikzpicture}%
		}%
	\fi
	\ifx\SMQuoteSource\@empty\else
		\source{\SMQuoteSource}%
	\fi
	\end{frame}%
	% Cleared here, *after* \end{frame}, not inside the frame body above
	% (2026-09-24): beamer re-runs the collected body once per reveal step,
	% so clearing it in the body left only the first step with a source
	% line -- the quote's \quotesource vanished on every later step of a
	% quote with \conditionalpause. Code after \end{frame} runs once.
	\global\let\SMQuoteSource\@empty
}
\makeatother

% figureframe: a slide that is just a figure (ExampleSlides.pptx slide 5) --
% same standardframe layout (title, footer, \source{} for the usual
% "Source: ..."/"Image source: ..." line), but with the figure positioned
% via \figurecenter (call before the figure, inside the frame body) rather
% than left-aligned, which -- like the source deck's own figure slide -- is
% the default:
%
%   \begin{figureframe}{Correlation between DJIA stock returns over 15 years}
%   	\includegraphics[width=0.5\linewidth]{correlation-matrix}
%   	\source{Source: Own illustration, data from Yahoo! Finance, 24.02.2026.}
%   \end{figureframe}
%
%   \begin{figureframe}{Title}
%   	\figurecenter
%   	\includegraphics[width=0.6\linewidth]{diagram}
%   \end{figureframe}
%
% Vertical placement is standardframe's own [t]/[c] (its existing optional
% argument, passed straight through as figureframe's own -- so
% \begin{figureframe}[c]{Title} vertically centers exactly like
% \begin{standardframe}[c]{Title} already does); \figurecenter only ever
% affects the horizontal axis.
%
% \figurecenter is plain \centering under a more discoverable name, nothing
% more -- *not* an opt-in-before-the-frame flag like \chapterlogo{}/
% \sidefootnotes (those must run before \begin{standardframe}, since they
% configure something the frame's own opening code reads). This is called
% *inside* the body instead, exactly where it takes effect, so it can just
% take effect immediately like any other formatting declaration -- no flag,
% no state to reset afterwards, and (checked by first writing it as a flag
% read in figureframe's own opening code, mirroring \sidefootnotes) no way
% for the frame's opening code to see it anyway, since that code runs
% before the body -- and \figurecenter -- does.
\newcommand{\figurecenter}{\centering}
\newenvironment{figureframe}[2][t]{%
	\begin{standardframe}[#1]{#2}%
}{%
	\end{standardframe}%
}

% highlightframe: a standardframe with two things about it made unmistakably
% different -- the title bar's *and* the footer's background color, both
% together -- so a presenter can never flip to it without immediately
% noticing (ExampleSlides.pptx slide 6 was the inspiration for wanting a
% standardframe variant here, though the source slide itself has no such
% color cue; this is a new addition, not a conversion). Everything else --
% body, logo, bullet styling -- is exactly standardframe's own, via the same
% \begingroup/\endgroup wrapper sectionframe above uses for the title-color
% half of this. \SMHighlightColor is a \colorlet, not a color baked
% directly into the environment, so retargeting every highlightframe in a
% deck to a different color -- title bar and footer both -- is still the
% same one-line change it was when this only touched the title bar.
%
% The footer box needs the \global\def dance (\SMFooterBoxColor, defined
% above with the frametitle hook's other opt-in flags, reset in
% standardframe's own closing code) that \setbeamercolor{frametitle} doesn't:
% that assigns through beamer's own color-management system, which handles
% the frame-closing group correctly on its own (checked -- see
% \SMChapterLogoPath's comment for why plain macros need \global\def but
% beamer's own set-commands don't), but the footer box's fill is read from a
% plain macro inside a tikz key, same category as \SMChapterLogoPath.
\colorlet{SMHighlightColor}{c_primary_gray}
\newenvironment{highlightframe}[2][t]{%
	\begingroup
	\setbeamercolor{frametitle}{bg=SMHighlightColor}
	\global\def\SMFooterBoxColor{SMHighlightColor}%
	\begin{standardframe}[#1]{#2}%
}{%
	\end{standardframe}%
	\endgroup
}


% The strip is an overlay (a zero-size box: the Top Margin measurement does not see it), so
% the body below it starts where the title bar's body would -- under the strip. The frame's
% own target line is therefore the Top Margin plus the strip's height
% \SMProgressStripHeight (the strip's bottom edge minus the title bar's: measured 11.08pt in
% the PDF; the node is 11.19pt minimum height): the first line sits as far below the strip
% as it sits below the title bar on any other frame, whatever the body starts with (v5.30.5; the
% frame used to be exempt from the measurement, with a fixed -3pt, which left a paragraph
% 10pt and a list 14pt below the strip against the usual 16pt -- direct instruction, slide 1j2).
\makeatletter
\newcommand{\SM@topstrip}{%
	\ifluatex
		\directlua{slidemage_measure.top_extra(\SM@topthis, \strip@pt\dimexpr\SMProgressStripHeight\relax)}%
	\fi
}
\makeatother
% \progressstrip{Step 1,Step 2,Step 3}{k} -- a full-width strip in the
% institution's tertiary yellow, flush below the title bar, that lists a
% short chain of steps joined by arrows and underlines the k-th one to mark
% where the slide sits in that chain (added 2026-09-19 for the Publishing
% masterdeck's "Field -> Topic -> Question" strip, which the source
% PowerPoint's "Progress" layout draws directly below the title bar; revised
% the same day, direct instruction: edge-to-edge, flush under the title bar,
% smaller font, tight vertical padding, secondary yellow). Written in the
% frame body, as the very first thing. The strip itself is a zero-space
% absolute overlay (so it can span the whole page width, which in-flow text
% cannot); the \vspace after it reserves its height so the body starts just
% below it.
\newcommand{\progressstrip}[2]{%
	\csname SM@topstrip\endcsname% the text below the strip: Top Margin + the strip's height (measured, like any frame)
	\vbox to 0pt{\vss
	\begin{tikzpicture}[overlay, remember picture]
		\node[fill=c_tertiary_yellow, anchor=north west, outer sep=0pt,
		      inner xsep=\SMFrametitleLeftPadding, inner ysep=1.5pt, text height=6.2pt, text depth=2pt,
		      minimum height=11.19pt,
		      text width=\dimexpr\paperwidth-2\SMFrametitleLeftPadding\relax,
		      align=left, text=c_primary_dark, font=\fontsize{6.3}{7.5}\selectfont]
			at ($(current page.north west)+(0,-\SMTitleBarBottomAlign)$)
			{\foreach \SMStepName [count=\SMStepIndex] in {#1}{%
				\ifnum\SMStepIndex>1\ $\rightarrow$\ \fi
				\ifnum\SMStepIndex=#2\underline{\SMStepName}\else\SMStepName\fi}};
	\end{tikzpicture}%
	}%
	\vspace{-3pt}%
}

% \SMCodeFile[listings keys]{path} -- LaTeX source shown on a slide, read from
% a file (2026-10-04, for the SlideMage manual: "show the result and the
% source together"). The code lives in its own small file, e.g.
% snippets/<slide id>-1.tex at the project root, so the slide itself needs no
% [fragile] frame and the code is exactly what was compiled elsewhere. Found on
% the TeX search path like any \input (the project root is on it, see
% .latexmkrc). Optional keys go to \lstinputlisting (firstline=, lastline=,
% basicstyle=...). Fira Mono Regular (the theme's \ttfamily is Fira Mono Bold),
% control sequences, comments and the panel in the SMCode* colours below
% (set per template). Write snippets in ASCII; the few non-ASCII characters below are mapped.
% SlideMage's transcript shows the file as a code block, and the pre-render
% check reports a missing file.
\usepackage{listings}
% Every control sequence, backslash included, is a keyword: the backslash is
% made a letter and is the keyword prefix (listings' prefix keywords need a
% non-empty keyword list to be active, hence the dummy). An escaped \% and a
% \\ are mapped literally (\% must not start a comment; \\ otherwise loses a
% backslash to the \% mapping) and reset the prefix flag, which would
% otherwise colour the word after them.
%
% Code colours (2026-10-05, direct instruction: "set them to something more
% readable in the institutions' template files"): three named colours a
% template defines before % SPDX-FileCopyrightText: 2026 Stefan Palan, Jürgen Fleiß
% SPDX-License-Identifier: LPPL-1.3c+
% A template \input{}s SlideMage_setup (the theme and every template setting
% with its default) before this file; a template from before SlideMage v9.0.0
% does not, and SlideMage's Project compatibility update fixes it.
\ifdefined\SMSetupLoaded\else
	\PackageError{SlideMage}{The template does not \string\input{SlideMage_setup}}{%
		Add \string\input{SlideMage_setup} at the top of the template, before its colors,
		or run Info -> Project compatibility in SlideMage.}
\fi
% \SMLuaOnce{<lua>}: run Lua now AND remember it, so a precompiled format
% (SlideMage's fast preview, see preview_format.py) can replay it. A LuaTeX
% format stores TeX state but not Lua state: every global function, table and
% callback a preamble sets up is gone when the format is loaded. The replay list
% is run by the format's \everyjob; in an ordinary build it is never used.
\newtoks\SMLuaReplayList
\newcommand{\SMLuaOnce}[1]{%
	\directlua{#1}%
	\global\SMLuaReplayList\expandafter{\the\SMLuaReplayList\directlua{#1}}%
}
% The packages this file uses, loaded here explicitly (the theme loads most of them already, which makes
% each \usepackage below a no-op): iftex for \ifluatex, fontspec for \newfontfamily, xcolor for the colors.
\usepackage{iftex}
\usepackage{fontspec}
\usepackage{xcolor}
\usepackage{etoolbox} % \AtEndEnvironment, used for the itemize level-drop spacing below
\usepackage{tikz} % For painting a white box over title bar
\usetikzlibrary{calc}
\usetikzlibrary{arrows.meta} % for e.g. arr/.style={-{Stealth[...]}...} on diagram slides
\usetikzlibrary{positioning} % for `right=<dist> of <node>`-style relative node placement
\usetikzlibrary{decorations.text} % text-along-a-path, for a circular photo's curved caption
\usetikzlibrary{decorations.pathreplacing} % the `brace` decoration shape, for \SMUnderbrace below
% TikZ reveal steps that keep the picture's size (2026-09-30): `visible on=<3->`
% draws an element on steps 3+ and draws it invisibly (opacity 0) before, so it
% still counts toward the bounding box -- unlike \only<3->{\draw ...;}, whose
% element is simply missing on steps 1-2 and can make the whole picture jump
% when it appears (the "diagram shifts downward in step 7" class of fix).
% Same key names as the aobs-tikz library (overlay-beamer-styles), defined here
% so no extra package is needed.
\tikzset{
	invisible/.style={opacity=0, text opacity=0},
	alt/.code args={<#1>#2#3}{\alt<#1>{\pgfkeysalso{#2}}{\pgfkeysalso{#3}}},
	visible on/.style={alt={#1{}{invisible}}},
}
\usepackage{booktabs} % \toprule/\midrule/\bottomrule for SMTable
\usepackage{colortbl} % \arrayrulecolor{}, used by SMTable to color those rules
\usepackage{pifont} % \ding{51}/\ding{55} checkmark/X-mark glyphs, e.g. slides/fr3-Forms-of-Capital-Market-Efficiency-Strong-Semistrong-Weak-current.tex's own validity ratings
\usepackage{pgfplots} % real-data line/bar charts (e.g. slides/3vp-Metallgesellschaft-Oil-Price-Chart-current.tex) -- confirmed safe alongside this deck's own \tableofcontents, see CONVERSION_NOTES.md's "beamer/TeX gotchas" entry
\pgfplotsset{compat=1.18}

% The math fonts (Computer Modern next to the sans text font) and the \text{} wrapper that
% goes with them are font handling: SlideMage_layout_\SMLayoutProfile.tex, input by SlideMage_setup.

% Colors, institute logo path, and the two logo-calibration lengths
% (\SMLogoHeight, \SMFrametitleRightPadding = \SMFrametitleRightPadding)
% are template settings: SlideMage_setup.tex defines each with its default
% (the hook catalog), the template changes what it wants, and this file only
% reads them -- shared machinery, no institution-specific branding. See e.g.
% SlideMage_template_uni-graz.tex for the Graz values. \SMLogoWidth is the one
% logo-derived length that *does* stay computed here (via \settowidth,
% further below) rather than hardcoded per institution, since it already
% adapts automatically to whatever image \SMInstituteLogoPath points at.

% Style hyperlinks (\href{url}{text}) consistently across all slides:
% underlined, so a link is visually recognizable as such against plain
% body text -- yellow (this project's usual accent color) was tried first
% for the *text* itself, but yellow-on-white body text is hard to read,
% unlike yellow used as an accent against dark/black text elsewhere in the
% theme. Went through a few more candidates on 2026-07-06 (all direct
% instruction) before settling on c_secondary_blue for both text
% and underline (a real color reads more clearly as "this is a link" than
% any gray does, without yellow's legibility problem).
%
% 2026-07-07, direct instruction: split text color and underline color
% instead of using one color for both -- link text back to
% c_custom_gray (the near-black dark gray also used for
% quoteframe's attribution text), underline in a separate accent color.
% Plain \underline can't do this: its rule is drawn in whatever color is
% current at that point, so a single \underline{...} call always
% underlines in its own text's color, never a second, different one.
% \SMHrefUnderline below reimplements it instead: measure the linked
% text's own width into a savebox, print that box in customgray, then
% back up by the box's own width and draw a plain \rule of the same
% width, raised to sit just below the text's baseline, in the accent
% color. Same single-line limitation as plain \underline (and the reason
% \url{} is used instead of \href on slides/acw-Some-Topics-in-Testing-the-EMH-current.tex --
% see that slide's own comment): a \rule, like \underline, can't span a
% line break either, so this only ever wraps normal, comparatively short
% \href display text, never something that needs mid-word breaking.
%
% First tried c_primary_highlight for the underline (same day) -- confirmed
% via the compiled PDF that the rule *was* being drawn (right color,
% position, width), so not a compile-time bug, but yellow (very high
% luminance, close to white) at this rule's 0.4pt thickness reads as
% essentially invisible against the white slide background, worst at
% footnote size (see e.g. slides/za2-Startup-Financing-Categorization-current.tex's
% own long inline URL) where it's easy to miss entirely. Switched
% 2026-07-07 (direct instruction) to c_secondary_blue instead --
% far more contrast against white at the same thickness -- and the rule
% was also thickened at the same time (0.4pt -> 0.6pt), on the reasoning
% that it might need it for the same contrast reason. It didn't: blue
% already reads clearly at the original thickness, and 0.6pt instead
% just looked heavier than intended (direct feedback 2026-07-08, e.g. on
% slides/6qi-Yield-Spreads-Corporate-Bonds-current.tex). Reverted to
% 0.4pt -- only the color was ever actually asked to change.
%
% `citecolor` (2026-07-23, direct instruction: "make sure the year in the
% short citation is printed in black", found on
% slides/64m-APT-Choosing-the-Factors-Chen-Roll-Ross-current.tex's own
% `\textcite{}` call, "and generally" -- a deck-wide fix, not a one-off):
% with `colorlinks=true` and no explicit `citecolor`, hyperref falls back
% to its own factory default (green) for every citation hyperlink --
% confirmed directly, rendering 64m with the real biblatex/biber setup
% (this deck's own `\hypersetup` alone, with no citecolor, has no other
% source that could be coloring it): `\textcite{chen1986}` printed as
% "Chen, Roll, Ross (\textcolor{green}{1986})", the year in green, nothing
% else on the line affected (`\textcite` only wraps the year/label itself
% in the actual hyperlink, not the surrounding author names, which are
% plain text). `citecolor=black` set explicitly here, alongside
% urlcolor/linkcolor, rather than left to hyperref's own default, so this
% applies deck-wide to every citation command (`\cite`, `\textcite`,
% `\parencite`, ...) that ends up as a colored link, not just 64m's own
% call site.
\hypersetup{colorlinks=true, urlcolor=c_custom_gray, linkcolor=c_custom_gray, citecolor=black}


\let\SMOrigHref\href
\newsavebox{\SMHrefTextBox}
%
% 2026-07-08 (direct instruction): the underline's own rule thickness now
% scales with the ambient font size instead of a single fixed 0.2pt for
% every link everywhere -- 0.2pt at this deck's footnote size, 0.4pt at its
% normal body-text size, interpolated (and extrapolated, for any other
% size a link happens to appear at) linearly in between.
%
% Deriving "current font size" from \f@size (the textbook NFSS way) was
% tried first and confirmed broken specifically *inside* a beamer frame --
% \typeout{\meaning\f@size} right after this preamble loads shows a normal,
% defined macro (`macro:->10.95`), but the exact same check inside
% `\begin{standardframe}...\end{standardframe}` reports `\f@size` as
% *undefined*, presumably from beamer's own overlay-aware font-switching
% machinery not keeping classic NFSS's own current-size bookkeeping in
% sync the way plain (non-beamer) LaTeX does. Since every real use of this
% macro happens inside a frame, that path was a dead end regardless of how
% well it worked in isolation.
%
% Measuring a reference glyph's own rendered height instead sidesteps this
% entirely: building a box and reading back its own height just measures
% whatever font is *actually* active at the call site, independent of
% whether beamer's own internal "current size" bookkeeping happens to
% agree -- and within one font
% family, every glyph's metrics scale by exactly the same ratio the point
% size does, so a measured height is just as valid a stand-in for "current
% size" as \f@size would have been, had it worked. "M" (present in every
% context this is used in) is the reference glyph; \SMHrefUnderlineFnHeight/
% \SMHrefUnderlineBodyHeight are its height at this deck's footnote font
% and normal body font respectively, each measured once, here, in the
% relevant font context.
% Deliberately built from raw TeX primitives (\newdimen/\setbox/\ht),
% *not* \newlength/\setlength/\settoheight -- found, while building this,
% that loading this deck's theme silently turns \setlength and
% \settoheight into something else entirely: `\meaning\setlength` reads
% `macro:->\protect \setlength` right after `\usetheme{gotham}` (still the
% ordinary kernel definition beforehand), and the same for `\settoheight`.
% Not chased to which file does this or why (bisected only as far as
% "plain beamer alone doesn't, beamer+gotham does" -- most likely
% `translator`, which beamer loads for localized strings and is known to
% robustify commands this way, but not confirmed) -- what matters
% practically is that a *global* assignment through either one silently
% no-ops instead of erroring, confirmed directly and in isolation (a
% minimal beamer+gotham file, no other content): `\global\setlength{\len}
% {\otherlen}` leaves `\len` at 0pt after the group closes, while the
% exact same pattern with plain `\newdimen`/raw `=` assignment propagates
% correctly. Every plain (non-global) `\setlength` elsewhere in this
% project's own files is unaffected -- it's specifically the *global*
% form that breaks, which is why this is the first place it's shown up
% despite `\setlength` being used constantly throughout this preamble.
\newdimen\SMHrefUnderlineFnHeight
\newdimen\SMHrefUnderlineBodyHeight
\newdimen\SMHrefUnderlineNowHeight
\newdimen\SMHrefUnderlineThickness
% A dedicated box register, not the commonly-borrowed box0 -- several
% slides in this deck (the "Long-term bonds" tables) already use box0
% themselves (`\sbox0{...}`/`\usebox0`) for their own, unrelated
% mid-table-layout purposes, and a link can appear anywhere, including
% inside one of those tables -- reusing box0 here would silently clobber
% whatever they were holding there.
\newbox\SMHrefUnderlineBox
% The two reference heights are (re)measured on every single \href call,
% not once at preamble-load time -- tried once-at-load first and found a
% second, independent problem with it: `\usebeamerfont{footnote}` right
% here (before `\begin{document}`) silently has no effect at all (footnote-
% size "M" and normal-size "M" measured identically, confirmed directly),
% presumably because beamer's font-theme machinery isn't fully installed
% until the document starts. Since every real `\href` call happens well
% inside the document (inside some frame), measuring fresh each time
% sidesteps that timing question entirely rather than needing to find the
% right hook to measure once, early enough to be free but late enough to
% be correct -- the extra cost is two throwaway one-glyph boxes per link,
% not something worth optimizing away in a slide deck.
\newcommand{\SMHrefUnderlineMeasureRefHeights}{%
	{\usebeamerfont{footnote}\setbox\SMHrefUnderlineBox=\hbox{M}\global\SMHrefUnderlineFnHeight=\ht\SMHrefUnderlineBox}%
	{\normalsize\setbox\SMHrefUnderlineBox=\hbox{M}\global\SMHrefUnderlineBodyHeight=\ht\SMHrefUnderlineBox}%
}
% A plain e-TeX \dimexpr can multiply a dimension by another dimension and
% then divide by a third, as one combined operation, as long as the
% *middle* operand of the chain is itself wrapped in its own nested
% \dimexpr...\relax -- confirmed directly (an expl3 \dim_set:Nn attempt
% first, using \seq_.../\fp_... the way \setfootnotesymbols above does,
% kept hitting "Missing ) inserted for expression"/type-mismatch errors:
% dividing one dimension by another is a *ratio*, not a dimension, and
% \dim_set:Nn's own expression grammar doesn't implicitly convert between
% the two the way this nested-\dimexpr form does). No expl3 needed here at
% all -- plain e-TeX arithmetic, already available in any current LaTeX.
\newcommand{\SMHrefUnderlineSetThickness}{%
	\SMHrefUnderlineThickness=\dimexpr 0.2pt +
		(\SMHrefUnderlineNowHeight - \SMHrefUnderlineFnHeight)
		* \dimexpr0.2pt\relax
		/ (\SMHrefUnderlineBodyHeight - \SMHrefUnderlineFnHeight)\relax
}
\newcommand{\SMHrefUnderline}[1]{%
	\sbox{\SMHrefTextBox}{#1}%
	\textcolor{c_custom_gray}{\usebox{\SMHrefTextBox}}%
	\hspace{-\wd\SMHrefTextBox}%
	\SMHrefUnderlineMeasureRefHeights
	\setbox\SMHrefUnderlineBox=\hbox{M}\SMHrefUnderlineNowHeight=\ht\SMHrefUnderlineBox
	\SMHrefUnderlineSetThickness
	\raisebox{-0.35ex}[0pt][0pt]{\textcolor{c_secondary_blue}{\rule{\wd\SMHrefTextBox}{\SMHrefUnderlineThickness}}}%
}
\renewcommand{\href}[2]{\SMOrigHref{#1}{\SMHrefUnderline{#2}}}
% \url{} (the `url` package hyperref depends on/patches) defaults to a
% monospace-ish "tt" style regardless of the above -- used deliberately in
% only one place in this deck,
% slides/acw-Some-Topics-in-Testing-the-EMH-current.tex, specifically
% *for* its line-breaking-at-slashes behavior (a plain \href with the
% URL as its own display text has no break points at all -- see that
% slide's own comment). Every other same-text-as-target URL in the deck
% (including slides/bb6-A-Perpetual-Bond-current.tex's own video link,
% which used to be a \url{} too) uses plain \href{url}{url} instead, since
% \url{} doesn't pick up this deck's custom colored underline below (its
% breakable content can't be measured into the savebox the underline
% needs -- see that mechanism's own comment) and most of these URLs are
% short enough to not need the breaking \url{} exists for anyway. 2026-07-07,
% direct instruction: that one genuinely-needed \url{} (on s-acw) should
% still look like this deck's other hyperlinks, i.e. the surrounding body
% font, not \url{}'s own default monospace.
% `\urlstyle{same}` (a standard url-package switch) does exactly that,
% with no effect on \url{}'s own breaking behavior. Set deck-wide, not
% only on that one slide, so any future \url{} use matches by default
% too, rather than needing to remember this each time. Coloring already
% matches automatically (hyperref's own \url respects
% \hypersetup{urlcolor=...} the same way \href does); the underline
% \href gets above doesn't carry over, since \underline's own argument
% can't contain line breaks -- wrapping \url{}'s necessarily-breakable
% text in \underline would silently undo the breaking this exists for,
% so that one difference is accepted rather than "fixed."
\urlstyle{same}



% The colour roles of the theme follow the template's palette: SMColorDark and SMColorA, the canvas
% default and the roles (beamerthemeSlideMage.sty). Needs to go after the template's color redefinitions.
\SMApplyColorRoles
% The hooks that go into registers of beamer's and LaTeX's own (text margins, footnote font
% and spacing): needs to go after the template as well.
\SMApplyLayoutHooks

% Derived from the template's hooks and the font (v9.7.0; they were hand-kept numbers of the catalog
% and the layout profile before). The setup layer holds a macro under each of these four names that
% stops with an error when a template or a profile sets one; a name still holding that macro is
% replaced by the length it stands for, and a name that was redefined is reported.
\makeatletter
\def\SM@derivedlength#1{%
	\expandafter\ifx\csname #1\expandafter\endcsname\csname SM@guard@#1\endcsname\else
		\PackageError{SlideMage}{\expandafter\string\csname #1\endcsname\space is derived, not set}{%
			A template or a layout profile redefined it. SlideMage_base derives it from the hooks and
			the font; remove that line from the template or the profile.}%
	\fi
	\expandafter\let\csname #1\endcsname\relax
	\expandafter\newlength\csname #1\endcsname}
\SM@derivedlength{SMContentTopAlign}
\SM@derivedlength{SMTitleBarBottomAlign}
\SM@derivedlength{SMFootnoteRightShift}
\SM@derivedlength{SMFooterRightShift}
% v9.8.0: the lengths below were numbers of the layout profile (measured once at Fira) or of the
% catalog; each is evaluated in the font where it is used. The footer's, the lists' and the section
% frame's follow with the code that uses them.
\SM@derivedlength{SMFooterTextDrop}
\SM@derivedlength{SMFootlineVOffset}
\SM@derivedlength{SMLabelItemWidth}
\SM@derivedlength{SMLabelItemSep}
\SM@derivedlength{SMEnumerateNestedMinNumber}
\SM@derivedlength{SMEnumerateNestedLabelSep}
\SM@derivedlength{SMEnumerateNestedLeftMargin}
\SM@derivedlength{SMEnumerateLabelBox}
\SM@derivedlength{SMSectionframeRefDepth}
\SM@derivedlength{SMSectionframeBottomGap}
\SM@derivedlength{SMSectionframeBottomPadCorrection}
\SM@derivedlength{SMSectionframeHeightCorrection}
% v9.9.0: the pull-ups of the frames the automatic Top Margin does not correct by their content were
% the last three numbers of the layout profiles. \SMContentTopShift is measured (the frame's origin,
% see AUTOMATIC TOP MARGIN; in a run without \jobname.top, from the first frame that counts on), the two columns
% corrections are measured where a columns environment starts such a frame.
\SM@derivedlength{SMContentTopShift}
\SM@derivedlength{SMColumnsTopCorrection}
\SM@derivedlength{SMColumnsTTopCompensation}
\makeatother
% - \SMContentTopAlign, the Top Margin as an overlay coordinate: the same line the automatic Top Margin
%   puts the first text line's box on (the Top Margin is in bp, the coordinate in pt: 41.58bp is 41.7359pt).
\setlength{\SMContentTopAlign}{\SMTopMargin}
% - \SMTitleBarBottomAlign, the visible bottom of the title bar: where the logo ends. The frame title
%   template centres the logo on the capital height of the title font (\raisebox{.5\fontcharht\font`E-0.5\height}
%   at the baseline, which is the top padding plus the height of an X below the bar's top), so the
%   bottom is the top padding, plus that X, minus half an E, plus half the logo.
\begingroup\usebeamerfont{frametitle}\settototalheight{\dimen0}{X}\dimen1=\fontcharht\font`E
\global\SMTitleBarBottomAlign=\dimexpr\SMFrametitleTopPadding+\dimen0-0.5\dimen1+0.5\SMLogoHeight\relax\endgroup
% - \SMFootnoteRightShift: a footnote ends where the text margin does; the logo ends its right padding
%   from the page's edge, so the footnotes are pushed on by the difference.
\setlength{\SMFootnoteRightShift}{\dimexpr\SMTextMarginRight-\SMFrametitleRightPadding\relax}
% (\SMFooterRightShift is derived with the footer, below.)
% - \SMFootlineVOffset (v9.8.0; the catalog's -0.5ex-0.77pt until then), the space the footline
%   template puts before the footline box. Beamer sets the frame text area to end where the footline
%   box begins (its height plus depth, and the offset, are its \footheight) plus a fixed 4pt
%   (beamerbaseframecomponents.sty, \beamer@calculateheadfoot), and the last footnote line's baseline
%   is the bottom of that area. So the baseline lies
%       \SMFootlineHeight + \SMFootlineDepth + offset + 4pt
%   above the page's bottom edge, and the offset that puts the descenders of that line (the depth of
%   "gjpqy" in the footnote font) \SMFootnoteFooterGap above the footer box follows.
\makeatletter
\begingroup\usebeamerfont{footnote}\settodepth{\dimen0}{gjpqy}
\global\SMFootlineVOffset=\dimexpr\SMFooterBoxHeight+\SMFootnoteFooterGap+\dimen0-4pt-\SMFootlineHeight-\SMFootlineDepth\relax
\xdef\SMPDFInfoFootnoteDescender{\strip@pt\dimexpr\dimen0*7200/7227\relax}\endgroup
\makeatother

% The effective margins, for tools/check_layout.py (SlideMage v9.5.0): custom keys of the PDF's
% info dictionary, in bp (PDF units, what the tool measures in). SMTopMargin is the Top Margin
% (\SMTopMargin), SMLeftMargin the frame title text's left edge (\SMFrametitleLeftPadding). A template
% has set its hooks by now; a slide cannot change them. SMTitleBarBottom (v9.6.0) is the bottom edge of
% the title bar (\SMTitleBarBottomAlign): the tool tells a frame with the ordinary title bar by it.
% SMFontAscender, the text font's ascender, follows once it is known (below, with the Lua file).
% v9.8.0: SMFootnoteFooterGap, SMFootnoteDescender (the depth of "gjpqy" in the footnote font) and
% SMFooterBoxHeight, which tools/calibrate_font_profile.py checks the footnotes' position against.
\ExplSyntaxOn
\edef\SMPDFInfoTopMargin{\dim_to_decimal_in_bp:n{\SMTopMargin}}
\edef\SMPDFInfoLeftMargin{\dim_to_decimal_in_bp:n{\SMFrametitleLeftPadding}}
\edef\SMPDFInfoTitleBarBottom{\dim_to_decimal_in_bp:n{\SMTitleBarBottomAlign}}
\edef\SMPDFInfoFootnoteGap{\dim_to_decimal_in_bp:n{\SMFootnoteFooterGap}}
\edef\SMPDFInfoFooterBoxHeight{\dim_to_decimal_in_bp:n{\SMFooterBoxHeight}}
\ExplSyntaxOff
\hypersetup{pdfinfo={SMTopMargin={\SMPDFInfoTopMargin}, SMLeftMargin={\SMPDFInfoLeftMargin},
	SMTitleBarBottom={\SMPDFInfoTitleBarBottom},
	SMFootnoteFooterGap={\SMPDFInfoFootnoteGap}, SMFootnoteDescender={\SMPDFInfoFootnoteDescender},
	SMFooterBoxHeight={\SMPDFInfoFooterBoxHeight}}}

% Shared diagram roles: templates set the colors once; slides use these styles.
\makeatletter
\@ifundefinedcolor{SMCardBackgroundColor}{\colorlet{SMCardBackgroundColor}{white}}{}
\@ifundefinedcolor{SMCardBorderColor}{\colorlet{SMCardBorderColor}{c_primary_dark}}{}
\@ifundefinedcolor{SMProcessBorderColor}{\colorlet{SMProcessBorderColor}{c_tertiary_gray}}{}
\makeatother
\tikzset{
	SMCard/.style={draw=SMCardBorderColor, fill=SMCardBackgroundColor, text=c_primary_dark},
	SMProcess/.style={draw=none, fill=c_tertiary_gray, text=c_primary_dark},
}

% The theme's divider frames for \part, \section, \subsection and \subsubsection
% (\SMDividerFrames in beamerthemeSlideMage.sty) are off, so \section{} (used
% for a deck's chapter grouping/table of contents) does not insert its own
% plain, unstyled divider frame plus a small "current section highlighted in
% the outline" frame -- redundant with (and visually inconsistent next to)
% this project's own hand-styled sectionframe (see README.md's "Special frame
% types" section), which every chapter in this deck already opens with as an
% explicit \slide{}. \section{} only feeds \tableofcontents (see
% slides/nkj-Agenda-current.tex).
%
% Same for \subsection{} (2026-07-15, direct instruction): \subsection{} calls
% still register in the PDF's own bookmark/outline structure (Acrobat's
% "chapters" pane -- a plain PDF sectioning feature, independent of whether
% beamer auto-inserts a frame for it) and still feed 	ableofcontents (see
% slides/nkj-Agenda-current.tex's own `	ableofcontents[hideallsubsections]`,
% which then hides them from that one slide specifically), without also
% spawning their own unstyled divider slide in the middle of a chapter every
% time one is added.


% A color leak from \titleframe/\quoteframe into whichever frame follows
% them is fixed in standardframe's own opening code below (search
% "Forces this frame's own body") rather than here -- an
% \AtBeginEnvironment{frame}{\color{c_primary_dark}} hook was tried first
% and confirmed, by isolated reproduction, to have *no effect at all*:
% beamer's own "frame" isn't a plain \newenvironment etoolbox's
% \AtBeginEnvironment can patch that way, so the hook silently never fired.


%\colorlet{colorPale}{gPaleYell} % BG in light/normal mode
%\colorlet{colorDark}{gDarkBlack} % FG in light/normal mode
%\colorlet{colorA}{gAnthracite} % frametitle, standin.out,
%\colorlet{colorAreversed}{gLightTeal} % frametitle, standin.in,
%\colorlet{colorB}{gMidGrey} % gray BG : progress bar, blocks
%\colorlet{colorC}{gDeepYellOr} % progress bar
%\colorlet{colorD}{gLightOrange} % alert
%\colorlet{colorE}{gLightGreen} % example

% Opt-in yellow border for an image placed via a tikz node (e.g.
% \node[SMImageBorder] {\includegraphics{...}};) -- matches the
% 3pt/accent1-yellow outline PowerPoint draws directly on a picture shape's
% own border (<a:ln w="38100"><a:solidFill><a:schemeClr val="accent1"/>,
% where accent1 turned out to already be this project's own
% c_primary_highlight, not a new color to map). A tikz *style* (a named
% bundle of node keys), not a wrapper command: adding or removing it is
% just adding or deleting one word from that one \node's own option list,
% no separate macro call/argument needed -- add `SMImageBorder` to draw
% the border, leave it off for a plain image, same as any other tikz key.
%
% Thickness is a single named length, \SMImageBorderWidth, not a bare
% "3pt" baked directly into the style (2026-07-15, direct instruction: "make
% sure the yellow image borders have consistent thickness throughout the
% project by defining a single variable") -- confirmed directly, by
% measuring the actual PDF content-stream "w" operator, that every
% SMImageBorder use already rendered at the exact same 3pt regardless
% (a genuinely shared style was already in effect, no numeric drift found),
% so this is a refactor for future adjustability/discoverability, not a fix
% for an actual inconsistency -- any perceived thickness difference between
% slides traces to something else instead (e.g. a border straddling flush
% image content on one slide and not another, which reads as "heavier" even
% at an identical stroke width -- see the SMImageBorder padding gotcha
% below).
% (\SMImageBorderWidth, 2.5pt, is a hook of SlideMage_setup.tex.)
% outer sep = half the stroke (2026-10-01, direct instruction: "if a figure
% has a border, that border should determine its position"): a bordered
% node's anchors (north east, west, ...) sit on the border's *outer* edge, so
% anchoring it at a margin puts the border on the margin. That is TikZ's own
% default (0.5\pgflinewidth, evaluated with the node's final line width, so a
% slide's line width=1.2pt override still gets 0.6pt); the style sets it
% explicitly so that an "outer sep=0pt" earlier in the option list
% (\cornerfigure's node, for an unbordered image) doesn't take it away. A
% slide that writes outer sep=0pt *after* the style opts out again -- ts8 did,
% and its border stuck out 1.25pt past the Right Margin.
\tikzset{SMImageBorder/.style={draw=c_primary_highlight, line width=\SMImageBorderWidth, outer sep=0.5\pgflinewidth}}

% Named styles for the drawing conventions slides repeat by hand
% (2026-09-30): a restyle -- "make the arrows a bit bolder", "all
% highlights dashed" -- is then one edit here instead of one per slide.
% Each is exactly the option list the slides used, so switching a slide to
% the name changes nothing on the page.
\tikzset{
	SMHighlight/.style={c_custom_red, very thick, rounded corners},   % red box around part of a figure/table (18x)
	SMDottedHighlight/.style={densely dotted, thick, c_custom_red},  % red dotted guide line (19x)
	SMArrow/.style={-{Stealth[length=2mm]}},                         % the usual arrow tip (54x)
}

% \SMTikzmark{name} -- a zero-size remember-picture node at the current
% text position (on the baseline, 2pt up), for arrows/braces that point at
% a word. Three slides defined their own identical copy (\tikzmark on 33r
% and ha9 -- two \newcommand's of one name, an error the moment both land in
% one deck -- and ebq's \ebqTikzmark).
\newcommand{\SMTikzmark}[1]{\tikz[remember picture, overlay, baseline] \node[inner sep=0pt, anchor=base, yshift=2pt] (#1) {};}

% \SMTikzword{name}{text} -- typesets text as usual, but as a named
% remember-picture node, so an overlay arrow can point at the word itself
% ((name.south), (name.north east), ...) and follows it when the layout
% moves (2026-10-01, direct instruction: 3ac's arrow at measured page
% coordinates had drifted off its words). Transcripts keep the text.
\newcommand{\SMTikzword}[2]{\tikz[remember picture, baseline=(#1.base)] \node[inner sep=0pt, outer sep=0pt] (#1) {#2};}

% \SMLineTop -- a zero-width strut as tall as a text line's *line box*
% (\SMFontAscender em above the baseline, the text font's ascender, 0.935 for Fira Sans: the box the Top Margin is
% measured on). Put at the start of the first line of a text column that sits
% next to a figure in top-aligned minipages,
%   \begin{minipage}[t]{0.4\linewidth}\vspace{0pt}
%   	\begin{itemize}\item \SMLineTop First item ...
%   \end{minipage}\hfill
%   \begin{minipage}[t]{0.56\linewidth}\vspace{0pt} <figure> \end{minipage}
% and the text's line box starts exactly where the figure starts -- so the
% automatic Top Margin puts both on it. Without it the first line's box is
% only as tall as its glyphs and the figure lands ~1.8pt below the margin
% (2026-10-01, m2k: its two minipages were vertically *centred*, LaTeX's
% default, and the text column started 12pt below the Top Margin;
% check_layout.py's `coltop` check now reports that pattern).
\newcommand{\SMLineTop}{\rule{0pt}{\SMFontAscender em}}

% \EUR -- the euro sign in Fira Sans, in text and in math alike. In math,
% \text{\texteuro} came out in the math font's style ("Euro symbols ...
% proper", 2026-07-25); slides wrote \text{\sffamily€} by hand (34x).
\newcommand{\EUR}{\ifmmode\text{\sffamily€}\else\texteuro\fi}

% \shrinklists{n} -- all list levels n font steps smaller, for a slide with
% more text than fits (put it inside the frame, like the \setbeamerfont
% lines it replaces; the frame's group resets it). The beamer defaults are
% normal/\small/\footnotesize for levels 1-3; n=1 gives
% \small/\footnotesize/\scriptsize, n=2 \footnotesize/\scriptsize/\scriptsize,
% n=3 \scriptsize on all three levels (the floor; v5.30.0 -- before, every
% n>=2 meant n=2). 13 slides wrote 2-3 \setbeamerfont lines for this, in five
% different combinations.
% The item gaps shrink with the font since 2026-10-01 -- see \SMNoListParsep
% next to the itemize code.
\newcommand{\shrinklists}[1]{%
	\ifcase#1\relax
	\or
		\setbeamerfont{itemize/enumerate body}{size=\small}%
		\setbeamerfont{itemize/enumerate subbody}{size=\footnotesize}%
		\setbeamerfont{itemize/enumerate subsubbody}{size=\scriptsize}%
	\or
		\setbeamerfont{itemize/enumerate body}{size=\footnotesize}%
		\setbeamerfont{itemize/enumerate subbody}{size=\scriptsize}%
		\setbeamerfont{itemize/enumerate subsubbody}{size=\scriptsize}%
	\else
		\setbeamerfont{itemize/enumerate body}{size=\scriptsize}%
		\setbeamerfont{itemize/enumerate subbody}{size=\scriptsize}%
		\setbeamerfont{itemize/enumerate subsubbody}{size=\scriptsize}%
	\fi
}

% The "padding gotcha" promised above, finally root-caused (2026-07-23,
% background-task follow-up to the IA2 fine-tuning batch #2 anomaly on
% slides/m2k-Exercise-CAPM-2-current.tex/
% slides/n8e-Tail-Risk-Variance-Clustering-current.tex/
% slides/fm6-Tail-Risk-Market-vs-Simulated-current.tex/
% slides/xqq-Fat-Tails-current.tex, all fixed there with a pragmatic
% `line width=1.2pt` override and an explicitly *unresolved* root cause at
% the time -- "declares an identical stroke width and CTM in the compiled
% PDF's own raw content stream, yet renders at a genuinely different pixel
% thickness", with CTM scale/clips/dash/cap/join all directly checked and
% ruled out). Confirmed via a controlled reproducer (four `SMImageBorder`
% nodes -- plain opaque image, nested pgfplots axis, plain image at
% `inner sep=3pt`, nested plain tikzpicture at `inner sep=3pt` -- rendered
% at 600dpi and measured by a half-max cross-section profile, not just
% eyeballed): the one variable that actually predicts the rendered width is
% `inner sep`, not content type at all. At `inner sep=0pt` a *fully opaque*
% image measured 1.2pt; the exact same image at `inner sep=3pt` measured
% 2.52pt (matching the declared 2.5pt); a nested pgfplots axis and a plain
% nested tikzpicture, both at `inner sep=3pt`, also measured 2.52pt. Two
% more points confirmed it's a smooth transition, not a binary switch:
% `inner sep=0.6pt` measured 1.8pt (= the 1.2pt floor plus that 0.6pt gap,
% exactly); `inner sep=1.25pt` (= \SMImageBorderWidth/2, the stroke's
% own half-width) already measured the full 2.52pt.
%
% The mechanism: tikz draws a node's border stroke centered on the node's
% own boundary path, then paints the node's *content* on top of it,
% afterwards. At `inner sep=0`, the content's own bounding box coincides
% exactly with that boundary path, so if the content is opaque all the way
% to its own edge, it paints directly over the stroke's inner half,
% visually leaving only the outer half (`\SMImageBorderWidth`/2,
% ~1.2pt) -- not a rendering bug, just tikz's normal draw order interacting
% with zero padding. A gap of at least `\SMImageBorderWidth`/2 (or
% content that's transparent right at its own edge -- see below) leaves the
% border with nothing opaque to cover it, so it renders at its true
% declared width instead.
%
% This also explains why the fix landed on exactly the 4 slides it did,
% each for a different reason: slides/m2k-Exercise-CAPM-2-current.tex
% used `inner sep=3pt` for unrelated layout reasons (room for the axis's
% own tick labels), well past the 1.25pt threshold, so nothing ever covered
% its border. slides/n8e-Tail-Risk-Variance-Clustering-current.tex/
% slides/fm6-Tail-Risk-Market-vs-Simulated-current.tex/
% slides/xqq-Fat-Tails-current.tex all use `inner sep=0pt` -- matching the
% deck's own overwhelming convention -- but their own chart PNGs
% (`chart-index-daily-returns-*.png`/`chart-index-fat-tails-histogram*.png`,
% matplotlib exports with a transparent background) measured alpha=0 at
% every edge pixel (confirmed directly, Pillow), so there was no opaque
% content to cover their own border either. Every other `SMImageBorder`
% use in this project that was never flagged uses an opaque image (a JPG,
% or a PNG with no transparent margin) at `inner sep=0pt` -- which
% *accidentally* renders at the ~1.2pt covered-half width instead of the
% full 2.5pt this constant nominally declares, and that accidental width is
% what reads as "the deck's normal border thickness" throughout, since it's
% by far the most common case.
%
% Practical upshot for a *new* slide: a `SMImageBorder` node with
% `inner sep=0pt` wrapping opaque content (the common case) needs no
% special handling -- it will match the rest of the deck automatically. A
% node that instead uses a nonzero `inner sep`, or wraps content with a
% transparent edge (a matplotlib export, an SVG-derived PDF with padding,
% ...), will render at the *full* \SMImageBorderWidth and look
% noticeably heavier next to the deck's other borders -- either add an
% explicit `line width=\SMImageBorderWidth/2` (matching the 4 fixed
% slides' own literal 1.2pt) to match the accidental deck-wide look, or
% reduce `inner sep` below \SMImageBorderWidth/2 so the same coverage
% effect kicks in naturally. The 4 existing per-slide overrides above are
% left as-is -- they're already the correct fix for what's now a fully
% understood cause, not a mystery workaround needing revisiting.

% Reusable "mugshot" -- a circular-cropped photo, bordered like any other
% SMImageBorder image, with a caption curving along its own bottom arc
% (2026-07-06, direct instruction, generalized from
% slides/4uq-Definition-of-Capital-Market-Efficiency-current.tex's own
% Eugene Fama photo -- see that slide's earlier comment for the source
% PPTX effect this reproduces, `<a:prstGeom prst="ellipse">` +
% `<a:prstTxWarp prst="textArchDown">`, and CONVERSION_NOTES.md's own
% "beamer/TeX gotchas" entry, search "circular photo crop", for the
% clip/border/decorations.text mechanics themselves -- unchanged here,
% just factored out so a second photo doesn't need to repeat them).
%
%   \coordinate (myMugshotCenter) at (...);
%   \SMMugshot{myMugshotCenter}{media/photo-example.jpeg}{2.6cm}{Some Person, NP 2013}
%
% #1 the *name* of a coordinate the caller has already placed (not a raw
% point expression -- needed as a plain name, not just a point, since the
% caption's own arc path re-references it by name via `++(...)`, which
% only works with a named coordinate/node, not an arbitrary `($...$)`
% calc expression); #2 the image path; #3 the desired final diameter;
% #4 the caption text.
%
% Scales the source photo so it fully covers the circle with no gap
% anywhere -- both dimensions end up at least as big as the diameter --
% while using as much of the source photo as possible (2026-07-07,
% direct instruction: "should not include whitespace... scale the photo
% such that its edges just touch the circle's inner border"). This
% briefly (same day) scaled the *other* way -- by the photo's own
% diagonal, the largest scale at which the *entire* rectangle fits
% inside the circle with nothing clipped at all -- which does avoid
% ever clipping real content, but leaves the circle's own four corners
% uncovered (visible as plain background there, since a non-circular
% photo can only touch an inscribing circle at its corners, not along
% the whole rim) -- exactly the "whitespace" this round's feedback
% ruled out in favor of full coverage. Covering scale is
% `diameter / min(w,h)`: whichever of the photo's own two dimensions is
% *smaller* becomes exactly the diameter (that edge pair touches the
% circle's own inner border, per the direct instruction), and the other,
% larger dimension overhangs the diameter and gets clipped by the
% circular `\clip` -- for a portrait photo (every one used here so far),
% that means width becomes the diameter and a sliver clips off the top
% and bottom, which is the origin of the clipping problem the diagonal
% approach above existed to solve in the first place. Reintroduced
% deliberately: showing as much of the photo as the "no whitespace"
% constraint allows is now the stated priority, over never clipping
% anything.
%
% Measuring the photo's own native size needs `\savebox` (into
% \SMMugshotBox, declared once below, not a fresh box register per
% call) -- `\includegraphics` with no size option at all renders at the
% image's native resolution, and `\wd`/`\ht` then read that box's actual
% dimensions. pgfmath (already loaded for tikz) converts any TeX
% dimension fed into it to a plain pt-count number automatically, so
% arithmetic on the two measured lengths just works. `min(1,ratio)`
% (ratio = height/width) picks out whichever of the two is smaller
% *relative to the other*, without needing a `sqrt`/`^` at all this time
% -- so this version doesn't hit the "! Dimension too large" pgfmath
% overflow the diagonal version's own arithmetic did (see this file's
% own comment history / CONVERSION_NOTES.md's "beamer/TeX gotchas" for
% that point (easy to misread as a harmless warning, since compilation
% continues and produces *a* number regardless). Reformulating via the
% aspect ratio sidesteps the whole problem: `scale = diameter /
% sqrt(1+ratio^2)` is the same formula algebraically (factor `w` out of
% `sqrt(w^2+h^2)` and it becomes `w*sqrt(1+(h/w)^2)`), but every
% the full story of that overflow, from when this macro's own arithmetic
% did use a `sqrt`).
%
% Caption font reverted to this deck's normal body font (2026-07-07,
% direct instruction: "change the mugshot template back to using regular
% font instead of monospace") -- `\ttfamily` (Fira Mono) was a direct
% instruction on 2026-07-06, one day earlier; \normalfont resets
% explicitly rather than just omitting \ttfamily, so this doesn't
% silently inherit whatever font happens to be active at the call site.
\newsavebox{\SMMugshotBox}
\newcommand{\SMMugshot}[4]{%
	\savebox{\SMMugshotBox}{\includegraphics{#2}}%
	\pgfmathsetmacro{\SMMugshotImgW}{\wd\SMMugshotBox}%
	\pgfmathsetmacro{\SMMugshotImgH}{\ht\SMMugshotBox}%
	\pgfmathsetmacro{\SMMugshotRatio}{\SMMugshotImgH/\SMMugshotImgW}%
	\pgfmathsetmacro{\SMMugshotScaledW}{(#3)/min(1,\SMMugshotRatio)}%
	\begin{scope}
		\clip (#1) circle (0.5*#3);
		\node[anchor=center, inner sep=0pt] at (#1)
			{\includegraphics[width=\SMMugshotScaledW pt]{#2}};
	\end{scope}
	\draw[SMImageBorder] (#1) circle (0.5*#3);
	{\normalfont\tiny
	\path[decoration={text along path, text={#4}, text align=center}, decorate]
		(#1) ++(200:0.5*#3+0.25cm) arc (200:340:0.5*#3+0.25cm);
	}
}

% Reusable "warning sign" -- a rotated, no-fill, thick-bordered box with
% centered bold text, generalized (2026-07-24, direct instruction: "build a
% GENERAL reusable command for placing such signs") from the red
% "Extremely skewed outcomes!" callout on slides/awe-Example-Investment-
% Return-Simulation-current.tex's own source PPTX slide (233) -- read
% directly from that slide's `ppt/slides/slide233.xml`: a rotated
% no-fill rectangle (`rot="20410887"`, i.e. 340.18 degrees clockwise =
% -19.82 degrees), a thick solid border in the same color as the text
% (`<a:ln w="57150">` = 4.5pt, color C00000 = this project's own
% `c_custom_red`), text centered inside.
%
% Usage: \SMWarningSign[color]{x}{y}{text}{aspect ratio}{rotation}
% #1 (optional) the border/text color, defaults to \c_custom_red; #2/#3
% the sign's own center, as a page-absolute (x,-y) offset from
% "current page.north west" in pt -- the same coordinate convention as
% this deck's other absolute-overlay figures (e.g.
% slides/ivw-Divestment-of-Polluting-Plants-Overview-current.tex's own
% photo nodes); #4 the sign's text; #5 the desired width:height aspect
% ratio of the box; #6 the rotation angle in degrees, tikz convention
% (counterclockwise positive) -- note OOXML's own `rot` attribute on a
% source PPTX shape is clockwise positive instead, so a value read
% directly off a source slide needs negating first.
%
% The box is sized to fit the text *wrapped* at the requested aspect
% ratio, not a caller-supplied width -- matching how the source PPTX's
% own sign isn't a single unwrapped line either (rendering the actual
% source deck's own slide 233 to a flat image for comparison confirmed
% its text wraps across 3 lines inside a squarish box, not one long
% strip). #5 drives that wrap width via a closed-form estimate (start
% from the unwrapped text's own natural width #4 and single-line height
% #4, i.e. `sqrt(ratio*width*height)`, the width at which a box of
% roughly that many lines would hit the target ratio if text-ink area
% stayed constant under rewrapping) plus one correction pass (rewrap at
% that estimate, remeasure the real height, then solve
% `sqrt(width*height*ratio)` again treating the *measured* width*height
% as the area instead of the rough estimate) -- not exact, since
% rewrapping to a handful of lines isn't a smooth function of width, but
% close enough for the short warning-sign text this is meant for; a
% caller can always nudge #5 if a specific case lands visibly off.
% (\SMWarningSignPadding, 6pt, and \SMWarningSignBorderWidth, 4.5pt, are hooks of SlideMage_setup.tex.)
\newsavebox{\SMWarningSignBox}
\newlength{\SMWarningSignW}
\newlength{\SMWarningSignH}
\newcommand{\SMWarningSign}[6][c_custom_red]{%
	\savebox{\SMWarningSignBox}{\bfseries\color{#1}#4}%
	\pgfmathsetlength{\SMWarningSignW}{sqrt((#5)*\wd\SMWarningSignBox*(\ht\SMWarningSignBox+\dp\SMWarningSignBox))}%
	\savebox{\SMWarningSignBox}{\parbox{\SMWarningSignW}{\centering\bfseries\color{#1}#4}}%
	\pgfmathsetlength{\SMWarningSignW}{sqrt((#5)*\SMWarningSignW*(\ht\SMWarningSignBox+\dp\SMWarningSignBox))}%
	\savebox{\SMWarningSignBox}{\parbox{\SMWarningSignW}{\centering\bfseries\color{#1}#4}}%
	\setlength{\SMWarningSignH}{\ht\SMWarningSignBox+\dp\SMWarningSignBox}%
	% \vbox to 0pt{\vss ...} keeps this at zero space in the normal vertical
	% list, the same reason every other absolute-page-coordinate tikz overlay
	% in this file uses it (footer/logo/title-bar boxes, y5z's/22s's own
	% corner images, ...): a bare tikzpicture here is otherwise sized from
	% its own local origin (wherever this macro happens to be called in the
	% flow) out to the far-away absolute page coordinate the node is drawn
	% at, so without this wrapper the "sign" silently pushes real, sizeable
	% blank space into the frame's own body -- confirmed directly on
	% slides/awe-Example-Investment-Return-Simulation-current.tex: with this
	% call gated behind `\deckreveal[3]{...}`, the reveal step where it first
	% turns on pushed the *unrelated* left-column text down by about 100pt
	% relative to the steps before it, a "content moves between reveal
	% steps" bug rather than anything wrong with that text itself.
	\vbox to 0pt{\vss
	\begin{tikzpicture}[overlay, remember picture]
		\node[rotate=#6, draw=#1, line width=\SMWarningSignBorderWidth,
		      minimum width=\dimexpr\SMWarningSignW+2\SMWarningSignPadding\relax,
		      minimum height=\dimexpr\SMWarningSignH+2\SMWarningSignPadding\relax,
		      inner sep=0pt] at ($(current page.north west)+(#2,-#3)$)
			{\usebox{\SMWarningSignBox}};
	\end{tikzpicture}}%
}

% A \frac/\dfrac-stacked formula reads as cramped against a plain-text
% neighbor in an itemize -- it's visibly taller than a normal line (both
% above *and* below the line, since the fraction bar sits on the text
% baseline with numerator/denominator stacked around it), so both the
% fraction's own item *and* the item that follows it need a bit of extra
% headroom that a plain-text \item doesn't. Started as a manually-repeated
% \vspace{3pt} before each such \item (first on
% slides/u74-Example-Dynamic-Mattress-Company-Periods-current.tex);
% \tallitem below is that same convention factored into a single drop-in
% \item replacement, so a future fraction-heavy slide (or a change to the
% spacing amount) doesn't mean re-editing several slide files by hand.
% Takes an optional leading multiplier (default 1, i.e. 3pt): \tallitem[1.5]
% for 4.5pt, for a formula tall enough that the default doesn't read as
% enough room. Usage: replace \item with \tallitem both on any itemize line
% containing a stacked fraction *and* on the item immediately following it
% (same nesting level; skip this second replacement if that next item is
% itself already a \tallitem for its own fraction, or if the fraction item
% is the last one in its list) -- it's a plain substitute, not an
% environment, so it doesn't change how the surrounding list itself is
% structured. (Named \SMFracItem originally; renamed 2026-07-21 to the
% more discoverable \tallitem, and given the optional multiplier and the
% following-item rule above at the same time -- no behavior change for any
% existing default-strength call site.)
% Ratio, not a frozen length (2026-07-22, direct instruction: "modify
% \tallitem and \tallitemnobullet to also use scaling relative to font
% size") -- the original \newlength+\setlength{...}{3pt} pair freezes to
% one fixed pt value the moment \setlength runs, the same issue (and same
% fix) as \SMItemsepRatioI/II/III/IV further below: converted to a
% bare ratio macro, with the "em" unit written back on inside \tallitem's
% and \tallitemnobullet's own bodies (not here), so it's recomputed fresh
% against whatever font is actually live at each call site -- correct
% regardless of nesting level or a \footnotesize/\scriptsize override,
% unlike the old fixed 3pt (calibrated against one specific level's font,
% silently wrong at any other). 0.3 reproduces the original 3pt exactly
% at level 2's own font (10pt) -- the level \tallitem was first used at,
% on slides/u74-Example-Dynamic-Mattress-Company-Periods-current.tex (see
% this macro's own comment above) -- so no visual change there, while now
% correctly scaling up/down at every other level and font override.
% (\SMTallItemBaseRatio, 0.3, is a hook of SlideMage_setup.tex.)
\newlength{\SMTallItemBaseUnit}
\newlength{\SMTallItemSpace}
% The optional argument is either a number (the multiplier above -- any
% argument made only of digits and dots) or anything else, which is taken as
% a bullet label and passed on as \item[<label>] with the default spacing:
% \tallitem[\SMBulletEllipsis] gives an ellipsis-labelled item with the
% usual headroom (2026-09-26, direct instruction), \tallitem[] a label-less
% one. (For a label-less item with a custom multiplier, use
% \tallitemnobullet[<multiplier>].)
\ExplSyntaxOn
\cs_new_protected:Npn \SMTallItemIfNumberTF #1
	{ \regex_match:nnTF { \A [0-9.]+ \Z } {#1} }
\ExplSyntaxOff
\newcommand{\tallitem}[1][1]{%
	\setlength{\SMTallItemBaseUnit}{\SMTallItemBaseRatio em}%
	\SMTallItemIfNumberTF{#1}%
		{\setlength{\SMTallItemSpace}{#1\SMTallItemBaseUnit}%
		 \vspace{\SMTallItemSpace}\item}%
		{\setlength{\SMTallItemSpace}{\SMTallItemBaseUnit}%
		 \vspace{\SMTallItemSpace}\item[#1]}}

% For a line that needs \tallitem's own extra headroom but no bullet at
% all -- a display equation continuing the previous \item's own sentence,
% rather than a new bulleted point of its own (see README.md's own
% \tallitem section, which already documents this exact combination as a
% rare, previously hand-written case: "\vspace{\SMTallItemBaseUnit}\item[]").
% \tallitemnobullet (2026-07-21, direct instruction) factors that hand-
% written idiom into the same kind of drop-in replacement \tallitem
% already is for its own case, same optional multiplier: \tallitemnobullet[1.5].
\newcommand{\tallitemnobullet}[1][1]{%
	% Recomputes \SMTallItemBaseUnit itself, fresh, rather than reusing
	% whatever \tallitem's own body last left it at -- since it's now a
	% per-call, live "em"-based value (see \tallitem's own comment above),
	% not the one-time preamble-load-time constant it used to be, relying
	% on call order/a preceding \tallitem to have set it correctly would
	% silently break the moment \tallitemnobullet is used on its own, or
	% at a different font size than whichever \tallitem last ran.
	\setlength{\SMTallItemBaseUnit}{\SMTallItemBaseRatio em}%
	\setlength{\SMTallItemSpace}{#1\SMTallItemBaseUnit}%
	\vspace{\SMTallItemSpace}\item[]}

% flalign* (used to left-align a display equation inside an itemize --
% see CONVERSION_NOTES.md's "flalign* for left-aligning..." entry) inherits
% plain LaTeX's \abovedisplayskip/\belowdisplayskip, sized for a full
% document page, not a compact slide itemize -- reads as far too much blank
% space above/below the equation next to the surrounding item text (found
% 2026-07-21, direct instruction; most visible on
% slides/bwi-Pearson-Correlation-Definition-current.tex's own single-line
% equation, but the same oversized gap is present everywhere flalign* is
% used this way, e.g. slides/7se-Portfolio-Risk-Three-Assets-current.tex
% and slides/rb3-Asset-Allocation-Example-CAL-current.tex). \SMTightDisplaySkips
% tightens both skips -- and their "short" variants, since amsmath picks
% between the two depending on how much text precedes the display, so
% setting only one pair left the gap unfixed whenever the preceding line
% happened to be short -- down to something that reads as part of the
% same list rather than a separate floating display. Usage: wrap the
% *real*, unmodified flalign* in a plain TeX group together with this
% command, `{\SMTightDisplaySkips\begin{flalign*} ... \end{flalign*}}`,
% rather than a new same-named environment -- tried that first
% (`\newenvironment{SMTightFlalign*}{...\begin{flalign*}}{\end{flalign*}}`)
% and it broke LaTeX's own environment-nesting tracking (`\@currenvir`),
% surfacing as "\begin{flalign*} ... ended by \end{SMTightFlalign*}"
% and a fatal, PDF-less compile on the very first slide that used it. A
% plain group sidesteps that class of bug entirely since it never
% introduces a same-named wrapper environment.
% (\SMTightDisplaySkip, 3pt, is a hook of SlideMage_setup.tex.)
\newcommand{\SMTightDisplaySkips}{%
	\abovedisplayskip=\SMTightDisplaySkip
	\abovedisplayshortskip=\SMTightDisplaySkip
	\belowdisplayskip=\SMTightDisplaySkip
	\belowdisplayshortskip=\SMTightDisplaySkip}


% \SMFrametitleRightPadding (left/right padding around the logo, chosen
% to match its top/bottom margins -- measured against the title colorbox's
% own top and bottom edges -- rather than the theme's 2.2ex default) is
% also set by the per-university preamble, not here: the top/bottom margins
% are a side effect of vertically centering a \SMLogoHeight-tall logo in
% the title bar, not something set directly, so this value has to be
% re-measured and set per institution whenever its own logo/height differs
% from another's.

% The logo macro \SMInstituteLogo, the frame title hook that paints the white patch behind the
% logo and redraws it (\SMFrametitleHookCode), \chapterlogo and the template without the inline
% logo (SMNoInlineLogo) are the theme's (beamerthemeSlideMage.sty). The logo width, which the
% hook uses to size the white patch, is measured here from the template's logo path and height.
\settowidth{\SMLogoWidth}{\includegraphics[height=\SMLogoHeight]{\SMInstituteLogoPath}}



% Define slideID as a simple value holder (initially empty)
\newcommand{\slideID}{}%

% Set (locally, by \slidestep) when this inclusion of a slide id must not define
% the id's \hypertarget again, because an earlier \slidestep of the same id
% already owns it (see \slidestep's own comment).
\newif\ifSMSuppressSlideTarget

% Define slideID as a simple value holder (not a command with arguments)
\makeatletter
\newcommand{\setSlideID}[1]{%
	% Grab everything up to the first '-' and store it in \slideID
	\edef\tempa{#1}%
	\expandafter\setSlideID@aux\tempa-\relax
}
\def\setSlideID@aux#1-#2\relax{%
	\renewcommand{\slideID}{#1}%
	% Non-linear navigation (direct instruction, 2026-09-15): every real
	% slide inclusion -- \slide{}, \slidefinal{} (delegates to \slide{}),
	% \slideversion{}, \omitslide{}, \omitslideversion{}, \conditionalslide{}
	% -- calls \setSlideID{} immediately before \input{}ing the slide file,
	% regardless of which frame type (standardframe, sideimageframe,
	% figureframe, sectionframe, highlightframe, ...) that file happens to
	% use internally. This is the one place all of them funnel through, so
	% it is the correct place for a single, frame-type-independent
	% \hypertarget{} -- giving every slide a free, permanent jump target
	% keyed to its own existing id, with no per-slide authoring and no new
	% bookkeeping. Deliberately \hypertarget (hyperref's own mechanism, not
	% plain LaTeX \label{}: slides/ZZM-EN-Slidemaster-colors-current.tex's
	% own comment documents a real, already-hit collision between \label{}
	% and pgffor's \foreach inside a tikz loop, which this sidesteps
	% entirely. Used by \hyperlink{<bare id>}{...} anywhere in the deck
	% (README.md's own "Non-linear navigation" section) to jump straight to
	% this slide -- e.g. an excursus link's own "Back" button, or a
	% generated appendix's thumbnail grid (SlideMage's "Generate
	% appendix TOC" checkbox, appendix_toc.py).
	%
	% Known, pre-existing limitation, not introduced by this: if the same
	% id is genuinely \slide{}'d twice in one deck (it happens -- e.g. a
	% \slidefinal{} recap of an earlier slide), hyperref just warns
	% ("multiply defined") and the *later* occurrence's target wins, so a
	% link to that id lands on whichever instance is last in the document
	% -- inherent to any permanent-id-keyed target scheme, not specific to
	% this one.
	\ifSMSuppressSlideTarget\else\hypertarget{\slideID}{}\SM@agenda@slide{#1}{#2}\fi%
}
\makeatother

% \setSlideVersion{X-Y-Z} -- footer's own version display (added
% 2026-08-20, direct instruction: show the slide's version next to its id
% in every compiled deck). Takes the same filename-safe hyphenated form a
% version file's own "-vX-Y-Z.tex" suffix and \slideversion{}'s own second
% argument already use, and reformats it with dots ("1.0.0") for display,
% the conventional way semver is written -- \setSlideVersion@aux splits it
% the same way \setSlideID@aux above splits an id-slug, just at both
% hyphens instead of the first one only. \ifSMSlideVersionSet tracks
% whether *this* frame actually has a tracked version to show at all --
% reset false at the top of every \slide{}/\omitslide{}/\conditionalslide{}
% call (current.tex may or may not embed a \setSlideVersion{} of its own,
% see SlideMage's own automatic-versioning save flow) so an untracked
% slide's footer never inherits whatever the immediately preceding frame
% last set. \slideversion{}/\omitslideversion{} instead set it directly
% from their own second argument -- already an explicit version number at
% the call site, so there is nothing to read from the pinned file itself.
\newcommand{\slideVersion}{}%
\newif\ifSMSlideVersionSet
\makeatletter
\newcommand{\setSlideVersion}[1]{%
	\edef\SMTempSlideVersion{#1}%
	% No trailing "-" appended before \relax here, unlike \setSlideID@aux's
	% own \tempa-\relax above -- that extra hyphen is harmless there since
	% \setSlideID@aux's own #2 (everything after the first "-") is thrown
	% away entirely, but this pattern's #3 is a *used* greedy match
	% ("everything up to \relax"), so an appended "-" would land inside it
	% verbatim (confirmed directly: produced "v2.3.1-" instead of
	% "v2.3.1" the first time this was compiled).
	\expandafter\setSlideVersion@aux\SMTempSlideVersion\relax
	\global\SMSlideVersionSettrue
}
\def\setSlideVersion@aux#1-#2-#3\relax{%
	\renewcommand{\slideVersion}{#1.#2.#3}%
}
\makeatother

% \setSlideLanguageFromSlug{ID-slug} -- multi-language support (added
% 2026-08-24). Language lives in the id-slug's own *second* hyphen segment
% rather than a separate in-body macro call, so it's visible directly in
% Explorer/grep and in every \slide{}/\slideversion{}/etc. call site, not
% just inside the file it names: slides/<id>-EN-<slug>-current.tex,
% slides/<id>-DE-<slug>-current.tex, slides/<id>-mult-<slug>-current.tex
% (mult = multilingual, i.e. a single file holding both languages via
% \multilingualtext{}{} below, picked at compile time by \ifSMGerman).
% A slide that predates this feature simply has no such segment (its own
% second hyphen segment is the first word of an ordinary slug, e.g.
% "neb-Market-Environment-...") -- \ifSMSlideLanguageSet stays false
% for it, exactly like an untracked \setSlideVersion{} today shows no
% "v..." footer segment, so the other ~665 files need no migration.
% Same three-way split idiom as \setSlideVersion@aux above, but the
% trailing padding ("-zzz-zzz") is load-bearing here in a way it isn't
% there: \setSlideVersion@aux's own caller (X-Y-Z) always has exactly two
% hyphens, but an id-slug can legally have as few as *one* (e.g. a
% single-word slug like "abc-Word"), which would leave #2-#3\relax with no
% closing delimiter to match at all ("Use of ... doesn't match its
% definition") without this padding. #3 (and the padding itself) is never
% read, so its exact content doesn't matter -- only that the two extra
% hyphens are guaranteed to exist.
\newcommand{\slideLanguage}{}%
\newif\ifSMSlideLanguageSet
% \ifSMSlideLanguageMult -- set alongside \slideLanguage whenever it's
% specifically "mult", so the footer tag below never has to run
% \ifstrequal against \slideLanguage itself: etoolbox's \ifstrequal
% compares via \detokenize, which does not expand its argument -- fed a
% bare control sequence like \slideLanguage it would stringify the
% *macro's name*, not its meaning, and always take the false branch (see
% \decklanguage's own comment above, where this was caught the same way).
% A dedicated boolean, set right here where the value is still a plain
% matched string (#2, not a macro), sidesteps the trap entirely.
\newif\ifSMSlideLanguageMult
\makeatletter
\newcommand{\setSlideLanguageFromSlug}[1]{%
	\global\SMSlideLanguageSetfalse
	\global\SMSlideLanguageMultfalse
	\edef\SMTempSlideLangSlug{#1}%
	\expandafter\setSlideLanguageFromSlug@aux\SMTempSlideLangSlug-zzz-zzz\relax
}
\def\setSlideLanguageFromSlug@aux#1-#2-#3\relax{%
	\ifstrequal{#2}{EN}{\renewcommand{\slideLanguage}{EN}\global\SMSlideLanguageSettrue}{%
	\ifstrequal{#2}{DE}{\renewcommand{\slideLanguage}{DE}\global\SMSlideLanguageSettrue}{%
	\ifstrequal{#2}{mult}{\renewcommand{\slideLanguage}{mult}\global\SMSlideLanguageSettrue\global\SMSlideLanguageMulttrue}{}}}%
}
\makeatother

% Define a new \slide command to automatically pick the current version of the slide
% 2026-07-08 (direct instruction): slide filenames dropped their leading
% "s-" (e.g. slides/ab3-Some-Slide-current.tex, not the old
% slides/ab3-Some-Slide-current.tex) so slide IDs are directly
% searchable in Explorer -- this macro (and the other three below) is the
% *only* place that ever needed to change for that: every deck's own
% \slide{}/\slideversion{}/\omitslide{}/\omitslideversion{} calls pass just
% "ID-slug", never the filename itself, so renaming every file on disk to
% match needed no changes anywhere else (not the deck files, not
% slide_ids.csv, which only ever stored bare IDs to begin with).
\newcommand{\slide}[1]{%
	\global\SMSlideVersionSetfalse
	\setSlideID{#1}%
	\setSlideLanguageFromSlug{#1}%
	\input{slides/#1-current.tex}%
}

% \slidefinal{ID-slug} -- like \slide{}, but renders the slide at its final,
% fully-revealed overlay state in every deck version, instructor included
% (2026-07-12, direct instruction: show jww's finished build in place of a
% separate recap slide, "irrespective of which deck version is being
% compiled" -- that recap slide has since been deleted entirely, its own ID
% no longer in use). A slide
% built with \conditionalpause[instructor] already renders fully built in
% the during-course/after-course compiles -- that pause is a no-op for
% them, so nothing there ever waits on a reveal step -- the only case
% actually needing help is the instructor deck, which would otherwise still
% step through the build. Locally shadowing \conditionalpause to a no-op
% (same optional-argument signature as the real one, so it still accepts
% e.g. \conditionalpause[instructor] without erroring) for the duration of
% one \slide{} call makes every \conditionalpause in that one inclusion
% inert, regardless of the actual compiled \SMDeckVersion -- scoped
% inside \begingroup/\endgroup, so it reverts immediately after, and has no
% effect on any *other* \slide{}/\slidefinal{} call to the same ID
% elsewhere in the same deck (e.g. an earlier, normally-building instance of
% the same slide). \deckonly/\deckuncover/\deckvisible are untouched, so
% deck-version-tagged content (as opposed to reveal-step-tagged content)
% still follows its usual per-version rules.
\newcommand{\slidefinal}[1]{%
	\begingroup
		\renewcommand{\conditionalpause}[1][instructor]{}%
		\slide{#1}%
	\endgroup
}

% Load a specific, pinned version of a slide instead of whichever is "current"
% Usage: \slideversion{ID-slug}{M-m-p} loads slides/ID-slug-vM-m-p.tex
\newcommand{\slideversion}[2]{%
	\setSlideID{#1}%
	\setSlideLanguageFromSlug{#1}%
	\setSlideVersion{#2}%
	\input{slides/#1-v#2.tex}%
}

% Student-version decks: \studentversion (called anywhere in a deck, e.g.
% right after \input{SlideMage_template_uni-graz}, similar to \sidefootnotes or
% \setfootnotesymbols) switches that compile into the "handout" variant --
% including, as of 2026-07-03, actually forcing beamer's own "handout" mode
% (\gdef\beamer@currentmode{handout}, normally only reachable via the
% \documentclass[handout]{beamer} *class option*, set once at class-loading
% time -- too early for a command like \studentversion, called later from a
% deck's own preamble, to use that route). This was checked, not assumed:
% until this point \studentversion only ever touched \omitslide{}, so
% nothing suppressed \pause/overlay specs -- true only by coincidence, since
% no slide converted so far happens to use them. Late-redefining
% \beamer@currentmode was verified (isolated minimal-beamer tests, plain
% \pause *and* inline specs like \item<2->) to reproduce the same collapse
% -to-one-page, show-the-final-revealed-state behavior as the class option,
% even set well after \documentclass{beamer} -- beamer's overlay decoder
% (beamerbasedecode.sty) reads \beamer@currentmode at frame-processing time,
% not just once at class-option time, so this is safe, not a hack that
% happens to work by luck of file ordering.
% \omitslide{}/\omitslideversion{} are drop-in replacements for
% \slide{}/\slideversion{} that tag a slide as excluded from student deck versions:
%
%   \omitslide{xyz-Answer-key}                 % instead of \slide{...}
%   \omitslideversion{xyz-Answer-key}{1-0-0}    % instead of \slideversion{...}{...}
%
% In a \studentversion compile or a deck version other than instructor,
% a tagged slide's \input{} is skipped entirely and its frame number still
% advances, keeping numbering aligned across deck versions. In an
% instructor compile (the default), a tagged slide loads exactly as it
% would via \slide{}/\slideversion{}, with
% one difference: \SMOmittedSlide is set for the duration of that one
% \input{}, so the footer template below can print a marker next to that
% frame's slide ID -- a deliberately visible reminder, while reviewing the
% instructor deck, of which slides won't make it to students. The same tag
% works for all deck versions because it lives at the call site in the *deck*
% file (not inside the slide's own content file): the identical slide list
% can serve the instructor deck and both student deck versions, e.g.
%
%   \input{SlideMage_template_uni-graz}
%   \FooterText{...}
%   \makeatletter\ifdefined\HANDOUT\studentversion\fi\makeatother  % or similar
%   \input{course-content.tex}   % the shared \slide{}/\omitslide{} list
%
% with two thin wrapper decks (one plain, one calling \studentversion first)
% both \input-ing the same shared content file -- see README.md.
\newif\ifSMStudentVersion
% \makeatletter must wrap the \newcommand itself, not sit *inside* the macro
% body: catcodes apply at tokenization time, which for a macro body happens
% once, immediately, when \newcommand reads it -- not deferred until the
% macro is later called. \makeatletter placed inside the body (tried first)
% was too late: \beamer@currentmode had already been mis-tokenized as
% "\beamer" + the six-way-otherwise-ordinary character "@" + eleven plain
% letters "currentmode", by the time that inner \makeatletter could act on
% it, so the \gdef silently created/redefined the wrong thing and never
% touched beamer's actual mode variable. Confirmed via \meaning\beamer@
% currentmode before/after \studentversion -- unchanged with the inside-the-
% body version, fixed once \makeatletter moved out here.
\makeatletter
\newcommand{\studentversion}{%
	\global\SMStudentVersiontrue
	\gdef\beamer@currentmode{handout}%
}
\makeatother

% Deck version system - new boolean flags and string for current version
\newif\ifSMInstructorVersion
\newif\ifSMDuringCourseVersion
\newif\ifSMAfterCourseVersion
\gdef\SMDeckVersion{instructor}
\gdef\SMDeckVersionText{}

% \deckversion[] - set the deck version(s)
% Usage: \deckversion[instructor], \deckversion[during-course], etc.
% Multiple versions: \deckversion[instructor,during-course,after-course]
% Parses #1 as a comma-separated list via etoolbox's \forcsvlist (rather
% than the single \ifstrequal{whole-list}{"instructor"} check this used to
% do, which only ever matched a single-version list -- a multi-version
% \deckversion[instructor,during-course] silently set none of the three
% booleans, since the compared string was the literal, unsplit
% "instructor,during-course"). Note \forcsvlist, not etoolbox's similarly-
% named \forlistloop: \forlistloop iterates etoolbox's own internal
% pipe-delimited list format (built via \listadd), not a plain comma string
% like this command's own argument -- confirmed directly (isolated test),
% passing a raw comma string to it doesn't error or split on the commas, it
% just treats the whole string as a single one-item "list".
\makeatletter
\newcommand{\deckversion}[1][instructor]{%
	% \xdef (not \gdef): #1 needs to be fully expanded here, not just
	% stored as whatever token(s) were passed -- \deckversion[\myvar]
	% (compile-all-versions.bat's own override mechanism, via a wrapper
	% file's \def\deckversionoverride{...}) passes a macro, not a literal
	% string, and everywhere else in this deck-version system that reads
	% \SMDeckVersion only ever applies one \expandafter to it (see
	% \ifdeckversioninlist's own comment) -- enough to unwrap \SMDeckVersion
	% itself, but not a further macro still sitting inside it. Confirmed
	% directly: with the old \gdef, \deckversion[\deckversionoverride]
	% left \SMDeckVersionText empty and every \ifdeckversioninlist
	% check false, since \SMDeckVersion's own stored content was the
	% single unexpanded token \deckversionoverride, not "during-course".
	\xdef\SMDeckVersion{#1}%
	\global\SMInstructorVersionfalse
	\global\SMDuringCourseVersionfalse
	\global\SMAfterCourseVersionfalse
	\gdef\SMDeckVersionText{}%
	\def\do@deckversions##1{%
		\ifstrequal{##1}{instructor}{\global\SMInstructorVersiontrue\gdef\SMDeckVersionText{instructor deck}}{}%
		\ifstrequal{##1}{during-course}{\global\SMDuringCourseVersiontrue\gdef\SMDeckVersionText{student course deck}}{}%
		\ifstrequal{##1}{after-course}{\global\SMAfterCourseVersiontrue\gdef\SMDeckVersionText{student final deck}}{}%
	}%
	% \forcsvlist itself splits on literal commas in its own raw,
	% unexpanded argument (confirmed directly against etoolbox's own
	% source, \etb@defparser@do@aux) -- passing #1 straight through here
	% has the exact same problem as the \gdef above did, one call site
	% later. Building the whole \forcsvlist{...}{...} call inside an
	% \edef (protecting \forcsvlist/\do@deckversions themselves with
	% \noexpand so they're invoked *later*, once this constructed call is
	% run, not expanded away here) expands only \SMDeckVersion --
	% already fully resolved by the \xdef above -- into the list argument.
	\edef\SM@deckversion@call{\noexpand\forcsvlist{\noexpand\do@deckversions}{\SMDeckVersion}}%
	\SM@deckversion@call
}
\makeatother

% \deckoutputfolder[version]{path} - declares where SlideMage should
% place this deck's own compiled PDF, read directly out of this file's own
% source text by SlideMage/deck_scan.py (OUTPUT_FOLDER_RE) at deck-scan
% time -- never expanded by LaTeX itself. A silent no-op here (2026-09-16,
% direct instruction) purely so a deck carrying this directive still
% compiles cleanly under a plain lualatex/TeXstudio run, the same reason
% \deckversionoverride's own \ifdefined door (used elsewhere in this file)
% never needs a real definition either. #1 (the optional [version]) is
% intentionally unused: unlike \deckversion above, this macro's own
% argument is never inspected by LaTeX at all, only by SlideMage's own
% Python-side regex -- see README.md's "Custom output folder" section for
% the actual directive syntax.
\newcommand{\deckoutputfolder}[2][]{}

% \decktags{tag1,tag2} - declares this deck's tags (2026-09-19), read
% directly out of this file's own source text by SlideMage/deck_scan.py
% (DECKTAGS_RE) and never expanded by LaTeX itself -- a silent no-op here
% for exactly the same reason \deckoutputfolder above is. SlideMage's deck
% tables group decks by their tag: work (the pinned workdeck), course,
% research, test; a deck with none of those is a course deck. Written and
% edited from SlideMage itself (the "Tags" button, and the Duplicate
% dialog), placed just before \begin{document}.
\newcommand{\decktags}[1]{}

% Deck language system (added 2026-08-24) -- fixed per deck, like
% \FooterText or the institution \input{} choice, not a per-compile-run
% toggle like \deckversion above: a German course is its own deck file
% with mostly \setSlideLanguageFromSlug-tagged German/multilingual slide
% IDs, not the same deck file recompiled twice. \decklanguage[en]/
% \SMDeckLanguage/\ifSMGerman mirror \deckversion/\SMDeckVersion
% mechanically (same \xdef for the same \deckversion[\deckversionoverride]-
% style macro-argument reason documented above), so every deck gets the
% same \ifdefined\decklanguageoverride ... \fi override door
% \deckversionoverride already has, for SlideMage's own one-off render
% preview (slide_preview.py) and for ad-hoc hand-testing via
% decks/#workdeck.tex -- never wired into the Active-decks table itself.
\newif\ifSMGerman
\gdef\SMDeckLanguage{en}
\makeatletter
\newcommand{\decklanguage}[1][en]{%
	\xdef\SMDeckLanguage{#1}%
	\global\SMGermanfalse
	% \ifstrequal{<macro>}{...} is a trap, not just here but anywhere in this
	% file: its etoolbox implementation compares via
	% \edef\tempX{\detokenize{<arg>}} -- \detokenize does NOT expand its
	% argument first, so handed a bare control sequence it stringifies the
	% *macro's name* ("\SMDeckLanguage"), never its meaning ("de"),
	% and the comparison silently always takes the false branch (confirmed
	% directly: \ifSMGerman never went true this way, even under
	% \decklanguage[de]). \deckversion's own \do@deckversions above sidesteps
	% this by building the comparison inside an \edef first, so \ifstrequal
	% only ever sees an already-expanded literal string -- same fix here.
	\def\do@decklanguage##1{%
		\gdef\SMPDFLanguage{##1}% any other language: its own code (fr, es, ...)
		\ifstrequal{##1}{en}{\gdef\SMPDFLanguage{en-US}}{}%
		\ifstrequal{##1}{de}{\global\SMGermantrue\gdef\SMPDFLanguage{de-DE}}{}}%
	\edef\SM@decklanguage@call{\noexpand\do@decklanguage{\SMDeckLanguage}}%
	\SM@decklanguage@call
	% Document-level PDF /Lang, independent of (and not blocked by) the
	% structural PDF/UA tagging this project can't use under beamer -- see
	% \SMArtifactBegin's own comment above for why full tagging is out;
	% a screen reader still uses this Catalog-level entry to pick
	% pronunciation for the whole document.
	\hypersetup{pdflang=\SMPDFLanguage}%
}
\makeatother

% \multilingualtext{EN text}{DE text} -- for the handful of light,
% image-only/title-only slides (id-slug's own language segment "mult",
% see \setSlideLanguageFromSlug above) that hold both languages in one
% file rather than being split into two separate slide IDs, switching on
% whichever language the *deck* is compiling as. The two-argument form is
% EN then DE, as it always was.
%
% \multilingualtext*{EN={Hello}, DE={Hallo}, FR={Bonjour}} -- the keyed form
% (2026-10-05, direct instruction), for any number of languages: one
% key=value pair per language, the key being the deck's \decklanguage code
% (upper or lower case: EN = en). The value is picked whose key is the
% deck's language; a deck in a language the list does not have gets the EN
% value, or the first one if there is no EN. Commas and equals signs inside a
% value need the braces shown. A deck in another language than en/de just says
% \decklanguage[fr] (fr, es, ...): \ifSMGerman stays false for it, and the
% PDF's /Lang is the code itself.
\ExplSyntaxOn
\tl_new:N \l__sm_mlt_pick_tl
\tl_new:N \l__sm_mlt_en_tl
\tl_new:N \l__sm_mlt_first_tl
\bool_new:N \l__sm_mlt_pick_bool
\bool_new:N \l__sm_mlt_en_bool
\bool_new:N \l__sm_mlt_first_bool
\cs_new:Npn \__sm_mlt_pair:nn #1#2 { \ifSMGerman #2 \else #1 \fi }
\cs_new_protected:Npn \__sm_mlt_item:nn #1#2
  {
    \bool_if:NF \l__sm_mlt_first_bool
      { \tl_set:Nn \l__sm_mlt_first_tl {#2} \bool_set_true:N \l__sm_mlt_first_bool }
    \str_if_eq:eeT { \str_casefold:n {#1} } { en }
      { \tl_set:Nn \l__sm_mlt_en_tl {#2} \bool_set_true:N \l__sm_mlt_en_bool }
    \str_if_eq:eeT { \str_casefold:n {#1} } { \str_casefold:V \SMDeckLanguage }
      { \tl_set:Nn \l__sm_mlt_pick_tl {#2} \bool_set_true:N \l__sm_mlt_pick_bool }
  }
\cs_new_protected:Npn \__sm_mlt_item_bad:n #1
  { \msg_warning:nnn { slidemage } { multilingual-key-without-value } {#1} }
\msg_new:nnn { slidemage } { multilingual-key-without-value }
  { \string\multilingualtext*:~"#1"~has~no~=~value;~ignored. }
\cs_new_protected:Npn \__sm_mlt_keyed:n #1
  {
    \group_begin:
      \bool_set_false:N \l__sm_mlt_pick_bool
      \bool_set_false:N \l__sm_mlt_en_bool
      \bool_set_false:N \l__sm_mlt_first_bool
      \keyval_parse:NNn \__sm_mlt_item_bad:n \__sm_mlt_item:nn {#1}
      \bool_if:NTF \l__sm_mlt_pick_bool
        { \exp_last_unbraced:NV \group_end: \l__sm_mlt_pick_tl }
        {
          \bool_if:NTF \l__sm_mlt_en_bool
            { \exp_last_unbraced:NV \group_end: \l__sm_mlt_en_tl }
            { \exp_last_unbraced:NV \group_end: \l__sm_mlt_first_tl }
        }
  }
\NewDocumentCommand{\multilingualtext}{s m}
  { \IfBooleanTF {#1} { \__sm_mlt_keyed:n {#2} } { \__sm_mlt_pair:nn {#2} } }
\ExplSyntaxOff

\newif\ifSMOmittedSlide
\newcommand{\SMOmittedSlideMarker}{ | omitted from student deck}
% \stepcounter{framenumber} in the skipped branches below keeps the student
% deck's slide numbers matching the instructor deck's, gap and all, so "see
% slide 12" means the same physical slide to both audiences: beamer steps
% "framenumber" exactly once per \begin{frame} (checked in
% beamerbaseframe.sty -- once per *frame*, not per overlay page, so this is
% right regardless of how many overlay steps the omitted slide would have
% had), and this is the one place that increment is skipped along with the
% frame itself, so it has to be replayed by hand.
% Both forms are instructor-only; \studentversion also skips them.
\newcommand{\omitslide}[1]{%
	\ifSMStudentVersion
		\stepcounter{framenumber}%
	\else
		\ifdeckversioninlist{\SMDeckVersion}{instructor}%
			{\global\SMOmittedSlidetrue
			 \global\SMSlideVersionSetfalse
			 \setSlideID{#1}%
			 \setSlideLanguageFromSlug{#1}%
			 \input{slides/#1-current.tex}%
			 \global\SMOmittedSlidefalse}%
			{\stepcounter{framenumber}}%
	\fi
}
\newcommand{\omitslideversion}[2]{%
	\ifSMStudentVersion
		\stepcounter{framenumber}%
	\else
		\ifdeckversioninlist{\SMDeckVersion}{instructor}%
			{\global\SMOmittedSlidetrue
			 \setSlideID{#1}%
			 \setSlideLanguageFromSlug{#1}%
			 \setSlideVersion{#2}%
			 \input{slides/#1-v#2.tex}%
			 \global\SMOmittedSlidefalse}%
			{\stepcounter{framenumber}}%
	\fi
}

% Deck version-aware overlay commands
% Syntax: \deckonly<overlay>[versions]{content}
% If [versions] omitted, defaults to all decks (instructor,during-course,after-course)
%
% Built on etoolbox's own \forcsvlist for the comma-splitting (rather than a
% hand-rolled parser: an earlier version tried to split a list at each comma
% using \newcommand's plain, non-delimited argument grabbing
% (\newcommand{\ifdeckversioninlist@split}[3]{...} called as
% \ifdeckversioninlist@split#1\@nil{#1}) -- but a plain \newcommand doesn't do
% delimited-parameter matching against \@nil the way a \def with a matching
% pattern would; it just grabs the next 3 tokens/groups one at a time. So
% e.g. "instructor" silently split into single characters ("i", "n", "s",
% ...) instead of being read as one list item, which eventually ran the
% parser off the end of its own input and surfaced as "Undefined control
% sequence" wherever the first \conditionalpause[...]/\deckonly[...] in a
% deck happened to be used -- confirmed directly against the compile log,
% traced to slide y5z-CoCo-Bonds's own \conditionalpause[instructor].
%
% NOT etoolbox's own \ifinlist{item}{list}, despite the name similarity: that
% one tests membership in etoolbox's own internal list format, built only via
% \listadd (items joined with a literal "|", not ","); fed a plain comma
% string instead, confirmed directly (isolated test) to silently return
% false for every item, including an exact single-item match.
%
% The membership test itself needs one more subtlety: #1 here is the deck
% version to search for -- always \SMDeckVersion at every call site
% below, i.e. a macro, not a literal string -- so comparing it via a plain
% \ifstrequal{##1}{#1} does not work either: etoolbox's \ifstrequal compares
% via \detokenize, which (confirmed directly, isolated test) stringifies an
% unexpanded macro token to its own control-sequence *name* ("SMDeckVersion"),
% not the current value that macro holds. \expandafter forces that expansion
% before \ifstrequal ever sees it.
\makeatletter

\let\SM@original@only\only
\let\SM@original@uncover\uncover
\let\SM@original@visible\visible

\newif\ifSM@deckversion@found

% Usage: \ifdeckversioninlist{version}{list}{true-branch}{false-branch}
\newcommand{\ifdeckversioninlist}[2]{%
	\global\SM@deckversion@foundfalse
	\def\SM@deckversion@check##1{%
		\expandafter\ifstrequal\expandafter{#1}{##1}{\global\SM@deckversion@foundtrue}{}%
	}%
	\forcsvlist{\SM@deckversion@check}{#2}%
	\ifSM@deckversion@found
		\expandafter\@firstoftwo
	\else
		\expandafter\@secondoftwo
	\fi
}

% \@ifnextchar-based parsing (the same idiom LaTeX2e's own kernel commands
% use for an optional argument that isn't a plain trailing [...]) so the
% overlay spec reads exactly like beamer's own \only<...>, with the deck-
% version list as a further, genuinely optional [...] afterwards:
%   \deckonly<2>{content}                              visible in all versions
%   \deckonly<2>[instructor]{content}                  instructor only
%   \deckonly<2>[during-course,after-course]{content}  during + after
\newcommand{\deckonly}{\@ifnextchar<{\SM@deckonly@overlay}{\SM@deckonly@overlay<1->}}
\def\SM@deckonly@overlay<#1>{\@ifnextchar[{\SM@deckonly@versions<#1>}{\SM@deckonly@versions<#1>[instructor,during-course,after-course]}}
\def\SM@deckonly@versions<#1>[#2]#3{%
	\ifSMInstructorVersion
		\SM@original@only<#1>{#3}%
	\else
		\ifdeckversioninlist{\SMDeckVersion}{#2}{\SM@original@only<#1>{#3}}{}%
	\fi
}

% Same interface as \deckonly, wrapping \uncover and \visible respectively
\newcommand{\deckuncover}{\@ifnextchar<{\SM@deckuncover@overlay}{\SM@deckuncover@overlay<1->}}
\def\SM@deckuncover@overlay<#1>{\@ifnextchar[{\SM@deckuncover@versions<#1>}{\SM@deckuncover@versions<#1>[instructor,during-course,after-course]}}
\def\SM@deckuncover@versions<#1>[#2]#3{%
	\ifSMInstructorVersion
		\SM@original@uncover<#1>{#3}%
	\else
		\ifdeckversioninlist{\SMDeckVersion}{#2}{\SM@original@uncover<#1>{#3}}{}%
	\fi
}

\newcommand{\deckvisible}{\@ifnextchar<{\SM@deckvisible@overlay}{\SM@deckvisible@overlay<1->}}
\def\SM@deckvisible@overlay<#1>{\@ifnextchar[{\SM@deckvisible@versions<#1>}{\SM@deckvisible@versions<#1>[instructor,during-course,after-course]}}
\def\SM@deckvisible@versions<#1>[#2]#3{%
	\ifSMInstructorVersion
		\SM@original@visible<#1>{#3}%
	\else
		\ifdeckversioninlist{\SMDeckVersion}{#2}{\SM@original@visible<#1>{#3}}{}%
	\fi
}

% \conditionalpause[versions] - \pause, but only takes effect in the listed
% deck versions (default: instructor only); a no-op, not even a skipped
% overlay step, in any other version
% Usage: \conditionalpause[instructor,during-course] or \conditionalpause[after-course]
\newcommand{\conditionalpause}[1][instructor]{%
	\ifdeckversioninlist{\SMDeckVersion}{#1}{\pause}{}%
}

% \deckreveal[step]{content} - the "hidden answer, click-revealed for the
% instructor, absent for during-course, pre-shown for after-course" pattern
% already documented by hand on several exercise slides (e.g. mbz/x7q/37s/
% 3cv's own header comments: "a slide/reveal-step that is hidden in the
% source PowerPoint should, by default, only be made visible on the
% instructor version"). A bare, unconditional \only<N->{content} looks like
% it does this, but doesn't: \only's overlay range registers with beamer's
% own overlay-counting pass regardless of \ifSMInstructorVersion/
% \conditionalpause -- neither of those touch \only itself, they only guard
% *other* content around it -- so a frame with nothing else establishing a
% second overlay step still gets a real second slide out of a bare
% \only<2->, in *every* deck version, not just the instructor one. Confirmed
% directly (2026-07-21, direct instruction) on 37s/3cv: both a) leaked the
% "hidden" content to during-course as a genuine second slide, the exact
% thing this pattern exists to prevent, and b) showed a blank second bullet/
% cell first on that leaked slide's own step 1 -- \only<2-> hides the
% *content*, not the space it would occupy, so during-course's own step 1
% still had nothing else deciding where the bullet/cell landed.
% \deckreveal fixes both by picking a genuinely different rendering per
% version rather than reusing the same overlay-restricted \only call
% everywhere: instructor gets the real overlay build (\only<step->,
% default step 2, matching \tallitem's own default-multiplier convention);
% during-course gets nothing at all (no token emitted, so beamer's overlay
% pass never sees an overlay spec to register); after-course gets the
% content unconditionally, with no overlay spec of its own either, so it
% doesn't silently force a second slide there either. Content-level only
% (wraps a value/expression, like \only itself) -- doesn't cover an
% itemize \item's own bullet marker, which needs its own overlay spec
% (`\item<step->`) directly on the \item, not wrapped from outside; see
% 37s's own comment for that variant, hand-written rather than factored
% into a second macro since only one call site needs it so far.
% (2026-09-30: that variant is now \deckitem<step->, below.)
\NewDocumentCommand{\deckreveal}{s O{2} m}{%
	\ifSMDuringCourseVersion
	\else
		\ifSMInstructorVersion
			\IfBooleanTF{#1}{\visible<#2->{#3}}{\only<#2->{#3}}%
		\else
			#3%
		\fi
	\fi
}

% \deckrevealfinal[step]{content} - reveal step for the instructor deck only,
% everything shown at once in every other deck version (added 2026-09-19,
% direct instruction: "make the student versions show the final slide
% version immediately"). Unlike \deckreveal, the during-course deck gets the
% content too (not omitted), and like it no overlay spec is emitted outside
% the instructor version, so a non-instructor compile has no extra pages.
% For plain \pause steps use \conditionalpause[instructor] instead.
\NewDocumentCommand{\deckrevealfinal}{s O{2} m}{%
	\ifSMInstructorVersion
		\IfBooleanTF{#1}{\visible<#2->{#3}}{\only<#2->{#3}}%
	\else
		#3%
	\fi
}

% Starred forms, \deckreveal*[step]{content} and \deckrevealfinal*[step]{content}
% (2026-09-30, direct instruction: systematic replacements for per-slide
% fixes): the instructor steps before `step` keep the content's space
% (\visible instead of \only), so nothing around it moves when it appears --
% the "the diagram/equation/text jumps between reveal steps" class of
% complaint. Replaces the per-slide helpers that did the same by hand
% (\only<N->{x}\only<1..N-1>{\phantom{x}}; vxd's \vxdRevealB/C/D, c4c's
% \cRevealVisible). Other versions behave exactly as the unstarred forms.

% \deckitem<spec> - a list item whose overlay spec only applies in the
% instructor version; every other version gets a plain \item (the
% "student versions show the final state" rule). Replaces the pattern
%   \ifSMInstructorVersion \item<3-> text \else \item text \fi
% which repeated each item's text twice (42 items in 12 slides by
% 2026-09-30). An item label may follow: \deckitem<2->[\SMBulletEllipsis].
% Without a <spec> it is a plain \item.
\newcommand{\deckitem}{\@ifnextchar<{\SM@deckitem}{\item}}
\def\SM@deckitem<#1>{%
	\ifSMInstructorVersion
		\def\SM@next{\item<#1>}%
	\else
		\let\SM@next\item
	\fi
	\SM@next
}

% \conditionalslide[versions]{ID-slug} - \slide{}, but only included in the
% listed deck versions; keeps the frame counter in step with the omitted
% version otherwise (same \stepcounter{framenumber} technique \omitslide{}
% already uses for the student deck), so "see slide N" still points at the
% same physical slide across every deck version.
% Usage: \conditionalslide[instructor,after-course]{xyz-Answer-key}
\newcommand{\conditionalslide}[2][instructor,during-course,after-course]{%
	\ifdeckversioninlist{\SMDeckVersion}{#1}%
		{\global\SMSlideVersionSetfalse
		 \setSlideID{#2}%
		 \setSlideLanguageFromSlug{#2}%
		 \input{slides/#2-current.tex}}%
		{\stepcounter{framenumber}}%
}

% \slidestep{ID-slug}[versions]{steps} - \conditionalslide{} that also restricts
% which reveal (overlay) steps of the slide are typeset (2026-09-20, direct
% instruction). One slide file with several reveal steps can then appear at
% several places of a deck, each time showing only the steps that belong
% there, and a different set of steps per deck version:
%   \slidestep{egk-EN-Structure-of-a-Scientific-Paper-Overview}[instructor]{1,2}
%   \slidestep{egk-EN-Structure-of-a-Scientific-Paper-Overview}[during-course,after-course]{1}
% [versions] defaults to instructor,during-course,after-course (same as
% \conditionalslide; a presenter compile runs as version "instructor", so it
% always shows what the instructor deck shows); {steps} is a beamer overlay
% specification without the angle brackets ("1,2", "3", "2-", "1-2,4").
% The slide file itself is written once, with plain overlay specifications
% (\textcolor<2>{...}{...}, \only<3>{...}) and no version conditionals: which
% steps end up in the PDF is decided here, by handing the specification to
% beamer's own frame-level overlay restriction, \begin{frame}<steps>. Every
% frame environment opens its frame through \SMOpenFrame below, which
% inserts it.
%
% - Versions not in the list: the call does nothing at all -- not even a
%   \stepcounter{framenumber}, unlike \conditionalslide{}. Two \slidestep calls
%   at one position (say [instructor]{1,2} and [during-course,after-course]{1})
%   are alternatives, so every deck version has exactly one frame there;
%   stepping the counter for the omitted one would shift every later slide
%   number by one.
% - Jump target: the first \slidestep{} of a given id-slug in the document
%   defines the id's \hypertarget (see \setSlideID); later ones do not, so
%   there is no "multiply defined" warning and \hyperlink{id}{...} lands on the
%   first place the slide is shown.
% - \SMSlideSteps is cleared again right after the slide file has been
%   read, so it can never leak into a following plain \slide{}.
\gdef\SMSlideSteps{}
% \begin{frame}<steps>[options]{title}: the <...> has to be the very next
% token after "\begin{frame}" (beamer looks for it with \@ifnextchar), so the
% whole opening is built by \edef and only then executed. #1 is everything that
% follows "\begin{frame}" in the caller ("[#1]{#2}", "[plain,t]", ...).
% \SMFinalStateForStudents -- put it in a slide file, before the frame environment: in the
% student deck versions (everything but "instructor") that frame then typesets only its last
% overlay page, i.e. the final state of all its reveal steps (2026-09-21, direct instruction:
% "student decks should only show the final state"). Instructor and presenter compiles are
% unaffected. Unlike \conditionalpause/\deckrevealfinal, which each turn one reveal
% construct off, this needs no change to the slide body, so it also covers overlay specs
% like \only<2>, \item<3->, \uncover, \textcolor<2>, \alt and reveal steps that hide things
% again (the final state is what is left at the last step).
% How: beamer's frame-level overlay restriction \begin{frame}<N> (see \slidestep) with N =
% the frame's overlay count, which a frame only knows after it was typeset once. Every
% student compile therefore records N per frame (keyed by beamer's frame number, which is
% stable) in \jobname.ovl and reads it back at the start of the next run; the first run of
% a fresh job still typesets all pages, latexmk's rerun then restricts them. Beamer reports
% the same N for a restricted frame, so the value is stable.
\newif\ifSMFinalOnly
\newcommand{\SMFinalStateForStudents}{\global\SMFinalOnlytrue}
\makeatletter
\newwrite\SMOverlayFile
\newcommand{\SMSetOverlays}[2]{\global\@namedef{SMOverlays@#1}{#2}}
% Build files written before v5.3.0 (2026-10-01: this framework's \gotham...
% names became \SM...) still use the old names: .ovl (\gotham...SetOverlays),
% .top (...SetTopCorrection). The old names are assembled from \SMOldPrefix so a
% rerun of slidemage.sm_rename leaves them alone. The next compile rewrites
% the files with the new names.
\def\SMOldPrefix{gotham}
\expandafter\def\csname\SMOldPrefix SetOverlays\endcsname{\SMSetOverlays}
\expandafter\def\csname\SMOldPrefix SetTopCorrection\endcsname{\SMSetTopCorrection}
% A .aux written before v9.1.0 holds each subsection title (\SMAgendaSubsection)
% as `\gotham@subsectiontitleformat {Title}`: the Gotham theme wrapped titles
% in that no-op switch, which the SlideMage theme no longer has.
\expandafter\let\csname\SMOldPrefix @subsectiontitleformat\endcsname\@empty
\newcommand{\SMRecordOverlays}{%
	\ifSMInstructorVersion\else
		\immediate\write\SMOverlayFile{\string\SMSetOverlays{\the\c@framenumber}{\the\beamer@slideinframe}}%
	\fi
}
\apptocmd{\beamer@writeslidentry}{\SMRecordOverlays}{}{}
% "No file" is the line LaTeX's own \@input prints for a missing .aux or .nav: latexmk reads it and then knows the
% file as an input from the first run, so a second run that reads an unchanged one asks for no third.
\AtBeginDocument{%
	\InputIfFileExists{\jobname.ovl}{}{\typeout{No file \jobname.ovl.}}%
	\immediate\openout\SMOverlayFile=\jobname.ovl\relax
}
% Expands to <N> when the frame that is about to open is flagged and N is known.
\newcommand{\SMFinalSpec}{%
	\ifSMFinalOnly
		\ifSMInstructorVersion\else
			\ifcsname SMOverlays@\the\numexpr\c@framenumber+1\relax\endcsname
				<\csname SMOverlays@\the\numexpr\c@framenumber+1\relax\endcsname>%
			\fi
		\fi
	\fi
}
\newcommand{\SMOpenFrame}[1]{%
	\edef\SMOpenFrame@cmd{%
		\noexpand\begin{frame}%
		\ifx\SMSlideSteps\empty\SMFinalSpec\else<\SMSlideSteps>\fi
		\unexpanded{#1}}%
	\global\SMFinalOnlyfalse
	\SMOpenFrame@cmd
}
\NewDocumentCommand{\slidestep}{m O{instructor,during-course,after-course} m}{%
	\ifdeckversioninlist{\SMDeckVersion}{#2}%
		{\ifcsname SMSlideStepSeen@#1\endcsname
			\global\SMSuppressSlideTargettrue
		 \else
			\global\SMSuppressSlideTargetfalse
			\expandafter\gdef\csname SMSlideStepSeen@#1\endcsname{}%
		 \fi
		 \gdef\SMSlideSteps{#3}%
		 \global\SMSlideVersionSetfalse
		 \setSlideID{#1}%
		 \global\SMSuppressSlideTargetfalse
		 \setSlideLanguageFromSlug{#1}%
		 \input{slides/#1-current.tex}%
		 \gdef\SMSlideSteps{}}%
		{}%
}

\makeatother

% \SMUnderbrace{math}{label} - seamless TikZ-drawn brace, replacing plain
% amsmath \underbrace (2026-07-23, direct instruction, after presenting
% options on slide ebq -- see SESSION_HANDOFF.md's own entry for that task):
% amsmath assembles \underbrace's glyph from separate piecewise brace-half
% characters, which visibly shows a seam/notch at narrow widths (ebq's own
% "risk-free rate" term is narrow enough to show it plainly); this instead
% draws one continuous vector brace path via TikZ's `decorate` +
% `decoration={brace}` (decorations.pathreplacing, loaded above), which has
% no such seam at any width. #1 is typeset in math mode automatically (do
% not wrap the call itself in $...$); #2 is the label, set in upright text
% (font=\small, align=center so a multi-line label via \\ stays centered
% under the brace, e.g. two-line labels converted from \substack{}). Each
% call gets its own auto-numbered internal node name (SMUnderbraceCounter)
% so multiple calls can coexist on one slide without a caller-supplied name
% argument colliding (unlike the original \tikzbrace{}{}{} demo macro this
% replaces, which took the node name as a third argument).
\newcounter{SMUnderbraceCounter}
\newcommand{\SMUnderbrace}[2]{%
	\stepcounter{SMUnderbraceCounter}%
	\tikz[baseline=(SMUnderbraceNode\theSMUnderbraceCounter.base)]{
		\node[inner sep=0pt] (SMUnderbraceNode\theSMUnderbraceCounter) {$#1$};
		\draw[decorate,decoration={brace,amplitude=4pt,mirror,raise=1pt}]
			(SMUnderbraceNode\theSMUnderbraceCounter.south west) -- (SMUnderbraceNode\theSMUnderbraceCounter.south east)
			node[midway,below=6pt,font=\small,align=center] {#2};
	}
}

% PDF accessibility: mark purely decorative/running content (footer, below;
% also reused on titleframe's/quoteframe's own decorative background boxes,
% and -- since 2026-08-14, via the \includegraphics wrapper further down --
% on every image the alt-text manifest classifies as decorative, the
% institute logo among them) as a PDF "Artifact", so screen readers skip it
% -- via
% LuaTeX's raw \pdfextension literal
% primitive. \ifluatex-guarded (iftex is already loaded, via beamer.cls
% itself) because \pdfextension is LuaTeX-only -- this project's README
% documents xelatex as a supported alternative engine, and compiling with it
% hit exactly this: "Undefined control sequence: \pdfextension", hard-failing
% every frame, since the first version of this didn't guard the primitive at
% all. Under any non-LuaTeX engine these two commands now silently do
% nothing instead: decorative content just isn't marked as an Artifact under
% xelatex/pdflatex, rather than breaking the compile -- a real loss of
% accessibility under those engines specifically, not a full equivalent fix,
% but not something to solve today (see also: compiling with plain pdflatex
% already silently mis-renders the whole deck's fonts, so xelatex is really
% the only *other* engine this needs to not break under; pdflatex was never
% viable here to begin with, per "Compiling" in README.md).
% Deliberately NOT tagpdf/kernel tagging (\DocumentMetadata):
% beamer 3.77 refuses to compile at all when that's active ("The beamer
% class is not compatible with \DocumentMetadata ... switch to the new
% ltx-talk class if you require tagged output"), and ltx-talk is a from-
% scratch class (not beamer-based, no theme/template system) that this
% framework can't realistically migrate to. Bypassing beamer's check by
% force-loading the underlying package before \documentclass also failed --
% it gets past hyperref but then hits a genuine internal error in tagpdf's
% LuaTeX backend from the non-standard load order, not something to build
% real coursework on.
%
% /Artifact is a standalone PDF marked-content tag (PDF spec 14.8.2.2):
% valid in *any* PDF, with or without a structure tree, no /Properties
% resource entry or MCID required -- unlike structure-tree tagging, this
% doesn't require kernel tagging support at all. It's a real, working
% partial fix, not full PDF/UA conformance: there's no structure tree, so
% no true semantic headings, and reading order elsewhere on the slide still
% depends on plain content-stream order (verified, not enforced, for
% standardframe -- see its own comment).
% #1 is the artifact's inline PDF dictionary content, e.g.
% "/Type /Pagination /Subtype /Footer" for a running footer, or
% "/Type /Background" for a decorative fill -- not hardcoded to one /Type,
% since not everything decorative here is "Pagination" in the PDF spec's own
% sense of that word (running headers/footers/page numbers/watermarks);
% titleframe's/quoteframe's plain background color fills are semantically
% "Background" artifacts instead, a distinct standard type.
\newcommand{\SMArtifactBegin}[1]{%
	\ifluatex\pdfextension literal direct{/Artifact <<#1>> BDC}\fi
}
\newcommand{\SMArtifactEnd}{\ifluatex\pdfextension literal direct{EMC}\fi}

% Per-frame "Image source(s):" auto-footnote (ALT_TEXT_NOTES.md phase 6 step
% 3). The \includegraphics wrapper further down calls
% slidemage_alt.recordSource() for every image (decorative or informative --
% source is independent of role, same reasoning as media/alt-text.json's own
% ImageEntry.source), accumulating each frame's distinct, order-of-first-
% -appearance source strings in Lua. These two commands are standardframe's
% own hooks into that: reset at frame open (same per-overlay-replay reset
% \SMClearSideFootnotesBox already does), emit at frame close.
%
% Defined here (outside/before the "Alt-text manifest" \ifluatex block
% further down that actually loads SlideMage_alt.lua and redefines
% \includegraphics) rather than inside it, matching \SMArtifactBegin/End
% just above: standardframe unconditionally calls both, so they need to
% exist regardless of engine, each independently \ifluatex-guarded inside
% -- under any non-LuaTeX engine \SMResetImageSources/
% \SMEmitImageSourceFootnote are both no-ops, same degradation as the
% rest of this accessibility system under a non-LuaTeX engine. Safe to
% define before SlideMage_alt.lua is actually loaded (\AtBeginDocument,
% further down): these are only ever *called* from inside a real frame,
% well after \begin{document} has fired.
\newif\ifSMImageSourcesHasText
\newcommand{\SMImageSourcesText}{}
\newcommand{\SMResetImageSources}{%
	\ifluatex\directlua{slidemage_alt.resetFrameSources()}\fi
}
% \source{} is itself implemented as an empty-marker \footnote{} (see its
% own definition further below) -- calling it here gives the auto-footnote
% the exact same rendering (font, spacing, right-alignment, stretch-to-
% -sit-above-the-footer, side-footnote redirection under \sidefootnotes) as
% a hand-authored \source{} line, for free, rather than reimplementing any
% of that. Coexists with a real hand-authored \source{} call elsewhere in
% the same frame -- both are just separate unmarked footnote-styled lines,
% stacking the same way multiple real \footnote{}s already do.
% \SMImageSourceLabel -- the "Image source(s):"/"Bildquelle(n):" label,
% driven by \ifSMGerman (set from the compiling deck's own
% \decklanguage) rather than hardcoded, so a German deck's auto-footnote
% reads naturally instead of showing an English label on German content.
\newcommand{\SMImageSourceLabel}{\ifSMGerman Bildquelle(n)\else Image source(s)\fi}
\newcommand{\SMEmitImageSourceFootnote}{%
	\ifluatex
	\SMImageSourcesHasTextfalse
	\directlua{slidemage_alt.emitSourceFootnote()}%
	\ifSMImageSourcesHasText
		\source{\SMImageSourceLabel: \SMImageSourcesText}%
	\fi
	\fi
}

% Invisible slide-id stamp for footer-less frame types. \titleframe and
% \fullbleedframe (both defined further down) render via beamer's "plain"
% frame option, which suppresses the footline entirely -- so the visible
% footer's own "(slide id: ...)" text (see the footer template below) never
% reaches those pages. SlideMage's own compile-time preview snapshot
% (slide_snapshot.py) identifies which slide a compiled page belongs to
% purely by scanning its extracted text for that "slide id: XXX" string, so
% a plain frame with none of that text never gets a snapshot captured at
% all (caught 2026-08-15: titleframe's own title slide and 25 straight
% fullbleedframe "Excursus" slides in the IA2 deck all had no preview).
% This stamps the same marker, invisibly (PDF text-render mode 3, the same
% technique \SMAltInvisible below uses for image alt text) and wrapped
% as a /Pagination/Footer artifact, same as the real footer's own slide-id
% text -- so it reaches neither a sighted reader nor a screen reader,
% purely a technical marker for the compile pipeline. Only \slideID itself
% is stamped, not the omitted-slide/version suffix the visible
% footer also carries: nothing here is ever read by a human, and
% slide_snapshot.py's own regex only looks for the id.
%
% \vbox to 0pt{\vss ...}: same reasoning as every other absolute/zero-space
% overlay in this file (see \titleframe's own comment on this) -- \sbox
% builds an hbox, and a bare hbox placed directly in a frame's vertical
% list silently opens an implicit paragraph, switching following material
% into horizontal mode. Wrapping in a vbox contains that entirely within
% its own internal vertical mode instead of leaking out to the frame body.
\newsavebox{\SMSlideIDStampBox}
\newcommand{\SMInvisibleSlideIDStamp}{%
	\ifluatex
	\vbox to 0pt{\vss
	\begingroup
	\hbadness=10000 \hfuzz=\maxdimen
	\sbox{\SMSlideIDStampBox}{%
		\parbox[t]{\textwidth}{%
			\tiny\raggedright
			\SMArtifactBegin{/Type /Pagination /Subtype /Footer}%
			\pdfextension literal direct{3 Tr}%
			% Trailing \space before the closing paren, not just "\slideID)":
			% slide_snapshot.py's own regex ("slide id:\s*(\S+)") reads
			% everything up to the next whitespace as the id, and without a
			% space there it would capture "gcs)" instead of "gcs" (caught
			% by testing the actual compiled+extracted text, not assumed --
			% the real footer's own "... \slideID ... | ..." format always
			% has a space after \slideID for the same reason).
			(slide id: \slideID\space)%
			\pdfextension literal direct{0 Tr}%
			\SMArtifactEnd
		}%
	}%
	\wd\SMSlideIDStampBox=0pt
	\ht\SMSlideIDStampBox=0pt
	\dp\SMSlideIDStampBox=0pt
	\usebox{\SMSlideIDStampBox}%
	\endgroup
	}%
	\fi
}

% ---------------------------------------------------------------------------
% Alt-text manifest: automatic /Artifact marking of decorative images
% ---------------------------------------------------------------------------
% See ALT_TEXT_NOTES.md for the whole design; this is phase 1 of it. Every
% image in the project has a row in media/alt-text.csv saying whether it is
% "decorative" (carries no information beyond the slide text -- Midjourney
% illustrations, icons, chapter logos, institute logos) or "informative" (a
% chart/screenshot/photo a reader actually needs described). Decorative ones
% get wrapped in /Artifact here, so a screen reader skips them instead of
% announcing a bare filename. Informative ones are left alone for now; their
% descriptions are authored per slide and emitted later (phases 4 and 5).
%
% Why wrap \includegraphics rather than change the 265 call sites: the call
% sites don't need to change at all this way, including the ones inside tikz
% \node{}s and the ones that pass the filename through a macro argument
% (\SMInstituteLogoPath, or a #1 inside another macro's definition) --
% those only know their filename at expansion time, which a source-level
% rewrite couldn't have handled anyway. An image with no manifest row simply
% renders exactly as before, so a missing or unreadable CSV degrades to
% today's behavior rather than breaking every image in the deck.
%
% LuaTeX-only, same \ifluatex reasoning as \SMArtifactBegin above -- and
% here the guard covers the whole redefinition, not just the marking, since
% the lookup itself is Lua. Under any other engine \includegraphics is left
% completely untouched.
\ifluatex
	% The Lua itself lives in SlideMage_alt.lua rather than inline here --
	% see that file's own header for why (short version: TeX tokenizes a
	% \directlua argument first, so every "%" in a Lua string pattern
	% silently comments out the rest of its line, and the first version of
	% this failed exactly that way). This loader snippet is deliberately
	% free of "%" and "#" so it is safe to keep inline. Found via kpathsea
	% (TEXINPUTS), since it lives in latex/, outside the project.
	\SMLuaOnce{
		local path = kpse.find_file("SlideMage_alt.lua", "lua")
			or kpse.find_file("SlideMage_alt.lua", "tex")
		if path then
			dofile(path)
		end
	}

	\newif\ifSMAltDecorative
	\newif\ifSMAltHasText
	\newsavebox{\SMAltTextBox}
	\newcommand{\SMAltShort}{}

	% Invisible alt text (ALT_TEXT_NOTES.md phase 5). An untagged PDF is read
	% by *text extraction*, so text placed in the content stream at the image's
	% position is what a screen reader encounters there -- text rendering mode
	% 3 ("neither fill nor stroke", PDF spec 9.3.6) makes it invisible on the
	% page while leaving it fully extractable.
	%
	% Two deliberate properties:
	%
	% - The box's three dimensions are zeroed after typesetting, so the text
	%   contributes nothing whatsoever to layout. Setting \wd/\ht/\dp of an
	%   already-built box is the only way to be sure of that: \rlap/\smash
	%   each zero only some of them, and a stray descender would shift a
	%   hand-positioned slide.
	% - \SMAltImagePrefix ("Image: "/"Bild: ") prefixes the description so a listener can tell it apart
	%   from the slide's own prose, which they otherwise cannot.
	%
	% Known and accepted side effect: this text is real text, so it appears in
	% copy-paste and in Ctrl-F. That was the explicit tradeoff for reaching a
	% screen reader at all without a structure tree, which beamer cannot give
	% us (see README.md's accessibility section).
	% \SMAltImagePrefix -- language-conditional counterpart to
	% \SMImageSourceLabel above, for this invisible layer's own "Image: "/
	% "Bild: " prefix.
	\newcommand{\SMAltImagePrefix}{\ifSMGerman Bild: \else Image: \fi}
	\newcommand{\SMAltInvisible}[1]{%
		\begingroup
		% A long alt string in an unbreakable hbox would otherwise flood the
		% log with overfull-box warnings on every image.
		\hbadness=10000 \hfuzz=\maxdimen
		% Wrapped in a \parbox, and set \tiny, so the text stays *inside the
		% page*. Neither is cosmetic: an unbreakable line starting at the
		% image's own x-position runs past the right page edge after about 80
		% characters, and every PDF text extractor tested (pdftotext with and
		% without -layout, PyMuPDF) silently drops whatever falls outside the
		% page rectangle -- so a longer alt was being truncated mid-sentence
		% with no error anywhere. \tiny is free here (the text is invisible)
		% and fits several times more characters per line, which keeps even a
		% long description to one or two lines.
		\sbox{\SMAltTextBox}{%
			\parbox[t]{\textwidth}{%
				\tiny\raggedright
				\pdfextension literal direct{3 Tr}%
				\SMAltImagePrefix#1%
				\pdfextension literal direct{0 Tr}%
			}%
		}%
		\wd\SMAltTextBox=0pt
		\ht\SMAltTextBox=0pt
		\dp\SMAltTextBox=0pt
		\usebox{\SMAltTextBox}%
		\endgroup
	}

	% Deferred to \AtBeginDocument so it cannot possibly run before beamer has
	% loaded graphicx, and so a deck that loads further graphics packages of
	% its own still ends up with this wrapper outermost.
	\AtBeginDocument{%
		\directlua{
			local loaded = slidemage_alt.load()
			if loaded == nil then
				texio.write_nl("slidemage_alt: media/alt-text.csv not found -- "
					.. "decorative images will not be marked as artifacts.")
			end
		}%
		% \NewCommandCopy (LaTeX 2020-10-01 and later) rather than \let: it
		% copies a robust command's full definition correctly, which \let
		% does not for commands whose top level is a \protected wrapper.
		\NewCommandCopy{\SMOrigIncludegraphics}{\includegraphics}%
		% Signature "s o o m": the starred form and the legacy two-optional
		% bounding-box form aren't used anywhere in this project today
		% (checked, 2026-08-14: zero occurrences of either), but graphicx
		% accepts them, so the wrapper accepts and forwards them rather than
		% silently breaking the day one gets used. "o" (not "O{}") because
		% only -NoValue- distinguishes an absent optional argument from an
		% empty one, and graphicx reads [a][b] and [a] as different things.
		\RenewDocumentCommand{\includegraphics}{s o o m}{%
			\begingroup
			% \edef so a filename reaching this through a macro
			% (\SMInstituteLogoPath, or a #1 in another macro's body) is
			% resolved to the actual path before the lookup. Every path-valued
			% macro in this project is a plain expandable string; a fragile one
			% would need \protected@edef here instead.
			\edef\SMAltReference{#4}%
			\SMAltDecorativefalse
			\directlua{slidemage_alt.mark("\luaescapestring{\SMAltReference}")}%
			% Unconditional -- unlike SMAltInvisible below (informative
			% images only), a source is recorded regardless of role.
			\directlua{slidemage_alt.recordSource("\luaescapestring{\SMAltReference}", "\luaescapestring{\SMDeckLanguage}")}%
			% /Type /Layout is the PDF spec's artifact type for cosmetic design
			% elements, which is what a decorative illustration is. Not
			% /Background: that one means a fill or image behind the whole
			% page, and is already used for exactly that by titleframe and
			% quoteframe above.
			\ifSMAltDecorative\SMArtifactBegin{/Type /Layout}\fi
			% Emitted *before* the image, so the description precedes it in the
			% content stream and therefore in extraction order -- a listener
			% hears what the image is before reaching it, rather than after.
			% slidemage_alt.altText prints nothing for a decorative image or one
			% with no short alt yet, so this line is inert until the manifest
			% is filled in.
			\SMAltHasTextfalse
			\ifSMAltDecorative\else
				\directlua{slidemage_alt.altText("\luaescapestring{\SMAltReference}", "\luaescapestring{\SMDeckLanguage}")}%
				\ifSMAltHasText\SMAltInvisible{\SMAltShort}\fi
			\fi
			\IfBooleanTF{#1}%
				{\IfNoValueTF{#2}%
					{\SMOrigIncludegraphics*{#4}}%
					{\IfNoValueTF{#3}%
						{\SMOrigIncludegraphics*[#2]{#4}}%
						{\SMOrigIncludegraphics*[#2][#3]{#4}}}}%
				{\IfNoValueTF{#2}%
					{\SMOrigIncludegraphics{#4}}%
					{\IfNoValueTF{#3}%
						{\SMOrigIncludegraphics[#2]{#4}}%
						{\SMOrigIncludegraphics[#2][#3]{#4}}}}%
			\ifSMAltDecorative\SMArtifactEnd\fi
			\endgroup
		}%
	}
\fi

% \tikzillustration{ID}: extends the alt-text system above -- which only
% ever sees a real \includegraphics file -- to a hand-drawn tikz diagram
% with no file behind it at all (ALT_TEXT_NOTES.md). #2 is the
% illustration's own opaque, SlideMage-assigned ID (never hand-invented
% -- same "permanent ID" convention slides themselves use, so a description
% survives edits to the drawing itself); #1 (optional) is passed straight
% through as the inner tikzpicture's own options. The body between
% \begin/\end is the tikzpicture's content, completely unmodified.
%
% Role (decorative/informative) is looked up exactly the way an image's is
% -- slidemage_alt.mark(), the same Lua call \includegraphics's own wrapper
% above makes, against the same media/alt-text.csv -- rather than a second
% bracket argument SlideMage would otherwise have to text-edit into this
% call in place every time the role changes. SlideMage writes this
% illustration's own row into media/alt-text.json (and its compile-time
% media/alt-text.csv projection) keyed "tikz:ID", never a real media/ path;
% slidemage_alt.mark() itself needed no changes, since it already keys purely
% on whatever string it is given (SlideMage/alt_images.py).
%
% A decorative illustration is wrapped in the same \SMArtifactBegin{/Type
% /Background}/\End pair sideimageframe's/fullbleedframe's own decorative
% filler content already uses, so a screen reader skips it entirely; an
% informative one is left completely unmarked, same as real content in
% figureframe/fullbleedframe. Unlike an informative *image*, it gets no
% invisible-text alt layer (phase 5) of its own -- that layer's positioning
% is built around \includegraphics's own bounding box, which a freeform
% tikzpicture doesn't have the same way. The per-deck accessible transcript
% (README.md) is this illustration's own route to a screen reader instead,
% the same status phase 5 already carries for every image.
%
% \ifluatex-guarded exactly like \SMArtifactBegin/\SMArtifactEnd
% above (both \SMAltDecorativefalse/\ifSMAltDecorative only exist
% inside that same guard) -- under any other engine this degrades to a
% plain, unmarked tikzpicture, the same "no accessibility marking, but
% still compiles and renders" degradation the rest of this system already
% has. Safe even though those two control sequences are otherwise undefined
% under a non-LuaTeX engine: TeX's own \iffalse-branch skipping never looks
% up a token's meaning, so an unrecognized name inside the untaken
% \ifluatex branch is simply skipped over, never executed.
\newenvironment{tikzillustration}[2][]{%
	\ifluatex
		\SMAltDecorativefalse
		\directlua{slidemage_alt.mark("\luaescapestring{tikz:#2}")}%
		\ifSMAltDecorative\SMArtifactBegin{/Type /Background}\fi
	\fi
	\begin{tikzpicture}[#1]
}{%
	\end{tikzpicture}%
	\ifluatex\ifSMAltDecorative\SMArtifactEnd\fi\fi
}

% Footer
% The footer box on the footline, the footnote template with its symbols, \sidefootnotes and the
% strip they are drawn in are the theme's (beamerthemeSlideMage.sty); the footer text, which
% depends on the deck version, is the "frame footer" template below.
% Smaller footer + a bit of padding under it
\setbeamerfont{frame footer}{size=\SMFooterTextSize} % the catalog: {font size}{baseline skip} in pt

% Shrinks the gap between footnotes and the footer box below them. The
% theme positions the footer's content via \vspace*{\SMFootlineVOffset}
% (a fixed, negative length) before it starts; making that more negative
% turns out to give footnotes (which are placed just before the footer, and
% stretch to fill whatever room remains via beamer's own
% \beamer@framebottomskip) more room to extend into, pulling them down
% towards the footer -- checked empirically against the PDF's own
% coordinates, this does *not* move the footer's own text at all (that
% still sits exactly where \SMFooterBoxHeight below was calibrated for),
% so the footer box itself needs no compensating adjustment.
%
% Partly given back (+1.4pt) 2026-07-05: a footnote/source containing an
% underlined \href (see the hyperlink-styling comment above) sits close
% enough to the footer at the original -1.3pt that the underline visually
% touches the footer box (checked directly against the compiled PDF's own
% pixel coordinates: a plain-text footnote's own gap to the footer measured
% ~1.9pt, but an underlined-link footnote's gap measured only ~0.6pt at
% the same setting, since the strut rebuilt for footnotes just below --
% correctly sized for plain text's own natural depth -- doesn't account
% for the underline rule extending further down). +1.4pt splits the
% difference: a linked footnote's gap now measures close to the plain-
% footnote reference (~1.9pt) instead of colliding with the footer, and a
% plain footnote's own gap just becomes a little more relaxed (~3.3pt)
% rather than stretched uncomfortably tight -- applied globally rather
% than only to footnotes with links, since there's no clean per-footnote
% hook for it and the plain-footnote side effect is minor.
%
% Nudged back down by 0.4pt the next day (direct instruction: move
% footnotes down "by the width of the underline line") -- the +1.4pt above
% left a bit more headroom than necessary; 0.4pt is the underline rule's
% own measured stroke thickness (a solid ~1px core at 600dpi, i.e.
% 1/600in \approx 0.12pt, plus antialiasing bleed either side consistent
% with the common 0.4pt LaTeX default rule width -- measured directly
% against the compiled PDF's own pixel values, not assumed).
%
% Reverted the very next round (direct instruction: move back *up* by that
% same underline thickness) -- net effect is back to the plain +1.4pt this
% section started with; kept as two separate historical adjustments below
% rather than collapsing them into one, so this back-and-forth stays
% visible rather than silently disappearing.
% (v9.4.0: the adjustments are one value now, the default of the catalog's
% \SMFootlineVOffset in SlideMage_setup.tex: -0.5ex-1.3pt+1.4pt-0.4pt+0.4pt-0.87pt
% = -0.5ex-0.77pt.)

% 2026-07-06 (direct instruction): move footnotes down by half the
% measured distance between a linked footnote's own underline and the
% footer box, checked on slides/3ic-Corporate-Venture-Capital-
% current.tex specifically (the underline sits at y=245.95pt, the footer
% box's own top edge at y=247.69pt, a 1.74pt gap) -- half of that is
% 0.87pt, applied as a further negative adjustment (more negative =
% footnotes get pulled further down, per this section's own comment
% above).

% \SMFooterTextDrop (v9.8.0: derived; it was a number of the layout profile since v9.4.0, -0.49pt for
% Fira Sans, -0.37pt for Latin Modern Sans) centres the ink of the frame number in the footer box. The
% template below starts the footer line with \vspace*{drop}, which in the footline box's horizontal
% mode is a \vadjust below the line, as is the \vspace*{0.4ex} that ends it; the box has the declared
% depth \SMFootlineDepth, so the number's baseline lies
%     \SMFootlineDepth + (the depth of the line) + 0.4ex + drop
% above the page's bottom edge. The depth of the line is that of the slide id's strut
% (\SMFooterIDStrut makes it the same on every slide). The ink of the digits is what luaotfload
% keeps as the height and depth of a digit's box (the glyph's bounding box: 676 and 19 units in Fira
% Sans Light), so centring it in the footer box of height \SMFooterBoxHeight gives the drop.
\makeatletter
\newlength{\SM@footerlinedepth}
\newcommand{\SM@footertextdrop}{%
	\setbox\z@\hbox{\SMFooterIDSize\SMFooterIDStrut}\setlength{\SM@footerlinedepth}{\dp\z@}%
	\setbox\z@\hbox{0123456789}%
	\setlength{\SMFooterTextDrop}{\dimexpr(\SMFooterBoxHeight-\ht\z@+\dp\z@)/2-\SMFootlineDepth-\SM@footerlinedepth-0.4ex\relax}}
\makeatother
%
% \SMFooterRightShift (v9.7.0: derived; it was a number of the layout profile, 4.6pt for Fira Sans,
% -0.33pt for Latin Modern Sans) puts the frame number's right edge on the logo's right edge above
% it, as the footnotes' lines end there. The footer line is typeset in the footline box, which is
% \hsize wide and ends its lines \rightskip before that; after the number come, in this order,
% \hspace*{-4em} and the shift (set here), then the box's \SMRightFooterPadding and its
% \SMFootlineHRightOffset. So the number's right edge is
%     \hsize - \rightskip - \SMRightFooterPadding - \SMFootlineHRightOffset + 4em + shift
% and the shift that puts it at \paperwidth - \SMFrametitleRightPadding follows. Every term is the
% one in force where the footer is typeset (the \rightskip is 3ex of the footline font, the 4em is
% the footer's), which is why the shift is computed in the frame footer template and not once: the
% two fonts differ in size, and a profile with another x-height has other paddings in ex.

% The footer no longer shows a compile date/time (removed in v8.25.0: a
% changing stamp made every wrapper and PDF differ between compiles). The
% two names below remain as harmless no-ops so decks that still call
% \SMShowCompileTimestampfalse/true or print \SMCompileTimestamp compile.
\newif\ifSMShowCompileTimestamp
\newcommand{\SMCompileTimestamp}{}

% \SMPresenterModeOverride: SlideMage's own hook for a "Presenter"
% compile -- a separate output alongside the plain instructor PDF, never a
% toggle on it (see SlideMage/compiler.py's own CompileJob.presenter
% field and SlideMage/qt_deck_table.py's own "Presenter" checkbox,
% independent of the three version checkboxes the same way "After-course,
% per chapter" already is). Same \ifdefined-hook pattern as
% \deckversionoverride: compiler.py's wrapper file
% \def's it, before \input{}ing the real deck, whenever that checkbox is
% ticked for a job.
%
% Turns on beamer's own built-in dual-screen presenter mechanism:
% \note{...} content (see "Speaker notes" below -- already authored
% throughout this project's own slides, previously never actually emitted
% anywhere at all) becomes a facing "notes page" after every real slide,
% which presenter software (pympress is this project's own recommended
% choice -- see README.md's "Presenting" section) auto-detects and splits
% across the projector (plain slides only) and the presenter's own screen
% (current slide, next slide, and the note text). \usepackage{pgfpages} is
% loaded explicitly rather than trusted to \setbeameroption's own
% auto-load, since which beamer versions actually do that on their own
% isn't worth depending on.
%
% A presenter compile's own CompileJob always carries version="instructor"
% (or the deck's own hardcoded content, for a deck with no
% \deckversionoverride at all) regardless of whether the plain "Instructor"
% checkbox is *also* ticked in the same batch -- notes and full \pause
% reveals only make sense against the instructor's own content, so
% Presenter never needs (or offers) During-course/After-course versions of
% itself. slide_snapshot.py's own compile-time preview refresh is
% deliberately skipped for a presenter job (compiler.py's own
% _relocate_pdf_and_finish()) -- its page count and per-page layout are
% nothing like a normal compile's, so reading it the normal way would
% silently corrupt the real preview cache with wrong or missing pages.
\ifdefined\SMPresenterModeOverride
	\usepackage{pgfpages}
	\setbeameroption{show notes on second screen=right}
	% Presenter-mode note page (direct instruction, 2026-09-10): beamer's
	% own default "note page" template (beamerouterthemedefault.sty,
	% \defbeamertemplate*{note page}{default}) reserves its top quarter for
	% a thumbnail of the current slide (\insertslideintonotes) plus a
	% rotated compile-date printed beside it, then gives the title band
	% (short title + section/subsection/frametitle breadcrumb) the same
	% 0.25\paperheight, leaving only the bottom 0.75 for \insertnote. Both
	% the thumbnail and date are redundant here -- pympress (this project's
	% own recommended presenter software, see README.md) already shows the
	% current and next slide on its own -- so this redefinition drops that
	% whole top block (thumbnail + date) outright and shrinks the title
	% band from 0.25\paperheight to 0.08\paperheight (still one line for
	% the short title plus room for the breadcrumb, per beamer's own
	% layout below, copied unchanged apart from widening each line from
	% beamer's own 0.71-0.75\textwidth, sized to leave room for the now-
	% removed thumbnail, to 0.85-0.95\textwidth), giving \insertnote the
	% remaining ~0.92\paperheight instead of 0.75. No color-theme "note
	% title"/"note page" background is set anywhere in this project (a
	% plain grep confirms it), so beamer's own \ifbeamercolorempty guards
	% around both of those bands are dropped too -- they were already
	% pure no-ops.
	%
	% Registered as a new named variant ("SMPresenter"), not a
	% redefinition of beamer's own "default" -- confirmed directly that
	% \defbeamertemplate*{note page}{default}{...} here errors ("Command
	% \beamer@@tmpop@note page@default already defined"): both the
	% starred and unstarred \defbeamertemplate forms register a variant
	% name via a plain \newcommand internally (beamerbasetemplates.sty),
	% which always errors on a name beamer's own outer theme already
	% registered, star or not. Same pattern this file's own footnote
	% template already uses just above (\defbeamertemplate{footnote}
	% {SMFootnote}{...} then \setbeamertemplate{footnote}[SMFootnote]) --
	% define-then-select instead of overwrite-in-place.
	%
	% \makeatletter/\makeatother wraps this specific block (not the whole
	% \ifdefined) because \beamer@tempbox below is an @-command: this
	% point in the preamble sits between two of this file's own
	% \makeatletter/\makeatother pairs (the previous one closes well
	% above, the next opens well below), so @ is still catcode-other here
	% otherwise -- confirmed directly rather than assumed.
	\makeatletter
	\defbeamertemplate{note page}{SMPresenter}{%
		{%
			\scriptsize
			% The breadcrumb typesets the frame title (twice: once to measure, once
			% to show it), and a \footnote{}/\source{} in a frame title would run each
			% time: the slide's last reveal step then carried its footnote three
			% times (direct report 2026-10-13, mbz). Footnotes are switched off in
			% this group; the slide itself keeps its own, single one.
			\renewcommand*{\footnote}[2][]{}%
			\renewcommand*{\footnotemark}[1][]{}%
			\usebeamerfont{note title}\usebeamercolor[fg]{note title}%
			% Natural-height \vbox (no "to <len>"/\vfil): this project's own
			% real decks call \section{}/subsection ones sometimes not (a
			% plain grep of decks/*.tex confirms both), so a fixed-height
			% band sized for the *shortest* case overlaps \insertnote below
			% on any slide needing all three breadcrumb lines -- confirmed
			% directly: a first attempt at a fixed 0.08\paperheight did
			% exactly that against a real slide's own \note{}. A natural-
			% height box is exactly as tall as whichever of these three
			% lines are non-empty for *this* slide, never more, so it can
			% never collide with the notes below it, and hands \insertnote
			% every bit of space a section/subsection-free slide doesn't
			% need for a breadcrumb. \insertshorttitle (beamer's own first
			% line here, normally the deck's own \title{}) is dropped
			% outright rather than shrunk -- confirmed via grep that
			% \title{} is never called anywhere in this project, so that
			% line has always been silently empty.
			\vbox{\vskip0.3em
				\setbox\beamer@tempbox=\hbox{\insertsection}%
				\hbox{\ifdim\wd\beamer@tempbox>1pt{\hskip4pt\raise3pt\hbox{\vrule
							width0.4pt height7pt\vrule width 9pt
							height0.4pt}}\hskip1pt\hbox{\begin{minipage}[t]{0.95\textwidth}\def\breakhere{}\insertsection\end{minipage}}\fi%
				}%
				\setbox\beamer@tempbox=\hbox{\insertsubsection}%
				\hbox{\ifdim\wd\beamer@tempbox>1pt{\hskip17.4pt\raise3pt\hbox{\vrule
							width0.4pt height7pt\vrule width 9pt
							height0.4pt}}\hskip1pt\hbox{\begin{minipage}[t]{0.95\textwidth}\def\breakhere{}\insertsubsection\end{minipage}}\fi%
				}%
				\setbox\beamer@tempbox=\hbox{\insertshortframetitle}%
				\hbox{\ifdim\wd\beamer@tempbox>1pt{\hskip30.8pt\raise3pt\hbox{\vrule
							width0.4pt height7pt\vrule width 9pt
							height0.4pt}}\hskip1pt\hbox{\insertshortframetitle[width=0.9\textwidth]}\fi%
				}%
				\vskip0.3em}%
		}%
		\nointerlineskip
		\insertnote
	}
	\setbeamertemplate{note page}[SMPresenter]
	\makeatother
\fi
% \SMHideFooterText: opt-in flag suppressing just \FooterText (the left
% part of the footer) -- default false. Set true by the standout template's
% own hook below, so `\begin{frame}[standout]` hides only that left part,
% keeping the "(slide id: ...)" block and frame number visible (2026-07-16,
% direct instruction). Reset immediately after being read, inside the
% footer template itself, rather than relying on a per-environment closing
% hook (standout frames are bare `\begin{frame}[standout]...\end{frame}`,
% not one of this project's own environments like standardframe, so there's
% no custom closing code to reset it in) -- safe because \SMHideFooterText
% is only ever consumed here, once, immediately before being cleared.
\newif\ifSMHideFooterText
% The theme's own \defbeamertemplate*{standout}{SMDefault}
% (beamerthemeSlideMage.sty) is what `\begin{frame}[standout]` actually
% activates; \addtobeamertemplate's "after" hook appends to it without
% needing to touch the theme.
\addtobeamertemplate{standout}{}{\global\SMHideFooterTexttrue}

% \SMSlideLanguageFooterTag -- resolved once here rather than inlined
% into the footer template below, same reason \SMDeckVersionText is
% precomputed rather than built inline in \deckversion. Only ever appears
% appended directly onto the version digits ("v1.1.0EN"), so it's a no-op
% whenever \ifSMSlideVersionSet is false -- ties the language tag's
% rollout to the same version-bump touchpoint slides already go through,
% instead of needing a one-time pass over every slide. For an "mult"
% slide (\multilingualtext above) this shows whichever language the deck
% itself is compiling as, not a fixed value -- \MakeUppercase{en} = "EN".
\newcommand{\SMSlideLanguageFooterTag}{%
	\ifSMSlideLanguageSet
		\ifSMSlideLanguageMult\MakeUppercase{\SMDeckLanguage}\else\slideLanguage\fi
	\fi
}

% \SMFooterIDStrut -- fixed-size strut for the footer's "(slide id: ...)"
% block (2026-09-19, direct instruction: footer text moved by ~0.27pt between
% slides). Cause: the footline's beamercolorbox only has a *minimum* height and
% depth, and the id block is the only footer text that differs from slide to
% slide -- an id containing a descender letter (g, j, p, q, y: "egk", "y89")
% is deeper than the box's minimum depth and lifts the whole footer, an id
% without one ("kaz", "c2t") leaves it at the minimum. The strut spans the
% largest ascender and descender the block can ever contain (measured with the
% block's own font size, at use time, so it follows a font change), which makes
% the footline's height and depth identical for every slide id. Must be
% called after the block's \fontsize.
\newsavebox{\SMFooterIDMeasureBox}
\newcommand{\SMFooterIDStrut}{%
	\sbox{\SMFooterIDMeasureBox}{gjpqy()[]|bdfhklt0123456789}%
	\rule[-\dp\SMFooterIDMeasureBox]{0pt}{\ht\SMFooterIDMeasureBox}%
}

% \SMFooterHide: which footer elements a deck leaves out, a comma list of
% text (the left \FooterText), id ("slide id: ... | v..."), deck ("instructor
% deck" etc.). Empty by default = everything shown. SlideMage's deck
% settings write it as \def\SMFooterHide{id,deck} in the deck's preamble;
% it is read once at \begin{document}, so the deck may set it before or
% after \input{SlideMage_template_...}. The frame number is always shown.
\providecommand{\SMFooterHide}{}
\newif\ifSMFooterShowText
\newif\ifSMFooterShowID
\newif\ifSMFooterShowDeck
\newcommand{\SMFooterHideItem}[1]{%
	\ifstrequal{#1}{text}{\SMFooterShowTextfalse}{}%
	\ifstrequal{#1}{id}{\SMFooterShowIDfalse}{}%
	\ifstrequal{#1}{deck}{\SMFooterShowDeckfalse}{}%
}
\AtBeginDocument{%
	\SMFooterShowTexttrue\SMFooterShowIDtrue\SMFooterShowDecktrue
	\expandafter\forcsvlist\expandafter{\expandafter\SMFooterHideItem\expandafter}\expandafter{\SMFooterHide}%
}

\newif\ifSMNavigationBar % \SMNavigationBar, see "Navigation for readers" at the end
\setbeamertemplate{frame footer}{%
	\csname SM@footertextdrop\endcsname%
	\vspace*{\SMFooterTextDrop}%
	\ifSMHideFooterText
		\global\SMHideFooterTextfalse
	\else
		\ifSMFooterShowText
			\FooterText\ifSMNavigationBar\SMNavigationLinks\fi
		\fi
	\fi
	\hfill%
	{\SMFooterIDSize \SMFooterIDStrut
		\ifSMFooterShowID(slide id: \slideID\ifSMOmittedSlide\SMOmittedSlideMarker\fi\ifSMSlideVersionSet\ | v\slideVersion\SMSlideLanguageFooterTag\fi\ifSMFooterShowDeck\ | \SMDeckVersionText\fi)%
		\else\ifSMFooterShowDeck(\SMDeckVersionText)\fi\fi}\hspace*{2em}\insertframenumber%
	\hspace*{-4em}%
	\setlength{\SMFooterRightShift}{\dimexpr\paperwidth-\SMFrametitleRightPadding-\hsize+\rightskip+\SMRightFooterPadding+\SMFootlineHRightOffset-4em\relax}%
	\hspace*{-\SMFooterRightShift}%
	\vspace*{0.4ex}%
}




% \source{text} -- for a "Source: ..."/"Image source: ..."/etc. citation line
% (common at the bottom/edge of a converted PowerPoint slide) that should look
% and sit *exactly* like the footnotes above -- same font, spacing,
% right-alignment, and the same stretch-to-sit-just-above-the-footer
% behaviour beamer already gives real footnotes -- but with no marker, neither
% inline in the body text nor next to the note itself, since a source line
% isn't anchored to one specific word. Implemented as a real \footnote{} (so
% it gets that positioning for free, rather than reimplementing it) with
% \thefootnote emptied locally; since \footnote{} still advances the footnote
% counter even with an empty marker, \addtocounter reverses that so a real
% \footnote{} elsewhere on the same slide doesn't skip a symbol in the
% \fnsymbol sequence. Unlike a real footnote, no label is added automatically
% here -- the label varies ("Source:", "Image source:", ...), so write it as
% part of #1.
\newcommand{\source}[1]{%
	\begingroup
	\renewcommand{\thefootnote}{}%
	\footnote{#1}%
	\endgroup
	\addtocounter{footnote}{-1}%
}

% A plain tabular with yellow \toprule/\midrule/\bottomrule, matching the
% frametitle/footer accent color. Takes the same column-spec argument as
% tabular itself, e.g. \begin{SMTable}{lccc} ... \end{SMTable}. Only
% booktabs rules are colored (via \arrayrulecolor, scoped to this table by
% \begingroup/\endgroup) -- if a table ever needs a plain black \hline too,
% that's unaffected since \arrayrulecolor only recolors rules drawn *after*
% it takes effect, and booktabs' rules read it fresh each time.
%
% RULE-BLEED WARNING (2026-07-12, direct instruction: "make sure the table
% lines do not extend beyond the text"), read this before writing a new
% SMTable: a plain column spec like `{lccc}` above leaves the default
% \tabcolsep (6pt, or whatever a slide-local \setlength last set it to) as
% *implicit* padding on the table's own outermost edges -- invisible to the
% table's own text (which starts/ends flush with each column's actual
% content), but very visible to booktabs' \toprule/\midrule/\bottomrule and
% to colortbl's \rowcolor/\cellcolor, both of which span the *declared*
% column width including that implicit padding. The result, confirmed
% directly and independently on at least three slides so far
% (slides/u5h-Long-term-Bonds-Provisions-current.tex's own v1-1-0/v1-4-0,
% slides/c4c-Macbeth-Spot-Removers-Proposed-Structure-current.tex's own
% v1-2-2, slides/tw5-Exercise-Leverage-and-ROE-Solution-current.tex's own
% v1-4-0): every rule/highlight bleeds \tabcolsep past the visible text on
% both the left and right edges, most noticeable on a left-aligned
% (non-\centering) table where the mismatch is easy to see against the
% title's own left margin.
%
% (2026-09-30: write new tables as \begin{SMTable}[gap=12pt]{...}
% instead -- see "FLUSH EDGES BUILT IN" below; the rest of this comment
% describes the hand-written form that option now produces.)
% Not fixed here in the environment definition itself: doing so would
% change every existing SMTable's own column spacing out from under it
% (this environment has no way to tell "plain default spacing was intended"
% apart from "the bleed was never noticed"). Fix per table instead, the
% same three-part pattern every slide above uses -- worth reaching for by
% default on any *new* SMTable, not just when the bleed is spotted:
%   {\setlength{\tabcolsep}{0pt}
%   \begin{SMTable}{@{}l<{\hspace{<gap>pt}}c<{\hspace{<gap>pt}}c@{}}
%   ...
%   \end{SMTable}}
% -- \tabcolsep zeroed (scoped locally, e.g. by wrapping the whole call in
% a `{...}` group, so it doesn't leak into later tables), explicit `@{}` at
% both outer edges of the column spec (so booktabs/colortbl have no
% implicit padding left to bleed into), and array's `<{\hspace{<gap>pt}}>`
% column modifier appended after every column but the last, to put back an
% explicit, *content-column* gap (matching whatever `2*\tabcolsep` the
% table used to have, if preserving its exact look, e.g. `<{\hspace{12pt}}>`
% for the ordinary 6pt-\tabcolsep default) -- `<{...}>` puts the gap
% *inside* the preceding column's own content, which `\rowcolor` (if used)
% colors along with everything else, unlike a plain `@{...}` separator,
% which `\rowcolor` deliberately leaves untinted (see u5h's own v1-6-0
% comment for that exact failure mode if a highlighted table ever needs
% this too).
%
% RE-CONFIRMED 2026-07-16, direct instruction, after this exact bleed
% reappeared on essentially every SMTable in a freshly-converted
% chapter (Capital Budgeting Techniques) despite the warning above already
% existing at the time -- the rule above is not optional/only-if-noticed,
% it's the default way to write *any* new SMTable, full stop. Check
% every new SMTable against it before considering a slide done, not
% just when a rendered page happens to make the bleed visually obvious.
%
% BLOCK-RULE-COLOR CONVENTION (2026-07-16, direct instruction): a
% SMTable placed *inside* any block environment (`infoblock`,
% `exampleblock`, `referenceblock`, `definitionblock`, plain `block`, or
% their `SMFloating*Block` wrappers) should use black rules, not this
% environment's own default yellow -- yellow-on-the-block's-own-tinted-
% background reads poorly, unlike yellow-on-plain-white outside a block.
% Override locally, *inside* `\begin{SMTable}{...}`, before `\toprule`
% (not before `\begin{SMTable}` itself -- the environment's own
% opening code sets yellow *after* its `\begingroup`, so a color set
% before that point is simply overwritten; confirmed directly, first
% attempt this same day):
%   \begin{SMTable}{...}
%   	\arrayrulecolor{c_primary_dark}
%   	\toprule
%   	...
%   \end{SMTable}
% Booktabs rules read `\arrayrulecolor` fresh at each rule, and
% `\endgroup` (SMTable's own closing code) resets it back to yellow
% afterward, so this never leaks into a later table. First applied on
% slides/6ux-Example-51-BMA-Payback-Table-current.tex's own
% `SMFloatingInfoBlock`.
%
% FLUSH EDGES BUILT IN (2026-09-30, direct instruction: "propose a more
% systematic solution" for per-slide fixes -- 53 tables carried the
% three-part pattern above by hand): give the column gap as an option and
% write the plain column spec,
%   \begin{SMTable}[gap=12pt]{lccc} ... \end{SMTable}
% which produces exactly the hand-written form -- \tabcolsep 0pt, @{} at
% both edges, <{\hspace{12pt}} after every column that is followed by
% another one (not before an explicit @{...}/!{...} separator), *{n}{...}
% expanded -- so a table rewritten this way renders identically. After a
% fixed-width p/m/b/w column the gap is an @{\hspace{...}} separator
% instead: array places a <{...} suffix inside such a column's box, where
% it only pads the end of the cell text and leaves the columns touching
% (the hand-written pattern had exactly that flaw for p columns). The
% expansion is done by SlideMage_measure.lua (slidemage_measure.gap_spec).
% gap=12pt is the ordinary look (2 x beamer's 6pt \tabcolsep). Without the
% option the environment behaves as it always has (plain tabular with
% yellow rules), so every existing table is untouched.
\makeatletter
\pgfqkeys{/SMTable}{
	gap/.store in=\SMTable@gap,
}
\ifluatex
	\SMLuaOnce{
		local path = kpse.find_file("SlideMage_measure.lua", "lua")
			or kpse.find_file("SlideMage_measure.lua", "tex")
		if path then
			dofile(path)
		end
	}
	% The two numbers the Top Margin measurement needs. The ascender is read from the text font in
	% force (luaotfload's OpenType metrics) and replaces the layout profile's \SMFontAscender where
	% the font gives one; the Top Margin is the template's (SlideMage_setup.tex). Not replayed by
	% a precompiled format: the macro it sets lives in the format.
	\begingroup\normalfont\selectfont\directlua{slidemage_measure.derive_ascender()}\endgroup
	\SMLuaOnce{
		slidemage_measure.configure{top_margin = "\SMTopMargin", ascender = \SMFontAscender}
	}
\fi
\hypersetup{pdfinfo={SMFontAscender={\SMFontAscender}}}
\def\SMTable@begin#1{\begin{tabular}{#1}}
% BLOCK-RULE-COLOR, automatic (2026-09-30): inside any block (the
% convention above) the rules default to c_primary_dark -- all 11 tables
% placed in a block so far set that by hand. An \arrayrulecolor inside the
% table still overrides it, as before.
\newif\ifSM@inblock
\AtBeginEnvironment{block}{\SM@inblocktrue}
\AtBeginEnvironment{exampleblock}{\SM@inblocktrue}
\AtBeginEnvironment{alertblock}{\SM@inblocktrue}
\AtBeginEnvironment{infoblock}{\SM@inblocktrue}
\AtBeginEnvironment{referenceblock}{\SM@inblocktrue}
\AtBeginEnvironment{definitionblock}{\SM@inblocktrue}
\NewDocumentEnvironment{SMTable}{O{}m}{%
	\begingroup
	\ifSM@inblock
		\arrayrulecolor{c_primary_dark}%
	\else
		\arrayrulecolor{c_primary_highlight}%
	\fi
	\let\SMTable@gap\@empty
	\pgfqkeys{/SMTable}{#1}%
	\ifx\SMTable@gap\@empty
		\def\SMTable@spec{#2}%
	\else
		\setlength{\tabcolsep}{0pt}%
		\directlua{slidemage_measure.set_gap_spec(
			"\luaescapestring{\detokenize{#2}}",
			"\luaescapestring{\detokenize\expandafter{\SMTable@gap}}")}%
	\fi
	\expandafter\SMTable@begin\expandafter{\SMTable@spec}%
}{%
	\end{tabular}%
	\endgroup
}
\makeatother

% For a side-corner image overlay (the usual \vbox to 0pt{...} tikz overlay
% pattern used throughout this deck's converted slides) whose *top edge*
% should line up with the content text's own top edge -- distance from the
% page's own top edge to the cap-height top of a standardframe's first body
% line, measured directly against the compiled PDF's own text coordinates,
% not derived from \SMContentTopShift or guessed. Previous per-slide
% image placements each guessed their own downward shift (e.g. a bare
% "-1.3cm") independently -- close, but not exact, landing the image's top
% a couple mm above the text's own top; use this constant instead of a
% fresh guess for any new "text top edge"-aligned image.
%
% 41.58pt (2026-07-22, re-measured): the previous value, 44.0pt, went stale
% the moment \SMContentTopShift itself was increased (7.33pt -> 11.5964pt,
% "reduce the Top Margin", see that length's own comment above) without a
% matching update here -- that change pulls plain body text up, but this
% length is a separate, hand-measured constant with no formula tying it to
% \SMContentTopShift, so it kept anchoring images at the *old*, lower
% position. Confirmed directly (pymupdf against this same freshly-compiled
% preview deck, not assumed from the \SMContentTopShift delta alone):
% three independent plain-first-line standardframes with no per-slide
% top-shift correction (slides/5v4-CML-vs-SML-current.tex,
% slides/u5j-Exercise-CAPM-4-current.tex,
% slides/nt4-Exercise-CAPM-5-current.tex) all land their first line's own
% top edge at y=41.58pt post-change, not 44.0pt. Re-measure and adjust
% again if \SMContentTopShift ever changes in the future.
% (v9.7.0: \SMContentTopAlign is derived from \SMTopMargin, further up; it was a hook of SlideMage_setup.tex, 41.58pt, in v9.4.0-v9.6.0.)

% For a "side filler image" -- a decorative image stuck flush to the page's
% own right edge (zero margin) with its top edge flush against the title
% bar's own bottom edge (zero gap), rather than sitting in the usual
% content-aligned corner inset -- distance from the page's own top edge
% down to the standardframe title bar's bottom edge, measured directly
% against the compiled PDF (a pixel scan for the yellow-to-white
% transition, at this deck's 16cm-wide/9cm-tall page size). First used on
% slides/dfc-Angel-Investors-current.tex.
%
% 25.43pt -- the raw measured value (originally 25.44pt; re-measured
% 2026-07-07 against the current compiled PDF via the title bar's own
% white-patch rectangle, which came back 25.427pt, close enough to the
% original measurement to be the same underlying edge, not a drift).
% Briefly set to 24.94pt instead (2026-07-06-ish), a deliberate 0.5pt
% overlap *into* the title bar, after a report of a visible hairline
% white gap right at this seam in one render -- not reproduced in this
% project's own pdftoppm-based rendering checks at the time (every pixel
% scan at the raw value showed a clean cut, no gap), so assumed to be a
% renderer-specific sub-pixel rounding artifact rather than a real error
% in the constant, and "fixed" by deliberately overlapping instead of
% chasing the specific viewer. That overlap turned out to be the wrong
% call: direct feedback 2026-07-07 reported the *opposite* symptom on
% `sideimageframe` (the image's own top edge visibly overlapping the
% title box) -- exactly what a real 0.5pt overlap looks like once it's
% actually noticed, rather than a rendering artifact. Reverted back to
% the precise raw measurement: a real, small, deliberate overlap chosen
% to paper over an unreproduced gap in an unspecified other renderer
% turns out to be a worse trade than just using the measured value and
% letting a genuine future gap report (if one ever recurs, in a
% specific, identified renderer) be diagnosed on its own terms.
%
% Re-measured 2026-07-16, 25.43pt -> 25.32pt: the frametitle hook's
% outer sep fix (see its own comment) moved the visible yellow bottom
% edge up by the 0.2pt the trim patch's fill had previously been inset
% below its SMLogoAnchor.south anchor -- the yellow now genuinely
% ends flush with the logo's bottom edge, as always intended. New raw
% measurement against the rebuilt PDF: the trim patch's fill top edge
% (= the visible yellow bottom) sits 25.228 PDF units below the page
% top, i.e. 25.3226pt (earlier measurements in this comment quoted PDF
% coordinate values directly as "pt"; PDF units are big points, 1bp =
% 1.00375pt -- a systematic ~0.1pt understatement at this magnitude,
% harmless then, but corrected now that the constant is being reset
% anyway). Note this value is also what the overlay-logo nodes
% (titleframe/quoteframe/the hook's fixed-position redraw) derive their
% vertical inset from: visible-bar-bottom minus \SMLogoHeight is
% exactly the inline logo's own vertically-centered top margin.
% (v9.7.0: \SMTitleBarBottomAlign is derived from the paddings, the frame title font and the logo height, further up; it was a hook of SlideMage_setup.tex, 25.32pt, in v9.4.0-v9.6.0.)

% Note on the [t] default below: it top-aligns content within the frame
% (beamer's own default is vertical centering, "c"), so slides don't have
% their content floating in the middle. It also sidesteps a real beamer bug:
% \begin{frame}[]{Title} (a literal empty options bracket) fails outright,
% even in plain beamer with no theme at all -- so the optional argument must
% never default to an empty string. The title is always passed as an
% explicit argument (rather than left for beamer to pick up from a trailing
% brace group), since anything placed between \begin{frame} and its title --
% like our \setbeamertemplate calls -- breaks that trailing-brace mechanism.

% Pulls the frame's content up closer to the title bar: measured (via the
% PDF's own coordinates) at about 22pt of gap by default between the title
% bar's bottom edge and the first line of body content: this is a third of
% that, shrinking the gap to about two thirds of its original size.
%
% Called "Top Margin" in conversation from 2026-07-12 onward (direct
% instruction) -- the deck-wide name for the resulting gap between the
% title bar and a standardframe's first line of body content, i.e. the
% vertical position a plain-text frame's first baseline naturally lands at
% (measured directly against a reference slide's own compiled PDF, e.g.
% slides/aii-Leverage-Effects-current.tex's first line). Any first content
% that isn't plain text (a table, a minipage, an itemize with an unusual
% first item, ...) needs its own per-slide \SMIncreaseContentTopShift/
% \SMReduceContentTopShift correction to land on this same baseline --
% not portable between slides even for superficially similar content (see
% slides/zgz-EBIT-Definition-current.tex's own v1-2-0 comment for a case
% where an earlier, wrong correction was carried over uncritically).
% (\SMContentTopShift was a number of the layout profile from v9.4.0 to v9.8.0 and is measured
% since v9.9.0, see AUTOMATIC TOP MARGIN; its history is below.)
% 11.5964pt (2026-07-22, direct instruction: "reduce the Top Margin"
% generally, tested at 20%/30%/40% smaller against
% slides/8s8-Unit-Investment-Trusts-current.tex before settling on 40% --
% up from the original 7.33pt so the *resulting* gap (this length pulls
% content up by, so a bigger value means a smaller visible margin) comes
% out to 60% of its old size: measured directly on 8s8 at 10.67pt before,
% 6.42pt after (target was 6.40pt, within 0.02pt). Since every
% standardframe applies this same pull-up before any per-slide
% \SMIncreaseContentTopShift/\SMReduceContentTopShift correction,
% changing it here shifts the baseline project-wide, not just on 8s8.

% A `columns` environment that is the first content of a frame starts its first line
% ~10.17pt lower than a plain list does (measured on 5s6/vbb/yxv/c2t/n83/tj4/eey/yvw: first
% line box at 51.7pt against the 41.58pt Top Margin). It used to be corrected slide by
% slide with \SMIncreaseContentTopShift{10.13..10.17pt}, each value re-measured (and
% re-tuned when \SMContentTopShift changed); now the environment corrects itself
% (2026-09-21, direct instruction: "if this is a general problem for two-column slides, fix
% it generally"). "First content" = the frame's vertical list holds no box taller than
% 0.5pt yet -- the frame's own zero-size stamp boxes, glue and a zero-height overlay picture
% (\vbox to 0pt{...}) in front of the columns do not count, so a slide that puts such a
% picture first needs no special treatment either. (\prevdepth cannot tell: those zero-size
% boxes already set it.) Checked in LuaTeX by walking the current list.
% (\SMColumnsTopCorrection was a number of the layout profile from v9.4.0 to v9.8.0.)
\SMLuaOnce{
	slidemage_list_has_no_visible_box = function()
		local head = tex.nest[tex.nest.ptr].head
		for n in node.traverse(head) do
			if (n.id == node.id("hlist") or n.id == node.id("vlist")) and n.height > 32768 then
				return false
			end
		end
		return true
	end
}
% Since 2026-09-30 only for frames the automatic Top Margin doesn't manage
% (see AUTOMATIC TOP MARGIN below: [c] frames, sectionframe,
% \SMManualTopMargin, a deck with \SMAutoTopMarginfalse) -- in every
% other frame the measurement puts a columns-first slide on the margin by
% itself, so this correction (and the [T] compensation below) stays out.
% v9.9.0: the two corrections are measured where they are used, not numbers of the layout
% profile (10.17pt and 13.76pt, fitted on Fira; with Latin Modern Sans under another title bar a
% [T] columns landed 5.2pt high). Three reference lists are typeset in boxes at this point -- a
% top-level itemize, the same itemize in the first column of a columns[t], and of a columns[T]
% -- and \SMColumnsTopCorrection is how much lower the [t] one starts, \SMColumnsTTopCompensation
% how much of that the [T] one has taken back; a list in a column of either kind then starts
% where the list would start without the columns. Both names stay readable (the last values
% measured, 0pt before the first columns of the deck).
\makeatletter
\newcommand{\SM@columnsprobe}{%
	\SM@toprefprobe{\begin{itemize}\item\SM@topreftext\end{itemize}}\let\SM@columnsref\SM@topref
	\SM@toprefprobe{\begin{columns}[t]\begin{column}{0.46\linewidth}\begin{itemize}\item\SM@topreftext\end{itemize}\end{column}\end{columns}}%
	\global\SMColumnsTopCorrection=\dimexpr\SM@topref sp-\SM@columnsref sp\relax
	\SM@toprefprobe{\begin{columns}[T]\begin{column}{0.46\linewidth}\begin{itemize}\item\SM@topreftext\end{itemize}\end{column}\end{columns}}%
	\global\SMColumnsTTopCompensation=\dimexpr\SMColumnsTopCorrection+\SM@columnsref sp-\SM@topref sp\relax
}
\BeforeBeginEnvironment{columns}{%
	\ifSM@topprobing\else
	\ifnum0\ifdefined\SMTopAutoActive\SMTopAutoActive\fi=1
		\def\SMColumnsAtTop{0}%
	\else
		\edef\SMColumnsAtTop{\directlua{tex.sprint(slidemage_list_has_no_visible_box() and "1" or "0")}}%
	\fi
	\ifnum\SMColumnsAtTop=1
		\ifluatex\ifvmode\SM@columnsprobe\fi\fi
		\vspace*{-\SMColumnsTopCorrection}%
	\fi
	\fi
}
\makeatother
% The correction above is measured for `columns[t]`. `columns[T]` (top-aligned, used on
% most text-next-to-image slides) already starts its columns higher by itself -- beamer's
% own T key adds \vskip-1ex\nointerlineskip at every column head -- so the same pull-up left
% its first line 13.70pt (itemize) / 13.82pt (plain text) above the Top Margin, flush
% against the title bar (2026-09-30, direct instruction, seen on c58-DE-Unikat-Tipps).
% Measured against a plain standardframe: itemize 41.58pt, [t] 41.58pt, [T] 27.89pt;
% plain text 38.60pt, [T] 24.78pt. The T key gives the difference back while the columns
% environment is still in vertical mode, i.e. only for the columns[T] call itself (the
% per-column \setkeys calls run inside the columns' hbox, where \ifvmode is false) and only
% when the hook above applied its pull-up. A block as first content of a [T] column lands
% ~7pt below its plain-frame position with this -- extra clearance, not a collision.
% The column head also sets LaTeX's "top of a minipage" flag again (v5.14.1, found by
% check_layout.py's `coltop` in the SlideMage manual): beamer opens each column with
% \leavevmode, which clears it, so a list at the top of a [T] column kept its \topsep --
% zero at normal size, but a smaller list (\shrinklists, \small) loads \@listi's 4pt and
% started that much below a figure beside it that sat on the Top Margin.
% (\SMColumnsTTopCompensation was a number of the layout profile from v9.4.0 to v9.8.0; it is
% measured since v9.9.0, see above, and the figures in this paragraph are history.)
\makeatletter
\define@key{beamer@col}{T}[true]{%
	\def\beamer@colmode{T}\def\beamer@colalign{t}\def\beamer@colheadskip{\vskip-1ex\nointerlineskip\@minipagetrue}%
	\ifvmode
		\ifnum0\ifdefined\SMColumnsAtTop\SMColumnsAtTop\fi=1
			\vspace*{\SMColumnsTTopCompensation}%
		\fi
	\fi
}
\makeatother


% "Top Margin" (above) is one of four deck-wide margin names introduced in
% conversation from 2026-07-12 onward (direct instruction) for the plain
% page-edge/title-bar/footer gaps every slide is expected to obey by
% default. The other three, for reference (each maps to a length already
% declared elsewhere, not a new one introduced here):
% - "Left Margin": the default gap from the page's own left edge to the
%   title text and body content -- \SMFrametitleLeftPadding
%   (beamerthemeSlideMage.sty, 2.2ex, measured
%   directly against the compiled PDF at 12.552pt in this deck's usual
%   fonts -- see e.g. slides/tw5-Exercise-Leverage-and-ROE-Solution-current.tex's
%   own v1-1-0 comment for that same measured value).
% - "Right Margin": the default gap from the page's own right edge to the
%   institute logo (and, by the same convention, to anything else
%   deliberately flush-right, e.g. a corner-anchored image) --
%   \SMFrametitleRightPadding (0.28cm/7.93pt, this file's own
%   \SMInstituteLogo placement above uses it directly as the
%   logo's own inner sep).
% - "Bottom Margin": the default gap between a frame's own body content
%   and the \source{}/footnote line below it. Unlike the other three, this
%   one has no single settable length behind it -- \source{}'s own
%   vertical position isn't fixed the way the footer box is; it tracks the
%   frame's actual content height once that content gets tall enough (see
%   slides/u5h-Long-term-Bonds-Provisions-current.tex's own v1-3-0 comment,
%   "past a certain table height it starts tracking the table's own growth
%   roughly 1:1"). "Respecting Bottom Margin" in practice means leaving
%   comparable clearance to whatever a reference slide with unremarkable
%   content already shows, the same empirical, measure-against-a-reference-
%   slide approach "Top Margin" itself uses, rather than a fixed constant.

% Compensates a gap that opens up specifically when a standardframe's body
% starts with a \begin{minipage}{...} (used to width-limit text next to an
% image/diagram overlay, e.g. slides/igr-Budgeting-current.tex,
% slides/bdh-Working-Capital-Example-current.tex -- see README.md's
% "Available layouts" section for why minipage instead of columns there) --
% a minipage is a *box*, and TeX's vertical glue insertion (\topskip
% specifically) treats the first box in a vertical list differently from
% the first line of plain paragraph text, so the same \SMContentTopShift
% pull-up above lands closer to the title for minipage-wrapped content than
% for plain content. Cannot be fixed by \vspace(*) right after
% \begin{minipage}{...}: both plain and starred forms were tried and
% measurably had zero effect (the star only protects glue from page-break
% discard, not from \topskip's own separate glue-insertion calculation,
% which absorbs leading glue at the start of a vertical list regardless).
% Fixed instead by shrinking standardframe's own pull-up for the affected
% frame -- see \SMContentTopShiftReduction below. The value below is
% calibrated empirically, not derived from first principles: an isolated
% two-frame test (single-level itemize, one plain, one minipage-wrapped)
% measured only a 4.4pt gap, but real slides (nested itemize inside the
% minipage) showed a larger ~7.5-7.6pt gap against a real plain-itemize
% slide (slides/jyn-Categorization-current.tex) -- the extra few points
% likely come from itemize's own topsep/partopsep being computed differently
% depending on whether the list opens a box or a paragraph, on top of the
% topskip/box effect alone. Tuned so slides/igr-Budgeting-current.tex and
% slides/bdh-Working-Capital-Example-current.tex's first content line
% lands within 0.1pt of slides/jyn-Categorization-current.tex's own
% (200.917pt vs. 200.871pt/200.980pt in the compiled PDF).
% v9.4.0: the four corrections below (SMMinipageTopCorrection, SMMinipageWithCornerImageExtraCorrection,
% SMSideImageFrameMinipageCorrection, SMFlatItemizeTopCorrection) are 0pt. The automatic Top Margin (below) moves
% every frame it manages onto the Top Margin, so the slides that pass them to \SMReduceContentTopShift or
% \SMIncreaseContentTopShift land where they did: a copy of a real project's 16 decks (5274 pages, 122 slides
% that read SMMinipageTopCorrection, 123 SMFlatItemizeTopCorrection) rendered pixel-identically with the
% values at 0. The names stay defined because slides read them. The numbers in the comments are history.
\newlength{\SMMinipageTopCorrection}
% 9.38pt (2026-07-22, re-measured): the previous 7.6pt value went stale
% the moment \SMContentTopShift itself was increased (7.33pt ->
% 11.5964pt, "reduce the Top Margin", see that length's own comment above)
% without a matching update here -- confirmed directly (pymupdf against a
% freshly-compiled preview deck, not assumed from the \SMContentTopShift
% delta alone, which doesn't translate 1:1 into a minipage's own \topskip
% behavior): four independent slides all using this constant completely
% unmodified (slides/fr3-Forms-of-Capital-Market-Efficiency-Strong-Semistrong-Weak-current.tex,
% slides/za2-Startup-Financing-Categorization-current.tex,
% slides/yke-Implications-of-Homemade-Dividends-current.tex,
% slides/6j4-Modigliani-Miller-Assumptions-current.tex) all land their
% own first line at an identical y=39.80pt post-change, 1.78pt short of
% this deck's current 41.58pt Top Margin reference -- not slide-specific
% noise, the same deterministic amount on every one of them. Re-measure
% and adjust again if \SMContentTopShift ever changes in the future.
\setlength{\SMMinipageTopCorrection}{0pt}% obsolete since v9.4.0 (9.38pt before): the automatic Top Margin corrects it; slides still read the name

% A *third* minipage-top correction, on top of \SMMinipageTopCorrection
% above rather than replacing it, for a minipage whose own first internal
% content is a zero-height `\vbox to 0pt{...}` tikz overlay (a
% corner-anchored decorative image, placed there specifically so it stays
% unconditionally visible even when the minipage's own later content
% includes a `\pause` -- see e.g.
% slides/y5z-CoCo-Bonds-current.tex's own "v1-2-0" note for the fuller
% pause-inheritance story behind why the image has to sit there at all).
% 2026-07-08, found while chasing "the top margin of the text on y5z/22s
% is too low" -- measured directly against the compiled PDF, y5h/22s both
% landed 16.5pt too far down despite already using
% \SMMinipageTopCorrection, with the zero-height image overlay sitting
% *before* the minipage entirely as the actual root cause (that ordering
% makes the image, not the minipage, the frame body's true first box,
% which \topskip's own special first-box glue handling reacts to
% differently than a plain first minipage). Moving the image inside the
% minipage as its own first content closed most of the gap on its own
% (restoring the minipage as the frame's true first box) but left a
% further 10.35pt uncorrected -- confirmed the same on two independent
% slides (y5z, 22s) with identical structure, not assumed transferable
% from one to the other. Call via
% `\SMReduceContentTopShift{\dimexpr\SMMinipageTopCorrection-\SMMinipageWithCornerImageExtraCorrection\relax}`
% for any slide combining this exact pattern (corner image as a minipage's
% own first content, ahead of a `\pause`) -- confirmed transferable as-is
% to two more slides with the same structure
% (slides/qbj-Bank-Loans-Credit-Lines-current.tex,
% slides/jkf-Bank-Loans-Term-Loan-current.tex), neither of which
% actually uses `\pause` itself, so this isn't specific to that -- the
% image-inside-minipage-as-first-content shape alone is what triggers it.
\newlength{\SMMinipageWithCornerImageExtraCorrection}
\setlength{\SMMinipageWithCornerImageExtraCorrection}{0pt}% obsolete since v9.4.0 (10.35pt before), see above

% A *different* minipage-top correction, specifically for a minipage+
% itemize sitting as the first content of a `sideimageframe` (2026-07-06,
% found while chasing "the top margin of the text on ezz/acw is again
% off" -- see CONVERSION_NOTES.md's own "beamer/TeX gotchas" entry for the
% full story). `\SMMinipageTopCorrection` above is calibrated for a
% *plain* standardframe, and simply does not transfer to `sideimageframe`:
% measured directly on three real sideimageframe slides
% (slides/xmp-Implications-of-Capital-Market-Efficiency-Considerations-current.tex,
% slides/dfc-Angel-Investors-current.tex,
% slides/ezz-Examples-of-Tests-for-Market-Efficiency-current.tex --
% one flat itemize, one flat itemize, one nested-first-item itemize
% respectively, so this isn't a flat-vs-nested effect either), all three
% landed within 0.6pt of each other at their own *uncorrected* natural
% position (52.03-52.57pt against this deck's usual 44.0pt first-line
% reference) regardless of which of the two *wrong* corrections
% (`\SMReduceContentTopShift{\SMMinipageTopCorrection}` or
% `\SMIncreaseContentTopShift{\SMFlatItemizeTopCorrection}`) had
% been applied to each beforehand -- i.e. `sideimageframe` itself pushes
% its own first-content minipage down by a consistent ~8.2pt relative to
% a plain standardframe, on top of whatever generic itemize-shape
% correction would otherwise apply, most likely from the extra
% \sidefootnotes/\settowidth setup work `sideimageframe`'s own opening
% code does before `\begin{standardframe}` ever runs. Use
% `\SMIncreaseContentTopShift{\SMSideImageFrameMinipageCorrection}`
% (note: *increase*, not reduce) as the *only* topshift call on any
% future sideimageframe slide whose first content is a minipage+itemize --
% not the generic `\SMMinipageTopCorrection`/
% `\SMFlatItemizeTopCorrection` this deck uses everywhere else.
\newlength{\SMSideImageFrameMinipageCorrection}

% 6.26pt (2026-07-22, re-measured): stale for the same reason
% \SMMinipageTopCorrection above went stale -- \SMContentTopShift
% itself changed (7.33pt -> 11.5964pt) without a matching update here.
% Re-measured directly against a freshly-compiled preview deck across six
% independent sideimageframe slides using this constant unmodified: four
% (ezz, xmp, hzx, u6n) land their own first line at an identical
% y=39.64pt, 1.94pt short of this deck's current 41.58pt Top Margin
% reference (i.e. already pulled up *too far*); the other two (c7w, dfc)
% are close but not identical (39.53pt/-2.05, 40.18pt/-1.4) -- no single
% value fixes all six exactly, so this recalibrates to the four-way
% cluster (the majority case) and accepts a small (<0.6pt) residual on the
% two outliers. Since this constant feeds
% `\SMIncreaseContentTopShift{}` (more pull-up, i.e. *smaller* y), and
% the observed y was already short of target, the fix is a *decrease*
% (8.2pt - 1.94pt), not an increase -- caught by re-measuring after a
% first, wrong-direction attempt at 10.14pt drove the four-slide cluster
% to y=37.70pt, further from target instead of closer.
\setlength{\SMSideImageFrameMinipageCorrection}{0pt}% obsolete since v9.4.0 (6.26pt before, 8.2pt originally), see above

% Opt-in override: reduces standardframe's \SMContentTopShift pull-up by
% this amount for the frame currently being opened. Empty-default-style flag
% (mirrors \SMChapterLogoPath etc.): 0pt by default (no effect), set via
% \SMReduceContentTopShift{<length>} called *before* \begin{standardframe},
% reset back to 0pt in standardframe's own closing code so it never leaks
% into the next frame. Used by minipage-based slides to counteract
% \SMMinipageTopCorrection above.
\newcommand{\SMContentTopShiftReduction}{0pt}
\newcommand{\SMReduceContentTopShift}[1]{\global\def\SMContentTopShiftReduction{#1}}

% A *second*, separate first-line-position quirk, found and fixed
% 2026-07-05 after a run of slides (batch 4 plus several earlier ones) were
% found to land at two visibly different heights depending purely on
% content shape, with no minipage involved at all: a plain (non-minipage)
% standardframe whose top-level itemize's very *first* \item is immediately
% followed by its own nested \begin{itemize}/\begin{enumerate} lands 1.834pt
% higher than one whose first \item has no such immediate child (checked
% against every content slide converted so far via each compiled page's own
% first-line PDF baseline coordinate, not eyeballed -- 13 "nested-first-item"
% slides all landed within 0.05pt of each other, and 9 "flat-first-item"
% slides all landed within 0.05pt of each other, at a clean +1.834pt offset
% from the first group; isolated two-frame test cases, differing only in
% presence/absence of that immediate nested list, reproduced the exact same
% 1.834pt gap independent of which text/list type/item count was used).
% Not chased to one named TeX primitive (most likely some interaction of
% \addvspace's glue-coalescing behaviour with \partopsep/\topsep at the
% \begin{itemize}/\begin{frame} boundary, given the existing
% \SMMinipageTopCorrection comment above already flags topsep/partopsep
% as varying with box vs. paragraph context) -- calibrated empirically like
% \SMMinipageTopCorrection, not derived from first principles.
% Interacts with minipage differently, not additively: inside a minipage,
% this nested-vs-flat distinction disappears entirely (both land within
% 0.05pt of each other) -- so this correction is for *plain* itemize-first
% standardframes only, never combined with \SMMinipageTopCorrection.
\newlength{\SMFlatItemizeTopCorrection}

% 0.01pt (2026-07-22, re-measured): stale for the same reason
% \SMMinipageTopCorrection above went stale -- \SMContentTopShift
% itself changed (7.33pt -> 11.5964pt) without a matching update here.
% Re-measured directly against three independent flat-itemize slides using
% this constant unmodified (3ce, 5cr, c8e): all three landed their own
% first line at an identical y=39.76pt, 1.82pt short of this deck's
% current 41.58pt Top Margin reference -- coincidentally, the recalibrated
% \SMContentTopAlign reference now already accounts for almost all of
% what this correction used to supply on its own, leaving next to nothing
% left to add. Effectively zeroed out rather than removing the call
% itself, so every flat-itemize slide's own
% `\SMIncreaseContentTopShift{\SMFlatItemizeTopCorrection}` line
% keeps working unchanged.
\setlength{\SMFlatItemizeTopCorrection}{0pt}% obsolete since v9.4.0 (0.01pt before, 1.834pt originally), see above

% Opposite-direction counterpart to \SMReduceContentTopShift: *increases*
% the pull-up instead of reducing it (needed for \SMFlatItemizeTopCorrection
% above, since a flat first item needs *more* pull-up, not less). Same
% storage variable, just negated -- \SMContentTopShiftReduction is
% allowed to go negative, since the frame-opening code only ever computes
% \SMContentTopShift-\SMContentTopShiftReduction, never checks its
% sign.
\newcommand{\SMIncreaseContentTopShift}[1]{\global\def\SMContentTopShiftReduction{-#1}}

% AUTOMATIC TOP MARGIN (2026-09-30, direct instruction: "propose a more
% systematic solution" -- 195 slides carried a hand-measured
% \SMIncrease/ReduceContentTopShift, 89 of them a raw pt value, every one
% of them re-measured by hand whenever \SMContentTopShift or a frame type
% changed, and a stale one survived a central change unnoticed more than
% once, e.g. six slides double-corrected after the 2026-09-21 columns fix).
% Where a frame's first line lands depends on what comes first -- a plain
% paragraph sits 3pt higher than a list, an image 3.5pt higher, a minipage,
% columns, a table each differently -- so no constant pull-up can be right
% for all. Instead every top-aligned standardframe (and figureframe,
% highlightframe, sideimageframe) now measures where its first visible
% content actually landed (at shipout, SlideMage_measure.lua; overlays at
% absolute positions don't count, they are placed on the margin already) and
% moves its body onto the Top Margin line on the next run, via \jobname.top
% -- the same record-then-apply-on-rerun pattern as \jobname.ovl above;
% latexmk reruns while the file changes, and one rerun converges.
% Any older \SMIncrease/ReduceContentTopShift call still works but no
% longer matters: the measurement corrects whatever is left.
% Exceptions: [c] frames (centred on purpose: no pull-up at all), sectionframe
% (own geometry), and any frame opened after
%   \SMManualTopMargin
% which keeps the frame's default position: no correction that depends on what
% the frame starts with. (A \progressstrip frame is measured since v5.30.5.)
% The default position (v9.9.0) is \SMContentTopShift, the pull-up that puts the
% first line of a top-level list on the Top Margin in a standard frame. It used to
% be a number of the layout profile (11.5964pt, fitted on Fira under the default
% title bar: with another title bar or font the frames above landed up to 2.9pt
% off); now every top-aligned frame records where its body begins (the origin
% marker) and how far below that a top-level itemize's first line would start
% (typeset in a box there), and the measurement writes the pull-up most frames
% share into \jobname.top. The corrections of the measured frames are relative
% to it. In a run without \jobname.top the first titled frame's own value is used
% from the next frame on, so the first run already places the deck as the next one
% will (no deck needs a run more than with the old stored number); a frame that
% started that run elsewhere (the first one) is corrected from its default position. \SMAutoTopMarginfalse in a deck's preamble places every frame like
% \SMManualTopMargin does (the default position, measured; no content-dependent
% correction anywhere in the deck); it does not switch the measurement off.
\newif\ifSMAutoTopMargin
\SMAutoTopMargintrue
\makeatletter
\newif\ifSM@topmanual
\newif\ifSM@topsection
\newcommand{\SMManualTopMargin}{\global\SM@topmanualtrue}
\newcount\SM@topindex
\newwrite\SM@topfile
\newcommand{\SMSetTopCorrection}[2]{\global\@namedef{SM@topcorr@#1}{#2}}
\newcommand{\SMSetContentTopShift}[1]{\global\SMContentTopShift=#1\relax
	\ifluatex\directlua{slidemage_measure.top_shift_read = \strip@pt\SMContentTopShift}\fi}
\ifluatex
	\SMLuaOnce{
		luatexbase.add_to_callback("pre_shipout_filter", function(head)
			if slidemage_measure and slidemage_measure.top_measure_page then
				slidemage_measure.top_measure_page(head)
			end
			return true
		end, "slidemage top margin")
	}
	\AtBeginDocument{%
		\InputIfFileExists{\jobname.top}{}{\typeout{No file \jobname.top.}}%     % "No file": see \jobname.ovl
		\immediate\openout\SM@topfile=\jobname.top\relax
	}
	\AtEndDocument{\directlua{slidemage_measure.top_emit()}\immediate\closeout\SM@topfile}
\fi
% Before \begin{frame} (runs once per frame, not once per overlay page):
% number the frame and fetch its correction from the previous run.
\newcommand{\SM@topprepare}{%
	\ifdefined\SM@topshiftrun% a first run: measured on the first frame that counts
		\global\SMContentTopShift=\SM@topshiftrun\relax
		\global\let\SM@topshiftrun\undefined
	\fi
	\global\advance\SM@topindex by 1
	\xdef\SM@topthis{\the\SM@topindex}%
	\ifcsname SM@topcorr@\SM@topthis\endcsname
		\xdef\SM@topcorr{\csname SM@topcorr@\SM@topthis\endcsname}%
	\else
		\gdef\SM@topcorr{0pt}%
	\fi
	\ifSMAutoTopMargin\ifSM@topmanual\gdef\SM@topcorr{0pt}\fi\else\gdef\SM@topcorr{0pt}\fi
}
% In the frame body, right after the pull-up: the start marker, then the origin marker.
\newcommand{\SM@topstart}{%
	\ifluatex
		\directlua{slidemage_measure.top_begin(\SM@topthis, \strip@pt\dimexpr\SM@topcorr\relax,
			\number\dimexpr\SMContentTopShift-\SMContentTopShiftReduction+\SM@topcorr\relax,
			\strip@pt\SMContentTopShift, \ifSMAutoTopMargin\ifSM@topmanual false\else true\fi\else false\fi)}%
		\ifSMAutoTopMargin\ifSM@topmanual\else
			\def\SMTopAutoActive{1}% this frame is measured: the columns hook stays out
		\fi\fi
		\SM@toprefprobe{\begin{itemize}\item\SM@topreftext\end{itemize}}%
		\directlua{slidemage_measure.top_origin(\SM@topthis, \SM@topref,
			\ifSM@topsection false\else\ifSM@toptitled true\else false\fi\fi)}%
	\fi
}
% The reference list's first line box, in sp below this point (\SM@topref): #1 typeset in
% a box that starts as the body does here (the same \prevdepth, font, \parskip and list
% settings), measured by SlideMage_measure.lua like a frame's first content. Nothing of it
% reaches the page; a deck's default overlay specification stays out (an \item<+-> would
% count a step).
% The reference line's text: a capital and an ascender, as tall as a line of running text (a
% [T] column aligns its first line by the line's height, not by \baselineskip).
\newcommand{\SM@topreftext}{Xh}
\newbox\SM@topprobebox
\newif\ifSM@topprobing
\newif\ifSM@toptitled
\newcommand{\SM@toprefprobe}[1]{%
	\begingroup
	\SM@topprobingtrue
	\def\beamer@defaultospec{<*>}% beamer's own default: every step
	\def\SMColumnsAtTop{0}%
	\ifvmode\dimen@=\prevdepth\else\dimen@=\z@\fi
	\setbox\SM@topprobebox=\vbox{\prevdepth=\dimen@#1\par}%
	\xdef\SM@topref{\directlua{tex.sprint(slidemage_measure.box_content_top(\number\SM@topprobebox))}}%
	\endgroup
}
% End of the body, or anywhere a frame wants to stop the measurement
% (\progressstrip): nothing after this marker counts as first content.
\newcommand{\SM@topstop}{%
	\ifluatex
		\directlua{slidemage_measure.top_end(\SM@topthis)}%
	\fi
}
\makeatother

% Define custom frame type with en-dash bullets
\newenvironment{standardframe}[2][t]{%
	% Restart footnote numbering (the *, dagger, ... symbols) fresh on every
	% slide, rather than letting it run on cumulatively across the whole deck
	% the way beamer does by default. Two reasons: \fnsymbol only defines 9
	% symbols before erroring out, which a real multi-slide deck would run
	% through in no time if every \footnote{}/\source{} call across the whole
	% talk drew from one shared counter; and each slide's footnotes are read
	% and referenced only on that one slide anyway, so there's no reason for
	% their numbering to depend on how many footnotes earlier slides happened
	% to have. This must run *before* \begin{frame}: beamer saves the
	% counter's current value as that frame's baseline as part of
	% \begin{frame}'s own internal setup (restoring it before every overlay
	% replay, so overlays within one frame don't double-count) -- resetting
	% after \begin{frame} would be reset too late, after that baseline was
	% already captured.
	\setcounter{footnote}{0}%
	\csname SM@topprepare\endcsname
	\ifstrempty{#2}{\global\csname SM@toptitledfalse\endcsname}{\global\csname SM@toptitledtrue\endcsname}%
	\SMOpenFrame{[#1]{#2}}%
	% Forces this frame's own body to start with the correct default text
	% color, regardless of what any *previous* frame left active
	% (2026-07-16, direct instruction, root-caused via an isolated minimal
	% reproduction). \titleframe's/\quoteframe's own
	% \color{\SMTitleTextColor} calls -- set inside content built via a
	% stored box register or a zero-space overlay, used/shipped out *later*
	% rather than immediately -- don't reliably get "undone": xcolor's
	% \color doesn't emit a PDF color operator immediately when called, only
	% when the *next* real glyph/path actually needs it, so whatever color
	% one of those frame types last set can end up as the active PDF color
	% well into the *next* frame's own first draw call. Invisible for Graz
	% (the leaked color, c_primary_dark, happens to match ordinary body text
	% anyway) but catastrophic for MCI/Innsbruck (c_primary_white, rendering
	% plain non-itemize body text invisible against its own white
	% background -- itemize items escape it, since beamer's own itemize
	% template sets color fresh per bullet). An \AtBeginEnvironment{frame}
	% hook (etoolbox) was tried first and confirmed *not* to fire at all --
	% beamer's own "frame" isn't a plain \newenvironment etoolbox can patch
	% that way -- so this resets directly in standardframe's own opening
	% code instead, exactly where this project already has full control.
	\color{c_primary_dark}%
	% Defensive: the tokens after \begin{frame} are collected into the frame
	% body and replayed once per overlay slide, so this clears the side-
	% footnote collector at the start of every slide's pass -- normally a
	% no-op (the previous slide's footline already consumed the box), but it
	% stops stale footnotes leaking in if a previous frame ever collected
	% without shipping a footline.
	\SMClearSideFootnotesBox%
	\SMResetImageSources%
	% The itemize bullets are the theme's (\SMBulletDot, beamerthemeSlideMage.sty).
	% Only for top-aligned frames (the default, #1=t): a [c]-centered frame
	% (2026-07-22, direct instruction, found while centering
	% slides/bw9-Pearson-Correlation-Illustration-current.tex/agc/crj) was
	% landing its content almost flush against the title bar instead of
	% truly centered -- root cause was this same pull-up firing
	% unconditionally regardless of #1. beamer's own [c] option already
	% centers the frame body within the available space; stacking this
	% deck's *own* top-alignment pull-up on top of that just drags the
	% already-centered box further upward by the same fixed amount, which
	% is exactly the asymmetry that was measured (a [c] figureframe's image
	% landing ~0pt from the title bar and ~10pt from the footer instead of
	% symmetric). Guarded with \ifstrequal (etoolbox, already loaded above)
	% rather than \ifx, matching this file's own \ifstrequal{##1}{...}
	% pattern elsewhere (\ifx would need #1 fully expanded first).
	\ifstrequal{#1}{c}{}{%
		\vspace*{-\dimexpr\SMContentTopShift-\SMContentTopShiftReduction+\csname SM@topcorr\endcsname\relax}%
		\csname SM@topstart\endcsname}%
}{%
	\csname SM@topstop\endcsname
	% Must run *before* \end{frame}: like any other frame content, a
	% \footnote{}/\source{} call has to happen while the frame body is
	% still open -- everything in this closing block from \end{frame}
	% onward runs *after* the frame has already closed (see README.md's
	% own \titleframe comment on beamer's \beamer@collect@@body for the
	% general shape of this constraint).
	\SMEmitImageSourceFootnote%
	\end{frame}%
	\global\def\SMChapterLogoPath{}% opt-in only for the frame just closed
	\global\def\SMContentTopShiftReduction{0pt}% opt-in only for the frame just closed
	\global\csname SM@topmanualfalse\endcsname% opt-in only for the frame just closed
	\global\csname SM@topsectionfalse\endcsname% set by sectionframe for the frame just closed
	\global\SMSideFootnotesfalse% opt-in only for the frame just closed
	\global\def\SMFrametitleFullHeightYellow{}% opt-in only for the frame just closed
	\global\def\SMFrametitleLogoFixed{}% opt-in only for the frame just closed
	\global\def\SMFooterBoxColor{\SMFooterBoxColorDefault}% opt-in only for the frame just closed
}


% "a)"/"b)"/"c)"-labeled list, for exercise sub-tasks that read better
% lettered than bulleted (2026-07-20, direct instruction, first used on
% slides/y3a-Exercise-Diversification-2-current.tex's own bottom three
% items).
%
% First attempt (superseded, kept as a documented wrong turn -- see
% README.md's own "lualatex vs latex" section for why this project keeps
% wrong write-ups rather than deleting them): reuse itemize's own top-level
% `\llap`-based label positioning by locally overriding
% `\setbeamertemplate{itemize item}` inside a plain `\begin{itemize}` --
% the reasoning being that `\llap` draws its content flush against the
% text's own left edge "regardless of the label's width." That reasoning
% turned out to be wrong in a way that only showed up once actually
% measured against the compiled PDF (not just visually eyeballed): a plain
% bullet's own `\llap`-drawn footprint (title/Left-Margin-relative-left-
% edge to text-start) measured only ~9.05pt end-to-end on this deck's real
% slides, matching the Left-Margin-to-text-start gap exactly -- but "a)"
% rendered through the *exact same* mechanism measured a ~13.97pt
% footprint for the same gap, not the ~8.5-9pt `\settowidth` reported for
% the bare glyphs in isolation. The extra ~5pt was never fully root-caused
% (this custom itemize's own `\list` call, further below, only overrides
% `\makelabel` -- it never sets `\leftmargin`/`\labelwidth` itself, so
% those are left at whatever this deck's global itemize defaults happen to
% be, and `\llap`'s actual box-width behavior inside that inherited,
% unexamined list-parameter state is where the extra width most likely
% comes from) -- rather than keep guessing at *why*, measured the actual
% failure directly (the compiled "a)" sat ~4.9pt left of the Left Margin,
% confirmed via exact PDF text-position extraction, not a visual guess)
% and abandoned the approach once even `\tiny` plus manual per-letter
% kerning only just barely closed the gap while reading as illegibly small
% -- a real dead end, not a tuning problem.
%
% Working fix: don't reuse itemize's `\llap` label mechanism at all -- use
% a real, width-reserving `\makebox[W][r]{...}` label instead (the same
% *kind* of mechanism `\SMEnumerateLabelWidth` above already uses for
% a nested enumerate, just with an explicit `\leftmargin`/`\labelwidth`/
% `\labelsep` this environment sets itself rather than inheriting), inside
% its own plain `\list`, not a reused `\itemize`. `\leftmargin` (label box
% + gap) is set to 11.1pt -- 9.1pt label box (fits "a)"/"b)"/"c)" at this
% deck's normal `\small` itemize size with a little clearance, confirmed
% directly: measured label-left-edge stays right of the 12.552pt Left
% Margin for all three letters, closest case (widest letter, "c)") 1.37pt
% clear) plus 2pt `\labelsep` for a visible gap before the text. Trade-off,
% accepted deliberately rather than found by accident: item text lands
% ~2.2pt right of where a plain bulleted item's text would (23.8pt vs.
% 21.6pt) -- fitting a legible "a)"/"b)"/"c)" at normal reading size plus
% a real gap simply doesn't fit in the ~9pt a plain bullet's tiny dot
% leaves, so *exact* text-position parity with plain bullets and a legible
% label are mutually exclusive here; legibility won given the two-point
% shift is visually negligible and there's no plain-bulleted item left on
% the same slide to compare against directly anyway (y3a's own first item
% moved into its infoblock). Re-derive `\leftmargin`/`\labelwidth` the same
% way (render, measure with exact text-position extraction, adjust) if
% this is ever used at a different font size than this deck's normal
% itemize `\small`.
% The two lengths of every list with "a)" / "1." / "1." / "a." labels (v9.8.0: derived; they were
% 9.1pt and 4.75pt in the Fira profile and 9.16pt and 5.20pt in the lmsans one). \SMLabelItemWidth
% is the box that holds the widest label of the first items, a)-c) and 1)-9) with a bracket or a
% period, in the font in force where the list begins, plus 0.1pt of clearance; \SMLabelItemSep is the
% hook \SMLabelItemGap (in em of that font). Called first in each list's setup, so that
% the lists follow \small, \footnotesize and \shrinklists, and once at \begin{document}
% (in \small) for a slide that reads the two lengths itself.
\makeatletter
\newcommand{\SM@labelitemmetrics}{%
	\setlength{\SMLabelItemWidth}{0pt}%
	\@for\SM@labelprobe:=a),b),c),a.,b.,c.,1),2),3),4),5),6),7),8),9),1.,2.,3.,4.,5.,6.,7.,8.,9.\do{%
		\settowidth{\dimen@}{\SM@labelprobe}%
		\ifdim\dimen@>\SMLabelItemWidth\SMLabelItemWidth=\dimen@\fi}%
	\addtolength{\SMLabelItemWidth}{0.1pt}%
	\setlength{\SMLabelItemSep}{\SMLabelItemGap}}
\makeatother
\newcounter{SMAlphaItemCounter}
\newcommand{\SMAlphaItem}{\stepcounter{SMAlphaItemCounter}\alph{SMAlphaItemCounter})}
\newenvironment{SMAlphaItemize}{%
	\setcounter{SMAlphaItemCounter}{0}%
	\list{\SMAlphaItem}{%
		% 4.75pt (2026-07-20, direct instruction: "half the width of the
		% 'C' in 'Calculate'" -- \settowidth{C} at this deck's normal
		% itemize \small measured 5.5pt, so 2pt original + 2.75pt half-C).
		% Only \labelsep changes the text's own start position -- the
		% label's own box stays exactly where it was (its left edge is
		% pinned to this list's origin regardless of \labelsep, since
		% \leftmargin=\labelwidth+\labelsep always cancels \labelsep back
		% out of the label box's own position), so this widens the gap
		% without reopening the Left-Margin question already solved above.
		\csname SM@labelitemmetrics\endcsname\setlength{\labelsep}{\SMLabelItemSep}%
		\setlength{\labelwidth}{\SMLabelItemWidth}%
		\setlength{\leftmargin}{\labelwidth}%
		\addtolength{\leftmargin}{\labelsep}%
		\setlength{\itemindent}{0pt}%
		% \itemsep/\parsep left at plain \list's own defaults otherwise --
		% since this environment is built directly on \list rather than
		% reusing beamer's own itemize template, it never got beamer's own
		% itemize spacing hooks at all, only whatever \list falls back to.
		% Measured directly (2026-07-21, direct instruction, found on
		% slides/nt4-Exercise-CAPM-5-current.tex): that default reads as
		% *less* space between two items (13.7pt, item-to-item) than
		% between two wrapped lines within one *single* item (15.6pt,
		% pure \baselineskip) -- backwards from how a list should read
		% (item boundaries should be at least as open as a same-item line
		% wrap, not tighter). \itemsep set explicitly so the total gap at
		% an item boundary (baselineskip + itemsep) reads as clearly more
		% open than a plain wrapped line, not accidentally tighter.
		% Reduced from the original 6pt to 2.5pt (2026-07-22, found during
		% a full-project check after the itemize-rules work, not part of
		% that work itself): 6pt overshot by more than needed to clear the
		% 1.9pt gap (15.6 vs. 13.7), and slides/kmf-Exercise-CAPM-6-current.tex's
		% own 5-item list -- longer than nt4's own -- compounded that
		% excess across 4 item boundaries into a ~25pt overfull frame, its
		% own last item cut off by the footer (confirmed directly,
		% rendered and compared against a pre-regression PDF). 2.5pt
		% keeps a real, visible margin over the 1.9pt needed (not just
		% barely crossing it) while re-measuring clean (no overfull) on
		% kmf; re-verify against both nt4 and kmf directly if this is
		% ever changed again -- they're this environment's own tightest
		% (most items) and originally-tuned-against cases respectively.
		\setlength{\itemsep}{\SMLabelItemItemsep}%
	}
	\raggedright% plain \list defaults to justified+hyphenated text, unlike
	% this deck's own itemize (which calls this too, further below) --
	% missed on the first pass, only caught by comparing a rendered a)/b)/c)
	% item side by side against a plain bulleted one ("stan-\ndard deviation"
	% hyphenating where a plain-bulleted version of the same sentence
	% elsewhere in this deck doesn't).
}{%
	\endlist
}

% Three siblings of \SMAlphaItemize above, same measured \list metrics
% (labelsep/labelwidth/leftmargin/itemindent/itemsep), only the label text
% itself differs -- added 2026-08-09 for slides whose numbering has to
% match a real external scheme (statute paragraph numbers, PowerPoint's
% own "1./2." or "a./b." auto-numbering) rather than plain letters-with-
% parenthesis. Not merged into one parameterized environment: each name is
% called directly from a slide file, and duplicating six short \setlength
% lines three times is cheaper to read at the call site than a macro
% layer over \newenvironment for four near-identical bodies.
\newcounter{SMNumItemCounter}
\newcommand{\SMNumItem}{\stepcounter{SMNumItemCounter}\arabic{SMNumItemCounter})}
\newenvironment{SMNumItemize}{%
	\setcounter{SMNumItemCounter}{0}%
	\list{\SMNumItem}{%
		\csname SM@labelitemmetrics\endcsname\setlength{\labelsep}{\SMLabelItemSep}%
		\setlength{\labelwidth}{\SMLabelItemWidth}%
		\setlength{\leftmargin}{\labelwidth}%
		\addtolength{\leftmargin}{\labelsep}%
		\setlength{\itemindent}{0pt}%
		\setlength{\itemsep}{\SMLabelItemItemsep}%
	}
	\raggedright%
}{%
	\endlist
}

\newcounter{SMAlphaPeriodItemCounter}
\newcommand{\SMAlphaPeriodItem}{\stepcounter{SMAlphaPeriodItemCounter}\alph{SMAlphaPeriodItemCounter}.}
\newenvironment{SMAlphaPeriodItemize}{%
	\setcounter{SMAlphaPeriodItemCounter}{0}%
	\list{\SMAlphaPeriodItem}{%
		\csname SM@labelitemmetrics\endcsname\setlength{\labelsep}{\SMLabelItemSep}%
		\setlength{\labelwidth}{\SMLabelItemWidth}%
		\setlength{\leftmargin}{\labelwidth}%
		\addtolength{\leftmargin}{\labelsep}%
		\setlength{\itemindent}{0pt}%
		\setlength{\itemsep}{\SMLabelItemItemsep}%
	}
	\raggedright%
}{%
	\endlist
}

\newcounter{SMNumPeriodItemCounter}
\newcommand{\SMNumPeriodItem}{\stepcounter{SMNumPeriodItemCounter}\arabic{SMNumPeriodItemCounter}.}
\newenvironment{SMNumPeriodItemize}{%
	\setcounter{SMNumPeriodItemCounter}{0}%
	\list{\SMNumPeriodItem}{%
		\csname SM@labelitemmetrics\endcsname\setlength{\labelsep}{\SMLabelItemSep}%
		\setlength{\labelwidth}{\SMLabelItemWidth}%
		\setlength{\leftmargin}{\labelwidth}%
		\addtolength{\leftmargin}{\labelsep}%
		\setlength{\itemindent}{0pt}%
		\setlength{\itemsep}{\SMLabelItemItemsep}%
	}
	\raggedright%
}{%
	\endlist
}

% Itemize indentation: level 1 bullets should sit flush with the frame title
% (no extra left margin), and each deeper level should be indented just
% enough for its bullet to line up under the *text* (not the bullet) of the
% level above. A nested list's \leftmarginN is added on top of its parent's
% (indentation is cumulative), so giving every level the same "step" --
% label width + labelsep, the width a bullet plus its gap to text takes up --
% makes level N's text sit at N*step and level N's bullet at (N-1)*step, i.e.
% exactly at level (N-1)'s text (0 for level 1, flush with the title).
\makeatletter
\newlength{\SMBulletWidth}
\settowidth{\SMBulletWidth}{\SMBulletDot}
\newlength{\SMItemStep}
\setlength{\SMItemStep}{\labelsep}
\addtolength{\SMItemStep}{\SMBulletWidth}
\setlength{\leftmargini}{\SMItemStep}
\setlength{\leftmarginii}{\SMItemStep}
\setlength{\leftmarginiii}{\SMItemStep}
\setlength{\leftmarginiv}{\SMItemStep}
\makeatother

% \item[\SMBulletEllipsis] -- an ellipsis used as the item label (the source's
% "...how you contribute" continuation items). A plain \item[\ldots] right-aligns the
% label to the bullet dot's right edge, so the wider ellipsis pokes out to the left
% of where a bullet starts; this one starts at the dot's left edge instead
% (2026-09-20, direct instruction: the ellipses were too far left).
\newcommand{\SMBulletEllipsis}{\hbox to \SMBulletWidth{\rlap{\ldots}\hss}}

% \setbeamertemplate{itemize subitem}{\SMBulletEllipsisScript} -- the same for a list whose *bullet itself* is a small ellipsis (the source's
% "..." sub-points; slides quy, 3a8, unt). The bare \scriptsize… bullet used to sit 1.6pt left of where a dot starts (2026-09-20).
\newcommand{\SMBulletEllipsisScript}{\hbox to \SMBulletWidth{\rlap{\scriptsize…}\hss}}

% A 4th itemize level -- beamer's own \itemize (beamerbaselocalstructure.sty)
% hard-caps nesting at 3 levels (\ifnum \@itemdepth >2\relax\@toodeep),
% independently of anything this deck sets: a 4th \begin{itemize} always
% raised "! LaTeX Error: Too deeply nested." even though plain LaTeX's own
% kernel itemize supports 4. Found 2026-07-15, direct instruction, on
% slides/hrg-Exchange-Rate-Quotation-Forward-current.tex -- genuinely
% 4 levels deep in the source pptx itself (confirmed via a PowerPoint COM
% render, not assumed), not a conversion artifact to design around.
% Since this deck already styles every existing level identically (same
% \SMBulletDot, same \SMItemStep margin -- see the comments right
% above/below), extending to a 4th level needs no new visual design, just
% patching beamer's own depth limit up by one and telling it what to call
% the new level ("subsubsub", the next name in beamer's own
% sub/subsub/... naming scheme) so its template/color lookups
% (\usebeamertemplate{itemize subsubsubitem}, \usebeamercolor{itemize/
% enumerate subsubsub body}, etc.) resolve to *something* instead of
% erroring -- undefined beamer templates/colors are expected to fall back
% to sensible defaults (no extra spacing hook, plain text color) rather
% than raising an error, same as any other beamer template/color that a
% given theme never bothers to set; the actual bullet symbol still needs
% its own \SMBulletDot template (in the theme), added explicitly the same way
% as the existing 3 levels, not left to whatever fallback beamer has for
% an unset "itemize subsubsubitem" template.
% \beamer@computepref is patched wholesale (not via \patchcmd) since its
% body is one single \ifcase chain, not a stable string to match against.
% \itemize itself was *first* tried via \patchcmd{\itemize}{...}
% (matching this deck's own established practice for beamer/theme
% internals, e.g. beamerthemeSlideMage.sty's own
% \patchcmd{\beamer@title}{...}) -- this failed outright, confirmed
% directly (\patchcmd's own failure branch fired, raising a clean
% "Patching failed" error rather than silently doing nothing): recent
% LaTeX2e kernels implement an optional-argument environment like
% beamer's own `\renewcommand{\itemize}[1][]{...}` via
% `\@protected@testopt`, which stores the *real* body not under `\itemize`
% itself but under a differently-named internal control sequence (`\show
% \itemize` confirms this directly: it prints as
% `\@protected@testopt \itemize \\itemize {}`, i.e. plain `\itemize` is
% only a thin dispatcher, and the actual `\ifnum \@itemdepth >2...` body
% \patchcmd was trying to match against was never there to find). Rather
% than chase that internal storage name (liable to differ across LaTeX
% kernel versions, which \patchcmd's own documented usage doesn't need to
% care about), \itemize is fully re-declared here instead, copied
% verbatim from beamer's own beamerbaselocalstructure.sty (this project's
% installed version) with only the depth bound itself changed from 2 to
% 3 -- the same declaration form beamer itself used
% (\renewcommand{\itemize}[1][]{...}), which regenerates both the
% dispatcher and the real body correctly under whatever internal name
% the current kernel expects, sidestepping the naming question entirely.
% Trade-off, stated plainly: if a future beamer release changes this
% macro's own body for an unrelated reason, this override will silently
% keep using today's (2026-07) copy instead of picking up that change --
% acceptable here since the override is one line different from upstream
% and clearly marked, but re-diff against beamerbaselocalstructure.sty's
% own \itemize if beamer is ever upgraded and itemize behavior looks off.
% \SMNoListParsep (2026-10-01, direct instruction: "Implement a general
% fix"): a level-1 list typeset at \small or \footnotesize -- an ambient size
% switch before it, \setbeamerfont, \shrinklists -- loads LaTeX's
% size-specific \@listi, which brings \parsep 3pt plus 2pt (\footnotesize:
% similar) and \parskip with it; a normal-size list has neither. So a smaller
% list's items sat *further* apart than at normal size (6j4: 19.9pt per item
% at normal size, 20.8pt at \small), against this deck's rule that list
% spacing is em-based and shrinks with the font (\SMItemsepRatioI etc.).
% Removed wherever \parsep has a natural (non-stretch) part, i.e. exactly in
% those lists; deeper levels (0pt plus 1pt) are left as they were.
\makeatletter
\newcommand{\SMNoListParsep}{%
	\ifdim\parsep>\z@
		\setlength{\parsep}{\z@}\setlength{\parskip}{\z@}%
	\fi
}
\long\def\beamer@computepref#1{%
	\let\beameritemnestingprefix\@empty%
	\ifcase#1\or\or\def\beameritemnestingprefix{sub}\or\def\beameritemnestingprefix{subsub}\or\def\beameritemnestingprefix{subsubsub}\or\@toodeep\fi%
}
\renewcommand{\itemize}[1][]{%
	\ifblank{#1}{}{\def\beamer@defaultospec{#1}}%
	\ifnum \@itemdepth >3\relax\@toodeep\else
		\advance\@itemdepth\@ne
		\beamer@computepref\@itemdepth% sets \beameritemnestingprefix
		\usebeamerfont{itemize/enumerate \beameritemnestingprefix body}%
		\usebeamercolor[fg]{itemize/enumerate \beameritemnestingprefix body}%
		\usebeamertemplate{itemize/enumerate \beameritemnestingprefix body begin}%
		\list
			{\usebeamertemplate{itemize \beameritemnestingprefix item}}
			{\def\makelabel##1{%
					{%
						\hss\llap{{%
								\usebeamerfont*{itemize \beameritemnestingprefix item}%
								\usebeamercolor[fg]{itemize \beameritemnestingprefix item}##1}}%
					}%
				}%
			}
	\fi%
	\beamer@cramped%
	% \beamer@cramped (beamerbaselocalstructure.sty, unmodified upstream
	% code) unconditionally does `\itemsep0pt\topsep0.5pt`, every single
	% time -- called right here, in this project's own copy of beamer's
	% own \itemize body, same position as upstream. Setting \itemsep any
	% earlier (first tried as a separate \AtBeginEnvironment{itemize} hook,
	% run *before* this whole macro body even starts) gets silently wiped
	% out by this same line a moment later -- confirmed directly via
	% \typeout debugging (not assumed): the hook's own \setlength really
	% did run with the right values, \itemsep still read back as 0pt
	% immediately after. Setting it *here*, right after \beamer@cramped
	% already ran (not before), is what actually survives.
	\ifcase\@itemdepth
	\or\setlength{\itemsep}{\SMItemsepRatioI em}% level 1
	\or\setlength{\itemsep}{\SMItemsepRatioII em}% level 2
	\or\setlength{\itemsep}{\SMItemsepRatioIII em}% level 3
	\or\setlength{\itemsep}{\SMItemsepRatioIV em}% level 4
	\fi%
	\SMNoListParsep
	\raggedright%
	\beamer@firstlineitemizeunskip%
}
\makeatother

% A nested \begin{enumerate} inherits the same \leftmarginN step as itemize
% above for *where its text starts* (matching itemize's own "one step past
% the parent level's text" rule -- confirmed directly against the compiled
% PDF: an enumerate's own body text, e.g. "Senior secured debt", already
% lands exactly one \SMItemStep past its parent item's text), but its
% *label* ("1.", "2.", ...) doesn't follow \labelwidth the way plain LaTeX's
% list mechanism would -- beamer draws enumerate labels through its own
% template machinery, which turned out not to respect \labelwidth at all:
% \setlength{\labelwidth}{...} right after \begin{enumerate} measurably took
% effect (checked with \typeout{\the\labelwidth} immediately after) but the
% compiled numeral's own position never moved. What *does* move it: directly
% overriding \makelabel with a fixed-width \makebox -- empirically
% calibrated (render, measure the numeral's own x-position in the compiled
% PDF, adjust, repeat) since the box's own left edge turned out to sit at a
% fixed position unrelated to \labelwidth or \leftmarginN, with the
% right-aligned numeral inside it moving in lockstep with the box's width
% (confirmed linear: +1pt width -> +~1pt numeral position, across several
% test widths). 19.23pt (at this deck's default itemize font size) lands a
% single-digit numeral flush with the parent item's own text -- only valid
% at that font size and for single-digit lists (1-9); re-calibrate the same
% way (not by recomputing from first principles) for a differently-sized or
% 10+-item nested enumerate. Call right after \begin{enumerate}, before the
% first \item.
%
% 20.33pt (2026-07-08, direct instruction: "the numbers still extend a
% little farther to the left than the text above them does" --
% slides/f6m-Recovery-Rates-current.tex's own "Pecking order:" list).
% 19.23pt landed *close* but not quite flush: measured directly against
% the compiled PDF, every numeral's own left edge (20.55pt-20.95pt,
% varying slightly by digit width even though all are right-aligned to
% the same box) still sat up to ~1.05pt left of "Pecking order:"'s own
% text start (21.6pt) -- the widest digit tried, "6.", was the binding
% case. +1.1pt of box width (confirmed +1pt width -> +1pt numeral
% position, per this comment's own note above) covers that with a small
% margin, re-confirmed against the recompiled PDF.
\newcommand{\SMEnumerateLabelWidth}{\renewcommand{\makelabel}[1]{\makebox[\SMEnumerateLabelBox][r]{##1}}}

% Automatic version of the above (2026-07-23, background-task follow-up
% flagged in slides/kfh-Asset-Allocation-Process-current.tex's own v1-1-0
% comment: "fix the enumerate behavior generally, if possible", task
% `task_bb3fbb66`). That comment's first attempt, an
% `\AtBeginEnvironment{enumerate}{...}` hook, failed because the hook fires
% *before* `\enumerate`'s own body runs, and that body's own internal
% `\list{...}{\def\makelabel...}` setup (the beamer copy directly below the
% `\itemize` override above) clobbers whatever a hook that early leaves
% behind -- the exact same "template runs its own setting, external pre-set
% gets overwritten" problem `\itemize`'s own `\itemsep` comment above
% already documents. The real fix, same idea `\itemize` itself used above:
% override beamer's own low-level piece directly rather than hooking
% \AtBeginEnvironment. Unlike `\itemize` (declared via
% `\@protected@testopt`, needing a full re-declaration to reach -- see the
% `\itemize` override's own comment for the full story), `\enumerate`'s
% relevant piece, `\beamer@enum@` (beamerbaselocalstructure.sty), is a
% plain `\def` with no optional-argument indirection, so it can be
% `\renewcommand`'d directly, verbatim from upstream except for the
% `\ifnum\@itemdepth>1` branch added below.
%
% Only fires for a *nested* enumerate (`\@itemdepth>1`, i.e. already inside
% at least one `\itemize`/`\enumerate` when this one opens) -- a top-level
% enumerate (`\@itemdepth=1`, e.g.
% slides/t9z-Outland-Steel-Hedging-Borrowing-Investing-current.tex/
% slides/g4c-Outland-Steel-Hedging-Forward-Sale-current.tex/
% slides/j84-Outland-Steel-No-Hedging-current.tex's own outermost numbered
% list, or slides/nkj-Agenda-current.tex's own \tableofcontents-driven one)
% is left at beamer's stock label position entirely untouched, matching
% this project's own established behavior for that shape (never reported
% as misaligned, unlike the nested case) -- and, for nkj specifically,
% doesn't matter either way, since that slide's own
% `\SMTocLabelFlushLeft` call (right after its own `\begin{enumerate}`)
% overrides `\makelabel` again regardless of whatever this default sets.
%
% Reuses the exact same 20.33pt calibration as the manual
% `\SMEnumerateLabelWidth` above, confirmed (pymupdf, compiled PDF) to
% land correctly not just at `\@itemdepth=2` (one `\itemize` above the
% enumerate -- slides/kfh-Asset-Allocation-Process-current.tex/
% slides/f6m-Recovery-Rates-current.tex/
% slides/6pz-Outland-Steel-Hedging-Strategies-Overview-current.tex, all
% previously needing the manual call) but also at `\@itemdepth=3` (two
% `\itemize`s above it -- slides/bdh-Working-Capital-Example-current.tex's
% own "Development of balance sheet:" list, also previously needing the
% manual call, with the identical 20.33pt value) -- this deck's itemize/
% enumerate templates don't shrink font size by nesting depth (only the
% bullet glyph/margin step does), so the numeral's own glyph width -- the
% only thing \SMEnumerateLabelWidth's own box width actually
% compensates for -- stays the same at every depth, and one constant covers
% both cases without needing a depth-dependent formula. Not verified beyond
% depth 3 (no slide currently nests an enumerate that deep); if one ever
% does, re-measure before trusting this constant that far down, same as any
% other empirical calibration in this file.
% (slides/scs-Political-Risk-Stability-Index-current.tex also matches
% "`\SMEnumerateLabelWidth`" in a grep -- but only inside its own prose
% comments describing the general technique; its actual call is the
% bespoke `\SMEnumerateLabelFlushLeftShifted{16pt}{10.2pt}` for its own
% alphabetic, left-aligned, narrower-box labels, unaffected by this
% override either way since that call runs after `\begin{enumerate}` opens,
% same ordering as every other post-open `\makelabel` override here.)
%
% Former per-slide manual `\SMEnumerateLabelWidth` calls removed from
% the 4 slides above now that this covers them automatically; the macro
% itself is kept defined, both for this override to reuse and as a manual
% escape hatch (e.g. a differently-calibrated width some future slide might
% need).
% (v9.8.0: \SMEnumerateTextExtra, 6.3pt, is a hook of the catalog; the four lengths below were numbers of the
% layout profile: 7.85pt, 13.1pt, 5.25pt and 20.33pt for Fira Sans, 8.18pt, 13.43pt, 5.25pt and 20.33pt for
% Latin Modern Sans.)
% The numeral column of a nested enumerate, from the widest numeral of the list's own font (the digits
% 1-9 with the dot, which is what the item template sets). The slot a numeral sits in is that width plus
% \SMEnumerateNestedSlotExtra, the text starts \SMEnumerateNestedLabelGap after it
% (\SMEnumerateNestedLeftMargin is both), and the box the numeral is right-aligned in only has to hold a
% numeral of three digits ("999."), so it takes the width of that. Called first in the list's setup of
% every enumerate, and once at \begin{document} (in \small) for a slide that reads the lengths itself.
\makeatletter
\newcommand{\SM@enumeratemetrics}{%
	\setlength{\SMEnumerateNestedMinNumber}{0pt}%
	\setlength{\SMEnumerateLabelBox}{0pt}%
	\@for\SM@numprobe:=1,2,3,4,5,6,7,8,9\do{%
		\settowidth{\dimen@}{\SM@numprobe.}%
		\ifdim\dimen@>\SMEnumerateNestedMinNumber\SMEnumerateNestedMinNumber=\dimen@\fi
		\settowidth{\dimen@}{\SM@numprobe\SM@numprobe\SM@numprobe.}%
		\ifdim\dimen@>\SMEnumerateLabelBox\SMEnumerateLabelBox=\dimen@\fi}%
	\addtolength{\SMEnumerateNestedMinNumber}{\SMEnumerateNestedSlotExtra}%
	\setlength{\SMEnumerateNestedLabelSep}{\SMEnumerateNestedLabelGap}%
	\setlength{\SMEnumerateNestedLeftMargin}{\dimexpr\SMEnumerateNestedMinNumber+\SMEnumerateNestedLabelSep\relax}}
\makeatother
\makeatletter
\renewcommand{\beamer@enum@}{%
	\beamer@computepref\@itemdepth% sets \beameritemnestingprefix
	\usebeamerfont{itemize/enumerate \beameritemnestingprefix body}%
	\usebeamercolor[fg]{itemize/enumerate \beameritemnestingprefix body}%
	\usebeamertemplate{itemize/enumerate \beameritemnestingprefix body begin}%
	% Level-1 enumerate text sits further right than bullet text, as in the source
	% (numeral at the bullet dot's left edge, item text 27.9pt from the page's left edge
	% against 21.6pt for a bullet's text; 2026-09-20). Nested lists follow it, as they do for
	% plain bullets: a level-2 bullet starts under the level-1 text (2026-09-21, direct
	% instruction -- an earlier version gave the extra back to the nested list, which left
	% level-2 bullets 6.3pt left of the level-1 text above them). Local to this list.
	\ifnum\@itemdepth=1
		\addtolength{\leftmargini}{\SMEnumerateTextExtra}%
	\fi
	\expandafter
		\list
			{\usebeamertemplate{\beamer@enumtempl}}
			{\usecounter\@enumctr%
				\SM@enumeratemetrics% the numeral column from this list's font (v9.8.0)
				\ifnum\@itemdepth>1
					% Nested numerals: same rule as level 1 -- left edge at the left
					% edge of this level's bullet dot (2026-09-20). The box below is
					% the old 20.33pt right-aligned one, whose right end the numerals
					% used to fill flush; giving the numeral a 7.85pt left-aligned
					% slot at that end puts every numeral's left edge exactly at the
					% dot's, instead of up to 1pt right of it for narrow digits
					% ("1." vs "4."). A wider numeral keeps the old right alignment.
					% Hanging indent (2026-09-21, direct instruction: the second line of
					% an item did not start under its first line's text). The 20.33pt label
					% box is wider than beamer's \labelwidth, so it pushed the first
					% line's text 4.3pt right of \leftmargin, where every later line
					% starts. Sizing \labelwidth to the box and \leftmargin to
					% "numeral slot + \labelsep" puts the numeral's left edge at the dot's
					% and the text of every line at the same x (13.1pt from the parent's
					% text, i.e. where the first line always was).
					\setlength{\labelwidth}{\SMEnumerateLabelBox}%
					\setlength{\labelsep}{\SMEnumerateNestedLabelSep}%
					\setlength{\leftmargin}{\SMEnumerateNestedLeftMargin}%
					\def\makelabel##1{\makebox[\SMEnumerateLabelBox][r]{{%
							\setbox\@tempboxa\hbox{##1}%
							\ifdim\wd\@tempboxa>\SMEnumerateNestedMinNumber
								\box\@tempboxa
							\else
								\hbox to \SMEnumerateNestedMinNumber{\unhbox\@tempboxa\hss}%
							\fi}}}%
				\else
					% Level-1 numerals start at the left edge of the level-1 bullet
					% dot, like the source's numerals (2026-09-20, direct
					% instruction: the numbers were too far left). They used to
					% be right-aligned to the dot's right edge, which for a
					% numeral wider than the dot (every "1.") puts its left edge
					% left of where a bullet would start. A numeral too wide to
					% keep 1pt clear of the item text ("10.") keeps the old
					% right alignment instead of running into the text.
					\def\makelabel##1{{\hss\llap{{%
							\usebeamerfont*{enumerate \beameritemnestingprefix item}%
							\usebeamercolor[fg]{enumerate \beameritemnestingprefix item}%
							\setbox\@tempboxa\hbox{##1}%
							\ifdim\dimexpr\wd\@tempboxa+1pt\relax>\dimexpr\SMItemStep+\SMEnumerateTextExtra\relax
								\box\@tempboxa
							\else
								\hbox to \dimexpr\SMBulletWidth+\SMEnumerateTextExtra\relax{\rlap{\unhbox\@tempboxa}\hss}%
							\fi}}}}%
				\fi}%
	\beamer@cramped%
	\SMNoListParsep
	\raggedright%
	\beamer@firstlineitemizeunskip%
}
% The lists' lengths as a slide that reads them outside a list sees them: those of the \small font.
\AtBeginDocument{%
	\begingroup\small\SM@labelitemmetrics\SM@enumeratemetrics
	\xdef\SM@listmetrics{%
		\global\SMLabelItemWidth=\the\SMLabelItemWidth\relax
		\global\SMLabelItemSep=\the\SMLabelItemSep\relax
		\global\SMEnumerateNestedMinNumber=\the\SMEnumerateNestedMinNumber\relax
		\global\SMEnumerateNestedLabelSep=\the\SMEnumerateNestedLabelSep\relax
		\global\SMEnumerateNestedLeftMargin=\the\SMEnumerateNestedLeftMargin\relax
		\global\SMEnumerateLabelBox=\the\SMEnumerateLabelBox\relax}%
	\expandafter\endgroup\SM@listmetrics}
\makeatother

% Extra breathing room when a nested itemize closes and control returns to
% its *parent* list, right before that parent's next \item -- e.g. sub-items
% under "Mid-Term Exam" followed by the next top-level "Final Exam" -- so a
% new top-level group reads as visually distinct from the one before it,
% rather than the same uniform gap used between every other item regardless
% of level. \@itemdepth (LaTeX's own nesting-level counter, already
% decremented back to the parent's level by the time \end{itemize}'s normal
% closing has run) says which level control returned to -- 0 means the
% outermost list itself just closed (nothing left to add space before,
% deliberately excluded), 1 means back to a level-1 item, 2 means back to a
% level-2 item.
%
% Level-dependent (2026-07-06, direct instruction) -- a flat, single
% constant regardless of *which* level control returned to (the original
% version of this hook) reads as "extra"/inconsistent precisely because a
% return-to-level-1 and a return-to-level-2 transition both got the exact
% same gap, rather than level 1 getting a visibly bigger gap than level 2 --
% found on slides/57r-Reasons-for-Leasing-current.tex, where a level-2
% sub-list closing back to a level-1 item looked like a random extra gap
% rather than "level 1 items are consistently more separated than level 2
% ones." \addvspace (not plain \vspace) matters too: it coalesces with
% whatever glue itemize's own \itemsep/\topsep already contributes at the
% closing boundary (taking whichever is larger) rather than *adding* on top
% of it unconditionally -- \vspace's own stacking behavior was the other
% half of why the old flat version could look like too much space.
\makeatletter
% Ratios, not lengths (2026-07-21, direct instruction: "all of these
% spacings should scale with font size") -- a \newlength+\setlength{...}{0.5em}
% pair, the way this used to be written, freezes to a fixed *pt* value the
% one time \setlength runs (at preamble-load time, using whatever font is
% active then), despite the "em" in its own definition -- it does NOT
% re-track whichever font size is actually active on a given slide/level
% later. Confirmed directly: a \typeout{\f@size} placed at the point each
% of these actually gets used (inside \AtEndEnvironment{itemize}'s own
% hook body, further below) showed the *closing* list's own font still
% active there (9pt/10pt/10.95pt for a level-4/3/2 list closing,
% respectively -- not yet restored to the parent's own font), confirming
% this hook fires early enough in `\end{itemize}`'s own processing that a
% live "0.5em" computed *there* would correctly track whichever level is
% actually involved. Fixed by storing the bare ratio as a plain macro
% instead, and writing the literal "em" back on at the point of use
% (`\addvspace{\SMItemizeLevelDropSpaceRatioI em}`), so it's
% recomputed fresh, against whatever font is live, every time.
% Level 3's own drop-space, for when a closed *4th*-level list returns
% control to a level-3 item -- added 2026-07-15 alongside the 4th itemize
% level itself (search "subsubsub" further above). Interpolated from the
% existing I/II progression (0.5, 0.3 -- each step roughly 60% of the one
% before) rather than empirically re-measured against a compiled PDF the
% way I/II themselves were: no slide used a 4th level before
% slides/hrg-Exchange-Rate-Quotation-Forward-current.tex to measure
% against. Re-derive properly (render, measure, adjust) if this ever
% looks visibly off on a real slide.
% Level 4's own drop-space, for when a closed level-4 list returns control
% to a level-3 item -- added 2026-07-21 alongside \SMItemsepRatioIV.
% \@itemdepth at this hook's fire time is the depth of the list that *just
% closed*, confirmed directly via a temporary \typeout{\the\@itemdepth}
% (not assumed): closing the level-4 list on
% slides/ZZO-Itemize-Spacing-Reference-current.tex measured itemdepth=4
% there, closing level-3 measured 3, level-2 measured 2, and the outermost
% level-1 list itself closing measured 1 (never 0 -- \@itemdepth's own
% restore to the parent's depth happens *after* this hook runs, so the "0:
% outermost list itself closed" case below is dead code, kept only because
% removing a genuinely unreachable \ifcase branch isn't worth the risk of
% getting the numbering wrong a second time). So this is very much a real,
% reachable case -- not decorative symmetry -- and needed its own
% \addvspace branch (previously missing entirely, silently leaving the
% level-4-to-level-3 closing transition with no drop-space at all).
% Interpolated from the existing I/II/III progression as a starting guess,
% then re-measured directly against the compiled ZZO reference and bumped:
% 0.13 measured at 12.6pt, still *less* than level 4's own same-level gap
% of 13.4pt (violates rule III -- a level's own same-level gap must be
% smaller than that level's own closing-transition gap); raised until
% re-measurement cleared it with a comfortable margin again, matching the
% margin every other level's own pair already has.
\AtEndEnvironment{itemize}{%
	% Force vertical mode before \addvspace below (2026-09-22, found while
	% isolating a "Something's wrong--perhaps a missing \item" error that
	% fired at the *frame's* own \end{...}, not the itemize's, on every
	% slide using a plain, single-level \begin{itemize}...\end{itemize}
	% with nothing else in the frame -- confirmed via a from-scratch minimal
	% reproduction outside this project's own slides entirely, then
	% bisected to this exact hook by diffing against the same reproduction
	% with this \AtEndEnvironment call disabled). \AtEndEnvironment inserts
	% this code at the very *start* of \enditemize, before beamer's own
	% `\ifhmode\unskip\fi\endlist` -- so without this guard, \addvspace
	% below runs while TeX is often still in horizontal mode (right after
	% the last \item's own text, mid-paragraph). \addvspace's own hmode
	% branch defers via `\vadjust{\addvspace{...}}` instead of inserting
	% the glue directly, and that deferred insertion colliding with
	% beamer's own frame-body/overlay-replay collection is what actually
	% produced the corruption -- not reproducible as a plain "wrong
	% \item" typo anywhere in the affected slides' own source, only via
	% this exact hmode-vs-vmode timing. \ifhmode\unskip\par\fi here forces
	% vertical mode *before* \addvspace ever runs, matching the standard,
	% documented way to call \addvspace (kernel doc: "should only be used
	% between paragraphs") -- same idiom beamer's own \enditemize already
	% uses \unskip for, just moved earlier so it also covers this hook.
	% Doesn't disturb the \@itemdepth-based \ifcase dispatch just below:
	% \par doesn't touch \@itemdepth, so the closing list's own depth is
	% still read correctly afterward, same as before this guard existed.
	\ifhmode\unskip\par\fi
	\ifcase\@itemdepth
		% 0: dead code -- @itemdepth is never actually 0 when this hook
		% fires (see \SMItemizeLevelDropSpaceRatioIV's own comment
		% just above); kept only for alignment with the \ifcase branches
		% below, each keyed on the depth of the list that just closed.
	\or\addvspace{\SMItemizeLevelDropSpaceRatioI em}% level-1 (outermost) list itself closed -- a pre-existing quirk (not introduced by this comment's own edit): there's no parent item above level 1 for this to visually separate from, so this only actually matters if something follows the outermost list directly on the same frame; harmless on every slide seen so far since nothing does. Out of scope to redesign here.
	\or\addvspace{\SMItemizeLevelDropSpaceRatioII em}% level-2 list closed, back to a level-1 item
	\or\addvspace{\SMItemizeLevelDropSpaceRatioIII em}% level-3 list closed, back to a level-2 item
	\or\addvspace{\SMItemizeLevelDropSpaceRatioIV em}% level-4 list closed, back to a level-3 item
	\fi
}

% Per-level \itemsep -- the item-to-item gap *within* one list, at one
% fixed nesting depth (as opposed to the drop-space mechanism just above,
% which only fires when a nested list *closes* back up to its parent).
% Ratios, not lengths -- same "scale with font size" reasoning and fix as
% \SMItemizeLevelDropSpaceRatioI/II/III above (that comment has the
% full story on why a plain \newlength+\setlength pair here would freeze
% to one fixed pt value instead of actually tracking each level's own
% font size). Values only, defined here for discoverability alongside the
% drop-space ones above -- actually *applied* inside the `\itemize`
% redefinition itself (search "\beamer@cramped (beamerbaselocalstructure.sty"
% further above), not via a separate hook here; see that comment for why.
% Found missing entirely (2026-07-21, direct instruction, see
% slides/ZZO-Itemize-Spacing-Reference-current.tex's own header comment
% for the rules this and the drop-space mechanism together need to
% satisfy): this deck's own itemize never set \itemsep itself at any
% level -- confirmed directly, grepping this whole file -- so every depth
% shared whatever bare default beamer's own theme happens to leave
% \itemsep at, uniformly. That produced two visible problems, both
% measured directly (pymupdf against slides/2c2-Topics-Around-the-CAPM-Empirical-Validity-current.tex
% and an isolated bisection test, not guessed): (a) level-2 and level-3
% same-level item-to-item gaps came out *identical* (12.6pt each,
% confirmed against an isolated reproduction), since nothing distinguished
% them by depth at all -- only a closing transition ever got a
% level-dependent gap, never a same-level one; (b) that shared default
% \itemsep read as *smaller* than a single item's own wrapped-line
% \baselineskip (13.7-13.8pt within a wrap vs. 12.6pt between items),
% backwards from how a list should read -- an item boundary should never
% be tighter than a same-item line wrap.
% Ratios themselves are not arbitrary -- each chosen, then re-measured
% against the compiled ZZO reference slide and adjusted, specifically to
% satisfy all three of that slide's own header-comment rules together
% (extended to level 4 the same way, 2026-07-21), not just "level 1
% biggest" in isolation:
% - Rule I (level-1 same-level gap > level-2's > level-3's > level-4's)
%   falls straight out of \SMItemsepRatioI > ...II > ...III > ...IV,
%   added on top of the same \baselineskip every level shares (itself
%   already smaller per level here, since this theme's own itemize/
%   enumerate body font shrinks with depth too -- confirmed directly,
%   10.95pt/10pt/9pt/9pt for levels 1-4 respectively, level 4 sharing
%   level 3's own font since this theme has no distinct "subsubsub" size
%   of its own).
% - Rule II (every level's own same-level gap > that level's own wrapped-
%   line gap) needs \SMItemsepRatioIV alone (the smallest, hence the
%   binding case) to still be positive and large enough that it clearly
%   reads as more open than a plain wrapped line once rendered, not just
%   bigger on paper.
% - Rule III (a level's own same-level gap < that same level's own
%   *closing*-transition gap) is *not* automatic just because
%   \SMItemizeLevelDropSpaceRatio* was already bigger than the old
%   shared default -- \addvspace's own coalesce-to-the-larger-value
%   behavior (documented above) means the closing gap is
%   max(natural glue, that *drop-space* ratio), not a sum of the two. Each
%   drop-space ratio's own \ifcase branch is keyed on the depth of the
%   list that *just closed* (see \SMItemizeLevelDropSpaceRatioIV's own
%   comment for the full story on why, confirmed via \typeout, not
%   assumed), so it lines up with the same-numbered itemsep ratio despite
%   describing a different transition -- \SMItemizeLevelDropSpaceRatioIV
%   governs the level-4-list-closing-back-to-level-3 transition and is the
%   one that has to individually stay bigger than \SMItemsepRatioIV.
%   This branch was missing entirely at first (only I/II/III existed, a
%   leftover from before level 4 existed) -- confirmed directly on ZZO: the
%   level-4-to-level-3 closing transition measured 12.6pt with no drop-space
%   branch at all, less than level 4's own same-level gap of 13.4pt, a real
%   rule III violation. Adding the branch (\SMItemizeLevelDropSpaceRatioIV,
%   starting at 0.5) fixed it. Re-verify rule III directly against ZZO's own
%   compiled PDF (don't assume it still holds on paper) if any of these ratios is
%   ever changed again.

% Per-level "opening" space -- the gap between a level-n item and the
% first item of a level-(n+1) list that immediately follows it, as
% opposed to \SMItemizeLevelDropSpaceRatio* (a nested list *closing*
% back up) or \SMItemsepRatio* (two items at the *same* depth).
% Added 2026-07-21, direct instruction, replacing the original three
% separate rules on slides/ZZO-Itemize-Spacing-Reference-current.tex's
% own header comment with one multi-part inequality stated per level n:
%   level-n -> level-(n-1) close  >  level-n same-level  >
%   level-n -> level-(n+1) open  >  level-n's own wrap gap
% The "open" term was entirely unhandled before this -- confirmed
% directly on ZZO: every opening transition (level 1 into level 2, level
% 2 into level 3, level 3 into level 4) measured *smaller* than that same
% level's own wrap gap, backwards from what the rule needs, since nothing
% but beamer's own bare default \topsep governed it.
% \AtBeginEnvironment{itemize}, not \topsep -- \topsep was tried first
% (set via the same \ifcase\@itemdepth block, right after \itemsep,
% inside the `\itemize` redefinition above) and measured *unreliable*:
% two of the three levels barely moved despite a sizeable ratio, while
% the third moved by roughly the right amount -- \topsep's own value is
% evidently consulted by `\list` at some point this override comes too
% late for, at least for some levels, not evenly. \addvspace via
% \AtBeginEnvironment instead measured consistently and predictably
% (each level's own change matched its own ratio's calculated pt value
% directly), the same mechanism \SMItemizeLevelDropSpaceRatio* already
% relies on for the closing side. This hook fires with \@itemdepth
% reflecting the *parent* item's own depth (confirmed via \typeout, not
% assumed: opening the outermost level-1 list itself measured 0, opening
% a level-2 list from within a level-1 item measured 1, level-3-from-
% level-2 measured 2, level-4-from-level-3 measured 3) -- the opposite of
% \SMItemizeLevelDropSpaceRatio*'s own hook, which keys on the depth
% of the list that *just closed*. That makes the \ifcase numbering here
% match the "n" in the rule directly: case 1 is level-1's own open-to-
% level-2 ratio, case 2 is level-2's own open-to-level-3, case 3 is
% level-3's own open-to-level-4; case 0 (opening the outermost list
% itself) is a genuine no-op, correctly -- there's no parent item above
% level 1 for that transition to separate from.
% Ratios chosen, then re-measured against the compiled ZZO reference and
% adjusted, the same way as every other ratio here: each has to
% individually land strictly between that level's own \SMItemsepRatio*
% (the ceiling) and its own wrap gap (the floor). Re-verify directly
% against ZZO's own compiled PDF (don't assume it still holds on paper)
% if any of these, or the itemsep/drop-space ratios they're measured
% against, are ever changed again.
% \SMItemizeLevelOpenSpaceRatioI trimmed from 0.28 to 0.02 (2026-07-21,
% direct instruction: "the level-1 to level-2 distance...seems a bit
% much"). Confirmed directly: even at 0, this transition still measured
% ~16.3pt -- \addvspace's own coalesce-to-the-larger-value behavior means
% the ratio can only ever push this *above* whatever glue the list
% mechanism inserts here regardless (that natural floor, not the ratio),
% so 0.28 down to 0 only ever moved the actual gap by ~2pt (18.5->16.3);
% 0.02 splits the difference, keeping the mechanism nominally "in effect"
% rather than reading as accidentally disabled. \SMItemizeLevelOpenSpaceRatioIII
% (level-3-to-level-4) had to come down too, from 0.08 to 0.02, purely to
% restore its own margin -- changing RatioI shifted that unrelated
% transition's own measured gap closer to level 3's own same-level gap
% (12.1 vs. 12.0, too thin), confirmed directly on ZZO, not assumed;
% ratios elsewhere in this file have shown the same small-ripple behavior
% before (see \SMItemizeLevelDropSpaceRatioIV's own comment) --
% re-measure the *whole* table after touching any single one of these,
% not just the transition being changed.
\AtBeginEnvironment{itemize}{%
	% Same \ifhmode\unskip\par\fi guard as \AtEndEnvironment{itemize} above,
	% and for the identical reason -- see that hook's own comment for the
	% full story. This hook fires when a *nested* itemize opens (\@itemdepth
	% here is the *parent* item's depth, per the comment below), i.e. right
	% after "\item ...text..." mid-paragraph, immediately before the
	% \addvspace calls further down -- the same hmode-vs-vmode timing that
	% broke the closing hook, just on the opening side. Case 0 (opening the
	% outermost level-1 list itself) never reaches an \addvspace call below
	% regardless, so a plain single-level itemize was never affected by
	% this half of the bug -- only a *nested* \begin{itemize} (depth >=1 at
	% open time) actually needed this guard, which is why the "missing
	% \item" error survived this fix's own first pass specifically on the
	% two demo slides with a nested itemize (sectionframe's own ZZR,
	% standardframe's own ZZZ) while every plain single-level-itemize slide
	% was already fixed by the closing-hook guard alone.
	\ifhmode\unskip\par\fi
	\ifcase\@itemdepth
		% 0: opening the outermost (level-1) list itself -- no parent item
		% precedes it, nothing to add.
	\or\addvspace{\SMItemizeLevelOpenSpaceRatioI em}% opening a level-2 list from a level-1 item
	\or\addvspace{\SMItemizeLevelOpenSpaceRatioII em}% opening a level-3 list from a level-2 item
	\or\addvspace{\SMItemizeLevelOpenSpaceRatioIII em}% opening a level-4 list from a level-3 item
	\fi
}
\makeatother


% Even with the plain-enumerate template (SMPlainEnumerate, defined in the theme), \tableofcontents still
% doesn't match this deck's own normal itemize spacing -- beamer's own
% \beamer@sectionintoc (beamerbasetoc.sty), which wraps *every* section
% entry regardless of which "section in toc" template is active,
% unconditionally inserts a fixed \vskip1.5em before each entry (an
% unbounded \vfill, even worse, if subsections aren't hidden -- neither is
% reachable via a public \tableofcontents option, just this show/hide
% binary), *and* wraps the entry itself in \hbox{\vbox{...}} -- box
% wrapping that (like \SMMinipageTopCorrection's own minipage case
% elsewhere in this file) changes how surrounding baselineskip/topskip
% glue gets inserted compared to plain paragraph text. Both together added
% roughly 13pt of extra gap
% between entries beyond this deck's own normal ~18.6pt level-1 item
% spacing (measured directly against
% slides/57r-Reasons-for-Leasing-current.tex's own compiled PDF).
% Neither is controllable via a public \tableofcontents option, so this
% redefines \beamer@sectionintoc itself -- scoped safely since
% \tableofcontents is only ever called on the one Agenda slide in this
% deck -- keeping its numbering/hyperlink/part-number/hide-checking logic
% intact but dropping the vskip/vfill and the hbox/vbox wrapping, so an
% entry behaves exactly like a plain \item once printed.
\makeatletter
\def\beamer@sectionintoc#1#2#3#4#5{%
	\ifnum\c@tocdepth>0%
	\ifnum#4=\beamer@showpartnumber%
	{%
	\beamer@saveanother%
	\gdef\beamer@todo{}%
	\beamer@slideinframe=#1\relax%
	\expandafter\only\beamer@tocsections{\gdef\beamer@todo{%
			\beamer@tempcount=#5\relax%
			\advance\beamer@tempcount by\beamer@sectionadjust%
			\ifnum\beamer@tempcount>0
				\ifnum\beamer@tempcount>\beamer@toclastsection
					\def\inserttocsectionnumber{}%
				\else
					\edef\inserttocsectionnumber{\the\beamer@tempcount}%
				\fi
			\else
				\def\inserttocsectionnumber{}%
			\fi%
			\def\inserttocsection{\hyperlink{Navigation#3}{#2}}%
			\beamer@tocifnothide{\ifnum\c@section=#1\beamer@toc@cs\else\beamer@toc@os\fi}%
			{%
				\beamer@tocact{\ifnum\c@section=#1\beamer@toc@cs\else\beamer@toc@os\fi}{section in toc}%
			}%
		}%
	}%
	\beamer@restoreanother%
	}%
	\beamer@todo%
	\fi\fi%
}
\makeatother

% Left-aligns an enumerate's labels within a caller-chosen \labelwidth --
% used by the \tableofcontents Agenda slide (see
% slides/nkj-Agenda-current.tex's own comment for why its labels need
% this at 3.6pt). Defined here as a named helper rather than written as a
% literal \renewcommand{\makelabel}[1]{...#1...} inside the slide's own
% frame body, because that literal form -- specifically its raw `#1`
% parameter token -- breaks inside any frame that \tableofcontents forces
% through beamer's overlay reprocessing: beamer replays such a frame's
% collected body via \def-based machinery (\loop's own \iterate, and
% \beamer@doifinframe), and a `#` in the material of any \def must be
% doubled -- so the compile hits "! Illegal parameter number in definition
% of \iterate." (twice) plus the same for \beamer@doifinframe, once per
% compile, at that frame's own \end{frame}. TeX's error recovery ("I'll
% assume you meant ##") happens to produce the intended definition anyway,
% which is why the rendered PDF was always correct despite the errors --
% but lualatex still exits nonzero, which is exactly what kept breaking
% latexmk's automated lualatex/biber/rerun sequencing on the main deck
% (the "Known open issue" long tracked in CONVERSION_NOTES.md, originally
% misattributed to a biblatex interaction -- see its "Resolved" entry).
% Hoisting the definition out here means the frame body carries only this
% macro's *call* token, with no `#` in it, which sidesteps the whole
% mechanism. The inner \makelabel parameter is written `##1` below since
% it sits inside another \newcommand's body -- standard nesting, and also,
% incidentally, exactly the doubling rule the inline version fell foul of.
\newcommand{\SMTocLabelFlushLeft}[1]{%
	\setlength{\labelwidth}{#1}%
	\renewcommand{\makelabel}[1]{\makebox[\labelwidth][l]{##1}}%
}

% Same idea as \SMTocLabelFlushLeft above, plus an extra fixed
% \hspace shift before the label (for a list that needs to start further
% right than the plain nested-list indent, e.g. to line up with a
% sibling item's own text rather than its bullet -- see
% slides/scs-Political-Risk-Stability-Index-current.tex's own v1-5-0/
% v1-6-0 comments for the worked example this was hoisted out of).
% Found 2026-07-20 that the raw-`#1`-in-frame-body bug documented above
% is *not* actually \tableofcontents-specific, contrary to that comment's
% own framing -- slides/scs-Political-Risk-Stability-Index-current.tex has
% no \tableofcontents anywhere in the deck-wide overlay chain leading to
% it, yet hit the identical "Illegal parameter number in definition of
% \iterate" at its own \end{standardframe} from the exact same literal-
% \renewcommand{\makelabel}[1]{...#1...}-in-frame-body pattern. Whatever
% the precise trigger condition actually is, the same hoist-to-preamble
% fix (a bare macro *call*, no `#` token, left in the frame body) sidesteps
% it either way.
\newcommand{\SMEnumerateLabelFlushLeftShifted}[2]{%
	\setlength{\labelwidth}{#2}%
	\renewcommand{\makelabel}[1]{\hspace{#1}\makebox[\labelwidth][l]{##1}}%
}

% Width of a single digit in this deck's normal content font, measured once
% here (safe, plain preamble-load-time \settowidth) rather than inline in
% the Agenda slide's own frame body -- \tableofcontents there puts the whole
% frame through beamer's overlay-replay machinery, which the comment above
% \SMTocLabelFlushLeft already found chokes on raw tokens (like a bare
% \settowidth call) sitting directly in a replayed frame body, the same
% "Missing number"/"Illegal unit of measure" class of error as that `#`
% problem, confirmed directly the same way. Used by
% slides/nkj-Agenda-current.tex's own "widen the number-to-title gap"
% fix, which only ever needs to *reference* this precomputed value inside
% the frame, not measure anything itself there.
\newlength{\SMTocLabelExtraGap}
\settowidth{\SMTocLabelExtraGap}{0}

% For a two-column layout (e.g. bullets next to a chart/formula), nest
% beamer's own "columns" directly inside standardframe -- do NOT bake
% \begin{columns}/\end{columns} into a custom environment's begin/end code:
% beamer collects each frame's body as a token list to replay it once per
% overlay, and a columns environment that is opened in the "before" half and
% closed in the "after" half of a *different* environment confuses that
% collection (it never finds a properly matched \end{frame}). Nesting it in
% the document body, as below, avoids the issue entirely:
%
%   \begin{standardframe}{Title}
%   	\begin{columns}[t,onlytextwidth]
%   		\begin{column}{0.48\linewidth} ... \end{column}
%   		\begin{column}{0.48\linewidth} ... \end{column}
%   	\end{columns}
%   \end{standardframe}

% For automatic section dividers and single-statement "standout" slides, use
% what the gotham theme already provides instead of a custom environment:
%   \section{...}                                                (section divider, automatic)
%   \begin{frame}[standout]{Key takeaway}\end{frame}             (big single statement)
%   \begin{block}{...}/\begin{alertblock}{...}/\begin{exampleblock}{...}/\begin{referenceblock}{...}  (definitions, formulas, examples, references)
% For a course's own title slide, section-topic overview, quote, figure-only,
% and presenter-attention slides, see the five environments below.

% The block templates (block, alerted, example, reference, info and definition begin and end), their
% colors and the helpers they use (\SMBlockSideMargin, \SMBlockTitleWithMarker, \SMSetBlockBullets,
% \SMAlertMarker) are the theme's (beamerthemeSlideMage.sty); the three block types the beamer
% environments do not have are defined here.
\newenvironment<>{referenceblock}[1]{%
	\begin{actionenv}#2%
		\def\insertblocktitle{#1}%
		\par%
		\usebeamertemplate{block reference begin}}
	{\par%
		\usebeamertemplate{block reference end}%
	\end{actionenv}}

\newenvironment<>{infoblock}[1]{%
	\begin{actionenv}#2%
		\def\insertblocktitle{#1}%
		\par%
		\usebeamertemplate{block info begin}}
	{\par%
		\usebeamertemplate{block info end}%
	\end{actionenv}}

\newenvironment<>{definitionblock}[1]{%
	\begin{actionenv}#2%
		\def\insertblocktitle{#1}%
		\par%
		\usebeamertemplate{block definition begin}}
	{\par%
		\usebeamertemplate{block definition end}%
	\end{actionenv}}


% A narrower, freely-positionable version of \begin{block}{...} (2026-07-06,
% direct instruction) -- a plain block always spans the *current*
% \linewidth (that's what block begin/end's own \hspace+minipage wrapper
% above is measured against), so it can't be dropped into a fixed-width
% spot (e.g. beside a table, inside a tikz overlay \node) without first
% narrowing \linewidth some other way. \SMFloatingBlock just wraps a
% plain block in one more, *outer* minipage of the caller's own chosen
% width -- block begin/end's own inner minipage then measures itself
% against that narrower width instead of the frame's full one, nesting
% correctly with no changes needed to the block templates themselves.
% Title/body colors, margins, and font are exactly \block's own (same
% "block title"/"block body" beamer color names) -- this only changes the
% overall width and (by taking it out of the normal document flow, left to
% the caller) the position, not the look. Font size is also left to the
% caller (wrap the whole environment in a size declaration, e.g.
% {\footnotesize \begin{SMFloatingBlock}{5cm}{Excel} ... \end{SMFloatingBlock}})
% rather than fixed here, since the right size depends entirely on how
% much room the block ends up placed in.
%
% INNER PADDING (2026-09-30, direct instruction: systematic replacements for
% per-slide fixes): 20 of 27 floating blocks opened with a hand-written
% \vspace{2mm} and 17 closed with \par\vspace{1.2mm} (or 1mm/3.4pt), to give
% the body text room from the title bar and the block's bottom edge. All
% three SMFloating*Block environments now do that themselves, via an
% optional key list:
%   \begin{SMFloatingInfoBlock}[top=2mm, bottom=1.2mm]{6cm}{Title}
% Defaults top=2mm, bottom=1.2mm (the common values); 0pt inserts nothing
% at all. Blocks that had no or other padding carry their own values
% explicitly (e.g. [top=0pt]), so every block kept its exact look.
\makeatletter
\pgfqkeys{/SMFloatingBlock}{
	top/.store in=\SM@fb@top,
	bottom/.store in=\SM@fb@bottom,
}
\newcommand{\SM@fb@begin}[1]{%
	\pgfqkeys{/SMFloatingBlock}{top=2mm, bottom=1.2mm, #1}%
	\ifdim\SM@fb@top=0pt \else\vspace{\SM@fb@top}\fi
}
\newcommand{\SM@fb@end}{%
	\ifdim\SM@fb@bottom=0pt \else\par\vspace{\SM@fb@bottom}\fi
}
\NewDocumentEnvironment{SMFloatingBlock}{O{}mm}{%
	\begin{minipage}{#2}
	\begin{block}{#3}
	\SM@fb@begin{#1}%
}{%
	\SM@fb@end
	\end{block}
	\end{minipage}
}
\makeatother

% Same idea, but wrapping \exampleblock instead of \block (2026-07-06,
% direct instruction, first used on
% slides/fr3-Forms-of-Capital-Market-Efficiency-Strong-Semistrong-Weak-current.tex
% and slides/ha9-Monacorp-Truck-Lease-Breakeven-Rent-current.tex, both
% switching an existing \SMFloatingBlock to the "example" green header
% instead of the plain violet one) -- not worth generalizing into a single
% macro parameterized over the block name: \SMFloatingBlock's own body
% is only two lines, so a third variant is cheaper to just write out than
% to build a block-name-dispatch layer for.
\makeatletter
\NewDocumentEnvironment{SMFloatingExampleBlock}{O{}mm}{%
	\begin{minipage}{#2}
	\begin{exampleblock}{#3}
	\SM@fb@begin{#1}%
}{%
	\SM@fb@end
	\end{exampleblock}
	\end{minipage}
}
\makeatother

% Same idea again, wrapping \infoblock (2026-07-06, direct instruction,
% first used on
% slides/fr3-Forms-of-Capital-Market-Efficiency-Strong-Semistrong-Weak-current.tex's
% own "Random walk" formula box, switched from the green
% \SMFloatingExampleBlock to this blue one the same day).
\makeatletter
\NewDocumentEnvironment{SMFloatingInfoBlock}{O{}mm}{%
	\begin{minipage}{#2}
	\begin{infoblock}{#3}
	\SM@fb@begin{#1}%
}{%
	\SM@fb@end
	\end{infoblock}
	\end{minipage}
}
\makeatother

% Every SMFloatingInfoBlock anchored top-right against the institute
% logo (the deck's standard
% "(current page.north east)+(-\SMFrametitleRightPadding,-\SMContentTopAlign)"
% corner pattern) lands its own visible right edge short of the logo's,
% by a fixed amount independent of the block's own chosen width (#1) --
% \infoblock's underlying template insets its colored fill by
% \SMBlockSideMargin from *both* sides of whatever width it's given
% (that margin exists to match a first-level bullet's own left-hand
% position, see its own comment above; the fact that it also narrows the
% *right* side is just a side effect of reusing one symmetric
% hspace+minipage trick for both, not a deliberate right-margin choice),
% partly offset by beamer's own beamercolorbox "colsep*" background bleed
% (colsep* pads the colored fill outward past its own nominal box edge
% without also insetting the text, unlike unstarred "colsep" -- see
% beamerbasecolor.sty's own \beamercolorbox definition). Root-caused and
% measured directly (pymupdf against a compiled probe deck, not derived
% purely on paper) 2026-07-21, after
% slides/j8y-Asset-Allocation-Example-Weights-current.tex's own one-off
% "+2.2pt" nudge (measured that same day against just the institute
% logo's yellow chip) turned out to disagree with
% slides/fr3-Forms-of-Capital-Market-Efficiency-Strong-Semistrong-Weak-current.tex's
% own older "+3.165pt" nudge (2026-07-06, measured against the logo's
% *full* gray-bordered footprint, shadow/card included) -- both slides'
% own nudges were each individually correct for their own chosen
% reference edge, they just didn't agree on what "the logo's own edge"
% means. Direct instruction, 2026-07-21: the *full* gray-bordered
% footprint is the deck's intended reference, so 3.165pt is the value
% kept here; j8y's own newer 2.2pt nudge was the one revised to match,
% not this one.
%
% Use "-\SMFloatingInfoblockRightAlign" (not the bare
% "-\SMFrametitleRightPadding") for a SMFloatingInfoBlock's own
% anchor node whenever its right edge is meant to align with the logo --
% every current use of this pattern across the deck has been updated to
% match; new slides using this same corner pattern should use this
% length from the start rather than re-measuring their own one-off nudge.
\newlength{\SMFloatingInfoblockRightAlign}
\setlength{\SMFloatingInfoblockRightAlign}{\dimexpr\SMFrametitleRightPadding-3.165pt\relax}

% An infoblock's colored title bar starts 6.02pt below the top of its own
% box (the block template's top padding); a floating infoblock anchored at
% the Top Margin therefore sits 6.02pt low unless its anchor is lifted by
% this much. Measured 2026-07-23 (y3a/gf2: title-bar rect at 47.60pt vs.
% the 41.58pt Top Margin), written as a literal "+6.02pt" on 15 anchors
% until 2026-09-30.
\newlength{\SMFloatingInfoblockTopAlign}
\setlength{\SMFloatingInfoblockTopAlign}{6.02pt}

% "Bottom Margin" as a length (2026-09-30): the gap between the footer box
% and content anchored at the bottom of a slide, the mirror image of the Top
% Margin gap between the title bar and the first line (direct instruction on
% bwi, 2026-07-2x: "the same clearance from the footer that the text has from
% the title bar"). Measured from the page's bottom edge, like
% \SMContentTopAlign is from its top edge.
\newlength{\SMContentBottomAlign}
\setlength{\SMContentBottomAlign}{\dimexpr\SMFooterBoxHeight+\SMContentTopAlign-\SMTitleBarBottomAlign\relax}

% \cornerfigure[keys]{image file or content} (2026-09-30, direct instruction:
% systematic replacements for per-slide fixes) -- a figure in a corner of the
% content area, as a zero-space overlay: the construct 105 slides wrote out
% by hand,
%   \vbox to 0pt{\vss\begin{tikzpicture}[overlay, remember picture]
%     \node[anchor=north east, inner sep=0pt] at ($(current page.north east)
%       +(-\SMFrametitleRightPadding,-\SMContentTopAlign)$) {...};
%   \end{tikzpicture}}
% The usual call (2026-10-01, direct instruction: "the width or height of
% the figure has to be specified in the code, that's all that should be
% needed") names the image and one size; the other follows from the aspect
% ratio:
%   \cornerfigure[height=5.6cm]{media/chart-x.png}
%   \cornerfigure[width=4cm, border]{media/photo-x.jpg}
% Keys:
%   width=, height=      the image's size; with either one the argument is an
%                        image file (both: fitted inside, aspect ratio kept).
%                        Without them the argument is arbitrary content (a
%                        block, a \parbox, an \includegraphics with options).
%   border               the yellow image border (SMImageBorder); the
%                        border's outer edge sits on the margins
%   corner=ne|nw|se|sw   which corner (default ne): Top/Bottom Margin
%                        (\SMContentTopAlign/\SMContentBottomAlign)
%                        and Right/Left Margin (\SMFrametitleRightPadding,
%                        \SMFrametitleLeftPadding)
%   corner=e|w           on the Right/Left Margin, centred vertically between
%                        Top and Bottom Margin
%   block                the content is a SMFloatingInfoBlock: lifts it by
%                        \SMFloatingInfoblockTopAlign and aligns its
%                        visible right edge with the logo
%                        (\SMFloatingInfoblockRightAlign)
%   xshift, yshift       a deliberate offset from that corner (right/up +)
%   node={...}           extra TikZ node options
% The node has no inner or outer sep, so the figure's own edge -- or, with a
% border, the border's outer edge (see SMImageBorder) -- is what lands on
% the margins. Like any remember-picture overlay it needs the usual second
% compile. Text that must not run under the figure: besidefigure below.
\makeatletter
\newif\ifSM@cf@block
\newif\ifSM@cf@border
\pgfqkeys{/SMCorner}{
	corner/.store in=\SM@cf@corner,
	xshift/.store in=\SM@cf@x,
	yshift/.store in=\SM@cf@y,
	node/.store in=\SM@cf@node,
	width/.store in=\SM@cf@w,
	height/.store in=\SM@cf@h,
	gap/.store in=\SM@bf@gap,
	block/.code={\SM@cf@blocktrue},
	border/.code={\SM@cf@bordertrue},
}
% all keys back to their defaults, then the caller's
\def\SM@cf@setkeys#1{%
	\SM@cf@blockfalse\SM@cf@borderfalse
	\pgfqkeys{/SMCorner}{corner=ne, xshift=0pt, yshift=0pt, node={},
		width={}, height={}, gap=12pt, #1}%
}
% \SM@cf@content: the figure itself -- \includegraphics[width/height]{#1}
% when a size key was given, #1 unchanged otherwise (\long: content such as
% a floating block may hold paragraphs)
\long\def\SM@cf@makecontent#1{%
	\def\SM@cf@gkeys{}%
	\ifx\SM@cf@w\@empty\else
		\edef\SM@cf@gkeys{width=\unexpanded\expandafter{\SM@cf@w}}%
	\fi
	\ifx\SM@cf@h\@empty\else
		\edef\SM@cf@gkeys{\ifx\SM@cf@gkeys\@empty\else\unexpanded\expandafter{\SM@cf@gkeys},keepaspectratio,\fi
			height=\unexpanded\expandafter{\SM@cf@h}}%
	\fi
	\ifx\SM@cf@gkeys\@empty
		\def\SM@cf@content{#1}%
	\else
		\edef\SM@cf@content{\noexpand\includegraphics[\unexpanded\expandafter{\SM@cf@gkeys}]{\unexpanded{#1}}}%
	\fi
}
\NewDocumentCommand{\cornerfigure}{O{} +m}{%
	\begingroup
	\SM@cf@setkeys{#1}%
	\SM@cf@makecontent{#2}%
	\ifSM@cf@block
		\edef\SM@cf@right{\SMFloatingInfoblockRightAlign}%
		\edef\SM@cf@top{\dimexpr\SMContentTopAlign-\SMFloatingInfoblockTopAlign\relax}%
	\else
		\edef\SM@cf@right{\SMFrametitleRightPadding}%
		\edef\SM@cf@top{\SMContentTopAlign}%
	\fi
	\def\SM@cf@anchor{north east}%
	\edef\SM@cf@at{(current page.north east)+(-\SM@cf@right+\SM@cf@x,-\SM@cf@top+\SM@cf@y)}%
	\ifdefstring{\SM@cf@corner}{nw}{%
		\def\SM@cf@anchor{north west}%
		\edef\SM@cf@at{(current page.north west)+(\SMFrametitleLeftPadding+\SM@cf@x,-\SM@cf@top+\SM@cf@y)}}{}%
	\ifdefstring{\SM@cf@corner}{se}{%
		\def\SM@cf@anchor{south east}%
		\edef\SM@cf@at{(current page.south east)+(-\SM@cf@right+\SM@cf@x,\SMContentBottomAlign+\SM@cf@y)}}{}%
	\ifdefstring{\SM@cf@corner}{sw}{%
		\def\SM@cf@anchor{south west}%
		\edef\SM@cf@at{(current page.south west)+(\SMFrametitleLeftPadding+\SM@cf@x,\SMContentBottomAlign+\SM@cf@y)}}{}%
	% e/w: on the Right/Left Margin, centred vertically between Top and Bottom Margin
	\ifdefstring{\SM@cf@corner}{e}{%
		\def\SM@cf@anchor{east}%
		\edef\SM@cf@at{(current page.north east)+(-\SM@cf@right+\SM@cf@x,-0.5\paperheight-0.5\SMContentTopAlign+0.5\SMContentBottomAlign+\SM@cf@y)}}{}%
	\ifdefstring{\SM@cf@corner}{w}{%
		\def\SM@cf@anchor{west}%
		\edef\SM@cf@at{(current page.north west)+(\SMFrametitleLeftPadding+\SM@cf@x,-0.5\paperheight-0.5\SMContentTopAlign+0.5\SMContentBottomAlign+\SM@cf@y)}}{}%
	% TikZ's calc parser must see the coordinate itself, not a macro: build
	% the \node call with everything but the content already expanded
	\edef\SM@cf@cmd{\noexpand\node[anchor=\SM@cf@anchor, inner sep=0pt, outer sep=0pt%
		\ifSM@cf@border,SMImageBorder\fi
		\ifx\SM@cf@node\@empty\else,\unexpanded\expandafter{\SM@cf@node}\fi]
		at ($\SM@cf@at$)}%
	\vbox to 0pt{\vss
		\begin{tikzpicture}[overlay, remember picture]
			\SM@cf@cmd {\SM@cf@content};
		\end{tikzpicture}%
	}%
	\endgroup
}

% \begin{besidefigure}[keys]{image file or content} text \end{besidefigure}
% (2026-09-30; 2026-10-01 the figure's width is measured, not passed) -- text
% next to a figure in the top-right corner: the figure goes into \cornerfigure
% (same keys: height=/width=, border, ...), the text into a top-aligned
% minipage exactly as wide as the room left of the figure minus `gap`
% (default 12pt). 84 slides built this by hand, each with its own guessed
% minipage width (9.3cm, 0.58\linewidth, ...).
%   \begin{besidefigure}[height=5.6cm, border]{media/chart-x.png}
%   	\begin{itemize} ... \end{itemize}
%   \end{besidefigure}
% Optical top alignment (v5.27.0): the text column's first line sits on the Top
% Margin by its line box (the ascender height, 0.935em), so its visible letters
% start about 0.2em lower; a figure placed on the margin itself looked that
% much too high next to it. besidefigure therefore lowers its figure by this
% much (\SMBesideFigureTopGap, default 0.2em of the surrounding text), so the
% figure's top edge meets the top of the first line's letters. Give
% yshift=0pt among the keys to put the figure on the margin again.
\newlength{\SM@bf@width}
\NewDocumentEnvironment{besidefigure}{O{} +m}{%
	\SM@cf@setkeys{#1}%
	\SM@cf@makecontent{#2}%
	\settowidth{\SM@bf@width}{\SM@cf@content}%
	\ifSM@cf@border\addtolength{\SM@bf@width}{\SMImageBorderWidth}\fi
	\cornerfigure[yshift=-\SMBesideFigureTopGap,#1]{#2}%
	% text column: from the text's left edge to the figure's left edge minus gap
	\setlength{\SM@bf@width}{\dimexpr\linewidth+\beamer@rightmargin-\SMFrametitleRightPadding-\SM@bf@width-\SM@bf@gap\relax}%
	\begin{minipage}[t]{\SM@bf@width}%
}{%
	\end{minipage}%
}

% \fitfigure[keys]{image file} (2026-09-30) -- an image as large as it fits,
% keeping its aspect ratio, in the box from a top line down to the Bottom
% Margin (\SMContentBottomAlign) and from the Left to the Right Margin,
% placed as a zero-space overlay. Replaces the per-slide
% "scale up ... as much as possible" arithmetic (hand-computed widths and
% \vspace offsets, often converged in two passes).
%   top=here|margin   where the box starts: 'here' (default) = the current
%                     position in the slide (below the text so far), plus
%                     `gap`; 'margin' = the Top Margin
%   gap=<length>      space below the text for top=here (default 6pt)
%   valign=top|center|bottom, halign=left|center|right  (default center, center)
%   bottom=<length>   the box's bottom, from the page's bottom edge (default
%                     \SMContentBottomAlign)
%   node={...}        extra TikZ node options (e.g. SMImageBorder)
% The position is a remember-picture coordinate, so like every overlay here
% it needs the second compile.
\newcount\SM@fit@count
\def\SM@fit@settop#1{\coordinate (SMFit@top) at #1;}
\pgfqkeys{/SMFit}{
	top/.store in=\SM@fit@top,
	gap/.store in=\SM@fit@gap,
	valign/.store in=\SM@fit@valign,
	halign/.store in=\SM@fit@halign,
	bottom/.store in=\SM@fit@bottom,
	node/.store in=\SM@fit@node,
}
\NewDocumentCommand{\fitfigure}{O{} m}{%
	\begingroup
	\pgfqkeys{/SMFit}{top=here, gap=6pt, valign=center, halign=center,
		bottom=\SMContentBottomAlign, node={}, #1}%
	\global\advance\SM@fit@count by 1
	\edef\SM@fit@name{SMFit\the\SM@fit@count}%
	\ifdefstring{\SM@fit@top}{here}{%
		\vbox to 0pt{\hbox{\tikz[remember picture, overlay]\coordinate (\SM@fit@name) at (0,0);}\vss}%
	}{}%
	% alignment choices, decided here (pgfmath can't evaluate \ifstrequal)
	\def\SM@fit@xexpr{(\x2+\x3)/2}\def\SM@fit@h@anchor{}%
	\ifdefstring{\SM@fit@halign}{left}{\def\SM@fit@xexpr{\x2}\def\SM@fit@h@anchor{west}}{}%
	\ifdefstring{\SM@fit@halign}{right}{\def\SM@fit@xexpr{\x3}\def\SM@fit@h@anchor{east}}{}%
	\def\SM@fit@yexpr{(\y1+\y2)/2}\def\SM@fit@v@anchor{}%
	\ifdefstring{\SM@fit@valign}{top}{\def\SM@fit@yexpr{\y1}\def\SM@fit@v@anchor{north}}{}%
	\ifdefstring{\SM@fit@valign}{bottom}{\def\SM@fit@yexpr{\y2}\def\SM@fit@v@anchor{south}}{}%
	\edef\SM@fit@anchor{\SM@fit@v@anchor
		\ifx\SM@fit@v@anchor\@empty\else\ifx\SM@fit@h@anchor\@empty\else\space\fi\fi
		\SM@fit@h@anchor}%
	\ifx\SM@fit@anchor\@empty\def\SM@fit@anchor{center}\fi
	\ifdefstring{\SM@fit@top}{here}{%
		\edef\SM@fit@topcoord{(\SM@fit@name)}%
		\edef\SM@fit@gapexpr{\SM@fit@gap}%
	}{%
		\def\SM@fit@topcoord{($(current page.north)+(0,-\SMContentTopAlign)$)}%
		\def\SM@fit@gapexpr{0pt}%
	}%
	\edef\SM@fit@nodeopts{anchor=\SM@fit@anchor, inner sep=0pt%
		\ifx\SM@fit@node\@empty\else,\unexpanded\expandafter{\SM@fit@node}\fi}%
	\vbox to 0pt{\vss
		\begin{tikzpicture}[overlay, remember picture]
			\expandafter\SM@fit@settop\expandafter{\SM@fit@topcoord}%
			\coordinate (SMFit@bl) at ($(current page.south west)+(\SMFrametitleLeftPadding,\SM@fit@bottom)$);
			\coordinate (SMFit@r) at ($(current page.east)+(-\SMFrametitleRightPadding,0)$);
			\path let \p1 = (SMFit@top), \p2 = (SMFit@bl), \p3 = (SMFit@r) in
				\pgfextra{%
					\pgfmathsetlengthmacro\SM@fit@h{max(\y1-\SM@fit@gapexpr-\y2,1pt)}%
					\pgfmathsetlengthmacro\SM@fit@w{max(\x3-\x2,1pt)}%
					\pgfmathsetlengthmacro\SM@fit@x{\SM@fit@xexpr}%
					\pgfmathsetlengthmacro\SM@fit@y{\SM@fit@yexpr}%
					\ifdefstring{\SM@fit@valign}{top}{\pgfmathsetlengthmacro\SM@fit@y{\y1-\SM@fit@gapexpr}}{}%
					\ifdefstring{\SM@fit@valign}{center}{\pgfmathsetlengthmacro\SM@fit@y{(\y1-\SM@fit@gapexpr+\y2)/2}}{}%
					\global\let\SM@fit@H\SM@fit@h \global\let\SM@fit@W\SM@fit@w
					\global\let\SM@fit@X\SM@fit@x \global\let\SM@fit@Y\SM@fit@y
				};
			\expandafter\node\expandafter[\SM@fit@nodeopts] at (\SM@fit@X,\SM@fit@Y)
				{\includegraphics[width=\SM@fit@W, height=\SM@fit@H, keepaspectratio]{#2}};
		\end{tikzpicture}%
	}%
	\endgroup
}
\makeatother

%%%%%%%%%%%%%%%%%%%%%%%%%%%%%%%%%%%%%%%%%%%%%%%%%%%%%%%%%%%%%%%%%%%%%%%%%%%%%
% Special frame types (converted from ExampleSlides.pptx; see README.md)
%
% All five below deliberately reuse as much of standardframe -- and of each
% other -- as possible, rather than reimplementing frame chrome from
% scratch: sectionframe, figureframe, and highlightframe are literally thin
% wrappers around standardframe itself (same title bar, logo, footer, bullet
% styling; only what's documented at each one differs). titleframe and
% quoteframe can't be standardframe wrappers -- both are title-less, and
% quoteframe replaces the whole body layout -- but still reuse its palette,
% fonts, and margins (\SMFrametitleLeftPadding, the institute logo
% drawing command, etc.) rather than inventing new ones.
%%%%%%%%%%%%%%%%%%%%%%%%%%%%%%%%%%%%%%%%%%%%%%%%%%%%%%%%%%%%%%%%%%%%%%%%%%%%%

% titleframe: the deck's own title slide (ExampleSlides.pptx slide 1) --
% typically used exactly once, at the very start of a course's first deck.
% One continuous yellow box -- like standardframe's own title bar -- flush
% with the page's top and left edges, stopping short of the institute logo
% on the right (the same margin already used around the logo elsewhere:
% \SMLogoWidth + 2*\SMFrametitleRightPadding), holding the title and,
% below it, a second line (author/term/whatever) in the same black text as
% standardframe's own frametitle. Usage:
%
%   \titleframe{Financial Markets}{Stefan Palan\\Summer term 2026}
%
% A *command*, not an environment, taking the second (lower) line's content
% as its second argument -- plain text, use \\ for line breaks same as any
% normal LaTeX paragraph, not an itemize, since the source deck's own
% subtitle block is just stacked lines. This has to be a command rather
% than an environment with the body as its content: beamer's frame body is
% captured by a delimited-argument macro (\beamer@collect@@body, looked up
% by triggering the error below on purpose) hunting for a literal
% \end{frame}, and a beamercolorbox -- needed by an earlier version of this
% to wrap the subtitle text in its own colored band -- opened in a *custom
% environment's* begin code and closed in its end code confuses that hunt
% exactly like the \columns example above (same root cause, different
% environment): with
% \newenvironment{titleframe}{...\begin{beamercolorbox}{...}}{\end{beamercolorbox}...},
% compiling ends in "! File ended while scanning use of \beamer@collect@@body."
% every time, however short the body. Taking the content as a normal brace
% argument instead sidesteps the whole issue. No footer/frame number
% (matches beamer's own convention for title slides, e.g. the old
% \begin{frame}[plain,noframenumbering]\titlepage\end{frame} this replaces --
% see README.md) and no chapter-logo support: both would be meaningless on a
% slide that appears once, before any chapter exists.
%
% Unlike the frametitle hook further up, neither the logo nor the yellow box
% here need the white-patch-and-redraw trick: the box's own width already
% stops short of the logo (rather than being painted full-width and then
% partly painted back over), so there's no yellow left to paint over in the
% first place. Both still need \tikz's "remember picture" on their
% "current page"-anchored nodes, though, and so still need the usual second
% compile pass (see "Title bar and logo" in README.md) -- checked
% empirically: without it, "current page" doesn't resolve to the actual page
% box on the first pass, and the logo/yellow box end up positioned nowhere
% near the page corners they're meant to be anchored to.
%
% \SMTitleframeBoxHeight is a fixed target for the box's *bottom* edge
% (matching what was actually asked for: "reaching down to where the lower
% grey bar currently ends" in the two-band design this replaces), not the
% blank space above the text -- that space is derived from it instead (box
% height minus the title+subtitle content's own measured height), so the
% bottom edge stays fixed even if a longer title ever wraps onto a second
% line; the blank space above would just shrink to compensate. Measured
% against the old two-band design's own rendered PDF coordinates (before
% this change): its second band's bottom edge sat 187.147pt below the page's
% top edge.
\newlength{\SMTitleframeBoxWidth}
\setlength{\SMTitleframeBoxWidth}{\dimexpr\paperwidth-\SMLogoWidth-2\SMFrametitleRightPadding\relax}
% Text inset from the box's own edges: \SMFrametitleLeftPadding, the same
% length standardframe's own title text uses, plus \SMTitleframeTextSep
% (2.5ex, matching this command's own former per-band beamercolorbox "sep").
% Because the leading blank space above the text (further down) is itself
% computed as "box height minus the content's own natural height", the
% content's *bottom* padding is the only thing that determines how far the
% last line sits above the box's bottom edge -- the leading space silently
% absorbs whatever that padding is, keeping the bottom edge fixed either
% way. (The relationship is a direct 1:1 offset, checked empirically by
% comparing two rendered values, but with a nonzero intercept from beamer's
% own default top-of-frame vskip -- there's no clean closed-form for the
% right value, hence the empirical calibration here rather than a formula.)
% Calibrated (against the rendered PDF's own text coordinates, not just
% visually) to about 1.5 line-heights of the subtitle text's own \Large
% size (~20.6pt/line here) above the box's bottom edge; re-measure and
% adjust if the subtitle's font size, or beamer's own default top-of-frame
% vskip, ever change.
\newbox\SMTitleframeBox
\makeatletter
\newcommand{\titleframe}[2]{%
	% No noframenumbering (2026-08-11, direct instruction): this frame has
	% no footer at all ("plain"), so removing it has no visible effect on
	% the title slide itself -- but it does mean the title slide now
	% consumes a real "framenumber" step, same as any other frame, so the
	% deck's first real content slide is numbered 2, not 1 (matching this
	% slide's own physical position in the compiled PDF -- SlideMage's
	% own deck-editor "Render" preview relies on this: it derives a
	% previewed slide's page number from its plain list position, which
	% only equals the real compiled frame number when every \slide{}-family
	% call -- including this one -- consumes exactly one framenumber step).
	\SMOpenFrame{[plain,t]}%
	% This frame has no footer to carry the usual "(slide id: ...)" marker
	% -- see \SMInvisibleSlideIDStamp's own comment above for why that
	% breaks SlideMage's compile-time preview snapshot, and what this
	% stamps instead.
	\SMInvisibleSlideIDStamp%
	% A bare tikzpicture behaves like an hbox (same reasoning as the footer
	% box's own \vbox-to-0pt wrapper further up): sitting directly in the
	% frame's vertical list, it silently opens an implicit paragraph,
	% switching into horizontal mode. Everything placed after it -- the
	% \vspace* further down included -- would then run in horizontal mode
	% instead of vertical mode, which is what a first version of this command
	% got wrong: with two bare tikzpictures here, \vspace*'s glue ended up as
	% *horizontal* space on an implicit line together with them (visible as
	% the title/subtitle shifting sideways, not down, and the vspace amount
	% appearing to have no effect at all -- it was moving the wrong axis).
	% \vbox to 0pt{\vss ...} keeps each overlay at zero space in the normal
	% flow without ever leaving vertical mode.
	%
	% \SMArtifactBegin/\SMArtifactEnd (defined with the footer, "PDF
	% accessibility" comment): this box is purely decorative background, same
	% kind of element as the footer's own yellow box -- marked the same way,
	% for a screen reader to skip. titleframe has no footer at all to worry
	% about separately (it's a "plain" frame, no footline), so this box is
	% the only thing here that needs it; the title/subtitle text right below
	% is real content and stays unmarked.
	\SMArtifactBegin{/Type /Background}%
	\vbox to 0pt{\vss
	\begin{tikzpicture}[overlay, remember picture]
		% \SMTitleBoxColor, not a literal color token -- same setting the
		% standard frametitle bar (colorA, above) uses, so the title slide's
		% own box matches every other title bar in the deck (2026-07-16,
		% direct instruction).
		% outer sep=0pt here and on the logo below: without it, tikz's
		% default outer sep insets the fill/logo 0.2pt from the anchored
		% page corner -- a hairline white line above/left of this box, and
		% its right edge/the logo 0.2-0.4pt off every other frame's. See
		% the frametitle hook's own comment for the full story.
		\node[fill=\SMTitleBoxColor, anchor=north west, inner sep=0pt, outer sep=0pt,
		      minimum width=\SMTitleframeBoxWidth,
		      minimum height=\SMTitleframeBoxHeight]
			at (current page.north west) {};
	\end{tikzpicture}%
	}%
	\SMArtifactEnd
	\vbox to 0pt{\vss
	\begin{tikzpicture}[overlay, remember picture]
		% inner ysep derived, not \SMFrametitleRightPadding -- see the
		% frametitle hook's fixed-position logo redraw for why (matches the
		% inline logo's vertically-centered position on every other frame).
		\node[anchor=north east, inner xsep=\SMFrametitleRightPadding,
		      inner ysep=\dimexpr\SMTitleBarBottomAlign-\SMLogoHeight\relax,
		      outer sep=0pt]
			at (current page.north east) {\SMInstituteLogo};
	\end{tikzpicture}%
	}%
	\setbox\SMTitleframeBox=\vbox{%
		\hsize=\SMTitleframeBoxWidth
		\parindent=0pt
		\leftskip=\SMFrametitleLeftPadding
		\advance\leftskip by \SMTitleframeTextSep
		\rightskip=\SMFrametitleLeftPadding
		\advance\rightskip by \SMTitleframeTextSep
		% \SMTitleTextColor, not a literal color token -- follows the
		% box color change right above (2026-07-16, direct instruction):
		% leaving this hardcoded dark would be unreadable once the box
		% itself switches to a dark blue for MCI/Innsbruck.
		\color{\SMTitleTextColor}%
		% Beamer redefines \\ inside every frame to peek for an overlay spec
		% (\\<2->, for line-by-line reveals in itemize/enumerate); that
		% redefinition assumes it's running in beamer's own frame-body
		% machinery and chokes ("Undefined control sequence", silently
		% swallowing the following word) when \\ is used for a plain line
		% break in a raw \vbox like this one. Restoring plain LaTeX's own
		% \\ (\@normalcr) for this group only fixes it without affecting \\
		% anywhere else in the deck. (The @ here needs catcode 11, which is
		% why this whole command definition -- not just this line -- is
		% wrapped in \makeatletter/\makeatother: macro bodies are tokenized
		% at *definition* time, so a \makeatletter placed inside the body
		% itself would take effect too late, after "\@normalcr" had already
		% been mis-tokenized as "\@" followed by four ordinary-text letters.)
		\let\\\@normalcr
		\vskip\SMTitleframeTextSep
		\noindent{\Huge\bfseries #1\par}%
		\vskip\SMTitleframeTextSep
		\vskip\SMTitleframeBandGap
		\vskip\SMTitleframeTextSep
		\noindent{\Large #2\par}%
		\vskip\SMTitleframeBottomPad
	}%
	\vspace*{\dimexpr\SMTitleframeBoxHeight-\ht\SMTitleframeBox-\dp\SMTitleframeBox\relax}%
	\box\SMTitleframeBox
	% This command's own \color{\SMTitleTextColor} above (set inside a
	% stored box register) can leak into the *next* frame's body text --
	% see standardframe's own opening code ("Forces this frame's own body")
	% for the real fix and the full story (a trailing reset attempted
	% here first turned out to be a silent no-op, since nothing draws after
	% it within this same frame to actually force the color operator to be
	% emitted).
	\end{frame}%
}
\makeatother

% fullbleedframe: a slide that is *only* a full-screen image (FM2_2026s_ps_deck.pptx
% slide 13, a recruitment flyer pasted in as one embedded picture on an
% otherwise blank layout -- no title, no PPTX text at all) -- no title bar,
% no footer, no logo, matching that source's own blank layout intent, unlike
% figureframe (which is standardframe underneath and always keeps the usual
% title/footer chrome). The image is scaled to fit the slide while keeping
% its own aspect ratio (\includegraphics's keepaspectratio, capped by
% \paperwidth *and* \paperheight together), centered; whatever's left over
% on the binding axis (e.g. a widescreen slide holding a taller/narrower
% image) is filled by a solid black backdrop rather than left showing the
% frame's usual white canvas -- the point of a full-*bleed* image is that
% nothing reads as "empty margin" even when the aspect ratios don't match.
%
%   \fullbleedframe{media/photo-datacamp-recruitment-flyer}
%
% \SMFullBleedBlack is a fixed HTML black, deliberately *not* a
% \colorlet of c_primary_dark: this frame type's backdrop should stay
% neutral true black regardless of whether c_primary_dark itself is ever
% restyled to something else (e.g. a softer near-black, the way
% c_custom_gray already is elsewhere in this file) for branding
% reasons -- those are two different jobs (brand color vs. letterbox
% backdrop) that only happen to want the same value today.
\definecolor{SMFullBleedBlack}{HTML}{000000}
\newcommand{\fullbleedframe}[1]{%
	% No noframenumbering -- see \titleframe's own comment above for why
	% (same reasoning, same 2026-08-11 change: this frame has no footer
	% either, so nothing here changes visibly, only later slides' numbers).
	\SMOpenFrame{[plain,t]}%
	% This frame has no footer either -- same \SMInvisibleSlideIDStamp
	% need as \titleframe above, see its own comment for why.
	\SMInvisibleSlideIDStamp%
	% \SMArtifactBegin/\SMArtifactEnd (defined with the footer, "PDF
	% accessibility" comment): the black backdrop is purely decorative
	% letterboxing, marked the same way as titleframe's/quoteframe's own
	% background boxes so a screen reader skips it. The image itself is left
	% unmarked *here* -- not because it goes undescribed, but because
	% \includegraphics's own global wrapper (this file's own
	% \RenewDocumentCommand{\includegraphics}, "PDF accessibility" section)
	% already does that job for every \includegraphics call site, this one
	% included: it records the image's role/source with
	% slidemage_alt.mark()/recordSource() and, for an informative image, emits
	% its short alt text as an invisible layer via \SMAltInvisible
	% before the image itself. Nothing extra is needed at this call site.
	\SMArtifactBegin{/Type /Background}%
	\vbox to 0pt{\vss
	\begin{tikzpicture}[overlay, remember picture]
		% outer sep=0pt: see the frametitle hook's comment -- without it the
		% backdrop sits 0.2pt in from the anchored page corner.
		\node[fill=SMFullBleedBlack, anchor=south west, inner sep=0pt, outer sep=0pt,
		      minimum width=\paperwidth, minimum height=\paperheight]
			at (current page.south west) {};
	\end{tikzpicture}%
	}%
	\SMArtifactEnd
	\SMFullBleedPick{#1}{\SMFullBleedLayers{#1}}{\SMFullBleedImage{#1}}%
	\end{frame}%
}
\makeatletter
\newcommand{\SMFullBleedPick}[3]{\in@{,}{#1}\ifin@#2\else#3\fi}
\makeatother

% \fullbleedframe{a.png,b.png,c.png} (2026-09-18, direct instruction:
% combine the Fink/Palan/Theissen excursus's build-step screenshots into
% one slide): a comma-separated list layers the images on top of each
% other, revealed one per overlay step (image k appears at step k and stays,
% each covering the previous one -- the screenshots are same-size, opaque
% and each already contains every earlier step's content, so stacking
% reproduces the source's own build animation). Instructor deck only --
% the during-course/after-course versions show just the last (final)
% image, in a single step. A single path behaves
% exactly as before. Once a layer is covered by the next one it is wrapped
% in the same /Background artifact as the backdrop above, so on step k a
% screen reader reads only image k's own alt text (each screenshot is
% described as its own "build step k of n") instead of every buried layer's.
\newcommand{\SMFullBleedImage}[1]{%
	\vbox to 0pt{\vss
	\begin{tikzpicture}[overlay, remember picture]
		\node[anchor=center] at (current page.center)
			{\includegraphics[width=\paperwidth, height=\paperheight, keepaspectratio]{#1}};
	\end{tikzpicture}%
	}%
}
\newcount\SMFBtotal
\newcommand{\SMFullBleedLayers}[1]{%
	\global\SMFBtotal=0
	\foreach \SMFBimg in {#1} {\global\advance\SMFBtotal by 1 }%
	\foreach \SMFBimg [count=\SMFBi] in {#1} {%
		\ifSMInstructorVersion
			\ifnum\SMFBi<\SMFBtotal
				\pgfmathtruncatemacro{\SMFBnext}{\SMFBi+1}%
				\only<\SMFBi>{\SMFullBleedImage{\SMFBimg}}%
				\only<\SMFBnext->{%
					\SMArtifactBegin{/Type /Background}%
					\SMFullBleedImage{\SMFBimg}%
					\SMArtifactEnd
				}%
			\else
				\only<\SMFBi>{\SMFullBleedImage{\SMFBimg}}%
			\fi
		\else
			\ifnum\SMFBi=\SMFBtotal
				\SMFullBleedImage{\SMFBimg}%
			\fi
		\fi
	}%
}

% sideimageframe: a standardframe that always carries one auto-placed
% "side filler image" (the user's own term, see
% slides/dfc-Angel-Investors-current.tex's own comment for the technique
% this formalizes into a reusable frame type) -- 2026-07-06, direct
% instruction. The image is positioned flush to three edges at once: its
% top touches the title bar's own bottom edge, its right touches the
% page's own right edge, and its bottom touches the footer box's own top
% edge -- so it's scaled by *height* to exactly
% \paperheight-\SMTitleBarBottomAlign-\SMFooterBoxHeight (the
% measured gap between those two), with `keepaspectratio` so the width
% simply follows from the image's own aspect ratio rather than being
% forced/distorted to some second, independent target -- anchoring the
% node at the page's own north-east corner then makes the top+right
% touches automatic, and the height exactly filling that gap makes the
% bottom touch automatic too, with no separate width constraint needed
% (only a *square* source image would ever visibly reach some further
% left boundary; this only guarantees top/right/bottom, matching what was
% actually asked for).
%
%   \begin{sideimageframe}{Title}{media/photo-example}
%   	\begin{itemize}
%   		\item ...
%   	\end{itemize}
%   	\source{...}
%   \end{sideimageframe}
%
% \sidefootnotes is switched on automatically (not left to the caller) --
% the whole point of this frame type is that its footnotes always end up
% rotated along the left edge of the image rather than colliding with it
% the way a full-width bottom footnote would (see
% slides/u29-Implications-of-Capital-Market-Efficiency-Prices-current.tex's
% own comment for exactly that collision, on a plain standardframe). Unlike
% every other \sidefootnotes use elsewhere in this deck, though, the strip
% can't sit at its usual fixed inset from the page's own right edge here --
% that would run it directly underneath the image, the same problem u29
% had. Fixed by *computing* the image's own rendered width (via
% \settowidth against the exact same \includegraphics call used to draw
% it) and using that to override \SMSideFootnoteRightInset for just
% this one frame -- moving the whole rotated strip to sit immediately left
% of the image instead of at the page's raw edge, entirely automatically,
% with no manual x-coordinate needed from the caller (the direct
% instruction offered a manual-x-coordinate fallback "if it is possible
% [to do automatically], that would be great" -- turned out to be
% possible). \SMSideImageFrameHeight/\SMSideImageFrameWidth are
% declared once, globally, and merely *set* (not re-declared) inside the
% environment body, so the frame type can be used more than once per deck;
% the \setlength override of \SMSideFootnoteRightInset is a plain,
% non-\global assignment, which -- like every other frame-local override in
% this file (e.g. slides/u29-...'s own \setbeamercolor{footnote}{fg=white})
% -- resets on its own once \begin{frame}...\end{frame}'s own TeX group
% closes, with no explicit reset needed.
%
% v5.30.5: 3.1pt became 2.1pt. Measured on a rendered slide whose rotated footnote
% text was finally visible (see the colour note at the strip's drawing code), 3.1pt
% left 3.9pt between the ink of the line nearest the image and the image -- the
% text's own depth strut adds ~0.8pt -- against the 2.93pt this was meant to match.
%
% The extra gap added on top of the image's own measured width (2026-07-06,
% direct instruction: "replicate ... the precise margin we have on regular
% slides between footnote and footer box") was originally an unmeasured
% 4pt guess -- replaced with 3.1pt, back-solved from the deck's own real
% reference gap: measured directly on a plain bottom-footnote slide
% (slides/xhe-Active-vs-Passive-Management-2026-First-current.tex's own
% \source{} line) that the real gap between a footnote's own text and the
% footer box's top edge is 2.93pt, while the original 4pt guess here was
% actually landing the rotated text 3.8pt from the image -- 4pt - (3.8-2.93)
% = 3.13pt, rounded to 3.1pt.
\newlength{\SMSideImageFrameHeight}
\newlength{\SMSideImageFrameWidth}
\newenvironment{sideimageframe}[3][t]{%
	\setlength{\SMSideImageFrameHeight}{\dimexpr\paperheight-\SMTitleBarBottomAlign-\SMFooterBoxHeight\relax}%
	\settowidth{\SMSideImageFrameWidth}{\includegraphics[height=\SMSideImageFrameHeight, keepaspectratio]{#3}}%
	\sidefootnotes
	\setlength{\SMSideFootnoteRightInset}{\dimexpr\SMSideImageFrameWidth+2.1pt\relax}%
	\begin{standardframe}[#1]{#2}%
	\vbox to 0pt{\vss
	\begin{tikzpicture}[overlay, remember picture]
		% outer sep=0pt: see the frametitle hook's comment -- without it the
		% image sits 0.2pt in from the page's right edge (and 0.2pt low).
		\node[anchor=north east, inner sep=0pt, outer sep=0pt] at ($(current page.north east)+(0,-\SMTitleBarBottomAlign)$)
			{\includegraphics[height=\SMSideImageFrameHeight, keepaspectratio]{#3}};
	\end{tikzpicture}%
	}%
}{%
	\end{standardframe}%
}

% sectionframe: a section-opening overview (ExampleSlides.pptx slide 3) --
% same layout as standardframe (footer, logo, bullet styling all completely
% unchanged), except the title bar itself is much taller (about a third of
% the page) and set in a bigger font, and the body text is a step bigger
% too, so it reads as "somewhat greater focus" and is unmistakable at a
% glance when flipping through slides -- while still keeping the deck's
% usual yellow-background/black-text look, unlike an earlier version of this
% that inverted to dark-background/yellow-text instead.
%
% The taller bar is NOT a separately-drawn box the way titleframe's is (see
% "Title slide" above): the gotham theme's own frametitle box already sizes
% itself from \SMFrametitleTopPadding/\SMFrametitleBottomPadding plus
% the title font's own strut height (beamerthemeSlideMage.sty, the
% "frametitle" template), so overriding those two lengths locally -- the
% same trick already used below for the font/color overrides -- grows the
% *theme's own* box instead of fighting it with a second one. This also
% means the white-patch-behind-the-logo/chapter-logo machinery (tied to the
% frametitle template via \addtobeamertemplate) keeps working unmodified,
% since it measures the box's actual rendered bottom edge rather than
% assuming any particular height.
%
% The title text is bottom-anchored within the bar -- sitting a fixed "one
% line" above the bar's bottom edge, growing *upward* for a title long
% enough to wrap onto a second line -- rather than top-anchored the way
% gotham's own frametitle template flows by default (which would instead
% leave a large, growing gap *below* a wrapped title, between it and the
% bar's bottom edge). gotham's template has no built-in way to do this: its
% only two knobs are the toppading/bottompading struts, both anchored to
% *one* end of a single top-down paragraph flow. So instead of leaving
% toppading fixed (as an earlier version of this did) the whole title text
% is measured into \SMSectionframeTitleBox first -- built with the same
% font and (approximately) the same usable width gotham's own template
% gives it, so its \ht+\dp closely track the title's real rendered size,
% including wrapping onto extra lines for a long title -- and toppading is
% then computed as "whatever's left" after that measured height and the
% fixed bottom gap are subtracted from the bar's total target height. A
% longer/wrapped title therefore automatically shrinks toppading (pushing
% the whole title block up, growing upward) rather than changing the gap
% below it. \SMSectionframeTitleWidth is a slight underestimate of
% gotham's actual per-line width (which is only narrowed by the logo on
% whichever line ends up next to it, typically just the last one) --
% applying that narrower width to *every* line is a deliberate, harmless
% simplification rather than replicating gotham's per-line logic exactly.
%
% The institute logo moves down along with the title now (it's vertically
% centered against the title's own baseline, gotham's instituteLogo
% template), rather than staying pinned where it sits on every other frame
% type, as an earlier version of this kept it -- an intentional change from
% that version: with the title itself now sitting near the bar's bottom
% edge, keeping the logo up at the bar's top would visually separate it
% from the title it's meant to sit beside.
%
% \begingroup/\endgroup (rather than the \global-set-and-reset dance
% \chapterlogo{}/\sidefootnotes need) is enough here: this only ever touches
% \setbeamercolor/\setbeamerfont/\setlength/\setbox, all of which assign
% locally under the hood (checked against beamerbasecolor.sty/
% beamerbasefont.sty directly for the first two -- their key handlers are
% plain \def, not \gdef -- and \setlength/\setbox are plain TeX assignments,
% local by default), so a plain surrounding TeX group reverts them
% correctly on its own.
%
% \SMSectionframeContentDrop adds a little extra room below the bar
% before the body content starts, on top of standardframe's own
% \SMContentTopShift pull-up (applied inside \begin{standardframe}
% already, before this code regains control). All lengths here calibrated
% against the rendered PDF's own coordinates; re-measure and adjust if the
% title/body font sizes change.
% \SMSectionframeTrueBoxBottom (2026-07-25, direct instruction: "the
% image extends into the title box" on neb/u97 -- followed by "tyn actually
% looks fine, find out what is really driving this issue/difference"):
% \SMSectionframeBoxHeight above is only the *target* the toppading
% formula below solves for -- by construction that formula makes
% toppading+title-height+title-depth+bottompading algebraically equal
% boxHeight+\SMSectionframeHeightCorrection regardless of title content,
% so the *rendered* bar height should be a fixed, title-independent
% 97.14pt... except it measurably isn't (confirmed directly, PyMuPDF vector
% fill): neb/u97 ("Market Environment"/"Household Finance") render at
% 89.19pt, while tyn/vtm ("Capital asset pricing model"/"Measuring and
% managing risk") render at 85.04pt -- a real, reproducible ~4.15pt
% difference driven entirely by the title text itself, not by anything
% neb/u97 did wrong.
%
% Root cause (traced into gotham.dtx's own
% \SM@frametitlestrut@start/@end, not just inferred): the toppading
% formula's own `-\dp\SMSectionframeTitleBox` term assumes the title's
% real descender depth gets reintroduced somewhere else in the box so it
% nets out -- it doesn't. `@start` is a zero-depth `\rule` whose height is
% `toppading+\totalheightof{X}` (a *fixed reference glyph*, not the real
% title), so the line's rendered height is toppading-driven, independent of
% the title's own real height. `@end` is a zero-height `\rule` of depth
% `\SMFrametitleBottomPadding` (19.90pt for sectionframe) -- also fixed,
% not the real title's own depth. The real title text sits between these
% two struts and contributes its own real height/depth to the same TeX
% line, but a line's rendered depth is the *max* of every contributor, not
% a sum -- and 19.90pt comfortably exceeds any real title's own descender
% depth (measured directly: 0.25pt with none, e.g. "Market
% Environment"/"Household Finance"; 4.42pt with any single descender
% present -- g/j/p/q/y, letter/count doesn't matter, e.g. "Capital asset
% pricing model"/"Measuring and managing risk"/a plain "Testing gjpqy"),
% so bottompading always wins there and the title's real depth is never
% actually given back. Net effect: every point of a title's own real
% descender depth shrinks toppading (and therefore the whole visible box)
% by that same amount, with nothing compensating -- confirmed exactly:
% 4.42pt-0.25pt=4.17pt predicted vs. 89.19pt-85.04pt=4.15pt measured, a
% match well within measurement noise. \SMSectionframeHeightCorrection
% (12.1pt) was presumably calibrated against a descender'd test title (it
% lands tyn/vtm almost exactly on the 85.04pt target), leaving every
% no-descender title's own bar ~4.15pt taller than that target instead.
% Full derivation, and why this means \SMSectionframeBoxHeight was
% never actually title-independent, in README.md's own "Section overview
% (sectionframe)" section.
%
% \SMSectionframeTrueBoxBottom below is the *measured* value for a
% no-descender title specifically (matching neb/u97, the only two slides
% using it) -- it is NOT a general replacement for
% \SMSectionframeBoxHeight: a sectionframe whose title has any
% descender (tyn/vtm, confirmed landing on the *other* value, 85.04pt) must
% keep anchoring against \SMSectionframeBoxHeight directly instead
% (already correct for them, so left untouched). Anything new needing the
% bar's own true rendered bottom should re-measure per its own actual title
% text (real vector fill, not a visual guess, and not an assumption that
% either existing constant transfers) rather than reuse this length
% unless its own title also has zero descenders.
% v5.28.0: the bar height no longer depends on the title (see
% \SMSectionframeRefDepth below), so the "true" bottom is the same 85.04 for
% every title; the old 89.19pt value only fitted titles without descenders.
\newlength{\SMSectionframeTrueBoxBottom}
\setlength{\SMSectionframeTrueBoxBottom}{\SMSectionframeBoxHeight}
% (v9.8.0: \SM@sectionbarmetrics below places the title now, from the bar's own material; the
% four profile numbers named in this and the paragraphs above and below, \SMSectionframeRefDepth,
% \SMSectionframeBottomGap, \SMSectionframeBottomPadCorrection and \SMSectionframeHeightCorrection,
% are gone, and so is the formula these paragraphs explain. They stay as the record of why it was
% done this way for ten versions.)
% \SMSectionframeRefDepth (v5.28.0): the depth the toppading formula assumes for
% the title. It used to subtract the title's own measured depth, which nothing
% adds back, so a title without descenders ("Slides", "Decks and deck
% versions") got a bar 4.15pt taller than one with them, and every other depth
% (p, y, ...) a height in between. With this fixed value (the depth of "p",
% 4.42pt) every title gets the 85.04bp bar the calibrated constants below were
% fitted to, and \SMSectionframeBarBottom (where \sectionfigure starts) is that
% same constant.
% \SMSectionframeBottomGap is the *visual* target -- one huge-font line
% height (baselineskip, measured directly from a deliberately 2-line test
% title: the gap between the two lines' own rendered top edges). It is NOT
% used directly as \SMFrametitleBottomPadding: gotham's bottom strut
% (\SM@frametitlestrut@end) consistently produces a visual gap about
% 5.58pt *smaller* than the padding value fed into it (checked empirically,
% same offset for both a one-line and a two-line test title, so it's a
% fixed per-box offset rather than something that scales with content) --
% \SMSectionframeBottomPadCorrection adds that back.
% Similarly, the box's total rendered height comes out consistently ~11.4pt
% *less* than toppading+measuredTitleHeight+bottompading would suggest --
% the measured title box (built as a plain \vbox{...\par}) doesn't reproduce
% gotham's own leading strut (\SM@frametitlestrut@start, a zero-width
% rule with an explicit height rather than natural interline glue) exactly,
% and this correction closes that gap. Checked the same way as the bottom
% one above: consistent across both the one-line and two-line test titles.
\newlength{\SMSectionframeTitleWidth}
\setlength{\SMSectionframeTitleWidth}{\dimexpr\paperwidth-\SMFrametitleLeftPadding-\SMLogoWidth-2\SMFrametitleRightPadding\relax}
\newbox\SMSectionframeTitleBox
\newlength{\SMSectionframeBarBottom}
% The title's place in the bar (v9.8.0: from the title's own material, no longer from four numbers
% measured at Fira): its last baseline lies the depth of a \strut of the title font plus
% \SMSectionframeTitleGap above the bar's bottom edge, and the top padding is what makes the bar
% \SMSectionframeBoxHeight high. The bar's content is typeset once into a box, by the same
% beamercolorbox the frame title template uses and with a top padding as large as the bar; the box's
% height is that padding plus a constant from there on (any title, any number of lines), so the
% padding that gives the bar its height follows. #1 is the title.
\makeatletter
\newbox\SM@sectionprobe
\newcommand{\SM@sectionbarmetrics}[1]{%
	\begingroup\usebeamerfont{frametitle}\setbox\z@\hbox{\strut}\setlength{\dimen0}{\dimexpr\dp\z@+\SMSectionframeTitleGap\relax}%
	\xdef\SM@sectionbottom{\the\dimen0}\endgroup
	\setlength{\SMFrametitleBottomPadding}{\SM@sectionbottom}%
	\setlength{\SMFrametitleTopPadding}{\SMSectionframeBoxHeight}%
	\setbox\SM@sectionprobe=\vbox{\usebeamerfont{frametitle}%
		\begin{beamercolorbox}[wd=\paperwidth, sep=0pt, leftskip=\SMFrametitleLeftPadding,
			rightskip=\dimexpr\SMLogoWidth+2\SMFrametitleRightPadding\relax plus 1fil]{frametitle}%
			\SM@frametitlestrut@start #1\SM@frametitlestrut@end
		\end{beamercolorbox}}%
	\setlength{\SMFrametitleTopPadding}{\dimexpr2\SMSectionframeBoxHeight-\ht\SM@sectionprobe-\dp\SM@sectionprobe\relax}}
\makeatother
\newenvironment{sectionframe}[2][t]{%
	\begingroup
	\SMManualTopMargin% own title-bar geometry: no automatic Top Margin
	\global\csname SM@topsectiontrue\endcsname% its taller bar: not counted for \SMContentTopShift
	% Disable hyphenation for the title -- both here, in the measurement
	% box below, and (since this is in effect for the rest of the group,
	% which \begin{standardframe} -- triggering the real \insertframetitle
	% -- runs inside) in the actual rendered title too, so a line break
	% never lands mid-word. Setting the penalties alone is the standard,
	% sufficient way to do this (TeX treats a penalty of 10000 at a
	% breakpoint as forbidden, not just heavily discouraged, so it's never
	% chosen even when the alternative is an overfull line); the hyphenmin
	% lengths are redundant insurance, not strictly needed on top of that.
	\hyphenpenalty=10000
	\exhyphenpenalty=10000
	\lefthyphenmin=100
	\righthyphenmin=100
	\setbeamerfont{frametitle}{size=\huge}
	\csname SM@sectionbarmetrics\endcsname{#2}%
	\global\def\SMFrametitleFullHeightYellow{1}%
	\global\def\SMFrametitleLogoFixed{1}%
	\setbeamertemplate{frametitle}[SMNoInlineLogo]%
	\setlength{\SMSectionframeBarBottom}{\dimexpr\SMSectionframeBoxHeight+0.12bp\relax}% 85.14bp for the default bar: the figure starts 0.12bp under it
	% \sectionfigure sets the side-footnote lengths \global (see there); the
	% values from before this frame come back at \end{sectionframe}.
	\edef\SMSectionframeRestoreSideFootnotes{%
		\global\SMSideFootnoteRightInset=\the\SMSideFootnoteRightInset\relax
		\global\SMSideFootnoteHsize=\the\SMSideFootnoteHsize\relax}%
	\begin{standardframe}[#1]{#2}%
	\vspace*{\SMSectionframeContentDrop}%
	\large
}{%
	\end{standardframe}%
	\SMSectionframeRestoreSideFootnotes
	\endgroup
}

% \sectionfigure[includegraphics keys]{image} -- a sectionframe's side image
% (2026-10-01, direct instruction: bij's image should "fit snugly between
% header box and footer and touch the right edge of the slide", for section
% slides with images in general, with the image's source visible). Put it
% anywhere in the sectionframe body. The image is scaled by height to fill
% the gap from the title bar's bottom edge to the footer box's top edge and
% sits flush with the page's right edge -- sideimageframe's geometry, but
% for sectionframe's taller bar. That bar's bottom edge depends on the
% title: the theme's strut gives the last line a fixed depth of ~4.4pt,
% but the padding subtracts the title's real depth -- so the bottom edge is
% 89.44bp minus that depth, but never above 85.04bp (a descender deeper
% than the strut deepens the bar again). Measured in the PDF, i.e. in bp:
% 89.19bp for "Market Environment" (depth 0.25pt), 86.01bp for "Market,
% Environment" (3.44pt), 85.04bp for "Capital asset pricing model"
% (4.42pt) and "Trends in publishing" (4.73pt). (The two
% constants \SMSectionframeTrueBoxBottom/\SMSectionframeBoxHeight
% hold the same numbers in TeX pt, 0.3pt short -- neb/u97 used to add
% that back by hand.) sectionframe computes it as
% \SMSectionframeBarBottom. (Not measured
% from the SMTitleBottom node: the title bar is typeset after the body,
% so that node still belongs to the previous frame while the body runs.)
% The frame's footnotes -- the automatic "Image source(s):" line included --
% go into the rotated side strip immediately left of the image, as on
% sideimageframe. The strip's inset is read at shipout, after the frame body,
% so it and the strip length are set \global here; \end{sectionframe} puts
% the previous values back. Replaces the per-slide \settowidth/
% \SMSideFootnoteRightInset/\SMSideFootnoteHsize set-up of neb, u97,
% tyn and vtm.
\newlength{\SMSectionFigureHeight}
\newlength{\SMSectionFigureWidth}
\newcommand{\sectionfigure}[2][]{%
	\setlength{\SMSectionFigureHeight}{\dimexpr\paperheight-\SMSectionframeBarBottom-\SMFooterBoxHeight\relax}%
	\settowidth{\SMSectionFigureWidth}{\includegraphics[height=\SMSectionFigureHeight, keepaspectratio, #1]{#2}}%
	\vbox to 0pt{\vss
	\begin{tikzpicture}[overlay, remember picture]
		\node[anchor=north east, inner sep=0pt, outer sep=0pt] at ($(current page.north east)+(0,-\SMSectionframeBarBottom)$)
			{\includegraphics[height=\SMSectionFigureHeight, keepaspectratio, #1]{#2}};
	\end{tikzpicture}%
	}%
	\global\SMSideFootnotestrue
	% 3.1pt: sideimageframe's measured gap between strip and image.
	% 2.31pt: the clearance above the footer box that \SMSideFootnoteHsize's
	% 207pt default leaves under a standard title bar.
	\global\SMSideFootnoteRightInset=\dimexpr\SMSectionFigureWidth+3.1pt\relax
	\global\SMSideFootnoteHsize=\dimexpr\SMSectionFigureHeight-\SMSideFootnoteTopSep-2.31pt\relax
}

% quoteframe: a quote slide (ExampleSlides.pptx slide 4) -- institute logo
% as usual, and the quote itself set large/italic/black on a yellow
% full-bleed block, same style as titleframe/sectionframe's own title boxes
% (flush top-left, stopping short of the logo on the right, same margin as
% elsewhere) rather than the dark block an earlier version of this used. The
% usual footer still shows (unlike titleframe, which drops it -- a title
% slide is a one-off outside any particular course content, while a quote
% is still one slide *among* a deck's others, worth being able to reference
% by its footer slide ID like any other). Opening/closing curly quotes are
% added automatically (write the quote text with none of its own); the
% second argument is the person it's from, set smaller/plain, below the
% yellow box and right-aligned to its own right edge (not the page's usual
% text margin) -- pass an empty second argument for an unattributed quote.
%
%   \quoteframe{One of the funny things about the stock market is that
%   	every time one person buys, another sells, and both think they
%   	are astute.}{William Feather, 1889-1981}
%
% \SMQuoteBoxHeight is a plain percentage of the page, unlike
% titleframe's/sectionframe's own (calibrated-to-land-on-an-exact-fraction)
% box heights, since nothing here needs to line up with another element's
% measured position the way those do. \SMTitleframeBoxWidth (defined
% above) is reused as-is for the box's width -- the same
% stop-short-of-the-logo formula, so no reason to duplicate it under a new
% name.
%
% Every piece here -- box, logo, quote text, attribution -- is drawn as a
% zero-space tikz overlay anchored to absolute page coordinates (the same
% \vbox to 0pt{\vss ...} pattern used elsewhere in this file for exactly
% this reason), rather than any of it flowing through the frame's normal
% vertical list. This sidesteps needing to know/counteract beamer's own
% default top-of-frame spacing (which standardframe explicitly measures and
% cancels via \SMContentTopShift) -- there's no such frame-relative
% positioning to get right here, since literally everything is placed by
% explicit coordinate instead. Like the other page-corner overlays in this
% file, "current page" doesn't resolve correctly on a document's very first
% compile pass, so (like those) this needs the usual second pass too.
%
% A *command*, not an environment -- see the identically-reasoned comment on
% titleframe just above for why: the same beamercolorbox-spanning-a-custom-
% environment's-begin/end problem would apply here too, for the same reason.
% \@empty (the empty-attribution check below) needs \makeatletter.
% \quotesource{text} -- opt-in \source{} for \quoteframe (call *before*
% \quoteframe{}{}, same set-then-consume pattern as \chapterlogo{}
% above), needed because \quoteframe is a *command* that auto-closes its
% own \end{frame} -- a bare \source{...} typed after \quoteframe{}{} in
% the calling slide file (as if it were still inside the frame, the way
% \source{} reads everywhere else in this deck) is actually issued
% *outside* any frame, and \source{}'s own \footnote{} then has no page to
% attach to yet. Confirmed directly (2026-07-12, direct instruction:
% traced a spurious extra blank leading overlay step on
% slides/c4c-Macbeth-Spot-Removers-Proposed-Structure-current.tex -- the
% very next multi-overlay frame after
% slides/w8w-Pie-Model-of-Capital-Structure-current.tex's own bare
% \source{} -- down to that orphaned footnote specifically, via a reduced
% two-slide test deck): the dangling footnote insert doesn't render
% anywhere visible on w8w itself, but silently corrupts beamer's own
% overlay-counting pre-pass on whichever frame happens to come right
% after -- reproduced with w8w+c4c alone (bug present) vs. c4c alone (bug
% absent), and confirmed it's specifically about the footnote being
% orphaned, not about w8w's content otherwise, by removing just its
% trailing \source{} line (bug also disappears). A plain \clearpage
% between the two frames also "fixes" c4c but only by giving the orphaned
% footnote a page of its own to attach to -- confirmed directly this
% splits w8w itself into two pages (real content, then a second page with
% nothing but the footnote/footer), not an acceptable trade. The only
% correct fix is keeping the footnote inside the frame that's actually
% citing something, i.e. before \quoteframe's own \end{frame} -- exactly
% what \quotesource{} below does.
\newcommand{\SMQuoteSource}{}
\newcommand{\quotesource}[1]{\gdef\SMQuoteSource{#1}}
\makeatletter
\newcommand{\quoteframe}[2]{%
	\SMOpenFrame{[t]}%
	% Disable hyphenation for the quote text (2026-07-26, direct
	% instruction: "do not allow words to be broken across lines" --
	% confirmed directly first, e.g. on `n3q`'s own long quote,
	% "limi-tations"/"finan-cial" both landed as real mid-word line
	% breaks). Same standard, sufficient technique already used for
	% `sectionframe`'s own title (see that environment's own comment,
	% right before \SMSectionframeTitleBox) -- setting the penalties
	% to 10000 makes TeX treat a hyphenated break as forbidden rather than
	% just discouraged, so one is never chosen even over an otherwise
	% ragged/short line; the hyphenmin lengths are redundant insurance on
	% top of that. Scoped to this frame only via the enclosing
	% \begin{frame}/\end{frame} group, same as sectionframe's own
	% \begingroup/\endgroup scoping -- doesn't leak into any other frame
	% type.
	\hyphenpenalty=10000
	\exhyphenpenalty=10000
	\lefthyphenmin=100
	\righthyphenmin=100
	% \SMArtifactBegin/\SMArtifactEnd (defined with the footer, "PDF
	% accessibility" comment): same reasoning as titleframe's own background
	% box above -- purely decorative, marked so a screen reader skips it. The
	% quote text and attribution below stay unmarked (real content); the
	% footer (this frame keeps one, unlike titleframe) is already marked via
	% the global footline hook.
	\SMArtifactBegin{/Type /Background}%
	\vbox to 0pt{\vss
	\begin{tikzpicture}[overlay, remember picture]
		% \SMTitleBoxColor, not a literal color token -- same setting
		% titleframe's own box and the standard frametitle bar (colorA) use
		% (2026-07-16, direct instruction).
		% outer sep=0pt here and on the logo/attribution below: same hairline-
		% inset-from-the-anchored-corner issue as titleframe's own box -- see
		% the frametitle hook's comment for the full story.
		\node[fill=\SMTitleBoxColor, anchor=north west, inner sep=0pt, outer sep=0pt,
		      minimum width=\SMTitleframeBoxWidth,
		      minimum height=\SMQuoteBoxHeight]
			at (current page.north west) {};
	\end{tikzpicture}%
	}%
	\SMArtifactEnd
	\vbox to 0pt{\vss
	\begin{tikzpicture}[overlay, remember picture]
		% inner ysep derived, not \SMFrametitleRightPadding -- see the
		% frametitle hook's fixed-position logo redraw for why (matches the
		% inline logo's vertically-centered position on every other frame).
		\node[anchor=north east, inner xsep=\SMFrametitleRightPadding,
		      inner ysep=\dimexpr\SMTitleBarBottomAlign-\SMLogoHeight\relax,
		      outer sep=0pt]
			at (current page.north east) {\SMInstituteLogo};
	\end{tikzpicture}%
	}%
	\vbox to 0pt{\vss
	\begin{tikzpicture}[overlay, remember picture]
		% \SMTitleTextColor, not a literal color token -- follows the box
		% color change right above (2026-07-16, direct instruction). Set
		% *inside* the node's own {...} content group (not a stored box
		% register, unlike titleframe's own vbox above), so this doesn't hit
		% the same "\color leaks out of a stored \vbox" pitfall documented on
		% \titleframe's own closing code -- confirmed directly by checking
		% the frame right after this one for the same symptom.
		% The quote text is set in a \parbox instead of the node's own
		% `text width` (2026-09-24), so #1 can hold reveal steps
		% (\conditionalpause, \uncover, \pause). Two separate failures
		% with those directly inside the node, both reproduced in plain
		% beamer and TikZ (no project code): (1) TikZ reads a `text width`
		% node's text through an \aftergroup hook that a \pause/\uncover
		% inside it breaks -- "Package tikz Error: Giving up on this path.
		% Did you forget a semicolon?" on every overlay page; (2) even
		% with (1) avoided, a text-less \pause starts a "covered" state
		% that is only ended by a *global* \beamer@endpause at the next
		% \pause or the end of the frame, so everything after the quote --
		% the footer included -- stayed covered until the last step.
		% \beamer@endpause is therefore run at the end of the quote (and
		% cleared), which confines the covering to the quote itself.
		%
		% The parbox reproduces what TikZ builds for `text width`
		% (minipage[t], \leavevmode\hbox{}, then \tikz@text@action --
		% ragged-right with its own \rightskip/\spaceskip -- and the
		% paragraph ended *after* the \Large group, i.e. with the normal
		% line spacing) so that existing quote slides are unchanged:
		% verified pixel-identical on all 18 other \quoteframe slides.
		\node[anchor=center, inner sep=0pt]
			at ($(current page.north west)+(0.5\SMTitleframeBoxWidth,-0.5\SMQuoteBoxHeight)$)
			{\parbox[t]{\dimexpr\SMTitleframeBoxWidth-2\SMFrametitleLeftPadding\relax}{%
				\leavevmode\hbox{}\tikz@textfont\tikz@text@action%
				{\color{\SMTitleTextColor}\Large\itshape\textquotedblleft #1\textquotedblright%
				\beamer@endpause\global\let\beamer@endpause\@empty}\par}};
	\end{tikzpicture}%
	}%
	\def\SMQuoteAttributionTemp{#2}%
	\ifx\SMQuoteAttributionTemp\@empty\else
		\vbox to 0pt{\vss
		\begin{tikzpicture}[overlay, remember picture]
			\node[anchor=north east, inner sep=0pt, outer sep=0pt]
				at ($(current page.north west)+(\SMTitleframeBoxWidth,-\SMQuoteBoxHeight-\SMQuoteAttributionGap)$)
				{\normalfont\itshape\color{c_custom_gray} #2};
		\end{tikzpicture}%
		}%
	\fi
	\ifx\SMQuoteSource\@empty\else
		\source{\SMQuoteSource}%
	\fi
	\end{frame}%
	% Cleared here, *after* \end{frame}, not inside the frame body above
	% (2026-09-24): beamer re-runs the collected body once per reveal step,
	% so clearing it in the body left only the first step with a source
	% line -- the quote's \quotesource vanished on every later step of a
	% quote with \conditionalpause. Code after \end{frame} runs once.
	\global\let\SMQuoteSource\@empty
}
\makeatother

% figureframe: a slide that is just a figure (ExampleSlides.pptx slide 5) --
% same standardframe layout (title, footer, \source{} for the usual
% "Source: ..."/"Image source: ..." line), but with the figure positioned
% via \figurecenter (call before the figure, inside the frame body) rather
% than left-aligned, which -- like the source deck's own figure slide -- is
% the default:
%
%   \begin{figureframe}{Correlation between DJIA stock returns over 15 years}
%   	\includegraphics[width=0.5\linewidth]{correlation-matrix}
%   	\source{Source: Own illustration, data from Yahoo! Finance, 24.02.2026.}
%   \end{figureframe}
%
%   \begin{figureframe}{Title}
%   	\figurecenter
%   	\includegraphics[width=0.6\linewidth]{diagram}
%   \end{figureframe}
%
% Vertical placement is standardframe's own [t]/[c] (its existing optional
% argument, passed straight through as figureframe's own -- so
% \begin{figureframe}[c]{Title} vertically centers exactly like
% \begin{standardframe}[c]{Title} already does); \figurecenter only ever
% affects the horizontal axis.
%
% \figurecenter is plain \centering under a more discoverable name, nothing
% more -- *not* an opt-in-before-the-frame flag like \chapterlogo{}/
% \sidefootnotes (those must run before \begin{standardframe}, since they
% configure something the frame's own opening code reads). This is called
% *inside* the body instead, exactly where it takes effect, so it can just
% take effect immediately like any other formatting declaration -- no flag,
% no state to reset afterwards, and (checked by first writing it as a flag
% read in figureframe's own opening code, mirroring \sidefootnotes) no way
% for the frame's opening code to see it anyway, since that code runs
% before the body -- and \figurecenter -- does.
\newcommand{\figurecenter}{\centering}
\newenvironment{figureframe}[2][t]{%
	\begin{standardframe}[#1]{#2}%
}{%
	\end{standardframe}%
}

% highlightframe: a standardframe with two things about it made unmistakably
% different -- the title bar's *and* the footer's background color, both
% together -- so a presenter can never flip to it without immediately
% noticing (ExampleSlides.pptx slide 6 was the inspiration for wanting a
% standardframe variant here, though the source slide itself has no such
% color cue; this is a new addition, not a conversion). Everything else --
% body, logo, bullet styling -- is exactly standardframe's own, via the same
% \begingroup/\endgroup wrapper sectionframe above uses for the title-color
% half of this. \SMHighlightColor is a \colorlet, not a color baked
% directly into the environment, so retargeting every highlightframe in a
% deck to a different color -- title bar and footer both -- is still the
% same one-line change it was when this only touched the title bar.
%
% The footer box needs the \global\def dance (\SMFooterBoxColor, defined
% above with the frametitle hook's other opt-in flags, reset in
% standardframe's own closing code) that \setbeamercolor{frametitle} doesn't:
% that assigns through beamer's own color-management system, which handles
% the frame-closing group correctly on its own (checked -- see
% \SMChapterLogoPath's comment for why plain macros need \global\def but
% beamer's own set-commands don't), but the footer box's fill is read from a
% plain macro inside a tikz key, same category as \SMChapterLogoPath.
\colorlet{SMHighlightColor}{c_primary_gray}
\newenvironment{highlightframe}[2][t]{%
	\begingroup
	\setbeamercolor{frametitle}{bg=SMHighlightColor}
	\global\def\SMFooterBoxColor{SMHighlightColor}%
	\begin{standardframe}[#1]{#2}%
}{%
	\end{standardframe}%
	\endgroup
}


% The strip is an overlay (a zero-size box: the Top Margin measurement does not see it), so
% the body below it starts where the title bar's body would -- under the strip. The frame's
% own target line is therefore the Top Margin plus the strip's height
% \SMProgressStripHeight (the strip's bottom edge minus the title bar's: measured 11.08pt in
% the PDF; the node is 11.19pt minimum height): the first line sits as far below the strip
% as it sits below the title bar on any other frame, whatever the body starts with (v5.30.5; the
% frame used to be exempt from the measurement, with a fixed -3pt, which left a paragraph
% 10pt and a list 14pt below the strip against the usual 16pt -- direct instruction, slide 1j2).
\makeatletter
\newcommand{\SM@topstrip}{%
	\ifluatex
		\directlua{slidemage_measure.top_extra(\SM@topthis, \strip@pt\dimexpr\SMProgressStripHeight\relax)}%
	\fi
}
\makeatother
% \progressstrip{Step 1,Step 2,Step 3}{k} -- a full-width strip in the
% institution's tertiary yellow, flush below the title bar, that lists a
% short chain of steps joined by arrows and underlines the k-th one to mark
% where the slide sits in that chain (added 2026-09-19 for the Publishing
% masterdeck's "Field -> Topic -> Question" strip, which the source
% PowerPoint's "Progress" layout draws directly below the title bar; revised
% the same day, direct instruction: edge-to-edge, flush under the title bar,
% smaller font, tight vertical padding, secondary yellow). Written in the
% frame body, as the very first thing. The strip itself is a zero-space
% absolute overlay (so it can span the whole page width, which in-flow text
% cannot); the \vspace after it reserves its height so the body starts just
% below it.
\newcommand{\progressstrip}[2]{%
	\csname SM@topstrip\endcsname% the text below the strip: Top Margin + the strip's height (measured, like any frame)
	\vbox to 0pt{\vss
	\begin{tikzpicture}[overlay, remember picture]
		\node[fill=c_tertiary_yellow, anchor=north west, outer sep=0pt,
		      inner xsep=\SMFrametitleLeftPadding, inner ysep=1.5pt, text height=6.2pt, text depth=2pt,
		      minimum height=11.19pt,
		      text width=\dimexpr\paperwidth-2\SMFrametitleLeftPadding\relax,
		      align=left, text=c_primary_dark, font=\fontsize{6.3}{7.5}\selectfont]
			at ($(current page.north west)+(0,-\SMTitleBarBottomAlign)$)
			{\foreach \SMStepName [count=\SMStepIndex] in {#1}{%
				\ifnum\SMStepIndex>1\ $\rightarrow$\ \fi
				\ifnum\SMStepIndex=#2\underline{\SMStepName}\else\SMStepName\fi}};
	\end{tikzpicture}%
	}%
	\vspace{-3pt}%
}

% \SMCodeFile[listings keys]{path} -- LaTeX source shown on a slide, read from
% a file (2026-10-04, for the SlideMage manual: "show the result and the
% source together"). The code lives in its own small file, e.g.
% snippets/<slide id>-1.tex at the project root, so the slide itself needs no
% [fragile] frame and the code is exactly what was compiled elsewhere. Found on
% the TeX search path like any \input (the project root is on it, see
% .latexmkrc). Optional keys go to \lstinputlisting (firstline=, lastline=,
% basicstyle=...). Fira Mono Regular (the theme's \ttfamily is Fira Mono Bold),
% control sequences, comments and the panel in the SMCode* colours below
% (set per template). Write snippets in ASCII; the few non-ASCII characters below are mapped.
% SlideMage's transcript shows the file as a code block, and the pre-render
% check reports a missing file.
\usepackage{listings}
% Every control sequence, backslash included, is a keyword: the backslash is
% made a letter and is the keyword prefix (listings' prefix keywords need a
% non-empty keyword list to be active, hence the dummy). An escaped \% and a
% \\ are mapped literally (\% must not start a comment; \\ otherwise loses a
% backslash to the \% mapping) and reset the prefix flag, which would
% otherwise colour the word after them.
%
% Code colours (2026-10-05, direct instruction: "set them to something more
% readable in the institutions' template files"): three named colours a
% template defines before \input{SlideMage_base} --
% SMCodeKeywordColor (control sequences), SMCodeCommentColor (% comments),
% SMCodeBackgroundColor (the panel), SMCodeTextColor (ordinary code). A template that doesn't keeps the
% original look: highlight colour, primary gray, tertiary gray -- which in
% the Graz/MCI/Innsbruck palettes is yellow/orange and light gray on light
% gray, about 1.1-1.3:1 contrast; their templates set readable colours.
\makeatletter
\@ifundefinedcolor{SMCodeKeywordColor}{\colorlet{SMCodeKeywordColor}{c_primary_highlight}}{}
\@ifundefinedcolor{SMCodeCommentColor}{\colorlet{SMCodeCommentColor}{c_primary_gray}}{}
\@ifundefinedcolor{SMCodeBackgroundColor}{\colorlet{SMCodeBackgroundColor}{c_tertiary_gray}}{}
\@ifundefinedcolor{SMCodeTextColor}{\colorlet{SMCodeTextColor}{c_primary_dark}}{}
\newcommand{\SMCodePercent}{\textbackslash\%\global\let\lst@prefixkeyword\relax}
\newcommand{\SMCodeNewline}{\textcolor{SMCodeKeywordColor}{\textbackslash\textbackslash}\global\let\lst@prefixkeyword\relax}
\makeatother
\lstdefinestyle{SMCode}{
	basicstyle=\SMCodeFont\SMCodeSize\color{SMCodeTextColor},
	columns=fullflexible, keepspaces=true, tabsize=2,
	keywords={SMCodeNoSuchKeyword},
	alsoletter={\\}, keywordsprefix=\\,
	keywordstyle=\color{SMCodeKeywordColor},
	morecomment=[l]\%, commentstyle=\color{SMCodeCommentColor},
	showstringspaces=false, upquote=true,
	breaklines=true, breakindent=1.5em,
	aboveskip=0pt, belowskip=0pt,
	literate={\\\%}{{\SMCodePercent}}2 {\\\\}{{\SMCodeNewline}}2 {ä}{{\"a}}1 {ö}{{\"o}}1 {ü}{{\"u}}1 {Ä}{{\"A}}1 {Ö}{{\"O}}1 {Ü}{{\"U}}1
		{ß}{{\ss}}1 {é}{{\'e}}1 {–}{{--}}1 {—}{{---}}1 {·}{{$\cdot$}}1 {…}{{...}}1
		{→}{{$\rightarrow$}}1,
}
% The panel is a \colorbox around the listing, not listings' own frame and
% background: listings draws its frame outside the box it typesets, so a
% panel stuck out above the top of its column by 4-8pt (depending on the font
% size) and missed the Top Margin (v5.14.1, found by check_layout.py in the
% SlideMage manual). The box holds everything it draws.
\newcommand{\SMCodeFile}[2][]{\par\noindent{\setlength{\fboxsep}{3pt}%
	\colorbox{SMCodeBackgroundColor}{\begin{minipage}{\dimexpr\linewidth-2\fboxsep\relax}%
		\lstinputlisting[style=SMCode,#1]{#2}%
	\end{minipage}}}\par}

% \SMCode{text} / \SMCmd{name} -- code in running text: Fira Mono Regular at
% 90% (\texttt is the theme's Fira Mono Bold, which reads heavy and wide in a
% sentence of Fira Sans Light). \SMCmd{slide} prints \slide. The argument is
% typeset, so write \_ \{ \} \# as in any text. The transcript shows both as
% `code`. Ligatures=TeXOff: code shows -- and `` as typed, not as an en dash
% and a curly quote (fontspec's default for \newfontfamily).
\newcommand{\SMInlineCodeStyle}{}
\newcommand{\SMCode}[1]{{\SMInlineCodeFont\SMInlineCodeStyle #1}}
\newcommand{\SMCmd}[1]{{\SMInlineCodeFont\SMInlineCodeStyle\textbackslash #1}}

% \SMQRCode[height]{image}{caption} -- a QR code SlideMage made (2026-10-25,
% direct instruction: the Source view's Insert QR, the assistant's
% make_qr_code; src/slidemage/qr_code.py), the code at [height] (default
% 0.6\textheight: a lecture hall's back row needs about that) with an
% optional caption centred under it, as wide as the code -- the link, or a
% text such as "QR 1" that the slide refers to; {} for none. An inline box:
% several codes go side by side or in columns, and
% \cornerfigure{\SMQRCode{...}{...}} puts one in a corner. The caption's
% look is \SMQRCodeCaptionStyle, which a template may redefine.
\newsavebox{\SMQRCodeBox}
\providecommand{\SMQRCodeCaptionStyle}{\scriptsize}
\NewDocumentCommand{\SMQRCode}{O{0.6\textheight} m m}{%
	\begingroup
	\sbox{\SMQRCodeBox}{\includegraphics[height=#1]{#2}}%
	\begin{minipage}[t]{\wd\SMQRCodeBox}%
		\centering\usebox{\SMQRCodeBox}%
		\if\relax\detokenize{#3}\relax\else
			\par{\SMQRCodeCaptionStyle #3\par}%     the code's own white border is the gap
		\fi
	\end{minipage}%
	\endgroup
}

% Navigation for readers (2026-10-04, SlideMage manual: readers must find
% their way without knowing that the logo is a hidden link). Three opt-in
% pieces, each usable in any deck:
%
% \SMChapterAgenda[width] -- the slides of the current \section as a
%   two-column list, each entry a link to that slide (its frame title, or the
%   slide's \SMAgendaTitle{...}, or else its slug with spaces). Meant for the
%   body of a chapter's sectionframe (the manual puts one on every divider).
%   [width] (default \linewidth) leaves room for e.g. a \sectionfigure. Set in
%   \normalsize, else \small, else \footnotesize, the first that fits
%   \SMChapterAgendaMaxHeight.
%   Works from the .aux file like a table of contents: every slide writes
%   \SMAgendaSlide{section}{id}{label} (in \setSlideID) and its frame title
%   (\SMAgendaEntryTitle, in the frametitle hook), and the agenda shows what
%   the previous pass recorded -- latexmk's rerun fills it in. A slide shown
%   twice (\slidestep) is listed once, at its first place; the agenda's own
%   slide is not listed.
% \SMSubsectionAgenda[width] -- the same, but one entry per \subsection of
%   the current \section (2026-10-12): the subsection's title, linked to its
%   first slide in this version. For a chapter divider that introduces the
%   subsections instead of listing every slide. Slides before the chapter's
%   first \subsection are not listed; it fills in on the rerun.
% \SMAgendaTitle{text} -- in a slide file, before the frame: the entry text
%   for a slide whose frame has no title (\quoteframe, \fullbleedframe, ...)
%   or should read differently in the agenda.
% \SMDeckAgenda -- marks the deck's contents slide (put it in that slide's
%   body, next to \tableofcontents) as the target of the navigation bar's
%   first link.
% \SMNavigationBar[deck,chapter,previews] -- in the deck preamble: the footer
%   then reads "\FooterText | Navigation: <deck> · <chapter> · <previews>",
%   three links to the \SMDeckAgenda slide, the current chapter's
%   \SMChapterAgenda (the contents slide before the first \section, or where
%   a chapter has none) and the generated appendix agenda with the slide
%   previews. A link whose target is not in this PDF is left out: the last
%   one exists only when SlideMage generates the appendix (the Generate
%   appendix TOC checkbox, or \deckappendixtoc in the deck). Default labels
%   "Contents,Chapter contents,Slide previews"; the word before them is
%   \SMNavigationLead ("Navigation:"). Off by default: other decks' footers
%   and the logo link are unchanged.
% \deckappendixtoc -- in the deck preamble, inert for LaTeX (like
%   \decktags): SlideMage then generates the appendix for this deck on every
%   whole-deck compile, whatever the toolbar checkbox says.
\newcommand{\deckappendixtoc}{}
\newlength{\SMChapterAgendaMaxHeight}
\newcommand{\SMNavigationLead}{Navigation:}
\newcommand{\SMNavigationLabelDeck}{Contents}
\newcommand{\SMNavigationLabelChapter}{Chapter contents}
\newcommand{\SMNavigationLabelPreviews}{Slide previews}
\makeatletter
% -- recording (writes to the .aux) --
\newcommand{\SM@agenda@slide}[2]{%
	\ifcsname SM@ag@written@#1\endcsname\else
		\expandafter\gdef\csname SM@ag@written@#1\endcsname{}%
		% label from the rest of the id-slug ("EN-Slide-IDs-"): language
		% segment and trailing hyphen dropped, hyphens to spaces. (No % in
		% the Lua code -- TeX would read it as a comment -- and no hash
		% sign, which is a macro parameter here.)
		\edef\SM@ag@label{\directlua{
			local s = "\luaescapestring{#2}"
			while string.sub(s, -1) == "-" do s = string.sub(s, 1, -2) end
			for _, p in ipairs({"EN-", "DE-", "mult-"}) do
				if string.sub(s, 1, string.len(p)) == p then s = string.sub(s, string.len(p) + 1) end
			end
			s = string.gsub(s, "[-]", " ")
			tex.sprint(-2, s)}}%
		\immediate\write\@auxout{\string\SMAgendaSlide{\the\c@section}{#1}{\SM@ag@label}}%
	\fi
	% the first slide of each \subsection and the subsection's title, for
	% \SMSubsectionAgenda (its own guard: a slide shown twice must not hide
	% that it also opens a later subsection)
	\ifnum\c@subsection>0
		\ifcsname SM@ag@subwritten@\the\c@section.\the\c@subsection\endcsname\else
			\expandafter\gdef\csname SM@ag@subwritten@\the\c@section.\the\c@subsection\endcsname{}%
			\immediate\write\@auxout{\string\SMAgendaSubsection{\the\c@section}{#1}{\detokenize\expandafter{\subsecname}}}%
		\fi
	\fi
}
\newcommand{\SM@agenda@title}[1]{%
	\ifcsname SM@ag@titled@\slideID\endcsname\else
		\expandafter\gdef\csname SM@ag@titled@\slideID\endcsname{}%
		\immediate\write\@auxout{\string\SMAgendaEntryTitle{\slideID}{\detokenize{#1}}}%
	\fi
}
\newcommand{\SMAgendaTitle}[1]{\SM@agenda@title{#1}}
\newcommand{\SM@agenda@frametitle}{%
	\ifx\beamer@frametitle\@empty\else\expandafter\SM@agenda@title\expandafter{\beamer@frametitle}\fi
}
\newcommand{\SMAgendaFrametitleHook}{\SM@agenda@frametitle}% called by \SMFrametitleHookCode
\newcommand{\SMDeckAgenda}{%
	\hypertarget{SMDeckAgenda}{}%
	\immediate\write\@auxout{\string\SMAgendaDeckTarget}%
}
% -- reading (the .aux, at \begin{document}; its second reading at
% \end{document} changes nothing) --
\newcommand{\SMAgendaSlide}[3]{%
	\ifcsname SM@ag@seen@#2\endcsname\else
		\expandafter\gdef\csname SM@ag@seen@#2\endcsname{}%
		\listcsgadd{SM@ag@list@#1}{#2}%
		\expandafter\gdef\csname SM@ag@label@#2\endcsname{#3}%
	\fi
}
\newcommand{\SMAgendaEntryTitle}[2]{%
	\ifcsname SM@ag@title@#1\endcsname\else
		\expandafter\gdef\csname SM@ag@title@#1\endcsname{#2}%
	\fi
}
% {section}{id of the subsection's first slide}{subsection title}
\newcommand{\SMAgendaSubsection}[3]{%
	\ifcsname SM@ag@subseen@#2\endcsname\else
		\expandafter\gdef\csname SM@ag@subseen@#2\endcsname{}%
		\listcsgadd{SM@ag@sublist@#1}{#2}%
		\expandafter\gdef\csname SM@ag@subtitle@#2\endcsname{#3}%
	\fi
}
\newcommand{\SMAgendaDeckTarget}{\gdef\SM@ag@decktarget{}}
\newcommand{\SMAgendaChapterTarget}[1]{\expandafter\gdef\csname SM@ag@chtarget@#1\endcsname{}}
% -- the chapter agenda --
\newcount\SM@ag@n
\newcount\SM@ag@i
\newbox\SM@ag@box
\newcommand{\SM@ag@entry}[1]{%
	\item\hyperlink{#1}{\ifcsname SM@ag@title@#1\endcsname\csname SM@ag@title@#1\endcsname
		\else\csname SM@ag@label@#1\endcsname\fi}%
}
% a subsection agenda's entry: #1 is the id of the subsection's first slide
\newcommand{\SM@ag@subentry}[1]{%
	\item\hyperlink{#1}{\csname SM@ag@subtitle@#1\endcsname}%
}
% (\def + \ifx, not \ifstrequal: that one would compare the *name* of
% \SM@ag@self, see \decklanguage)
\newcommand{\SM@ag@count}[1]{\def\SM@ag@this{#1}\ifx\SM@ag@this\SM@ag@self\else\advance\SM@ag@n by 1 \fi}
\newcommand{\SM@ag@split}[1]{%
	\def\SM@ag@this{#1}%
	\ifx\SM@ag@this\SM@ag@self\else
		\advance\SM@ag@i by 1
		\ifnum\SM@ag@i>\SM@ag@half
			\gappto\SM@ag@right{\SM@ag@entrycmd{#1}}%
		\else
			\gappto\SM@ag@left{\SM@ag@entrycmd{#1}}%
		\fi
	\fi
}
\newlength{\SM@ag@maxh}
\newcommand{\SM@ag@typeset}[2]{% #1 width, #2 font switch
	\setbox\SM@ag@box=\hbox{\begin{minipage}[t]{#1}%
		#2%
		\begin{columns}[t,onlytextwidth]
			\begin{column}{0.48\linewidth}
				\ifx\SM@ag@left\@empty\else\begin{itemize}\SM@ag@left\end{itemize}\fi
			\end{column}
			\begin{column}{0.48\linewidth}
				\ifx\SM@ag@right\@empty\else\begin{itemize}\SM@ag@right\end{itemize}\fi
			\end{column}
		\end{columns}%
	\end{minipage}}%
}
% \SMChapterAgenda lists the chapter's slides, \SMSubsectionAgenda its
% subsections (2026-10-12, direct instruction: "an agenda slide that lists
% all the subsections in a chapter"; each entry links to the subsection's
% first slide, titled as its \subsection{}). Same two columns, font steps and
% chapter target ("Chapter contents" in the navigation bar).
\newcommand{\SMChapterAgenda}[1][\linewidth]{%
	\SM@ag@main{SM@ag@list@}{\SM@ag@entry}{#1}%
}
\newcommand{\SMSubsectionAgenda}[1][\linewidth]{%
	\SM@ag@main{SM@ag@sublist@}{\SM@ag@subentry}{#1}%
}
% #1 the name prefix of the recorded list (+ section number), #2 the entry
% macro, #3 the width
\newcommand{\SM@ag@main}[3]{%
	\let\SM@ag@entrycmd#2%
	\edef\SM@ag@sec{\the\c@section}%
	\edef\SM@ag@self{\slideID}%
	\hypertarget{SMChapterAgenda\SM@ag@sec}{}%
	\immediate\write\@auxout{\string\SMAgendaChapterTarget{\SM@ag@sec}}%
	\SM@ag@n=0
	\ifcsdef{#1\SM@ag@sec}{\forlistcsloop{\SM@ag@count}{#1\SM@ag@sec}}{}%
	\ifnum\SM@ag@n>0
		\SM@ag@i=0
		\edef\SM@ag@half{\the\numexpr(\SM@ag@n+1)/2\relax}% ceil(n/2); \numexpr rounds
		\ifodd\SM@ag@n\else\edef\SM@ag@half{\the\numexpr\SM@ag@n/2\relax}\fi
		\gdef\SM@ag@left{}\gdef\SM@ag@right{}%
		\forlistcsloop{\SM@ag@split}{#1\SM@ag@sec}%
		\ifdim\SMChapterAgendaMaxHeight=0pt
			% the room below a sectionframe's title bar, down to the Bottom Margin
			\setlength{\SM@ag@maxh}{\dimexpr\paperheight-\SMSectionframeBarBottom
				-\SMSectionframeContentDrop-\SMContentBottomAlign-14pt\relax}%
		\else
			\setlength{\SM@ag@maxh}{\SMChapterAgendaMaxHeight}%
		\fi
		\SM@ag@typeset{#3}{\normalsize}%
		\ifdim\dimexpr\ht\SM@ag@box+\dp\SM@ag@box\relax>\SM@ag@maxh
			\SM@ag@typeset{#3}{\small}%
			\ifdim\dimexpr\ht\SM@ag@box+\dp\SM@ag@box\relax>\SM@ag@maxh
				\SM@ag@typeset{#3}{\footnotesize}%
			\fi
		\fi
		\noindent\box\SM@ag@box\par
	\fi
}
% -- the navigation bar (footer) --
\newcommand{\SMNavigationBar}[1][\SMNavigationLabelDeck,\SMNavigationLabelChapter,\SMNavigationLabelPreviews]{%
	\global\SMNavigationBartrue
	\SM@nav@labels#1,,,\@nil
}
\def\SM@nav@labels#1,#2,#3,#4\@nil{%
	\gdef\SMNavigationLabelDeck{#1}\gdef\SMNavigationLabelChapter{#2}\gdef\SMNavigationLabelPreviews{#3}%
}
% Link text in the footer: the footer's own text colour, underlined in
% c_secondary_blue (the \href underline, see \SMHrefUnderline).
\newcommand{\SMLinkUnderline}[1]{%
	\sbox{\SMHrefTextBox}{#1}%
	\usebox{\SMHrefTextBox}%
	\hspace{-\wd\SMHrefTextBox}%
	\SMHrefUnderlineMeasureRefHeights
	\setbox\SMHrefUnderlineBox=\hbox{M}\SMHrefUnderlineNowHeight=\ht\SMHrefUnderlineBox
	\SMHrefUnderlineSetThickness
	\raisebox{-0.35ex}[0pt][0pt]{\textcolor{c_secondary_blue}{\rule{\wd\SMHrefTextBox}{\SMHrefUnderlineThickness}}}%
}
\newcommand{\SM@nav@link}[2]{{\hypersetup{linkcolor=\SMFooterTextColor}\hyperlink{#1}{\SMLinkUnderline{#2}}}}
% \SMLink{target}{text} -- a link to a slide (or any \hypertarget) in running
% text, underlined like an \href so a reader sees it is a link: a bare
% \hyperlink{target}{text} has no underline. One line only, like \href's
% underline: keep the link text short. Buttons still use \hyperlink with
% \beamergotobutton / \beamerreturnbutton.
\newcommand{\SMLink}[2]{\hyperlink{#1}{\SMLinkUnderline{#2}}}
\newcommand{\SM@nav@sep}{\ \textperiodcentered\ }
\newcommand{\SMNavigationLinks}{%
	\def\SM@nav@chapter{}%
	\ifnum\c@section>0
		\ifcsname SM@ag@chtarget@\the\c@section\endcsname
			\edef\SM@nav@chapter{SMChapterAgenda\the\c@section}%
		\fi
	\fi
	\ifx\SM@nav@chapter\@empty
		\ifdefined\SM@ag@decktarget\def\SM@nav@chapter{SMDeckAgenda}\fi
	\fi
	\def\SM@nav@any{}%
	\ \textbar\ \SMNavigationLead\ %
	\ifdefined\SM@ag@decktarget
		\SM@nav@link{SMDeckAgenda}{\SMNavigationLabelDeck}\def\SM@nav@any{1}%
	\fi
	\ifx\SM@nav@chapter\@empty\else
		\ifx\SM@nav@any\@empty\else\SM@nav@sep\fi
		\SM@nav@link{\SM@nav@chapter}{\SMNavigationLabelChapter}\def\SM@nav@any{1}%
	\fi
	\ifdefined\SMAppendixTocOverride
		\ifx\SM@nav@any\@empty\else\SM@nav@sep\fi
		\SM@nav@link{SMAppendixAgenda}{\SMNavigationLabelPreviews}%
	\fi
}
\makeatother

% ---- Bibliography (biblatex + biber, opt-in per deck) ----
% \SMUseBibliography[file.bib] -- in a deck's preamble (after the institution
% template is \input) instead of \usepackage{biblatex} + \addbibresource: loads
% biblatex with \SMBiblatexOptions, registers the .bib file (default
% SlideMage.bib), then runs \SMBibliographyStyle. The template decides the
% citation format by redefining the two macros (2026-10-07, direct instruction:
% the house format had been copied into every course deck); a deck that wants
% another format redefines them before the call, or does not use the command.
% Defaults: biblatex's stock authoryear style.
%
% The .bib lookup: \IfFileExists below is the biber equivalent of the
% \input@path/\graphicspath fallback: biber runs as its own process and reads
% the literal path recorded in the .bcf at compile time, so it cannot follow
% \input@path -- resolving the correct relative path once, here, at LaTeX time
% (when the real cwd is known) is what makes biber find the .bib file whether
% compiling from the project root or from inside decks/. \input@path is
% emptied for the check itself: \IfFileExists honors \input@path too, so
% without this it reports the file as "found" via the ../ fallback and then
% writes that unresolved literal path into the .bcf, which biber (a separate
% process that never sees \input@path) then cannot find.
\makeatletter
\newcommand{\SMBiblatexOptions}{style=authoryear, backend=biber}
\newcommand{\SMBibliographyStyle}{}
\newcommand{\SMUseBibliography}[1][SlideMage.bib]{%
	\expandafter\usepackage\expandafter[\SMBiblatexOptions]{biblatex}% (options are not expanded otherwise)
	\let\SM@SavedInputPath\input@path
	\def\input@path{}%
	\IfFileExists{#1}{\addbibresource{#1}}{\addbibresource{../#1}}%
	\let\input@path\SM@SavedInputPath
	\SMBibliographyStyle}
\makeatother
 --
% SMCodeKeywordColor (control sequences), SMCodeCommentColor (% comments),
% SMCodeBackgroundColor (the panel), SMCodeTextColor (ordinary code). A template that doesn't keeps the
% original look: highlight colour, primary gray, tertiary gray -- which in
% the Graz/MCI/Innsbruck palettes is yellow/orange and light gray on light
% gray, about 1.1-1.3:1 contrast; their templates set readable colours.
\makeatletter
\@ifundefinedcolor{SMCodeKeywordColor}{\colorlet{SMCodeKeywordColor}{c_primary_highlight}}{}
\@ifundefinedcolor{SMCodeCommentColor}{\colorlet{SMCodeCommentColor}{c_primary_gray}}{}
\@ifundefinedcolor{SMCodeBackgroundColor}{\colorlet{SMCodeBackgroundColor}{c_tertiary_gray}}{}
\@ifundefinedcolor{SMCodeTextColor}{\colorlet{SMCodeTextColor}{c_primary_dark}}{}
\newcommand{\SMCodePercent}{\textbackslash\%\global\let\lst@prefixkeyword\relax}
\newcommand{\SMCodeNewline}{\textcolor{SMCodeKeywordColor}{\textbackslash\textbackslash}\global\let\lst@prefixkeyword\relax}
\makeatother
\lstdefinestyle{SMCode}{
	basicstyle=\SMCodeFont\SMCodeSize\color{SMCodeTextColor},
	columns=fullflexible, keepspaces=true, tabsize=2,
	keywords={SMCodeNoSuchKeyword},
	alsoletter={\\}, keywordsprefix=\\,
	keywordstyle=\color{SMCodeKeywordColor},
	morecomment=[l]\%, commentstyle=\color{SMCodeCommentColor},
	showstringspaces=false, upquote=true,
	breaklines=true, breakindent=1.5em,
	aboveskip=0pt, belowskip=0pt,
	literate={\\\%}{{\SMCodePercent}}2 {\\\\}{{\SMCodeNewline}}2 {ä}{{\"a}}1 {ö}{{\"o}}1 {ü}{{\"u}}1 {Ä}{{\"A}}1 {Ö}{{\"O}}1 {Ü}{{\"U}}1
		{ß}{{\ss}}1 {é}{{\'e}}1 {–}{{--}}1 {—}{{---}}1 {·}{{$\cdot$}}1 {…}{{...}}1
		{→}{{$\rightarrow$}}1,
}
% The panel is a \colorbox around the listing, not listings' own frame and
% background: listings draws its frame outside the box it typesets, so a
% panel stuck out above the top of its column by 4-8pt (depending on the font
% size) and missed the Top Margin (v5.14.1, found by check_layout.py in the
% SlideMage manual). The box holds everything it draws.
\newcommand{\SMCodeFile}[2][]{\par\noindent{\setlength{\fboxsep}{3pt}%
	\colorbox{SMCodeBackgroundColor}{\begin{minipage}{\dimexpr\linewidth-2\fboxsep\relax}%
		\lstinputlisting[style=SMCode,#1]{#2}%
	\end{minipage}}}\par}

% \SMCode{text} / \SMCmd{name} -- code in running text: Fira Mono Regular at
% 90% (\texttt is the theme's Fira Mono Bold, which reads heavy and wide in a
% sentence of Fira Sans Light). \SMCmd{slide} prints \slide. The argument is
% typeset, so write \_ \{ \} \# as in any text. The transcript shows both as
% `code`. Ligatures=TeXOff: code shows -- and `` as typed, not as an en dash
% and a curly quote (fontspec's default for \newfontfamily).
\newcommand{\SMInlineCodeStyle}{}
\newcommand{\SMCode}[1]{{\SMInlineCodeFont\SMInlineCodeStyle #1}}
\newcommand{\SMCmd}[1]{{\SMInlineCodeFont\SMInlineCodeStyle\textbackslash #1}}

% \SMQRCode[height]{image}{caption} -- a QR code SlideMage made (2026-10-25,
% direct instruction: the Source view's Insert QR, the assistant's
% make_qr_code; src/slidemage/qr_code.py), the code at [height] (default
% 0.6\textheight: a lecture hall's back row needs about that) with an
% optional caption centred under it, as wide as the code -- the link, or a
% text such as "QR 1" that the slide refers to; {} for none. An inline box:
% several codes go side by side or in columns, and
% \cornerfigure{\SMQRCode{...}{...}} puts one in a corner. The caption's
% look is \SMQRCodeCaptionStyle, which a template may redefine.
\newsavebox{\SMQRCodeBox}
\providecommand{\SMQRCodeCaptionStyle}{\scriptsize}
\NewDocumentCommand{\SMQRCode}{O{0.6\textheight} m m}{%
	\begingroup
	\sbox{\SMQRCodeBox}{\includegraphics[height=#1]{#2}}%
	\begin{minipage}[t]{\wd\SMQRCodeBox}%
		\centering\usebox{\SMQRCodeBox}%
		\if\relax\detokenize{#3}\relax\else
			\par{\SMQRCodeCaptionStyle #3\par}%     the code's own white border is the gap
		\fi
	\end{minipage}%
	\endgroup
}

% Navigation for readers (2026-10-04, SlideMage manual: readers must find
% their way without knowing that the logo is a hidden link). Three opt-in
% pieces, each usable in any deck:
%
% \SMChapterAgenda[width] -- the slides of the current \section as a
%   two-column list, each entry a link to that slide (its frame title, or the
%   slide's \SMAgendaTitle{...}, or else its slug with spaces). Meant for the
%   body of a chapter's sectionframe (the manual puts one on every divider).
%   [width] (default \linewidth) leaves room for e.g. a \sectionfigure. Set in
%   \normalsize, else \small, else \footnotesize, the first that fits
%   \SMChapterAgendaMaxHeight.
%   Works from the .aux file like a table of contents: every slide writes
%   \SMAgendaSlide{section}{id}{label} (in \setSlideID) and its frame title
%   (\SMAgendaEntryTitle, in the frametitle hook), and the agenda shows what
%   the previous pass recorded -- latexmk's rerun fills it in. A slide shown
%   twice (\slidestep) is listed once, at its first place; the agenda's own
%   slide is not listed.
% \SMSubsectionAgenda[width] -- the same, but one entry per \subsection of
%   the current \section (2026-10-12): the subsection's title, linked to its
%   first slide in this version. For a chapter divider that introduces the
%   subsections instead of listing every slide. Slides before the chapter's
%   first \subsection are not listed; it fills in on the rerun.
% \SMAgendaTitle{text} -- in a slide file, before the frame: the entry text
%   for a slide whose frame has no title (\quoteframe, \fullbleedframe, ...)
%   or should read differently in the agenda.
% \SMDeckAgenda -- marks the deck's contents slide (put it in that slide's
%   body, next to \tableofcontents) as the target of the navigation bar's
%   first link.
% \SMNavigationBar[deck,chapter,previews] -- in the deck preamble: the footer
%   then reads "\FooterText | Navigation: <deck> · <chapter> · <previews>",
%   three links to the \SMDeckAgenda slide, the current chapter's
%   \SMChapterAgenda (the contents slide before the first \section, or where
%   a chapter has none) and the generated appendix agenda with the slide
%   previews. A link whose target is not in this PDF is left out: the last
%   one exists only when SlideMage generates the appendix (the Generate
%   appendix TOC checkbox, or \deckappendixtoc in the deck). Default labels
%   "Contents,Chapter contents,Slide previews"; the word before them is
%   \SMNavigationLead ("Navigation:"). Off by default: other decks' footers
%   and the logo link are unchanged.
% \deckappendixtoc -- in the deck preamble, inert for LaTeX (like
%   \decktags): SlideMage then generates the appendix for this deck on every
%   whole-deck compile, whatever the toolbar checkbox says.
\newcommand{\deckappendixtoc}{}
\newlength{\SMChapterAgendaMaxHeight}
\newcommand{\SMNavigationLead}{Navigation:}
\newcommand{\SMNavigationLabelDeck}{Contents}
\newcommand{\SMNavigationLabelChapter}{Chapter contents}
\newcommand{\SMNavigationLabelPreviews}{Slide previews}
\makeatletter
% -- recording (writes to the .aux) --
\newcommand{\SM@agenda@slide}[2]{%
	\ifcsname SM@ag@written@#1\endcsname\else
		\expandafter\gdef\csname SM@ag@written@#1\endcsname{}%
		% label from the rest of the id-slug ("EN-Slide-IDs-"): language
		% segment and trailing hyphen dropped, hyphens to spaces. (No % in
		% the Lua code -- TeX would read it as a comment -- and no hash
		% sign, which is a macro parameter here.)
		\edef\SM@ag@label{\directlua{
			local s = "\luaescapestring{#2}"
			while string.sub(s, -1) == "-" do s = string.sub(s, 1, -2) end
			for _, p in ipairs({"EN-", "DE-", "mult-"}) do
				if string.sub(s, 1, string.len(p)) == p then s = string.sub(s, string.len(p) + 1) end
			end
			s = string.gsub(s, "[-]", " ")
			tex.sprint(-2, s)}}%
		\immediate\write\@auxout{\string\SMAgendaSlide{\the\c@section}{#1}{\SM@ag@label}}%
	\fi
	% the first slide of each \subsection and the subsection's title, for
	% \SMSubsectionAgenda (its own guard: a slide shown twice must not hide
	% that it also opens a later subsection)
	\ifnum\c@subsection>0
		\ifcsname SM@ag@subwritten@\the\c@section.\the\c@subsection\endcsname\else
			\expandafter\gdef\csname SM@ag@subwritten@\the\c@section.\the\c@subsection\endcsname{}%
			\immediate\write\@auxout{\string\SMAgendaSubsection{\the\c@section}{#1}{\detokenize\expandafter{\subsecname}}}%
		\fi
	\fi
}
\newcommand{\SM@agenda@title}[1]{%
	\ifcsname SM@ag@titled@\slideID\endcsname\else
		\expandafter\gdef\csname SM@ag@titled@\slideID\endcsname{}%
		\immediate\write\@auxout{\string\SMAgendaEntryTitle{\slideID}{\detokenize{#1}}}%
	\fi
}
\newcommand{\SMAgendaTitle}[1]{\SM@agenda@title{#1}}
\newcommand{\SM@agenda@frametitle}{%
	\ifx\beamer@frametitle\@empty\else\expandafter\SM@agenda@title\expandafter{\beamer@frametitle}\fi
}
\newcommand{\SMAgendaFrametitleHook}{\SM@agenda@frametitle}% called by \SMFrametitleHookCode
\newcommand{\SMDeckAgenda}{%
	\hypertarget{SMDeckAgenda}{}%
	\immediate\write\@auxout{\string\SMAgendaDeckTarget}%
}
% -- reading (the .aux, at \begin{document}; its second reading at
% \end{document} changes nothing) --
\newcommand{\SMAgendaSlide}[3]{%
	\ifcsname SM@ag@seen@#2\endcsname\else
		\expandafter\gdef\csname SM@ag@seen@#2\endcsname{}%
		\listcsgadd{SM@ag@list@#1}{#2}%
		\expandafter\gdef\csname SM@ag@label@#2\endcsname{#3}%
	\fi
}
\newcommand{\SMAgendaEntryTitle}[2]{%
	\ifcsname SM@ag@title@#1\endcsname\else
		\expandafter\gdef\csname SM@ag@title@#1\endcsname{#2}%
	\fi
}
% {section}{id of the subsection's first slide}{subsection title}
\newcommand{\SMAgendaSubsection}[3]{%
	\ifcsname SM@ag@subseen@#2\endcsname\else
		\expandafter\gdef\csname SM@ag@subseen@#2\endcsname{}%
		\listcsgadd{SM@ag@sublist@#1}{#2}%
		\expandafter\gdef\csname SM@ag@subtitle@#2\endcsname{#3}%
	\fi
}
\newcommand{\SMAgendaDeckTarget}{\gdef\SM@ag@decktarget{}}
\newcommand{\SMAgendaChapterTarget}[1]{\expandafter\gdef\csname SM@ag@chtarget@#1\endcsname{}}
% -- the chapter agenda --
\newcount\SM@ag@n
\newcount\SM@ag@i
\newbox\SM@ag@box
\newcommand{\SM@ag@entry}[1]{%
	\item\hyperlink{#1}{\ifcsname SM@ag@title@#1\endcsname\csname SM@ag@title@#1\endcsname
		\else\csname SM@ag@label@#1\endcsname\fi}%
}
% a subsection agenda's entry: #1 is the id of the subsection's first slide
\newcommand{\SM@ag@subentry}[1]{%
	\item\hyperlink{#1}{\csname SM@ag@subtitle@#1\endcsname}%
}
% (\def + \ifx, not \ifstrequal: that one would compare the *name* of
% \SM@ag@self, see \decklanguage)
\newcommand{\SM@ag@count}[1]{\def\SM@ag@this{#1}\ifx\SM@ag@this\SM@ag@self\else\advance\SM@ag@n by 1 \fi}
\newcommand{\SM@ag@split}[1]{%
	\def\SM@ag@this{#1}%
	\ifx\SM@ag@this\SM@ag@self\else
		\advance\SM@ag@i by 1
		\ifnum\SM@ag@i>\SM@ag@half
			\gappto\SM@ag@right{\SM@ag@entrycmd{#1}}%
		\else
			\gappto\SM@ag@left{\SM@ag@entrycmd{#1}}%
		\fi
	\fi
}
\newlength{\SM@ag@maxh}
\newcommand{\SM@ag@typeset}[2]{% #1 width, #2 font switch
	\setbox\SM@ag@box=\hbox{\begin{minipage}[t]{#1}%
		#2%
		\begin{columns}[t,onlytextwidth]
			\begin{column}{0.48\linewidth}
				\ifx\SM@ag@left\@empty\else\begin{itemize}\SM@ag@left\end{itemize}\fi
			\end{column}
			\begin{column}{0.48\linewidth}
				\ifx\SM@ag@right\@empty\else\begin{itemize}\SM@ag@right\end{itemize}\fi
			\end{column}
		\end{columns}%
	\end{minipage}}%
}
% \SMChapterAgenda lists the chapter's slides, \SMSubsectionAgenda its
% subsections (2026-10-12, direct instruction: "an agenda slide that lists
% all the subsections in a chapter"; each entry links to the subsection's
% first slide, titled as its \subsection{}). Same two columns, font steps and
% chapter target ("Chapter contents" in the navigation bar).
\newcommand{\SMChapterAgenda}[1][\linewidth]{%
	\SM@ag@main{SM@ag@list@}{\SM@ag@entry}{#1}%
}
\newcommand{\SMSubsectionAgenda}[1][\linewidth]{%
	\SM@ag@main{SM@ag@sublist@}{\SM@ag@subentry}{#1}%
}
% #1 the name prefix of the recorded list (+ section number), #2 the entry
% macro, #3 the width
\newcommand{\SM@ag@main}[3]{%
	\let\SM@ag@entrycmd#2%
	\edef\SM@ag@sec{\the\c@section}%
	\edef\SM@ag@self{\slideID}%
	\hypertarget{SMChapterAgenda\SM@ag@sec}{}%
	\immediate\write\@auxout{\string\SMAgendaChapterTarget{\SM@ag@sec}}%
	\SM@ag@n=0
	\ifcsdef{#1\SM@ag@sec}{\forlistcsloop{\SM@ag@count}{#1\SM@ag@sec}}{}%
	\ifnum\SM@ag@n>0
		\SM@ag@i=0
		\edef\SM@ag@half{\the\numexpr(\SM@ag@n+1)/2\relax}% ceil(n/2); \numexpr rounds
		\ifodd\SM@ag@n\else\edef\SM@ag@half{\the\numexpr\SM@ag@n/2\relax}\fi
		\gdef\SM@ag@left{}\gdef\SM@ag@right{}%
		\forlistcsloop{\SM@ag@split}{#1\SM@ag@sec}%
		\ifdim\SMChapterAgendaMaxHeight=0pt
			% the room below a sectionframe's title bar, down to the Bottom Margin
			\setlength{\SM@ag@maxh}{\dimexpr\paperheight-\SMSectionframeBarBottom
				-\SMSectionframeContentDrop-\SMContentBottomAlign-14pt\relax}%
		\else
			\setlength{\SM@ag@maxh}{\SMChapterAgendaMaxHeight}%
		\fi
		\SM@ag@typeset{#3}{\normalsize}%
		\ifdim\dimexpr\ht\SM@ag@box+\dp\SM@ag@box\relax>\SM@ag@maxh
			\SM@ag@typeset{#3}{\small}%
			\ifdim\dimexpr\ht\SM@ag@box+\dp\SM@ag@box\relax>\SM@ag@maxh
				\SM@ag@typeset{#3}{\footnotesize}%
			\fi
		\fi
		\noindent\box\SM@ag@box\par
	\fi
}
% -- the navigation bar (footer) --
\newcommand{\SMNavigationBar}[1][\SMNavigationLabelDeck,\SMNavigationLabelChapter,\SMNavigationLabelPreviews]{%
	\global\SMNavigationBartrue
	\SM@nav@labels#1,,,\@nil
}
\def\SM@nav@labels#1,#2,#3,#4\@nil{%
	\gdef\SMNavigationLabelDeck{#1}\gdef\SMNavigationLabelChapter{#2}\gdef\SMNavigationLabelPreviews{#3}%
}
% Link text in the footer: the footer's own text colour, underlined in
% c_secondary_blue (the \href underline, see \SMHrefUnderline).
\newcommand{\SMLinkUnderline}[1]{%
	\sbox{\SMHrefTextBox}{#1}%
	\usebox{\SMHrefTextBox}%
	\hspace{-\wd\SMHrefTextBox}%
	\SMHrefUnderlineMeasureRefHeights
	\setbox\SMHrefUnderlineBox=\hbox{M}\SMHrefUnderlineNowHeight=\ht\SMHrefUnderlineBox
	\SMHrefUnderlineSetThickness
	\raisebox{-0.35ex}[0pt][0pt]{\textcolor{c_secondary_blue}{\rule{\wd\SMHrefTextBox}{\SMHrefUnderlineThickness}}}%
}
\newcommand{\SM@nav@link}[2]{{\hypersetup{linkcolor=\SMFooterTextColor}\hyperlink{#1}{\SMLinkUnderline{#2}}}}
% \SMLink{target}{text} -- a link to a slide (or any \hypertarget) in running
% text, underlined like an \href so a reader sees it is a link: a bare
% \hyperlink{target}{text} has no underline. One line only, like \href's
% underline: keep the link text short. Buttons still use \hyperlink with
% \beamergotobutton / \beamerreturnbutton.
\newcommand{\SMLink}[2]{\hyperlink{#1}{\SMLinkUnderline{#2}}}
\newcommand{\SM@nav@sep}{\ \textperiodcentered\ }
\newcommand{\SMNavigationLinks}{%
	\def\SM@nav@chapter{}%
	\ifnum\c@section>0
		\ifcsname SM@ag@chtarget@\the\c@section\endcsname
			\edef\SM@nav@chapter{SMChapterAgenda\the\c@section}%
		\fi
	\fi
	\ifx\SM@nav@chapter\@empty
		\ifdefined\SM@ag@decktarget\def\SM@nav@chapter{SMDeckAgenda}\fi
	\fi
	\def\SM@nav@any{}%
	\ \textbar\ \SMNavigationLead\ %
	\ifdefined\SM@ag@decktarget
		\SM@nav@link{SMDeckAgenda}{\SMNavigationLabelDeck}\def\SM@nav@any{1}%
	\fi
	\ifx\SM@nav@chapter\@empty\else
		\ifx\SM@nav@any\@empty\else\SM@nav@sep\fi
		\SM@nav@link{\SM@nav@chapter}{\SMNavigationLabelChapter}\def\SM@nav@any{1}%
	\fi
	\ifdefined\SMAppendixTocOverride
		\ifx\SM@nav@any\@empty\else\SM@nav@sep\fi
		\SM@nav@link{SMAppendixAgenda}{\SMNavigationLabelPreviews}%
	\fi
}
\makeatother

% ---- Bibliography (biblatex + biber, opt-in per deck) ----
% \SMUseBibliography[file.bib] -- in a deck's preamble (after the institution
% template is \input) instead of \usepackage{biblatex} + \addbibresource: loads
% biblatex with \SMBiblatexOptions, registers the .bib file (default
% SlideMage.bib), then runs \SMBibliographyStyle. The template decides the
% citation format by redefining the two macros (2026-10-07, direct instruction:
% the house format had been copied into every course deck); a deck that wants
% another format redefines them before the call, or does not use the command.
% Defaults: biblatex's stock authoryear style.
%
% The .bib lookup: \IfFileExists below is the biber equivalent of the
% \input@path/\graphicspath fallback: biber runs as its own process and reads
% the literal path recorded in the .bcf at compile time, so it cannot follow
% \input@path -- resolving the correct relative path once, here, at LaTeX time
% (when the real cwd is known) is what makes biber find the .bib file whether
% compiling from the project root or from inside decks/. \input@path is
% emptied for the check itself: \IfFileExists honors \input@path too, so
% without this it reports the file as "found" via the ../ fallback and then
% writes that unresolved literal path into the .bcf, which biber (a separate
% process that never sees \input@path) then cannot find.
\makeatletter
\newcommand{\SMBiblatexOptions}{style=authoryear, backend=biber}
\newcommand{\SMBibliographyStyle}{}
\newcommand{\SMUseBibliography}[1][SlideMage.bib]{%
	\expandafter\usepackage\expandafter[\SMBiblatexOptions]{biblatex}% (options are not expanded otherwise)
	\let\SM@SavedInputPath\input@path
	\def\input@path{}%
	\IfFileExists{#1}{\addbibresource{#1}}{\addbibresource{../#1}}%
	\let\input@path\SM@SavedInputPath
	\SMBibliographyStyle}
\makeatother

%   ... additions only: dark-mode beamer colors, citation format, project macros
%
% Rules for a template:
%   1. Before this file: only \SMDarkMode (written into the deck by SlideMage).
%   2. Between this file and SlideMage_base: \setlength, \renewcommand,
%      \definecolor and \colorlet on the names in this catalog. Every name
%      exists already, so a misspelt \renewcommand is an error instead of a
%      setting that silently does nothing.
%   3. After SlideMage_base: additions only.
%
% Decks never load the beamer theme themselves; this file does. A length given
% in ex or em is evaluated where the template sets it, in the normal body font,
% after the theme has loaded its fonts.

\newcommand{\SMSetupLoaded}{}

% The beamer theme, with the body fonts.
\usetheme{gotham}

% ---------------------------------------------------------------------------
% Palette
% ---------------------------------------------------------------------------
% The SlideMage look; a template redefines what it changes (usually all of
% them). The framework combines these pairs, so a template should keep them
% readable (WCAG 2.1): body text c_primary_dark on each c_secondary_* block
% title color >= 4.5:1, title/footer text on its box color >= 4.5:1, white text
% on c_primary_gray (highlightframe), link underline c_secondary_blue on the
% canvas >= 3:1.
%
% SlideMage's own values: wherever possible a color the app itself uses
% (src/slidemage/palette.py, qt_deck_table.py, assets/splash.png). White
% title/footer text on c_primary_dark is 16.7:1, white on c_primary_gray 5.0:1.
\definecolor{c_primary_dark}{HTML}{1e1e1e}      % body text, dark boxes; splash screen background (palette dark "bg")
\definecolor{c_primary_highlight}{HTML}{0066cc} % accent: bullets, highlights; SlideMage button blue (Create PDF, progress bar, links)
\definecolor{c_primary_gray}{HTML}{6f6f6f}      % highlightframe background; palette dark "disabled_fg", dark enough for white text
\definecolor{c_primary_white}{HTML}{FFFFFF}
\definecolor{c_secondary_violet}{HTML}{ba68c8}  % block titles (block); palette dark "tab_deck_editor"
\definecolor{c_tertiary_violet}{HTML}{dcb4e4}   % block bodies; derived (50% white-blend)
% Derived, not an app color: the button blue is too dark for the dark text
% an infoblock title puts on it, the app's dark-mode link blue (4da6ff) too
% light for a link underline on white; this one sits between them and passes
% both (4.9:1 and 3.4:1).
\definecolor{c_secondary_blue}{HTML}{3d8fe0}    % infoblock titles, link underline
\definecolor{c_tertiary_blue}{HTML}{9ec7f0}     % derived (50% white-blend)
\definecolor{c_secondary_green}{HTML}{4caf50}   % exampleblock titles; palette dark "status_pass"
\definecolor{c_tertiary_green}{HTML}{81c784}    % palette dark "tab_slide_editor"
\definecolor{c_secondary_red}{HTML}{f04848}     % alertblock titles; palette dark "status_fail"
\definecolor{c_tertiary_red}{HTML}{f8a4a4}      % derived (50% white-blend)
\definecolor{c_secondary_orange}{HTML}{ffb74d}  % reference block titles; palette dark "tab_archived"
\definecolor{c_tertiary_orange}{HTML}{ffdba6}   % derived (50% white-blend)
\definecolor{c_secondary_yellow}{HTML}{ffd500}  % logo/splash yellow
\definecolor{c_tertiary_yellow}{HTML}{ffe066}   % deck editor's highlight yellow
\definecolor{c_tertiary_gray}{HTML}{e4e4e4}     % block bodies, process boxes; palette light "panel_bg"
% Supporting colors ("custom": not part of the palette proper).
% c_custom_gray is hyperlink text, button fill and the quote attribution.
\definecolor{c_custom_gray}{HTML}{2d2d2d}       % palette dark "panel_bg"
\definecolor{c_custom_red}{HTML}{cc0000}        % palette light "status_fail"
\definecolor{c_custom_green}{HTML}{1a8a1a}      % palette light "status_pass"
\definecolor{c_custom_brown}{HTML}{996600}      % palette light "syntax_brace"

% ---------------------------------------------------------------------------
% Roles: which palette color a part of the slide uses (color names, read
% when used, so they follow a redefined palette color)
% ---------------------------------------------------------------------------
\newcommand{\SMBulletColor}{c_primary_highlight}       % itemize bullets
\newcommand{\SMTitleBoxColor}{c_primary_dark}          % frame title bar
\newcommand{\SMTitleTextColor}{c_primary_white}        % frame title text
\newcommand{\SMFooterBoxColorDefault}{c_primary_dark}  % footer box
\newcommand{\SMFooterTextColor}{c_primary_white}       % footer text, slide id

% Optional colors a template may define; SlideMage_base derives each one from
% the palette when the template leaves it undefined:
%   SMCanvasColor          slide background                  (white)
%   SMCardBackgroundColor  SMCard fill                       (white)
%   SMCardBorderColor      SMCard border                     (c_primary_dark)
%   SMProcessBorderColor   SMProcess border, dark mode       (c_tertiary_gray)
%   SMCodeKeywordColor     \SMCodeFile commands              (c_primary_highlight)
%   SMCodeCommentColor     \SMCodeFile comments              (c_primary_gray)
%   SMCodeBackgroundColor  \SMCodeFile panel                 (c_tertiary_gray)
%   SMCodeTextColor        \SMCodeFile text                  (c_primary_dark)

% ---------------------------------------------------------------------------
% Logo
% ---------------------------------------------------------------------------
% A vector PDF cropped to its ink, found on the TeX search path (a project
% path such as schemes/<name>/logo.pdf, or a file next to the template).
\newcommand{\SMInstituteLogoPath}{slidemage-logo.pdf}
\newlength{\SMLogoHeight}
\setlength{\SMLogoHeight}{3ex}                         % logo height in the title bar

% ---------------------------------------------------------------------------
% Frame title bar
% ---------------------------------------------------------------------------
% Space around the frame title (the theme's registers under their SlideMage
% names). The right padding is measured against the compiled PDF for the logo:
% see SlideMage_base.tex's comment on \gothamFrametitleRightpading.
\let\SMFrametitleTopPadding\gothamFrametitleToppading
\let\SMFrametitleBottomPadding\gothamFrametitleBottompading
\let\SMFrametitleLeftPadding\gothamFrametitleLeftpading
\let\SMFrametitleRightPadding\gothamFrametitleRightpading
\setlength{\SMFrametitleRightPadding}{0.28cm}

%
% and fira is the default. A project may keep its own profile next to its
% templates (SlideMage_layout_<name>.tex); the TeX search path finds it. The
% recipe for a new profile is in SlideMage_setup.tex, in the header.
%
% What a profile must define:
%   fonts                   the sans text font (\sfdefault, with the weight
%                           mapping), the mono font (\ttfamily), \SMCodeFont,
%                           \SMInlineCodeFont, \SMCardHeadingFont
%   \SMFontAscender         the ascender of the text font in em (see below)
%   \SMLinespread           the \linespread of running text
%   \SMParskip              the \parskip of running text; the theme applies both
%                           where it is loaded, in the normal font
% The theme sets only font sizes, never font files.

\RequirePackage{etoolbox}

\ExplSyntaxOn
\makeatletter

\msg_new:nnn { SlideMage } { font-not-found }
  { Could~not~find~font~'#1'.\c_space_tl }
\msg_new:nnn { SlideMage } { fonts-not-available }
  { No~suitable~fonts~found.~Using~fallback~fonts. }
\msg_new:nnn { SlideMage } { requires-lualatex }
  { The~SlideMage~layout~profile~needs~LuaLaTeX~(Fira~fonts~and~Lua~code). }

\int_new:N \g__SM_fonts_not_found_int

\cs_new_protected:Nn \__SM_check_font:n
  {
    \fontspec_font_if_exist:nF { #1 }
      {
        \fontspec_font_if_exist:nF { #1~Regular }
          {
            \int_incr:N \g__SM_fonts_not_found_int
            \msg_warning:nnn { SlideMage } { font-not-found } { #1 }
          }
      }
  }

\cs_new_protected:Nn \__SM_if_fonts_available:nnn
  {
    \int_zero:N \g__SM_fonts_not_found_int
    \clist_map_function:nN { #1 } \__SM_check_font:n
    \int_compare:nTF { \g__SM_fonts_not_found_int = 0 } { #2 } { #3 }
  }

% Text and code fonts. fontspec loads with no-math, so that the sans font is in
% charge of text only; the math fonts are restored further down.
\sys_if_engine_luatex:TF
  {
    \@ifpackageloaded { fontspec }
      { \PassOptionsToPackage { no-math } { fontspec } }
      { \RequirePackage [ no-math ] { fontspec } }
    \__SM_if_fonts_available:nnn
      {
        FiraSans-Light.otf, FiraSans-LightItalic.otf,
        FiraSans-Regular.otf, FiraSans-Italic.otf
      }
      {
        \setsansfont
          [
            ItalicFont = { FiraSans-LightItalic.otf },
            BoldFont = { FiraSans-Regular.otf },
            BoldItalicFont = { FiraSans-Italic.otf }
          ]
          { FiraSans-Light.otf }
        \renewcommand { \ttdefault } { lmtt }
      }
      { \msg_warning:nn { SlideMage } { fonts-not-available } }
    \__SM_if_fonts_available:nnn { FiraMono-Medium.otf, FiraMono-Bold.otf }
      { \setmonofont [ BoldFont = { FiraMono-Medium.otf } ] { FiraMono-Bold.otf } }
      { \msg_warning:nn { SlideMage } { fonts-not-available } }
    % Digits in tables line up in columns.
    \AtBeginEnvironment { tabular } { \addfontfeature { Numbers = { Monospaced } } }
  }
  { \msg_warning:nn { SlideMage } { requires-lualatex } }

\makeatother
\ExplSyntaxOff

% Metrics -----------------------------------------------------------------------
% The ascender of the text font in em: the height of a text line's line box above
% its baseline, which is what the automatic Top Margin measures to (a line's top
% is its baseline minus this much) and what \SMLineTop is as tall as.
% SlideMage_setup.tex reads the ascender from the loaded font (luaotfload's
% OpenType metrics: ascender / units per em), rounds it to four digits and keeps
% this number only where the font gives none.
\newcommand{\SMFontAscender}{0.935}
% Running text: \linespread and \parskip. The theme applies them where it is
% loaded, with the normal font in force, so the em is this font's.
\newcommand{\SMLinespread}{1.15}
\newcommand{\SMParskip}{0.5em}

% Calibrations -------------------------------------------------------------------------
% Lengths that were measured against the rendered PDF at Fira Sans Light and Fira Mono
% (SlideMage v9.4.0 collects them here; they sat in SlideMage_base.tex). They are numbers
% of this font, so they belong with it: a template does not change them, and a profile for
% another font gives its own (tools/calibrate_font_profile.py measures them on the calibration deck
% of tests/fixtures/template_stress/; for this profile it reads back the values below).
%
% The footer. The text in its box: the frame footer sat noticeably closer to the top edge
% of the box than to the bottom (0.14pt above against 1.12pt below).
\newlength{\SMFooterTextDrop}
\setlength{\SMFooterTextDrop}{-0.49pt}
% (The right ends of the footer line and of the footnotes against the logo's right edge are
% derived by SlideMage_base since v9.7.0: they were two numbers here.)
% Numbered lists. A nested enumerate sets its numeral in a right-aligned box whose width puts
% a single-digit numeral flush with the text of the item above (20.33pt; 19.23pt left the
% widest digit, "6.", up to 1.05pt left of it); the numeral may be wider than the dot
% column (7.85pt) and the list keeps a left margin and label gap of its own. The outer
% enumerate's numeral column is wider than a bullet by the last length.
\newlength{\SMEnumerateLabelBox}
\setlength{\SMEnumerateLabelBox}{20.33pt}
\newlength{\SMEnumerateNestedLabelSep}
\setlength{\SMEnumerateNestedLabelSep}{5.25pt}
\newlength{\SMEnumerateNestedLeftMargin}
\setlength{\SMEnumerateNestedLeftMargin}{13.1pt}
\newlength{\SMEnumerateNestedMinNumber}
\setlength{\SMEnumerateNestedMinNumber}{7.85pt}
\newlength{\SMEnumerateTextExtra}
\setlength{\SMEnumerateTextExtra}{6.3pt}
% Lists with letter or number labels ("a)", "1."): the box that holds a label at the small
% itemize size with a little clearance (9.1pt), and the gap between it and the text: 2pt
% plus half the width of a "C" (4.75pt).
\newlength{\SMLabelItemWidth}
\setlength{\SMLabelItemWidth}{9.1pt}
\newlength{\SMLabelItemSep}
\setlength{\SMLabelItemSep}{4.75pt}
% The section divider frame measures its title with a \strut and corrects what remains:
% the depth of the reference text, the gap under the title and its correction, and the
% height correction.
\newlength{\SMSectionframeRefDepth}
\setlength{\SMSectionframeRefDepth}{4.42pt}
\newlength{\SMSectionframeBottomGap}
\setlength{\SMSectionframeBottomGap}{14.32pt}
\newlength{\SMSectionframeBottomPadCorrection}
\setlength{\SMSectionframeBottomPadCorrection}{5.58pt}
\newlength{\SMSectionframeHeightCorrection}
\setlength{\SMSectionframeHeightCorrection}{16.03pt}

% The pull-up of a top-aligned frame (standardframe) and the corrections of the frames the
% automatic Top Margin does not manage (the [c] frame, sectionframe, \SMManualTopMargin,
% \SMAutoTopMarginfalse). In every other frame the measurement corrects what is left, but not
% to the last digit: with these at 0 the frames of the manual land a thousandth of a point
% elsewhere and 11 of 92 pages differ in a few pixels (v9.4.0 trial), and a \begin{columns} below
% a paragraph still needs the 10.17pt, so none of the three goes. 11.5964pt is 40% more than
% the original 7.33pt (the Top Margin was reduced to 60%).
\newlength{\SMContentTopShift}
\setlength{\SMContentTopShift}{11.5964pt}
\newlength{\SMColumnsTopCorrection}
\setlength{\SMColumnsTopCorrection}{10.17pt}      % columns[t] starts 10.17pt below a list
\newlength{\SMColumnsTTopCompensation}
\setlength{\SMColumnsTTopCompensation}{13.76pt}  % columns[T] has taken part of that back

% Code fonts ----------------------------------------------------------------------
% \SMCodeFile (listings) and the like use Fira Mono Regular; the theme's \ttfamily
% is the heavier Fira Mono Bold.
\newfontfamily\SMCodeFont{FiraMono-Regular.otf}[BoldFont=FiraMono-Medium.otf]
% Code in running text (\SMCode, \SMCmd): Fira Mono Regular at 90%, which reads
% lighter than \texttt next to Fira Sans Light. Ligatures=TeXOff: code shows -- and
% `` as typed, not as an en dash and a curly quote (fontspec's default for
% \newfontfamily).
\newfontfamily\SMInlineCodeFont{FiraMono-Regular.otf}[Scale=0.9, Ligatures=TeXOff]
% The heading of a card in a diagram (\SMCard): a template with a medium weight
% redefines it, e.g. \renewcommand{\SMCardHeadingFont}{\SMCardMediumFont}.
\newcommand{\SMCardHeadingFont}{\bfseries}

% Math ------------------------------------------------------------------------------
% Restores genuine Computer Modern math fonts. Found 2026-07-15, direct
% instruction ("math equations do not have serifs, look at vcv"): despite
% fontspec being loaded with `no-math`, math still wasn't landing on Computer
% Modern -- checked directly against a compiled PDF's own per-glyph font names
% (pymupdf), not just eyeballed: digits/math-italic letters were coming out as
% FiraSans-Light(Italic) (this deck's own body-text sans font) and
% operators/parens as CMSS10 (Computer Modern's *sans-serif* shape, not its
% serif one). Root cause: beamer's own base classes redefine the math alphabets
% to follow the presentation's sans-serif text font by default, specifically so
% math matches an all-sans deck -- `no-math` only stops *fontspec* from doing
% this, it has no say over beamer's own separate override. Beamer ships the
% standard fix for exactly this case: a font theme named "serif" with an
% "onlymath" option that restores ordinary (Latin Modern/Computer-Modern-
% equivalent) math fonts without touching any of the surrounding sans-serif
% text. Confirmed directly afterward: every glyph in a re-checked formula now
% reports as LMRoman*/CMR10/CMMI8 (all genuinely serif), not FiraSans/CMSS10.
\usefonttheme[onlymath]{serif}

% \text{} (amsmath, used throughout for plain-word operands inside math, e.g.
% \text{E}[\tilde{R}_M]) is a separate mechanism from the math alphabets
% \usefonttheme[onlymath]{serif} above just fixed: unlike
% \mathrm/\mathit/etc., \text{} is defined to typeset in whatever the
% *surrounding running text* font is at that point in the document, not a math
% alphabet -- so on this deck's all-Fira body text, \text{} inside a formula
% still comes out in Fira Sans (and picks up italic if the enclosing text
% happens to be italicized), even though the digits/operators around it are
% correctly Computer Modern. Found 2026-07-20 on
% slides/2gd-Analyzing-Reward-to-Risk-Ratio-CAPM-current.tex's own
% \text{E}/\text{Cov} operands. Fixed generally, project-wide, by wrapping
% \text{}'s argument in an explicit upright-Computer-Modern font switch --
% \usefont{OT1}{cmr}{m}{n} rather than \mathrm{} (which would pull in
% \text{}'s own math-alphabet substitution logic unnecessarily) or \textrm{}
% (which switches *family* only, so still comes out italic inside an italicized
% context, which \text{} should never do here) -- deliberately hardcoding "not
% italic" (shape n) rather than inheriting the ambient shape, since \text{}'s
% whole purpose in this deck's formulas is plain upright labels (E, Cov, "the"
% CAPM, ...), regardless of whatever font context happens to surround the
% formula itself.
\let\SMOrigText\text
\renewcommand{\text}[1]{\SMOrigText{\usefont{OT1}{cmr}{m}{n}#1}}
